% Research contribution of agent /root/monodromy.  Input into a parent document.
% Status is local to the assertions below; this is not a claim to solve all of A.3--A.5.
\section{A.3--A.5: proved statements and remaining gates}
\label{sec:monodromy-contribution}
\status{Incomplete universal determination; the individually labelled results below are proved.  Lead A's full-linear assertion is refuted, its positivity assertion is proved with an explicit threshold, and its plane-curve point-count consequence is proved.}

\subsection{Conventions and the generic rank for projective hypersurfaces}

Let $S_d=\mathbf P(H^0(\mathbf P^r,\mathcal O(d)))$, and let
$\pi:\mathcal Y_d\to S_d$ be the universal degree-$d$ hypersurface.
For $r\geq2$, the sheaf $\mathcal H=R^{r-1}\pi_*\mathcal O_{\mathcal Y_d}$
is locally free of rank
\[
 h_d=\begin{cases}\binom{d-1}{r}&d\geq r+1,\\0&d\leq r.\end{cases}
\]
This follows from the universal divisor exact sequence and the cohomology
of line bundles on projective space; formation of $\mathcal H$ commutes
with all base changes.  Absolute Frobenius induces a map
$F_S^*\mathcal H\to\mathcal H$.  On its invertibility locus $W_d$, the
kernel of $F-1$ is a locally constant $\mathbf F_p$-sheaf of rank $h_d$.
On the smooth locus it is $R^{r-1}\pi_*\mathbf F_p$, by the
Artin--Schreier sequence.  For $X=\mathbf P^r$, this is the vanishing
local system in the question, since $H^{r-1}(\mathbf P^r,\mathcal O)=0$.

\begin{proposition}[All degrees, all characteristics]\label{prop:projective-generic-hw-rank}
\status{Proved under stated hypotheses}
The generic degree-$d$ hypersurface of $\mathbf P^r$ has bijective
Frobenius on $H^{r-1}(\mathcal O)$, and its generic stable rank is $h_d$.
In particular, for plane curves the generic rank is
$g_d=(d-1)(d-2)/2$, and the smallest positivity threshold is $d_0=3$.
\end{proposition}

\begin{proof}
Work over an algebraic closure and take $d$ hyperplanes in general
position.  Their union $D$ is a simple normal crossings divisor.
The alternating resolution of $\mathcal O_D$ by the structure sheaves
of intersections has only $H^0$ in its first hypercohomology page,
because every nonempty intersection is a projective space.  Frobenius
is bijective on each of these $H^0$'s and hence on the cohomology of
their complex.  It is therefore bijective on
$H^{r-1}(D,\mathcal O_D)$.  Thus $W_d$ is nonempty.  The smooth locus
is a nonempty dense open of the irreducible scheme $S_d$, so it meets
$W_d$.  Invertibility is open, which proves the assertion.  This proof
establishes only bijectivity on the specified coherent cohomology;
it does not conflate this condition with every formulation of
ordinarity for higher dimensional varieties.
\end{proof}

\par\smallskip\noindent\textbf{Comparison with Illusie.}
Illusie (1990), Theorem~0.1, p.~375, proves the stronger statement that
the general smooth complete intersection of every specified dimension
and multidegree is ordinary: its Newton and Hodge polygons coincide.
The theorem and the explanation of the earlier Koblitz argument on
pp.~375--376 were inspected in the publisher's two-page preview of
the reprint.  In particular, generic Hasse--Witt bijectivity follows;
the proposition above gives a separate complete proof of exactly the
part needed here.  The parameter space in Illusie's theorem varies all
defining equations.  It does not, by restriction of an open set alone,
prove generic ordinarity for hypersurface sections of every fixed $X$.

\subsection{The determinant character and the Hasse invariant}
\label{subsec:hasse-determinant}

Let $S$ be connected and let $(\mathcal E,F)$ be a vector bundle of
rank $h$ in characteristic $p$, with $F:F_S^*\mathcal E\simeq\mathcal E$.
Write $V=\ker(F-1)$ and $\mathcal L=\det\mathcal E$.  Taking
determinants gives an invertible section
\[
 \operatorname{Ha}_F\in
 H^0\bigl(S,\operatorname{Hom}(F_S^*\mathcal L,\mathcal L)\bigr)
 =H^0(S,\mathcal L^{\otimes(1-p)}).
\]
This is the Hasse invariant of the specified Frobenius module.  It
determines the determinant character exactly, including its arithmetic
part: $\det V$ is the rank-one Lang local system
$\ker(F_{\mathcal L}-1)$.  If $e$ trivializes $\mathcal L$ and
$F_{\mathcal L}(e)=h e$, a fixed nonzero vector $ze$ satisfies
\begin{equation}\label{eq:determinant-kummer}
 z^{p-1}=h^{-1}.
\end{equation}
Thus the $\mathbf F_p^\times$-torsor of frames of $\det V$ is the
Kummer torsor \eqref{eq:determinant-kummer}; a deck transformation
$z\mapsto az$ acts on the determinant fibre by $a$.  Changing the
trivialization of $\mathcal L$ changes this equation by a
$(p-1)$st power, as required.  Over $\mathbf F_q$, geometric Frobenius
acts on $\det V$ by the determinant of the $q$-power coherent
Frobenius; arithmetic Frobenius is its inverse.  This description
holds for the stable bundle of the relative vanishing system on its
constant-rank locus as well as for the full ordinary coherent module.
It does not assert that the image of the Kummer character is always
all of $\mathbf F_p^\times$: that depends on its Kummer class.

\subsection{The exact plane-curve monodromy in degrees at most four}
\label{subsec:small-plane-monodromy}

For $d=1,2$, the rank is zero and both monodromy groups are trivial.
For $d=3$, both images are $\mathrm{GL}_1(\mathbf F_p)$: for odd $p$
this follows from the Legendre subfamily and its simple Hasse roots,
proved below; for $p=2$ the assertion means the trivial group on a
one-dimensional space.  A nonzero rank-one representation is
irreducible.

For $d=4$, both images are $\mathrm{GL}_3(\mathbf F_p)$, and the
representation is irreducible.  Here is why the moduli-space theorem
actually applies.  Every nonhyperelliptic smooth genus-three curve is
canonically a smooth plane quartic, and conversely.  The stack quotient
of the parameter space of smooth plane quartics by $\mathrm{PGL}_3$
is the nonhyperelliptic open substack of $\mathcal M_3$.
The map from the parameter space to this quotient is a
$\mathrm{PGL}_3$-torsor in the stack sense, with geometrically connected
fibres, and is surjective on etale fundamental groups: a connected
finite etale cover of the quotient remains connected after pullback
along a torsor under a connected algebraic group.  Restriction to
the ordinary locus preserves this assertion.  That locus is a
nonempty open of the ordinary locus of $\mathcal M_3$, hence its
fundamental group has the same finite monodromy image.
Achter--Pries (2008), Proposition~4.8, says that this image on the
etale $p$-torsion of the Jacobian is $\mathrm{GL}_3(\mathbf F_p)$.
Dualizing gives the same full group on $V$.
The arithmetic image contains the geometric image and is already
a subgroup of the same finite general linear group, so it too is full.
No analogous moduli-space surjectivity assertion is used for degrees
$d\geq5$.  In characteristic two the separate contact-curve argument
and group induction instead determine every degree, in
Theorem~\ref{thm:plane-char2-all-groups}.

\subsection{A universal obstruction in characteristic two}
\label{subsec:odd-plane-obstruction}

Let $p=2$ and let $d\geq3$ be odd.  Write
\[
 f=\sum_{|a|=d}c_aX^{a_1}Y^{a_2}Z^{a_3},\qquad
 I_d=\{u\in\mathbf Z_{>0}^3:|u|=d\}.
\]
Put $c_a=0$ if a coordinate of $a$ is negative.  On a coefficient
chart of $S_d$, use the Cech basis of
$H^2(\mathbf P^2,\mathcal O(-d))$ indexed by $I_d$.
The matrix of the semilinear Frobenius on $H^1(C_f,\mathcal O)$ is
\begin{equation}\label{eq:plane-hw-char2}
 H_{u,v}=c_{2v-u}\quad(u,v\in I_d),\qquad F(x)=Hx^{[2]}.
\end{equation}
Indeed, the connecting map from
$0\to\mathcal O(-d)\xrightarrow{f}\mathcal O\to\mathcal O_{C_f}\to0$
identifies Frobenius with multiplication by $f$ after squaring the
Cech representative $X^{-v_1}Y^{-v_2}Z^{-v_3}$.

Define a column $w=(w_u)$ by
\[
 w_u=\begin{cases}c_u&u_1,u_2,u_3\text{ all odd},\\0&\text{otherwise}.
 \end{cases}
\]
Then the following is an identity over the polynomial ring in the
coefficients, not merely a numerical identity at selected curves:
\begin{equation}\label{eq:universal-covector}
 H^t w=w^{[2]}.
\end{equation}
For its $v$-coordinate, pair the term $c_{2v-u}c_u$ indexed by $u$
with the term indexed by $2v-u$.  A nonzero term has both indices
positive and all odd, so the involution preserves the indexing set.
All two-element orbits cancel in characteristic two.  The only possible
fixed index is $u=v$, contributing $c_v^2$ exactly when $v$ is all odd.

If $H$ is invertible, then $w\ne0$.  Otherwise all coefficients with
three odd exponents vanish, and every row with all-odd index $u$
in \eqref{eq:plane-hw-char2} vanishes.  Such an index exists because
$(1,1,d-2)\in I_d$.  This contradicts invertibility.

For a Lang fixed vector $x$, namely $Hx^{[2]}=x$, equation
\eqref{eq:universal-covector} gives
\[
 (w^tx)^2=(w^{[2]})^tx^{[2]}=w^tHx^{[2]}=w^tx.
\]
Thus $x\mapsto w^tx$ is an $\mathbf F_2$-valued linear functional on
the local system $V=\ker(F-1)$.  It is nonzero, since the Lang fixed
vectors form a basis after an etale cover.  It is a section of
$V^\vee$, so it is fixed by both geometric and arithmetic monodromy.
Working on a dense coefficient chart suffices: its intersection with
the ordinary smooth locus is dense, and the map on fundamental groups
from a nonempty open of a normal integral scheme is surjective.

\begin{theorem}[Refutation of full linear monodromy]\label{thm:odd-plane-char2-obstruction}
\status{Proved under stated hypotheses}
For every finite field of characteristic two and every odd $d\geq5$,
let $\bar\eta$ be a geometric generic point of the ordinary smooth
degree-$d$ plane-curve locus $U_d$, and put
$V=H^1_{\mathrm{et}}(C_{\bar\eta},\mathbf F_2)$.
Both images
\[
 G_{\mathrm{geom}}=\mathrm{im}\bigl(\pi_1(U_d\otimes\bar k,\bar\eta)
 \to\mathrm{GL}(V)\bigr),\qquad
 G_{\mathrm{arith}}=\mathrm{im}\bigl(\pi_1(U_d,\bar\eta)
 \to\mathrm{GL}(V)\bigr)
\]
fix a nonzero covector.  Consequently both are contained in the
proper subgroup
\[
 \mathrm{Stab}(\ell)\simeq
 \mathbf F_2^{\,g_d-1}\rtimes\mathrm{GL}_{g_d-1}(\mathbf F_2),
\]
and $V$ is reducible, with invariant hyperplane $\ker\ell$.
There is no integer $d_1(\mathbf P^2)$ for which the proposed full
linear geometric or arithmetic monodromy holds for every $d\geq d_1$
in characteristic two.
\end{theorem}

The theorem proves containment, not equality with the displayed
stabilizer.  Determining the exact image remains a separate gate.
The determinant character in this characteristic is automatically
trivial, because $\mathbf F_2^\times=\{1\}$.

\par\smallskip\noindent\textbf{Geometric interpretation and literature check.}
The obstruction agrees with the difference between the effective
geometric theta characteristic
$\mathcal O_C((d-3)/2)$ and the canonical characteristic-two theta
characteristic, which is non-effective for an ordinary curve.
Cais, Ellenberg, and Zureick-Brown (2013, Theorem~4.2) prove the resulting
rational two-torsion statement for odd degree plane curves.
Castryck, Streng, and Testa (2014, Corollary~13, Remark~14, and
Theorem~16) give the ordinary and Hasse--Witt formulations.
The explicit identity above avoids any descent ambiguity for theta
characteristics over imperfect coefficient fields.

\par\smallskip\noindent\textbf{Explicit checked fibre.}
Over $\mathbf F_2$, take
\begin{align*}
f={}&x^5+x^4y+x^3y^2+x^2y^3+x^2y^2z+xy^4+xyz^3+xz^4\\
 &+y^4z+y^3z^2+yz^4.
\end{align*}
In the ordered basis
$(1,1,3),(1,2,2),(1,3,1),(2,1,2),(2,2,1),(3,1,1)$, its matrix is
\[
H=\begin{pmatrix}
1&0&0&0&0&0\\1&0&1&0&1&1\\0&1&0&0&0&0\\
1&1&0&0&1&1\\0&0&1&0&1&0\\0&0&0&1&0&0
\end{pmatrix},\qquad w=(1,0,0,0,0,0)^t.
\]
Its rank is six.  The ideal $(f,f_x,f_y,f_z)$ has Groebner basis
$\{1\}$ on each of the three standard projective charts, proving
geometric smoothness.  Direct enumeration over
$\mathbf F_{2^m}$ for $m=1,\ldots,6$ gives respectively
$3,3,9,23,38,33$ points; the traces of $H^m$ modulo two are
$0,0,0,0,1,0$.  The program \texttt{code/plane\_char2.py} and its JSON
output record the input, field arithmetic, certificates, and timings.
The fibre is a check of the formulas; the universal identity, rather
than this one fibre, proves the generic monodromy obstruction.

\subsection{Plane curves: every point-count residue in every open}
\label{subsec:plane-residues}

\begin{theorem}\label{thm:plane-all-residues}
\status{Proved under stated hypotheses}
Let $k$ be a finite field of characteristic $p$.  For every $d\geq4$,
every nonempty Zariski open $U\subset S_d$, and every $a\in\mathbf F_p$,
there exist a finite extension $k'/k$ and a smooth plane curve
$C\in U(k')$ of degree $d$ with $\#C(k')\equiv a\pmod p$.
For $d=1,2$ the only residue for smooth curves is $1$.
For $d=3$, the residues on the ordinary smooth locus are exactly
$\mathbf F_p\setminus\{1\}$, and every nonempty open of that locus
attains every such residue after a finite extension.  Supersingular
smooth cubics give residue $1$, so all residues occur among all smooth
cubics; the every-open assertion for all residues fails in degree three.
\end{theorem}
\begin{proof}

\par\smallskip\noindent\textbf{Elliptic input.}
For every $\lambda\in\mathbf F_p^\times$, there is an ordinary elliptic
curve $E/\mathbf F_Q$, with $\mathbf F_Q/k$ finite, for which geometric
Frobenius on $H^1_{\mathrm{et}}(E_{\bar k},\mathbf F_p)$ is multiplication
by $\lambda$.  For $p=2$ take any ordinary elliptic curve.  For odd $p$,
write $m=(p-1)/2$.  The Legendre family
$y^2=x(x-1)(x-t)$ has Hasse polynomial
\[
 A(t)=(-1)^m\sum_{i=0}^{m}\binom mi^2t^i.
\]
Its roots are simple and avoid $0,1$: its values there are nonzero,
and it satisfies
$t(1-t)A''+(1-2t)A'-m^2A=0$ modulo $p$.
A repeated root away from $0,1$ would, by successive differentiation,
force all derivatives through degree $m<p$ to vanish, a contradiction.
At a simple root, the Kummer cover $z^{p-1}=A(t)^{-1}$ is totally
ramified of degree $p-1$.  Its geometric monodromy is therefore the
whole $\mathbf F_p^\times$.  Finite-field Chebotarev on the ordinary
Legendre locus gives every scalar $\lambda$.
Moreover we may replace $Q$ by $Q^e$ with arbitrarily large
$e\equiv1\pmod{p-1}$ without changing $\lambda$.

\par\smallskip\noindent\textbf{A nodal curve with prescribed residue.}
Put $r=d-3\geq1$ and embed $E$ as a plane cubic by $|3O|$.
Choose rational lines $L_1,\ldots,L_r$ in general position such that
all intersections are transverse, no three components meet, the
intersection $E\cap L_i$ is split for $i\geq2$, and
$E\cap L_1$ has either cycle type $3$ or cycle type $1+2$ under
Frobenius.  Write $t=0$ or $1$ for the number of rational points in
this last intersection.  Such lines exist after an extension of the
above kind, as follows.

Split lines arise from distinct $P,Q,R\in E(\mathbf F_Q)$ with
$P+Q+R=0$, giving order $Q^2$ choices.  Lines of type $1+2$ arise from
$P\in E(\mathbf F_{Q^2})\setminus E(\mathbf F_Q)$ and
$R=-(P+\pi P)\in E(\mathbf F_Q)$; the three points are distinct,
collinear, and Frobenius invariant.  There are order $Q^2$ choices.
For type $3$, use nonrational points in
$\ker(1+\pi+\pi^2)$.  This is the kernel of a separable isogeny,
is contained in $E(\mathbf F_{Q^3})$, and has cardinality
$\#E(\mathbf F_{Q^3})/\#E(\mathbf F_Q)$, of order $Q^2$.
At most nine of these points are rational, since rational points in
the kernel lie in $E[3]$.  Thus there are order $Q^2$ suitable lines
of each required type.  The finitely many incidence and tangency
conditions excluded at each choice remove only order $Q$ rational
lines.  This proves existence for sufficiently large $Q$.

The connected nodal curve
$C_0=E\cup L_1\cup\cdots\cup L_r$ has degree $d$.
Its normalization sequence
\[
0\longrightarrow\mathcal O_{C_0}\longrightarrow
\mathcal O_E\oplus\bigoplus_i\mathcal O_{L_i}
\longrightarrow\bigoplus_{P\text{ node}}k(P)\longrightarrow0
\]
shows that Frobenius on $H^1(C_0,\mathcal O_{C_0})$ is bijective:
it is bijective on the finite reduced node scheme, on constants of
the normalization, and on $H^1(E,\mathcal O_E)$.
Inclusion--exclusion gives
\begin{align*}
 \#C_0(\mathbf F_Q)
 &\equiv (1-\lambda)+r-\{t+3(r-1)\}-\binom r2\\
 &\equiv 4-2r-\binom r2-t-\lambda\pmod p.
\end{align*}
As $t\in\{0,1\}$ and $\lambda\in\mathbf F_p^\times$ vary, these
residues fill $\mathbf F_p$, including when $p=2$.

\par\smallskip\noindent\textbf{Passage to an arbitrary smooth open.}
The ordinary coherent-cohomology locus $W_d\subset S_d$ introduced
above is geometrically integral and contains $C_0$.
The finite etale local system $\ker(F-1)$ is defined over all of $W_d$,
including its singular fibres.  After base changing to $\mathbf F_Q$,
the Frobenius class of $C_0$ belongs to the finite arithmetic
monodromy group.  Chebotarev gives that same class at a closed point
of any nonempty open of $W_d$.  Apply this to
$U\cap W_d\cap S_d^{\mathrm{smooth}}$.
The coherent Frobenius trace formula
\[
 \#C(\mathbf F_{Q^e})\equiv
 1-\operatorname{Tr}\bigl(\operatorname{Frob}_{Q^e}^{\mathrm{geom}},
 H^1_{\mathrm{et}}(C_{\bar k},\mathbf F_p)\bigr)\pmod p
\]
also applies to the proper nodal curve $C_0$, so equality of the
classes gives the desired residue.  More precisely, the class at
$C_0$ maps to degree one in the finite constant-field quotient, and
Chebotarev supplies it at degrees congruent to one in that quotient;
the permitted residue fields are finite extensions of $k$.
This argument never assumes full general linear monodromy for
degree-$d$ plane curves.
\end{proof}

\par\smallskip\noindent\textbf{Primary-source verification of the two inputs.}
Katz--Sarnak (1999), Theorem~9.7.13, p.~290, is the finite-monodromy
equidistribution theorem, with the constant-field coset condition made
explicit in 9.7.1--9.7.9.  To apply its characteristic-zero formulation
to our finite $\mathbf F_p$-representation, take a faithful finite
dimensional characteristic-zero representation of its finite image
(for example the regular representation), choose $\ell\ne p$, and
regard it as a lisse $\overline{\mathbf Q}_\ell$-sheaf.  It is pure of
weight zero, and its finite geometric and arithmetic groups have the
required normality.  The theorem gives a positive limiting proportion
for every nonempty conjugacy class in the required coset.  Restriction
to a nonempty open preserves the finite monodromy image because the
base is normal.  Katz (1973), SGA~7~II, Expose~XXII, Corollary~3.2,
formulas (3.2.1)--(3.2.2), pp.~410--411, gives precisely the trace
congruence for every proper scheme over a finite field, including
our nodal curve.  Proposition~2.0 and formula (2.2.4) of the same
expose identify its coherent and etale formulations.

\subsection{An explicit positivity threshold for every smooth ambient variety}
\label{subsec:explicit-positivity}

Here $X\subset\mathbf P^N$ is smooth projective of pure dimension
$n\geq2$ over a perfect field of characteristic $p$, and $L=\mathcal O_X(1)$.
Geometric connectedness is not needed for this proposition.  All
dimensions in the following definitions may be computed after extending
the ground field to an algebraic closure.  Put
\[
 D=\deg_L X,\qquad
 a=\dim H^{n-1}(X,\mathcal O_X)_{\mathrm{ss}},
\]
and, for $1\leq m\leq n-2$, define the explicit cohomological bounds
\[
 B_m=\sum_{j=0}^{m}\binom mj
       h^{n-m-1+j}(X,L^{-j}),\qquad
 B=D+a+\sum_{m=1}^{n-2}B_m.
\]
Let $r_\omega$ be any Castelnuovo--Mumford regularity bound for
$i_*\omega_X$ on $\mathbf P^N$, and put
\begin{equation}\label{eq:explicit-d0}
 d_0=\max\left\{2n,1,r_\omega-1,
       \left\lfloor\frac{2^n n!B}{D}\right\rfloor+1\right\}.
\end{equation}
Every term is determined by the embedding and finitely many coherent
cohomology groups of $X$; this is a computable bound, not an unspecified
``sufficiently large'' constant.

\begin{theorem}\label{thm:all-X-explicit-positive-rank}
\status{Proved under stated hypotheses}
For every $d\geq d_0$, the generic smooth hypersurface section $Y$ of
degree $d$ has
$\dim E(Y,\mathcal O_Y)_{\mathrm{ss}}>0$.
More precisely, for every $d\geq\max\{1,r_\omega-1\}$, its generic
rank $f_{X,d}$ satisfies
\begin{equation}\label{eq:generic-rank-lower-bound}
 f_{X,d}\ \geq\
 D\binom dn-D\binom d{n-1}
 -\sum_{m=1}^{n-2}B_m\binom dm-a,
\end{equation}
where a binomial coefficient with lower index greater than $d$ is zero.
Consequently
\[
 f_{X,d}=\frac{D}{n!}d^n+O_X(d^{n-1})\quad(d\longrightarrow\infty).
\]
\end{theorem}

\begin{proof}
Take $d$ hyperplanes in simultaneous general position over the
algebraic closure and set $Y_0=\bigcup_{i=1}^d(X\cap H_i)$.
Every intersection $Z_I=X\cap\bigcap_{i\in I}H_i$ with $|I|\leq n$
is smooth of dimension $n-|I|$, and intersections of $n+1$ components
are empty.  Intersections with $|I|=n$ consist of $D$ reduced points.
The alternating normal-crossings resolution gives the Frobenius
equivariant spectral sequence
\[
 E_1^{s,t}=\bigoplus_{|I|=s+1}H^t(Z_I,\mathcal O_{Z_I})
 \ \Longrightarrow\ H^{s+t}(Y_0,\mathcal O_{Y_0}).
\]
The term $E_1^{n-1,0}$ has dimension $D\binom dn$ and bijective
Frobenius.  It has no outgoing differential on any page.  The image
of the incoming $d_1$ has dimension at most
$D\binom d{n-1}$: each smooth positive dimensional intersection has
at most $D$ connected components, since each has positive integral
degree.  For $r=2,\ldots,n-1$, the possible incoming $d_r$ has source
$E_r^{n-1-r,r-1}$.  Writing $m=n-r$, its dimension is at most
$\binom dm B_m$.  Indeed, the Koszul resolution of a transverse
$m$-fold hyperplane section bounds
\[
 h^{n-m-1}(Z_I,\mathcal O_{Z_I})
 \leq\sum_{j=0}^{m}\binom mj
 h^{n-m-1+j}(X,L^{-j})=B_m.
\]
Every quotient of the top term by an incoming image still has
bijective Frobenius: a Frobenius-stable subspace of a finite dimensional
bijective semilinear space over a perfect field itself has bijective
Frobenius.  Hence $E_\infty^{n-1,0}$ is a bijective-Frobenius subquotient
of $H^{n-1}(Y_0,\mathcal O_{Y_0})$ of dimension at least the first three
terms of the right side of \eqref{eq:generic-rank-lower-bound}.
Exactness of the stable-part functor gives the same lower bound for
$\dim H^{n-1}(Y_0,\mathcal O_{Y_0})_{\mathrm{ss}}$.

For $d\geq\max\{1,r_\omega-1\}$, regularity and Serre duality give
$H^{n-1}(X,L^{-d})=H^1(X,\omega_X\otimes L^d)^\vee=0$.
The divisor sequence, including for the normal-crossings divisor,
then makes $h^{n-1}(Y,\mathcal O_Y)$ independent of $Y$ in the space of
properly intersecting degree-$d$ sections: the terminal map
$H^n(X,L^{-d})\to H^n(X,\mathcal O_X)$ is surjective because
$H^n(Y,\mathcal O_Y)=0$.  The universal family is
flat, and cohomology and base change therefore gives a vector bundle
with Frobenius.  If its rank is $h$, the stable rank at a perfect
geometric point is the rank of $F^h$; thus its maximal value is attained
on an open, and the generic value is at least the value at $Y_0$.
The smooth locus is dense and meets this open.

For a smooth $Y$, R1 identifies the stable ambient image with
$H^{n-1}(X,\mathcal O_X)_{\mathrm{ss}}$, of dimension $a$; in fact the
coherent restriction is injective in the present range of $d$.
Exactness of the stable-part functor subtracts $a$ and proves
\eqref{eq:generic-rank-lower-bound}.
For $d\geq2n$ one has
$\binom dn\geq(d/2)^n/n!$ and $\binom dm\leq d^{n-1}$ for $m<n$.
The right side of \eqref{eq:generic-rank-lower-bound} is therefore at
least $d^{n-1}(Dd/(2^nn!)-B)$, which is positive under
\eqref{eq:explicit-d0}.  Finally, Serre duality, Riemann--Roch, and
Serre vanishing give
$\dim E(Y,\mathcal O_Y)=Dd^n/n!+O_X(d^{n-1})$; together with the lower
bound this proves the asymptotic formula.
\end{proof}

\begin{corollary}[The eventual parenthesis of 8.5.8]
\label{cor:eventual-hodge-parenthesis}
\status{Proved under stated hypotheses}
For the explicit $d_0$ of \eqref{eq:explicit-d0}, every $d\geq d_0$
has the following properties.  Every smooth degree-$d$ section $Y$
satisfies
\[
 \dim E(Y,\mathcal O_Y)
 =h^0(X,\omega_X\otimes L^d)-h^0(X,\omega_X)>0.
\]
For every $Y$ in the maximal-stable-rank open $U_d$, the number of
$p$-adic unit roots in vanishing cohomology is positive.  Thus, on
this open and in this stated eventual range, unit-root nonvanishing
and the maximal Hodge-level condition $E(Y,\mathcal O_Y)\ne0$ are
equivalent.  No assumption on crystalline $p$-torsion of $X$ is needed.
\end{corollary}

\begin{proof}
The coherent divisor sequence and the vanishings used in the preceding
proof identify $E(Y,\mathcal O_Y)$ with
$\ker(H^n(X,L^{-d})\to H^n(X,\mathcal O_X))$, whose dimension is the
displayed, $Y$-independent difference by Serre duality.  It is positive
because it is at least the positive generic stable rank of
Theorem~\ref{thm:all-X-explicit-positive-rank}.
For $Y\in U_d$, Theorems~\ref{thm:A-exact85} and
\ref{thm:A-unit-count} identify the unit-root count with
\[
 f_{X,d}+\dim H^n_{\mathrm{cris}}(X/W)[p]_{\mathrm{ss}},
\]
which is positive.  Both properties therefore hold for every point
under these quantifiers.
\end{proof}

\par\smallskip\noindent\textbf{Scope of the manuscript's parenthesis.}
The supplied transcription of cote~4, scan~8, first specifies a
section for which the semisimple rank has its maximum value.  Scan~9,
continuing 8.5.8, then says ``pour tous ces $Y$'' and, equivalently,
``pour $Y$ g\'en\'erique''.  Corollary~\ref{cor:eventual-hodge-parenthesis}
answers this eventual generic assertion, including varieties with
$p$-torsion.  It does not assert that an arbitrary smooth section
outside $U_d$ has a unit root whenever its coherent vanishing space
is nonzero.  For example, a supersingular smooth plane cubic has
one-dimensional $H^1(\mathcal O)$ and no unit root.  That example
distinguishes the stronger pointwise assertion from the quantifiers
actually written in 8.5.8; it is not substituted for the requested
eventual assertion.

\par\smallskip\noindent\textbf{Relation to Koblitz.}
Koblitz (1975), I.\S5, Theorems~1--2, pp.~139--141, proves the
asymptotic invertibility estimate with an error of order $d^{n-1}$.
The argument above gives explicit constants in terms of the embedding
and the displayed cohomology groups.  Koblitz's I.\S6, Theorem~3,
p.~142, concerns generic complete intersections, not every hypersurface
section of an arbitrary fixed ambient variety.  His revised conjecture
in II.\S3, pp.~149--151, must not be cited as an established exact rank
formula.

\subsection{All residues for hypersurfaces in every projective space}

\begin{theorem}\label{thm:projective-all-residues}
\status{Proved under stated hypotheses}
Let $k$ be finite of characteristic $p$, let $n\geq2$, and let
$d\geq pn+1$.  Every nonempty open in the degree-$d$ hypersurface
parameter space of $\mathbf P^n_k$ contains, over a finite extension,
a smooth hypersurface with any prescribed point-count residue
modulo $p$.
\end{theorem}

\begin{proof}
Choose the integer representative $t\in\{0,\ldots,p-1\}$ of the
desired residue.  Over a sufficiently large finite extension
$\mathbf F_Q/k$, choose a general-position collection of $d$
geometric hyperplanes on which Frobenius has $t$ orbits of length $n$
and one orbit of length $d-nt>n$.  Existence follows by taking a
rational point in the open set of general-position configurations
in the product of Weil restrictions of the dual projective space
from the indicated extension degrees.  This parameter space is
geometrically integral, since over the algebraic closure it is a
product of projective spaces; the general-position open is nonempty.
Lang--Weil supplies rational points over arbitrarily large fields.
One may first form these Weil restrictions over a fixed finite base
field and then take extension degrees congruent to one modulo the
least common multiple of the prescribed orbit lengths; this preserves
their field factors while making the field arbitrarily large.
Distinctness of every geometric hyperplane forces the specified
orbit lengths to be exact.

The union $Y_0$ of these hyperplanes is defined over $\mathbf F_Q$.
A rational point lying on one hyperplane lies on its entire Frobenius
orbit.  An orbit of length $n$ has exactly one common intersection
point in $\mathbf P^n$, and that point is rational.  An orbit of
length greater than $n$ has empty intersection.  The $t$ points from
the shorter orbits are distinct, because any $n+1$ distinct
hyperplanes have empty intersection.  Therefore
\[
 \#Y_0(\mathbf F_Q)=t
\]
as an equality of integers.
The general-position hyperplane-union argument proves that
$H^{n-1}(Y_0,\mathcal O_{Y_0})$ has bijective Frobenius, of the full
generic rank $\binom{d-1}{n}$.  Thus $Y_0$ belongs to the same open
$W_d$ on which the entire Lang local system is defined.
The coherent trace congruence, together with the standard vanishing
of intermediate coherent cohomology for a projective hypersurface,
reads
\[
 \#Y(\mathbf F_{Q^e})\equiv
 1+(-1)^{n-1}\operatorname{Tr}(\operatorname{Frob}_{Q^e}^{\mathrm{geom}},V_Y)
 \pmod p.
\]
Apply finite Chebotarev on $W_d$ to the Frobenius class of $Y_0$ and
the given smooth open.  The resulting smooth hypersurface has the
same residue $t$.  No exact monodromy-group computation is required.
\end{proof}

\subsection{Relation to the subsequent classifications}
The odd-degree characteristic-two obstruction refutes the universal
full-$\mathrm{GL}$ expectation, even for geometrically connected
$X=\mathbf P^2$.  The all-degree characteristic-two theorem identifies
the exact geometric and arithmetic images: full $\mathrm{GL}_g(2)$
for even $d\ge4$, and the full stabilizer
$2^{g-1}\!:\mathrm{GL}_{g-1}(2)$ for odd $d\ge5$.
The subsequent projective-space theorems settle every remaining
plane degree and every dimension, in all characteristics.
For arbitrary $X$, the explicit asymptotic positivity bound above is
proved using a normal-crossings union of hyperplane sections.
The later all-ambient theorem also proves every point-count residue
with an explicit degree bound. An effective finite generic-rank
calculation and exact determinant image are provided below.
A structural generic-rank formula and the complete nonabelian
monodromy group for arbitrary $X$ remain unresolved here.

% /root/monodromy, 19 November 2026.
% Complete mathematical construction; independently audited by /root.
\subsection{A smooth ambient surface whose relative Hasse--Witt map is always singular}
\label{subsec:relative-ordinarity-counterexample}
\status{Proved; independently audited by /root.  The finite Hopf-algebra checks
were executed; the final invariant-ring/Bertini search was not executed.}

The following construction refutes eventual bijectivity on
$E(Y,\mathcal O_Y)$ for arbitrary smooth ambient varieties.  The
ambient variety is smooth.  In particular, this does not reuse the
singular cones in Koblitz's counterexamples.

\begin{lemma}\label{lem:picard-finite-quotient}
\status{Proved under stated hypotheses}
Let $G$ be a finite commutative group scheme over a field $k$, and
let $Z\to X$ be a $G$-torsor, where $Z$ is a positive dimensional
projective complete intersection of dimension at least two and $X$
is smooth and geometrically connected.  If $G^D$ is connected, then
\[
 \operatorname{Pic}_{X/k}^0=G^D.
\]
\end{lemma}

\begin{proof}
The Koszul resolution gives $H^0(Z,\mathcal O_Z)=k$ and
$H^1(Z,\mathcal O_Z)=0$, compatibly with every extension of scalars;
smoothness or reducedness of $Z$ is unnecessary.  Therefore the
identity component of $\operatorname{Pic}_{Z/k}$ is trivial: its
Zariski tangent space is zero.  The kernel, as an fppf sheaf, of
\[
 \operatorname{Pic}_{X/k}\longrightarrow\operatorname{Pic}_{Z/k}
\]
is $G^D$.  Indeed, after trivializing the pullback of a line bundle,
its descent datum is a $G$-linearization of $\mathcal O_Z$.
All functions on $G\times Z$ come from $G$, because
$H^0(Z_R,\mathcal O)=R$ for every $k$-algebra $R$.  The descent
cocycle condition is precisely the condition that this unit be a
character $G_R\to\mathbf G_{m,R}$.  This proves the kernel assertion
without a surjectivity claim about the entire Picard scheme.
The connected group $\operatorname{Pic}_{X/k}^0$ maps to the trivial
identity component of $\operatorname{Pic}_{Z/k}$ and hence lies in
$G^D$.  Conversely the connected kernel $G^D$ lies in
$\operatorname{Pic}_{X/k}^0$.
\end{proof}

\par\smallskip\noindent\textbf{A specified finite group scheme and linear action.}
Work first over $\mathbf F_2$ and put
\[
 G=\operatorname{Spec}\mathbf F_2[a,b]/(a^2,b^2),\qquad
 \Delta(a)=a\otimes1+1\otimes a,\quad
 \Delta(b)=b\otimes1+1\otimes b+a\otimes a.
\]
Its group law is
$(a,b)(c,e)=(a+c,b+e+ac)$.  The basis $1,a,b,ab$ is the basis dual
to $1,t,t^2,t^3$ in the Hopf algebra of $\alpha_4$; the multiplication
rule is the divided-power rule
$f_if_j=\binom{i+j}{i}f_{i+j}$, truncated at index four.  Thus
$G=(\alpha_4)^D$.  Its Frobenius is zero, so
$V_{\alpha_4}=D(F_G)=0$, by Cartier duality for Verschiebung
\cite[Proposition~5.19]{EdixhovenVanDerGeerMoonenAV}.

Let $J$ be the direct sum of two nilpotent Jordan blocks of size
four on $\mathbf F_2^8$, with $J e_i=e_{i+1}$ in each block.
The formula
\begin{equation}\label{eq:alpha4-regular-action}
 (a,b)\longmapsto I+aJ+bJ^2+abJ^3
\end{equation}
is a faithful representation on $\mathbf A^8$ and acts on
$\mathbf P^7$.  Multiplication of these matrices gives exactly the
displayed group law.
The only subgroup of $G$ of order two is
$H=\{a=0\}\simeq\alpha_2$, and every nontrivial subgroup contains it.
One may see this directly from the restricted Lie algebra, with
basis $J,J^2$ and square operation $J^{[2]}=J^2$, $(J^2)^{[2]}=0$:
its only one-dimensional restricted subalgebra is $\langle J^2\rangle$.
Since $G$ has height one, its subgroup schemes correspond to these
restricted Lie subalgebras.  A projective line $\langle v\rangle$
is fixed by $H$ precisely when $J^2v=0$; a nilpotent operator has
no nonzero eigenvalue.  Consequently the nonfree locus is exactly
\[
 Z_{\rm bad}=\mathbf P(\ker J^2)=\mathbf P^3\subset\mathbf P^7.
\]
In particular, it has codimension four.

\par\smallskip\noindent\textbf{The surface and its fixed embedding.}
Let $Q=\mathbf P^7/G$ be the projective quotient and let
$Q^\circ=(\mathbf P^7\setminus Z_{\rm bad})/G$.  The quotient exists,
the map $\mathbf P^7\to Q$ is finite, and $Q^\circ$ is smooth;
these are Kothari's Theorem~3.3, p.~6
\cite{Kothari2026}.  The image of $Z_{\rm bad}$ has dimension three.
Fix a projective embedding $Q\hookrightarrow\mathbf P^M$ supplied
by a sufficiently divisible invariant Veronese ring.  Five general
hyperplanes of $\mathbf P^M$ cut out a smooth surface
\[
 X=Q\cap H_1\cap\cdots\cap H_5\subset Q^\circ.
\]
Their intersection avoids the three-dimensional bad locus and is
smooth by Bertini on $Q^\circ$.  The five pulled-back equations cut
out a surface $Z\subset\mathbf P^7$ which is a complete intersection;
the equality of the number of equations and the codimension makes
them a regular sequence in the polynomial ring.  The action on $Z$
is free and $Z/G=X$.  Geometric connectedness of $Z$, and then of
$X$, follows from the Koszul computation $H^0(Z,\mathcal O_Z)=k$.
This is also the particular dimension-two case of Kothari's
Proposition~3.4, pp.~6--7, and Corollary~3.7, p.~9; those statements
explicitly require smoothness of the quotient, not of $Z$.

This construction can be made unambiguous over a specified finite
field without leaving an unspecified ``general'' choice.  Use the
given ordered coordinates of $\mathbf P^7$ and their graded
lexicographic monomial order.  Compute the homogeneous invariant
ring by solving the coaction equations.  Every coordinate square is
invariant, because $a^2=b^2=0$, so this invariant ring is a finite
module over $\mathbf F_2[x_0^2,\ldots,x_7^2]$.  Its computation is
therefore a finite module-kernel calculation over an explicit
polynomial ring, rather than an unbounded search through degrees.
Take the least positive
Veronese degree which is generated in degree one.  Use its reduced
row-echelon monomial basis to embed $Q$.  Enumerate finite extensions
of $\mathbf F_2$ in increasing degree, each presented by the first
monic irreducible polynomial in lexicographic order, and then
enumerate ordered five-tuples of hyperplanes.  Retain the first
tuple for which the intersection is two dimensional, smooth, and
disjoint from the bad locus.  These conditions are decidable by
elimination and the Jacobian criterion.  The preceding nonempty
Bertini open proves termination.  Fix this output field $k$, this
surface $X$, and the resulting very ample bundle
$L=\mathcal O_{\mathbf P^M}(1)|_X$.
The description therefore defines one embedded surface by a finite
algorithm; the proof applies to every tuple passing these tests.
We do not claim to have numerically executed the invariant-ring and
Bertini search to identify the first tuple.

Lemma~\ref{lem:picard-finite-quotient} gives
\begin{equation}\label{eq:alpha4-picard-surface}
 \operatorname{Pic}_{X/k}^0=\alpha_4,\qquad
 h^1(X,\mathcal O_X)=1.
\end{equation}
Frobenius on this one-dimensional coherent cohomology is zero.
For a direct verification, the Cartan--Leray spectral sequence for
$Z\to X$ and $H^1(Z,\mathcal O_Z)=0$ identifies
$H^1(X,\mathcal O_X)$ with $H^1(G,k)$, the primitive elements of
$\mathcal O(G)$.  Those primitives are exactly $k a$; $b$ has the
additional term $a\otimes a$ in its coproduct, and $ab$ has the
terms $a\otimes b+b\otimes a$.  Absolute Frobenius sends $a$ to zero.

\begin{theorem}\label{thm:relative-ordinarity-fails-smooth}
\status{Proved under stated hypotheses}
For the fixed smooth projective geometrically connected surface
$(X,L)/k$ just constructed, let $r_\omega$ be a regularity bound for
its embedded dualizing sheaf, and put
$d_* =\max\{1,r_\omega-1\}$.
For every $d\geq d_*$ and every smooth section $Y\in|L^d|$, of
genus $g_Y$, one has
\[
 \dim E(Y,\mathcal O_Y)=g_Y-1,
 \qquad
 \dim E(Y,\mathcal O_Y)_{\mathrm{ss}}
       =f_Y\leq g_Y-2,
\]
where $f_Y$ is the $2$-rank of its Jacobian.  In particular the
Frobenius on $E(Y,\mathcal O_Y)$ is never bijective in this entire
range, including at the generic point.
\end{theorem}

\begin{proof}
Regularity and Serre duality give
$H^1(X,L^{-d})=H^1(X,\omega_X\otimes L^d)^\vee=0$.
Thus $H^1(X,\mathcal O_X)\to H^1(Y,\mathcal O_Y)$ is injective.
The induced homomorphism of group schemes
$\alpha_4=\operatorname{Pic}_{X/k}^0\to J_Y$ has zero Lie kernel.
Its finite kernel is therefore etale, but a subgroup of the local
group scheme $\alpha_4$ is local.  The kernel is trivial, and this
finite monomorphism is a closed immersion.

Naturality of Verschiebung and $V_{\alpha_4}=0$ put
$\alpha_4^{(2)}$ inside $\ker(V_{J_Y}:J_Y^{(2)}\to J_Y)$.
For an abelian variety of dimension $g_Y$ and $2$-rank $f_Y$, this
kernel has order $2^{g_Y}$ and its etale quotient has order $2^{f_Y}$;
its connected component therefore has order $2^{g_Y-f_Y}$.
Here the degree and point-count assertions are respectively
Propositions~5.20 and~5.22 of
\cite{EdixhovenVanDerGeerMoonenAV}: Frobenius is bijective on geometric
points, so $V\circ F=[2]$ identifies the geometric points of $\ker V$
with $J_Y[2](\bar k)$.
The connected subgroup $\alpha_4^{(2)}$ has order four, so
$g_Y-f_Y\geq2$.
Equation~\eqref{eq:alpha4-picard-surface} and coherent injectivity
give the dimension $g_Y-1$.  Since the ambient Frobenius is zero,
exactness of stable parts gives
$\dim E(Y,\mathcal O_Y)_{\mathrm{ss}}=f_Y$.
This proves the strict inequality and failure of bijectivity.
\end{proof}

\par\smallskip\noindent\textbf{What is and is not determined.}
This is a smooth counterexample to the universal eventual-relative-
ordinarity assertion.  It is compatible with the explicit positive
generic-rank theorem: for still larger degrees its generic rank is
positive and at least one less than its coherent dimension.  The
argument gives an exact structural obstruction and the displayed
upper bound; it does not assert an exact generic $2$-rank for this
surface.  The exact generic value for arbitrary $X$ remains a
separate determination.

\par\smallskip\noindent\textbf{Retained finite verification.}
The program \texttt{code/odd\_plane\_and\_alpha4\_checks.py} verifies
the Hopf group law, associativity, the matrix representation, its
divided-power multiplication table, the one-dimensional primitive
space and its zero Frobenius, and $\operatorname{rank}(J^2)=4$ on
the eight-dimensional representation.  The output is retained in
\texttt{outputs/odd\_plane\_and\_alpha4\_checks.json}.  These are
checks of the finite algebraic construction; existence and
smoothness of the surface are proved by the specified quotient and
Bertini construction, not inferred from a sample point calculation.

% Exact quintic case; independently audited by parent /root.
\subsection{The exact group for ordinary plane quintics in characteristic two}
\label{subsec:quintic-exact}

\begin{theorem}\label{thm:quintic-char2-monodromy}
\status{Proved under stated hypotheses}
For the universal ordinary smooth plane quintic in characteristic two,
the geometric and arithmetic monodromy groups on
$V=H^1_{\mathrm{et}}(C_{\bar\eta},\mathbf F_2)$ are both the full
stabilizer of the canonical nonzero covector:
\[
 G_{\mathrm{geom}}=G_{\mathrm{arith}}
 =\operatorname{Stab}_{\mathrm{GL}_6(\mathbf F_2)}(\ell)
 \simeq\mathbf F_2^5\rtimes\mathrm{GL}_5(\mathbf F_2).
\]
In particular their order is $319979520$, the representation is
reducible with invariant hyperplane of dimension five, and the
determinant character is trivial.
\end{theorem}
\begin{proof}

\par\smallskip\noindent\textbf{Theta-characteristic input, with a proof.}
For a smooth curve over an algebraically closed field of characteristic
two, let $B=F_*\mathcal O_C/\mathcal O_C$ for absolute Frobenius.
This is a line bundle with $B^{\otimes2}\simeq\omega_C$.
Locally a separating parameter $t$ gives the basis $1,t$ of
$F_*\mathcal O_C$ over $\mathcal O_C$; a change of separating parameter
$u=a^2+tb^2$ sends the generator of $B$ to $b$ times that generator,
whereas $du=b^2dt$.  This proves the claimed square relation.
The long exact sequence of
$0\to\mathcal O_C\to F_*\mathcal O_C\to B\to0$ shows
\[
 H^0(C,B)=\ker\bigl(F:H^1(C,\mathcal O_C)\to H^1(C,\mathcal O_C)\bigr).
\]
Thus $B$ is non-effective for an ordinary curve.
Every other theta characteristic is $B\otimes T$ for a nontrivial
two-torsion line bundle $T$.  Tensoring this sequence by $T$ and using
the projection formula and $F^*T=T^{\otimes2}\simeq\mathcal O_C$
gives an injection
\[
 H^0(C,\mathcal O_C)=H^0(C,F_*F^*T)
 \lhook\joinrel\longrightarrow H^0(C,B\otimes T),
\]
because $H^0(C,T)=0$.  Hence every theta characteristic other than $B$
is effective.  This is the statement recorded as St\"ohr--Voloch's
theorem in Ishitsuka--Ito, Theorem~3.1; Mumford (1971), p.~191,
describes the canonical characteristic-two theta characteristic.
For the geometric generic curve, the same reasoning may be performed
after perfecting its coefficient field.  Perfection does not change
the absolute Galois group, and $B$ is functorial under that group.

An ordinary curve of genus six has exactly $2^6=64$ geometric theta
characteristics.  On a plane quintic, $L=\mathcal O_C(1)$ is one of
them and $h^0(L)=3$, while $B$ is the unique non-effective one.
Under the identification of their set with the geometric points of
$J[2]$ by $T\mapsto L\otimes T$, the two distinguished points are
$0$ and $v=B\otimes L^{-1}\ne0$.  This is the dual representation
$V^\vee$, and $v$ is fixed by monodromy.

\par\smallskip\noindent\textbf{An irreducible incidence correspondence of contact conics.}
Let $S=\mathbf P(H^0(\mathbf P^2,\mathcal O(5)))$, of dimension twenty.
Let $I$ parametrize pairs $(C,Q)$ where $Q$ is a smooth plane conic
and the divisor cut out by $C$ on $Q$ is twice a divisor of degree five.
For a fixed conic $Q\simeq\mathbf P^1$, this means that the restriction
of its quintic equation, a binary form of degree ten, is a square.
The square forms constitute a vector subspace of dimension six inside
$H^0(Q,\mathcal O_Q(5))$, of dimension eleven.  The restriction map
from ternary quintics is surjective, with kernel
$Q\cdot H^0(\mathbf P^2,\mathcal O(3))$ of dimension ten.
Consequently the allowed quintic equations form a vector space of
dimension sixteen, or a projective space of dimension fifteen.
The smooth-conic parameter space has dimension five, so $\dim I=20$.
It is irreducible: fix $Q_0=(XZ+Y^2=0)$, take the irreducible product
$\mathrm{PGL}_3\times\mathbf P^{15}$ for its allowed quintics, and
map it surjectively to $I$ by changing projective coordinates.

The locus of pairs with a reduced reducible conic has dimension at
most nineteen.  Indeed, write $Q=L_1\cup L_2$ and $P=L_1\cap L_2$.
If $C$ is smooth and $C\cap Q$ is even, then $P\in C$: away from $P$
each of the two degree-five line sections must have even multiplicity,
so its multiplicity at $P$ is odd.  On each line the restriction is
therefore its linear equation for $P$ times the square of a binary
quadratic.  Each such restriction space has dimension three, and the
two restrictions automatically agree with value zero at $P$.
The combined restriction space has dimension six, so the quintic
equations again form a sixteen-dimensional vector space.
Pairs of distinct lines have dimension four, proving the bound
$4+15=19$.  Double-line conics cut exactly the divisors corresponding
to $L=\mathcal O_C(1)$.

Every effective theta characteristic $\theta$ and nonzero section $s$
gives a contact conic: $s^2\in H^0(C,\omega_C)=H^0(C,\mathcal O_C(2))$
lifts uniquely to a quadratic form on $\mathbf P^2$.
If that conic is a double line, then $\theta=L$.
Outside the dimension-at-most-nineteen image just described, each of
the other sixty-two theta characteristics therefore supplies a smooth
contact conic.  Thus $I\to S$ is dominant.  Since both dimensions are
twenty, its geometric generic fibre is zero dimensional.
It follows also that every such theta characteristic has $h^0=1$:
otherwise its projective space of sections would give a positive
dimensional family of conics in the same fibre.  Distinct theta
characteristics give distinct conics, because the half of an even
divisor on a smooth curve is unique.  The geometric generic fibre
therefore has exactly sixty-two points.

Irreducibility of $I$ implies transitivity of the geometric Galois
group on these sixty-two points.  This remains valid if the incidence
map has a purely inseparable part: the number of geometric generic
points is its separable degree, and a field's embeddings into a
separable closure form one orbit.  The Galois-equivariant
theta-characteristic interpretation proves that
$G_{\mathrm{geom}}$ on $V^\vee\setminus\{0,v\}$ is transitive.

\par\smallskip\noindent\textbf{The factor seven supplied by a quartic boundary.}
Take an ordinary smooth plane quartic $C_4$ and a line meeting it
transversely in four distinct points.  Their union $C_0$ is an ordinary
nodal quintic: the normalization sequence gives bijective Frobenius
on its $H^1(\mathcal O)$, of dimension six.  Over the parameter space
of such pairs, the etale normalization sequence gives a surjection
of local systems
\[
 H^1_{\mathrm{et}}(C_0,\mathbf F_2)
 \longrightarrow H^1_{\mathrm{et}}(C_4,\mathbf F_2).
\]
Forgetting the line is a smooth surjection with geometrically
connected fibres (nonempty opens of the dual plane), hence is
surjective on fundamental groups.  The quartic local system has
geometric monodromy $\mathrm{GL}_3(\mathbf F_2)$ by the already proved
quartic case.  Its order is divisible by seven, so the monodromy on
the nodal quintic family, and consequently the generic quintic group,
has order divisible by seven.  This specialization is legitimate
because the entire rank-six Lang system extends across $C_0$ on the
invertible-Hasse--Witt locus $W_5$.  The dense smooth open and $W_5$
have the same monodromy image for that extended finite etale system.

\par\smallskip\noindent\textbf{Finite group conclusion.}
Write $M=\operatorname{Stab}(v)\simeq2^5\!:\!\mathrm{GL}_5(2)$
in its action on $V^\vee$, and let $G\subset M$ be the geometric image.
Its transitive orbit of size sixty-two implies $31\mid|G|$; the
boundary argument gives $7\mid|G|$.
The complete maximal-subgroup list for $M$ is given in
Basheer--Moori (2021), Table~1, p.~3:
\[
\begin{array}{c|r|r}
\text{maximal subgroup}&\text{order}&\text{index in }M\\\hline
2_+^{1+8}:\mathrm{GL}_4(2)&10321920&31\\
2^{4+5}:\mathrm{GL}_4(2)&10321920&31\\
\mathrm{GL}_5(2)&9999360&32\\
2^{2+9}:(\mathrm{GL}_3(2)\times S_3)&2064384&155\\
2^9:(\mathrm{GL}_3(2)\times S_4)&2064384&155\\
(2^5:31):5&4960&64512.
\end{array}
\]
The first, second, fourth, and fifth orders are not divisible by
thirty-one.  The last is not divisible by seven.  The remaining
maximal subgroup is a complement $\mathrm{GL}_5(2)$; its action on
$V^\vee\setminus\{0,v\}$ has two orbits of length thirty-one, so it
cannot contain a group transitive on all sixty-two points.
No proper maximal subgroup can contain $G$, and therefore $G=M$.
The arithmetic image is contained in $M$ and contains $G$, hence
equals it.
\end{proof}

\par\smallskip\noindent\textbf{Independent elementary check of the complement issue.}
The same conclusion can be obtained from the maximal-subgroup table
for $\mathrm{GL}_5(2)$ in the ATLAS and the elementary fact
$H^1(\mathrm{GL}_m(2),\mathbf F_2^m)=0$ for $m\geq4$.
For completeness, let $c$ be a cocycle and put
$t_{ij}=1+E_{ij}$.  The relation $t_{ij}^2=1$ says that the $j$th
coordinate of $c(t_{ij})$ is zero.  Commuting $t_{ij}$ with
$t_{k\ell}$ for an auxiliary index $k$ distinct from $i,j,\ell$
shows that every coordinate except the $i$th is zero.
Thus $c(t_{ij})=a_{ij}e_i$.  The relation
$t_{ij}t_{jk}t_{ij}=t_{jk}t_{ik}$, for distinct indices,
gives $a_{ik}=a_{jk}$, so $a_{ij}=b_j$ depends only on $j$.
The coboundary associated with $b=(b_j)$ agrees with $c$ on all
elementary transvections, which generate the group.  Hence $c$ is
a coboundary.  The same result holds for the dual module by the
inverse-transpose automorphism.  Complements in $2^5:\mathrm{GL}_5(2)$
are therefore conjugate to the standard one, justifying the asserted
two-orbit description independently of the affine maximal-subgroup
classification.

% Exact characteristic-two plane sextic case; proof for independent audit.
\subsection{The exact group for ordinary plane sextics in characteristic two}
\label{subsec:sextic-exact}

\begin{theorem}\label{thm:sextic-full-monodromy}
\status{Proved under stated hypotheses}
For the universal ordinary smooth plane sextic in characteristic two,
the geometric and arithmetic monodromy images on
$V=H^1_{\mathrm{et}}(C_{\bar\eta},\mathbf F_2)$ are both
$\mathrm{GL}_{10}(\mathbf F_2)$.  Thus $V$ is irreducible and its
determinant character is trivial.
\end{theorem}

\begin{proof}
The generic genus is ten, and ordinarity gives $2^{10}=1024$
geometric theta characteristics.  The canonical characteristic $B_C$
is the unique non-effective one.  It is invariant under the
geometric Galois group, also when a purely inseparable extension is
needed to define it as a line bundle.  Identifying the set of theta
characteristics with $V^\vee$ using $B_C$ as origin, the remaining
1023 effective theta characteristics correspond to the nonzero
vectors.  We first prove that monodromy is transitive on them.

\par\smallskip\noindent\textbf{The contact reciprocity calculation.}
Let $C$ be a smooth plane curve of degree $d$, let $Q$ be a smooth
plane curve of degree $e=d-3$, and suppose $C\cap Q=2D$.
The half-intersection divisor defines the same zero-dimensional
subscheme $D$ on both curves: locally, in coordinates in which one
curve is $y=0$, the other has contact $y=u x^{2m}+\cdots$, and both
half divisors have ideal $(y,x^m)$.
Put
\[
 \theta_C=\mathcal O_C(D),\qquad
 \theta_Q=\mathcal O_Q(D-3H).
\]
These are theta characteristics.  The two ideal-sheaf sequences,
twisted by $e$, are
\[
 0\to\mathcal O_{\mathbf P^2}(-3)\to\mathcal I_D(e)
       \to\theta_C\to0,
 \qquad
 0\to\mathcal O_{\mathbf P^2}\to\mathcal I_D(e)
       \to\theta_Q\to0.
\]
Since $H^1(\mathbf P^2,\mathcal O(-3))=H^1(\mathbf P^2,\mathcal O)=0$,
they give the exact numerical relation
\begin{equation}\label{eq:contact-reciprocity}
 h^0(\theta_C)=1+h^0(\theta_Q).
\end{equation}

\par\smallskip\noindent\textbf{Smooth contact cubics.}
For $d=6$, every effective $\theta_C$ with section $s$ supplies a
contact cubic, because $s^2\in H^0(C,\omega_C)=H^0(C,\mathcal O_C(3))$
lifts uniquely to a ternary cubic.
If $Q$ is smooth, its theta characteristics are degree-zero line
bundles.  For an ordinary cubic, the only non-effective one is $B_Q$;
the other is $\mathcal O_Q$.  For a supersingular cubic, the sole
geometric theta characteristic is $\mathcal O_Q$.

The incidence corresponding to $\theta_Q=B_Q$ over the ordinary
smooth-cubic parameter space is irreducible of dimension 27.
To see this without a descent assumption on $B_Q$, pull the
coefficient space back along Frobenius (and work in an affine
coefficient chart).  The quotient of the relative Frobenius pushforward
by the structure sheaf supplies $B_Q$ on the Frobenius-twisted universal
curve.  This purely inseparable base change is irreducible and changes
no geometric Galois group.  The spaces
$H^0(Q,B_Q(3))$ form a vector bundle of rank nine, since its degree is
nine and the genus is one.  Squaring their sections gives the allowed
restrictions in $H^0(Q,\mathcal O_Q(6))$.
The restriction of ternary sextics is surjective, with kernel
$Q\cdot H^0(\mathbf P^2,\mathcal O(3))$ of dimension ten.
Thus the allowed sextics form, after this purely inseparable cover,
a nineteen-dimensional affine space, or a projective space of
dimension eighteen.  Smooth cubics have parameter dimension nine,
so the total is $9+18=27$, equal to the sextic parameter dimension.
The image of an irreducible space is irreducible.

The smooth-contact-cubic loci with effective $\theta_Q$ do not
dominate the sextic parameter space.  Their incidence dimensions are
at most 27 by the same degree-nine Riemann--Roch calculation; by
\eqref{eq:contact-reciprocity}, every such fibre over a sextic contains
the positive dimensional projective linear system of its
$\theta_C$, which has at least two sections.  Their images therefore
have dimension at most 26.  This remains true for the supersingular
cubic locus, whose base dimension is smaller.

\par\smallskip\noindent\textbf{Singular contact cubics.}
We show that these also occur only for sextics in a proper closed
subset of the 27-dimensional parameter space.
For every fixed cubic, the kernel of restriction of sextic forms has
dimension ten.  The following bounds are therefore bounds on the
dimension of the allowed restriction space, before adding ten and
projectivizing.

An irreducible nodal cubic has normalization $\mathbf P^1$ and two
points $P_+,P_-$ above the node.  Its stratum has dimension eight.
The pullback of $\mathcal O_Q(6)$ has degree eighteen.
Away from $P_+,P_-$, the orders of the restricted sextic are even.
At these two points the parities are equal.  If both are even, the
binary form is a square of a degree-nine form; the node gluing
condition gives one independent linear condition on that square root,
leaving dimension nine.  If both are odd, it is a fixed binary
quadratic with zeros $P_+,P_-$ times a square of degree eight, again
dimension nine; the gluing values are then both zero.  The incidence
dimension is consequently at most $8+9+10-1=26$.

An irreducible cuspidal cubic has normalization $\mathbf P^1$, one
point above its cusp, and parameter dimension at most seven.  All
orders of the pulled-back sextic are even, since its total degree is
eighteen and all orders away from that point are even.  It is therefore
a square of a degree-nine binary form, a space of dimension ten.
Any descent condition can only decrease the dimension.  The resulting
bound is $7+10+10-1=26$.  The dimensions eight and at most seven of
these strata follow directly by fixing a singular point: singularity
imposes three independent linear conditions on cubic forms, and a
repeated tangent cone imposes at least one further condition.

If the cubic is a transverse union of a line and a smooth conic, its
parameter dimension is seven.  On the two normalization components,
the sextic has degrees six and twelve.  Its odd orders, if any, must
occur at both of their two intersection points.  With no odd orders,
the square-root spaces have dimensions four and seven, subject to
two independent gluing conditions, giving nine.  With odd orders at
both points, extracting their fixed divisors leaves square-root
spaces of dimensions three and six, also giving nine.  The total
incidence dimension is at most 25.  If the line is tangent to the
conic, the cubic stratum has dimension at most six; each component's
restriction must have even order at the unique intersection point,
and the two square-root spaces satisfy at least their common-value
condition.  Their combined dimension is at most ten, again giving
at most 25.

For three distinct lines in a triangle, the parameter dimension is
six.  The possible odd-node patterns are either no nodes or all three
nodes.  In the first case the three four-dimensional square-root
spaces have three independent common-value conditions; in the second,
the three spaces have dimension three each.  Both bounds are nine,
giving total dimension at most 24.  For three concurrent lines the
parameter dimension is five.  Each line restriction has even order
at their sole common point, and the three square-root spaces satisfy
two common-value conditions.  Their combined dimension is at most
ten, again giving at most 24.

Finally a nonreduced cubic is $L^2M$, allowing $L=M$.
Its restriction to a smooth sextic can be even only if the line
section $M|C$ is even.  Sextics with such a line lie in dimension at
most 26: for a fixed line, binary sextics that are squares form a
four-dimensional subspace of a seven-dimensional restriction space;
the ternary restriction kernel has dimension 21.  Projectivizing and
allowing the line to vary gives $(4+21-1)+2=26$.
All possible singular cubics have now been covered.

Outside the proper loci just bounded, every one of the 1023 effective
theta characteristics of an ordinary sextic has a smooth contact
cubic, and this cubic corresponds to $B_Q$ by
\eqref{eq:contact-reciprocity}.  The irreducible incidence therefore
dominates the sextic space.  Its dimension is 27, so its generic
fibre is finite; every effective theta has $h^0=1$, and the fibre has
exactly 1023 geometric points.  As in the quintic case, irreducibility
implies geometric Galois transitivity on these points even if the
incidence has a purely inseparable part.  Thus $G_{\mathrm{geom}}$
is transitive on $V^\vee\setminus\{0\}$.

\par\smallskip\noindent\textbf{A boundary element that excludes the smaller transitive groups.}
Degenerate a sextic to an ordinary smooth quintic and a transverse
line, meeting in five distinct points.  This is an ordinary nodal
sextic.  Its rank-ten local system has a four-dimensional graph
submodule and the rank-six quintic local system as quotient.
The node permutation group is a subgroup of $S_5$; the graph
representation is its augmentation representation.
The exact quintic monodromy theorem supplies a quotient with order
divisible by seven and an element of order thirty-one acting on the
six-dimensional quotient as a one-dimensional identity plus an
irreducible five-dimensional summand.  Hence $7\mid|G_{\mathrm{geom}}|$.
Take the inverse image of the cyclic order-thirty-one subgroup in
the boundary monodromy group.  A Sylow thirty-one subgroup of this
inverse image maps onto the cyclic quotient.  There are no elements
of order $31^2$ in $\mathrm{GL}_{10}(2)$, since the multiplicative
order of $2$ modulo $31^2$ is 155.  Every nonidentity element of this
Sylow subgroup therefore has order thirty-one; choose one with
nonidentity quotient image and denote it by $g$.
It acts trivially on the graph submodule, since $31\nmid|S_5|$.
Its characteristic polynomial consequently has the form
\begin{equation}\label{eq:sextic-order31}
 \chi_g(T)=(T-1)^5 f(T),
\end{equation}
where $f$ is an irreducible polynomial of degree five whose roots
are primitive thirty-first roots of unity.
This polynomial is not reciprocal: inversion does not belong to
the orbit of squaring on primitive thirty-first roots, since
$\langle2\rangle\subset(\mathbf Z/31)^\times$ has odd order five.
As before, the specialization is inside the full-rank Lang locus,
so these are elements of the same generic monodromy image.

Apply Hering's classification of transitive linear groups in the
form of Churikov--Vasil'ev--Zvezdina (2022), Lemma~3.3, arXiv p.~6,
which was opened and checked.  In dimension ten over $\mathbf F_2$,
its alternatives, unless the group contains $\mathrm{SL}_{10}(2)$,
place the group in one of
\[
 N(\mathrm{Sp}_{10}(2)),\qquad
 \Gamma\mathrm L_5(4),\qquad
 \Gamma\mathrm L_2(32),\qquad
 \Gamma\mathrm L_1(1024).
\]
The exceptional alternatives have dimensions four or six or odd
characteristic and do not occur here; the $G_2$ alternative has
dimension divisible by six.
The symplectic normalizer preserves its unique invariant alternating
form up to a scalar, and over $\mathbf F_2$ the scalar is one.
Every one of its characteristic polynomials is reciprocal, contrary
to \eqref{eq:sextic-order31}.
An element of order thirty-one in $\Gamma\mathrm L_5(4)$ lies in
$\mathrm{GL}_5(4)$, so its fixed space has even dimension over
$\mathbf F_2$, contrary to the fixed dimension five in
\eqref{eq:sextic-order31}.
Finally, neither $\Gamma\mathrm L_2(32)$ nor
$\Gamma\mathrm L_1(1024)$ has order divisible by seven, which has
already been forced by the quintic quotient.
All proper alternatives are excluded.  Therefore
$G_{\mathrm{geom}}=\mathrm{SL}_{10}(2)=\mathrm{GL}_{10}(2)$.
The arithmetic image contains this group and is contained in it,
so it is the same group.
\end{proof}

% Characteristic-two plane septic case: separate pending independent audit.
\subsection{The exact group for ordinary plane septics in characteristic two}
\label{subsec:septic-exact}

\begin{theorem}\label{thm:septic-stabilizer-monodromy}
\status{Proved under stated hypotheses}
For the universal ordinary smooth plane septic in characteristic two,
the geometric and arithmetic monodromy images on
$V=H^1_{\mathrm{et}}(C_{\bar\eta},\mathbf F_2)$ are the full stabilizer
of the distinguished nonzero covector supplied by
$\mathcal O_C(2)$ and the canonical theta characteristic:
\[
 G_{\mathrm{geom}}=G_{\mathrm{arith}}
 =\operatorname{Stab}(\ell)
 \simeq 2^{14}\!:\mathrm{GL}_{14}(\mathbf F_2).
\]
In particular, the rank-fifteen representation is reducible, with its
distinguished invariant hyperplane, and its determinant is trivial.
\end{theorem}

\begin{proof}
We work first over an algebraic closure.  An ordinary genus-fifteen
curve has $2^{15}$ geometric theta characteristics.  Its canonical
theta characteristic $B_C$ is the unique non-effective one.  The
plane theta characteristic $L_C=\mathcal O_C(2)$ has six sections.
The class $v=L_C\otimes B_C^{-1}$ is the distinguished nonzero element
of $V^\vee$, fixed by arithmetic as well as geometric monodromy.
We prove transitivity on the other $2^{15}-2=32766$ elements of
$V^\vee$, then identify the whole stabilizer using a boundary subgroup.

\par\smallskip\noindent\textbf{Smooth contact quartics.}
Every nonzero section $s$ of an effective theta characteristic
$\theta_C$ supplies a unique quartic $Q$ with
$Q\cap C=2\operatorname{div}(s)$: the section $s^2$ belongs to
$H^0(C,\omega_C)=H^0(C,\mathcal O_C(4))$, and restriction from plane
quartics is an isomorphism.  If $Q$ is smooth, set
$D=\operatorname{div}(s)$ and
$\theta_Q=\mathcal O_Q(D-3H)$.  The contact reciprocity calculation
\eqref{eq:contact-reciprocity} gives
\[
 \theta_Q^2\simeq\omega_Q,\qquad
 h^0(\theta_C)=1+h^0(\theta_Q).
\]
For every smooth quartic and every theta characteristic on it,
$\theta_Q(3)$ has degree fourteen and exactly twelve sections, by
Riemann--Roch and Serre duality.  Restriction of septic forms to a
quartic is surjective and has kernel
$Q\cdot H^0(\mathbf P^2,\mathcal O(3))$, of dimension ten.

Consider first ordinary smooth quartics with $\theta_Q=B_Q$.
After a purely inseparable Frobenius base change, $B_Q$ is supplied
by the relative Frobenius sequence, so its twelve-dimensional spaces
of sections form a vector bundle.  Squaring their sections and lifting
through restriction gives an irreducible incidence of dimension
\[
 14+(12+10-1)=35,
\]
the dimension of the projective space of septics.  This incidence is
geometrically irreducible: the base is a nonempty open in the
geometrically irreducible quartic parameter space, the Frobenius
base change is a universal homeomorphism, and the section and lifting
spaces have geometrically irreducible affine fibres before
projectivizing.

For all smooth quartics and all theta characteristics, the incidence
dimension is at most thirty-five, since the scheme of roots is
finite over the quartic parameter space and the same twelve-section
calculation applies.  Those with effective $\theta_Q$ have images
of dimension at most thirty-four.  Indeed, reciprocity gives
$h^0(\theta_C)\ge2$, and the open part of
$\mathbf P H^0(C,\theta_C)$ near a smooth contact quartic is a
positive-dimensional fibre of this incidence over $C$.  This covers
the nonordinary smooth quartics as well, since all their geometric
theta characteristics are effective.  Thus a generic septic can
have a smooth contact quartic only through the canonical
non-effective theta characteristic of an ordinary quartic.

We next exclude singular contact quartics for the generic septic,
apart from the squares that give $L_C$.

\par\smallskip\noindent\textbf{Nonreduced quartics.}
A nonreduced quartic is either a square of a conic, or has the form
$Q=L^2R$ with $R$ a reduced conic.  A square gives exactly
$\theta_C=\mathcal O_C(2)$.  In the other case, evenness of $Q|C$
implies evenness of $R|C$.
For a fixed smooth conic, the restrictions of septics have binary
degree fourteen.  Those with even zero divisor are squares and form
an eight-dimensional space; the restriction kernel has dimension
$h^0(\mathbf P^2,\mathcal O(5))=21$.
Allowing the conic to vary gives dimension
$5+8+21-1=33$ for the image in the septic parameter space.
For a reduced pair of lines, each restricted binary form has odd
degree seven.  Its sole possible odd-order zero must be the node,
so each restriction is the fixed linear node factor times the
square of a cubic.  This leaves four parameters on each line;
both node values vanish, so their combined dimension is eight.
The conic parameter dimension is four, giving the smaller bound
$4+8+21-1=32$.  Consequently all non-square nonreduced contact
quartics occur over a proper subset of the septic parameter space.

\par\smallskip\noindent\textbf{Reduced reducible quartics.}
We give upper bounds that allow all coincidences and higher contacts
among components.  For a fixed reduced quartic, normalize each
component.  Away from the preimages of singular points, the divisor
of the restricted septic is even.  For each choice of the finite
odd-parity divisor $\epsilon$ on these preimages, its restriction
is a fixed section with divisor $\epsilon$ times a square.
The square-root line bundles form a finite set.  On a normalization
component of degree $a$ and genus $h$, the root has degree
$(7a-|\epsilon_i|)/2$ and, in the cases below, vanishing $H^1$;
hence its section dimension is
$(7a-|\epsilon_i|)/2+1-h$.
These degrees exceed $2h-2$: the only positive-genus component is
a smooth cubic, where the root degree is at least nine.
On a line the parity divisor has at most three points, and on a
conic at most four, so the genus-zero assertion is immediate.

For a smooth cubic and a line, the base dimension is eleven.
Both components have odd total intersection degree, so
$|\epsilon|\ge2$.  The sum of root dimensions is at most
$14+2-1-1=14$.  Adding the ten-dimensional restriction kernel and
projectivizing gives at most $11+14+10-1=34$.
If the cubic is irreducibly singular, its base dimension decreases
by at least one while the normalization genus decreases by one;
the same upper bound thirty-four follows.  Cases where the cubic
itself is reducible are included below.

For two smooth conics the base dimension is ten and the unrestricted
sum of root dimensions is sixteen.  If $\epsilon\ne0$, its total
degree is at least two, giving at most fifteen.  If
$\epsilon=0$, choose any intersection point of the two conics.
Descent of the septic restriction imposes equality of its two values,
and, after taking square roots over the algebraically closed field,
this is one independent linear condition on the root sections.
Independence follows from global generation on both normalization
components.  Thus the root dimension is again at most fifteen,
and the total bound is $10+15+10-1=34$.

For a smooth conic and two distinct lines the base dimension is
nine.  The two odd-degree line restrictions force
$|\epsilon|\ge2$, so the sum of root dimensions is at most
$14+3-1=16$.  The total is at most $9+16+10-1=34$.
For four distinct lines the base dimension is eight.  Each line
forces an odd parity divisor, so $|\epsilon|\ge4$ and the sum of
root dimensions is at most $14+4-2=16$; the total is at most
$8+16+10-1=33$.  No transversality assumption was used in these
upper bounds.

\par\smallskip\noindent\textbf{Integral singular quartics: normalization parameter spaces.}
An integral singular quartic has normalization of genus
$h\in\{0,1,2\}$.  We need two parameter bounds, and establish them
directly in characteristic two instead of using a Severi-variety
dimension assertion:
\begin{align*}
 \dim\{Q:\text{normalization genus }h\}&\le11+h,\\
 \dim\{Q:\text{normalization genus }h,
          \text{all singularities unibranch}\}&\le10+h.
\end{align*}
For $h=0$, a normalization map is given by a base-point-free
three-dimensional subspace of $H^0(\mathbf P^1,\mathcal O(4))$,
together with a projective frame.  The parameter dimension is
$\dim\operatorname{Gr}(3,5)+\dim\operatorname{PGL}_3
 -\dim\operatorname{PGL}_2=6+8-3=11$.
For $h=1$, temporarily mark an origin on the normalization.  Its
moduli have dimension one, its degree-four Picard variety dimension
one, the three-dimensional subspace of its four-dimensional section
space varies in dimension three, and a projective frame contributes
eight.  Forgetting the arbitrary origin subtracts one, giving
$1+1+3+8-1=12$.
For $h=2$, smooth-curve moduli have dimension three, the degree-four
Picard variety dimension two, and every degree-four line bundle has
three sections, giving $3+2+8=13$.
The dimension of genus-two moduli follows directly from curve
deformations: $H^0(T)=0$, $h^1(T)=3$ by Riemann--Roch, and $H^2(T)=0$.
The images under consideration are birational normalization maps,
so quotienting by reparametrization introduces no positive-dimensional
stabilizer.  These constructions cover all integral quartics of the
indicated normalization genus.

If all singularities are unibranch, the normalization map is
nonimmersive somewhere.  Otherwise at the preimage of each singular
point, one plane coordinate would be a uniformizer, making that
branch smooth.  For $h=0$ or $1$, and a degree-four line bundle $A$,
the map
$H^0(A)\to H^0(A|_{2P})$ is onto for every $P$, by Serre duality.
Among three-dimensional subspaces of $H^0(A)$, failure to have rank
two under this map has codimension two: for $h=1$ the rank-one locus
has dimension one in $\operatorname{Gr}(3,4)$ of dimension three;
for $h=0$ it has dimension four in
$\operatorname{Gr}(3,5)$ of dimension six.
Allowing $P$ to vary restores at most one dimension, so the
nonimmersive locus has codimension at least one in the parameter
space above.  This argument is invariant under reparametrization,
and the arbitrary elliptic origin remains a one-dimensional fibre.
For $h=2$, the system is the complete system $|A|$.
Failure to separate $2P$ is possible exactly when
$H^1(A(-2P))\ne0$, or $A\simeq\omega(2P)$.
As $P$ varies this is a subset of dimension at most one in the
two-dimensional degree-four Picard variety.  This remains a
dimension bound in characteristic two even if the map
$P\mapsto\omega(2P)$ is inseparable.  Thus the second bound follows
for all three genera.

Return to the restricted septic on such a quartic.  Its pullback
to the normalization has degree twenty-eight.  The odd-parity
divisor $\epsilon$ is supported above the finitely many singular
points.  There are at most six such branches, since the arithmetic
genus defect is at most three and the number of branches is at
most twice that defect.  A root after removing $\epsilon$ therefore
has degree at least eleven and vanishing $H^1$.  Its section
dimension is exactly
\[
 15-h-\frac{|\epsilon|}{2}.
\]
If $\epsilon\ne0$, then $|\epsilon|\ge2$; together with the
parameter bound $11+h$ this gives incidence dimension at most
$(11+h)+(14-h)+10-1=34$.
If $\epsilon=0$ and a singular point has two different branches,
descent imposes equality of the values at two distinct preimages.
The root line bundle separates these points, again by its degree,
so this is one independent linear condition after taking square
roots; the same bound thirty-four follows.
If all singularities are unibranch, the improved parameter bound
$10+h$ instead gives
$(10+h)+(15-h)+10-1=34$ without using any descent condition.
This excludes every integral singular contact quartic over a
generic septic.

\par\smallskip\noindent\textbf{Transitivity.}
Outside the proper closed subsets just bounded, every effective
theta characteristic other than $L_C$ has a smooth contact quartic,
and that contact quartic arises from its canonical non-effective
theta characteristic.  All $32766$ such theta characteristics
therefore occur in the geometrically irreducible canonical incidence
of dimension thirty-five.  This proves dominance.  Its generic
fibre is finite; in particular each of these theta characteristics
has exactly one section up to scalar, since otherwise its projective
section space would give a positive-dimensional fibre.  Consequently
the generic geometric fibre has precisely $32766$ points.
Irreducibility gives geometric Galois transitivity on these points;
a purely inseparable part of the incidence does not affect this
conclusion.  Hence $G_{\mathrm{geom}}$ is transitive on
$V^\vee\setminus\{0,v\}$.

\par\smallskip\noindent\textbf{The full quotient on the fourteen-dimensional space.}
Let $W=V^\vee/\langle v\rangle$, and let $H$ be the image on $W$.
Transitivity just proved implies transitivity of $H$ on
$W\setminus\{0\}$, of size $2^{14}-1$.
Degenerate to an ordinary sextic and a transverse line meeting it
in six points.  The rank-fifteen local system is still defined on
this nodal septic because Frobenius is bijective; its graph part
has rank five and its curve quotient has rank ten.  The projection
from this boundary family to the ordinary sextic parameter space
has geometrically connected fibres, so the quotient monodromy is
the full $\mathrm{GL}_{10}(2)$ established above.
In particular $11\mid|G_{\mathrm{geom}}|$.
The kernel of $\operatorname{Stab}(v)\to\mathrm{GL}(W)$ is a
two-group; therefore $11\mid|H|$.

Choose in the sextic quotient an order-thirty-one element acting
as a primitive five-dimensional block plus an identity of rank five.
An order-thirty-one lift exists by the Sylow argument used for
sextics: $\mathrm{GL}_{15}(2)$ has no element of order $31^2$, since
$\operatorname{ord}_{31^2}(2)=155$.
Its action on the graph part is trivial, since the node action
factors through $S_6$.  On $V^\vee$ it therefore has characteristic
polynomial $(T-1)^{10}f(T)$ for a primitive irreducible
degree-five factor $f$.  Quotienting by the fixed $v$ gives on $W$
\[
 (T-1)^9 f(T).
\]
This polynomial is not reciprocal, and the fixed-space dimension
is nine.

Hering's classification, in the opened form of
Churikov--Vasil'ev--Zvezdina (2022), Lemma~3.3, gives, in dimension
fourteen over $\mathbf F_2$, the alternatives
$\mathrm{GL}_{14}(2)$ or containment in
\[
 N(\mathrm{Sp}_{14}(2)),\quad
 \Gamma\mathrm L_7(4),\quad
 \Gamma\mathrm L_2(128),\quad
 \Gamma\mathrm L_1(16384).
\]
The symplectic alternative forces a reciprocal polynomial and is
excluded.  In $\Gamma\mathrm L_7(4)$ an order-thirty-one element
is $\mathbf F_4$-linear, so its fixed dimension over $\mathbf F_2$
is even, excluding nine.  The orders of the last two groups are
$1864254336$ and $229362$, respectively, neither divisible by eleven.
Thus $H=\mathrm{GL}_{14}(2)$.

\par\smallskip\noindent\textbf{The unipotent kernel and arithmetic equality.}
The intersection of $G_{\mathrm{geom}}$ with the kernel
$2^{14}$ of $\operatorname{Stab}(v)\to\mathrm{GL}(W)$ is an invariant
subspace for the full natural dual $\mathrm{GL}_{14}(2)$ module.
It is therefore zero or the whole kernel.  In the former case,
$G_{\mathrm{geom}}$ is a complement.  The elementary transvection
calculation $H^1(\mathrm{GL}_m(2),\mathbf F_2^m)=0$ for $m\ge4$,
proved in the quintic section (and equally valid for the dual),
shows that every such complement is conjugate to the standard one.
It has two orbits of size $2^{14}-1$ on
$V^\vee\setminus\{0,v\}$, contradicting transitivity.
Hence the whole kernel occurs, and the geometric group is the
full stabilizer.  Arithmetic monodromy also fixes $v$ and contains
the geometric group, so the two groups are equal.
\end{proof}

\par\smallskip\noindent\textbf{Finite certificate and scope.}
The retained program \texttt{sextic\_septic\_group\_checks.py}
checks the order of $2$ modulo $31^2$, the two distinct inverse
Frobenius orbits of primitive thirty-first roots, the semilinear
group orders and their missing prime divisors, and the two orbit
sizes used above.  It does not recompute Hering's classification
or replace the incidence argument by experimental Frobenius data.
This theorem settles degree seven in characteristic two only;
it does not assert the same stabilizer in every odd degree.

% All-degree characteristic-two contact argument; independent audit requested.
\subsection{Contact curves in every degree in characteristic two}
\label{subsec:all-degree-char2-contact}

The following argument removes the degree restriction from the
contact-incidence method.  Its global dimension input is
Christ--He--Tyomkin (2023), Proposition~2.7 (p.~8) and
Corollary~5.3 (p.~27), in the corrected publisher version of
\emph{On the Severi problem in arbitrary characteristic}.
Both statements apply over every algebraically closed field,
including characteristic two.  If a reduced plane curve of degree
$s$ has $c$ normalization components of genera $h_1,\ldots,h_c$,
their geometric genus convention is $1-c+\sum h_i$.  Consequently
the locus of such curves has dimension at most
\begin{equation}\label{eq:contact-severi-bound}
 3s+\sum_i h_i-c.
\end{equation}
If a family attains this dimension, its generic member is nodal.
The publisher's correction concerns the article history and missing
references; the opened corrected text contains these statements.

\begin{lemma}[The conductor inclusion for an even contact]
\label{lem:contact-conductor}
\status{Proved under stated hypotheses}
Let $C$ be a smooth plane curve of degree $d$ over an algebraically
closed field of characteristic two, and let $Q$ be a reduced plane
curve of degree $s<d$ with $Q\cap C=2D$ as divisors on $C$.
Write $\nu:\widetilde Q=\coprod Q_i\to Q$ for the normalization,
and write
\[
 \operatorname{div}(f|_{Q_i})=2A_i+\epsilon_i,
 \qquad M_i=\mathcal O_{Q_i}(A_i),
\]
where $f$ defines $C$ and $\epsilon_i$ is the reduced odd-order
divisor.  It is supported above the singular points of $Q$.
Set
\[
 J=\mathcal I_{D/Q}(s),\qquad
 L_C=\mathcal O_C(D),\qquad
 b=d-s.
\]
Then
\begin{equation}\label{eq:general-contact-reciprocity}
 h^0(C,L_C)=1+h^0(Q,J),
\end{equation}
and there is an inclusion
\begin{equation}\label{eq:conductor-section-inclusion}
 \nu_*\left(\bigoplus_i
       \omega_{Q_i}\otimes M_i^{-1}(3)\right)\hookrightarrow J.
\end{equation}
In particular, with $t=\sum_i h^1(Q_i,M_i)$,
\begin{equation}\label{eq:contact-h1-cancellation}
 h^0(Q,J)\ge t,
 \qquad t>0\Longrightarrow h^0(Q,J)\ge t+1.
\end{equation}
\end{lemma}

\begin{proof}
The same subscheme $D\subset\mathbf P^2$ is used on both curves.
Its ideal-sheaf sequences, twisted by $s$, are
\[
 0\to\mathcal O_{\mathbf P^2}(s-d)
  \to\mathcal I_D(s)\to L_C\to0,
 \qquad
 0\to\mathcal O_{\mathbf P^2}
  \to\mathcal I_D(s)\to J\to0.
\]
Here $\mathcal O_C(s)(-D)\simeq\mathcal O_C(D)$ because
$2D=Q|C$.  Since $s<d$ and the first cohomology of every line
bundle on $\mathbf P^2$ vanishes, these sequences prove
\eqref{eq:general-contact-reciprocity}.

We prove the inclusion locally and then express it globally.
At a point of $D$, complete the local ring of the plane.  Since
$C$ is smooth, choose its equation $f$ as one parameter and let
$b_0$ cut out the half divisor on $C$.  After multiplying the local
equation of $Q$ by a unit, that equation is
$b_0^2-fa_0$.  In the reduced one-dimensional ring
$A=\widehat{\mathcal O}_{Q,P}$, the ideal of $D$ is
$I=(f,b_0)$ and $f$ is a non-zero-divisor.
Reduction modulo $f$ gives
$A/(f)\simeq k[[x]]/(b_0^2)$, where the annihilator of $b_0$ is
$(b_0)$.  Hence
\[
 I^*:=\operatorname{Hom}_A(I,A)
 =\{z:zI\subset A\}=f^{-1}I,
 \qquad I^{**}=I.
\]
These equalities can be read directly in the total quotient ring:
$zf\in A$ and $zb_0\in A$ say that $zf$ belongs to
$(f):(b_0)=I$.

Let $v_i$ be the valuation on a normalization branch.  Put
\[
 r_i=v_i(f),\qquad j_i=\min\{v_i(f),v_i(b_0)\},
 \qquad\gamma_i=2j_i-r_i.
\]
The relation $b_0^2=fa_0$ gives $\gamma_i\ge0$; moreover
$\gamma_i\equiv r_i\pmod2$.  Thus
$\gamma_i=\epsilon_i(P_i)+2k_i$ for an integer $k_i\ge0$.
Write $c_i$ for the conductor exponent on this branch, so that
the conductor ideal on the normalization is generated there by
a parameter to the power $c_i$.
If a tuple $z$ of branch functions satisfies
\[
 v_i(z)\ge c_i+r_i-j_i\quad\hbox{for every branch},
\]
then its product with every member of $I^*=I/f$ lies in the
conductor, hence in $A$.  It follows that $z\in I^{**}=I$.
No separate assumption on a tame singularity, on its number of
branches, or on its analytic type is used here.

Let $\Delta$ be the conductor divisor on $\widetilde Q$, let
$\Gamma=\sum\gamma_iP_i$, and let $K_0=\sum k_iP_i$.
If $N=\nu^*J/\text{torsion}$, the valuation calculation gives
\[
 N|_{Q_i}\simeq M_i(-b)(-K_0|_{Q_i}),\qquad
 \nu_*N(\Gamma-\Delta)\subset J.
\]
Indeed the normalized ideal has vanishing order $j_i$, and
$\gamma_i=2j_i-r_i$ is precisely the adjustment from its square
to the divisor of $f$.

For a reduced plane curve, finite duality (Stacks Project,
Remark~48.20.12, Tag~0AX4, opened and checked) and adjunction give
\[
 \omega_{\widetilde Q}
 \simeq\nu^*\omega_Q(-\Delta)
 =\mathcal O_{\widetilde Q}(s-3)(-\Delta).
\]
For completeness, the conductor occurs because
$\operatorname{Hom}_{\mathcal O_Q}
(\nu_*\mathcal O_{\widetilde Q},\mathcal O_Q)$ is the conductor
ideal; finite duality identifies its tensor product with
$\omega_Q$ with $\nu_*\omega_{\widetilde Q}$.
Using $M_i^2\simeq\mathcal O_{Q_i}(d)(-\epsilon_i)$ now yields
\[
 \omega_{Q_i}\otimes M_i^{-1}(3)
 \simeq M_i(-b)(\epsilon_i-\Delta_i)
 =N|_{Q_i}(K_{0,i}+\epsilon_i-\Delta_i).
\]
Since $\Gamma=\epsilon+2K_0$, this line bundle is a subsheaf
of $N(\Gamma-\Delta)$, proving
\eqref{eq:conductor-section-inclusion}.

By Serre duality, $h^1(M_i)=h^0(\omega_{Q_i}M_i^{-1})$.
Twisting by the positive line bundle $\mathcal O_{Q_i}(3)$
does not decrease this dimension.  If the initial space is nonzero,
the increase is strict: choose a point where some section does not
vanish, and choose a plane cubic through its image whose pullback
is a nonzero divisor.  Multiplication by that cubic has image
consisting of sections vanishing there, whereas multiplying the
chosen initial section by a cubic not through the point gives a
section outside the image.  Summing over the components proves
\eqref{eq:contact-h1-cancellation}.
\end{proof}

\begin{lemma}[No lower-degree reduced contact for a generic curve]
\label{lem:no-low-degree-contact}
\status{Proved under stated hypotheses}
Fix $d\ge4$ and $0<s<d-3$ in characteristic two.
A generic smooth plane curve $C$ of degree $d$ has no reduced
plane curve $Q$ of degree $s$ with $Q|C$ even.
\end{lemma}

\begin{proof}
Let $N_d=\binom{d+2}{2}-1$.
For fixed normalization genera, component degrees, and odd-parity
divisors $\epsilon_i$, the line bundles
$M_i^2=\mathcal O_{Q_i}(d)(-\epsilon_i)$ have finite choices.
They can be parameterized after a finite base change of the
singularity stratum; inseparable base changes do not affect
dimension.  The normalized restrictions of $f$ are squares of
sections of these $M_i$, times the fixed odd-parity factors.
Ignoring descent conditions only increases the dimension.
Riemann--Roch gives
\[
 \sum_i h^0(M_i)
 =\frac{ds}{2}-\frac{|\epsilon|}{2}
   +c-\sum_i h_i+t,
 \qquad t=\sum_i h^1(M_i).
\]
The kernel of restricting degree-$d$ forms to $Q$ has dimension
$\binom{d-s+2}{2}$.  Combining this with
\eqref{eq:contact-severi-bound}, and projectivizing, gives the
incidence bound
\begin{equation}\label{eq:general-contact-dimension}
 \dim\mathcal I_{d,s}
 \le N_d-\frac{s(d-s-3)}2-\frac{|\epsilon|}2+t
\end{equation}
on each stratum with fixed $t$.

Here is the precise fibre argument used with this bound.
Suppose a component of the even-contact incidence dominates the
ordinary smooth degree-$d$ parameter space.  Over its generic
curve, roots $L_C^2=\mathcal O_C(s)$ form a finite scheme.
For each geometric root the contact forms are the projective
space of its nonzero sections, followed by squaring.
Restriction of degree-$s$ plane forms to $C$ is an isomorphism,
so no further choices of $Q$ occur.  The reduced-contact condition
is open in this projective section space.  The generic fibre of
our component therefore has dimension
$h^0(L_C)-1=h^0(J)$.
Its generic point belongs to a single stratum of the finite
stratification by the data above, and an open subset of the same
generic fibre remains in that stratum.  Thus the fibre dimension
may be subtracted in \eqref{eq:general-contact-dimension}; we are
not subtracting a family of contact curves from a special stratum
that it need not preserve.

The conductor lemma gives $h^0(J)\ge t$.  Dominance would imply
\[
 N_d\le N_d-\frac{s(d-s-3)}2-\frac{|\epsilon|}2<N_d,
\]
a contradiction.  There are finitely many relevant component
degrees, genera, parity choices in finite covers, and cohomology
dimensions.  Thus their images have proper closure, proving the
claim for a generic curve.
\end{proof}

\begin{theorem}[The theta orbits of generic plane curves in characteristic two]
\label{thm:all-char2-theta-orbits}
\status{Proved under stated hypotheses}
Let $d\ge4$, $g=(d-1)(d-2)/2$, and let $C$ be the geometric
generic ordinary smooth plane curve of degree $d$ in characteristic
two.
If $d$ is even, every theta characteristic other than the canonical
one $B_C$ has exactly one section, and geometric monodromy is
transitive on these $2^g-1$ characteristics.
If $d$ is odd, every theta characteristic other than $B_C$ and
$\mathcal O_C((d-3)/2)$ has exactly one section, and geometric
monodromy is transitive on these $2^g-2$ characteristics.
Equivalently, on $V^\vee$ geometric monodromy is transitive on
all nonzero vectors in even degree, and on
$V^\vee\setminus\{0,v\}$ in odd degree, where $v$ is the
distinguished nonzero invariant vector.
\end{theorem}

\begin{proof}
Put $e=d-3$.  A section of an effective theta characteristic
$\theta_C$ gives a contact form $Q$ of degree $e$, and conversely
such a form gives an effective theta characteristic and its
section up to scalar.  Work throughout over the ordinary smooth
degree-$d$ locus, so theta characteristics have finitely many
geometric choices, with constant count $2^g$.  Passing to a
purely inseparable cover if necessary changes neither dimensions
nor geometric Galois orbits.

\par\smallskip\noindent\textbf{The unique dominant smooth-contact incidence.}
For a smooth $Q$, contact reciprocity gives
$h^0(\theta_C)=1+h^0(\theta_Q)$, where
$\theta_Q=\mathcal O_Q(D-3)$.
Riemann--Roch gives
\[
 h^0(\theta_Q(3))=3e+h^0(\theta_Q(-3)).
\]
Thus the incidence dimension, for a stratum with these dimensions
constant, is at most $N_d+h^0(\theta_Q(-3))$.
If $\theta_Q$ is effective, then
$h^0(\theta_Q(-3))<h^0(\theta_Q)$: multiplication by a suitable
cubic gives a proper subspace, as in the final paragraph of the
conductor lemma.  The fibre dimension $h^0(\theta_C)-1$
is $h^0(\theta_Q)$, so such an incidence cannot dominate.
For an ordinary smooth $Q$, its canonical $B_Q$ is the unique
non-effective theta characteristic.  It has
$h^0(B_Q(3))=3e$.  For a nonordinary smooth $Q$, all its geometric
theta characteristics are effective.  Consequently the only
possible dominant smooth-contact incidence is the canonical one
over ordinary smooth $Q$.

That canonical incidence is geometrically irreducible of dimension
$N_d$: the ordinary degree-$e$ base is geometrically irreducible;
after Frobenius base change the relative canonical theta
characteristic exists; its $3e$ sections form a vector bundle;
and lifts of their squares through restriction form affine spaces
with the constant ten-dimensional kernel
$Q\cdot H^0(\mathbf P^2,\mathcal O(3))$.
Projectivization gives total dimension
$\binom{e+2}{2}-1+3e+10-1=N_d$.
For $e=1,2$ the genus-zero canonical theta characteristic is used,
and the same computation applies.

\par\smallskip\noindent\textbf{Reduced singular contact forms cannot dominate.}
Apply the notation and incidence bound of the preceding lemmas
with $s=e$.  A hypothetical dominating component has dimension
at most $N_d-|\epsilon|/2+t$, and generic fibre dimension
$h^0(J)$.  If $t>0$, the strict part of
\eqref{eq:contact-h1-cancellation} gives a contradiction.
If $|\epsilon|>0$, its total degree is even and at least two,
so the non-strict inequality also gives a contradiction.
If $h^0(J)>0$ while $t=|\epsilon|=0$, the fibre dimension again
gives a contradiction.  Only
\[
 t=0,\qquad \epsilon=0,\qquad h^0(J)=0
\]
could remain for a dominant component.

In this case its image in the degree-$e$ curve parameter space
must attain the maximal Severi dimension
$3e+\sum h_i-c$; otherwise the incidence bound loses at least
one further dimension.  By Christ--He--Tyomkin Corollary~5.3,
the generic $Q$ is therefore nodal.  Since it is singular, choose
one node.  The conductor inclusion and $h^0(J)=0$ imply
\[
 H^0(Q_i,\omega_{Q_i}M_i^{-1}(3))=0
 \quad\hbox{on every component}.
\]
If the two preimages $P_+,P_-$ of the node lie on one component,
there exists a plane cubic whose pullback contains both points.
It follows that
$H^0(\omega M_i^{-1}(P_++P_-))=0$, hence
$H^1(M_i(-P_+-P_-))=0$.  The evaluation map at the two points
is therefore onto.  If they lie on different components, the
same argument with one point proves surjectivity of each
one-point evaluation.  In either case, equality of the two values
needed for descent of the squared restriction imposes one
nonzero linear condition on the root sections after taking square
roots of its scalar coefficients.  This lowers the incidence
dimension by at least one, the final contradiction.

\par\smallskip\noindent\textbf{Nonreduced contact forms.}
Every nonreduced form of degree $e$ has a unique expression up to
scalars $Q=R^2S$, with $S$ reduced and $\deg S=s<e$.
If $s>0$, evenness of $Q|C$ implies evenness of $S|C$, excluded
for a generic $C$ by Lemma~\ref{lem:no-low-degree-contact}.
If $s=0$, then $e$ is even, and the half divisor of $Q|C$ is
the divisor of $R|C$.  It gives exactly the plane theta
characteristic $\mathcal O_C(e/2)$.  This case occurs precisely
when $d$ is odd and is the sole nonreduced contact case retained.

\par\smallskip\noindent\textbf{Counting the generic fibre and obtaining the orbit.}
Every noncanonical theta characteristic is effective.
All except the plane characteristic in odd degree must therefore
occur in the geometrically irreducible canonical smooth-contact
incidence.  This proves its dominance.  Since it has dimension
$N_d$, its generic fibre is finite, so each of these theta
characteristics has exactly one section up to scalar.
The generic geometric fibre has respectively $2^g-1$ or
$2^g-2$ points.  Irreducibility makes geometric Galois monodromy
transitive on them; purely inseparable factors of the incidence
do not change this orbit statement.  Translating theta
characteristics by the canonical $B_C$ identifies their set
with $V^\vee$, proving the asserted vector orbits.
\end{proof}


\subsection{Exact characteristic-two plane monodromy in every degree}
\label{subsec:char2-uniform}

\begin{lemma}[Two elementary group-theoretic inputs]
\label{lem:char2-group-inputs}
\status{Proved under stated hypotheses}
For every integer $a\ge8$, there is an odd prime $\ell_a$ with
$\operatorname{ord}_{\ell_a}(2)=a$.
For $N\ge4$, the first cohomology of $\mathrm{GL}_N(\mathbf F_2)$
with coefficients in its natural module or its dual is zero.
\end{lemma}
\begin{proof}
We include a proof of the needed special case of the primitive-divisor
theorem. If an odd prime $\ell$ divides $\Phi_a(2)$ without having
order $a$, the cyclotomic factorization and the lifting-the-exponent
identity show that
$a=\operatorname{ord}_{\ell}(2)\ell^j$ with $j\ge1$, and that
$v_\ell(\Phi_a(2))=1$. Indeed, substitute
$v_\ell(2^{fm}-1)=v_\ell(2^f-1)+v_\ell(m)$, where
$f=\operatorname{ord}_\ell(2)$, in the M\"obius formula for
$v_\ell(\Phi_a(2))$; the sum is zero unless $a/f$ is a power of
$\ell$, when it is one. Since $f\mid\ell-1$, such an $\ell$ is the
largest prime factor $P$ of $a$. Thus, if no primitive prime occurs,
$\Phi_a(2)$ is either $1$ or $P$.
On the other hand the cyclotomic product gives
\[
 \Phi_a(2)>2^{\varphi(a)}\prod_{j\ge1}(1-2^{-j})
 >2^{\varphi(a)-2}.
\]
For the last bound, the first three factors have product $21/64$,
and the remaining product is at least
$1-\sum_{j\ge4}2^{-j}=7/8$; their product exceeds $1/4$.
If $P\ge7$, then $\varphi(a)\ge P-1$ and $2^{P-3}>P$.
If $P\le3$ and $a\ge8$, then $\varphi(a)\ge4$ and the same
conclusion holds. If $P=5$, either $\varphi(a)\ge8$ or $a=10$;
the exceptional calculation is $\Phi_{10}(2)=11$.
This contradiction proves the first assertion. The cyclotomic
valuation calculation also agrees with
\cite[Lemma~3.1(b)--(c), pp.~125--126]{GlasbyEtAl2017}, whose
statement and proof were opened.

For the second assertion let $z$ be a $1$-cocycle with values in
$\mathbf F_2^N$ and put $t_{ij}=1+E_{ij}$.
The relation $t_{ij}^2=1$ forces the $j$th coordinate of
$z(t_{ij})$ to vanish. For $l\ne i,j$, choose
$k\notin\{i,j,l\}$. Comparing the cocycle on the commuting pair
$t_{ij},t_{kl}$ shows that its $l$th coordinate also vanishes.
Thus $z(t_{ij})=a_{ij}e_i$.
For distinct $i,j,k$, apply the cocycle to
$[t_{ij},t_{jk}]=t_{ik}$; it gives $a_{ik}=a_{jk}$.
Consequently $a_{ij}=b_j$ for a single vector $b$, and
$z(t_{ij})=(t_{ij}-1)b$. Since these transvections generate the
group, $z$ is a coboundary. Composing with inverse transpose
gives the dual-module assertion.
\end{proof}

\begin{theorem}[All plane degrees in characteristic two]
\label{thm:plane-char2-all-groups}
\status{Proved under stated hypotheses}
Let $k$ be any perfect field of characteristic two, let $U_d$ be
the ordinary smooth locus in the projective space of degree-$d$
plane forms, and choose a geometric generic base point $\bar\eta$.
Write
\[
 g=(d-1)(d-2)/2,\qquad
 V=H^1_{\mathrm{et}}(C_{\bar\eta},\mathbf F_2).
\]
For $d\ge4$, the images of $\pi_1(U_d\otimes_k\bar k,\bar\eta)$
and $\pi_1(U_d,\bar\eta)$ on $V$ are equal and are
\[
 \begin{cases}
 \mathrm{GL}_g(\mathbf F_2),&d\text{ even},\\
 \operatorname{Stab}(\ell)
   \simeq(\mathbf F_2)^{g-1}\rtimes
                  \mathrm{GL}_{g-1}(\mathbf F_2),&d\text{ odd},
 \end{cases}
\]
where $\ell\ne0$ is the canonical covector determined by
$\mathcal O_C((d-3)/2)\otimes B_C^{-1}$ in odd degree.
The nonzero representation is irreducible exactly in even degree;
its determinant is always trivial. In degrees one and two $V=0$;
in degree three it has rank one and trivial, hence irreducible,
monodromy. Thus there is no threshold after which all plane
degrees have full linear or irreducible monodromy.
\end{theorem}
\begin{proof}
The degrees at most four, the quintic stabilizer, and the sextic
full group have been proved above. We induct on $d\ge7$.
Work first over $\bar k$ and on the dual space $W_0=V^\vee$.
Theorem~\ref{thm:all-char2-theta-orbits} proves transitivity on
$W_0\setminus\{0\}$ when $d$ is even and on
$W_0\setminus\{0,\ell\}$ when $d$ is odd.
In the odd case put $W=W_0/\langle\ell\rangle$; in the even
case put $W=W_0$. The induced group on $W$ is therefore transitive
on its nonzero vectors. Its dimension is
\[
 N=\begin{cases}g,&d\text{ even},\\g-1,&d\text{ odd}.\end{cases}
\]

We specify the boundary subgroup used in the induction. The
universal plane-curve coherent $H^1$ is a vector bundle of rank
$g$, including singular fibres. The locus where its Frobenius is
bijective is open, and the Lang construction gives an etale
$\mathbf F_2$-local system on this whole locus. It contains the
nodal curves $C'\cup L$, where $C'$ is ordinary smooth of degree
$d-1$ and $L$ meets it transversely in $d-1$ distinct points.
The smooth ordinary locus is dense in this normal parameter
space. Restriction of connected finite etale covers to a dense
open remains connected, so its fundamental group maps onto the
fundamental group of the larger locus. Boundary monodromy thus
lies in the group being determined.

The normalization sequence for $C'\cup L$ gives
\[
 0\longrightarrow H^1(\Gamma,\mathbf F_2)
 \longrightarrow H^1_{\mathrm{et}}(C'\cup L,\mathbf F_2)
 \longrightarrow H^1_{\mathrm{et}}(C',\mathbf F_2)
 \longrightarrow0.
\]
Here the graph space has dimension $d-2$ and its monodromy factors
through the permutation group $S_{d-1}$ of the nodes. The
projection from the parameter space of transverse pairs $(C',L)$
to that of $C'$ is smooth, surjective, and has geometrically
connected fibres, namely nonempty opens of the dual plane.
It therefore surjects on fundamental groups and supplies the
entire already determined monodromy of $C'$.

Write $r=(d-2)(d-3)/2$. If $d$ is even, the preceding odd-degree
group has a quotient $\mathrm{GL}_{r-1}(2)$; the kernel of this
quotient in the boundary group has prime divisors only $2$ and
those dividing $(d-1)!$. If $d$ is odd, the dual boundary
sequence has an invariant subspace of dimension $r$ with full
$\mathrm{GL}_r(2)$ action. The nonzero invariant vector $\ell$
cannot lie in that subspace, since it has no nonzero invariant
vector. It remains nonzero at the boundary: an invariant section
of a finite etale sheaf on the dense open extends uniquely over
the normal larger locus, and the zero and nonzero sections are
disjoint. Thus the rank-$r$ subspace injects into $W$, with quotient
a subquotient of the graph permutation space. In this case also,
the kernel of the boundary quotient onto $\mathrm{GL}_r(2)$ has
prime divisors only $2$ and those dividing $(d-1)!$. These kernel
claims follow by writing the representation in block triangular
form: its off-diagonal part is a $2$-group, the preceding odd-degree
unipotent radical is a $2$-group, and its remaining diagonal
action factors through $S_{d-1}$.

Set $a=\lfloor N/2\rfloor+1$. The elementary inequalities
\[
 8\le a<a+1<N-1,\qquad a>d-1,\qquad
 a+1\le\begin{cases}r-1,&d\text{ even},\\r,&d\text{ odd}
 \end{cases}
\]
hold for even $d\ge8$ and odd $d\ge7$. For example the two initial
cases have $(N,r,a)=(21,15,11)$ and $(14,10,8)$.
For the upper bound, using $\lfloor N/2\rfloor\le N/2$
reduces it to $g\ge2d+2$ in even degree and $g\ge2d-1$ in odd
degree; both hold at the indicated initial degree and their
left-minus-right differences increase. The other inequalities
follow directly from the same genus formula.
For $j=a,a+1$, choose the prime $\ell_j$ in
Lemma~\ref{lem:char2-group-inputs}. We have $\ell_j>j>d-1$.
Also $\ell_j\nmid N$: otherwise $\ell_j>N/2$ forces $N=\ell_j$,
whereas $j\mid N-1$ and $N/2<j<N-1$, an impossibility.

In the boundary general linear quotient choose an element of
order $\ell_j$ with one irreducible $j$-dimensional block and
identity on its complement. The kernel has order prime to
$\ell_j$, so a Sylow subgroup of its inverse image supplies an
element of order $\ell_j$ with that image. It acts trivially on
the graph quotient, because $\ell_j\nmid(d-1)!$.
Semisimplicity of this odd-order action shows that its
characteristic polynomial on $W$ is
\[
                    (T-1)^{N-j}f_j(T),
 \qquad f_j\text{ irreducible of degree }j.
\]

Apply the opened classification in
\cite[Lemma~3.3, p.~6]{ChurikovVasilevZvezdina2022}.
Every proper extension-field alternative lies inside
$\Gamma\mathrm L_{N/b}(2^b)$ for some divisor $b>1$ of $N$.
Our order-$\ell_j$ elements belong to its linear subgroup,
since $\ell_j\nmid b$. Their fixed-space dimensions are therefore
divisible by $b$. The consecutive dimensions $N-a$ and $N-a-1$
force $b=1$, excluding all these alternatives, including the
one-dimensional semilinear case. Over the original field the
remaining symplectic alternative lies in $\mathrm{Sp}_N(2)$:
its normalizer preserves the unique invariant alternating form,
and the only nonzero scalar is one. One of $a,a+1$ is odd.
For this $j$, the irreducible factor $f_j$ is not reciprocal,
since inversion is not in the odd-order subgroup generated by
$2$ modulo $\ell_j$. Its characteristic polynomial contradicts
preservation of a nondegenerate alternating form. The $G_2$
alternative over a larger field was already excluded, and its
field-two exceptions have dimension six. The other exceptional
field-two cases have dimension four. Here $N\ge14$.
The classification therefore forces the image on $W$ to be
$\mathrm{SL}_N(2)=\mathrm{GL}_N(2)$.

For even $d$ this finishes the geometric assertion. For odd $d$
the kernel of the image in $\mathrm{GL}(W)$ is an invariant
subgroup of $\operatorname{Hom}(W,\langle\ell\rangle)$.
This is the irreducible dual natural module, so the kernel is
zero or the whole module. If it were zero, the group would be a
complement. Lemma~\ref{lem:char2-group-inputs} makes every such
complement conjugate to the block diagonal group. That group has
two separate orbits on $W_0\setminus\{0,\ell\}$, contradicting
the theta-orbit theorem. Thus the kernel is the whole module,
giving the full stabilizer and completing the induction.

The distinguished odd-degree covector is defined over the original
perfect field, as proved in the explicit Hasse--Witt construction.
Arithmetic monodromy is contained in its stabilizer and contains
the geometric group just determined; in even degree it is
contained in the full general linear group. Thus the two images
are equal in both cases. Irreducibility and the determinant
assertions follow immediately; $\mathbf F_2^\times$ is trivial.
\end{proof}

% Agent /root/monodromy; complete argument, 19 November 2026.
% Independently audited in full by /root, including the opened classification source.
\subsection{Plane curves in odd characteristic: the full group in every degree}
\label{sec:odd-plane-all-degrees}
\status{Proved; independently audited by /root, including the primary classification statement.}

Throughout this section $p$ is odd, $k$ is an algebraic closure of
$\mathbf F_p$, $d\geq4$, and $g=(d-1)(d-2)/2$.  Let $W_d$ be the
Hasse--Witt invertibility open in the projective space of degree-$d$
plane equations, including singular equations, and let $U_d\subset W_d$
be its smooth open.  The vector bundle $R^1\pi_*\mathcal O$ has constant
rank $g$ on the entire projective parameter space.  Its Lang fixed
sheaf $V$ is finite locally constant of rank $g$ on $W_d$.
The map $\pi_1(U_d)\to\pi_1(W_d)$ is surjective.  Thus computing its
image on $W_d$ computes the desired geometric image on $U_d$.
Choose as base point a union of $d$ lines with no triple intersection.
All groups below act on this same fibre, after choosing paths in the
specified connected parameter spaces.  This is permissible because
every boundary family used below is contained in $W_d$.

\subsubsection{The line-arrangement representation}

Put $P=\mathbf F_p^d$, with its permutation basis $e_1,\ldots,e_d$,
$u=e_1+\cdots+e_d$, and $W=P/\langle u\rangle$.  The graph of the
line arrangement is the complete graph $K_d$.  Its first cochain space
is $\bigwedge^2P$, with $e_i\wedge e_j$ representing an oriented edge.
The coboundaries are $u\wedge P$.  Consequently
\begin{equation}\label{eq:odd-arrangement-module}
 V=H^1(K_d,\mathbf F_p)
   =\bigwedge^2P/(u\wedge P)=\bigwedge^2W.
\end{equation}
This is cohomology, rather than the dual cycle space.  The coherent
normalization sequence gives the same identification after tensoring
with $k$, with bijective Frobenius on all the graph terms.

The space of ordered general-position lines is irreducible; its
quotient by the freely acting group $S_d$ is the unordered arrangement
space.  Its ordered cover is connected, so its permutation monodromy
is $S_d$.  Over the ordered cover the graph local system is constant.
Thus the image $G\subseteq\mathrm{GL}(V)$ contains exactly the
permutation action in \eqref{eq:odd-arrangement-module}.

Let $U=\{(x_i)\in P:\sum x_i=0\}$, let $D$ be its image in $W$, and
put $A=\bigwedge^2D\subseteq\bigwedge^2W$.

\begin{lemma}\label{lem:odd-wedge-submodules}
\status{Proved under stated hypotheses}
Every nonzero $S_d$-submodule of $\bigwedge^2W$ contains $A$.
If $p\nmid d$, then $A=V$ and $V$ is simple.  If $p\mid d$, then
$A$ is the unique proper nonzero submodule, and $V/A\simeq D$ is
simple of dimension $d-2$.
\end{lemma}

\begin{proof}
For four distinct indices $i,j,k,l$, put $r=e_i-e_j$ and $s=e_k-e_l$.
The group-algebra operator
\[
 (1-(ij))(1-(kl))
\]
on $\bigwedge^2W$ has image $\langle r\wedge s\rangle$ and is the
rank-one map obtained by evaluating the functional
$(x_i-x_j)\wedge(x_k-x_l)$ and multiplying by $r\wedge s$.
The latter vector is nonzero: the line $\langle u\rangle$ does not
meet $\langle r,s\rangle$, including when $d=4$, since $p\ne2$.

The vectors $r\wedge s$ for disjoint pairs span $\bigwedge^2U$ before
passage to $W$, and hence span $A$ afterwards.  Explicitly, on four
indices take $f_i=e_i-e_l$, $f_j=e_j-e_l$, $f_k=e_k-e_l$.  The three
perfect matchings give the columns
\[
 \begin{pmatrix}0&1&1\\1&0&-1\\-1&1&0\end{pmatrix}
\]
in the basis $f_i\wedge f_j,f_i\wedge f_k,f_j\wedge f_k$; their
determinant is $2$.  Every pair of difference-basis vectors belongs to
such a four-index set.  The same argument in the dual permutation
space shows that the corresponding functionals span
$\bigwedge^2W^*$, because $W^*$ is the sum-zero subspace of $P^*$.
For every nonzero vector of $V$, therefore, one of these rank-one
operators has nonzero value.  Its entire $S_d$-orbit spans $A$.
This proves the first assertion and also proves that $A$ is simple.

If $p\nmid d$, then $U\to W$ is an isomorphism.  If $p\mid d$,
$D=U/\langle u\rangle$ has dimension $d-2$, and the sum functional
$\sigma$ descends to $W$ with kernel $D$.  Contraction with $\sigma$
gives the exact sequence
\[
 0\longrightarrow\bigwedge^2D\longrightarrow\bigwedge^2W
 \xrightarrow{\ \iota_\sigma\ }D\longrightarrow0.
\]
The module $D$ is simple: a nonzero class has a representative whose
coordinates are not all equal; applying $1-(ij)$ for two unequal
coordinates produces a nonzero multiple of $e_i-e_j$, and the
translates of this difference span $D$.  Every proper nonzero
submodule of $V$ contains $A$ and has zero image in the simple
quotient unless it is all of $V$.  This proves uniqueness.
\end{proof}

\begin{proposition}\label{prop:odd-plane-irreducible}
\status{Proved under stated hypotheses}
The geometric monodromy group $G$ is irreducible on $V$.
\end{proposition}

\begin{proof}
Only the case $p\mid d$ needs another family.  Take a smooth ordinary
curve of degree $d-2$ and a smooth conic meeting it transversely in
$m=2d-4$ points.  The graph contribution is a sub-local-system
\[
 B=\mathbf F_p^m/\langle(1,\ldots,1)\rangle\subset V
\]
of rank $m-1=2d-5$.  Its node monodromy is the full symmetric group
$S_m$.  To verify this, fix the conic.  Restriction of degree-$(d-2)$
plane forms to it is surjective onto binary forms of degree $m$.
Over the squarefree binary forms, the ordered-root cover is connected.
Pulling it back gives an affine bundle after choosing a coefficient
chart (equivalently, an open of a projective bundle).  Imposing
smoothness and ordinarity of the plane curve removes an open subset
from this irreducible total space, and the open is nonempty.  The
ordered-root cover remains connected.  Thus the induced graph image
is $S_m$.

Here $p\nmid m$, since $p\mid d$ and $p$ is odd.  The deleted
permutation module $B$ is simple, by the same difference argument
used above.  If the only possible proper $G$-invariant subspace $A$
were invariant, the map $B\to V/A$ would be zero: its source is a
simple module of dimension $2d-5>d-2=\dim(V/A)$.  Hence $B\subset A$.

This inclusion can be tested at the common line arrangement.  Allow
the conic to become the union of the first two lines and the other
curve to become the remaining $d-2$ lines.  Keep the intersections
between the two subcurves transverse, disjoint from their internal
nodes.  Their ordinary coherent cohomology and the graph sequence
extend over an open containing both this point and the smooth
subcurve locus.  At this arrangement $B$ is the kernel of restriction
to the complete subgraph on vertices $3,\ldots,d$.  In particular the
edge cochain $e_1\wedge e_3$ represents an element of $B$.  But
\[
 \iota_\sigma(e_1\wedge e_3)=e_3-e_1\ne0\quad\hbox{in }D,
\]
so it does not belong to $A$.  This contradiction proves the claim.
No specialization through a drop of the Hasse--Witt rank has been
used.
\end{proof}

\subsubsection{Quartic boundaries and an elementary cocycle calculation}

\begin{lemma}\label{lem:sl3-natural-h1-odd}
\status{Proved under stated hypotheses}
For every odd prime $p$, the natural three-dimensional module $Q$
and its dual satisfy
$H^1(\mathrm{SL}_3(\mathbf F_p),Q)=
H^1(\mathrm{SL}_3(\mathbf F_p),Q^*)=0$.
\end{lemma}

\begin{proof}
Let $c$ be a cocycle.  The diagonal sign subgroup of determinant one
has order four, prime to $p$, so subtracting a coboundary makes $c$
zero on it.  Each $t_{ij}(a)=1+aE_{ij}$ commutes with the sign matrix
which is $-1$ in positions $i,j$ and $+1$ in the remaining position
$k$.  The cocycle relation then forces
$c(t_{ij}(a))\in\langle e_k\rangle$.
Write $x=t_{12}(a)$, $y=t_{23}(b)$, with $a,b\ne0$, and
$c(x)=Ae_3$, $c(y)=Be_1$.  Since $[x,y]=t_{13}(ab)$, the cocycle
identity gives
\[
 c([x,y])=-bA e_2-abA e_1.
\]
Its left side belongs to $\langle e_2\rangle$, hence $A=0$.
Permuting the three indices proves vanishing on every elementary
transvection.  They generate $\mathrm{SL}_3(\mathbf F_p)$, so the
normalized cocycle is zero.  Applying the inverse-transpose
automorphism proves the dual assertion.
\end{proof}

For any chosen four of the $d$ lines, fix the other $r=d-4$ lines and
vary an ordinary quartic, transverse to the fixed lines and avoiding
their pairwise intersections.  Include ordinary singular quartics
whenever the total degree-$d$ curve stays in $W_d$.  This connected
open contains the original arrangement.  Mayer--Vietoris for the
quartic and the fixed lines gives
\begin{equation}\label{eq:odd-quartic-boundary}
 0\longrightarrow T\longrightarrow V\longrightarrow Q
 \longrightarrow0,\qquad \dim Q=3,
\end{equation}
where $Q$ is the quartic Hasse--Witt local system.  Its image is
$\mathrm{GL}_3(\mathbf F_p)$ by the canonical quartic moduli argument
and Achter--Pries, Proposition~4.8, already checked above.  The image
on $T$ is solvable: its only moving graph edges are the four
intersections with each fixed line, so it factors through a subgroup
of $(S_4)^r$.  The nodes between the fixed lines are fixed.

Let $B$ be the image of this boundary family and let $D_T$ be its
image on $T$.  The image of $B$ in
$\mathrm{GL}(Q)\times D_T$ is a subdirect product.  Goursat's lemma,
or the elementary description of a subdirect product as equality in
a common quotient, shows that it contains
$\mathrm{SL}(Q)\times1$: every common quotient is solvable, whereas
$\mathrm{SL}_3(\mathbf F_p)$ is perfect.  Let $B_0$ be the inverse
image of $\mathrm{SL}(Q)\times1$.  It is normal in $B$, fixes $T$
pointwise, and maps onto $\mathrm{SL}(Q)$.
Its kernel is an additive submodule
\begin{equation}\label{eq:odd-quartic-kernel}
 K\subseteq\operatorname{Hom}(Q,T)=Q^*\otimes T.
\end{equation}
The natural dual is absolutely simple, with endomorphism ring
$\mathbf F_p$.  A submodule of this direct sum of copies of $Q^*$
therefore has the form $K=Q^*\otimes T'$ for a subspace $T'\subseteq T$.
The subspace $T'$ is $B$-invariant because $K$ is normal in $B$.
If $K\ne0$, it contains rank-one maps $w\otimes\phi$ and hence
transvections $1+w\otimes\phi$ of $V$.  If $K=0$, $B_0$ is a
complement to $\operatorname{Hom}(Q,T)$ in the indicated semidirect
product.  Lemma~\ref{lem:sl3-natural-h1-odd}, applied to each copy of
$Q^*$, makes that complement conjugate to the block-diagonal
$\mathrm{SL}(Q)$.  In this case too it contains transvections of $V$.
For $d=4$ the statement simply recovers the full quartic group.

\subsubsection{The normal subgroup generated by transvections}

\begin{lemma}\label{lem:transvection-clifford-blocks}
\status{Proved under stated hypotheses}
Let an irreducible finite group $G\subseteq\mathrm{GL}(V)$ contain a
transvection, and let $N\triangleleft G$ be the subgroup generated by
all transvections in $G$.  Then $V$ is a direct sum of pairwise
nonisomorphic simple $N$-modules, and $G$ permutes these summands
transitively.  Every transvection has its image line in exactly one
summand and fixes all the others pointwise.
\end{lemma}

\begin{proof}
The sum of the simple $N$-submodules, $\operatorname{soc}_N(V)$, is
nonzero and $G$-stable, since $N$ is normal.  Irreducibility of $G$
therefore makes it all of $V$.  The canonical isotypic summands are
permuted transitively by $G$.  Write any one of them as $S^{\oplus e}$.
For an element of $N$, the rank of its difference from the identity
on this summand is a multiple of $e$.  Some transvection has nonzero
action on some summand, and its total rank is one, so $e=1$ there;
transitivity gives $e=1$ for all summands.  Thus the canonical
isotypic summands are themselves simple.  Each transvection is
block diagonal because it belongs to $N$, and its rank-one
difference is nonzero on only one block.  This proves the assertion.
(The same argument with the finite division endomorphism ring of a
block shows it is absolutely simple, but that strengthening is not
needed below.)
\end{proof}

For a four-element subset $I\subset\{1,\ldots,d\}$, let
$U_I\subset P$ be its three-dimensional sum-zero permutation
subspace, supported on $I$.  Because $4\ne0$ in $\mathbf F_p$,
$U_I$ embeds in $W$.  Set
\[
 S_I=\bigwedge^2 U_I\subset V.
\]
The subgroup $S_I^{\rm perm}\simeq S_4$ permuting these four lines
lies in the corresponding quartic boundary group $B$.  Restricted
to it, the decomposition of the permutation space gives
\begin{equation}\label{eq:s4-quartic-decomposition}
 W\simeq U_I\oplus\mathbf F_p^{d-4},\qquad
 V\simeq S_I\oplus U_I^{\oplus(d-4)}
             \oplus\mathbf F_p^{\binom{d-4}{2}}.
\end{equation}
Both $U_I$ and $S_I$ are simple.  For the former use the elementary
deleted-permutation argument with $p\nmid4$; for the latter use
Lemma~\ref{lem:odd-wedge-submodules} with four indices.  They are
nonisomorphic: a transposition has trace $1$ on $U_I$ and trace $-1$
on $\bigwedge^2U_I$, distinct in every odd characteristic.  Neither
is trivial.  Thus \eqref{eq:s4-quartic-decomposition} is a semisimple
decomposition even when $p=3$, and $S_I$ occurs with multiplicity one.
The quotient $Q$ in \eqref{eq:odd-quartic-boundary} is isomorphic to
$S_I$ as an $S_4$-module, and $T$ contains no constituent of this type.

\begin{proposition}\label{prop:odd-transvection-normal-irreducible}
\status{Proved under stated hypotheses}
The normal subgroup $N\triangleleft G$ generated by all transvections
is irreducible on $V$.
\end{proposition}

\begin{proof}
Use the canonical $N$-block decomposition of
Lemma~\ref{lem:transvection-clifford-blocks}.  We first prove that for
each quartet $I$ its subspace $S_I$ lies entirely in one block.

If the kernel $K$ of its boundary group $B_0$ is zero, the cocycle
calculation makes $B_0$ a conjugate block-diagonal $\mathrm{SL}_3$.
Its commutator subspace $P_I=[B_0,V]$ is therefore three dimensional,
simple as a $B_0$-module, and maps isomorphically to $Q$.  Normality
of $B_0$ in $B$ makes $P_I$ stable under the quartet's $S_4$.
The unique-copy statement in \eqref{eq:s4-quartic-decomposition}
forces $P_I=S_I$.  This copy of $\mathrm{SL}_3$ is generated by
rank-one transvections of $V$, so $B_0\subseteq N$.  Consequently
each $N$-block is $B_0$-stable and
$[B_0,V]=\bigoplus_j[B_0,V_j]$.  Simplicity of $P_I$ puts it in one
$N$-block.

If $K\ne0$, write it as $Q^*\otimes T'$, with $T'\ne0$, and fix
one nonzero $\phi\in Q^*$.  For every nonzero $w\in T'$, the
transvection $1+w\otimes\phi$ is supported on a single $N$-block.
The functional here means the pullback of $\phi$ along $V\to Q$.
It is nonzero on that block and zero on every other block.  Since
the same nonzero functional is used for all $w$, that block must
be the same block $L$ for every $w$.  Thus $T'\subset L$.
Now fix one nonzero $w$ and vary $\phi$.  Every $\phi\in Q^*$
vanishes on every other $N$-block.  It follows that $L\to Q$ is
surjective.  Since $T'$ is stable under the quartet's $S_4$, this
$S_4$ must preserve $L$: its translate is another $N$-block sharing
the nonzero subspace $T'$.  As an $S_4$-submodule of the semisimple
module \eqref{eq:s4-quartic-decomposition}, $L$ is semisimple.
Its surjection onto $Q\simeq S_I$ forces it to contain the unique
copy $S_I$ in $V$.  This proves the claim in both cases.

When two quartets have three common indices, their subspaces $S_I$
intersect in the nonzero line $\bigwedge^2U_{I\cap J}$.
The graph of four-element subsets joined by such overlaps is
connected.  Hence all $S_I$ belong to the same $N$-block $L$.
Their span is $A=\bigwedge^2D$: the spanning disjoint-difference
wedges in Lemma~\ref{lem:odd-wedge-submodules} are supported on
quartets.  Thus $A\subset L$.

The subspace $A$ is $S_d$-stable.  Every translate of $L$ by $S_d$
is an $N$-block containing $A$, hence is $L$ itself.  Thus $S_d$
preserves $L$ and the sum $L'$ of the other blocks.  If $L'\ne0$,
Lemma~\ref{lem:odd-wedge-submodules} forces $A\subset L'$, which
contradicts $A\subset L$ and $L\cap L'=0$.  There are no other
blocks, and $N$ is irreducible.  The base case $d=4$ was already
settled by the quartic monodromy theorem.
\end{proof}

\subsubsection{Classification, exclusion of the alternating form, and determinant}

McLaughlin's transvection theorem, in the precise form of
Kantor--Seress (2001), Theorem~3.6, p.~22 \cite{KantorSeress2001},
states that for $n\geq3$
an irreducible subgroup of $\mathrm{SL}_n(\mathbf F_q)$ generated by
full transvection groups is $\mathrm{SL}_n(\mathbf F_q)$ or, when
$n$ is even, $\mathrm{Sp}_n(\mathbf F_q)$; the additional alternatives
in that theorem occur only when $q=2$.  Here $q=p$ is prime, and
the powers of $1+v\phi$ are precisely $1+a v\phi$, $a\in\mathbf F_p$.
Thus every generating transvection supplies its full transvection
group.  The theorem applies to $N$, which has determinant one and
is irreducible by Proposition~\ref{prop:odd-transvection-normal-irreducible}.
It leaves $N=\mathrm{SL}_g(\mathbf F_p)$ or $N=\mathrm{Sp}_g(\mathbf F_p)$.

The symplectic possibility is impossible.  If $N$ is symplectic,
normality puts $G$ in its group of similitudes: the invariant
alternating form is unique up to scalar.  For any quartet, the
multiplier on $B_0$ is trivial.  Indeed it is trivial on its
$p$-group kernel $K$ and on its perfect quotient
$\mathrm{SL}_3(\mathbf F_p)$.  Therefore $B_0$ preserves the
nondegenerate alternating form.  Since it fixes $T$ pointwise and
has fixed-point-free natural quotient $Q$, its fixed space is
$V^{B_0}=T$.  For a group preserving a nondegenerate bilinear form,
\[
 [B_0,V]=(V^{B_0})^\perp=T^\perp.
\]
The latter space has dimension three.  Its image in $Q$ is
$[B_0,Q]=Q$, so it maps isomorphically to $Q$ and meets $T$ trivially.
The alternating form is consequently nondegenerate on this
three-dimensional space $T^\perp$, a contradiction.  We have proved
\begin{equation}\label{eq:odd-plane-sl-contained}
 \mathrm{SL}_g(\mathbf F_p)\subseteq G.
\end{equation}

It remains to show that the geometric determinant is all of
$\mathbf F_p^\times$.  Take a Legendre elliptic cubic $E_t$, near a
supersingular parameter $t_0$, and add $d-3$ fixed lines chosen in
general position relative to $E_{t_0}$.  The Hasse polynomial of
the Legendre family has a simple root at $t_0\notin\{0,1\}$, as
proved in the elliptic calculation in
\S\ref{subsec:plane-residues}.  Over the strict local
trait at $t_0$, all graph contributions have bijective Frobenius
and determinant a unit.  The normalization sequence therefore
makes the total Hasse determinant a unit times the elliptic Hasse
invariant.  It has valuation exactly one at $t_0$.
On the punctured trait the total curve is ordinary, hence belongs
to $W_d$.  The determinant Lang torsor is
$z^{p-1}=h^{-1}$ by \eqref{eq:determinant-kummer}.  A valuation-one
$h$ gives tame inertia image of exact order $p-1$.  This is
geometric inertia, since the residue field is algebraically closed.
Thus $\det G=\mathbf F_p^\times$.  Combining this with
\eqref{eq:odd-plane-sl-contained} proves the theorem.

\begin{theorem}\label{thm:plane-odd-all-groups}
\status{Proved under stated hypotheses}
Let $p$ be odd and let $U_d$ be the ordinary smooth plane-curve locus
of degree $d$ over any perfect field $k$ of characteristic $p$.  At a
geometric generic point put
$V=H^1_{\mathrm{et}}(C_{\bar\eta},\mathbf F_p)$, of dimension
$g=(d-1)(d-2)/2$.  For every $d\geq3$, both the geometric and the
arithmetic monodromy images on $V$ are
$\mathrm{GL}_g(\mathbf F_p)$, with full determinant character.
For $d=1,2$ both images are trivial on the zero vector space.
For positive rank the representation is irreducible.
\end{theorem}

\begin{proof}
The preceding argument proves the geometric assertion for $d\geq4$
over $\overline{\mathbf F}_p$.  The connected finite etale frame cover
which it identifies is geometrically connected and remains connected
after every algebraically closed field extension.  Since the family
and its Frobenius local system are defined over $\mathbf F_p$, this
gives the same geometric image over the algebraic closure of any
perfect field of characteristic $p$.
The arithmetic image contains that image and is a subgroup of the
same finite general linear group, so it is equal to it.  The cubic
case follows from the same Legendre determinant calculation, since
its rank is one.  Degrees one and two have rank zero.
\end{proof}

% /root/monodromy, 19 November 2026.
% Independently audited in full by /root, including the determinant seeds.
\subsection{Hypersurfaces in projective space in odd characteristic}
\label{subsec:projective-space-odd-monodromy}
\status{Proved; independently audited by /root, including both determinant seeds.}

Fix an odd prime $p$ and an algebraically closed field $k$ of
characteristic $p$.  For degree $d$ hypersurfaces in $\mathbf P^n$,
let $W_{n,d}$ be the Hasse--Witt invertibility open in the full
projective space of equations, and let
\[
 h_{n,d}=\binom{d-1}{n},\qquad
 V_{n,d}=\left(H^{n-1}(\mathcal O_Y),F=1\right).
\]
The binomial coefficient is zero for $d\leq n$.
The full coherent bundle has this constant rank for all nonzero
equations, including singular equations.  A general-position union
of $d$ hyperplanes belongs to $W_{n,d}$, so this open is nonempty.
The desired smooth ordinary locus is a dense open in $W_{n,d}$;
passing to it does not change the finite monodromy image.
All the comparison families below therefore take place inside
one explicit constant-rank Lang local system, except for a
valuation-one Hasse degeneration used solely to compute its tame
determinant inertia.

\subsubsection{A connectedness lemma for linear projection}

\begin{lemma}\label{lem:linear-projection-monodromy}
\status{Proved under stated hypotheses}
Let $U$ be a nonempty open in $\mathbf P^N_k$, and let
$T\to U$ be a connected finite etale cover.  Let
$f:\mathbf P^N\dashrightarrow\mathbf P^M$ be a linear projection
whose center $Z$ meets $U$.  The restriction of $T$ to the
geometric generic fibre of $f|_{U\setminus Z}$ is connected.
Consequently a finite local system on $U$ retains its full
monodromy after restriction to that geometric generic fibre.
\end{lemma}

\begin{proof}
The open $U$ is smooth and integral, so $T$ is smooth and integral.
Let $B=\operatorname{Bl}_Z(\mathbf P^N)$ and put
$T'=T\times_U B_U$, where $B_U$ is the inverse image of $U$.
This is the blowup of $T$ in the smooth etale inverse image of
$Z\cap U$; in particular it is smooth and integral.  The resolved
projection gives a morphism $T'\to\mathbf P^M$.
Choose $z\in Z\cap U$.  The exceptional fibre above $z$ is
$E_z\simeq\mathbf P^M$, mapping isomorphically to the target.
Its pullback to $T'$ is the disjoint union of copies of $E_z$,
since it is the product of $E_z$ with the finite reduced set $T_z$.

Put $K=k(\mathbf P^M)$, and let $L$ be the algebraic closure of $K$
inside $k(T')$.  It is a finite extension of $K$.  Let $C$ be the
normalization of $\mathbf P^M$ in $L$; it is finite and integral.
There is a morphism $T'\to C$ over $\mathbf P^M$.  Indeed, on an
affine open of the target, each element of its integral closure
in $L$ is integral over the local ring at every point of the
normal scheme $T'$ lying above this open.  It therefore belongs
to each such local ring and defines a regular function there.
These morphisms glue.  One copy of $E_z$ in $T'$ consequently
gives a section $\mathbf P^M\to C$ of the finite map.
Its image is closed, has dimension $M$, and lies in the integral
scheme $C$ of dimension $M$; it is all of $C$.  Thus $C\simeq
\mathbf P^M$ and $L=K$.

On the nonempty open $T|_{U\setminus Z}$, the map to $\mathbf P^M$
is a composition of a finite etale map and a smooth linear
projection.  Hence $k(T')/K$ is separably generated.  The equality
$L=K$ implies that this extension is regular, so its geometric
generic fibre is integral.  Removing the exceptional divisor
does not alter that assertion.  For the local-system consequence,
apply the proved assertion to a connected component of its finite
frame cover, which is a torsor for its full monodromy image.
\end{proof}

\begin{corollary}\label{cor:restriction-independent-monodromy}
\status{Proved under stated hypotheses}
For $e\geq n+1$, fix a hyperplane $H\subset\mathbf P^n$.
The local system $V_{n,e}$ retains its full geometric monodromy
when the restriction of the equation to $H$ is prescribed to be
geometric generic.  On the parameter space where both the
equation and its restriction have invertible Hasse--Witt maps,
the joint monodromy image on $V_{n,e}$ and $V_{n-1,e}$ is the
product of their two full images.
\end{corollary}

\begin{proof}
Restriction of equations is a linear projection.  Its center
consists of equations divisible by the fixed equation of $H$.
It meets $W_{n,e}$: take $H$ and $e-1$ other general-position
hyperplanes.  Lemma~\ref{lem:linear-projection-monodromy} shows
that the joint image contains the first full group times $1$.
Its projection to the second group is surjective, since the
restriction map has geometrically integral generic fibre, an
open of an affine space.  The joint image is therefore the
product.  Removing further nonempty opens does not change
these finite images.
\end{proof}

\subsubsection{The arrangement module and independent quadratic reflections}

Take a general-position union of $d$ hyperplanes as base point.
Put $P=\mathbf F_p^d$, $u=\sum_i e_i$, $W=P/\langle u\rangle$,
$U=\ker(\sigma:P\to\mathbf F_p)$ with $\sigma(x)=\sum_i x_i$,
and $D=\operatorname{im}(U\to W)$.  The Cech resolution of the
union, whose strata are projective spaces, identifies its top
coherent cohomology and its Lang fibre with
\begin{equation}\label{eq:pn-arrangement-module}
 V_{n,d}=\bigwedge^nP/(u\wedge\bigwedge^{n-1}P)
       =\bigwedge^nW.
\end{equation}
The connected ordered-hyperplane cover of the unordered
arrangement locus gives the full permutation action of $S_d$
on this space.  Put $A=\bigwedge^nD$.

\begin{lemma}\label{lem:odd-general-hook-module}
\status{Proved under stated hypotheses}
Assume $1\leq n\leq d-2$.  If $p\nmid d$, then
$\bigwedge^nW$ is simple as an $S_d$-module.  If $p\mid d$, then
$A=\bigwedge^nD$ is its unique proper nonzero submodule; both
$A$ and the quotient $\bigwedge^{n-1}D$ are simple.
\end{lemma}

\begin{proof}
First suppose $2n\leq d$.  For $n$ pairwise disjoint transpositions
$\tau_j=(a_jb_j)$, the group-algebra operator
$\prod_j(1-\tau_j)$ has rank one on $\bigwedge^nW$.  Its image
is the wedge of the $n$ differences $e_{a_j}-e_{b_j}$, and its
functional is the wedge of the corresponding coordinate
differences.  This follows by simultaneously decomposing $W$
into the $n$ negative eigenlines and their common positive
eigenspace; $2$ is invertible.

These matching wedges span $\bigwedge^nU$.  To check this without
a modular representation-theory citation, use the difference
basis $f_i=e_i-e_o$ and the identity, for four distinct indices,
\[
 2f_a\wedge f_b
 =(e_a-e_b)\wedge(e_c-e_o)
  +(e_a-e_c)\wedge(e_b-e_o)
  +(e_a-e_o)\wedge(e_b-e_c).
\]
Beginning with a wedge of $n$ differences sharing $o$, repeatedly
introduce one unused index $c$.  Each application leaves one
fewer difference sharing $o$ and produces one disjoint pair.
There are $n-1$ available new indices because $d\geq2n$.
The resulting wedges are perfect matchings on $2n$ indices.
The same proof in the dual permutation space shows that the
matching functionals span $\bigwedge^nW^*$.
Thus every nonzero submodule contains a matching image and hence
all of $A$.  This proves simplicity of $A$, and proves simplicity
of the whole space when $p\nmid d$.

When $p\mid d$, the dot product gives a nondegenerate invariant
form on $D=U/\langle u\rangle$.  The same matching argument for
$n-1$ shows that $\bigwedge^{n-1}D$ is simple: matching images
span it, and their functionals span its dual under this form.
Contraction with the sum functional gives an exact sequence
\[
 0\longrightarrow\bigwedge^nD\longrightarrow\bigwedge^nW
 \longrightarrow\bigwedge^{n-1}D\longrightarrow0.
 \tag{*}
\]
Every nonzero submodule contains its kernel, and a submodule
which also surjects to the simple quotient is the whole space.
This proves the unique-submodule assertion in the stated range.

For $2n>d$, put $r=d-1-n$, so $1\leq r<d/2$.
The natural exterior pairing gives
$(\bigwedge^nW)^*\simeq\bigwedge^rW\otimes\operatorname{sgn}$.
If $p\nmid d$, the already proved simplicity for $r$ gives the
claim.  If $p\mid d$, duality preserves the property of having
exactly one proper nonzero submodule.  The displayed subspace
$\bigwedge^nD$ is nonzero and proper, so it must be that unique
submodule.  Its simplicity and that of the quotient follow
from the corresponding two simple factors for $r$.
\end{proof}

\begin{lemma}\label{lem:quadric-reflections}
\status{Proved under stated hypotheses}
For $n\geq2$ and $d\geq n+2$, the geometric group on
$V_{n,d}$ contains commuting rank-one reflections
$R_I$, indexed by $(n-1)$-subsets
$I\subset\{3,\ldots,d\}$.  At the arrangement base point,
their negative eigenlines are
\[
 \left\langle (e_1-e_2)\wedge\bigwedge_{i\in I}e_i\right\rangle,
\]
and their product is the permutation transposition $(12)$.
Here a rank-one reflection has eigenvalue $-1$ once and
eigenvalue $1$ on a codimension-one subspace.
\end{lemma}

\begin{proof}
Fix the last $d-2$ hyperplanes and replace the first two by a
variable quadric $Q$.  Require every positive-dimensional
quadric stratum to be connected with no higher coherent
cohomology, and every intersection with a line
$L_I=\bigcap_{i\in I}H_i$ to consist of two distinct points.
This is a nonempty open containing a smooth general quadric
and the reducible quadric $H_1H_2$.  A positive-dimensional
quadric hypersurface has $H^0(\mathcal O)=k$ and
$H^{>0}(\mathcal O)=0$, so the Cech complex computes all the
cohomology from its $H^0$ terms.  Only its top term has
nonconstant point sets: the two points on each $L_I$.
Swapping one pair acts on cohomology as a rank-one reflection,
since every coboundary has equal values on that pair.

All these swaps occur independently.  The restriction from
quadratic forms to $L_I\simeq\mathbf P^1$ is surjective.  The
discriminant of a binary quadratic is an irreducible polynomial
in its three coefficients in odd characteristic.  Its pullback
to the space of all quadratic forms is still irreducible, and
the discriminants for distinct lines are nonassociate.  To see
the last assertion, choose a point of $L_I\setminus L_J$ and
two hyperplanes through that point which meet $L_J$ in distinct
points.  Their product is tangent to $L_I$ and transverse to
$L_J$.  Their classes in the rational function field
modulo squares are therefore independent, by valuations at
these distinct prime divisors.  The ordered-pair cover has
Galois group $(\mathbf Z/2)^{\binom{d-2}{n-1}}$.
Removing the indicated open conditions leaves it connected.

The pair differences at $Q=H_1H_2$ are precisely the displayed
cochains, using the refinement of the cover with component $Q$
by its two hyperplanes.  Interchanging these two hyperplanes
interchanges every pair and fixes every other Cech term.
The transposition $(12)$ is thus the product of all the $R_I$.
Every fibre here has bijective Frobenius on the full top
coherent cohomology, because all terms are finite sets with
ordinary constant Frobenius.  The entire comparison is inside
$W_{n,d}$.
\end{proof}

\begin{proposition}\label{prop:pn-odd-irreducible}
\status{Proved under stated hypotheses}
For $n\geq2$ and $d\geq n+2$, the geometric monodromy group
$G_{n,d}$ is irreducible.
\end{proposition}

\begin{proof}
Only $p\mid d$ needs consideration, by
Lemma~\ref{lem:odd-general-hook-module}.  In that case any
proper invariant subspace must be $A=\bigwedge^nD$.
The negative eigenline of every $R_I$ lies outside $A$:
contraction with $\sigma$ sends its displayed generator to
\[
 -(e_1-e_2)\wedge\iota_\sigma
                  \left(\bigwedge_{i\in I}e_i\right)\ne0.
 \tag{**}
\]
Indeed the involved $n$ vectors remain linearly independent
modulo $u$: at least one index is outside $\{1,2\}\cup I$.
If $A$ were invariant, each $R_I$ would consequently fix $A$
pointwise.  Their product $(12)$ would do so too.  But on $D$
this transposition has a one-dimensional negative eigenspace
and a positive eigenspace of dimension $d-3$; on
$\bigwedge^nD$ its negative eigenspace has dimension
$\binom{d-3}{n-1}>0$.  This contradicts pointwise fixation.
\end{proof}


\subsubsection{The degree $n+2$ seed}

We use the following elementary observation about the Cech
computations, also in the determinant calculation below.  A family
consisting of coordinate hyperplanes and one further hypersurface
has the usual exact alternating resolution by its intersections
provided the latter equation restricts to a nonzero equation on
every coordinate stratum.  Locally its vanishing equation and any
of the coordinate equations form a regular sequence.  The
alternating union complex for a regular sequence is exact, by
induction using $(f)\cap(g)=(fg)$ and its versions modulo the
other equations.  Thus smoothness of the further hypersurface is
unnecessary.  Higher coherent cohomology of any hypersurface
stratum is given by its usual projective-space exact sequence.

\begin{proposition}\label{prop:pn-odd-seed-sl}
\status{Proved under stated hypotheses}
For $n\geq3$ and $d=n+2$, the group $G_{n,d}$ contains
$\mathrm{SL}_{n+1}(\mathbf F_p)$.
\end{proposition}

\begin{proof}
Let $N$ be the normal subgroup of $G_{n,d}$ generated by all its
transvections.  Replace any four of the arrangement hyperplanes,
indexed by $I$, by a varying quartic $Q$, and fix the other
$n-2$ hyperplanes.  The Cech spectral sequence has precisely
the following positive-cohomology terms contributing to total
degree $n-1$:
the quartic plane curve obtained by intersecting all the fixed
hyperplanes, of rank three, and the $n-2$ quartic surface strata,
each of rank one.  Higher-dimensional quartic strata have no
positive coherent cohomology.  The $H^0$ row has no positive
cohomology because the nerve is a full simplex.  There are no
nonzero incoming or outgoing differentials for the displayed
terms.  We consequently obtain an exact sequence of Frobenius
modules, and of Lang local systems on the invertibility open,
\[
 0\longrightarrow Q_I\longrightarrow V_{n,n+2}
      \longrightarrow T_I\longrightarrow0,
 \qquad \dim Q_I=3,\quad\dim T_I=n-2.
 \tag{1}
\]
Here $T_I$ is the direct sum of the rank-one surface local systems.
The open where all these terms are ordinary contains the original
arrangement.  Its image on $Q_I$ is $\mathrm{GL}_3(\mathbf F_p)$:
restriction of quartics to the fixed plane is a linear projection
with affine generic fibre, and the ordinary plane-quartic group
is already proved.  The inverse image of its connected frame
cover remains integral; imposing the other nonempty open
conditions preserves connectedness.  The image on $T_I$ is
abelian, being a subgroup of $(\mathbf F_p^\times)^{n-2}$.

Apply the quartic-boundary argument of the preceding plane
section to the dual of (1).  Goursat's lemma gives a normal
subgroup $B_0$ fixing $T_I^*$ pointwise and mapping onto
$\mathrm{SL}(Q_I^*)$.  Its kernel is
$Q_I\otimes T'$ for a subspace $T'\subseteq T_I^*$.
If the kernel is nonzero, it contains rank-one transvections.
If it is zero, Lemma~\ref{lem:sl3-natural-h1-odd} splits the
extension and again gives transvections.  Thus $N$ is nontrivial.
By Proposition~\ref{prop:pn-odd-irreducible} and
Lemma~\ref{lem:transvection-clifford-blocks}, $V_{n,n+2}^*$ is a
direct sum of canonical simple $N$-blocks permuted transitively
by $G_{n,n+2}$.

In the arrangement fibre there is the $S_d$-equivariant
identification
\[
 V_{n,n+2}^* \simeq W\otimes\operatorname{sgn},
 \qquad d=n+2.
 \tag{2}
\]
For a quartet $I$, let $U_I$ denote its three-dimensional
sum-zero subspace in $P$.  Since $4$ is invertible, it embeds
in $W$.  The restriction of (2) to its $S_4$ is the direct sum
of the unique simple module
$S_I=U_I\otimes\operatorname{sgn}$, of dimension three, and
$d-4$ copies of the sign character.  This is an actual
semisimple decomposition in every odd characteristic.
The quotient $Q_I^*$ in the dual of (1) is this unique
three-dimensional constituent.

The two cases in the proof of
Proposition~\ref{prop:odd-transvection-normal-irreducible} now
apply verbatim, with this specified decomposition: if the kernel
is zero, the commutator space of the split $\mathrm{SL}_3$ is
$S_I$ and belongs to one $N$-block; if it is nonzero, the
rank-one maps with a common quotient functional put $T'$ in one
block, and varying the functional makes that block surject to
$Q_I^*$.  The block is $S_4$-stable because $T'$ is stable, so
it contains the unique constituent $S_I$.

Quartets with three common indices give $S_I$'s with nonzero
intersection.  Their connected overlap graph puts all these
subspaces in one block $L$.  They span
$D\otimes\operatorname{sgn}$.  Every translate of $L$ by $S_d$
contains this same nonzero subspace and must equal $L$.
Thus $S_d$ also preserves the sum of the remaining blocks.
If that sum were nonzero, the degree-one case of
Lemma~\ref{lem:odd-general-hook-module}, with the sign twist,
would force it to contain $D\otimes\operatorname{sgn}$ too.
This contradicts the direct sum.  Therefore $N$ is irreducible.

Kantor--Seress, Theorem~3.6 \cite{KantorSeress2001}, applies:
over the prime field, powers of each transvection supply its
full transvection group.  It leaves
$N=\mathrm{SL}_{n+1}(\mathbf F_p)$ or a symplectic group.
The latter is excluded by
Lemma~\ref{lem:pn-reflection-excludes-symplectic} below.
\end{proof}

\begin{lemma}\label{lem:pn-reflection-excludes-symplectic}
\status{Proved under stated hypotheses}
If an irreducible normal subgroup $N\triangleleft G$ is the
natural symplectic group in dimension $h\geq3$ over an odd prime
field, then $G$ cannot contain a rank-one reflection.
\end{lemma}

\begin{proof}
The space of $N$-invariant alternating forms is one dimensional,
so normality puts $G$ in the symplectic similitude group.
Let $R$ be such a reflection, with positive eigenspace of
dimension $h-1$ and negative eigenline.  If its multiplier
were different from one, its positive eigenspace would be
totally isotropic, contradicting $h-1>h/2$.
Its multiplier is therefore one.  The negative eigenline is
then orthogonal to the positive eigenspace, since $p\ne2$,
and to itself because the form is alternating.  It is in the
radical of the form, a contradiction.
\end{proof}

\subsubsection{Induction in the degree}

\begin{proposition}\label{prop:pn-odd-degree-induction}
\status{Proved under stated hypotheses}
Suppose the groups for degree $d-1$ in $\mathbf P^n$ and
$\mathbf P^{n-1}$ contain their full special linear groups,
where $n\geq3$ and $d\geq n+3$.  Then
$G_{n,d}$ contains $\mathrm{SL}_{h_{n,d}}(\mathbf F_p)$.
If in addition the first of these two groups has full
determinant, then so does $G_{n,d}$.
\end{proposition}

\begin{proof}
Use the boundary $Y=A\cup H$, with $A$ a degree-$(d-1)$
hypersurface and $H$ a fixed hyperplane.  The usual exact
sequence of the union gives, on the simultaneous ordinary open,
\begin{equation}\label{eq:pn-two-piece-boundary}
 0\longrightarrow S\longrightarrow V_{n,d}
       \longrightarrow Q\longrightarrow0,
 \qquad
 S=V_{n-1,d-1},\quad Q=V_{n,d-1}.
\end{equation}
There are no adjacent positive coherent groups, since a
hypersurface in projective space has no intermediate coherent
cohomology.  Put
\[
 r=\dim Q=\binom{d-2}{n},\qquad
 s=\dim S=\binom{d-2}{n-1}.
\]
Both are at least three in the present range.
Corollary~\ref{cor:restriction-independent-monodromy} shows
that the graded image is the product of the full groups.
Take the inverse image $B_0$ of
$L=\mathrm{SL}(S)\times\mathrm{SL}(Q)$.
Its kernel $K$ is an additive $L$-submodule of
$M=\operatorname{Hom}(Q,S)$.  This outer tensor product of
the two absolutely simple natural modules is absolutely simple,
so $K=0$ or $K=M$.

If $K=M$, its rank-one elements are transvections of $V$.
Let $N$ be their full normal closure in $G_{n,d}$ and use
the canonical $N$-blocks, justified by irreducibility of $G$.
For every $w\in S$ and $\phi\in Q^*$, the operator
$1+w\otimes\phi$ is a transvection when both are nonzero.
Fixing $\phi$ puts all of $S$ in one block; fixing a nonzero
$w$ and varying $\phi$ makes that block surject to $Q$.
It contains $S$ and surjects to $V/S$, so it is all of $V$.
Thus $N$ is irreducible.

If $K=0$, the complement is conjugate to the split block
diagonal group $L$.  Here is the needed cohomology computation,
which requires no vanishing theorem for either natural module:
\[
 H^1(\mathrm{SL}(S)\times\mathrm{SL}(Q),S\otimes Q^*)=0.
\]
The first factor has no invariants in $S$.  In the low-degree
Hochschild--Serre sequence the remaining term is
$H^1(\mathrm{SL}(S),S)\otimes Q^*$, whose invariants under the
second factor are zero.  This proves the displayed vanishing
and hence the conjugacy of complements.
The split $L$ is generated by rank-one transvections on $V$,
so is contained in $N$.  Every canonical $N$-block is therefore
an $L$-submodule of the semisimple, multiplicity-free module
$S\oplus Q$.  If there is more than one block, they are exactly
these two summands.  Transitivity of $G$ on the blocks forces
$r=s$, equivalently $d=2n+1$.

Even this last possibility is impossible.  A rank-one reflection
cannot interchange two subspaces of common dimension $r>1$:
its difference from the identity would have rank at least $r$
on either one.  Thus every reflection of
Lemma~\ref{lem:quadric-reflections} preserves both blocks.
Their product $(12)$ does so too.  Conjugating within the
arrangement $S_d$ shows that every transposition, hence all
of $S_d$, preserves both blocks.  This contradicts
Lemma~\ref{lem:odd-general-hook-module}: a simple module, or a
module with a unique proper nonzero submodule, has no such
nontrivial direct sum decomposition.  We have proved that
$N$ is irreducible also when $K=0$.

The transvection classification again gives $N=\mathrm{SL}(V)$
or a symplectic group.  The latter is excluded by
Lemma~\ref{lem:pn-reflection-excludes-symplectic}, using the
quadric reflection already present in $G$.
Finally the full product assertion for the graded image
contains the old group on $Q$ times $1$ on $S$.  Its lifts
have the same determinant as their action on $Q$, because
the extension kernel is unipotent.  Full determinant for
that old group therefore implies full determinant for $G$.
\end{proof}

\subsubsection{Determinants in the two initial degrees}

\begin{lemma}\label{lem:elliptic-fixed-three-points}
\status{Proved under stated hypotheses}
There is a family of plane cubics over a strictly henselian
smooth trait over $k$ whose special fibre is smooth supersingular,
whose Hasse invariant has valuation one, and whose restriction
to a fixed line is the constant squarefree binary cubic
$g=x_0x_1(x_0-x_1)$.  There are also smooth ordinary plane
cubics with this same restriction.
\end{lemma}

\begin{proof}
Take the Legendre family at a simple supersingular zero of its
Hasse polynomial, as proved in the earlier elliptic calculation.
Choose a line transverse to the special fibre and identify it
with $\mathbf P^1$.  Transversality persists on the trait, and
the three intersection points form a finite etale scheme.
They split after strict henselization.  The unique projective
linear transformation on the line taking their chosen order
to $0,1,\infty$ varies regularly with the trait parameter.
Lift this transformation to the ambient projective plane and
rescale the cubic equation by a unit.  Its restriction to the
line is then exactly $g$.  Isomorphisms and unit changes of
cohomology trivialization multiply the Hasse invariant by a
unit, so its valuation is still one.  The same construction
with any ordinary elliptic plane cubic gives the last assertion.
\end{proof}

\begin{proposition}\label{prop:pn-seed-full-determinant}
\status{Proved under stated hypotheses}
For every $n\geq2$, the geometric determinant images in degrees
$n+1$ and $n+2$ are the full group $\mathbf F_p^\times$.
\end{proposition}

\begin{proof}
For degree $n+1$, fix $n-2$ coordinate hyperplanes and let their
intersection be a plane.  Lift the cubic family of
Lemma~\ref{lem:elliptic-fixed-three-points} from this plane to
$\mathbf P^n$, for instance by its same equation independent
of the other coordinates, and take its union with the fixed
hyperplanes.  Their intersections are proper.  All
positive-dimensional cubic strata of ambient dimension at
least three have no positive coherent cohomology; the only
nonzero term in total degree $n-1$ is the elliptic curve term.
The $H^0$ nerve is a full simplex and has no higher cohomology.
Thus the rank-one total Hasse invariant is a unit times the
elliptic one and has valuation exactly one.
For $n=2$ this is simply the elliptic family itself.

For degree $n+2$, use coordinates
$x_0,x_1,y_1,\ldots,y_{n-1}$ and hyperplanes $H_i=(y_i=0)$.
Their common line is $\ell=\mathbf P(x_0,x_1)$.
Let $f_1(x_0,x_1,y_1)$ be the family in the lemma, and for
$i>1$ let $f_i(x_0,x_1,y_i)$ be a fixed ordinary plane cubic,
all restricting to $g$ on $\ell$.  Define explicitly
\[
 q=g+\sum_{i=1}^{n-1}(f_i-g),\qquad
 Y=(q\,y_1\cdots y_{n-1}=0)\subset\mathbf P^n.
 \tag{3}
\]
Restriction to the plane where all $y_j$, $j\ne i$, vanish
is exactly $f_i$.  Every coordinate stratum has a nonzero
restricted cubic equation, because its restriction to $\ell$
is $g$.  The alternating union resolution described above
is therefore exact, even if the higher-dimensional cubic
strata happen to be singular.

Those higher-dimensional cubic strata have no positive
coherent cohomology, and every positive-dimensional stratum
has $H^0(\mathcal O)=k$.  The common zero-dimensional stratum
consists of the three constant, distinct roots of $g$.
The $H^0$ row is the full-simplex complex except that its last
copy of $k$ is replaced by $k^3$; it contributes an ordinary
rank-two piece in degree $n-1$.  In the positive-cohomology
rows the $n-1$ elliptic curves contribute their rank-one
cohomology, again in total degree $n-1$.  There are no
possible nonzero differentials.  This gives a Frobenius-stable
filtration on the total cohomology with ordinary graph piece
of rank two and the $n-1$ elliptic pieces.  Its total rank is
$2+(n-1)=n+1=\binom{n+1}{n}$.
The Hasse determinant is a unit times the product of their
Hasse invariants.  Exactly $f_1$ has valuation one and all the
other factors are units, so the total determinant has
valuation one.

In both constructions the punctured trait lies in $W_{n,d}$.
The determinant Lang torsor is
$z^{p-1}=h^{-1}$, as in \eqref{eq:determinant-kummer}.
A valuation-one $h$ gives tame geometric inertia of order
$p-1$.  This proves the assertion.  The special fibre is
allowed to have smaller Hasse--Witt rank: only its valuation
is used, and no stable local system is extended across it.
\end{proof}

\begin{theorem}\label{thm:pn-odd-full-groups}
\status{Proved under stated hypotheses}
Let $p$ be odd, let $k$ be any perfect field of characteristic
$p$, and let $n\geq2$.  On the ordinary smooth degree-$d$
hypersurface locus in $\mathbf P^n_k$, the geometric and
arithmetic monodromy images on
$H^{n-1}_{\mathrm{et}}(Y_{\bar\eta},\mathbf F_p)$ are both
\[
 \mathrm{GL}_{\binom{d-1}{n}}(\mathbf F_p).
\]
For $d\leq n$ the vector space is zero and the image is trivial.
For $d=n+1$ it has rank one and full determinant.
For every $d\geq n+2$ it is irreducible with full determinant.
\end{theorem}

\begin{proof}
Dimension two is the already proved plane theorem.
Induct on $n\geq3$.  The zero-rank degrees need no argument,
and Proposition~\ref{prop:pn-seed-full-determinant} proves the
rank-one degree $n+1$.  In degree $n+2$,
Proposition~\ref{prop:pn-odd-seed-sl} and the determinant
proposition give the full general linear group.
Now induct on $d\geq n+3$ using
Proposition~\ref{prop:pn-odd-degree-induction}.
The two smaller groups are known by the dimension and degree
inductions, and the old determinant propagates as proved there.
This proves the geometric image over $\overline{\mathbf F}_p$.
The connected finite frame cover remains geometrically
connected under algebraically closed extension; since the
family is defined over $\mathbf F_p$, the result holds over
the algebraic closure of any perfect $k$.  The arithmetic
image contains the full geometric general linear group, so
is equal to it.
\end{proof}

\begin{corollary}\label{cor:pn-odd-all-residues-sharp}
\status{Proved under stated hypotheses}
Let $k$ be finite of odd characteristic $p$ and $n\geq2$.
For every $d\geq n+2$, every nonempty open in the smooth
degree-$d$ hypersurface locus of $\mathbf P^n_k$ contains,
over suitable finite extensions of $k$, hypersurfaces with
each of the $p$ point-count residues.
For $d\leq n$ every smooth hypersurface has residue $1$.
For $d=n+1$ the ordinary open misses residue $1$.
Thus $d=n+2$ is the exact threshold for the assertion involving
every nonempty open.
\end{corollary}

\begin{proof}
The finite-field coherent trace formula gives
\[
 \#Y(k')\equiv 1+(-1)^{n-1}
       \operatorname{Tr}(F_{k'}\mid H^{n-1}(Y,\mathcal O_Y)_{\rm ss})
       \pmod p,
\]
because a projective hypersurface has no other positive
coherent cohomology.  In the zero-rank range the trace is zero.
In rank one on the ordinary open it is nonzero.
In rank $h\geq2$, traces of invertible matrices take every
value in $\mathbf F_p$: the block
$\begin{pmatrix}0&-1\\1&a\end{pmatrix}$ has determinant one
and trace $a$, and adjoining an identity block allows any
specified total trace.  By the theorem the finite geometric
and arithmetic image is the full group.  The finite
Chebotarev theorem already cited in the plane calculation
therefore realizes each such conjugacy class in every
nonempty open after suitable finite extension of $k$.
The intersection of that open with the ordinary locus is
nonempty because the latter is dense.
\end{proof}

% /root/monodromy, 19 November 2026.
% Elementary all-dimension characteristic-two computation; root independently audited.
\subsection{Hypersurfaces in characteristic two: an elementary exact group calculation}
\label{subsec:pn-char2-exact}
\status{Proved; independently audited.}

Fix an algebraically closed field $k$ of characteristic two,
$n\geq2$, and $d\geq n+2$.  Put $m=d-1$ and
$h=\binom{m}{n}$.  Work on the full Hasse--Witt invertibility
open $W_{n,d}$ as in the preceding section.  Its Lang fibre
at a general-position hyperplane arrangement is
$V=\bigwedge^n W$, where
$W=\mathbf F_2^d/\langle e_0+\cdots+e_m\rangle$.
Use $e_1,\ldots,e_m$ as a basis of $W$, so
$e_0=e_1+\cdots+e_m$.
Let $b_A=\bigwedge_{i\in A}e_i$, indexed by the $n$-subsets
$A\subset\{1,\ldots,m\}$, be the exterior basis of $V$.
Signs disappear in this characteristic.

\begin{lemma}\label{lem:char2-independent-quadric-swaps}
\status{Proved under stated hypotheses}
The geometric monodromy contains, for each
$(n-1)$-subset $I\subset\{2,\ldots,m\}$, the transvection
\begin{equation}\label{eq:char2-quadric-transvection}
 R_I=1+w_I\phi_I,\qquad
 w_I=\sum_{\substack{a\in\{2,\ldots,m\}\\a\notin I}}b_{I\cup\{a\}},
 \quad
 \phi_I=b_{\{1\}\cup I}^*.
\end{equation}
It fixes every exterior basis vector except
$b_{\{1\}\cup I}$, which it sends to
$\sum_{a\notin I}b_{I\cup\{a\}}$.
The swaps for different $I$ are independent monodromy elements.
\end{lemma}

\begin{proof}
Replace $H_0H_1$ by a varying quadric and keep the remaining
hyperplanes fixed.  The same alternating resolution as in
Lemma~\ref{lem:quadric-reflections} has constant lower terms
and a pair of points for every line
$L_I=\bigcap_{i\in I}H_i$.  On an open containing the original
arrangement these pairs are distinct, all positive-dimensional
quadric strata are connected and have no higher coherent
cohomology, and the entire top coherent Frobenius is bijective.

We check independence of the point-pair covers in characteristic
two, working on the affine space of quadric equations.  This
cover is the pullback of the projective parameter cover, so
its full image is already a sufficient geometric lower bound.
On a chosen basis for a line $L$, the restricted quadratic
form is $a_LX^2+b_LXY+c_LY^2$.  Its ordered roots give the
Artin--Schreier extension
\[
 z^2+z=a_Lc_L/b_L^2
\]
of the rational function field of the space of quadrics.
The restriction map to any line is surjective, so
$a_L,b_L,c_L$ are independent linear functions.
Moreover $b_L$ is evaluation of the polar alternating form
on the Pluecker vector of $L$.  For distinct lines these
linear functions are nonassociate, by injectivity of the
Pluecker embedding.

At the prime divisor $b_{L_I}=0$, the displayed function for
$I$ has pole order two and leading residue $a_{L_I}c_{L_I}$.
This residue is not a square in the divisor's rational
function field, since $a_{L_I},c_{L_I}$ remain independent
variables there.  The corresponding functions for all
$J\ne I$ have no pole at this divisor.
No Artin--Schreier coboundary $q^2+q$ can have that leading
term: a pole of order two would require a pole of order one
for $q$ and would have square leading residue.  Thus no
nonempty sum of these functions is an Artin--Schreier
coboundary.  Their classes are independent, and the compositum
has Galois group $(\mathbf Z/2)^{\binom{d-2}{n-1}}$.
Passing to the nonempty ordinary and transverse open preserves
connectedness of this cover and hence all the independent swaps.

At the arrangement, swapping the two points on $L_I$ exchanges
the cochains $e_0\wedge e_I$ and $e_1\wedge e_I$, and fixes
all other top cochains not involving both $0$ and $1$.
These latter cochains span the arrangement cohomology; those
involving both are eliminated by the coboundary relation
$u\wedge\bigwedge^{n-1}P=0$.
In the chosen exterior basis the only changed column is
$b_{\{1\}\cup I}$, and its difference is
$(e_0+e_1)\wedge e_I=w_I$.
This proves \eqref{eq:char2-quadric-transvection}.
The vector $w_I$ is nonzero because
$d-n-1\geq1$, and its $\phi_I$ value is zero.
Thus the operator is a transvection.
\end{proof}

\begin{lemma}\label{lem:char2-coordinate-generation}
\status{Proved under stated hypotheses}
Let $G_0$ be generated by the arrangement permutation group
and the swaps of Lemma~\ref{lem:char2-independent-quadric-swaps}.
Writing $\ell=\sum_A b_A^*$, one has
\[
 G_0=
 \begin{cases}
 \mathrm{GL}_{h}(\mathbf F_2),&d-n-1\ \text{odd},\\
 \operatorname{Stab}_{\mathrm{GL}_h(\mathbf F_2)}(\ell),
       &d-n-1\ \text{even}.
 \end{cases}
\]
In the second case the stabilizer is
$\mathbf F_2^{h-1}\rtimes\mathrm{GL}_{h-1}(\mathbf F_2)$.
\end{lemma}

\begin{proof}
Permuting the labels $1,\ldots,m$ and fixing label $0$ supplies
the following operators for every $(n-1)$-subset $I$ and
every $j\notin I$:
\[
 T^I_j(b_{I\cup\{j\}})=S_I,\qquad
 S_I=\sum_{a\notin I}b_{I\cup\{a\}},
\]
with every other exterior basis vector fixed.  This operator
exchanges the two vectors $b_{I\cup\{j\}}$ and $S_I$ and
fixes each of the other vectors in this sum.
For $j,t\notin I$, $j\ne t$, the conjugate
$T^I_jT^I_tT^I_j$ is therefore the pure coordinate
transposition exchanging $b_{I\cup\{j\}}$ and
$b_{I\cup\{t\}}$ and fixing all other exterior basis vectors.
The graph of $n$-subsets joined when they share $n-1$ elements
is connected.  Its edge transpositions generate the full
coordinate symmetric group $S_h$, so $G_0$ contains that group.

Put $l=d-n-1$.  Each $T^I_j$ is a transvection whose image
vector has exactly $l$ nonzero exterior coordinates, all
different from its input coordinate.  Conjugating by $S_h$
gives every transvection
$1+v b_a^*$ with $v$ supported on any $l$ coordinates other
than $a$.  If $l=1$, these are all elementary transvections
and generate $\mathrm{GL}_h(\mathbf F_2)$.
If $l\geq2$, choose two such supports with $l-1$ common
coordinates and differing in any desired pair $b,c$, both
different from $a$.  This is possible because
$h=\binom{d-1}{n}\geq d-1\geq d-n+1=l+2$
for $n\geq2,d\geq n+2$.
The product, with the same input functional, is
\[
 1+(b_b+b_c)b_a^*.
 \tag{4}
\]
If $l$ is odd, adding pairs from an $l$-element support reduces
it to a singleton.  Thus once again all elementary
transvections occur and $G_0=\mathrm{GL}_h(\mathbf F_2)$.

Suppose $l$ is even.  Every generator fixes $\ell$:
the coordinate transvections change a column by an even-weight
vector, and an arrangement permutation fixes $\ell$ as well.
For the latter assertion, permutations of $1,\ldots,m$ act
by coordinate permutations, while $(01)$ is the product of
the commuting swaps in
Lemma~\ref{lem:char2-independent-quadric-swaps}.
These permutations generate the full arrangement $S_d$.
Consequently $G_0\subseteq\operatorname{Stab}(\ell)$.

Relabel the exterior basis as $b_1,\ldots,b_h$ and put
$f_i=b_i+b_h$ for $1\leq i<h$.
The $f_i$ form a basis of $\ker\ell$, with complement
$\langle b_h\rangle$.  Operators (4), with image $f_i$ and
input functional $b_j^*$ for $i\ne j<h$, give the elementary
matrices on $\ker\ell$ and fix $b_h$.  They generate its
full $\mathrm{GL}_{h-1}(\mathbf F_2)$.
The coordinate transposition $(i\,h)$ is
$1+f_i(b_i^*+b_h^*)$.  Multiplying it by the commuting
operators $1+f_i b_j^*$ for all $j\ne i,h$ gives
$1+f_i\ell$.  These generate all the translations of the
complement by $\ker\ell$.  Hence the full stabilizer is
contained in $G_0$, proving equality.
\end{proof}

\begin{lemma}\label{lem:pn-char2-canonical-covector}
\status{Proved under stated hypotheses}
If $d\equiv n+1\pmod2$ and $d\geq n+1$, the universal ordinary
Hasse--Witt local system has a nowhere-zero invariant
$\mathbf F_2$-valued covector, defined over $\mathbf F_2$.
\end{lemma}

\begin{proof}
Write a degree-$d$ equation as $F=\sum_{|\alpha|=d}c_\alpha x^\alpha$.
Use the standard top-cohomology basis indexed by positive
$(n+1)$-tuples $u$ with $|u|=d$.  The Hasse--Witt matrix is
$H_{uv}=c_{2v-u}$, with an out-of-range coefficient understood
as zero.  Put
\[
 w_u=\begin{cases}c_u,&\text{every coordinate of }u\text{ is odd},\\
                   0,&\text{otherwise}.\end{cases}
\]
Then
\[
 H^{\mathsf T}w=w^{[2]}.
 \tag{5}
\]
Indeed, in the $v$ coordinate the sum is
$\sum_{u\ {\rm odd}}c_{2v-u}c_u$.  Whenever the coefficient
is nonzero, $2v-u$ is also a positive all-odd tuple.
The involution $u\leftrightarrow2v-u$ cancels all paired
terms.  Its only possible fixed term is $u=v$, giving
$c_v^2$ when $v$ is all odd and zero otherwise.

There exists a positive all-odd tuple of sum $d$ under the
assumed parity and bound.  If every $w_u$ vanished, the row
of $H$ indexed by any such tuple would vanish: every
possibly nonzero coefficient in that row is itself indexed
by an all-odd tuple.  Thus $w\ne0$ wherever $H$ is invertible.
The map $x\mapsto w^{\mathsf T}x$ is Frobenius equivariant
by (5), and restricts on Lang fixed vectors to an
$\mathbf F_2$-valued covector.  It is nonzero because those
fixed vectors span the coherent fibre.
This construction is global on projective parameter space:
the coherent bundle is $\mathcal O(-1)$ tensored with the
fixed top-cohomology vector space, and the entries $w_u$
are degree-one functions of the equation.  They define an
honest morphism of this bundle to $\mathcal O$, not merely
a covector up to scalar.  Everything is defined over
$\mathbf F_2$.
\end{proof}

\begin{theorem}\label{thm:pn-char2-all-groups}
\status{Proved under stated hypotheses}
Let $k$ be any perfect field of characteristic two and
$n\geq2$.  Put $h=\binom{d-1}{n}$.  The geometric and
arithmetic monodromy groups on the ordinary smooth
degree-$d$ hypersurface locus of $\mathbf P^n_k$ are:
\[
 \begin{array}{c|c}
 \text{degree range}&\text{both monodromy images}\\ \hline
 d\leq n&1\text{ on the zero vector space}\\
 d=n+1&\mathrm{GL}_1(\mathbf F_2)=1\\
 d\geq n+2,\ d\equiv n\pmod2&\mathrm{GL}_h(\mathbf F_2)\\
 d\geq n+3,\ d\equiv n+1\pmod2&
    \operatorname{Stab}(\lambda)
       \simeq\mathbf F_2^{h-1}\rtimes
                         \mathrm{GL}_{h-1}(\mathbf F_2),
 \end{array}
\]
where $\lambda$ is the nonzero covector in
Lemma~\ref{lem:pn-char2-canonical-covector}.
In the last case $\ker\lambda$ is a proper nonzero invariant
hyperplane, so the representation is reducible.
\end{theorem}

\begin{proof}
The two initial ranks are zero and one as stated.
In all higher degrees
Lemma~\ref{lem:char2-coordinate-generation} gives the required
geometric lower bound.  In the parity with a canonical
covector, Lemma~\ref{lem:pn-char2-canonical-covector} gives
the matching upper bound, even for the arithmetic image.
For precision in comparing covectors: the subgroup $G_0$
has exactly the one nonzero invariant covector $\ell$.
The canonical $\lambda$ is nonzero and $G_0$-invariant,
so its value on the arrangement fibre is $\ell$.
Thus the two stabilizers are the same.
In the other parity the lower bound is the entire general
linear group.  The result over arbitrary perfect $k$
follows by geometric connectedness of the corresponding
finite cover under algebraically closed extension, exactly
as in the odd-characteristic theorem.
\end{proof}

\begin{corollary}\label{cor:pn-all-p-sharp-residue-threshold}
\status{Proved under stated hypotheses}
For every finite field $k$, every $n\geq2$, and every
$d\geq n+2$, each nonempty open in the smooth degree-$d$
hypersurface parameter space of $\mathbf P^n_k$ contains
hypersurfaces, over suitable finite extensions of $k$,
with every point-count residue modulo the characteristic.
The threshold $n+2$ is exact: for $d\leq n$ the only
residue is $1$, and for $d=n+1$ the ordinary open misses $1$.
\end{corollary}

\begin{proof}
Odd characteristic was settled in
Corollary~\ref{cor:pn-odd-all-residues-sharp}.  In characteristic
two, the full general linear group in rank at least two
has both trace values.  In the covector-stabilizer case,
the trace is $1+\operatorname{Tr}(B)$ with
$B\in\mathrm{GL}_{h-1}(\mathbf F_2)$; here $h-1\geq2$,
so again both values occur.  The same finite Chebotarev
and coherent trace argument gives both point-count
residues in every open.  The zero- and one-rank failures
are as in the preceding corollary.
\end{proof}

\par\smallskip\noindent\textbf{Finite checks.}
The retained program
\path{code/projective_quadric_generators.py} constructs
these exact exterior matrices over $\mathbf F_2$.
It checks that every swap has rank-one difference and
order two, that their product is the hyperplane
transposition, and that the common fixed covector space
has dimension one exactly in the indicated parity.
For the tested ranks at most fifteen it additionally
enumerates all nonzero vector orbits, obtaining either
$2^h-1$ or $2^{h-1}-1$ and $2^{h-1}$ as predicted.
These finite checks support the formulas; the uniform
group-generation proof above is independent of sampling.

% /root/monodromy with independent suggestions and audit by /root.
% General-X A.5 proof; independently audited by /root on 19 November 2026.
\subsection{All point-count residues for arbitrary smooth ambient varieties}
\label{subsec:general-X-residues}
\status{Proved under the stated hypotheses; independently audited.}

The full generic monodromy group is unnecessary for this
consequence.  We isolate one alternating permutation factor
whose action on every other contribution has trivial trace
variation.  The argument permits a drop of the stable rank
at the reducible boundary, and includes its missing generic
unit roots in an explicit error representation.

\subsubsection{Independent node covers}

\begin{lemma}\label{lem:all-X-independent-nodes}
\status{Proved under stated hypotheses}
Let $X\subset\mathbf P^r_k$ be smooth projective of pure
dimension $n\geq2$ over an algebraically closed field.
Its connected components are denoted $X_1,\ldots,X_c$ and
have degrees $D_1,\ldots,D_c$.  Fix $e\geq3$ and $m\geq n+1$.
Take $m-1$ ordered variable ambient degree-$e$ hypersurfaces
and one of degree $e+\rho$, where $0\leq\rho<e$.
There is a nonempty open parameter space $S$ on which their
restrictions $D_i$ form a simple normal crossing divisor,
all positive-dimensional intersections on each $X_j$ are
smooth and connected, and all zero-dimensional intersections
are reduced.  On this space the joint monodromy of the
node sets
\[
 X_j\cap\bigcap_{i\in I}D_i,\qquad
 1\leq j\leq c,\quad
 I\subset\{1,\ldots,m-1\},\quad |I|=n,
\]
is the full product
\[
 \prod_{j=1}^c\prod_{|I|=n} S_{N_j},
 \qquad N_j=D_j e^n.
 \tag{6}
\]
It remains this product after restriction to any nonempty
open of $S$.
\end{lemma}

\begin{proof}
The required simultaneous smoothness and proper intersection
conditions hold on a nonempty open by Bertini, applied to
the finitely many subsets of the ordered divisors.
Positive-dimensional ample complete intersections in each
$X_j$ are geometrically connected.  This also follows by
iterating the degree-zero stable Lefschetz assertion R1.

We check the precise conditions of Entin's
Proposition~3.1, arXiv version~3, pp.~7--8
\cite{Entin2021}, for the disjoint union of these node covers
over a larger open on which they are finite.
The incidence cover for one pair $(j,I)$ is irreducible and
smooth: over $X_j$ it is a product of projective bundles
given by the linear conditions that the $n$ specified
equations vanish at the point, together with the unrestricted
other parameters.  Its fibre square off the diagonal is
irreducible for the same reason over
$X_j\times X_j\setminus\Delta$; evaluation at two distinct
points is surjective for ambient degree-$e$ forms.
Thus condition (i) of that proposition holds, and the
single-factor point action is doubly transitive.

For condition (ii), fix $(j,I)$.  Choose $n-1$ of these
divisors generally so that their intersection on $X_j$ is
a smooth connected curve $C$.  Choose a point $P\in C$.
There is a degree-$e$ equation whose restriction to $C$
has a zero of order exactly two at $P$ and only simple
zeros elsewhere, while all its intersections on the other
components are transverse.  Here are the required
jet-separation details.  In local coordinates from the
given projective embedding, the functions $1,t,t^2$ occur
as restrictions of degree-$e$ forms after multiplying by a
nonvanishing linear form to make the degrees equal.
Thus the second-order jet at $P$ can be prescribed.
The linear subsystem vanishing twice at $P$ separates
first-order jets at every other point $Q$: multiply a
degree-$(e-2)$ form by the square of a hyperplane equation
which vanishes at $P$ but not at $Q$.  Since $e-2\geq1$,
these residual forms separate tangent directions.
Bertini for this subsystem therefore gives just the one
double zero and no other nontransverse intersection.
All equations not indexed by $I$ may now be chosen generally
and not to vanish at $P$.  Every other node set stays
transverse.  The same imposed-jet Bertini argument makes
all proper partial intersections smooth.  At $P$, choose
the differential of the last equation as a linear combination
of the first $n-1$ independent differentials with every
coefficient nonzero, and prescribe its quadratic coefficient
along the tangent direction to $C$ to be nonzero.  Ambient
forms of degree at least two realize these jets.  Every
proper subtuple is then transverse at $P$, and only the
full $n$-tuple loses transversality.

The base parameter space is smooth.  The point-incidence
source is smooth at its point above $P$, hence formally
irreducible there.  Its fibre has exactly $N_j-1$ geometric
points, with all but the one double point etale.
All other node covers are etale over this parameter.
These are exactly Entin's condition (ii), valid also in
characteristic two.  Concretely, the colliding length-two
algebra has the local quadratic deformation
$z^2+az+b$; the independent constant and linear jet
parameters give the transposition even in characteristic
two.  The proposition consequently gives the full product
(6).  Removing nonempty opens preserves connectedness of
each component of its finite Galois cover, proving the
last assertion.
\end{proof}

\subsubsection{The finite part of the Lang kernel at a boundary}

\begin{lemma}\label{lem:boundary-Lang-finite-part}
\status{Proved under stated hypotheses}
Let $B$ be a smooth integral scheme over an algebraically
closed field of characteristic $p$, let $\mathcal E$ be a
vector bundle on $B$ with a Frobenius-semilinear endomorphism,
and let $U\subset B$ be a dense open on which its stable
rank is the generic rank $r$.
Let $S$ be a smooth connected affine scheme with a morphism
$f:S\to B$, such that the pulled-back stable rank is a
constant $r_0$.  After shrinking $S$ if necessary, the
special Lang local system $W$ of rank $r_0$ and any specified
finite collection of finite local systems on $S$ occur as
representations of one group $H$ which also acts on a
generic fibre $V$ of rank $r$.  There is an $H$-equivariant
inclusion $W\hookrightarrow V$, the image on the specified
finite local systems is their full joint image on $S$,
and the image on $V$ is a subgroup of the geometric
monodromy image of $\mathcal E|_U$.
\end{lemma}

\begin{proof}
Shrink to affine opens so that $f$ takes values in a smooth
affine open $B_0\subset B$.  Use the graph of $f$ in
$S\times B_0$ and henselize this smooth affine scheme along
the graph.  Write $T=\operatorname{Spec}(A)$ for the result,
with closed subscheme $\operatorname{Spec}(A/I)=S$.
The ring is Noetherian and normal, and is connected, hence integral:
henselization of a Noetherian ring is Noetherian and is a
filtered etale limit of the smooth ambient ring.  Every
connected component meets the connected closed subscheme
because $I$ is in the Jacobson radical.  The nonzero flat
map from the ambient domain to its henselization is injective,
so $T$ dominates $S\times B_0$.  The inverse image $T_U$
of $U$ is consequently a nonempty dense open.

The kernel scheme of $F-1$ on the additive vector bundle
is affine, separated, of finite presentation, and etale:
in a local trivialization its equations are
$x-Ax^{[p]}=0$, whose Jacobian is the identity.
Its fibres are finite, so it is quasi-finite.
Its restriction to $S$ is finite etale after shrinking
$S$, because the stable rank is constant.
The finite-part theorem for henselian pairs,
Stacks Lemma~41.20.5, Tag~0BTN \cite{Stacks0BTN},
splits this kernel over $T$ as a finite etale part and an
open-and-closed part disjoint from $S$.
The finite part is a subgroup: sums and inverses of its
sections stay in it. Indeed the inverse image of the away part
under multiplication or inversion is open and closed in a finite
etale scheme over the henselian pair and has no closed-fibre points;
the Jacobson-radical property makes it empty. Its restriction to
$T_U$ is a subgroup of the generic Lang local system $V$,
and its restriction to $S$ is exactly $W$.

By Stacks Lemma~58.8.2, Tag~09ZS \cite{Stacks09ZS},
finite etale covers, and their morphisms and group laws,
extend uniquely from $S$ to $T$.
Thus all the specified finite local systems extend.
The map $\pi_1(T_U)\to\pi_1(T)\simeq\pi_1(S)$ is
surjective: a connected finite etale cover of the normal
integral $T$ remains connected after restriction to a
nonempty open.  Take the joint finite image of
$\pi_1(T_U)$ on $V$ and these extended systems to be $H$.
All schemes were taken over an algebraically closed field,
so its image on $V$ lies in the desired geometric image.
No full stable-rank local system has been extended across
the rank drop; only its finite boundary part was extended.
\end{proof}

\begin{lemma}\label{lem:alternating-factor-isolation}
\status{Proved under stated hypotheses}
Let a finite group $H$ act on $\mathbf F_p$-vector spaces
$P,V,W,R,Q_0,Q_1$, with exact sequences
\[
 0\to R\to P\to M\to0,\qquad
 0\to M\to W\to Q_0\to0,\qquad
 0\to W\to V\to Q_1\to0.
 \tag{7}
\]
Suppose the image on $P$ is the product
$A=\prod_{\alpha=1}^t A_{N_\alpha}$ in its direct sum
of natural permutation representations, with every
$N_\alpha\geq\max(5,p+2)$.
Set $c=\dim R+\dim Q_0+\dim Q_1$ and
\[
 K_p=
 \begin{cases}
 2,&p=2,\\
 v_2(p^2-1)+1,&p\text{ odd}.
 \end{cases}
\]
If $t>K_pc$, then traces of elements of the image of $H$
on $V$ take every value in $\mathbf F_p$.
\end{lemma}

\begin{proof}
Let $K$ be the kernel of the action on
$R\oplus Q_0\oplus Q_1$.  Its image in the product $A$
is a normal subgroup.  Since each alternating factor is
nonabelian simple, it is the product of a subset of the
factors.  The product of the omitted factors is a quotient
of $H/K$, a subgroup of $\mathrm{GL}_c(\mathbf F_p)$.
Choose $\ell=3$ for $p=2$ and $\ell=2$ otherwise.
Each omitted factor has order divisible by $\ell$.
The number of omitted factors is therefore at most
$v_\ell(|\mathrm{GL}_c(\mathbf F_p)|)\leq K_pc$.

For completeness, the last estimate is elementary.
For $p=2$,
\[
 v_3(|\mathrm{GL}_c(\mathbf F_2)|)
 =\lfloor c/2\rfloor+
             v_3(\lfloor c/2\rfloor!)\leq2c.
\]
For odd $p$, the lifting-the-exponent formulas give
\[
 v_2(|\mathrm{GL}_c(\mathbf F_p)|)
 =c\,v_2(p-1)+\lfloor c/2\rfloor v_2(p+1)
                 +v_2(\lfloor c/2\rfloor!)
 \leq K_pc.
\]
The estimates also hold when $c=0$.
At least one factor $A_{N_\alpha}$ thus survives in the
image of $K$.  For any element $a$ of that single factor,
choose a lift $g\in K$ acting as $a$ on its permutation
summand of $P$ and identically on every other summand.
Trace additivity in (7) gives
\[
 \operatorname{Tr}(g\mid V)
  =\dim V-N_\alpha+\#\operatorname{Fix}(a)
       \quad\text{in }\mathbf F_p.
 \tag{8}
\]
Every integer $b=0,\ldots,p-1$ is the fixed-point count of
an even permutation on $N_\alpha$ letters.  Fix $b$ letters
and derange the remaining $q=N_\alpha-b\geq3$.
For odd $q$ use one $q$-cycle, and for even $q\geq4$ use
disjoint cycles of lengths $2$ and $q-2$.
These permutations are even and have exactly the required
fixed points.  Equation (8) proves the claim.
\end{proof}

\subsubsection{Explicit bounds for the error representation}

Return to a smooth projective connected
$X\subset\mathbf P^r_k$ over a finite field of characteristic
$p$, and put $L=\mathcal O_X(1)$.
The following numerical data are all computable from this
fixed embedded variety.  Cohomology dimensions are unchanged
after passage to $\bar k$ and include all its geometric
components.
Let
\[
 D=\deg_L(X),\quad A=h^{n-1}(X,\mathcal O_X),\quad
 P_\omega(t)=\chi(X,\omega_X\otimes L^t).
\]
Choose a regularity bound $r_\omega$ for the embedded
dualizing sheaf and put $d_{\rm coh}=\max(1,r_\omega-1)$.
For every $d\geq d_{\rm coh}$ the coherent relative bundle
has rank
\begin{equation}\label{eq:all-X-relative-coherent-rank}
 b_d=P_\omega(d)-h^0(X,\omega_X).
\end{equation}
Indeed Serre duality and regularity give
$H^{n-1}(X,L^{-d})=0$, and the divisor sequence identifies
the coherent cokernel with the kernel of the surjection
$H^n(X,L^{-d})\to H^n(X,\mathcal O_X)$.
This last map is onto because $H^n(Y,\mathcal O_Y)=0$.
Serre duality and the higher-cohomology vanishing for
$\omega_X\otimes L^d$ give \eqref{eq:all-X-relative-coherent-rank}.
The same construction works in families of arbitrary
effective Cartier divisors, including the boundary unions:
the dual map is multiplication by the nonzerodivisor
defining the divisor, so the displayed kernel is a vector
bundle and carries its natural semilinear Frobenius.

Choose an integer
\[
 e\geq3,\qquad e^n\geq\max(5,p+2),
 \qquad N=De^n .
 \tag{9}
\]
This guarantees that the node count on every geometric
component is at least $\max(5,p+2)$.
For $1\leq a\leq n-1$ and $t\geq0$ define the finite
cohomological upper bounds
\begin{equation}\label{eq:all-X-B-at}
 B_{a,t}=
 \sum_{j=0}^a\binom aj
 \max_{je\leq v\leq j(2e-1)}
       h^{t+j}(X,L^{-v}),
\end{equation}
where $v$ is integral and a cohomology group in degree
greater than $n$ is zero.
The Koszul resolution shows that $B_{a,t}$ bounds
$h^t(Z,\mathcal O_Z)$ for every proper complete intersection
$Z$ of $a$ divisors whose degrees lie in $[e,2e-1]$.
Its hypercohomology spectral sequence has terms
$H^{t+j}(X,L^{-\sum d_i})$ for subsets of $j$ equations,
which gives precisely the bound
\eqref{eq:all-X-B-at}.  The dimension of
$H^t_{\mathrm{et}}(Z_{\bar k},\mathbf F_p)$ is at most this
coherent dimension, by Artin--Schreier/Lang.

Set
\[
 R_0=A+\sum_{a=1}^{n-1}B_{a,n-a-1},
 \qquad
 Q_0=\sum_{a=1}^{n-1}B_{a,n-a}
           +D e^{n-1}(2e-1).
 \tag{10}
\]
For each $\rho=0,\ldots,e-1$, form the polynomial
\[
 q_\rho(z)=P_\omega(ez+\rho)-h^0(X,\omega_X)
                           -N\binom{z-1}{n}.
 \tag{11}
\]
Its degree is at most $n-1$, since the leading terms cancel.
Let $H_0$ be the maximum, over these $e$ polynomials, of
the sum of the absolute values of their coefficients.
Put
\begin{equation}\label{eq:all-X-C-bound}
\begin{aligned}
 C&=\left\lceil2R_0+Q_0+H_0\right\rceil,
       \qquad d_{\rm res}=eM_0,\\
 M_0&=\max\left\{2(n+1),\left\lceil d_{\rm coh}/e\right\rceil,
       \left\lfloor2\cdot3^{n-1}n!K_pC\right\rfloor+1\right\}.
\end{aligned}
\end{equation}
The constants can be very large; no optimality is asserted.

\begin{lemma}\label{lem:all-X-boundary-errors}
\status{Proved under stated hypotheses}
For $d=em+\rho\geq d_{\rm coh}$, with $m\geq n+1$ and
$0\leq\rho<e$, the construction of
Lemma~\ref{lem:all-X-independent-nodes} and
Lemma~\ref{lem:boundary-Lang-finite-part} gives exact
sequences of the form (7).  Here $P$ is the direct sum of
the node permutation spaces excluding the one exceptional
degree-$(e+\rho)$ component, so
$\dim P=N\binom{m-1}{n}$; $V$ is the generic relative
Lang fibre; and
\[
 \dim R+\dim Q_0^{\rm err}+\dim Q_1^{\rm err}
       \leq C(m+1)^{n-1}.
 \tag{12}
\]
The superscript distinguishes these error spaces from the
numerical constant $Q_0$ in (10).
\end{lemma}

\begin{proof}
Work first over $\bar k$ and shrink the arrangement
parameter space $S$ so that the relative stable rank is
constant and its Lang sheaf $W$ is finite locally constant.
The normal crossing cover of $Y=\bigcup_{i=1}^mD_i$
has an exact alternating resolution of its constant
$\mathbf F_p$-sheaf by the constant sheaves of its
intersections: exactness is checked on a geometric
stalk, where it is the augmented complex of a simplex.
Its cohomology spectral sequence is
\[
 E_1^{s,t}=
 \bigoplus_{|I|=s+1}
 H^t_{\mathrm{et}}(D_I,\mathbf F_p)
 \Longrightarrow H^{s+t}_{\mathrm{et}}(Y,\mathbf F_p).
 \tag{13}
\]
The last cochain term $E_1^{n-1,0}$ maps naturally to the
abutment in degree $n-1$.  Its summands from intersections
not involving the exceptional component form $P$.
Compose this map with the quotient by the ambient
$H^{n-1}_{\mathrm{et}}(X,\mathbf F_p)$ to obtain
$P\to W$, and write its kernel as $R$ and image as $M$.
This is a map of finite local systems on $S$ after the
chosen shrinking, and therefore its kernel and cokernel
are also finite local systems.

The kernel of the full last-cochain map to the abutment
is bounded by all incoming terms in the spectral sequence:
for $r=1,\ldots,n-1$ they originate at
$(n-1-r,r-1)$, of total degree $n-2$.
Taking the ambient quotient adds at most $A$ to the
kernel dimension.  Thus
\[
 \dim R\leq A+
       \sum_{a=1}^{n-1}\binom ma B_{a,n-a-1}
       \leq R_0(m+1)^{n-1}.
 \tag{14}
\]
There is no assumption that these higher differentials
vanish.

The quotient of the abutment by the full last-cochain
image is bounded by its other degree-$(n-1)$ graded terms,
namely the terms with $t>0$ in (13).  Removing from the
last-cochain space the intersections involving the
exceptional component adds at most
$D e^{n-1}(e+\rho)\binom{m-1}{n-1}$.
Passing to the relative ambient quotient cannot increase
this cokernel.  Consequently
\[
 \dim(W/M)\leq
 \sum_{a=1}^{n-1}\binom ma B_{a,n-a}
 +D e^{n-1}(2e-1)\binom{m-1}{n-1}
 \leq Q_0(m+1)^{n-1}.
 \tag{15}
\]

The finite-boundary-part lemma supplies $W\subset V$,
while $\dim V\leq b_d$.  Since
$\dim W\geq\dim P-\dim R$, we obtain
\[
 \dim(V/W)
 \leq b_d-N\binom{m-1}{n}+\dim R
 \leq(H_0+R_0)(m+1)^{n-1}.
 \tag{16}
\]
The last inequality is the coefficient bound from (11).
Equations (14)--(16) prove (12).
The exact sequences (7) now hold with
$Q_0^{\rm err}=W/M$ and $Q_1^{\rm err}=V/W$.
All node covers and these special finite systems are
included in the joint group of the boundary-part lemma;
their full independent node image is retained.
\end{proof}

\begin{theorem}\label{thm:all-X-residue-density}
\status{Proved under stated hypotheses}
Let $k$ be a finite field of characteristic $p>0$, let
$X\subset\mathbf P^r_k$ be smooth projective connected of
dimension $n\geq2$, and define the explicit integer
$d_{\rm res}$ by \eqref{eq:all-X-C-bound}.
For every $d\geq d_{\rm res}$ and every nonempty open
$\mathcal U$ in the degree-$d$ ambient hypersurface
parameter space, each $a\in\mathbf F_p$ occurs as
\[
 \#(X\cap H)(k')\equiv a\pmod p
\]
for some finite extension $k'/k$ and some
$H\in\mathcal U(k')$ with $X\cap H$ smooth.
Moreover, traces of the geometric generic relative
Hasse--Witt monodromy group take every value in
$\mathbf F_p$ in this degree range.
Geometric connectedness of $X$ is not required.
\end{theorem}

\begin{proof}
Write $d=em+\rho$ with $0\leq\rho<e$, so $m\geq M_0$.
The number of independent alternating factors is
$c_X\binom{m-1}{n}$, where $c_X\geq1$ is the number of
geometric connected components.  In particular it is
at least $\binom{m-1}{n}$.
Since $m\geq2(n+1)$,
\[
 \binom{m-1}{n}\geq\frac{m^n}{2^n n!},
 \qquad
 (m+1)^{n-1}\leq(3m/2)^{n-1}.
 \]
The strict inequality in the definition of $M_0$ therefore
gives
\[
 \binom{m-1}{n}>K_pC(m+1)^{n-1}.
 \tag{17}
\]
Restrict the joint group of the preceding lemma to the
inverse image of the product of alternating node groups.
It still surjects to that product.  Lemma
\ref{lem:alternating-factor-isolation}, with the bound (12),
shows that its traces on $V$ take every value in
$\mathbf F_p$.  This group maps into geometric monodromy
on the full generic smooth maximal-stable-rank open, so
the latter has the same trace-surjectivity.

It remains to pass from traces to the stated rational
points, with the exact extension-of-constants direction.
Replace $k$ by a finite extension splitting the geometric
components of $X$ and killing the finite Frobenius actions
on all constant spaces
$H^i_{\mathrm{et}}(X_{\bar k},\mathbf F_p)$.
We may further kill the finite arithmetic quotient of
the generic finite monodromy image.  A rational point
over any further finite extension is also a point over a
finite extension of the original $k$, so this replacement
preserves the required existential conclusion.
Every original nonempty open remains nonempty after this
base change and meets the dense smooth maximal-rank open.

Apply the finite Chebotarev theorem used earlier
\cite[Theorem~9.7.13]{KatzSarnak1999} to the finite
representation over this open.  Each geometric conjugacy
class giving a specified trace occurs over a suitable
finite extension, inside the prescribed open.
The Artin--Schreier form of weak Lefschetz R1 identifies
the lower stable cohomology of the section with that of
$X$, and identifies the remaining top quotient with $V$.
The coherent trace congruence
\cite[Corollary~3.2]{Katz1973Congruence} consequently gives
\[
 \#Y(k')\equiv
  \sum_{i=0}^{n-1}(-1)^i
       \dim_{\mathbf F_p}H^i_{\mathrm{et}}(X_{\bar k},\mathbf F_p)
  +(-1)^{n-1}\operatorname{Tr}(F_{k'}\mid V)\pmod p.
 \tag{18}
\]
Its first term is constant by the chosen field extension.
The second takes every value, which proves the claim.
If one uses arithmetic rather than geometric Frobenius
in defining the monodromy representation, replace the
chosen group element by its inverse; all conjugacy classes
are still available.  This keeps the coherent-Frobenius
trace convention in (18) unchanged.
\end{proof}

\par\smallskip\noindent\textbf{Scope of the conclusion.}
This proves the residue consequence A.5 with a specified
bound for every ambient variety in the question, including
the non-geometrically-connected case.  It does not identify
the full monodromy group in A.4 or the exact generic stable
rank in A.3.  In particular it never assumes generic
relative ordinarity, which the preceding smooth
$\alpha_4$ example disproves.



% /root/monodromy, 19 November 2026; independent parent audit passed.
\subsection{Effective generic-rank calculation and exact determinant images}
\label{subsec:effective-rank-determinant}
\status{Proved under stated hypotheses; independently audited.
These do not constitute a structural classification of general ambient monodromy.}

\begin{proposition}\label{prop:effective-generic-rank}
\status{Proved under stated hypotheses}
Let $k$ be a finite field of characteristic $p$ and let
$(X,L)$ be specified by homogeneous equations and its
projective embedding.  Fix a degree $d$ for which the
coherent relative group is the vector bundle of rank
$b=h^0(X,\omega_X\otimes L^d)-h^0(X,\omega_X)$ described
above; in particular every $d\geq d_{\rm coh}$ is allowed.
There is a finite exact algorithm to compute the generic
stable rank $r_d$ and equations for a dense open on
which that rank is attained.

More precisely, take a rational basis of the relative
coherent bundle over the parameter function field
$K=k(t_1,\ldots,t_M)$, and let $A\in M_b(K)$ be its
semilinear Frobenius matrix.  With $A^{[p^j]}$ denoting
entrywise $p^j$-th powers, one has
\begin{equation}\label{eq:generic-rank-finite-product}
 r_d=\operatorname{rank}_K
       \left(A A^{[p]}\cdots A^{[p^{b-1}]}\right).
\end{equation}
For $b=0$ this rank is zero.  The nonvanishing of the
maximal nonzero minors of this displayed product gives
the maximal-stable-rank open inside the chosen basis
chart, after also imposing smoothness.
\end{proposition}

\begin{proof}
There is no search through hypersurfaces of increasing
finite fields in this calculation.  The cohomological
matrix itself can be computed over the rational function
field from the input equations.  A finite graded free
resolution computes the coherent cohomology spaces, and
the standard affine Cech presentation computes the
induced maps by taking $p$-th powers and multiplying by
the universal section to the power $p-1$.
More intrinsically, the relative bundle is the kernel of
\[
 H^n(X,L^{-d})\otimes\mathcal O(-1)
       \longrightarrow H^n(X,\mathcal O_X)\otimes\mathcal O
\]
given by multiplication by the universal section.
The induced Frobenius on the source is the composite
\[
 H^n(X,L^{-d})
 \xrightarrow{\ F\ }H^n(X,L^{-pd})
 \xrightarrow{\ s^{p-1}\ }H^n(X,L^{-d});
\]
it preserves that kernel.  Kernel bases and the matrices
of these maps are finite linear-algebra computations
on the graded-resolution/Cech representatives; polynomial
ideal membership is decidable by Groebner bases.
Thus an explicit rational matrix $A$ is obtained by a
finite procedure.

For the stabilization bound, put
$B_0=1$ and
$B_j=A A^{[p]}\cdots A^{[p^{j-1}]}$ for $j>0$.
The spaces $I_j=\operatorname{im}_K B_j$ are decreasing.
If $\dim I_j=\dim I_{j+1}$, then they are equal and
the Frobenius-linearized map
$F^*I_j\to I_j$ is surjective, hence an isomorphism:
its image is exactly $I_{j+1}$ and both have the same
dimension.  All subsequent images are the same.
Every strict decrease lowers dimension by at least one,
so stabilization occurs by $j=b$.  Extension to the
perfect closure of $K$ preserves all these ranks, and
the stable part there has dimension $\operatorname{rank}B_b$.
This proves \eqref{eq:generic-rank-finite-product}.
The same argument applies in every geometric fibre
where the rational basis and its matrix are regular.
Nonvanishing of its generic maximal minors consequently
cuts out the asserted open.
\end{proof}

\begin{remark}\label{rem:algorithm-versus-generic-formula}
\status{Proved under stated hypotheses}
Proposition~\ref{prop:effective-generic-rank} is an exact
algorithm for each specified input.  We distinguish it
from a closed structural formula for $r_d$ in terms of
fixed invariants of $X$, or a uniform description of its
defect from coherent dimension as $d$ varies.
The latter determination for arbitrary $X$ remains
unsettled in this dossier.  No claim that the algorithm
was executed for every possible input is intended.
\end{remark}

\begin{proposition}\label{prop:exact-general-determinant-image}
\status{Proved under stated hypotheses}
Let $k$ be a perfect field of characteristic $p$.
Let $\mathcal V$ be the rank-$r$ Lang local system of an
invertible Frobenius bundle on a nonempty open
$U\subset\mathbf P^M_k$, as for the stable relative
Hasse--Witt bundle.  Choose a rational trivialization of
its determinant line and write its Frobenius as
$e\mapsto h e$, where
$h\in K^\times$ and $K=k(\mathbf P^M)$.
Put $N=p-1$ and
\begin{equation}\label{eq:det-valuation-gcd}
 g=\gcd\bigl(N,\ v_D(h)\text{ for every prime divisor }
                         D\subset\mathbf P^M_{\bar k}\bigr),
 \qquad m=N/g.
\end{equation}
Then $h=c f^g$ for some $f\in K^\times$ and $c\in k^\times$.
The geometric determinant image is the unique subgroup
of $\mathbf F_p^\times$ of order $m$.

The arithmetic determinant image is the unique subgroup
of order $m\delta$, where
$\delta$ is the order of the class of $c$ in
$k^\times/(k^\times)^g$.  In particular, for
$k=\mathbf F_q$,
\begin{equation}\label{eq:det-arithmetic-explicit}
 \delta=\operatorname{ord}_{\mu_g}
                   \left(c^{(q-1)/g}\right),
 \qquad
 |\det G_{\rm arith}|=\frac{p-1}{g}\,
                   \operatorname{ord}_{\mu_g}
                   \left(c^{(q-1)/g}\right).
\end{equation}
Both determinant images are trivial when $p=2$, or when
the local system has rank zero.
For finite $k$, factorization of the numerator and
denominator of $h$ computes these quantities exactly.
\end{proposition}

\begin{proof}
The determinant line $\mathcal L=\det\mathcal E_{\rm ss}$
has a nowhere-vanishing Frobenius section in
$\mathcal L^{1-p}$ on $U$.  In the chosen rational basis
it is $h$.  A Lang fixed basis $ze$ satisfies
$h z^p=z$, so the determinant frame torsor is
\begin{equation}\label{eq:det-kummer-general-image}
 z^N=h^{-1}.
\end{equation}
All $N$-th roots of unity lie in $\mathbf F_p^\times$.
The determinant image is therefore the Galois group
of the Kummer extension cut out by this equation.

Over $\bar k$, an element of the rational function field
is an $a$-th power, for $a\mid N$, exactly when all its
divisor valuations are divisible by $a$.
One direction is immediate.  Conversely, if
$\operatorname{div}(h)=aD$, then $a[D]=0$ in
$\operatorname{Pic}(\mathbf P^M_{\bar k})=\mathbf Z$.
Thus $D$ is principal, and the remaining constant has
an $a$-th root in $\bar k$.
The Kummer class of $h$ consequently has exact order
$N/g=m$ over the geometric function field, proving the
geometric assertion.

Because $k$ is perfect, its prime divisors become
geometrically reduced after algebraic extension.
The divisor $\operatorname{div}(h)/g$ is therefore an
integral divisor defined over $k$.
Its class in $\operatorname{Pic}(\mathbf P^M_k)=\mathbf Z$
is zero.  It is the divisor of a rational function
$f\in K^\times$, and $h/f^g=c\in k^\times$.
In the Kummer extension of \eqref{eq:det-kummer-general-image}
the element $w=z^m f$ satisfies $w^g=c^{-1}$.
Thus the extension contains the constant field
$k'=k(c^{1/g})$, whose degree is $\delta$ by elementary
Kummer theory.  Over $k'K$, the remaining equation has
degree exactly $m$: its divisor valuations, divided by
$g$, have gcd one with $m$ by the definition of $g$.
This degree remains $m$ over $\bar k K$, so $k'$ is
the entire algebraic constant extension.
The arithmetic extension therefore has degree
$m[k':k]=m\delta$, as claimed.

For a finite field, the map
$k^\times/(k^\times)^g\to\mu_g$ sending
$c$ to $c^{(q-1)/g}$ is an isomorphism of cyclic groups.
This proves \eqref{eq:det-arithmetic-explicit}.
Changing $f$ multiplies $c$ by a $g$-th power and does
not change $\delta$.  Changing the rational determinant
basis multiplies $h$ by an $N$-th power and does not
change any of the images.
For rank zero the determinant is $1$; for $p=2$ the
target $\mathbf F_2^\times$ is trivial.
\end{proof}

\begin{remark}\label{rem:det-frobenius-convention}
\status{Proved under stated hypotheses}
At a rational point over $\mathbf F_q$ where a determinant
basis is regular, the determinant of coherent $q$-Frobenius
is $h^{(q-1)/(p-1)}$.  On the Lang frame torsor,
arithmetic Frobenius acts by its inverse, since
$z^q/z=h^{-(q-1)/(p-1)}$.
Thus the geometric-Frobenius determinant character is
the coherent Frobenius determinant; the arithmetic
Frobenius character is its inverse.  The subgroup orders
in Proposition~\ref{prop:exact-general-determinant-image}
are unaffected by this convention.
\end{remark}


% /root/monodromy, 19 November 2026, 05:39 UTC.
\subsection{An ambient adjunction obstruction in characteristic two}
\label{subsec:general-adjunction-covector}
\status{Proved under the stated hypotheses; independently audited.}

\begin{theorem}\label{thm:general-adjunction-covector}
\status{Proved under stated hypotheses}
Let $X/k$ be smooth, projective and geometrically connected, of
dimension $n\geq2$, over a perfect field of characteristic two.
Let $L=\mathcal O_X(1)$ for the fixed embedding $X\subset\mathbf P^r$.
Write $R_I=\operatorname{reg}(\mathcal I_X)$.
Suppose that, for some
integer $d_*$, there is a line bundle $M_*$ defined over $k$ with
\[
 M_*^{\otimes2}\simeq\omega_X\otimes L^{d_*}.
\]
Put $d=d_*+2t$ and $M=M_*\otimes L^t$.
Write $F=F_{X/k}:X\to X^{(2)}$ and
$\mathcal B=\operatorname{coker}(\mathcal O_{X^{(2)}}\to F_*\mathcal O_X)$.
Choose Castelnuovo--Mumford regularities
\[
 R_M=\operatorname{reg}_L(M_*),\qquad
 R_B=\operatorname{reg}_{L^{(2)}}
       (\mathcal B^\vee\otimes M_*^{(2)}),
 \qquad
 \kappa=\frac{c_1(\omega_X)c_1(L)^{n-1}}{c_1(L)^n}.
\]
Assume
\begin{equation}\label{eq:adjunction-covector-threshold}
 t\geq\max\{0,R_M,R_B-1\},\qquad
 d\geq\max\{d_{\rm coh},R_I-1\},\qquad d>\kappa.
\end{equation}
On the maximal-stable-rank smooth-hypersurface open $U_d$,
the relative Hasse--Witt local system $\mathcal V_d$ has a
canonical nonzero morphism
\begin{equation}\label{eq:adjunction-invariant-covector}
 \lambda_M:\mathcal V_d\longrightarrow\mathbf F_2.
\end{equation}
Here ``canonical'' is relative to the given square root and its
isomorphism with $\omega_X L^d$; the construction is described below.
In particular, both the geometric and arithmetic images fix a
nonzero covector.  If $\operatorname{rank}\mathcal V_d\geq2$,
they are proper subgroups of $\GL(\mathcal V_{d,\bar\eta})$ and
the representation is reducible.

If the square root is defined only after an algebraic extension
of $k$, the conclusion about the geometric image still holds.
There are arbitrarily large degrees satisfying
\eqref{eq:adjunction-covector-threshold}; the rank is at least two
for all sufficiently large such degrees, by the previously proved
explicit stable-rank lower bound (or by the all-residue theorem
when $k$ is finite).
\end{theorem}

\begin{proof}
The morphism $F$ is finite flat of degree $2^n$.  Finite duality,
or the Cartier trace calculation of
\cite[Lemma~1.3.6, Proposition~1.3.7, Theorem~1.3.8, pp.~26--28]
{BrionKumar2005}, gives
\begin{align}
 F^!M^{(2)}
 &\simeq \omega_X\otimes F^*(\omega_{X^{(2)}}^{-1}\otimes M^{(2)})
       \simeq\omega_X^{-1}\otimes M^{\otimes2}
       \simeq L^d,\label{eq:adjunction-finite-duality}\\
 H^0(X,L^d)
 &\simeq\Hom_{X^{(2)}}(F_*\mathcal O_X,M^{(2)}).
 \label{eq:adjunction-section-map}
\end{align}
These are relative-Frobenius identifications and are $k$-linear.
Given $s\in H^0(X,L^d)$, let $\tau_s$ be the corresponding map.
Set
\begin{equation}\label{eq:adjunction-w-definition}
 v_s=\tau_s(1)\in H^0(X^{(2)},M^{(2)}),\qquad
 w_s=F^*v_s\in H^0(X,M^{\otimes2})
                  =H^0(X,\omega_X L^d).
\end{equation}
The map $s\mapsto w_s$ is $k$-linear.  Relative Frobenius is
essential in this assertion: its pullback fixes base-field
coefficients.  On a perfect geometric fibre one can express
$w_s$ as a square using the absolute Cartier operator, but
one must not square the universal parameter coefficients.

We check the Frobenius identity on geometric fibres.  Locally
choose a generator $m$ of $M$ and etale coordinates whose top
form is $\eta$.  The induced generator of $L^d$ is $m^2/\eta$;
write $s=fm^2/\eta$.  Write $T$ for the coefficient form of
the absolute Cartier trace, so that
\[
 C(f\eta)=T(f)\eta,\qquad T(a^2b)=aT(b).
\]
Formula~\eqref{eq:adjunction-w-definition} becomes
$w_s=T(f)^2m^2$.  The Cartier operator dual to
\[
 F_s:H^n(X,L^{-d})\longrightarrow H^n(X,L^{-d}),
       \quad \xi\longmapsto sF_X(\xi),
\]
sends $gm^2$ to $T(fg)m^2$.  Consequently
\begin{equation}\label{eq:adjunction-covector-Cartier}
 C_s(w_s)=T(fT(f)^2)m^2=T(f)^2m^2=w_s.
\end{equation}
Under Serre duality this is exactly
\[
 \langle F_s\xi,w_s\rangle=\langle\xi,w_s\rangle^2.
\]
Thus $w_s$ is a Frobenius-equivariant scalar-valued covector.
In matrix notation this is $A^tw_s=w_s^{[2]}$, consistently
with the explicit projective-space formula.

Because $w_s$ is linear in $s$, it defines a global covector
on $H^n(X,L^{-d})\otimes\mathcal O(-1)$ over
$\mathbf P H^0(X,L^d)$.  It restricts to the relative coherent
bundle
\[
 \mathcal E_d=\ker\bigl(H^n(X,L^{-d})\otimes\mathcal O(-1)
       \xrightarrow{\ s\ }H^n(X,\mathcal O_X)\otimes\mathcal O\bigr).
\]
The Frobenius identity shows that the restricted covector
kills the Frobenius-nilpotent part and induces a map from
the stable Lang local system to the constant sheaf
$\mathbf F_2$.

It remains to prove nonvanishing.  The sheaf $\mathcal B$ is
locally free: etale locally the standard square-free
monomials form a basis of $F_*\mathcal O_X$ containing $1$.
Dualizing its defining exact sequence and tensoring with
$M^{(2)}$ gives
\[
 0\longrightarrow\mathcal B^\vee\otimes M^{(2)}
 \longrightarrow\mathcal Hom(F_*\mathcal O_X,M^{(2)})
 \xrightarrow{\ \tau\mapsto\tau(1)\ } M^{(2)}
 \longrightarrow0.
\]
The regularity bounds in
\eqref{eq:adjunction-covector-threshold} give
\[
 H^1(X^{(2)},\mathcal B^\vee\otimes M^{(2)})=0,
       \qquad H^0(X^{(2)},M^{(2)})\ne0.
\]
Indeed, the first is the degree-one regularity vanishing,
and the second follows from global generation of the
nonzero line bundle $M$.  Thus $s\mapsto v_s$ is surjective
onto a nonzero space, and $w_s\ne0$ for generic $s$.

Take such an $s$ with $Y=(s=0)$ smooth.  If its covector
vanished on $E_s$, duality for the displayed kernel would
give $w_s=s\eta_0$ for some $\eta_0\in H^0(X,\omega_X)$.
Every valuation of $w_s=F^*v_s$ is even.  The divisor of
$s$ is reduced, so the valuation of $\eta_0$ along each
component of $Y$ would be odd and positive.  Hence
$\eta_0\in H^0(X,\omega_X L^{-d})$.  This last space is zero:
a nonzero section would give an effective divisor with
intersection number $(\kappa-d)c_1(L)^n<0$ against
$L^{n-1}$, an impossibility.  Thus the restricted covector
is nonzero generically.  A morphism of finite local
systems on the connected $U_d$ has locally constant rank;
it is therefore nonzero everywhere on the maximal-rank open
of the complete linear system.

For the prescribed ambient degree-$d$ hypersurface parameter
space, the additional bound $d\geq R_I-1$ gives
$H^1(\mathbf P^r,\mathcal I_X(d))=0$.  The restriction
$H^0(\mathbf P^r,\mathcal O(d))\to H^0(X,L^d)$ is therefore
surjective.  Its induced linear projection, away from its
center, pulls back the construction just given.  The smooth
ambient open avoids the center; its maximal stable rank is
the same as in the complete linear system, and the pulled-back
covector is nonzero on it.  This is the $U_d$ in the statement.

Everything was defined over $k$ if the square root was,
so arithmetic as well as geometric monodromy fixes it.
After extending the base field the same construction
gives the asserted geometric conclusion.  The kernel
of a nonzero covector is a nonzero proper invariant
subspace when the rank is at least two.
\end{proof}

\begin{remark}\label{rem:adjunction-square-Pn}
\status{Proved under stated hypotheses}
For $X=\mathbf P^n$ and $L=\mathcal O(1)$, the square-root
condition is $d-n-1\equiv0\pmod2$.  In homogeneous
monomial coordinates the construction recovers the
covector whose entry indexed by $u$ is the defining
coefficient $c_u$ when all coordinates of $u$ are odd,
and zero otherwise.  Thus this obstruction is an
ambient adjunction phenomenon, rather than a property
limited to plane theta characteristics.
The theorem gives a subgroup constraint for general
$X$, not equality with the full covector stabilizer.
\end{remark}


% /root/monodromy, 19 November 2026; independently audited by /root.
\subsection{Many fixed covectors for ordinary abelian ambient varieties}
\label{subsec:ordinary-abelian-covectors}
\status{Proved; independent parent proof and source audit passed.}

\begin{proposition}\label{prop:ordinary-abelian-fixed-quotient}
\status{Proved under stated hypotheses}
Let $A$ be an ordinary abelian variety of dimension $n\geq2$
over a perfect field of characteristic two, with fixed very
ample $L=\mathcal O_A(1)$.  For every sufficiently large even
degree $d=2t$, the geometric monodromy of $\mathcal V_d$
acts trivially on a quotient of dimension $2^n-1$.
More precisely, write $F:A\to A^{(2)}$ and, after a field
extension splitting $\ker F$, write
\[
 F_*\mathcal O_A=\bigoplus_{\chi\in H}\mathcal N_\chi,
 \qquad H=\operatorname{Hom}(\ker F,\mathbf G_m)
                 \simeq(\mathbf Z/2)^n.
\]
Set
\[
 R=\max_{\chi\in H}
         \operatorname{reg}_{L^{(2)}}(\mathcal N_\chi).
\]
It suffices that $t\geq\max\{1,R\}$ and
$2t\geq\max\{d_{\rm coh},\operatorname{reg}(\mathcal I_A)-1\}$.
After the same splitting extension, the arithmetic image
also fixes this quotient pointwise.
\end{proposition}

\begin{proof}
The character decomposition is the decomposition of the
algebra of a torsor under the diagonalizable group
$\ker F\simeq\mu_2^n$.  In particular,
$\mathcal N_0=\mathcal O_{A^{(2)}}$ and multiplication gives
$\mathcal N_\chi\otimes\mathcal N_\psi\simeq
\mathcal N_{\chi+\psi}$.
The line bundles are pairwise nonisomorphic two-torsion
bundles in $\Pic^0(A^{(2)})$.
This decomposition, including all the two-torsion line
bundles, is also stated and proved in
\cite[Lemma~4.5, p.~11]{SannaiTanaka2016}.

Choose the unique nonzero invariant top differential
$\omega$ normalized by $C(\omega)=\omega$, where $C$
is absolute Cartier.  Such a normalization exists because
$A$ is ordinary and the characteristic is two: if
$C(\omega_0)=c\omega_0$ with $c\ne0$, take
$\omega=c^2\omega_0$.
The coefficient trace $F_*\mathcal O_A\to\mathcal O_{A^{(2)}}$
defined by this differential is the projection to the
trivial character.  Indeed its restriction to every
nontrivial character line is zero, since a nontrivial
degree-zero line bundle has no global sections; on
$\mathcal N_0$ the normalization makes it the identity.

Because $L^{2t}=F^*(L^{(2)t})$, the projection formula
gives a direct sum decomposition
\begin{equation}\label{eq:abelian-section-character-decomposition}
 W:=H^0(A,L^{2t})
    =\bigoplus_{\chi\in H}
       H^0(A^{(2)},L^{(2)t}\otimes\mathcal N_\chi).
\end{equation}
Write $s=\sum_\chi s_\chi$ for the corresponding
character components, regarded again as sections on $A$.
Each summand in \eqref{eq:abelian-section-character-decomposition}
is nonzero under the stated regularity bound.  Thus for
generic $s$ all its $2^n$ components are nonzero, and
these components are linearly independent.

We claim that each
\[
 w_\chi(s)=s_\chi\omega\in H^0(A,\omega_A L^{2t})
\]
is a Frobenius-equivariant covector for $F_s=sF_A$.
The twisted Cartier operator dual to $F_s$ sends it to
$C(s s_\chi\omega)$.  In the sum
$s s_\chi=\sum_\psi s_\psi s_\chi$, every term with
$\psi\ne\chi$ has nontrivial character and is killed
by Cartier, whereas
\[
 C(s_\chi^2\omega)=s_\chi C(\omega)=s_\chi\omega.
\]
Consequently $C_s(w_\chi)=w_\chi$.
The character projections are linear in $s$, so these
are global covectors on the tautologically twisted
coherent bundle, just as in
Theorem~\ref{thm:general-adjunction-covector}.

By Serre duality the dual of the relative coherent
fibre is $W/k s$, since $H^0(A,\omega_A)=k\omega$.
The $2^n$ independent vectors $s_\chi$ span a subspace
whose intersection with $ks$ is exactly $ks$.
Their classes in $W/ks$ therefore have dimension
$2^n-1$, with the single relation
\begin{equation}\label{eq:abelian-covector-single-relation}
 \sum_{\chi\in H}w_\chi=0
       \quad\hbox{on the relative coherent fibre}.
\end{equation}
They kill the nilpotent part by Frobenius equivariance,
so remain independent on the stable part.  Lang descent
turns them into $2^n-1$ independent $\mathbf F_2$-valued
covectors; independence can be checked after tensoring
with the algebraic closure.  Their rank is constant on
the connected maximal-rank open.  The ambient restriction
bound in the statement permits pullback from the complete
linear system to the prescribed hypersurface parameter
space.  This proves the asserted fixed quotient.
\end{proof}

\begin{remark}\label{rem:abelian-obstruction-not-exact-group}
\status{Proved under stated hypotheses}
Already for an ordinary abelian surface there is a fixed
quotient of dimension three.  Thus even the proposed
replacement ``full stabilizer of one adjunction covector''
would be too large for general ambient varieties.
The proposition gives a lower bound for the full fixed
covector space; it does not assert that no further
invariants occur, nor determine the whole monodromy group.
\end{remark}

\begin{example}\label{ex:ordinary-abelian-surface-explicit}
\status{Proved under stated hypotheses}
Here is a completely specified ambient surface, with an executed
ordinary-hypersurface calculation.  Over $\mathbf F_2$, let
\[
 E:\quad Y^2Z+XYZ+X^3+Z^3=0,\qquad
 A=E\times E,\qquad L=\mathcal O_E(3O)\boxtimes\mathcal O_E(3O),
\]
where $O=[0:1:0]$, and embed $A$ in $\mathbf P^8$ by the
Segre embedding of these two plane cubics.  The derivatives of
the displayed cubic have no common projective zero on $E$;
its four $\mathbf F_2$-points are
$[0:1:1],[1:0:1],[1:1:1],O$.
Thus $E$ is smooth and ordinary (its Frobenius trace is $-1$).
The degree-one section
\begin{equation}\label{eq:abelian-degree-one-example}
 Z_1Z_2+Z_1Y_2+X_1Y_2+Y_1X_2=0
\end{equation}
cuts out a smooth ordinary genus-ten curve on $A$.

The exact programs \path{ordinary_abelian_surface.py} and
\path{ordinary_abelian_degree_one.py}, with their retained
JSON outputs, verify this as follows.  On $E\setminus\{O\}$
write $x=X/Z,y=Y/Z$ and $\omega=dx/x$.
The coefficient Cartier map satisfies
\[
 T(x^{2a})=x^a,\quad T(x^{2a+1})=0,\quad
 T(x^{2a}y)=x^{a+1},\quad T(x^{2a+1}y)=x^a(y+1).
\]
These formulas follow directly from
$((y+1)/x)^2=y/x+x$ and the decomposition
$k(E)=k(E)^2\oplus xk(E)^2$.
Multiplication uses $y^2=x^3+xy+1$.
On the tensor basis $(1,x,y)\otimes(1,x,y)$ the Cartier
map $g\mapsto T_1T_2(sg)$ has the following nine columns,
encoded as binary integers from the first basis vector:
\[
 (3,21,46,16,170,357,154,493,116).
\]
Its rank over $\mathbf F_2$ is nine.  Quotienting the fixed
line $ks$ gives relative rank eight.  Together with the
ordinary ambient $H^1(A,\mathcal O_A)$ of rank two, this
proves $p$-rank ten.

Smoothness is checked on all four product charts
$Z_i\ne0$ or $Y_i\ne0$.
The invariant tangent vector field has components
$(x,x^2+y)$ on the first chart and
$(1+u+v^2,u^2+v)$ on $u=X/Y,v=Z/Y$.
On each product chart the ideal generated by the two
elliptic equations, the section, and its two tangent
derivatives has Groebner basis $[1]$ over $\mathbf F_2$.
These charts cover $A$ and the vector fields never vanish,
so these certificates prove smoothness over the algebraic
closure.  The direct count is eight rational points,
consistent with the computed Cartier trace zero.

It follows that the generic degree-$d$ section on this
\emph{fixed} embedded surface is ordinary for every $d\geq1$.
Indeed choose $d$ general ordinary members of $|L|$;
they meet transversely, with no triple intersections.
Each has genus ten and each pair meets in $L^2=18$ points.
Their connected nodal union has
\[
 h^1(\mathcal O)=10d+18\binom d2-d+1=9d^2+1.
\]
The normalization sequence expresses this cohomology
as an extension of the ordinary component cohomology
by the Frobenius-bijective graph cohomology.  Its
Frobenius is therefore bijective.  Openness of that
condition and the existence of smooth members give
ordinary generic smooth sections.  Thus
\begin{equation}\label{eq:abelian-example-generic-rank}
 r_d=9d^2-1\qquad\text{for every }d\geq1.
\end{equation}
The ambient restriction map is surjective in every
degree: restriction of plane forms to each cubic is
surjective, and Segre monomials span all forms of
bidegree $(d,d)$.

For every even $d\geq2$, the four character summands
have nonzero spaces of sections (their factors on each
elliptic curve have degree $3d/2$), so the fixed quotient
of Proposition~\ref{prop:ordinary-abelian-fixed-quotient}
has dimension three with no additional degree bound.
The degree-two program also gives an explicit smooth
ordinary genus-37 section: its 36-dimensional Cartier
matrix has rank 36; the relative rank is 35, and its
four character covectors have rank four before taking
the single relation and rank three afterwards.
All defining coefficients, matrix columns, covectors,
chart certificates and the count of seven rational
points are retained in its JSON output.
\end{example}


% /root/monodromy, 19 November 2026; independently audited by /root.
\subsection{An odd-characteristic ambient obstruction in every degree}
\label{subsec:product-surface-obstruction}
\status{Proved; independent parent proof and source audit passed.}

\begin{proposition}\label{prop:product-surface-norm-obstruction}
\status{Proved under stated hypotheses}
Let $E_1,E_2$ be ordinary elliptic curves over a perfect
field of characteristic $p$, and let
\[
 A=E_1\times E_2,\qquad
 L=\mathcal O_{E_1}(aO_1)\boxtimes\mathcal O_{E_2}(bO_2),
       \qquad a,b\geq3.
\]
For every $d\geq1$ with $p\mid ad$ and $p\mid bd$, the
geometric relative Hasse--Witt monodromy of smooth members
of $|L^d|$ fixes pointwise a quotient of dimension two.
The same is true of arithmetic monodromy if both
$E_i[p]^{\et}$ are constant over the base field.
The statement applies equally to ambient degree-$d$
hypersurfaces for the product embedding by $L$.

In particular, over $\mathbf F_3$ take
\begin{equation}\label{eq:odd-product-explicit-elliptic}
 E:\quad y^2=x^3+x^2+1,
 \qquad A=E\times E,\qquad
 L=\mathcal O_E(3O)\boxtimes\mathcal O_E(3O).
\end{equation}
For this fixed smooth embedded surface, both monodromy
images have a fixed quotient of dimension two for
\emph{every} $d\geq1$.  In particular no eventual
full-linear-group assertion holds for this $X$, even
in odd characteristic.
\end{proposition}

\begin{proof}
Work first over the algebraic closure.  For a smooth
ample member $i:Y\hookrightarrow A$, put
$\pi_j=\operatorname{pr}_j|_Y$.
The restriction and Albanese maps are
\[
 j:\Pic^0(A)=E_1\times E_2\longrightarrow J_Y,
 \quad (P,Q)\longmapsto\pi_1^*P\otimes\pi_2^*Q,
\]
and
\[
 u:J_Y\longrightarrow A,
       \quad N\longmapsto
            (\operatorname{Nm}_{\pi_1}N,
             \operatorname{Nm}_{\pi_2}N),
\]
using the usual principal identifications on the elliptic
curves and the Jacobian.

The map $j$ is injective on the geometric points of
$\Pic^0(A)[p]$.  To see this directly, let
$P\boxtimes Q$ be a nontrivial degree-zero line bundle
on $A$.  It has no global sections, whereas
\[
 H^1\bigl(A,(P\boxtimes Q)\otimes L^{-d}\bigr)=0
\]
by the Kunneth formula: both factors have strictly
negative degree and hence no $H^0$.
The restriction exact sequence consequently gives
$H^0(Y,i^*(P\boxtimes Q))=0$.
Thus its restriction cannot be trivial.  Passing to
the etale $p$-torsion quotient of $J_Y$ preserves this
injection: over an algebraically closed field that
quotient induces a bijection on the geometric points
of the finite $p$-torsion group scheme.

We calculate the composite $u\circ j$.
The projection degrees are
$\deg\pi_1=bd$ and $\deg\pi_2=ad$.
The norm of a pulled-back line bundle is its power
by the degree, so the two diagonal entries are
$[bd]$ and $[ad]$.
The off-diagonal entries are zero.  For example,
represent $Q\in\Pic^0(E_2)$ as $[q]-[q_0]$.
Pushing its pullback to $Y$ along $\pi_1$ gives the
difference of the intersection divisors of $Y$ with
$E_1\times\{q\}$ and $E_1\times\{q_0\}$, viewed on
$E_1$.  Each divisor has line bundle
$\mathcal O_{E_1}(adO_1)$, because the section defining
$Y$ has bidegree $(ad,bd)$.  Their difference is
therefore trivial in $\Pic^0(E_1)$.
This calculation holds for generic $q,q_0$ and hence
for the entire homomorphism.  Thus
\begin{equation}\label{eq:product-norm-composition}
 u\circ j=\begin{pmatrix}[bd]&0\\0&[ad]\end{pmatrix}.
\end{equation}
Under the divisibility hypotheses it vanishes on
$\Pic^0(A)[p]^{\et}$.

The Artin--Schreier/Lang comparison identifies
\[
 \mathcal V_Y=\operatorname{coker}
       \bigl(H^1_{\et}(A,\mathbf F_p)
                 \longrightarrow H^1_{\et}(Y,\mathbf F_p)\bigr).
\]
For an abelian variety, the dual of $H^1_{\et}(-,\mathbf F_p)$
is its etale $p$-torsion; for a curve it is the etale
$p$-torsion of its Jacobian.  Functoriality identifies
the dual of restriction with $u[p]^{\et}$.
It follows that
\[
 \mathcal V_Y^\vee
   =\ker\bigl(J_Y[p]^{\et}\longrightarrow A[p]^{\et}\bigr)
\]
contains the injected copy of $\Pic^0(A)[p]^{\et}$
provided by $j$ and \eqref{eq:product-norm-composition}.
Both elliptic factors are ordinary, so this copy has
dimension two over $\mathbf F_p$.
These maps exist throughout the smooth family; on its
maximal-rank open they give an injection of the constant
ambient geometric local system into $\mathcal V_d^\vee$.
Thus geometric monodromy fixes a two-dimensional quotient
of $\mathcal V_d$.  If the ambient etale torsion is
constant over $k$, the same holds arithmetically.
Pulling the family back to ambient hypersurface
parameters preserves the constraint.  Smooth ambient
members form a nonempty open since $L^d$ is very ample.

For \eqref{eq:odd-product-explicit-elliptic}, the affine
partial derivative of $x^3+x^2+1$ is $2x$; at its only
zero the polynomial equals one, and the point at
infinity is smooth.  The numbers of affine points
above $x=0,1,2$ are respectively $2,1,2$.
Hence $\#E(\mathbf F_3)=6$ and the Frobenius trace
is $-2$, nonzero modulo three.  The curve is ordinary,
and its etale three-torsion is rational because the
group of six rational points contains all three
geometric points of that torsion.
Now $a=b=3$, so the divisibility conditions hold for
every $d$.  The two-dimensional fixed quotient is
incompatible with the full linear group in any
dimension at least two; its existence itself forces
the relative rank to be at least two.
\end{proof}


% /root/monodromy, 19 November 2026; independently audited by /root.
\subsection{The exact fixed quotient on products of ordinary elliptic curves}
\label{subsec:product-exact-fixed-quotient}
\status{Proved; independent parent proof and source audit passed.
This determines all invariant covectors, not the full group on their kernel.}

\begin{theorem}\label{thm:product-exact-fixed-covectors}
\status{Proved under stated hypotheses}
Let $E_1,\ldots,E_n$ be ordinary elliptic curves over a
perfect field of characteristic $p$, where $n\geq2$, and set
\[
 X=\prod_{i=1}^nE_i,\qquad
 L=\boxtimes_{i=1}^n\mathcal O_{E_i}(a_iO_i),\quad a_i\geq3,
 \qquad h_d=d^n\prod_i a_i.
\]
Use the product of the indicated complete embeddings and
the Segre embedding as the fixed projective embedding of $X$.
Then the generic relative Hasse--Witt rank is
\begin{equation}\label{eq:product-elliptic-generic-rank}
 r_d=h_d-1\qquad\text{for every }d\geq1.
\end{equation}
In particular the generic relative coherent Frobenius is
invertible in every degree.

Suppose in addition that $p\mid a_i d$ for every $i$.
For geometric monodromy on its maximal-rank open one has
\begin{equation}\label{eq:product-entire-fixed-space}
 \dim_{\mathbf F_p}(\mathcal V_{d,\bar\eta}^{\vee})^{G_{\rm geom}}
 =\begin{cases}
       2^n-1,&p=2,\\
       n,&p\text{ odd}.
   \end{cases}
\end{equation}
After a finite extension splitting the Frobenius kernel
and normalizing its invariant top Cartier differential,
the arithmetic invariant space is the same space.
Thus these numbers give the dimension of the entire
trivial quotient, rather than a lower bound for it.
\end{theorem}

\begin{proof}
First we prove generic ordinarity.  For each factor
choose a section of $\mathcal O_{E_i}(a_i dO_i)$ with
distinct zeros.  It can be chosen as a product of $d$
sections of $\mathcal O_{E_i}(a_iO_i)$ with disjoint
reduced zero sets.  The resulting product section on
$X$ cuts out an SNC union of coordinate fibres.
Every nonempty stratum is a product of ordinary elliptic
curves, so Frobenius is bijective on every coherent
cohomology group of that stratum.  In the structure-sheaf
Cech spectral sequence, Frobenius is consequently
bijective on every page and on the total cohomology.

For every nonzero defining section, the Kunneth formula
gives $H^{n-1}(X,L^{-d})=0$ and
$h^n(X,L^{-d})=h_d$.  Also $h^n(X,\mathcal O_X)=1$.
The multiplication map to $H^n(X,\mathcal O_X)$ is
surjective because its cokernel lies in
$H^n(Y,\mathcal O_Y)=0$.  Hence the relative coherent
rank is constantly $h_d-1$ on the whole projective
linear system.  Invertibility at the SNC section is
an open condition, and the smooth open is nonempty.
This proves \eqref{eq:product-elliptic-generic-rank}.

The ambient hypersurface family is the same complete
linear system up to a surjective linear projection.
For completeness, a line bundle $N$ of degree at least
three on an elliptic curve is normally generated.
For multiplication $H^0(N)\otimes H^0(N^j)\to H^0(N^{j+1})$
with $j\geq2$, a base-point-free pencil in $H^0(N)$
and $H^1(N^{j-1})=0$ prove surjectivity.
For $j=1$, choose a point $q$ and a base-point-free
pencil in $H^0(N(-q))$.  The pencil exact sequence
after tensoring with $N$ has kernel $\mathcal O(q)$,
whose $H^1$ vanishes, so its products span
$H^0(N^2(-q))$.  The square of a section nonzero at
$q$ spans the remaining one-dimensional quotient.
This proves normal generation.  Tensoring the factors
and grouping Segre monomials proves the required
surjectivity for $X$ in every degree.

We now work over the algebraic closure.  Put
$W=H^0(X,L^d)$ and $P=\mathbf P(W)$, and normalize an
invariant top form $\omega$ by $C(\omega)=\omega$.
The coefficient can be normalized since $X$ is ordinary
and the base field is algebraically closed.
The relative coherent bundle fits in the exact sequence
\begin{equation}\label{eq:product-global-bundle}
 0\longrightarrow\mathcal E_d\longrightarrow
       H^n(X,L^{-d})\otimes\mathcal O_P(-1)
       \longrightarrow\mathcal O_P\longrightarrow0.
\end{equation}
It is a vector bundle on all of $P$, not only on the
ordinary or smooth locus.  Frobenius extends there by
$F_s=s^{p-1}F_X$.

Any geometric monodromy-invariant covector corresponds,
on the invertibility open, to a Frobenius-equivariant
map $\lambda:\mathcal E_d\to\mathcal O$.
It extends over all of $P$.  Indeed, in a regular local
frame its equation is $A^t\lambda=\lambda^{[p]}$ with
$A$ regular.  If its largest pole along any prime
divisor had order $m>0$, the right side would have
a coordinate of valuation $-pm$, whereas every coordinate
on the left has valuation at least $-m$, a contradiction.
There are no divisorial poles; normality of $P$ and local
freeness extend the section also across codimension two.

Dualizing \eqref{eq:product-global-bundle} and taking
global sections, using $H^1(P,\mathcal O_P)=0$, shows
that such a covector is represented by
\begin{equation}\label{eq:product-linear-lift}
 w_s=D(s)\omega,\qquad D\in\End(W),
\end{equation}
modulo scalar multiples of the identity.
Its Frobenius defect lies in the constant one-dimensional
kernel in the Frobenius pullback of the dual sequence.
It is therefore $c s^{[p]}$ for a constant $c$.
Adding $\alpha\operatorname{id}_W$ to $D$ changes this
constant by $\alpha-\alpha^p$; an Artin--Schreier equation
over the algebraically closed field removes it.
Thus every invariant covector lifts to a $D$ satisfying
\begin{equation}\label{eq:product-exact-Cartier-operator}
 C\bigl(s^{p-1}D(s)\omega\bigr)=D(s)\omega
            \quad\text{for all }s\in W.
\end{equation}
Two such exact lifts differ by $c\operatorname{id}$
with $c\in\mathbf F_p$.

The divisibility assumptions give
$L^d=F^*M$ with
$M=\boxtimes_i\mathcal O_{E_i^{(p)}}((a_i d/p)O_i^{(p)})$.
Since $\ker F\simeq\mu_p^n$, its character algebra gives
\begin{equation}\label{eq:product-character-W}
 W=\bigoplus_{\chi\in H}W_\chi,
 \qquad H\simeq\mathbf F_p^n,
 \qquad W_\chi=H^0(X^{(p)},M\otimes\mathcal N_\chi).
\end{equation}
Every $W_\chi$ is nonzero: each factor is a positive
degree line bundle on an elliptic curve.
The normalized Cartier trace is projection to the
trivial character followed by taking the Frobenius
root.  Nontrivial character terms are killed.

Apply \eqref{eq:product-exact-Cartier-operator} to a
section $s\in W_\chi$, and put
$T=\operatorname{pr}_\chi\circ D|_{W_\chi}$.
Only this component contributes to Cartier.  Raising
the equation to the $p$-th power gives
\[
 D(s)^p=s^{p-1}T(s).
\]
We do not assume that the Cartier root preserves the
character space.  Instead choose a nonzero eigenvector
$s_0$ of the finite-dimensional operator $T$, say
$T(s_0)=a s_0$.  The displayed equation implies
$D(s_0)=a^{1/p}s_0$; projecting to $W_\chi$ then gives
$a=a^{1/p}$, so $a\in\mathbf F_p$.
For arbitrary $s\in W_\chi$, substitute $s_0+z s$
and compare the coefficient of $z$.
The left side has no such coefficient, and the right
side gives $s_0^{p-1}(T(s)-a s)=0$.
Since $X$ is integral, $T(s)=a s$.
The same equation now gives $D(s)^p=a s^p$, hence
$D(s)=a s$.  Thus $D$ does preserve every character
space and acts there by a scalar in $\mathbf F_p$:
\begin{equation}\label{eq:product-character-scalar}
 D|_{W_\chi}=c_\chi\operatorname{id},\qquad c_\chi\in\mathbf F_p.
\end{equation}

It remains to determine the allowed functions
$c:H\to\mathbf F_p$.
Choose one nonzero section in each character space
and vary its coefficient independently.
For a nonpure multi-index $(m_\chi)$ with
$\sum m_\chi=p$, the coefficient of the corresponding
monomial in $s^{p-1}D(s)$ is the nonzero factor
$(p-1)!/\prod m_\chi!$ times
$\sum m_\chi c_\chi$.
Cartier kills that monomial unless
$\sum m_\chi\chi=0$ in $H$; if this weight is zero,
the product of the chosen nonzero sections is nonzero
and its Frobenius root is nonzero.
Independence of the coefficient monomials therefore
shows that the exact conditions are
\begin{equation}\label{eq:product-scalar-relations}
 \sum_\chi m_\chi c_\chi=0
 \quad\text{whenever}\quad
 \sum m_\chi=p,\quad\sum m_\chi\chi=0,
 \quad(m_\chi)\text{ is not pure}.
\end{equation}
Pure monomials already give the right side of
\eqref{eq:product-exact-Cartier-operator}.

When $p=2$, every zero-weight pair consists of two
equal characters, so there are no conditions in
\eqref{eq:product-scalar-relations}.  The space of
exact lifts has dimension $2^n$ over $\mathbf F_2$.
When $p$ is odd, put $a_x=c_x-c_0$.
The multiset consisting of $p-2$ zeros, $x$, and $-x$
gives $a_{-x}=-a_x$.
The multiset consisting of $p-3$ zeros, $x$, $y$, and
$-x-y$ then gives $a_{x+y}=a_x+a_y$.
The only nontrivial pure exception is $p=3$, $x=y$;
there the same additivity follows already from
$a_{-x}=-a_x$, since $2x=-x$.
Thus $a$ is an $\mathbf F_p$-linear functional on $H$.
Conversely every affine-linear function
$c_x=c_0+\ell(x)$ satisfies all displayed relations.
The space of exact lifts has dimension $n+1$.

Finally quotient the exact lifts by the one-dimensional
space $\mathbf F_p\operatorname{id}_W$.
This gives precisely \eqref{eq:product-entire-fixed-space}.
The character projectors, their indicated combinations,
and the normalized top differential are defined over
a finite field extension when the original base is
finite; over any perfect base they descend after a
finite algebraic extension.  Over that extension the
whole geometric invariant space is already arithmetic.
\end{proof}

\begin{corollary}\label{cor:product-two-explicit-fixed-spaces}
\status{Proved under stated hypotheses}
For the explicit surface over $\mathbf F_3$ in
\eqref{eq:odd-product-explicit-elliptic}, the relative
rank is $9d^2-1$ and the full geometric and arithmetic
invariant covector spaces both have dimension two,
for every $d\geq1$.
For the surface over $\mathbf F_2$ in
Example~\ref{ex:ordinary-abelian-surface-explicit}, the
relative rank is also $9d^2-1$, and both invariant
covector spaces have dimension three for every even
$d\geq2$.
The exact full images inside the pointwise stabilizers
of these quotients are not determined by this theorem.
\end{corollary}


% /root/monodromy, 19 November 2026; independently audited by /root.
\subsection{A terminating finite-field algorithm for the entire monodromy image}
\label{subsec:effective-entire-monodromy}
\status{Proved; independent parent proof audit and exact quartic computation passed.
This is an effective answer for specified finite-field inputs, not a structural
classification for arbitrary perfect base fields.}

\begin{theorem}\label{thm:effective-monodromy-matrices}
\status{Proved under stated hypotheses}
Suppose $k=\mathbf F_q$ and the fixed smooth projective
embedding $X\subset\mathbf P^m_k$ is given by homogeneous
equations.  For every
\[
 d\geq d_{\rm coh}=\max\{1,\operatorname{reg}(\omega_X)-1\},
\]
there is an explicit terminating algorithm that outputs,
up to a simultaneous choice of basis, the complete finite
lists of matrices
\[
 G_{\rm geom},G_{\rm arith}\subset\GL_{r_d}(\mathbf F_p)
\]
for the local system in A.4.  It also decides geometric
and arithmetic irreducibility and outputs the determinant
images.  The algorithm uses only finite linear algebra,
Groebner-basis computations, construction of finite fields,
and bounded enumeration of polynomials and matrices.
It does not call an oracle for Galois groups or for
factorization over an algebraic closure.
\end{theorem}

\begin{proof}
We give the steps and their termination proofs.

\emph{Step 1: the stable matrix.}
Use Proposition~\ref{prop:effective-generic-rank} over
the rational function field $K=\mathbf F_q(t_1,\ldots,t_M)$
of an affine chart of the ambient degree-$d$ parameter
space.  Let $B$ be its coherent Frobenius matrix and
$b$ the coherent dimension.  Compute the image of
$B B^{[p]}\cdots B^{[p^{b-1}]}$ by Gaussian elimination.
The Frobenius-linearized map on that image is invertible;
a chosen basis gives an explicit $A\in\GL_r(K)$, where
$r=r_d$.  If $r=0$, output the trivial groups and the zero
representation (which is not called irreducible), and stop.
This construction is valid even though $K$ is imperfect.

\emph{Step 2: the frame algebra.}
Let $T=(T_{ij})$ be an $r\times r$ matrix of indeterminates
and introduce one more indeterminate $z$.  Form
\begin{equation}\label{eq:effective-frame-algebra}
 R=K[T_{ij},z]/
        (A T^{[p]}-T,\ z\det T-1),\qquad
 D=\prod_{i=0}^{r-1}(p^r-p^i)=|\GL_r(\mathbf F_p)|.
\end{equation}
This is a finite etale $K$-algebra of dimension $D$.
Indeed the equations for one column, $Av^{[p]}-v=0$,
have Jacobian $-1$.  Their leading homogeneous forms
are an invertible linear combination of the $p$-th
powers of the coordinates, so they have no common
projective zero at infinity and define a scheme of
length $p^r$.  Its geometric points form an
$\mathbf F_p$-space of dimension $r$ and span the
whole $r$-dimensional vector space: otherwise their
span, of dimension $s<r$, would itself have only
$p^s$ fixed points by the same argument.
The invertible frames are thus a torsor under
$\GL_r(\mathbf F_p)$ and number exactly $D$.
Compute a Groebner basis of the displayed ideal,
a $K$-basis of $R$, and its multiplication table.

\emph{Step 3: a primitive element from a finite list.}
Let $e_1,\ldots,e_D$ be the computed basis and put
$N=D(D-1)/2$.  Choose the finite set
$S=\{1,t_1,t_1^2,\ldots,t_1^N\}\subset K$.
Test the elements
\[
 \theta=\sum_{j=1}^D a_j e_j,\qquad (a_1,\ldots,a_D)\in S^D,
\]
until the vectors $1,\theta,\ldots,\theta^{D-1}$ are
linearly independent.  At least one test succeeds.
Over a separable closure, the algebra has $D$ distinct
embeddings.  The product of the pairwise differences
of their values at $\sum a_je_j$ is a nonzero polynomial
of degree $N$ in the $a_j$.  A nonzero polynomial of
degree at most $N$ cannot vanish on a grid having more
than $N$ distinct values in every coordinate, by
induction on the number of variables.  The successful
$\theta$ has a square-free monic minimal polynomial
$f\in K[Z]$ of degree $D$, computed by linear algebra,
and $R=K[\theta]\simeq K[Z]/(f)$.

\emph{Step 4: factorization by bounded enumeration.}
Clear the denominators of $f$ to obtain a nonzero
$F\in\mathbf F_q[t_1,\ldots,t_M,Z]$.
Enumerate all normalized polynomials over $\mathbf F_q$
of total degree at most $\deg F$, and test divisibility
by polynomial division or Groebner reduction.
This is a finite list.  The irreducible factors of
$F$ are among it, since total degrees add in a
polynomial domain.  Repeated exact division therefore
factors $F$ completely.  Discard factors independent
of $Z$, which are units over $K$, and normalize the
remaining ones to be monic in $Z$.
Gauss's lemma proves that these are precisely the
irreducible factors of $f$ over $K$.
Thus this step is a fully specified finite algorithm;
practical multivariate factorization can replace the
enumeration without changing its correctness.

Choose one factor $f_0$ by a fixed lexicographic rule.
It defines a field component $R_0$ of $R$ and the
corresponding maximal ideal
\[
 I_0=(A T^{[p]}-T,\ z\det T-1,\ f_0(\theta))
       \subset K[T_{ij},z].
\]

\emph{Step 5: the arithmetic matrices.}
Enumerate all $g\in\GL_r(\mathbf F_p)$.
Its action on the frame algebra is the substitution
\[
 T\longmapsto Tg,\qquad z\longmapsto(\det g)^{-1}z.
\]
Retain exactly those $g$ for which $f_0(\theta)$ after
this substitution reduces to zero modulo $I_0$.
This is a finite collection of Groebner reductions.
The retained list is the stabilizer of the chosen
field component, hence exactly $G_{\rm arith}$ in
the corresponding frame.  Preservation of the one
extra equation suffices because the frame ideal is
already preserved; inclusion between the resulting
maximal ideals is equality.

To justify the monodromy identification explicitly,
choose a geometric frame in the chosen component.
Its Galois orbit is the image of $\operatorname{Gal}(K)$
acting on the frames.  Under the simply transitive
right frame action, the stabilizer of its connected
component is the same subgroup of $\GL_r(\mathbf F_p)$,
up to the common initial choice of frame.
Restricting the original local system to any dense
open does not change its image: the fundamental group
of a dense open in a normal connected scheme surjects
onto that of the scheme.  Thus the generic calculation
computes the image over the stated $U_d$.

\emph{Step 6: the geometric matrices with a finite
constant-field bound.}
Construct the finite field $k_D=\mathbf F_{q^D}$ and
put $K_D=k_D(t_1,\ldots,t_M)$.
A finite field of this degree is constructed by
enumerating monic degree-$D$ polynomials and testing
irreducibility; the elementary existence theorem
for finite fields guarantees termination.
Factor $f_0$ over $K_D$ by Step 4, now over $k_D$,
choose one factor $f_{00}$, and repeat the matrix
preservation test with its component ideal.
The resulting subgroup is $G_{\rm geom}$, and it is
a subgroup of the already chosen arithmetic image.

Here is the reason the finite extension degree $D$
suffices.  The selected field component is a connected
torsor under its stabilizer $H$, so its degree over
$K$ is $|H|$, a divisor of $D$.  Its full algebraic
constant field is some $\mathbf F_{q^c}$, where
$c\mid |H|$, hence $c\mid D$.
After extending to $k_D$, every resulting field
component has its constant field algebraically closed
in it.  A finitely generated field over a perfect
field is separably generated; therefore each such
component is geometrically integral.
Consequently these are exactly the geometric
components, and their stabilizers give the geometric
image, not an upper approximation to it.

\emph{Step 7: irreducibility and determinant.}
For each of the two finite matrix lists, enumerate
all row-echelon matrices representing the proper
nonzero subspaces of $\mathbf F_p^r$ and test
invariance under every listed element.
There is an invariant subspace if and only if the
representation is reducible.  Taking determinants
of the listed matrices gives both determinant
images.  Proposition~\ref{prop:exact-general-determinant-image}
is a much smaller independent calculation of the
same images and provides a consistency check.

Every step has a finite bound or an explicitly proved
termination criterion.  This proves the theorem.
\end{proof}

\begin{remark}\label{rem:effective-monodromy-scope}
\status{Proved under stated hypotheses}
The theorem outputs the whole finite group for each
specified finite-field input, with a proved finite
algorithm.  Its worst-case complexity is enormous:
even the frame dimension is $|\GL_r(\mathbf F_p)|$.
It does not supply a uniform structural description
of the groups as $X$ and $d$ vary.  Nor does it supply
a uniform algorithm for arithmetic monodromy over
an arbitrary perfect field presented without effective
field operations and factorization data.
Those distinctions are retained in the result ledger.
\end{remark}

\begin{example}\label{ex:effective-quartic-frame-algebra}
\status{Proved under stated hypotheses}
The retained program \path{lang_frame_quartic.py}
executes a nontrivial small frame calculation for a genuine
smooth plane-quartic family.  Over $\mathbf F_2(t)$ take
\[
 f_t=x^4+y^4+z^4+x^2yz+xz^3+txy^3+xyz^2.
\]
In the standard genus-three monomial basis its Hasse matrix is
\[
 A_t=\begin{pmatrix}1&0&0\\0&0&1\\0&t&1\end{pmatrix}.
\]
The program extracts these coefficients from $f_t$ and checks
that all three affine-chart singularity ideals have Groebner
basis $[1]$ over $\mathbf F_2(t)$.
Thus the generic curve is smooth and ordinary, since
$\det A_t=t\ne0$.

The first coordinate is a constant summand.  A frame of the
remaining rank-two Lang module has columns $(c^2,c)^t$ and
$(d^2,d)^t$, with algebra
\[
 R=\mathbf F_2(t)[c,d]/
 \left(c^3+\frac{c+1}{t},\quad
              d^2+cd+c^2+\frac1t\right).
\]
Its basis is $1,c,c^2,d,cd,c^2d$, and its frame determinant
is $cd(c+d)=1/t$.
The program verifies all six substitutions from
$\GL_2(\mathbf F_2)$ preserve this algebra.
It also finds the primitive element $\theta=c+td$,
whose minimal polynomial, after clearing denominators, is
\[
 t^2 Z^6+(t^3+t^2)Z^3+(t^4+t^2+1)Z^2
       +(t^4+t)Z+t^6+t^5+1.
\]
Linear independence of $1,\theta,\ldots,\theta^5$ and
annihilation by this polynomial are checked exactly.

Its geometric irreducibility has a direct certificate.
The first equation gives $t=(c+1)/c^3$, a separable
degree-three map of rational function fields.
Setting $u=d/c$ transforms the second equation into
\[
 u^2+u=\frac1{c+1}.
\]
The right side has a simple pole and is not an
Artin--Schreier coboundary over $\overline{\mathbf F}_2(c)$:
every pole of $v^2+v$ has even order.
Thus the frame algebra remains a field of degree six
after extending constants to the algebraic closure.
Both monodromy images of this particular curve family
are exactly
$\operatorname{diag}(1,\GL_2(\mathbf F_2))\subset\GL_3(\mathbf F_2)$.
This is a restricted one-parameter family, not the full
space of plane quartics, whose image was determined
earlier.  Its exact coefficient matrix, basis, six group
elements, minimal polynomial and smoothness certificates
are retained in \texttt{lang\_frame\_quartic.json}.
\end{example}


% /root/monodromy with suggestions by /root, 19 November 2026.
\subsection{Finite constant data for an arbitrary perfect base field}
\label{subsec:perfect-field-monodromy-data}
\status{Proved; independent parent proof and original-source audit passed.
The output is finite polynomial data and matrix lists, rather than a
uniform structural classification of the resulting finite groups.}

\begin{lemma}\label{lem:effective-separable-factorization}
\status{Proved under stated hypotheses}
Let $k_0$ be a finitely generated field over $\mathbf F_p$,
given by a finite field presentation.  Arithmetic and
factorization of separable univariate polynomials over
$k_0(t_1,\ldots,t_M)$ are terminating exact procedures.
For a polynomial that defines a finite separable field
extension of this rational function field, its absolutely
irreducible factors over $\bar k_0(t_1,\ldots,t_M)$ can
also be found by a terminating procedure, with their
coefficients in explicitly presented finite separable
extensions of $k_0$.
\end{lemma}

\begin{proof}
We give details so that no absolute-factorization or
Galois-group oracle is concealed in the assertion.

A separating transcendence basis of $k_0/\mathbf F_p$
and a finite separable primitive presentation over the
corresponding rational function field can be found
by elimination, linear algebra and enumeration.
Explicitly, enumerate candidate rational functions
for a transcendence basis, check algebraic independence
by elimination, and check that the remaining algebra
is finite by a zero-dimensional Groebner computation.
Separability is checked by the rank of the differentials,
or by the nonzero derivative of a primitive minimal
polynomial.  Such a candidate exists because the
ground field is perfect.  For a finite separable
extension, the primitive element is found by testing
linear combinations of a finite basis, as in the
finite-grid argument of the previous theorem.
Thus we may write
\[
 k_0(t)=R(\alpha),\qquad
 R=\mathbf F_q(u_1,\ldots,u_s,t_1,\ldots,t_M),
\]
with a separable minimal polynomial for $\alpha$.
All arithmetic is then polynomial reduction and linear
algebra over a rational function field.

Here is a direct factorization procedure over $R(\alpha)$.
For a monic separable polynomial $f(Z)$, try translates
$g_c(Z)=f(Z+c\alpha)$ for successively enumerated
$c\in R$, until
\[
 N_c(Z)=\operatorname{Norm}_{R(\alpha)/R}(g_c(Z))
\]
is square-free.  Its norm is computed by a resultant.
Only finitely many $c$ are bad: over an algebraic
closure its roots are $\beta-c\sigma(\alpha)$, where
$\beta$ ranges over the roots of $f^\sigma$ and
$\sigma$ over the distinct embeddings of $R(\alpha)$.
Within one embedding they are distinct; across two
different embeddings an equality excludes at most
one value of $c$.  The rational function field $R$
is infinite here.  If there are no transcendental
variables, the field is finite and bounded polynomial
enumeration directly factors $f$ instead.

Factor $N_c$ over $R$ by clearing denominators and
the bounded finite-field polynomial enumeration in
Step 4 of Theorem~\ref{thm:effective-monodromy-matrices}.
For each irreducible factor $h$ of $N_c$, compute
$\gcd(g_c,h)$ over $R(\alpha)$ and discard the units.
These gcds are irreducible and give the factors of
$g_c$.  To verify irreducibility, two roots of $g_c$
in the same $R$-Galois orbit are carried to one another
by an automorphism fixing $R(\alpha)$: otherwise the
target root would occur in two different conjugates
of $g_c$, contradicting square-freeness of $N_c$.
Undoing the translate factors $f$.  Thus all separable
factorization steps used below are effective.

For absolute factorization, clear denominators of the
given separable polynomial and remove its content
in $k_0[t]$, obtaining a primitive polynomial
$F(t,Z)$.  Its total degree bounds the total degree
of every polynomial factor over $\bar k_0$.
For each possible pair of leading monomials and
factor degrees, write two polynomials with unknown
coefficients, normalize the leading coefficient of
one to one, and equate their product with $F$.
Impose the selected $Z$-degree of the proposed factor
to be positive and require the selected leading
coefficients to be nonzero (an inverse variable
imposes this as an equation).  Permit the complementary
factor to be the unit polynomial; in particular the
full-degree factor is always one of the cases.
There are finitely many cases.  The coefficient
ideal is proper if and only if that case has a
factorization over the algebraic closure, by the
Nullstellensatz.  Properness is decided by whether
its Groebner basis contains one.

Select the least positive $Z$-degree for which a
factor exists.  A factor of that degree is absolutely
irreducible: a nontrivial splitting would yield a
smaller positive $Z$-degree.  There are no factors
independent of $Z$, by primitivity and Gauss's lemma.
Thus a finite list of coefficient tests proves what
degree to seek; it is not a search whose negative
outcome is undecidable.

We can find the factor's coefficients explicitly.
The field defined by the original irreducible factor
is separably generated over $k_0$.  Its algebraic
constant field is finite separable, and after adjoining
it its components are geometrically integral.
Indeed finiteness and geometric irreducibility are
\cite[Lemmas~10.47.8 and 10.47.13]{Stacks037P,Stacks037Q};
geometric reducedness follows from separability by
\cite[Lemmas~10.44.2 and 10.44.3]{Stacks030W,Stacks030X}.
Consequently the normalized absolute factors have
coefficients in finite separable extensions of $k_0$.
Enumerate separable irreducible polynomials over
$k_0$, testing irreducibility by the procedure above,
and enumerate coefficient tuples in their finite
extensions.  Dovetail these enumerations and test
the finitely many coefficient equations.  A solution
exists in one such extension, so the search terminates.
Repeat division and the finite coefficient tests
until all factors have been found; their positive
$Z$-degrees sum to the original degree.
This proves the absolute-factorization assertion
without assuming a preexisting algebraic closure
computing oracle.
\end{proof}

\begin{theorem}\label{thm:perfect-field-exact-monodromy-data}
\status{Proved under stated hypotheses}
Let $k$ be any perfect field of characteristic $p$,
and let the fixed smooth projective embedding $X/k$
be defined over a finitely generated subfield
$k_0\subset k$.  Fix $d\geq d_{\rm coh}$.
There is a terminating construction from a finite
presentation of this field of definition and the
equations of $X$ of the following finite data:
\[
 H_{\rm geom}\triangleleft H_0\subset\GL_{r_d}(\mathbf F_p),
 \qquad E/k_0\text{ finite Galois},\qquad
 \rho:H_0\twoheadrightarrow\operatorname{Gal}(E/k_0),
 \quad\ker\rho=H_{\rm geom}.
\]
The two groups are output as complete lists of matrices,
$E$ by a finite field presentation, and $\rho$ by its
action on generators of $E$.
In a compatible geometric frame, the requested images
over $k$ are exactly
\begin{equation}\label{eq:perfect-field-monodromy-formula}
 G_{\rm geom}=H_{\rm geom},\qquad
 G_{\rm arith}=\rho^{-1}
       \bigl(\operatorname{Gal}(E/(E\cap k))\bigr).
\end{equation}
The intersection is taken using the chosen embedding
of $E$ into a separable algebraic closure of $k_0$
inside an algebraic closure of $k$.
Changing the initial frame conjugates all the displayed
matrix groups simultaneously.

Geometric irreducibility is a finite subspace test on
$H_{\rm geom}$.  Arithmetic irreducibility is the same
test on the explicit preimage in
\eqref{eq:perfect-field-monodromy-formula}; the only
additional base-field information is the finite
intersection $E\cap k$.
\end{theorem}

\begin{proof}
The cohomological matrix and stable-image algorithm
of Proposition~\ref{prop:effective-generic-rank} work
over $k_0(t)$ as well: they use rational linear
algebra and Frobenius linearization, not perfection
of this intermediate field.  Form the frame algebra
and find a primitive element as before.
Use Lemma~\ref{lem:effective-separable-factorization}
to factor its polynomial over $k_0(t)$ and choose
one field component $L$.
The finite list of matrix substitutions preserving
that component is $H_0$.
It acts as the full Galois group of $L/k_0(t)$:
the component is a connected torsor under its
stabilizer, whose order equals its degree.

Let $E$ be the algebraic closure of $k_0$ in $L$.
It is finite separable by the field facts cited
in the preceding lemma, and it is Galois over $k_0$
because $L/k_0(t)$ is Galois.
The field $L/E$ is separably generated with $E$
algebraically closed in it, hence geometrically
integral by the same cited lemmas.
It follows that the components after extending
constants to $\bar k_0$ are indexed by the embeddings
of $E$; their stabilizer in $H_0$ is
\[
 H_{\rm geom}=\operatorname{Gal}(L/E(t)).
\]
The restriction map to $E$ is $\rho$.
Both matrix lists are obtained by the finite ideal
preservation tests, now using one absolute factor
for $H_{\rm geom}$.

We explain how the field $E$ itself is recovered,
rather than leaving it as an unspecified relative
algebraic closure.  Let $f_0$ be the primitive
polynomial of $L$, and choose a monic absolutely
irreducible factor $f_{00}$ over $\bar k_0(t)$.
Normalize every rational-function coefficient by
coprime numerator and denominator with one leading
coefficient equal to one.
Let $E_0$ be the field over $k_0$ generated by all
the scalar coefficients in these normalized expressions.
The factorization algorithm gives them in a finite
explicit extension, so $E_0$ is computed by adjoining
them and closing their vector span under multiplication.
This terminates inside that finite extension.

In fact $E_0=E$ with the selected embedding.
Since $L/E$ is geometrically integral, the geometric
components are indexed by the embeddings of $E$.
The stabilizer of one component under
$\operatorname{Gal}(\bar k_0/k_0)$ is exactly
$\operatorname{Gal}(\bar k_0/E)$.
On the other hand, that automorphism fixes the
chosen factor precisely when it fixes all of its
normalized coefficients, so this stabilizer is
$\operatorname{Gal}(\bar k_0/E_0)$.
The finite separable fields are equal.
Thus $E$ has been explicitly obtained.

Its embedding into $L$ and the restriction map
$\rho$ can also be calculated by linear algebra.
The selected factor realizes $L$ inside
$E(t)[Z]/(f_{00})$.
This latter field has degree
$[E:k_0]\deg f_{00}=[L:k_0(t)]$ over $k_0(t)$,
so it is the same field $L$.
Express the computed generators of $E$ in the
basis $1,\theta,\ldots,\theta^{[L:k_0(t)]-1}$
by solving a linear system.  Apply every retained
matrix substitution to these expressions and
reduce.  This outputs $\rho$ explicitly.

Finally extend the constants to $k$.
Put $F=E\cap k$.  Since $E/k_0$ is finite Galois,
$\operatorname{Gal}(Ek/k)=\operatorname{Gal}(E/F)$.
Since $L/E$ is geometrically integral, extending
$E$ to $Ek$ does not change its geometric degree
or create new components.  The component over
$k(t)$ therefore has degree
\[
 |H_{\rm geom}|[E:F],
\]
and its automorphism group is exactly the preimage
of $\operatorname{Gal}(E/F)$ under $\rho$.
This proves \eqref{eq:perfect-field-monodromy-formula}.
The geometric image is unchanged by any extension
of the algebraically closed constant field, giving
the first formula.  Subspace enumeration gives the
irreducibility assertions exactly as before.
\end{proof}

\begin{remark}\label{rem:perfect-field-data-scope}
\status{Proved under stated hypotheses}
The theorem is an exact finite-data description for
an arbitrary perfect field, and an algorithm whenever
the finite field-of-definition presentation and the
finite constant intersection can be evaluated.
It makes no claim that membership in an arbitrary
unpresented, possibly noncomputable perfect field
can be decided by a universal program.
For a finite base field, the intersection is computed
by finite-field arithmetic and the previous theorem
already gives a complete unconditional algorithm.
The structural question of identifying these outputs
uniformly as familiar finite groups is separate;
for projective space it was answered above, whereas
for general ambient varieties it is not inferred
from this construction.
\end{remark}


% /root/monodromy, 19 November 2026; independently audited by /root.
\subsection{Two exact full monodromy groups on a product surface}
\label{subsec:product-full-groups}
\status{Proved; independent parent proof, primary-source audit and exact-program rerun passed.}

The following computation goes beyond determining invariant covectors.
It determines the entire images in two degrees for one fixed embedded
ordinary abelian surface.  Its use of a projective parameter space is
essential: it permits a short, checkable certificate for all divisorial
inertia, including divisors whose general point parametrizes a singular
curve.

\begin{lemma}[A line certificate for generic stable defect one]
\label{lem:hasse-line-certificate}
\status{Proved under stated hypotheses}
Let $M(s)$ be a $b\times b$ matrix of homogeneous linear forms over
$\mathbf F_2$ in homogeneous coordinates on $P=\mathbf P^m$.
Suppose $h(s)=\det M(s)$ is nonzero.  Let $a,b_0$ be coefficient
vectors with $\det M(b_0)=1$.  If, at every root $\alpha$ of
$\det M(a+t b_0)$ over $\overline{\mathbf F}_2$, one has
\[
 \rank\bigl(M(a+\alpha b_0)^{[2]}M(a+\alpha b_0)\bigr)=b-1,
\]
then at the geometric generic point of every irreducible component
of $h=0$ the semilinear map with matrix $M(s)^t$ has stable rank
$b-1$.
\end{lemma}

\begin{proof}
The determinant is homogeneous of degree $b$, and its restriction to
the indicated affine line has degree exactly $b$.  For every geometric
irreducible homogeneous factor $q$ of $h$, one has $q(b_0)\ne0$.
Consequently $q(a+t b_0)$ has its full positive degree and has a root
over the algebraic closure.  If the rank of $M^{[2]}M$ were below
$b-1$ generically on $q=0$, then $q$ would divide every $(b-1)$-minor
of that matrix.  All these minors would vanish at every root of
$q(a+t b_0)$, contrary to the hypothesis.  Thus the first and second
Frobenius images both have dimension $b-1$ at that generic point.
Over a perfect geometric field their equality makes this image
Frobenius-stable and bijective.  This is the asserted stable rank.
The argument does not require the determinant divisor to be reduced.
\end{proof}

\begin{lemma}[Rank-one loss and transvection generation]
\label{lem:char2-defect-one-inertia}
\status{Proved under stated hypotheses}
Let $\mathcal E$ be a vector bundle of rank $r$ on a projective space
$P$ over an algebraically closed field of characteristic two, with
a Frobenius-linear endomorphism.  Suppose this map is generically
bijective, and at the generic point of every component of its
noninvertibility divisor its stable rank is $r-1$.
Then the geometric image of the Lang local system on the
invertibility open is generated by transvections.
\end{lemma}

\begin{proof}
At such a divisor, pass to the strict henselization of the discrete
valuation ring.  The Lang scheme $x-F(x)=0$ is separated, affine,
quasi-finite and etale: its Jacobian is the identity, and its
geometric fibres are finite.  Its finite etale part has rank
$2^{r-1}$ and is a subgroup, by
Lemma~\ref{lem:boundary-Lang-finite-part}, or directly by
\cite[Lemma~41.20.5]{Stacks0BTN} and
\cite[Lemma~58.8.2]{Stacks09ZS}.
Since the local ring is strictly henselian, inertia fixes its
$(r-1)$-dimensional generic fibre pointwise.  The quotient has
dimension one over $\mathbf F_2$, so inertia also acts trivially on
that quotient.  Every nonidentity inertia element therefore has
the form $1+u\otimes\lambda$, where $u$ lies in the fixed hyperplane
and $\lambda$ annihilates it.  Its difference from the identity has
rank one and square zero: it is a transvection.

Let $G$ be the geometric image and divide it by the normal subgroup
generated by these divisorial inertia groups.  Normalize $P$ in the
resulting finite Galois extension.  This finite normal cover is
unramified in codimension one and hence is etale everywhere by
purity, \cite[Lemma~58.21.4]{Stacks0BMB}.  Projective space over an
algebraically closed field has no nontrivial connected finite etale
cover, so the quotient is trivial.

For clarity, the last fact can also be proved without a fundamental
group calculation.  For $\mathbf P^1$ it is Riemann--Hurwitz.
For higher-dimensional projective space, pull a cover back to the
blowup of a point.  This blowup is a $\mathbf P^1$-bundle over
$\mathbf P^{m-1}$.  On each fibre the pulled-back cover is trivial;
cohomology and base change applied to its structure sheaf identifies
it with the pullback of a finite etale cover of $\mathbf P^{m-1}$.
The exceptional section is mapped to the original point, so the
restriction of the cover to that section is trivial.  Thus the
cover of $\mathbf P^{m-1}$, and then the pulled-back cover, is
trivial.  The blowup is birational, so the original connected
function-field extension also has degree one.
\end{proof}

\begin{lemma}[A small Frobenius certificate for a transvection group]
\label{lem:product-frobenius-group-test}
\status{Proved under stated hypotheses}
Let $N$ be either $8$ or $32$, and let
$H\triangleleft G\subseteq\GL_N(\mathbf F_2)$.
Suppose $H$ is nontrivial and generated by transvections, and
$G$ contains an element whose characteristic polynomial is
irreducible of degree $N$ and is not reciprocal.  Then
$H=G=\GL_N(\mathbf F_2)$.
\end{lemma}

\begin{proof}
The irreducible polynomial makes $G$ irreducible.  The socle of its
restriction to $H$ is nonzero and $G$-stable, so it is the whole
space.  The isotypic $H$-summands are permuted transitively by $G$.
If a simple constituent occurred with multiplicity $e$, the rank
of $h-1$ on its isotypic summand would be a multiple of $e$.
A nontrivial generating transvection has total rank one, so
$e=1$ on one summand and, by transitivity, on all of them.
Thus the space is a direct sum of simple $H$-modules of a common
dimension $c$, and $G$ embeds in
$\GL_c(\mathbf F_2)\wr S_{N/c}$ if there is more than one summand.
This is precisely the elementary argument of
Lemma~\ref{lem:transvection-clifford-blocks}; it includes the
possibility of imprimitive groups rather than assuming it away.

For $N=8$, an irreducible element has order divisible by $17$:
otherwise its order divides $15=2^4-1$, contradicting its degree.
For $N=32$, its order is divisible by $65537$: otherwise it divides
$65535=2^{16}-1$.  Here $65537$ is prime, as direct trial division
up to $257$ verifies, and
\[
 \operatorname{ord}_{17}(2)=8,\qquad
 \operatorname{ord}_{65537}(2)=32.
\]
The indicated prime does not divide $|\GL_c(\mathbf F_2)|$ for
$c<N$, nor $|S_{N/c}|$.  It therefore cannot divide the order of
the displayed proper wreath product.  Thus $H$ is irreducible.

We now use exactly one external finite-group classification:
\cite[Theorem~3.6, printed p.~22]{KantorSeress2001}.
Over $\mathbf F_2$, an irreducible group of dimension at least three
generated by transvection groups is the full special linear group,
or a symplectic group, or one of the listed orthogonal and symmetric
groups in their even-dimensional natural representations.
Each non-linear alternative preserves a nondegenerate alternating
bilinear form.  This form is unique: the endomorphism ring of the
irreducible $H$-module is a finite field, whose degree over
$\mathbf F_2$ divides the rank one of a transvection, and hence is
one.  Schur's lemma then gives a one-dimensional space of invariant
forms.  Every normalizing element of $G$ preserves that form, since
$\mathbf F_2^\times=1$.  Its characteristic polynomial would be
reciprocal, contrary to the hypothesis.  Only
$H=\mathrm{SL}_N(\mathbf F_2)=\GL_N(\mathbf F_2)$ remains.
\end{proof}

\begin{theorem}[Exact groups in degrees one and two on one fixed surface]
\label{thm:product-char2-two-full-groups}
\status{Proved under stated hypotheses}
Let
\[
 E:\quad Y^2Z+XYZ+X^3+Z^3=0\quad\text{over }\mathbf F_2,
 \qquad X=E\times E,
 \qquad L=\mathcal O_E(3O)\boxtimes\mathcal O_E(3O),
\]
with the fixed Segre embedding of the two plane cubics.
For the relative Hasse--Witt local system on the ordinary smooth
degree-$d$ hypersurface-section open, both geometric and arithmetic
images are, up to a common choice of basis,
\[
 \begin{array}{c|c|c}
 d&\rank\mathcal V_d&G_{\rm geom}=G_{\rm arith}\\ \hline
 1&8&\GL_8(\mathbf F_2)\\
 2&35&\left\{
       \begin{pmatrix}A&B\\0&I_3\end{pmatrix}:
       A\in\GL_{32}(\mathbf F_2),\ B\in M_{32\times3}(\mathbf F_2)
       \right\}.
 \end{array}
\]
In particular the second image is the full pointwise stabilizer of
the three independent covectors already constructed, of order
$2^{96}|\GL_{32}(\mathbf F_2)|$.  The same statements hold after
base change to any perfect field of characteristic two.
\end{theorem}

\begin{proof}
The generic ranks and the equality between the ambient hypersurface
family and the complete linear system have already been established
in Theorem~\ref{thm:product-exact-fixed-covectors}.  We recall the
matrix model, to make the finite computations reproducible.
Write $x=X/Z$, $y=Y/Z$ and $\omega=dx/x$, normalized by
$C(\omega)=\omega$.  Put $C(f\omega)=T(f)\omega$.
The equation $y^2+xy=x^3+1$ gives
\[
 T(x^{2a})=x^a,\quad T(x^{2a+1})=0,\quad
 T(x^{2a}y)=x^{a+1},\quad
 T(x^{2a+1}y)=x^a(y+1).
\]
For $d=1$ use the ordered basis $(1,x,y)$ of
$H^0(E,\mathcal O(3O))$; for $d=2$ use
$(1,x,x^2,x^3,y,xy)$ of $H^0(E,\mathcal O(6O))$.
Use the lexicographically ordered tensor basis on $X$.
An integer $s=\sum s_j2^j$ encodes the section with coefficient
$s_j\in\mathbf F_2$ on the $j$th tensor basis vector, beginning
with index zero.  Write $W=H^0(X,L^d)$ and $b=\dim W=9d^2$.

For a universal coefficient vector $s$, let $M(s)$ be the linear
matrix whose $j$th column is the coefficient vector of
$T\otimes T(s e_j)$, with the scalar Frobenius root removed.
Thus the actual Cartier map is
$v\mapsto M(s)^{[1/2]}v^{[1/2]}$, and the dual coherent
Frobenius matrix is $M(s)^t$.
The full coherent bundle is $W^*\otimes\mathcal O(-1)$, and its
kernel under evaluation at $s$ is the relative bundle
$\mathcal E_d=\Omega^1_{\mathbf P(W)}$.
Its quotient is $H^2(X,\mathcal O_X)\otimes\mathcal O$, with
bijective Frobenius.  The identity
\[
                       M(s)s=s^{[2]}
\]
expresses this quotient.  Therefore $\det M(s)$ has exactly the
same invertibility open as the relative Frobenius.

This open is defined on the whole projective parameter space,
including singular sections.  The Lang local system extends over
all of it.  Removing the smooth-discriminant locus does not change
its image: a dense open in a normal integral scheme surjects onto
its etale fundamental group.  In particular every ordinary
$\mathbf F_2$ coefficient vector below supplies an actual
arithmetic Frobenius class for the image in question; it need not
itself define a smooth curve.  On the dual representation its
characteristic polynomial is that of the induced Cartier matrix;
on the original representation it is reciprocal.  Our irreducibility,
nonreciprocity, order and fixed-space tests are unchanged by this
duality convention.

\par\smallskip\noindent\textbf{Divisorial inertia certificates.}
The retained exact program
\path{code/product_char2_group_certificates.py}
constructs all matrices from the four displayed Cartier formulas.
Its full output, including every binary column and every finite-field
rank, is
\path{outputs/product_char2_group_certificates.json}.
The following two affine lines have ordinary leading coefficient:
\[
 \begin{array}{c|r|r|l}
 d&a&b_0&\det M(a+t b_0)\\ \hline
 1&265&303&t^2(t^2+t+1)(t^5+t^4+t^2+t+1)\\
 2&56866400210&54262944944&(t+1)^4q_7(t)q_9(t)q_{16}(t),
 \end{array}
\]
where
\[
\begin{split}
 q_7&=t^7+t^6+t^5+t^3+t^2+t+1,\\
 q_9&=t^9+t^8+t^6+t^3+1,\\
 q_{16}&=t^{16}+t^{15}+t^{13}+t^{10}+t^6+t^3+t^2+t+1.
\end{split}
\]
The listed factors are irreducible over $\mathbf F_2$.
At a root of each listed factor, both $M$ and $M^{[2]}M$
have rank $8$ for $d=1$, respectively $35$ for $d=2$.
The program verifies irreducibility by the finite-field Frobenius
criterion, multiplies the factors back to the determinant, and
performs Gaussian elimination in the respective finite fields.
The determinants have full degrees $9$ and $36$.
Lemma~\ref{lem:hasse-line-certificate} therefore gives stable rank
$b-1$ for the full coherent module at every geometric generic
Hasse divisor.  Removing its bijective rank-one quotient gives
stable relative rank $b-2$, a loss of exactly one.
Lemma~\ref{lem:char2-defect-one-inertia} proves that both geometric
images are generated by transvections.  This argument includes
all Hasse components without asserting their irreducibility or
reducedness.

\par\smallskip\noindent\textbf{Degree one.}
For the ordinary section encoded by $258$, the matrix on $W$
has characteristic polynomial
\[
 (T+1)(T^8+T^4+T^3+T+1).
\]
The degree-eight factor is irreducible and not reciprocal.
For the ordinary section encoded by $27$, the characteristic
polynomial is $(T+1)^9$.
Thus the geometric image is not trivial: if it were trivial,
the local system would be constant over the algebraic closure,
and all rational points would have the same arithmetic Frobenius
characteristic polynomial.  The arithmetic image is irreducible
because of the first section.  Apply
Lemma~\ref{lem:product-frobenius-group-test} with $N=8$.
It gives both images equal to $\GL_8(\mathbf F_2)$.

\par\smallskip\noindent\textbf{The action on the common kernel in degree two.}
The four character projections $s_\chi$ of $s$ give four fixed
vectors in $W$, with sum $s$; in $W/\langle s\rangle$ they give
three independent invariant covectors.  All are defined over
$\mathbf F_2$, by the explicit Cartier calculation above and
Theorem~\ref{thm:product-exact-fixed-covectors}.
Let $K\subset V$ be their common kernel.  Then $\dim K=32$ and
the quotient $Q=V/K$ is the constant three-dimensional module
for both geometric and arithmetic monodromy.

At the ordinary section encoded by $17707720254$, the characteristic
polynomial on $W$ is $(T+1)^4f_{32}(T)$, where
\[
\begin{split}
 f_{32}(T)={}&T^{32}+T^{28}+T^{27}+T^{26}+T^{25}+T^{24}
             +T^{23}+T^{22}+T^{20}+T^{17}\\
            &+T^{14}+T^{11}+T^8+T^6+T^5+T^3+T^2+T+1.
\end{split}
\]
The program verifies that this polynomial is irreducible and not
reciprocal.  Consequently the arithmetic action on $K$ is
irreducible.  At the second ordinary section, encoded by
$65866883812$, the polynomial on $K$ has a root at $1$.
Thus the geometric image on $K$ is nontrivial, by the same
constant-local-system argument as in degree one.

The restriction to $K$ of a transvection of $V$ is either the
identity or a transvection: every image line lies in $K$, since
the three quotient covectors are fixed.  Hence the geometric
image on $K$ is generated by transvections.  Applying
Lemma~\ref{lem:product-frobenius-group-test} with $N=32$ shows
that this image is the entire $\GL(K)$.

\par\smallskip\noindent\textbf{The translation subgroup in degree two.}
It remains to prove that every map $Q\to K$ occurs; surjectivity
onto $\GL(K)$ alone does not prove this.
Choose a vector-space splitting $V=K\oplus Q$ and let $G_0$ be
the geometric image.  Its kernel $R$ in
\[
 1\longrightarrow R\longrightarrow G_0
          \longrightarrow\GL(K)\longrightarrow1
\]
is a $\GL(K)$-submodule of $\Hom(Q,K)=K\otimes Q^*$.
The natural module is absolutely simple with endomorphism ring
$\mathbf F_2$.  Thus $R=K\otimes U$ for a subspace $U\subseteq Q^*$.
Modulo $R$, the subgroup is the graph of a cocycle in
$H^1(\GL(K),K\otimes(Q^*/U))$.
This cohomology is zero by the elementary transvection calculation
in Lemma~\ref{lem:char2-group-inputs}, valid in every dimension
at least four.  Conjugating the splitting therefore makes
$G_0=R\rtimes\GL(K)$, and its invariant-vector space projects
isomorphically onto $U^\perp\subseteq Q$.

That invariant-vector space is stable under arithmetic monodromy.
Its projection to the constant quotient $Q$ is injective, so
arithmetic monodromy acts trivially on it as well.
We now use the additional exact certificate at $s=65866883812$:
if $C=M(s)$ and the four character vectors are $s_\chi$, then
\begin{equation}\label{eq:product-no-complement-certificate}
 \dim\bigl(\im(C-I)+\langle s\rangle\bigr)
 =\dim\bigl(\im(C-I)+\langle s,s_\chi:\chi\rangle\bigr)=33.
\end{equation}
The full columns and the four vectors are retained in the output.
On the relative coherent fibre $V\otimes\bar k=\operatorname{ann}(s)$,
a vector fixed by this rational-point Frobenius lies in
$\ker(C^t-I)\cap\operatorname{ann}(s)$.
Equation~\eqref{eq:product-no-complement-certificate} forces it
to annihilate every $s_\chi$ and hence to lie in $K$.
Every geometric invariant vector is arithmetic invariant, so this
applies to it.  But $K^{\GL(K)}=0$.  Therefore $V^{G_0}=0$,
$U^\perp=0$, and $U=Q^*$.
The kernel is the whole $\Hom(Q,K)$, proving the asserted full
pointwise stabilizer.  The arithmetic image is contained in that
same stabilizer and contains the geometric image, so the two agree.

Finally the geometric covers remain connected under algebraically
closed extension.  Both displayed groups are already defined over
$\mathbf F_2$, and the three degree-two covectors are rational
over that field.  After any perfect base extension the arithmetic
image is therefore still between the same equal geometric image
and the same upper bound.  This proves the final assertion.
\end{proof}

\begin{remark}[Verification scope]\label{auto:A_monodromy:9178c750cd75}
\status{Proved under stated hypotheses}

The group conclusions use the lemmas proved here, the elementary
cocycle and socle arguments already proved, and the single opened
finite-group classification specified above.  The retained program
provides exact polynomial, rank and Frobenius certificates.  It does
not infer monodromy from a random sample.  The exploratory search
chooses certificates, after which every statement needed from each
certificate is verified exactly.  The high-dimensional characteristic
polynomials are also recomputed independently by SymPy's integer
characteristic-polynomial routine and reduced modulo two.
No assertion for degree $d\ge3$ on this surface follows just from
these two degree calculations.
\end{remark}


% /root/monodromy, 19 November 2026; independently audited by /root and /root/reduction_descent.
\subsection{Uniform structure on the fixed ordinary product surface}
\label{subsec:product-uniform-structure}
\status{Proved; parent proof, primary-source and code audits passed.
The uniform jet, incidence and union argument was also independently audited by the reduction/descent agent.}

We first record the general finite-group form of the prime-order
test used above.  We then use one further computed base degree to
obtain a statement in every degree; the latter statement does not
identify the entire group in every degree.

\begin{lemma}\label{lem:primitive-prime-transvection-test}
\status{Proved under stated hypotheses}
Let $N\ge4$ and $H\triangleleft G\subseteq\GL_N(\mathbf F_2)$.
Suppose that $H$ is nontrivial and generated by transvections.
Assume $G$ contains an element $g$ whose characteristic polynomial
is irreducible and not reciprocal, and whose order is divisible
by a prime $\ell$ with $\operatorname{ord}_{\ell}(2)=N$.
Then $H=G=\GL_N(\mathbf F_2)$.
\end{lemma}

\begin{proof}
Use the socle and multiplicity argument in
Lemma~\ref{lem:product-frobenius-group-test}.
If $H$ were reducible, $G$ would be contained in
$\GL_c(\mathbf F_2)\wr S_{N/c}$ for a proper divisor $c$ of $N$.
The order condition gives $\ell\equiv1\pmod N$, hence $\ell>N$;
and $\ell$ divides no $2^i-1$ with $i<N$.
It therefore divides neither factor of this wreath-product order.
This contradicts the order of $g$.
Thus $H$ is irreducible.  The same opened
\cite[Theorem~3.6, p.~22]{KantorSeress2001}, uniqueness of the
invariant alternating form in its non-linear alternatives, and
nonreciprocity exclude those alternatives exactly as before.
\end{proof}

\begin{corollary}\label{cor:product-degree-three-full-group}
\status{Proved under stated hypotheses}
For the fixed surface and embedding in
Theorem~\ref{thm:product-char2-two-full-groups}, the geometric and
arithmetic groups in degree three are both
$\GL_{80}(\mathbf F_2)$, over every perfect extension of
$\mathbf F_2$.
\end{corollary}

\begin{proof}
Here $W$ has dimension $81$.  The one-factor ordered basis is
\[
                 (1,x,x^2,x^3,x^4,y,xy,x^2y,x^3y).
\]
Use the same tensor and binary coefficient conventions as before.
The ordinary section encoded by $1284714998244677440413117$
has characteristic polynomial $(T+1)f_{80}(T)$, where the binary
coefficient encoding of $f_{80}$ is
\[
                         1592084315789834604401113.
\]
The retained program verifies that it is irreducible and not
reciprocal.  Its root order is divisible by
$\ell=4278255361$, and
$\operatorname{ord}_{\ell}(2)=80$.
The independent verifier checks primality by trial division
through $\lfloor\sqrt\ell\rfloor$, checks the order by the
exponents $80,40,16$, and checks the root-order assertion by
modular polynomial exponentiation.  The ordinary section
$810021528320842360421666$ has a different relative characteristic
polynomial, so the geometric image is nontrivial.

For the line with
\[
 a=998130104707720808648124,
 \qquad b_0=2032072755533295013746953,
\]
the determinant is monic of degree $81$.  Its six irreducible
factors have degrees $5,10,10,11,11,34$, and at a root of every
factor both $M$ and $M^{[2]}M$ have rank $80$.
Thus Lemmas~\ref{lem:hasse-line-certificate} and
\ref{lem:char2-defect-one-inertia} give transvection generation.
Apply Lemma~\ref{lem:primitive-prime-transvection-test}.

The complete matrices, factors, ranks and section coefficients
are in \path{outputs/product_degree3_certificates.json};
\path{code/verify_product_degree3.py} reconstructs them, independently
computes both $81\times81$ characteristic polynomials over the
integers before reduction modulo two, and repeats all the line
rank calculations.  Its exact output is
\path{outputs/verify_product_degree3.json}.
\end{proof}

\begin{lemma}[The required jet separation]
\label{lem:cubic-power-jets}
\status{Proved under stated hypotheses}
If $L$ is very ample on a smooth variety, $L^3$ separates second
jets at one point and first jets at two distinct points.
On a smooth projective surface its complete linear system has an irreducible
discriminant whose generic member has exactly one ordinary node;
the nonreduced locus has codimension at least two.
\end{lemma}

\begin{proof}
Ratios of sections of $L$ give local coordinates, so products of
at most two such ratios span the local ring modulo the cube of
the maximal ideal.  Multiplication by a section nonzero at the
point identifies this with second-jet generation by $L^3$.
For distinct points $P,Q$, take a section $m$ of $L$ vanishing
at $Q$ and nonzero at $P$.  Sections $m^2t$ with $t\in H^0(L)$
vanish to order at least two at $Q$ and give arbitrary first
jets at $P$.  Interchanging the points proves the assertion.

The incidence of a section with a specified singular point is
a projective bundle over the surface, since value and first
derivative impose three independent linear conditions.  It is
irreducible of dimension one less than the parameter space.
Second-jet separation makes an ordinary quadratic node a nonempty
open condition in this incidence, including in characteristic two.
Two specified distinct singular points impose six independent
conditions; varying the two points adds four dimensions.
Thus the locus with two distinct singular points has codimension
at least two in the parameter space.  The generic singular member
has exactly one ordinary node, and the discriminant is an
irreducible divisor.  A nonreduced divisor is singular along a
positive-dimensional locus, so belongs to the codimension-two
locus just described.
\end{proof}

\begin{theorem}[Transvection generation in every degree]
\label{thm:product-all-degree-transvections}
\status{Proved under stated hypotheses}
For the fixed $X=E\times E$ and $L$ of
Theorem~\ref{thm:product-char2-two-full-groups}, every geometric
monodromy group $G_{{\rm geom},d}$, $d\ge1$, is generated by
transvections.  Equivalently for the proof, every geometric
irreducible component of the full degree-$d$ Hasse divisor has
generic stable rank exactly $9d^2-2$ on the relative coherent
bundle of rank $9d^2-1$.
\end{theorem}

\begin{proof}
The divisor statement has been checked in degrees one, two and
three by the exact line certificates above.  Fix $d\ge4$, put
$e=d-3$, and work over the algebraic closure.
Choose a general smooth ordinary curve $D\in|L^e|$.
Such curves exist for every $e\ge1$ by
Theorem~\ref{thm:product-exact-fixed-covectors}.
Write $P_3=\mathbf P H^0(X,L^3)$ and let $C$ vary in this space.

We describe the bad loci for a general pencil in $P_3$.
The base Hasse divisor has generic stable rank loss one, and its
exceptional stable-rank locus has codimension at least two in
$P_3$.  By Lemma~\ref{lem:cubic-power-jets}, nonreduced curves and
curves worse than the generic one-node member form loci of
codimension at least two in the discriminant analysis relevant
here.  In particular every divisorial Hasse component has a
reduced generic member with at most one node.

For fixed smooth $D$, the incidence
\[
 \{(P,C): P\in D,\ s_C(P)=0,\ ds_C|_{T_PD}=0\}
\]
is a projective bundle over $D$ and is irreducible.  Its image
$\Delta_D\subset P_3$ is an irreducible divisor.  At its generic
point, $C$ is smooth and $C\cap D$ has one point of length two
and is otherwise reduced.  Indeed, contact of order at least
three imposes three conditions at one point of $D$, giving
codimension at least two after that point varies.  Two distinct
tangencies impose four independent conditions and the two points
vary in dimension two, again giving codimension at least two.
A singularity of $C$ at the tangency point adds a further linear
condition.  A singularity elsewhere, together with the tangency,
imposes five independent conditions at two points, which vary in
dimension three.  All these independence claims follow from
the two jet-separation assertions of
Lemma~\ref{lem:cubic-power-jets}.

For a general ordinary $D$, the generic $C$ on $\Delta_D$ is
ordinary.  Here is a dominance check to justify this assertion,
rather than a transversality assumption on Hasse strata.
Consider the joint incidence $(P,C,D)$ with both curves smooth
at $P$ and with equal tangent lines there.  Value and first-jet
generation for $L^3$ and $L^e$ identify it, locally over $P\in X$,
with the irreducible rank-one locus of two nonzero proportional
cotangent vectors, times affine kernel spaces.  Hence the joint
incidence is irreducible.  It dominates each of the two complete
linear systems: for a smooth curve in either system choose a
point on it and prescribe a proportional nonzero first jet for
the other section.  The inverse images of the two ordinary opens
are therefore nonempty dense opens of this same irreducible
incidence, so their intersection is dense.  Passing to a general
$D$ gives the asserted ordinary generic member of $\Delta_D$.
Consequently $\Delta_D$ is not a component of the base Hasse
divisor.  Their intersections and their exceptional loci have
codimension at least two in $P_3$.

The locus of $C$ containing $D$ is also of codimension at least
two.  For $e\le3$ it is the linear subspace
$s_D H^0(L^{3-e})$, whose codimension is respectively $45,72,80$;
for $e>3$ it is empty.  We may therefore choose a projective
pencil $\Lambda\subset P_3$ avoiding all the codimension-two
loci above.  Its point at infinity is chosen ordinary and
transverse to $D$.

For every member of this pencil, put $Y=C\cup D$.
There are no common components.  The exact sequence
\[
 0\longrightarrow\mathcal O_Y
 \longrightarrow\mathcal O_C\oplus\mathcal O_D
 \longrightarrow\mathcal O_{C\cap D}\longrightarrow0
\]
and exactness of stable parts give
\begin{equation}\label{eq:product-union-stable-rank}
 f(Y)=f(C)+f(D)+\#(C\cap D)(\bar k)-1,
\end{equation}
where $f$ denotes the stable dimension of $H^1(\mathcal O)$.
To see the finite-scheme term explicitly, Frobenius kills the
nilradical of $H^0(\mathcal O_{C\cap D})$ and is bijective on
its reduced geometric quotient.  Also $H^0(\mathcal O_C)=
H^0(\mathcal O_D)=\bar k$.  The coherent analogue uses the
scheme length instead of the number of points:
\[
 p_a(Y)=p_a(C)+p_a(D)+\operatorname{length}(C\cap D)-1.
\]
At a Hasse point of $C$, the intersection is reduced and $C$
loses exactly one stable dimension.  At a point of $\Delta_D$,
both components are ordinary and the number of intersection
points is one less than their scheme length.  These are distinct
sets of pencil parameters.  At every other parameter the union
is ordinary.  Thus every nonordinary union on the pencil has
stable rank loss exactly one.

Multiplication by the fixed defining section $s_D$ embeds this
pencil as a genuine projective line in $\mathbf P H^0(X,L^d)$.
Its point at infinity is ordinary, so every geometric component
of the full Hasse divisor meets the line at a finite parameter.
If such a component had generic stable rank loss at least two,
the closed stable-rank locus defined by the appropriate minors
of an iterated Frobenius matrix would contain the component and
all of those intersection points.  This contradicts
\eqref{eq:product-union-stable-rank}.  Hence every full-system
Hasse component has generic loss one.  The relative and absolute
losses agree, since $H^1(X,\mathcal O_X)$ has bijective Frobenius.
Lemma~\ref{lem:char2-defect-one-inertia} now proves transvection
generation in every degree.
\end{proof}

\begin{theorem}[All characteristic-two invariant covectors on products]
\label{thm:product-char2-all-parities-fixed-quotient}
\status{Proved under stated hypotheses}
Let $X=\prod_{i=1}^n E_i$ be a product of ordinary elliptic
curves in characteristic two, $n\ge2$, and
$L=\boxtimes_i\mathcal O_{E_i}(a_iO_i)$ with $a_i\ge3$.
For any $d\ge1$, let
$e=\#\{i: a_i d\text{ is even}\}$.
Then the entire geometric invariant covector space has dimension
\[
             \dim(V_d^\vee)^{G_{{\rm geom},d}}=2^e-1.
\]
After a finite extension splitting the ordinary Frobenius kernels
and normalizing the invariant top Cartier differential, the
arithmetic invariant space is the same space.
In particular, on the fixed $E\times E$ with $a_1=a_2=3$,
this dimension is zero for odd $d$ and three for even $d$.
\end{theorem}

\begin{proof}
Work over an algebraic closure and normalize the invariant top
form by Cartier.  The proof of
Theorem~\ref{thm:product-exact-fixed-covectors}, through
\eqref{eq:product-exact-Cartier-operator}, used no divisibility
assumption.  It identifies invariant covectors with exact lifts
$D\in\End(W)$, modulo $\mathbf F_2\operatorname{id}$, satisfying
\begin{equation}\label{eq:char2-projector-exact}
        \beta(s,D(s))=D(s)^{[2]}\quad\text{for every }s\in W,
\end{equation}
where $\beta$ is the symmetric relative Cartier pairing on
global sections and $\beta(s,s)=s^{[2]}$.

Every such $D$ is an idempotent.  An eigenvector with eigenvalue
$a$ in \eqref{eq:char2-projector-exact} gives $a=a^2$.
If a nontrivial Jordan block remained, subtract its eigenvalue
times the identity and choose $v,w\ne0$ with $D(v)=w$ and
$D(w)=0$.  Substituting $v+tw$ in the same equation makes the
coefficient of $t$ on the left equal to $\beta(w,w)=w^{[2]}$,
whereas the right side is independent of $t$, a contradiction.
Thus $W=A\oplus B$ is its $1$- and $0$-eigenspace decomposition.
Substituting $a+tb$ shows that
\begin{equation}\label{eq:char2-projector-orthogonality}
                    \beta(A,B)=0.
\end{equation}

Let $F:X\to X^{(2)}$ be relative Frobenius and
$\mathcal F=F_*L^d$.  Finite duality gives a perfect symmetric
pairing
\[
 \mathcal F\otimes\mathcal F\longrightarrow L^{d,(2)},
\]
obtained by multiplying sections and applying the normalized
Cartier trace.  Its pairing on global sections is $\beta$.
The bundle $\mathcal F$ is globally generated.  On a factor
$E_i$, evaluation on a Frobenius fibre is evaluation on a
length-two subscheme; $H^1(\mathcal O_{E_i}(a_idO_i)(-Z))=0$
because its degree is $a_i d-2>0$.  Tensoring these surjections
proves global generation on the product.

The projector $D$ preserves every evaluation kernel.
Indeed, write a global section whose value in a fibre is zero
as $a+b$ with $a\in A,b\in B$.  Its pairing with every section
is zero in that fibre.  By
\eqref{eq:char2-projector-orthogonality}, the value of $a$ pairs
to zero with $A$ and with $B$, hence with the whole fibre by
global generation.  Perfectness makes that value zero, and the
same holds for $b$.  The evaluation kernel is a vector bundle;
vanishing on all closed fibres therefore gives a descended
bundle endomorphism $\mathcal D\in\End(\mathcal F)$.
Its action on global sections is the original $D$.

We compute this endomorphism algebra directly.  For an ordinary
elliptic curve and a line bundle $N$ of odd degree $m$, the
rank-two bundle $F_*N$ has degree $m$ and is stable.
If a line subbundle $A\subset F_*N$ existed with degree at least
$m/2$, adjunction would give a nonzero map $F^*A\to N$ and
$2\deg A\le m$, impossible for the proposed integral degree.
Thus every line subbundle has degree strictly below the slope.
A noninvertible nonzero endomorphism would have a rank-one
kernel and image, both of degree below the slope (saturating the
image only raises its degree), contradicting additivity of
degree.  Over the algebraically closed field it follows that
\[
                     \End(F_*N)=\bar k.
\]

For even degree, $N=F^*N'$ and the ordinary splitting gives
$F_*N=N'\oplus(N'\otimes T)$, where $T$ is the nontrivial
two-torsion character line bundle.  The two summands are
nonisomorphic and have no homomorphisms between them.
This is the dimension-one case of the opened ordinary-Frobenius
splitting in \cite[Lemma~4.5, p.~11]{SannaiTanaka2016}.
Their cross Cartier pairing is zero, since the trace kills the
nontrivial character.

Relative Frobenius on the product is the product of the relative
Frobenii.  The bundle $\mathcal F$ is therefore a direct sum of
$2^e$ external products, one for each choice of a character on
the even-degree factors.  Each external product has endomorphism
algebra $\bar k$, by the degree-zero Kunneth formula.  Distinct
summands have no homomorphisms, because on one even-degree factor
the corresponding degree-zero nontrivial character line has no
sections.  Thus
\[
                      \End(\mathcal F)=\bar k^{\,2^e}.
\]
Every exact $D$ is consequently a sum of some of these $2^e$
projectors.  Conversely the summands are orthogonal for the
Cartier pairing; every such projector satisfies
\eqref{eq:char2-projector-exact}.  The exact-lift space therefore
has dimension $2^e$ over $\mathbf F_2$.
Quotienting by the identity gives the asserted $2^e-1$.
All splittings and the normalized top form descend after a
finite extension, which also proves the arithmetic statement.
\end{proof}

\begin{remark}[What the uniform results leave open]\label{auto:A_monodromy:1386149fe4fc}
\status{Proved under stated hypotheses}

For the fixed surface, we now know the exact generic rank, the
entire invariant quotient in every degree, and that geometric
monodromy is generated by transvections in every degree.
The complete images are identified above in degrees one, two
and three.  A transvection-generated group with the stated
invariant quotient can still have further invariant subspaces
or a proper translation subgroup.  The two uniform theorems do
not remove those possibilities for all $d$, and no all-degree
full-stabilizer assertion is made here.
\end{remark}


% /root/monodromy, 19 November 2026; complete conditional reduction independently audited by /root.
\subsection{A single geometric gate for the full product groups}
\label{subsec:product-irreducible-hasse-reduction}
\status{Reducedness and the complete reduction proved and independently audited by the parent. The geometric hypothesis is discharged for every degree at least four in Theorem~\ref{thm:product-Hasse-irred}; the resulting unconditional all-degree images are Corollary~\ref{thm:product-char2-all-full-groups}.}

\begin{proposition}\label{prop:product-hasse-reduced}
\status{Proved under stated hypotheses}
For the fixed surface $E\times E$ and embedding considered above,
the full Hasse polynomial $h_d(s)=\det M_d(s)$ is geometrically
reduced for every $d\ge1$.
\end{proposition}

\begin{proof}
At the geometric generic point of any irreducible Hasse component,
Theorem~\ref{thm:product-all-degree-transvections} gives stable rank
$b-1$ for the full coherent matrix, where $b=9d^2$.
Thus $M$ and $M^{[2]}M$ both have rank $b-1$.
Let $u$ span $\ker M$ and let $\ell$ span the left kernel.
Since $\ker M^{[2]}$ is spanned by $u^{[2]}$, equality of the two
ranks says that $u^{[2]}$ is not in $\im M$, or equivalently
$\ell(u^{[2]})\ne0$.

The symmetric Cartier multiplication pairing gives
$M(s)w=M(w)s$ and $M(w)w=w^{[2]}$.
The adjugate at this rank-$b-1$ point is a nonzero scalar multiple
of $u\ell$.  Differentiating the determinant in the coefficient
direction $u$ therefore gives a nonzero scalar multiple of
\[
                     \ell M(u)u=\ell(u^{[2]})\ne0.
\]
Each component is consequently generically smooth and occurs
with multiplicity one.  The repeated factors in some retained
pencil determinants reflect tangency of those pencils, not a
nonreduced global Hasse divisor.
\end{proof}

\begin{lemma}[Frobenius-periodic reflexive sheaves]
\label{lem:periodic-reflexive-projective-trivial}
\status{Proved under stated hypotheses}
Let $P$ be projective space over an algebraically closed field of
characteristic two.  If a coherent reflexive sheaf $\mathcal S$
has an isomorphism $F^*\mathcal S\simeq\mathcal S$, then it is a
trivial vector bundle with its constant unit-Frobenius structure.
\end{lemma}

\begin{proof}
Fitting ideals commute with base change.  In a local ring of $P$,
each Fitting ideal $I$ therefore satisfies $I=I^{[2]}$.
If $I$ is proper, this gives $I\subseteq\mathfrak m I$, so
Nakayama makes $I=0$.  The usual Fitting-ideal criterion, or a
minimal presentation, now shows that the sheaf is locally free
of its generic rank.  Its Frobenius isomorphism gives a finite
etale Lang frame cover.  Projective space has trivial geometric
fundamental group, as proved in
Lemma~\ref{lem:char2-defect-one-inertia}, so this frame cover has
a section.  The resulting frame is Frobenius-fixed.
\end{proof}

\begin{lemma}[Determinant allocation]
\label{lem:irreducible-hasse-quotients-trivial}
\status{Proved under stated hypotheses}
Let $P=\mathbf P^m$ over an algebraically closed field of
characteristic two, and let $\mathcal E=\Omega_P^1$ carry a
generically invertible Frobenius-linear map.  Suppose its Hasse
divisor is reduced and irreducible.
Every nonzero proper invariant subspace of its Lang representation
has trivial quotient representation.
\end{lemma}

\begin{proof}
Let $U$ be the invertibility open and extend the unit-Frobenius
subbundle of an invariant subspace by saturation in $\mathcal E$.
This gives a coherent exact sequence
\[
 0\longrightarrow\mathcal S\longrightarrow\mathcal E
       \longrightarrow\mathcal Q\longrightarrow0,
\]
where $\mathcal Q$ is torsion-free and $\mathcal S$ is reflexive.
The latter assertion follows locally from the depth lemma:
the kernel of a vector bundle mapping to a torsion-free sheaf
on a regular scheme satisfies the depth-two condition.
Frobenius preserves the saturated subsheaf, because it does so
at the generic point.  Flatness of Frobenius on the regular
scheme preserves reflexivity (regular sequences may equivalently
be replaced by their squares).  It therefore induces regular
maps on $\mathcal S$ and, after reflexive extension, on
$\mathcal Q^{**}$.

In codimension one all terms are vector bundles, and determinants
of these maps give effective divisors satisfying
\[
                     D_{\mathcal E}=D_{\mathcal S}+D_{\mathcal Q}.
\]
If $D_{\mathcal S}=0$, the map
$F^*\mathcal S\to\mathcal S$ is an isomorphism in codimension
one.  Both sheaves are reflexive, so it is an isomorphism
everywhere.  Lemma~\ref{lem:periodic-reflexive-projective-trivial}
makes $\mathcal S$ a nonzero trivial bundle.  Its inclusion in
$\Omega_P^1$ contradicts $H^0(P,\Omega_P^1)=0$, which follows
at once from the Euler sequence.

The reduced irreducible divisor must consequently be allocated
entirely to $\mathcal S$, and $D_{\mathcal Q}=0$.
The same argument makes $\mathcal Q^{**}$ a trivial unit-Frobenius
bundle.  On $U$ it is the original quotient bundle, because
that quotient is already locally free there.  The composite
$\mathcal E\to\mathcal Q\to\mathcal Q^{**}$ is Frobenius-equivariant
and exhibits the quotient local system as constant.
\end{proof}

\begin{lemma}[A fixed natural block in every higher degree]
\label{lem:product-degree-three-block}
\status{Proved under stated hypotheses}
For every $d\ge3$, the geometric group on the fixed surface
contains a subgroup acting as the natural
$\GL_{80}(\mathbf F_2)$ on an $80$-dimensional direct summand
and trivially on a complementary summand.
\end{lemma}

\begin{proof}
The assertion for $d=3$ has already been proved.  For $d>3$ fix
$d-3$ ordinary smooth members $D_i\in|L|$ in general transverse
position, with no triple points, and put $D=\bigcup D_i$.
This fixed connected nodal curve is ordinary.
Vary an ordinary degree-three curve $C$, transverse to every
$D_i$ and avoiding all nodes of $D$.
The resulting curves $C\cup D$ are ordinary by the union sequence,
so their Lang sheaf is a restriction of the full ordinary
coherent family in degree $d$.  Its image is a subgroup of the
degree-$d$ geometric group.  The allowed $C$ form a dense open
of the degree-three ordinary parameter space; hence their
relative degree-three image remains $\GL_{80}(\mathbf F_2)$.

The union sequence and the constant ambient
$H^1(X,\mathbf F_2)$ give a filtration on this boundary
representation $V$,
\[
 0\subset T\subset T_0\subset V,
 \qquad T_0/T\text{ constant},\qquad V/T_0=Q,
 \quad\dim Q=80.
\]
Here $T$ is the graph term from $C\cap D$ and $Q$ is the relative
degree-three local system.  More explicitly, the quotient by
$T$ is
$(H^1(C,\mathbf F_2)\oplus H^1(D,\mathbf F_2))/H^1(X,\mathbf F_2)$;
the inverse image of the constant $H^1(D,\mathbf F_2)$ is $T_0$.
The ambient maps are injective, so these identifications are exact.

Each $C\cap D_i$ consists of $3L^2=54$ points.
Thus the image on $T$ is a subgroup of a product of $S_{54}$'s.
The image on $T_0$ is an extension of such a subgroup by a
subgroup of the additive two-group
$\Hom(T_0/T,T)$.  Every prime in its order is at most $53$.
In contrast $\GL_{80}(\mathbf F_2)$ is simple and its order is
divisible by $4278255361$.

We recall an elementary simplicity proof to specify this input.
The group $\GL_N(\mathbf F_2)$ for $N\ge3$ is perfect, since
$1+E_{ij}=[1+E_{ik},1+E_{kj}]$ for distinct indices, and the
elementary matrices generate it.  It acts faithfully and doubly
transitively on nonzero vectors.  A nontrivial normal subgroup
is therefore transitive.  The stabilizer of a vector $v$ has
the abelian normal subgroup
$U_v=\{1+v\otimes\phi:\phi(v)=0\}$, and the conjugates of $U_v$
generate the whole group.  Modulo a transitive normal subgroup,
every conjugate of $U_v$ has the same image: move $v$ by an
element of the normal subgroup and then use its stabilizer.
The quotient is thus abelian, and perfectness makes it trivial.

Goursat's elementary common-quotient description now says that
the boundary image on $Q\oplus T_0$ contains
$\GL(Q)\times1$: a nontrivial common quotient would be the
simple group $\GL(Q)$, whose large prime cannot occur on $T_0$.
Let $B$ be the inverse image of $\GL(Q)\times1$.
It fixes $T_0$ pointwise and maps onto $\GL(Q)$.
Its kernel $R$ is a submodule of
$\Hom(Q,T_0)=Q^*\otimes T_0$, hence
$R=Q^*\otimes T'$ for a subspace $T'\subseteq T_0$.
Modulo $R$, the subgroup is a graph of a cocycle with values in
$Q^*\otimes(T_0/T')$.  The elementary vanishing
$H^1(\GL_{80}(\mathbf F_2),Q^*)=0$ from
Lemma~\ref{lem:char2-group-inputs} makes it conjugate to the
zero cocycle.  Its inverse image is then
$R\rtimes\GL(Q)$, and in particular contains the required
block-diagonal $\GL(Q)$.  A quotient surjection alone would
not have supplied this last subgroup; the cocycle step does.
\end{proof}

\begin{theorem}[Conditional exact uniform product groups]
\label{thm:product-full-group-if-hasse-irreducible}
\status{Proved under stated hypotheses}
Fix $d\ge3$ on the surface $E\times E$ above.
If the degree-$d$ Hasse divisor is geometrically irreducible,
then both geometric and arithmetic groups are
\[
 \begin{cases}
 \GL_{9d^2-1}(\mathbf F_2),&d\text{ odd},\\
 \mathbf F_2^{\,3(9d^2-4)}\rtimes\GL_{9d^2-4}(\mathbf F_2),
       &d\text{ even},
 \end{cases}
\]
where in the second line the group is the full pointwise
stabilizer of the three known independent covectors.
Theorem~\ref{thm:product-Hasse-irred} below proves the required
geometric irreducibility for every $d\ge4$, thereby completing
the entire fixed-surface full-image problem together with the
three certified initial degrees.
\end{theorem}

\begin{proof}
Put $q=0$ for odd $d$ and $q=3$ for even $d$; these are the
entire invariant-covector dimensions by
Theorem~\ref{thm:product-char2-all-parities-fixed-quotient}.
Let $K$ be the common kernel and $N=\dim K=9d^2-1-q$.
Proposition~\ref{prop:product-hasse-reduced} and the hypothesis
permit Lemma~\ref{lem:irreducible-hasse-quotients-trivial}.
Every nonzero proper invariant subspace therefore contains $K$.
It follows that $K$ is irreducible (or that $V=K$ is irreducible
when $q=0$), and that $V$ has no invariant vector: an invariant
line could not contain $K$, whose dimension is at least $80$.

The image on $K$ is irreducible and generated by transvections,
by Theorem~\ref{thm:product-all-degree-transvections}.
The natural $80$-dimensional block of
Lemma~\ref{lem:product-degree-three-block} lies in $K$, because
a nontrivial natural module has no equivariant map to a trivial
module.  On $K$ that subgroup acts naturally on this block and
trivially on its complement.
No nondegenerate alternating form on $K$ can be invariant under
it.  The natural module of $\GL_{80}(\mathbf F_2)$ has no
invariant bilinear form: elementary transvections first kill
every off-diagonal matrix entry, using a third index, and then
every diagonal entry.  Its pairing with a trivial module is
also zero.  Hence the entire natural block would lie in the
radical of any such form.

The alternatives in the opened
\cite[Theorem~3.6, p.~22]{KantorSeress2001} are therefore excluded,
and the image on $K$ is $\GL(K)$.
If $q=0$ this is the desired full group.
If $q=3$, the translation-kernel and $H^1$ argument in
Theorem~\ref{thm:product-char2-two-full-groups} applies verbatim.
The kernel is $K\otimes U$ for some $U\subseteq Q^*$, where
$Q=V/K$ is constant; vanishing of the natural-module first
cohomology conjugates the group to
$(K\otimes U)\rtimes\GL(K)$.
Its invariant-vector space is $U^\perp$.
The already proved vanishing $V^G=0$ forces $U=Q^*$, so all
translations occur.  The arithmetic image is between this
geometric group and the same pointwise stabilizer, because the
three covectors are defined over $\mathbf F_2$.
It is consequently the same group.
\end{proof}

\begin{remark}[Why reducedness alone was insufficient]\label{auto:A_monodromy:24716fc9b916}
\status{Proved under stated hypotheses}

Generic stable loss one proves reducedness, but does not by
itself exclude two reduced components meeting in codimension
two. The subsequent proof of
Theorem~\ref{thm:product-Hasse-irred} supplies the missing
irreducibility by mixed tacnodal unions and smooth points
joining all restriction factors. It does not assume a generic
symmetric-matrix rank bound for the special Cartier matrices.
\end{remark}


% /root/monodromy, 2026-11-19 08:15 UTC. Complete proof independently audited by root and RD.
\subsection{Irreducible product Hasse divisors and all full images}
\label{subsec:product-hasse-tacnode}
\status{Proved. The complete Cartier, tacnode, small-degree jet and component-merging arguments were independently audited by the parent and the reduction-descent agent. This completes every degree for the fixed surface; it is not a classification for arbitrary ambient varieties.}

\begin{lemma}[A pointed Cartier calculation]
\label{lem:pointed-Cartier}
\status{Proved under stated hypotheses}
Let $C$ be a smooth projective curve in characteristic two with
one-dimensional Cartier kernel, generated by $\omega$.
If $P\in C$ satisfies $\omega(P)\ne0$, then Cartier induces a
surjection
\[
 H^0(C,\omega_C(2P))\longrightarrow H^0(C,\omega_C),
\]
and its kernel is $k\omega$. If $D$ is ordinary, the corresponding
map for $D$ is surjective with a one-dimensional kernel, whose
nonzero members have a nonzero double-pole coefficient at $P$.
\end{lemma}

\begin{proof}
There is no residue for a meromorphic differential with its only
possible pole at $P$, by the residue theorem. Cartier kills its
double-pole term, so the displayed target is correct.
The source has dimension $g(C)+1$ by Riemann--Roch.
For an ordinary curve, Cartier is already bijective on its
$g(C)$-dimensional holomorphic subspace, which proves the second
assertion.

For the first assertion, Serre duality and finite Frobenius
duality identify the displayed Cartier map with Frobenius
\[
 H^1(C,\mathcal O_C(-P))\longrightarrow
 H^1(C,\mathcal O_C(-2P)).
\]
The duality is the curve case of the local Cartier trace and
finite-duality formula in
\cite[Lemma~1.3.6 and Proposition~1.3.7, pp.~26--27]{BrionKumar2005};
its twists are exactly $F^*\mathcal O_C(-P)=\mathcal O_C(-2P)$.
The first space maps isomorphically to $H^1(C,\mathcal O_C)$,
while the kernel of the second space mapping to
$H^1(C,\mathcal O_C)$ is
$\mathfrak m_P/\mathfrak m_P^2$, the first-jet line.

Here is the boundary coefficient explicitly. Use an affine cover
with only one member $U_0$ containing $P$. A Cech representative
$c_{ij}$ of a nonzero Frobenius-kernel class satisfies
$c_{ij}^2=b_i-b_j$. The differentials $db_i$ glue, because
$d(c_{ij}^2)=0$, and give a nonzero global Cartier-killed
form $\omega$, up to scale. This also follows directly from
the local exact sequence
$0\to\mathcal O_C\xrightarrow{F}F_*\mathcal O_C
\xrightarrow{d} B^1_C\to0$: in a uniformizer, the exact
differentials are exactly those killed by Cartier.
The map on global constants is bijective, so the connecting
map identifies $H^0(B^1_C)$ with the Frobenius kernel on
$H^1(\mathcal O_C)$. In particular the glued differential is
nonzero.

Subtract the same constant from all the $b_i$ so that $b_0(P)=0$.
The Cech intersections avoid $P$, so the original cocycle also
represents its unique lift to $H^1(\mathcal O_C(-P))$.
Its Frobenius image is the boundary of
$b_0\bmod\mathfrak m_P^2\in\mathfrak m_P/\mathfrak m_P^2$.
The coefficient of a local parameter in this class is
$db_0(P)=\omega(P)$, up to the harmless sign of the connecting
homomorphism. Hence the lift of the one-dimensional Frobenius
kernel has nonzero image in this jet line. Any kernel of the
displayed twisted Frobenius would already lie in that original
kernel; it is therefore zero. Duality proves the surjectivity,
and the dimension calculation gives its Cartier kernel.
\end{proof}

\begin{lemma}[Smoothness at a mixed tacnodal union]
\label{lem:mixed-tacnode-Hasse}
\status{Proved under stated hypotheses}
Let $X=E\times E$ be the fixed ordinary surface above.
Let $C\in|L^a|$ and $D\in|L^e|$ be smooth curves without common
components, ordinary except that $C$ has a one-dimensional
Cartier kernel $k\omega$. Suppose they have a single simple
tangency at $P$, are otherwise transverse, and $\omega(P)\ne0$.
Then the full degree-$(a+e)$ Hasse divisor is smooth at $C\cup D$.
\end{lemma}

\begin{proof}
The normalization condition can be read directly from the
hypersurface dualizing generator, so no assertion about a smooth
normalization being a smooth family is required. In completed
coordinates the two branches have equations $y=0$ and
$y=\phi(x)$, with $\phi(x)=c x^2+O(x^3)$ and $c\ne0$.
The normalization ring is the pair of power-series rings, and
its subring for the union consists of pairs congruent modulo
$\phi(x)$, equivalently modulo $x^2$.
The adjunction generator $dx/(\partial(y(y-\phi))/\partial y)$
restricts on each branch to $dx/\phi(x)$, up to a sign, which
is immaterial in characteristic two. Multiplying by the
normalization subring shows that dualizing forms have at most
double poles, with matching double-pole coefficients and
matching residues. Conversely these two equalities are precisely
the congruence modulo $x^2$. At an ordinary node the same
calculation gives simple poles and opposite residues.

Cartier on dualizing forms acts componentwise on these rational
differentials. This follows either from the same local trace
formula or its compatibility with residues: the residue of a
Cartier image is the square root of the original residue.
The local matching conditions are preserved; the double-pole
term is killed and the simple-pole term becomes its square root.
A Cartier-killed dualizing form therefore has zero residues,
so it consists of Cartier-killed elements of
$H^0(\omega_C(2P))$ and $H^0(\omega_D(2P))$ with matching
double-pole coefficients. By
Lemma~\ref{lem:pointed-Cartier}, the first is a multiple
of the holomorphic $\omega$, while a nonzero element of the
second has nonzero polar coefficient. The matching condition
forces the second to be zero. Thus the Cartier kernel on
$C\cup D$ is exactly $k(\omega,0)$.

The exact adjunction sequence
$0\to\mathcal O_X\xrightarrow{s}L^{a+e}\to\omega_{C\cup D}\to0$
has an ordinary two-dimensional ambient Cartier cokernel on
global sections, because $H^1(X,\mathcal O_X)$ is ordinary.
Its kernel is the Cartier-unit line $ks$. Consequently its
Cartier kernels identify canonically: the one-dimensional
kernel just computed lifts uniquely, after adjusting by $ks$,
to
$\ker M(s)$, where $s=f_Cg_D$. Let $u$ be its generator.
Since its differential vanishes on $D$, the section is divisible
by $g_D$: $u=g_Da_0$ with $a_0\in H^0(L^a)$ and $a_0(P)\ne0$.

Choose local coordinates so $C$ is $y=0$ and $D$ has equation
$y+x^2+$ terms of higher order, allowing a nonzero coefficient
of $x^2$ and multiplying by a unit. The section $s$ has quadratic
term proportional to $y^2$, with no $xy$ term. The local Cartier
trace of a top differential at $P$ extracts its $xy$ coefficient.
Consequently the Frobenius-twisted evaluation at $P$ annihilates
$M(s)$ and spans its left kernel. Explicitly, for the matrix
convention with the output scalar square root removed, this is
the functional with coefficient vector $(e_i(P)^2)_i$ in a
local frame: $\mathcal C(sv)(P)^2=\operatorname{ev}^{[2]}_P M(s)v$.
This twist does not affect vanishing or the following nonzero
local coefficient. The section $u$ has linear term a nonzero multiple
of $y$. Choose a global deformation $\dot s$ with zero value at
$P$ and first jet $x$; very ampleness supplies it.
Then
\[
       \operatorname{ev}^{[2]}_P\bigl(M(\dot s)u\bigr)\ne0.
\]
The rank-one adjugate formula makes this a nonzero scalar
multiple of $dh_{a+e}(\dot s)$. Hence the Hasse divisor is
smooth at this union.
\end{proof}

\begin{lemma}[The degree-one jet exception is small]
\label{lem:product-small-degree-jets}
\status{Proved under stated hypotheses}
On the fixed product, $L^e$ for $e\ge2$ separates second jets at
one point and first jets at any two distinct points.
For $L$ itself, second jets are separated at a point
$P=(P_1,P_2)$ when neither $P_i$ is a flex in its plane cubic.
For every $Q\ne P$, the joint first-jet evaluation has rank at
least five; for fixed $P$, rank six fails only on the two fibres
through $P$ and a finite set of other points.
\end{lemma}

\begin{proof}
An elliptic line bundle of degree at least six interpolates
any subscheme of length at most four, since the residual line
bundle has positive degree and hence vanishing first cohomology.
Tensoring the interpolation maps on the two factors gives the
first assertion. For degree three, the map to length-three
jets at $P_i$ is an isomorphism unless
$\mathcal O_E(3P_i)\simeq\mathcal O_E(3O)$, precisely the flex
condition. Tensoring again gives second jets on the surface.

Write $V_i=H^0(E_i,\mathcal O(3O))^*$, let $e_P$ be the
one-dimensional evaluation direction in its projective plane,
and let $S_P\subset V_i$ be the two-dimensional vector tangent
plane to the embedded cubic. The dual of the surface first-jet
map at $P$ has image
\[
 A_P=e_{P_1}\otimes S_{P_2}+S_{P_1}\otimes e_{P_2},
 \qquad\dim A_P=3.
\]
Distinct points of a smooth plane cubic have distinct tangent
lines: a common tangent line would have total intersection
multiplicity at least four with a cubic.
If $P_1\ne Q_1$ and $P_2\ne Q_2$, then
\[
 A_P\cap A_Q\subseteq
 (S_{P_1}\cap S_{Q_1})\otimes(S_{P_2}\cap S_{Q_2}),
\]
whose dimension is one. A nonzero decomposable tensor in $A_P$
has either its first factor $e_{P_1}$ or its second factor
$e_{P_2}$; applying the same observation at $Q$ leaves only
\[
 e_{P_1}\in S_{Q_1},\ e_{Q_2}\in S_{P_2},
 \quad\text{or}\quad
 e_{Q_1}\in S_{P_1},\ e_{P_2}\in S_{Q_2}.
\]
For fixed $P$ these conditions permit only finitely many $Q$:
a tangent meets a cubic in a finite subscheme. For completeness,
tangents through a fixed point are finite for the specific
cubic here without any separability assumption on a Gauss map.
For $F=Y^2Z+XYZ+X^3+Z^3$, its partial derivatives are
\[
 F_X=YZ+X^2,\qquad F_Y=XZ,\qquad F_Z=Y^2+XY+Z^2.
\]
No nonzero constant linear combination is zero: compare the
coefficients of $X^2$, $Y^2$, and $XZ$. A nonzero quadratic
cannot vanish identically on the irreducible cubic. Hence for
every projective point $(a:b:c)$ the equation
$aF_X+bF_Y+cF_Z=0$ cuts a finite subscheme on $E$. Its points
are exactly those whose tangent contains $(a:b:c)$.
If, for example, $P_1=Q_1$ and $P_2\ne Q_2$, expanding in a
basis of $S_{P_1}$ beginning with $e_{P_1}$ shows
\[
 A_P\cap A_Q=e_{P_1}\otimes(S_{P_2}\cap S_{Q_2}),
\]
again of dimension one. This proves all stated rank bounds.
\end{proof}

\begin{theorem}[Uniform Hasse irreducibility]
\label{thm:product-Hasse-irred}
\status{Proved under stated hypotheses}
For the fixed surface and embedding, the degree-$d$ Hasse divisor
is geometrically irreducible for every $d\ge4$.
\end{theorem}

\begin{proof}
Work over an algebraic closure and put $e=d-3\ge1$.
Let $H_1,\ldots,H_t$ be the irreducible components of the
already reduced degree-three Hasse divisor in
$P_3=\mathbf P H^0(X,L^3)$.
Its general member on each $H_i$ is smooth, has stable loss one,
and has a one-dimensional Cartier kernel. Smoothness follows
because the singular-curve discriminant is irreducible and its
general member is ordinary: it contains an ordinary transverse
union of an ordinary degree-one and an ordinary degree-two
curve, so its nonempty ordinary open is dense.

For each $H_i$, choose a smooth general point $C$ and a point
$P\in C$ where its unique Cartier-killed differential is nonzero.
The tangent linear functional $dh_3|_C$ fails to be independent
of the first-jet functionals at at most one point of $C$.
Indeed first-jet evaluations at two distinct points are jointly
surjective by Lemma~\ref{lem:cubic-power-jets}; their dual spans
have zero intersection. Since $dh_3|_C\ne0$, it cannot belong
to both spans. We may therefore also impose this independence
at $P$, and require its two elliptic coordinates to avoid the
finite flex sets. These exclusions remove only finitely many
points of the smooth irreducible ample curve $C$.

Consider the incidence of $(C,P,D)$ over this open of pairs
$(C,P)$, with $D\in|L^e|$ tangent to $C$ at $P$.
The universal smooth curve over the irreducible open of $H_i$
is irreducible, since its fibres are connected smooth curves.
The two conditions on $D$ are independent linear first-jet
conditions, so the incidence is a projective bundle over an
irreducible base before removing open conditions.

The projection to $|L^e|$ is submersive on the chosen locus.
For precision, work on the affine cone over $H_i$: its tangent
space is $\ker dh_3$, and independence of $dh_3$ from the
first-jet functionals makes its map to the three-dimensional
first-jet space surjective. Projective scaling is harmless.
For an arbitrary infinitesimal change of $D$, move $P$ to
cancel its value change, using $dD(P)\ne0$. Then choose the
first jet of the change of $C$ inside $\ker dh_3$ to preserve
$C(P)=0$ and equality of tangent directions. These are exactly
the three incidence equations. The projection is consequently
dominant.

We also check nonemptiness of its good open, including $e=1$.
Fix a good $(C,P)$. If $e\ge2$,
Lemma~\ref{lem:product-small-degree-jets} gives the
usual independent conditions: second jets permit a nonzero
normal derivative and contact of exactly two at $P$;
a singularity of $D$ at another $Q$ imposes three further
conditions, while $Q$ varies in dimension two; another contact
with $C$ imposes two further conditions, while $Q\in C$ varies
in dimension one. Thus both bad loci are proper.

For $e=1$, second jets are still separated at the chosen $P$.
Outside the exceptional two fibres and finite set of
Lemma~\ref{lem:product-small-degree-jets}, the same
counts apply. On that exceptional locus, joint first-jet rank
is at least five. The two contact conditions at $P$ together
with three singularity conditions at $Q$ therefore have total
rank at least four; here $Q$ varies in dimension at most one.
Relative to the two fixed conditions, the bad locus again has
codimension at least one. For an additional contact at $Q\in C$,
the four conditions have rank at least three. The exceptional
$Q$ on $C$ are finite, since an irreducible ample curve contains
neither fibre through $P$. This bad locus is likewise proper.
These estimates prove that smooth $D$ with just the chosen
simple tangency and all other intersections transverse form a
nonempty open in the same irreducible incidence.

A nonempty open in an irreducible dominant incidence remains
dominant. It intersects the inverse image of the ordinary open
of $|L^e|$. Its image therefore contains a nonempty open of that
linear system. There are finitely many $H_i$, so a single general
smooth ordinary $D$ works for every $i$.

Fix this $D$, with equation $g$. On the multiplication subspace
$gP_3\subset\mathbf P H^0(L^d)$, the support of the Hasse divisor
is precisely
\[
                      H_1\cup\cdots\cup H_t\cup\Delta_D,
\]
as a divisorial statement. Here $\Delta_D$ is the irreducible
simple-contact divisor of
Theorem~\ref{thm:product-all-degree-transvections}.
This follows from the union stable-rank formula; exceptional
nonreduced members and members containing $D$ have codimension
at least two. The generic points of these displayed divisors
are distinct, and there is an ordinary member outside them.

For each $i$ the incidence just constructed gives a point of
$H_i\cap\Delta_D$ where the full Hasse divisor is smooth, by
Lemma~\ref{lem:mixed-tacnode-Hasse}. Any global Hasse component whose restriction contains either
$H_i$ or $\Delta_D$ contains this chosen intersection point,
because it contains that entire closed prime divisor on $gP_3$.
At a smooth point of the full Hasse divisor there is only one
such global component. Thus all components assigned to either
prime are the same. This argument allows arbitrary pullback
multiplicities and does not assume unique assignments in
advance.
Finally, every global Hasse component restricts nontrivially
to $gP_3$: a nonzero homogeneous polynomial of positive degree
cannot become a nonzero constant on this projective linear
subspace, and the subspace contains an ordinary point.
Thus its restriction contains one of the displayed prime
divisors. All have just been assigned to a single component,
so the full Hasse divisor is irreducible.
\end{proof}


\begin{corollary}[All full groups on the fixed ordinary product]
\label{thm:product-char2-all-full-groups}
\status{Proved under stated hypotheses}
Let $k$ be any perfect field of characteristic two, let
\[
 E:\ Y^2Z+XYZ+X^3+Z^3=0,\qquad
 X=E\times E,\qquad
 L=\mathcal O_E(3O)\boxtimes\mathcal O_E(3O),
\]
and use the fixed Segre embedding. For every $d\ge1$, the
maximal relative Hasse--Witt rank is $r_d=9d^2-1$. On the smooth
ordinary degree-$d$ hypersurface-section open, the entire
geometric and arithmetic images coincide and are
\[
 G_{{\rm geom},d}=G_{{\rm arith},d}=
 \begin{cases}
 \GL_{9d^2-1}(\mathbf F_2),&d\text{ odd},\\[2mm]
 \left\{\begin{pmatrix}A&B\\0&I_3\end{pmatrix}:
 A\in\GL_{9d^2-4}(\mathbf F_2),\quad
 B\in M_{(9d^2-4)\times3}(\mathbf F_2)\right\},&d\text{ even}.
 \end{cases}
\]
In even degree this is exactly the pointwise stabilizer of all
three independent invariant covectors, and its order is
$2^{3(9d^2-4)}|\GL_{9d^2-4}(\mathbf F_2)|$.
The representation is irreducible exactly in odd degree.
\end{corollary}

\begin{proof}
The generic rank and exact invariant covectors were proved in
Theorem~\ref{thm:product-char2-all-parities-fixed-quotient} and
the preceding product rank calculation. Degrees one and two
are Theorem~\ref{thm:product-char2-two-full-groups}; degree three
is Corollary~\ref{cor:product-degree-three-full-group}.
For $d\ge4$, Theorem~\ref{thm:product-Hasse-irred} discharges the
geometric hypothesis in
Theorem~\ref{thm:product-full-group-if-hasse-irreducible}.
Its proof supplies the entire linear group in odd degree and
all translations as well as the full linear kernel group in
even degree. The arithmetic and geometric equality persists
under any perfect extension of $\mathbf F_2$: the finite
geometrically connected covers giving the geometric images
remain connected after algebraically closed field extension,
and the three covectors in even degree are already defined
over $\mathbf F_2$. Finally the natural full linear module is
irreducible, while the even-degree common kernel of the three
covectors is a nonzero proper invariant subspace.
\end{proof}


% /root/monodromy, 2026-11-19. Proved; independent root and reduction_descent audits passed.
\subsection{Uniform Hasse geometry for every ordinary elliptic product}
\label{subsec:ordinary-products-general-structure}
\status{Proved and checked. The parent and reduction/descent agent independently audited the complete proof. The universal polynomial symmetrizer calculation was independently rerun by both agents. Full group exclusion is treated separately.}

\begin{lemma}[A universal elliptic Cartier symmetrizer]
\label{lem:ordinary-elliptic-universal-symmetrizer}
\status{Proved under stated hypotheses}
Let $E:y^2+xy=x^3+A x^2+B$ over a perfect field of characteristic
two, with $B\ne0$, and put $N=\mathcal O_E(3O)$.
In the basis $(1,x,y)$ and with the Cartier-invariant differential
$\eta=dx/x$, the root-removed Cartier multiplication matrix is
\[
 M(s)=\begin{pmatrix}
 s_0&B s_2&B s_1\\
 s_2&s_1+A s_2&s_0+A s_1\\
 0&s_2&s_1+s_2
 \end{pmatrix}.
\]
The constant invertible matrix
\[
 R=\begin{pmatrix}1&0&B\\0&0&B\\0&B&(1+A)B\end{pmatrix},
 \qquad\det R=B^2,
\]
satisfies
\[
 RM(s)=\begin{pmatrix}
 s_0&0&B s_2\\
 0&B s_2&B(s_1+s_2)\\
 B s_2&B(s_1+s_2)&B s_0+B s_1+(1+A)B s_2
 \end{pmatrix}.
\]
Thus $RM(s)$ is symmetric and its diagonal is $R^t s$.
For any pair of ordinary elliptic curves with degree-three
line bundles, the analogous nine-dimensional Cartier pencil
has a constant invertible symmetrizer and an invertible
linear map to its nine diagonal entries.
\end{lemma}

\begin{proof}
The smooth ordinary Weierstrass equation can be put in the
displayed form by first normalizing its nonzero $xy$
coefficient, then translating $x$ to remove the $y$ term and
translating $y$ by a constant to remove the $x$ term.
These changes preserve $O$ and the degree-three line bundle.
The discriminant in the resulting equation is $B$, so $B\ne0$.
A degree-three line bundle is a translate of $\mathcal O(3O)$
over the algebraic closure, which suffices for the geometric
assertion.

Writing a rational function as $a^2+x b^2$ gives
$\mathcal C((a^2+x b^2)\eta)=a\eta$.
The Weierstrass relation consequently gives
\[
 \mathcal C(1)=1,\quad \mathcal C(x)=0,\quad
 \mathcal C(y)=x,\quad
 \mathcal C(xy)=\sqrt B+\sqrt A\,x+y.
\]
Together with $\mathcal C(x^2)=x$ and
$\mathcal C(y^2)=y$, this is precisely the matrix written above
when the output scalar square root is removed.
Direct multiplication proves the formulas for $R$.
The matrix $R$ itself need not be symmetric; only $RM(s)$ is.

On a product the Cartier multiplication coefficients are tensor
products of the two one-factor coefficient matrices. Therefore
$R_1\otimes R_2$ symmetrizes all nine coefficient matrices.
Its diagonal map is $R_1^t\otimes R_2^t$, and its determinant
is $B_1^6B_2^6\ne0$.
The universal polynomial identities over
$\mathbf F_2[A_1,B_1,A_2,B_2]$ are independently reproducible in
\path{agents/monodromy/code/ordinary_product_universal_symmetrizer.py};
the complete retained output is
\path{agents/monodromy/outputs/ordinary_product_universal_symmetrizer.json}.
\end{proof}

\begin{lemma}[Squarefreeness from independent symmetric diagonals]
\label{lem:symmetric-independent-diagonal-squarefree}
\status{Proved under stated hypotheses}
Let $S(z)$ be an $m\times m$ symmetric matrix of linear forms
in characteristic two, and suppose its diagonal entries are
independent linear coordinates $z_1,\ldots,z_m$.
Then $\det S$ is geometrically squarefree.
\end{lemma}

\begin{proof}
In the determinant expansion, permutations having a cycle of
length at least three cancel with the inverse of the first such
cycle. The surviving terms are indexed by matchings:
\[
 \det S=\sum_{\mathcal M}
 \prod_{\{i,j\}\in\mathcal M}S_{ij}^2
 \prod_{i\text{ unmatched}}z_i.
\]
Only the empty matching contributes to the monomial
$z_1\cdots z_m$, and its coefficient is one. Every term of a
nonempty matching includes a square of a linear form and hence
has some exponent at least two in each of its monomials.
If a nonconstant polynomial $q$ had $q^2\mid\det S$ over an
algebraic closure, write $\det S=q^2r$ with $\deg r<m$.
Every exponent in $q^2$ is even. A term contributing to
$z_1\cdots z_m$ would therefore require a monomial of $r$
with all $m$ exponents odd, whose degree is at least $m$.
This is impossible. Thus no irreducible factor is repeated.
\end{proof}

\begin{lemma}[Jet bounds for arbitrary ordinary cubic factors]
\label{lem:ordinary-product-uniform-jets}
\status{Proved under stated hypotheses}
Let $X=E_1\times E_2$ with both elliptic curves ordinary in
characteristic two and $L=\mathcal O_{E_1}(3O)\boxtimes
\mathcal O_{E_2}(3O)$. The conclusions of
Lemma~\ref{lem:product-small-degree-jets} hold on $X$.
In addition, the singular-curve discriminant of $|L|$ is
irreducible with a general one-node member, and its general
member is ordinary.
\end{lemma}

\begin{proof}
The tensor first-jet proof of that lemma uses only the plane
cubic embeddings. Its one finiteness point is also uniform:
for $F=Y^2Z+XYZ+X^3+A X^2Z+B Z^3$ the partials are
\[
 F_X=YZ+X^2,\qquad F_Y=XZ,\qquad
 F_Z=Y^2+XY+A X^2+B Z^2.
\]
Their constant linear combinations are linearly independent,
by comparing $Y^2,X^2,XZ$. No nonzero quadratic vanishes
identically on the cubic, so the tangents through any fixed
point form a finite set. Thus the exceptional rank-five pairs
are the pairs sharing one elliptic coordinate, of dimension
three in $X\times X$, together with the crossed incidence
locus, which is finite over the first point and has dimension
at most two. Elsewhere the two first jets have rank six.

A specified singular point imposes three independent conditions,
so its incidence is an irreducible projective bundle over $X$.
An ordinary node occurs in a nonempty open of every fibre:
products of sections vanishing simply in the two elliptic
directions provide an arbitrary nonzero mixed quadratic
coefficient. For two singular points, the rank-six locus gives
codimension $6-4=2$ after the points vary. The exceptional
rank-five loci give codimensions at least $5-3=2$ and $5-2=3$.
Hence the general discriminant member has exactly one node,
and nonreduced members have codimension at least two.

Finally $|L|$ contains an ordinary grid: take three distinct
vertical fibres with sum $3O$ on $E_1$ and three distinct
horizontal fibres with sum $3O$ on $E_2$.
Its normalization components are ordinary elliptic curves and
its graph term is Frobenius-bijective, so the grid is ordinary.
The irreducible discriminant thus contains an ordinary point,
and its ordinary open is nonempty and dense.
\end{proof}

\begin{theorem}[Uniform irreducible Hasse divisors on ordinary products]
\label{thm:ordinary-products-Hasse-irreducible}
\status{Proved under stated hypotheses}
Let $E_1,E_2$ be arbitrary ordinary elliptic curves over a
perfect field of characteristic two, and let
$L=\mathcal O_{E_1}(3O)\boxtimes\mathcal O_{E_2}(3O)$.
The full Hasse divisor in $|L|$ is geometrically reduced and
has generic Cartier corank one on every component.
For every $d\ge2$, the full Hasse divisor in $|L^d|$ is
geometrically irreducible and reduced, with generic stable
rank loss exactly one. Consequently the geometric relative
Hasse--Witt group in degree $d\ge2$ is generated by
transvections.
\end{theorem}

\begin{proof}
The two preceding algebraic lemmas make the degree-one Hasse
polynomial geometrically squarefree. Its derivative is nonzero
at each generic component point, so its matrix has corank one
there. By Lemma~\ref{lem:ordinary-product-uniform-jets}, no
component is the singular-curve discriminant, whose generic
point is ordinary. Thus the generic curve on each degree-one
Hasse component is smooth with a one-dimensional Cartier
kernel. Stable loss one in this degree is not assumed or
needed in the following argument.

Fix $d\ge2$, put $e=d-1$, and repeat the mixed-tacnode proof
with $C\in|L|$ and a general ordinary $D\in|L^e|$.
We spell out the differences caused by using $|L|$ in place
of $|L^3|$.
For a smooth general $C$ on a base Hasse component, the
nonzero tangent functional $dh_1$ belongs to the dual
first-jet space at only finitely many points of $C$.
Indeed, if there were infinitely many, choose one such $P$;
all but finitely many other points $Q$ on $C$ avoid the two
fibres and the finite crossed set through $P$.
The first-jet dual spaces at $P$ and such a $Q$ are disjoint,
contradicting that the same nonzero $dh_1$ lies in both.
We can therefore choose $P$ avoiding those finitely many
points, the zeros of the Cartier-killed differential, and the
flex fibres. The affine tangent cone $\ker dh_1$ maps
surjectively to first jets at this $P$.

The incidence with $D$ tangent to $C$ at $P$ consequently
projects dominantly to $|L^e|$, by the same three-equation
infinitesimal calculation in
Theorem~\ref{thm:product-Hasse-irred}. The good open of smooth
$D$ with exactly one simple contact is nonempty. For $e\ge2$
use full two-point first jets; for $e=1$ use rank at least five
and its exceptional one-dimensional locus through $P$.
The latter meets $C$ in only finitely many points. Exactly the
same dimension counts as in that theorem exclude a second
contact or a singularity of $D$.
Hence a single general ordinary $D$ admits such a mixed
contact with every base Hasse component.

For completeness the fixed-$D$ contact divisor is still
irreducible with generic ordinary $C$ despite the weaker
base jets. Its incidence over $P\in D$ is a projective bundle
cut out by two independent conditions. Two contacts impose
four conditions on the rank-six locus of pairs in $D^2$;
the rank-five exceptional pairs have dimension at most one
and impose at least three. Both give codimension at least
two in $|L|$. A singularity of $C$ at a second point imposes
five conditions generically, with the point pair varying in
dimension three; on the exceptional locus the rank is at
least four and the dimension is at most two. A singularity
at the contact point also has codimension at least two.
Order-three contact imposes three conditions at a nonflex
point; the flex fibres meet $D$ in finitely many points, so
even there the two contact conditions give codimension two.
The generic contact is therefore simple with $C$ smooth.
The joint tangency incidence for $(P,C,D)$ is irreducible
and dominates both complete linear systems by first jets.
Its two ordinary inverse images are nonempty dense opens,
so the generic contact $C$ is ordinary for a general ordinary
$D$.

On the multiplication subspace $g_D|L|$, the divisorial
nonordinary support is the union of the base Hasse components
and this contact divisor, by the exact union formula.
The exceptional nonreduced base locus has codimension at least
two; when $e=1$, $C=D$ is a single point of $\mathbf P^8$,
and when $e>1$ no base member contains $D$.
At each mixed contact, Lemmas~\ref{lem:pointed-Cartier} and
\ref{lem:mixed-tacnode-Hasse} make the full Hasse divisor
smooth. Their proofs use only ordinarity of the ambient product,
not the particular elliptic equation, and require only Cartier
corank one on $C$, not stable rank loss one on $C$. Therefore every base
Hasse component and the contact divisor are assigned to the
same global Hasse component. Every global component has a
positive-degree nonzero restriction to this linear subspace,
which contains an ordinary point. The component-merging
argument of Theorem~\ref{thm:product-Hasse-irred} proves that
there is exactly one global component. The smooth mixed
point also makes its multiplicity one.

There is an ordinary pair $(C,D)$ with a single simple
contact. The union formula gives stable loss one at that
point of the now irreducible Hasse divisor. The closed
iterated-Frobenius rank loci imply generic stable loss at
most one, and noninvertibility makes it exactly one.
Lemma~\ref{lem:char2-defect-one-inertia} proves the last assertion.
\end{proof}

\begin{corollary}[An irreducible kernel for every ordinary pair]
\label{cor:ordinary-products-irreducible-kernel}
\status{Proved under stated hypotheses}
In the preceding theorem let $d\ge2$, put $b=9d^2$, and set
$q=0$ for odd $d$ and $q=3$ for even $d$.
The entire invariant-covector space has dimension $q$.
Its common kernel $K\subset V_d$ is irreducible of dimension
$b-1-q$, and $V_d$ has no invariant vector.
The image on $K$ is irreducible and generated by transvections.
\end{corollary}

\begin{proof}
The fixed-covector calculation is already
Theorem~\ref{thm:product-char2-all-parities-fixed-quotient}.
Its statement and proof allow distinct ordinary factors: the
one-factor Frobenius pushforward calculations and their external
tensor product do not require the elliptic curves to be isomorphic.
The relative coherent bundle on $\mathbf P^{b-1}$ is
$\Omega^1_{\mathbf P^{b-1}}$, by the Euler/adjunction sequence
and the ordinary one-dimensional ambient $H^2(\mathcal O_X)$.
Theorem~\ref{thm:ordinary-products-Hasse-irreducible} therefore
permits Lemma~\ref{lem:irreducible-hasse-quotients-trivial}.
Every nonzero proper invariant subspace has trivial quotient,
so contains the common kernel. This makes $K$ irreducible and
rules out invariant lines, since $\dim K\ge32$.
The transvection assertion follows by restriction; a nontrivial
generator restricts either to a transvection or to the identity.
\end{proof}


% /root/monodromy, 2026-11-19. Proved; independent root and reduction_descent audits passed.
\subsection{An isogeny construction of a uniform linear block}
\label{subsec:ordinary-products-isogeny-full}
\status{Proved and checked. The parent and reduction/descent agent independently audited the geometric generic-fibre argument, nodal-image incidence, group extraction, arithmetic fixed covectors, and effective isogeny threshold. The original Chai theorem and Stacks generic-fibre lemma were opened by both auditors; the supporting exact computations were independently rerun.}

\begin{lemma}[An open locus of full degree-one images]
\label{lem:ordinary-pair-GL8-open}
\status{Proved under stated hypotheses}
There is a nonempty Zariski open subset $B^\circ$ of the space
$B=\mathbf G_m^2$ of pairs
\[
 E_{b_i}:y^2+xy=x^3+b_i
\]
on which the geometric relative degree-one image for
$\mathcal O(3O)\boxtimes\mathcal O(3O)$ is
$\GL_8(\mathbf F_2)$.
\end{lemma}

\begin{proof}
The nine sections $(1,x,y)\otimes(1,x,y)$ trivialize the
coefficient projective bundle over $B$.
Let $U\subset\mathbf P^8\times B$ be the smooth ordinary open.
It has nonempty geometrically integral fibres.
The relative rank-eight Lang frame cover $Z\to U$ is finite
etale of degree $|\GL_8(\mathbf F_2)|$.
At $b_1=b_2=1$ its geometric fibre is connected by
Theorem~\ref{thm:product-char2-two-full-groups}.

We justify passage to the geometric generic fibre, since $U/B$
is not proper. Suppose $Z_{\bar\eta}$ were disconnected.
After a finite separable extension of $k(B)$ all its components
are defined. Normalize $B$ in that extension; the resulting
integral normal scheme $B'$ is finite surjective over $B$: the
base is of finite type over a field, hence excellent and has
finite normalization in a finite field extension.
The scheme $U'=U\times_B B'$ is a nonempty open of a projective
bundle over a normal integral base, hence normal and integral.
The finite etale cover $Z'=Z\times_B B'$ is normal, so its
irreducible components are disjoint open-and-closed subschemes.
Each component is finite etale over $U'$ and has nonempty image,
which is both open and closed; it therefore surjects onto $U'$.
The distinct geometric generic components consequently all
meet the fibre over a point of $B'$ lying above $(1,1)$.
This contradicts geometric connectedness of that fibre.
Thus $Z_{\bar\eta}$ is connected. It is smooth, being finite
etale over an open of geometric projective space, so it is
normal; distinct irreducible components of a normal Noetherian
scheme are disjoint. It is therefore irreducible, as required.

By \cite[Lemma~37.27.5]{Stacks0559}, a finite-type morphism with
geometrically irreducible generic fibre has geometrically
irreducible fibres after restricting the integral base to a
nonempty open. Apply this to $Z\to B$.
Connectedness of a full frame cover is precisely transitivity
on all frames, hence equality of the monodromy image with
$\GL_8(\mathbf F_2)$. This gives $B^\circ$.
\end{proof}

\begin{lemma}[Ordinary nodal images under odd product isogenies]
\label{lem:ordinary-product-isogeny-nodal-images}
\status{Proved under stated hypotheses}
Let $\phi_i:E_i'\to E_i$ be odd-degree isogenies of ordinary
elliptic curves in characteristic two, of degrees $m_i$, and
put $\phi=\phi_1\times\phi_2$.
For a general smooth ordinary
$C'\in|\mathcal O_{E_1'}(3O)\boxtimes\mathcal O_{E_2'}(3O)|$,
its image $C=\phi(C')$ is an ordinary nodal integral curve,
$C'\to C$ is its normalization, and
\[
 \mathcal O_X(C)=
 \mathcal O_{E_1}(3m_2O)\boxtimes\mathcal O_{E_2}(3m_1O).
\]
Over a suitable dense open of the $C'$ parameter space, the
permutation image on the nodes of $C$ is a subquotient of a
product of copies of $S_{18}$.
\end{lemma}

\begin{proof}
Work over an algebraic closure, so the odd finite etale kernel
of $\phi$ is a constant group. Its translations act freely.
A general $C'$ is not preserved by any nonzero kernel
translation. If such a translation preserves the linear system,
its projective action on it is nontrivial: a scalar action
would act identically on the very ample embedding of the
surface, contradicting that the translation is nonzero.
Its fixed locus is therefore proper. If it does not preserve
the line bundle, it cannot preserve a member.
Thus $\phi|_{C'}$ is generically one-to-one.

For a nonzero kernel point $g$, the possible double points
are parametrized by $P\in C'\cap(C'-g)$, and this intersection
has degree $L'^2=18$. We show it is transverse for a general
$C'$. The two conditions $C'(P)=C'(P+g)=0$ are independent.
If both coordinates of $g$ are nonzero, the two first jets
have rank six except at finitely many $P$: the crossed tangent
conditions reduce to equations $3P_i=\text{constant}$.
On the rank-six locus, proportionality of the two nonzero
gradients is one further independent equation. The bad
incidence has codimension three with $P$ varying in dimension
two, and hence has proper image in $\mathbf P^8$.
The finitely many exceptional $P$ already impose two independent
value conditions, so their image is proper as well.

If exactly one coordinate of $g$ is zero, the joint first-jet
rank is five everywhere. Its sole relation involves only the
values and the two derivatives along the factor on which
$g$ translates; the transverse derivatives remain independent.
After the two values are set to zero, the relation on the two
parallel derivatives is nonzero, since two distinct evaluation
directions are independent. Thus the determinant of the two
gradients is not identically zero on their three-dimensional
allowed derivative space: one may choose a nonzero parallel
derivative and then an independent transverse derivative.
Tangency is again one extra equation. The same codimension
three, two-dimensional-point count excludes it generically.

Triple points are excluded by three-point values.
Three distinct Segre rank-one tensors are collinear only when
one factor is constant on all three and their other factor
points are collinear. For $P,P+g,P+h$ with distinct kernel
translations, this exceptional case requires both translations
to occur in the same elliptic factor, and then collinearity is
$3P_i+g_i+h_i=0$. Its locus in the surface has dimension one.
Outside it the three evaluations have rank three, while on it
they have rank two. The corresponding incidences have dimensions
at most $8-3+2=7$ and $8-2+1=7$, respectively.
Only finitely many kernel pairs are involved, so a general
$C'$ has no triple identifications. Etaleness of $\phi$ and
transversality now show that its image has only ordinary nodes.

The image class is the norm of $L'$.
Pushing the divisor $3(\{O\}\times E_2')+
3(E_1'\times\{O\})$ multiplies the first component by $m_2$
and the second by $m_1$, which gives the displayed exact
line bundle, with no unspecified degree-zero twist.
Ordinarity follows from the normalization exact sequence:
its graph term is Frobenius-bijective and $C'$ is ordinary.
Finally each $g$ gives an etale intersection cover of degree
$18$. The unordered nodes are quotients of the union of these
covers, with the constant identification of $g$ and $-g$.
Therefore the node permutation image is a subquotient of a
product of $S_{18}$'s, as claimed.
\end{proof}

\begin{lemma}[A natural eight-dimensional block after an isogeny]
\label{lem:ordinary-product-isogeny-GL8-block}
\status{Proved under stated hypotheses}
Suppose the degree-one relative image on $E_1'\times E_2'$
is $\GL_8(\mathbf F_2)$, and let $\phi_i$ be as above.
For every
\[
                    d\ge \max(m_1,m_2)+1,
\]
the degree-$d$ geometric group on $E_1\times E_2$ contains
an actual natural $\GL_8(\mathbf F_2)$ block fixing a
complementary summand pointwise.
\end{lemma}

\begin{proof}
Choose a fixed reduced vertical divisor of degree $3(d-m_2)$
linearly equivalent to $3(d-m_2)O$ on $E_1$, and a fixed
reduced horizontal divisor of degree $3(d-m_1)$ similarly.
These degrees are at least three, so such divisors can be
chosen in general position. Let $D$ be the resulting grid
of elliptic fibres on $X$.
It is ordinary and nodal. By the preceding line-bundle formula,
$Y=\phi(C')\cup D$ is a member of $|L^d|$.
For $C'$ in a dense open, it meets $D$ transversely, avoids its
nodes, and its own image nodes avoid $D$.
The last avoidance is also an elementary incidence condition:
for a fixed image fibre, the two value conditions at $P,P+g$
have $P$ varying in dimension one, and so have proper image
in the coefficient projective space.
All these $Y$ are ordinary nodal curves.

On this boundary family the relative Lang representation has
a filtration
\[
       0\subset T\subset T_0\subset V_d,
 \qquad T_0/T\text{ constant},\qquad
 V_d/T_0=Q,\quad\dim Q=8,
\]
where $T$ is the graph contribution and $Q$ is the degree-one
relative system of $C'$.
Indeed odd-degree isogenies identify the ambient
$H^1(X,\mathbf F_2)$ with $H^1(X',\mathbf F_2)$, and the
normalization exact sequence gives this filtration.
The image on $Q$ remains full after shrinking the parameter
space.

The image nodes of $C'$ give degree-$18$ permutation covers.
An intersection with a fixed image fibre splits upstairs into
fixed fibres of the corresponding odd isogeny; on each such
fibre, $C'$ cuts a divisor of degree three.
Their permutation images are therefore subgroups of products
of $S_3$. The nodes of $D$ itself are constant.
It follows that the image on $T$ has order with prime factors
at most $17$. The whole $T_0$ image is an extension of this
by a subgroup of the additive two-group
$\Hom(T_0/T,T)$, so the same prime bound holds.

The simple group $\GL_8(\mathbf F_2)$ has order divisible by
$31$ (the factor $2^5-1$ in its order).
Goursat therefore makes the kernel of the entire $T_0$ action
still surject onto $\GL(Q)$.
Inside its block parabolic, the translation kernel is
$Q^*\otimes T'$ for a subspace $T'\subset T_0$.
The elementary $H^1(\GL_8(\mathbf F_2),Q^*)=0$ calculation
of Lemma~\ref{lem:char2-group-inputs} splits the remaining
cocycle after conjugation, exactly as in
Lemma~\ref{lem:product-degree-three-block}.
Thus this subgroup contains the stated block-diagonal
$\GL_8$, not merely a quotient isomorphic to it.
The boundary lies in the full coherent ordinary open; its
image is a subgroup of the same monodromy image as the dense
smooth ordinary open.
\end{proof}

\begin{theorem}[Eventual full images for every ordinary pair]
\label{thm:ordinary-products-eventual-full}
\status{Proved under stated hypotheses}
Let $E_1,E_2$ be any ordinary elliptic curves over a perfect
field of characteristic two, with
$L=\mathcal O_{E_1}(3O)\boxtimes\mathcal O_{E_2}(3O)$.
There exist odd-degree isogenies
$\phi_i:E_i'\to E_i$, of degrees $m_i$, such that for every
$d\ge d_1=\max(m_1,m_2)+1$ the entire geometric and arithmetic
images are
\[
 \begin{cases}
 \GL_{9d^2-1}(\mathbf F_2),&d\text{ odd},\\
 \mathbf F_2^{3(9d^2-4)}\rtimes\GL_{9d^2-4}(\mathbf F_2),
       &d\text{ even},
 \end{cases}
\]
where the second line is the full pointwise stabilizer of the
three independent invariant covectors.
The isogenies, and hence the displayed integer $d_1$, can be
found by a terminating finite procedure over a finite field or
an effectively presented finitely generated field of definition.
No degree-one full-image assertion for an arbitrary pair is
being inferred from the generic pair.
\end{theorem}

\begin{proof}
By \cite[Theorem~2, printed p.~477]{Chai1995}, for any prime
$\ell\ne2$ the $\ell$-power Hecke orbit of an ordinary elliptic
curve is Zariski dense in its one-dimensional moduli space.
We use only the case $g=1$, with $\ell=3$.
The product of the two such dense orbits meets the nonempty
open $B^\circ$ of Lemma~\ref{lem:ordinary-pair-GL8-open}.
Choose a pair there; reversing an isogeny by its dual gives
isogenies $\phi_i:E_i'\to E_i$ of odd degrees $m_i$.
Geometric base extensions needed to define them do not change
the geometric image.

For $d\ge d_1$, Corollary~\ref{cor:ordinary-products-irreducible-kernel}
gives an irreducible common kernel $K$, no invariant vectors,
and transvection generation.
Lemma~\ref{lem:ordinary-product-isogeny-GL8-block} provides an
actual natural $\GL_8$ summand inside $K$, fixing its complement.
This excludes a nondegenerate invariant alternating form on
$K$: the natural $\GL_8$ has no invariant bilinear form and
no pairing with a trivial module, so this summand would lie
in the radical. The opened transvection classification
\cite[Theorem~3.6, printed p.~22]{KantorSeress2001} therefore
forces the full $\GL(K)$.
The translation-kernel and $H^1$ argument already used in
Theorem~\ref{thm:product-full-group-if-hasse-irreducible} gives
all translations when the fixed quotient has dimension three.
This proves the geometric images.

The three even-degree covectors are defined over the original
perfect field. Indeed the root line bundles
$\mathcal O_{E_i}(3dO/2)$ are defined there, and the ordinary
characteristic-two splitting $F_*\mathcal O_{E_i}=\mathcal O\oplus T_i$
is defined there: its nontrivial two-torsion character line is
unique. Normalize the invariant differential to $dx/x$ in the
ordinary Weierstrass form, for which Cartier is the identity.
The resulting four product projectors and their three quotient
covectors are then arithmetic invariants. The arithmetic image
is between the proved geometric group and its same pointwise
stabilizer, so the two images coincide.

For the effective assertion, enumerate finite chains of
three-isogenies from each input elliptic curve. This is a finite
branching enumeration: at each step solve the three-division
polynomial, enumerate its finite etale cyclic subgroups, and
form their quotients; reverse the resulting chains by duals.
Run the exact Lang-frame geometric-image algorithm of
Sections~\ref{subsec:effective-entire-monodromy} and
\ref{subsec:perfect-field-monodromy-data} on each resulting pair
in degree one, and stop at the first full $\GL_8$ image.
Dovetail the finite computations. Every individual test
terminates, and Hecke density plus the nonempty open proves
that some enumerated pair passes. Its explicitly retained
isogeny degrees give $d_1$. This is a terminating construction,
not an unspecified Galois-group oracle. For an arbitrary
perfect field with no effective coefficient presentation the
same argument proves existence; the calculation is performed
on a finitely generated field of definition when such a
presentation is supplied.
\end{proof}


\begin{remark}[Finite supporting checks]\label{auto:A_monodromy:3a04f9cdd890}
\status{Proved under stated hypotheses}

The exact program
\path{agents/monodromy/code/ordinary_product_isogeny_incidence_checks.py}
checks all $18424$ triples of distinct rank-one $3\times3$
matrices over $\mathbf F_2$, including all $98$ collinear
triples, against the common-factor criterion. It checks that
each of the three nonzero parallel-gradient relations admits
a transverse pair of gradients, and checks the prime $31$
separation from degree-$18$ node covers. The output is
\path{agents/monodromy/outputs/ordinary_product_isogeny_incidence_checks.json}.
These are finite supporting calculations; the universal
incidence assertions are proved in the preceding lemmas.
\end{remark}


\subsection{An exact parity formula for supersingular elliptic products}
\label{sec:DE-A3-ss-product}

The following family gives an explicit answer in every degree, and shows that
eventual relative ordinarity can fail even when the ambient Picard scheme is
smooth and its crystalline cohomology is torsion-free.  It does not purport to
determine the generic rank for an arbitrary fixed ambient variety.

Throughout this subsection let $k=\overline{\mathbf F}_2$, let
\[
 E_0:\quad y^2+y=x^3,
 \qquad X=E_0\times E_0,
 \qquad L=\OO_{E_0}(3O)\boxtimes\OO_{E_0}(3O),
\]
and embed $X\hookrightarrow\mathbf P^8_k$ by the product of the plane cubic
embeddings followed by the Segre embedding.  Write
$N_m=\OO_{E_0}(mO)\boxtimes\OO_{E_0}(mO)$, so $L^d=N_{3d}$.
For a smooth member $Y\in|N_m|$, put
\[
 Q_Y=\coker\bigl(H^1(X,\OO_X)\longrightarrow H^1(Y,\OO_Y)\bigr).
\]
Frobenius on these vector spaces is $a\mapsto a^2$-semilinear; ``stable rank''
means the dimension of the eventual image of its iterates over the perfect
field $k$.

\begin{lemma}[Explicit endomorphisms of every positive degree]
\label{lem:DE-A3-ss-norm}
\status{Proved}
For every integer $n\ge1$ there is an explicitly specified endomorphism
$b:E_0\to E_0$ with
$b^\dagger b=bb^\dagger=[n]$, where $\dagger$ is the Rosati involution for
$\OO_{E_0}(O)$.  In particular $\deg b=n$.
\end{lemma}
\begin{proof}
Choose $\omega\in\mathbf F_4$ with $\omega^2+\omega+1=0$, and define
\[
 i(x,y)=(x+1,y+x+\omega),\qquad
 j(x,y)=(x+\omega,y+\omega^2x+\omega).
\]
For $a\in\mathbf F_4^\times$ and $c^2+c=1$, substitution shows that
$(x,y)\mapsto(x+a,y+a^2x+c)$ preserves the equation, because $a^4=a$ and
$a^3=1$.  These maps extend to automorphisms fixing $O$.  Direct composition
gives $i^2=j^2=[-1]$ and $ij=-ji$, with $[-1](x,y)=(x,y+1)$.
Since they fix the polarizing divisor $O$, their Rosati adjoints are their
inverses.  Consequently
\[
 (a+bi+cj+dij)^\dagger(a+bi+cj+dij)
 =[a^2+b^2+c^2+d^2]
\]
for integers $a,b,c,d$; the same identity holds in the opposite order.

We recall a short proof that every positive integer is a sum of four squares,
so the construction does not depend on an unproved norm-surjectivity claim.
The quaternion norm identity makes products of such sums sums of four
squares.  For an odd prime $\ell$, the two sets of residues $a^2$ and
$-1-b^2$ each have $(\ell+1)/2$ elements.  They intersect.  Choosing
$|a|,|b|\le(\ell-1)/2$ gives
$a^2+b^2+1=k\ell$ with $0<k<\ell$.
Choose the smallest positive multiplier $k$ for which $k\ell$ is a sum of
four squares, and suppose $k>1$.  Write $k\ell=\sum a_i^2$ and choose
$x_i\equiv a_i\pmod k$ with $|x_i|\le k/2$.  Then
$\sum x_i^2=kt$ with $0\le t\le k$.  The case $t=0$ would imply
$k\mid\ell$, which is impossible.  If $t<k$, all four coefficients of the
quaternion product $a\overline x$ are divisible by $k$, since modulo $k$
this product is $x\overline x=kt$.  Dividing those coefficients by $k$
expresses $t\ell$ as a sum of four squares, contradicting minimality.
If $t=k$, equality in the four bounds forces $k$ even and every
$a_i\equiv k/2\pmod k$.  Hence $k\mid4\ell$, so $k=2$ or $4$.
For $k=2$, the four odd squares have sum $4\pmod8$, whereas $2\ell$ is
$2$ or $6\pmod8$.  For $k=4$, the four squares are each $4\pmod{16}$,
whereas $4\ell$ is $4$ or $12\pmod{16}$.  Both are impossible.  Thus
$k=1$.  The prime $2$ and the integer $1$ are immediate, and multiplicativity
finishes the proof.  To specify $b$ without choices, take the first
nonnegative quadruple in lexicographic order with sum of squares $n$ and
use the displayed endomorphism.  Finally
$(\deg b)^2=\deg[n]=n^2$, so $\deg b=n$.
\end{proof}

\begin{lemma}[The ambient quotient and the even-degree obstruction]
\label{lem:DE-A3-ss-prym}
\status{Proved}
For every smooth $Y\in|N_m|$, $m\ge3$, its genus is $g=m^2+1$ and
$\dim Q_Y=m^2-1$.  Its $2$-rank $f_Y$ satisfies
\[
 f_Y\le m^2-1,
 \qquad
 f_Y\le m^2-3\quad\hbox{if }2\mid m.
\]
Moreover $\dim(Q_Y)_{\ss}=f_Y$.
\end{lemma}
\begin{proof}
The plane cubic is smooth: on the affine chart the partial derivative with
respect to $y$ is $1$, and at its unique point at infinity the homogeneous
equation has a nonzero partial derivative.  Its inverse map is
$(x,y)\mapsto(x,y+1)$, so $E_0[2](k)=0$.  Thus its $2$-rank is zero and
Frobenius on $H^1(E_0,\OO_{E_0})$ is zero.  The latter also follows from
Cartier: $\omega_0=dx$ is a regular invariant differential and
$C(dx)=0$.  Hence Frobenius is zero on $H^1(X,\OO_X)$, a space of
dimension two.

K\"unneth gives $H^1(X,P\otimes N_m^{-1})=0$ for every
$P\in\Pic^0(X)$: on each elliptic factor the two line bundles in question
have negative degree, so both possible $H^1$ tensor summands have a zero
$H^0$ factor.  In particular $H^0(Y,\OO_Y)=k$ and
$H^1(X,\OO_X)\hookrightarrow H^1(Y,\OO_Y)$.  Adjunction, or the same
exact sequence and Serre duality, gives $g=m^2+1$.

Restriction
\[
 j:\Pic^0(X)=E_0^2\longrightarrow J_Y
\]
is a closed immersion of group schemes.  Indeed a nontrivial
$P\in\Pic^0(X)(k)$ has $H^0(X,P)=0$.  The preceding vanishing and the
restriction sequence then give $H^0(Y,P|_Y)=0$, so $P|_Y$ is not trivial.
The tangent map of $j$ is the already proved injection on $H^1(\OO)$.
Its proper kernel therefore has only the identity as a geometric point and
has zero tangent space.  It is a finite local group scheme with zero
cotangent space, and Nakayama's lemma makes it trivial.

Identify $X$ with its dual by the product principal polarization.  The map
dual to $j$ is
\[
 q:J_Y\longrightarrow X,\qquad
 M\longmapsto
 \bigl(\operatorname{Nm}_{\pi_1}M,
       \operatorname{Nm}_{\pi_2}M\bigr),
\]
where $\pi_i:Y\to E_0$ are the projections.  They are finite of degree
$m$: $Y$ is smooth and connected, hence integral, and cannot contain an
elliptic fibre.  The diagonal entries of $qj$ are $[m]$ by norm and
pullback.  Each off-diagonal entry is zero.  For example the intersection
divisor on the first $E_0$ obtained by fixing a point on the second factor
always belongs to $|\OO_{E_0}(mO)|$, so the norms of the pullbacks of two
point divisors on the second factor are linearly equivalent.  We obtain
\[
 qj=[m]_X.                                                     \tag{\ast}
\]
In particular $q$ is surjective.

For completeness its kernel is connected and smooth, not merely connected
up to isogeny.  Let $P=(\ker q)^0_{\mathrm{red}}$.  Over the perfect field
$k$ this is an abelian subvariety.  The quotient $J_Y/P$ is an abelian
variety and $q$ factors through an isogeny $\nu:J_Y/P\to X$.
Dualizing, the kernel of $\nu^t$ lies in the kernel of $q^t=j$, which is
trivial.  The isogeny duality theorem gives
$\ker(\nu^t)\simeq(\ker\nu)^D$, so $\nu$ is an isomorphism.
Thus $\ker q=P$.  Here the exact source statements are the quotient
existence theorem and its abelian-variety example, and the isogeny duality
theorem \cite[Theorem~4.38, Example~4.40, Theorem~7.5]{EdixhovenVanDerGeerMoonenAV}.

The $2$-rank is additive in an exact sequence of abelian varieties.
One can see this without choosing an isogeny complement: on $k$-points,
the sequence of $2^a$-torsion groups is exact, since a lift of a torsion
point can be adjusted by a point of $P(k)$ using the surjectivity of
$[2^a]:P(k)\to P(k)$.  Counting points gives additivity.
Since $f_X=0$, we have $f_Y=f_P$ and $\dim P=g-2$.  This proves the first
bound.

If $m$ is even, (\ast) restricts to a closed immersion of finite group
schemes $j_2:X[2]\hookrightarrow P[2]$.  Its Frobenius twist maps
$X^{(2)}[2]$ into $P^{(2)}[2]$.  The group $\ker V_X$ lies in
$X^{(2)}[2]$, and finite-group Verschiebung is functorial for $j_2$
(or use $V_{J_Y}j^{(2)}=jV_X$ and the restriction of $V_{J_Y}$ to
$P^{(2)}$).  Consequently it embeds
\[
 \ker(V_X:X^{(2)}\to X)\lhook\joinrel\longrightarrow
 \ker(V_P:P^{(2)}\to P).
\]
The source is connected of order $4$, since $\deg V_X=2^2$ and $f_X=0$.
For any abelian variety $A$ of dimension $t$, $\deg V_A=2^t$ and
$\#\ker V_A(k)=2^{f_A}$: relative Frobenius is bijective on $k$-points
and $V_AF_A=[2]$.  The identity component of $\ker V_A$ therefore has
order $2^{t-f_A}$.  The displayed injection yields
$\dim P-f_P\ge2$, hence $f_Y\le g-4=m^2-3$.
These degree and Frobenius identities are exactly
\cite[Propositions~5.19--5.22]{EdixhovenVanDerGeerMoonenAV}; the argument
uses the connected subgroup, so it is not just a bound on geometric
$2$-torsion points.

Finally Artin--Schreier and the semilinear Fitting decomposition identify
the stable rank of $H^1(Y,\OO_Y)$ with $f_Y$.  The stable-part functor on
finite-dimensional semilinear vector spaces over a perfect field is exact:
split each space into its nilpotent and bijective parts, or apply a
sufficiently high Frobenius iterate to an exact sequence.  Quotienting by
the Frobenius-zero space $H^1(X,\OO_X)$ therefore leaves the stable rank
unchanged.  This proves the last assertion.
\end{proof}

\begin{theorem}[Exact generic rank in every hypersurface degree]
\label{thm:DE-A3-ss-product-rank}
\status{Proved}
For the specified $X\subset\mathbf P^8_{\mathbf F_2}$, the geometric
generic stable rank in degree $d\ge1$ is
\[
 \boxed{\quad
 \dim E(Y,\OO_Y)_{\ss,\mathrm{gen}}
 =\begin{cases}
 9d^2-1,&d\ \text{odd},\\
 9d^2-3,&d\ \text{even}.
 \end{cases}\quad}                                           \tag{\dagger}
\]
The coherent relative space has dimension $9d^2-1$ in every degree.
Consequently its generic Frobenius is bijective exactly for odd $d$.
There is no integer $d_0$ making it generically bijective for every
$d\ge d_0$.  The generic stable rank is positive in every degree, so
for this example the positivity threshold is exactly $d_0=1$.
\end{theorem}
\begin{proof}
First let $m\ge3$ be odd.  Choose the endomorphism $b$ in
Lemma~\ref{lem:DE-A3-ss-norm} with $b^\dagger b=bb^\dagger=[m-1]$ and put
\[
 \phi=\begin{pmatrix}1&b\\-b^\dagger&1\end{pmatrix}:X\longrightarrow X.
\]
Its Rosati adjoint satisfies $\phi^\dagger\phi=[m]I$; hence $\phi$ is an
isogeny of degree $m^2$ with kernel killed by the odd integer $m$.
It is therefore \'etale.  The inverse images of the two coordinate
elliptic divisors are
\[
 C_1=\{x+by=0\},\qquad C_2=\{-b^\dagger x+y=0\}.
\]
Each is a graph isomorphic to $E_0$.  They meet transversely in
$\ker\phi$, which consists of exactly $m^2$ reduced points.

We need linear equivalence with the fixed $N_m$.  The polarization of
$\OO(C_1+C_2)=\phi^*(\OO_{E_0}(O)\boxtimes\OO_{E_0}(O))$ is $[m]I$,
by the pullback formula for polarizations
\cite[Proposition~7.6]{EdixhovenVanDerGeerMoonenAV}.
This is also the polarization of $N_m$, so their difference belongs to
$\Pic^0(X)$ by \cite[Corollary~7.22]{EdixhovenVanDerGeerMoonenAV}.  Both line bundles are symmetric; on $\Pic^0(X)$ inversion
pulls a line bundle back to its inverse.  Their difference is thus a
$2$-torsion $k$-point of $\Pic^0(X)$.  There are no such nontrivial points,
since $E_0[2](k)=0$.  Hence $C_1+C_2\in|N_m|$ exactly.

For a connected nodal curve $D$ with smooth components $D_i$, the
normalization sequence gives a Frobenius-equivariant exact sequence
\[
 0\longrightarrow H^1(\Gamma_D,k)
 \longrightarrow H^1(D,\OO_D)
 \longrightarrow\bigoplus_iH^1(D_i,\OO_{D_i})\longrightarrow0,
\]
where $\Gamma_D$ is its dual graph and Frobenius on the first term is
bijective.  Thus its stable rank is
$b_1(\Gamma_D)+\sum_i f_{D_i}$, and its genus minus stable rank is
$\sum_i(g(D_i)-f_{D_i})$.  Applied to $C_1\cup C_2$, this gives genus
$m^2+1$ and stable rank $m^2-1$.

The universal Cartier divisor over the projective space $|N_m|$ is proper
and flat.  Every nonzero section has connected divisor and
$h^1(\OO_D)=m^2+1$: this follows from
$H^1(N_m^{-1})=0$ and the surjection
$H^2(N_m^{-1})\to H^2(\OO_X)$ dual to the nonzero section.
Cohomology and base change therefore make $H^1(\OO_D)$ a vector bundle
on the parameter space.  Its relative Frobenius is a bundle map from the
Frobenius pullback.  For vector spaces of dimension $g=m^2+1$, the stable
rank is the rank of the $g$th iterate.  Rank at least a fixed integer is
an open condition.  The nodal member just constructed consequently
supplies a nonempty open locus of rank at least $m^2-1$.
Since $N_m$ is very ample for $m\ge3$, the smooth-member locus is also
nonempty open by Bertini.  These opens meet in the irreducible projective
space $|N_m|$.  Lemma~\ref{lem:DE-A3-ss-prym} proves that the generic
smooth member has rank exactly $m^2-1$.

Now let $m=3d$ be even, so $m\ge6$.  Write $m=3+(m-3)$, with both
summands odd and at least three.  Choose smooth members
$C\in|N_3|$ and $D\in|N_{m-3}|$ in the just proved maximal-rank opens.
They may be chosen transverse: both systems are very ample, so a general
member of the second system meets a fixed smooth member of the first
transversely.  Their union is a connected nodal member of $|N_m|$.
Each component has genus minus stable rank two.  The normalization
formula therefore gives
\[
 g(C\cup D)-f_{C\cup D}=4,\qquad f_{C\cup D}=m^2-3.
\]
The same open-rank and smoothness argument gives generic smooth rank at
least $m^2-3$.  The even-degree upper bound of
Lemma~\ref{lem:DE-A3-ss-prym} makes this an equality.

Finally these are precisely the complete linear systems cut out by
ambient hypersurfaces of degree $d$.  For the plane cubic,
$H^0(\mathbf P^2,\OO(d))\to H^0(E_0,\OO(3dO))$ is onto, by
$H^1(\mathbf P^2,\OO(d-3))=0$.  The Segre coordinate ring has all its
bihomogeneous pieces of bidegree $(d,d)$ generated by degree-$d$
polynomials in the pairwise products of coordinates.  Consequently
\[
 H^0(\mathbf P^8,\OO(d))\twoheadrightarrow H^0(X,L^d).
\]
The nonempty open rank and smoothness loci pull back to nonempty open
loci in the ambient hypersurface parameter space.  Taking $m=3d$ gives
(\dagger).  The entire construction is defined over an algebraic
extension of $\mathbf F_2$ and computes the geometric generic ranks for
the displayed variety over $\mathbf F_2$ itself.
\end{proof}

\begin{corollary}[An ambient Cartier-nilpotent dimension is insufficient]
\label{cor:DE-A3-ss-product-defect}
\status{Proved}
For this abelian surface, $\Pic^0(X)$ is smooth and
$H^*_{\cris}(X/W(k))$ is torsion-free.  All its global one-forms are
Cartier-killed, so the space of global forms eventually killed by Cartier
has dimension two.  Nevertheless the generic curve has
$g(Y)-f_Y=2$ in odd degree and $g(Y)-f_Y=4$ in even degree.
Thus neither eventual relative ordinarity nor a formula equating the
generic full defect with this ambient Cartier-nilpotent dimension holds
for every sufficiently large degree.
\end{corollary}
\begin{proof}
The Picard variety is the smooth abelian surface $E_0^2$.  Crystalline
cohomology of a product of elliptic curves is the exterior algebra on
their free first crystalline cohomology, hence is torsion-free
\cite[II, (7.1.1), p.~651]{Illusie1979DRW}.
The two pullbacks of $dx$ span the global one-forms and have Cartier zero.
The defect assertion is Theorem~\ref{thm:DE-A3-ss-product-rank} and
$g(Y)=9d^2+1$.
\end{proof}

\begin{remark}[Computations and exact scope]
\label{rem:DE-A3-ss-product-computation}
\status{Proved under stated hypotheses}
The exact finite-field probe
\texttt{A3\_supersingular\_product\_probe.py} uses the basis
$\{x^a:2a\le m\}\cup\{x^ay:2a+3\le m\}$ of
$H^0(E_0,\OO(mO))$ and the Cartier identities
\[
 C(x^{2a}dx)=0,\quad C(x^{2a+1}dx)=x^a dx,\quad
 C(x^{2a}y\,dx)=x^{a+1}dx,\quad
 C(x^{2a+1}y\,dx)=x^ay\,dx.
\]
Its recorded best stable matrix ranks for $d=1,\ldots,6$ are
$8,33,80,141,224,321$; elapsed time is
$29.180787964000046$ seconds.  It retains the exact section coefficients,
matrix columns, rank sequences and seeds.  These finite probes motivated
the parity formula; the nodal deformation and Prym arguments prove it
uniformly, without inferring either smoothness or an all-degree theorem
from random samples.  The additional certificate
\texttt{A3\_supersingular\_product\_checks.py} checks the explicit
automorphisms, both symbolic quaternion norm identities, four-square
witnesses for $1\le n\le4096$, the degeneration arithmetic for
$1\le d\le10000$, and all six retained matrix witnesses.  It passed
in $0.027462210000521736$ seconds; its complete code and output are retained.
The generic-rank problem for arbitrary $X$ in A.3 remains a separate
question.  No priority claim for this explicit formula is made here.
\end{remark}


% /root/monodromy, 2026-11-19. Proved; independent root and reduction_descent audits passed.
\subsection{A supersingular product: a double Hasse component and a full image}
\label{subsec:supersingular-product-degree-one}
\status{Proved and checked. The parent and reduction/descent agent independently
audited the complete determinant, local wild-inertia and full-group proofs.
Both independently read and reran the polynomial/Pfaffian certificate and
the separate finite-field group verifier; all checks passed.}

Let
\[
 E:\ y^2+y=x^3\quad\text{over }\mathbf F_2,
 \qquad X=E\times E,\qquad
 L=\mathcal O_E(3O)\boxtimes\mathcal O_E(3O).
\]
The indicated Segre embedding identifies ambient hyperplane sections with
the complete system $P=\mathbf P(W)$, where $W=H^0(X,L)$ has dimension nine.
Write $s=(s_{ij})_{0\le i,j\le2}$, in lexicographic order, in the tensor
basis of $(1,x,y)$, and write $S=(s_{ij})$ for the corresponding
$3\times3$ matrix. Put $J_3=(\delta_{i,2-j})$ and $J=J_3\otimes J_3$.

\begin{proposition}[An explicit nonreduced relative Hasse divisor]
\label{prop:ss-product-relative-Hasse}
\status{Proved under stated hypotheses}
The relative coherent Frobenius bundle is $\mathcal E=\Omega^1_P$.
Its determinant section, up to the constant trivialization implicit in
the chosen basis, is
\begin{equation}\label{eq:ss-product-Hasse-factorization}
 h(s)=D(s)^2Q(s),\qquad
 D(s)=\det S,\qquad Q(s)=s^{[2]t}Js.
\end{equation}
The cubic $D$ is geometrically irreducible; the cubic $Q$ is
geometrically smooth and irreducible; they are distinct. Thus the
relative Hasse divisor has a component of multiplicity two and a
different component of multiplicity one. The generic relative stable
rank is eight. A smooth generic section has genus ten and $2$-rank eight.
\end{proposition}

\begin{proof}
Use $\eta=dx$ on $E$. Writing a rational function as $a^2+xb^2$
gives $C((a^2+xb^2)\,dx)=b\,dx$. The curve equation gives
\[
 C(1)=0,\quad C(x)=1,\quad C(y)=x,\quad
 C(x^2)=C(y^2)=0,\quad C(xy)=y.
\]
Consequently the root-removed Cartier multiplication matrix on the
degree-three one-factor space is
\[
 M_E(u)=\begin{pmatrix}u_1&u_0&0\\u_2&0&u_0\\0&u_2&u_1\end{pmatrix},
 \qquad
 J_3M_E(u)=\begin{pmatrix}0&u_2&u_1\\u_2&0&u_0\\u_1&u_0&0\end{pmatrix}.
\]
If $M_i$ is the coefficient of $u_i$, then the full tensor matrix is
$M(s)=\sum s_{ij}M_i\otimes M_j$. Direct differentiation of the
six-term determinant shows
\begin{equation}\label{eq:ss-Hessian}
 A(s):=JM(s)=\operatorname{Hess}_s(D),\qquad
 A(s)^t=A(s),\quad A(s)_{ii}=0,\quad M(s)s=0.
\end{equation}
The eight-dimensional Pfaffians of the alternating $9\times9$
matrix $A(s)$ form the vector
\[
       \operatorname{Pf}_8(A(s))=D(s)s.
\]
This identity may be verified by expansion of the nine Pfaffians.
The retained exact polynomial program checks all nine over
$\mathbf F_2[s_0,\ldots,s_8]$, without specialization. The standard
Pfaffian adjugate identity, or the same expansion of cofactors, gives
\[
 \operatorname{adj}M(s)=D(s)^2s\,(Js)^t,
 \qquad
 \operatorname{adj}(M(s)^t)=D(s)^2(Js)s^t.
\]

Kunneth gives $H^1(X,L^{-1})=0$, and the hypersurface exact sequence
gives the Euler sequence
\[
 0\longrightarrow\mathcal E\longrightarrow
 W^*\otimes\mathcal O_P(-1)\longrightarrow\mathcal O_P\longrightarrow0.
\]
The quotient Frobenius is zero because $X$ is supersingular.
In affine coefficient coordinates the relative Frobenius is therefore
the restriction of $M(s)^t$ from $(s^{[2]})^\perp$ to $s^\perp$.
The cofactor formula identifies the determinant of this restriction
with $D(s)^2(s^{[2]})^tJs$: indeed the exterior-product vector of a
basis of the source hyperplane is its normal $s^{[2]}$, and the
cofactor matrix carries it to
$D(s)^2s\,(Js)^ts^{[2]}$, a scalar multiple of the target normal $s$.
The projective twists agree: $\det\mathcal E=\mathcal O_P(-9)$, so
the Frobenius determinant is a section of $\mathcal O_P(9)$.
The displayed homogeneous polynomial has degree nine, proving
\eqref{eq:ss-product-Hasse-factorization} globally.

The rank-at-most-two matrix locus is the image closure of the
irreducible space of products of a $3\times2$ and a $2\times3$
matrix, and is the determinant hypersurface. Thus $D$ is
geometrically irreducible. Moreover
\[
                  \nabla Q(s)=Js^{[2]},
\]
which never vanishes at a projective point. Thus $Q$ is smooth;
two distinct positive-degree projective hypersurfaces in $\mathbf P^8$
intersect, so smoothness also implies geometric irreducibility.
At $S=\operatorname{diag}(0,1,1)$ one has $D=0,Q=1$, proving that
the two cubics are distinct. At $S=I_3$ both $D$ and $Q$ are one.
Thus the relative Frobenius is generically invertible.
The smooth open is nonempty by the very ample embedding, so it meets
this invertibility open. Adjunction gives $g=1+L^2/2=10$.
The ambient $H^1(X,\mathcal O_X)$ has nilpotent Frobenius and the
relative stable rank is eight, so the smooth generic curve has
$2$-rank eight.
\end{proof}

\begin{lemma}[A first-derivative criterion for wild transvection inertia]
\label{lem:ss-simple-pole-transvection}
\status{Proved under stated hypotheses}
Let $R=k[[t]]$ with $k$ algebraically closed of characteristic two,
and let a rank-$r$ vector bundle over $R$ have regular semilinear
Frobenius $F$. Suppose that its special fibre has stable rank $r-1$,
its linearized kernel has dimension one, and its determinant has
valuation one. Let $u$ generate the semilinear kernel in the special
fibre, and let $v$ be the first derivative of $F(\widetilde u)$ at
$t=0$, for a regular lift $\widetilde u$. If $v\notin ku$, the
generic Lang local system has a nontrivial transvection in its inertia.
\end{lemma}

\begin{proof}
The finite part of the Lang scheme supplies $r-1$ regular fixed
sections, whose special values span the stable subspace. This is the
finite-part lemma already proved using Stacks~\cite{Stacks0BTN,Stacks09ZS}.
They form a Frobenius-fixed direct summand. Complete them by a lift of
$u$. In this frame
\[
 F=\begin{pmatrix}I_{r-1}&b(t)\\0&a(t)\end{pmatrix},
 \qquad a(t)=a_1t+O(t^2),\qquad b(t)=b_1t+O(t^2),\qquad a_1\ne0.
\]
The constant column is zero because $u$ is in the special kernel.
Changing the lift changes its Frobenius image only in order at least
two, so its first derivative is well defined in the special fibre.
The hypothesis $v\notin ku$ is precisely $b_1\ne0$.

The nonzero fixed vector in the one-dimensional generic quotient has
coordinate $z=a(t)^{-1}$. Lifting it requires solving
\[
                  y_i^2+y_i=b_i(t)/a(t)^2.
\]
At least one right side has a nonzero simple pole. Such an
Artin--Schreier equation is ramified: if $y$ lay in an unramified
extension, its negative valuation would give an even pole order
for $y^2+y$, a contradiction; a nonnegative valuation gives no pole.
Over the strictly henselian field all inertia fixes the regular
$(r-1)$-dimensional Lang subspace and its one-dimensional quotient.
A nonidentity element therefore differs from the identity by a rank-one
square-zero map, which is a transvection.
\end{proof}

\begin{theorem}[Full degree-one image on a supersingular product]
\label{thm:ss-product-degree-one-full-group}
\status{Proved under stated hypotheses}
For $X,L$ above, over every perfect field of characteristic two,
the degree-one relative Hasse--Witt local system on its maximal-rank
smooth open has
\[
                   G_{\rm geom}=G_{\rm arith}=\GL_8(\mathbf F_2).
\]
In particular its geometric representation is irreducible, despite
the double component of its coherent Hasse divisor.
\end{theorem}

\begin{proof}
We give two small exact certificates over
$k_8=\mathbf F_2[\alpha]/(\alpha^3+\alpha+1)$.
For readability an integer $c_0+2c_1+4c_2$ denotes the field element
$c_0+c_1\alpha+c_2\alpha^2$; all coefficient lists below use this
encoding and the fixed lexicographic tensor basis.

First take
\[
                s=(6,7,0,5,4,1,6,0,4).
\]
The relative Frobenius is invertible and the characteristic polynomial
of its third iterate, up to replacing it by its reciprocal according
to the chosen Frobenius convention, is
\begin{equation}\label{eq:ss-degree-one-Frobenius-polynomial}
                    T^8+T^7+T^5+T+1.
\end{equation}
The retained verifier checks its irreducibility over $\mathbf F_2$
by the Frobenius-power and gcd criterion and checks that it is not
reciprocal. Thus the arithmetic image contains the element required
in Lemma~\ref{lem:product-frobenius-group-test}.
The verifier uses an independent restriction-of-scalars calculation:
if $B=A A^{[2]}A^{[4]}$ is the full third Frobenius iterate, then
$B^{[2]}=A^{[2]}A^{[4]}A$ has the same characteristic polynomial,
because cyclic products of square matrices do. Thus that polynomial
has coefficients in $\mathbf F_2$. Restriction from $\mathbf F_8$ to
$\mathbf F_2$ cubes it. The directly computed $27\times27$ integer
characteristic polynomial, reduced modulo two, is
$(T(T^8+T^7+T^5+T+1))^3$; unique factorization, with monic factors,
therefore proves the asserted nine-dimensional polynomial and its
eight-dimensional relative factor.

For a geometric transvection, take
\[
 s_0=(0,6,4,5,5,5,3,3,4),\qquad\delta=(1,0,\ldots,0),
 \qquad s(t)=s_0+t\delta.
\]
Here $D(s_0)\ne0$, $Q(s_0)=0$, and $dQ_{s_0}(\delta)=6\ne0$.
The full linearized matrix $M(s_0)^t$ has rank eight; its semilinear
stable rank is seven. A generator of its semilinear kernel is
\[
 u=(2,7,7,3,3,3,2,4,0)=(Js_0)^{[1/2]}.
\]
The verifier checks $s_0^tu=0$, so this is also the relative kernel.
The derivative on this kernel along the displayed affine line is
\[
 v=M(\delta)^tu^{[2]}=(0,0,0,0,4,3,0,5,5).
\]
These two vectors are independent. The determinant factorization
shows that the relative determinant has valuation one along this
line at $t=0$.

For completeness, this ambient calculation is also the derivative
in a regular frame of the relative bundle. A lift of $u$ in
$s(t)^\perp$ has square $u^{[2]}+O(t^2)$; hence its image derivative
is exactly $M(\delta)^tu^{[2]}$. A change of output frame has no
first-order effect because its image at $t=0$ is zero. Passing
from the affine coefficient cone to a projective chart only subtracts
a radial tangent vector; that adds a multiple of
$M(s_0)^tu^{[2]}=0$ and changes neither derivative test.
Lemma~\ref{lem:ss-simple-pole-transvection} therefore supplies a
nontrivial geometric transvection.

Let $H$ be its normal closure in the arithmetic image $G$.
Then $H$ is nontrivial, generated by transvections, normal in $G$,
and lies in the geometric image. Applying
Lemma~\ref{lem:product-frobenius-group-test} with $N=8$ proves
$H=G=\GL_8(\mathbf F_2)$. That lemma explicitly handles Clifford
blocks using the primitive prime $17$ and excludes all form-preserving
alternatives by nonreciprocity; its only classification input is
\cite[Theorem~3.6, printed p.~22]{KantorSeress2001}.

All calculations were made on the invertibility open of the coherent
bundle over the full projective parameter space. Restriction to the
nonempty smooth open does not change the image: a connected finite
etale cover of a normal integral scheme remains connected on a
dense open. Thus the use of coefficient points or a local line does
not require a separate smoothness assertion for those particular
sections. Finally geometric connectedness of the finite frame cover
is preserved by algebraically closed base extension. The same
geometric full group, and hence the same arithmetic full group, holds
over every perfect base field of characteristic two.
\end{proof}

\begin{remark}[Retained exact calculations]\label{auto:A_monodromy:3502a4b6e144}
\status{Proved under stated hypotheses}

The polynomial certificate is
\path{agents/monodromy/code/supersingular_product_degree_one.py}, with
output \path{agents/monodromy/outputs/supersingular_product_degree_one.json}.
It verifies all nine Pfaffian identities and all $511$ projective
$\mathbf F_2$ coefficient vectors. The extension-field discovery
calculation is
\path{agents/monodromy/code/supersingular_product_frobenius_probe.py};
its output retains the field polynomials, coefficients, characteristic
polynomials and local derivative data. A separate certificate verifier
\path{agents/monodromy/code/verify_supersingular_degree_one_group.py}
checks only the two stated $\mathbf F_8$ points and does not depend
on the discovery search. Its output is
\path{agents/monodromy/outputs/verify_supersingular_degree_one_group.json}.
\end{remark}


% /root/monodromy, 2026-11-19. Proved; independent root, reduction_descent and obstructions_crystals proof audits passed.
\subsection{All invariant covectors on the supersingular product}
\label{subsec:ss-exact-covectors}
\status{Proved and checked. The parent and two other agents independently
audited the complete proof; the local Cartier certificate was independently
rerun by the parent. The persistent coherent nilpotent part is handled by a
bounded Frobenius-root extension and a homogeneous-coefficient argument.}

\begin{theorem}\label{thm:ss-product-exact-covectors}
\status{Proved under stated hypotheses}
Let $E_0:y^2+y=x^3$ over a perfect field of characteristic two, let
$X=E_0\times E_0$ and
$L=\mathcal O_{E_0}(3O)\boxtimes\mathcal O_{E_0}(3O)$.
For every $d\ge1$, the entire geometric and arithmetic invariant
covector spaces of the relative Hasse--Witt local system have dimension
\[
 \dim(V_d^\vee)^{G_{\rm geom}}
 =\dim(V_d^\vee)^{G_{\rm arith}}
 =\begin{cases}0,&d\text{ odd},\\1,&d\text{ even}.
 \end{cases}
\]
In even degree the unique nonzero covector is represented under
adjunction by $D(s)\omega$, where
$\omega=dx_1\wedge dx_2$, $D=\delta_1\delta_2$, and $\delta_i$
is the invariant derivation dual to $dx_i$, lifted to $L^d$ by its
canonical Frobenius-pullback connection. Thus every even-degree image
is contained in the pointwise stabilizer of one nonzero covector.
This theorem determines the entire invariant quotient, not the full group.
\end{theorem}

\begin{proof}
First work over an algebraic closure. Put $m=3d$,
$N=\mathcal O(mO)\boxtimes\mathcal O(mO)$,
$W=H^0(X,N)$ and $b=m^2$. On $P=\mathbf P(W)$ the full coherent
bundle is $\mathcal B=W^*\otimes\mathcal O_P(-1)$, with semilinear
Frobenius matrix $M(s)^t$. Here $M(s)$ is linear in the coefficient
vector $s$ and is the root-removed Cartier multiplication matrix.
Its relative subbundle is $\mathcal E=\Omega_P^1$; the quotient
$\mathcal B/\mathcal E=H^2(X,\mathcal O_X)\otimes\mathcal O_P$
has zero Frobenius. Thus the unit Lang local systems of the full
and relative bundles coincide.

\par\smallskip\noindent\textbf{The generic linearized image is already the stable image.}
The identity $C(\omega)=0$ gives $M(s)s=0$ in every degree.
For odd $m$, Theorem~\ref{thm:DE-A3-ss-product-rank} gives stable
rank $b-1$, so the generic linearized rank is also $b-1$.
For even $m$, $N$ is the relative Frobenius pullback of $N_{m/2}$.
The invariant derivations act on its sections, since they kill the
square transition functions of pullback frames. They commute and
satisfy $\delta_1^2=\delta_2^2=0$. A local $2$-basis calculation gives
\[
 C(s\delta_1(s)\omega)=C(s\delta_2(s)\omega)=0.
\]
Thus $s,\delta_1(s),\delta_2(s)$ lie in $\ker M(s)$.
They are generically independent: for the section $s=x_1x_2$
they are $x_1x_2,x_2,x_1$. Hence the generic linearized rank is
at most $b-3$. The exact stable-rank theorem gives $b-3$, so
equality holds. In both parities the image of the first linearized
Frobenius at the generic point is therefore its unit subspace.

\par\smallskip\noindent\textbf{A purely inseparable lift of an arbitrary invariant covector.}
Let $K=k(P)$ and let $I\subset\mathcal B_K$ be that unit image.
The restriction $F^*I\to I$ is invertible. The finite Lang-frame
equivalence identifies a geometric invariant covector with a
$K$-linear map $f:I\to K$ satisfying
$f(Fx)=f(x)^2$ on $I$. Define a functional on the full generic
bundle, with coefficients in $K^{1/2}$, by
\begin{equation}\label{eq:ss-covector-root-extension}
                  \lambda(x)^2=f(Fx).
\end{equation}
The right side is linear in the squares of the coordinates of $x$,
so this defines a linear functional over $K^{1/2}$. On $I$ it
equals $f$ by uniqueness of square roots. Since $Fx$ lies in $I$,
it satisfies $\lambda(Fx)=\lambda(x)^2$ on the whole bundle.
This extension is unique: any equivariant functional zero on $I$
has square zero by \eqref{eq:ss-covector-root-extension}.

In coordinates, $M(s)\lambda=\lambda^{[2]}$. The squared coefficient
vector $w=\lambda^{[2]}$ is $K$-rational and is a rational section of
$F_P^*\mathcal B^*=W\otimes\mathcal O_P(2)$. It satisfies
\begin{equation}\label{eq:ss-square-covector-equation}
                       M(s)^{[2]}w=w^{[2]}.
\end{equation}
It has no divisorial poles. In a regular local frame, if the largest
pole order among its coordinates were $a>0$, the right side would
have a coordinate of pole order $2a$, whereas the left side has
pole order at most $a$. Normality and local freeness extend the
section also across codimension two. Thus $w$ is a vector of
homogeneous quadratic polynomials in the coefficients $s_i$.

\par\smallskip\noindent\textbf{Cross terms cannot occur.}
Write $M(s)=\sum s_i M_i$. The common kernel of the $M_i$ is zero.
Indeed, relative finite-Frobenius duality, with the invariant top
form trivializing the canonical bundle, gives a perfect symmetric
Cartier pairing
\begin{equation}\label{eq:ss-perfect-Cartier-pairing}
 F_{X/k*}N\otimes F_{X/k*}N\longrightarrow N^{(2)}.
\end{equation}
The finite-duality and coefficient-trace formulas used here are
\cite[Lemma~1.3.6, Proposition~1.3.7, Theorem~1.3.8, pp.~26--28]
{BrionKumar2005}, already used in the general adjunction proof.
Its pairing on global sections is the root-removed bilinear map
$\beta(s,t)=M(s)t$. The bundle $F_{X/k*}N$ is globally generated.
On either elliptic factor, evaluation on a Frobenius fibre is
evaluation on a length-two subscheme, and the obstruction group is
$H^1(\mathcal O(mO)(-Z))=0$ because $m-2>0$.
Tensoring the two evaluation surjections proves the assertion on
the product. A global section pairing to zero with every global
section therefore vanishes in every fibre, by perfectness, and
is zero. This proves the common-kernel assertion, also after
squaring the matrices.

For $j<k$, let $c_{jk}$ be the coefficient of $s_js_k$ in $w$.
The right side of \eqref{eq:ss-square-covector-equation} has only
even exponents. The monomials with precisely $s_j,s_k$ occurring
to odd powers on its left side give
\[
              \sum_i s_i^2M_i^{[2]}c_{jk}=0.
\]
Thus every $M_i^{[2]}c_{jk}=0$, and $c_{jk}=0$.
All cross terms vanish. Since $k$ is perfect,
\[
                         w=D(s)^{[2]}
\]
for a $k$-linear endomorphism $D$ of $W$. Taking square roots of
\eqref{eq:ss-square-covector-equation} shows
\begin{equation}\label{eq:ss-exact-linear-operator}
                 \beta(s,D(s))=D(s)^{[2]}.
\end{equation}
The original functional is exactly $D(s)$, not merely its
purely inseparable extension. Conversely every $D$ satisfying
\eqref{eq:ss-exact-linear-operator} is a regular equivariant
covector. No quotient by scalar identities is taken here:
the full-to-relative quotient has Frobenius zero, so it has no
nonzero equivariant covectors.

\par\smallskip\noindent\textbf{Descent to a small endomorphism algebra.}
Polarizing \eqref{eq:ss-exact-linear-operator} gives
\[
                   \beta(Ds,t)=\beta(s,Dt).
\]
If a global section $s$ has zero value in a fibre of $F_*N$,
the right side pairs that zero value with $Dt$, so the left
side vanishes for every $t$. Global generation and the perfect
pairing force $Ds$ to have zero value there too.
Hence $D$ preserves every evaluation kernel and descends to an
$\mathcal O_{X^{(2)}}$-linear endomorphism of $F_*N$.
Formally the kernel of the global evaluation map is locally free;
the induced map from it to $F_*N$ vanishes on every closed fibre
and therefore vanishes identically. The descended endomorphism
acts on global sections as the original $D$.

For a line bundle $A$ of odd degree $m$ on an elliptic curve,
$F_*A$ is stable of rank two and degree $m$.
A line subbundle $T$ of degree at least $m/2$ would, by adjunction,
give a nonzero map $F^*T\to A$, hence $2\deg T\le m$,
contradicting the integral degree assumption. A noninvertible
nonzero endomorphism would have rank-one kernel and image; stability
puts both their degrees strictly below $m/2$, contradicting
additivity of degree. Every nonzero endomorphism is therefore
invertible, and over the algebraically closed field its endomorphism
algebra is $k$. Kunneth consequently gives
$\operatorname{End}(F_*N)=k$ when $m$ is odd.
Equation~\eqref{eq:ss-exact-linear-operator}, together with
$\beta(s,s)=0$, forces the scalar $D$ to be zero.

For even $m$, projection formula gives
$F_*\mathcal O(mO)=\mathcal O_{E_0^{(2)}}((m/2)O)\otimes U$,
where $U=F_*\mathcal O_{E_0}$. Its rank-one quotient in
$0\to\mathcal O\to U\to B\to0$ has degree zero. The cohomology
sequence and the zero Frobenius on $H^1(E_0,\mathcal O)$ give
$h^0(B)=1$, so $B\simeq\mathcal O$.
The extension is nonsplit because $h^0(U)=1$.
Every endomorphism preserves its unique global section and induces
scalars $a,b$ on the subline and quotient. Compatibility with its
nonzero extension class gives $a=b$. Subtracting $a$ times the
identity leaves a multiple of the quotient-followed-by-inclusion
map, which is nonzero and square zero. The invariant derivation
$\delta$ is precisely a nonzero multiple of this map. Thus
\[
 \operatorname{End}(U)=k[\delta]/(\delta^2),\qquad
 \operatorname{End}(F_*N)
       =k[\delta_1,\delta_2]/(\delta_1^2,\delta_2^2).
\]

It remains to impose the Cartier equation on
$D=a+b\delta_1+c\delta_2+e\delta_1\delta_2$.
In a pullback frame, write
$s=u^2+x_1v^2+x_2w^2+x_1x_2z^2$.
Taking the coefficient of $x_1x_2$ gives
\[
 C(s^2\omega)=C(s\delta_1s\,\omega)
 =C(s\delta_2s\,\omega)=0,
 \qquad C(s\delta_1\delta_2s\,\omega)=\delta_1\delta_2s\,\omega.
\]
Consequently $C(sD(s)\omega)=D(s)\omega$ is equivalent to
$a=b=c=0$ and $e^{1/2}=e$, or $e\in\mathbf F_2$.
The four operators are independent on global sections because
the evaluation map is surjective, so no further solutions arise.
The nonzero solution $D=\delta_1\delta_2$ exists and is nonzero:
$D(x_1x_2)=1$. This proves the geometric dimensions in both parities.

Finally $D$, the invariant differential, the line-bundle square root
and its Cartier connection are defined over $\mathbf F_2$.
The unique nonzero even-degree covector is therefore arithmetic
invariant over the stated perfect base field. In odd degree there
are no geometric invariant covectors and hence no arithmetic ones.
This proves the asserted arithmetic dimensions as well.
\end{proof}

\begin{remark}[Exact local checks and scope]\label{auto:A_monodromy:9cc9a63fa2c6}
\status{Proved under stated hypotheses}

The local four-operator identities are reproduced symbolically in
\path{agents/monodromy/code/supersingular_cartier_operator_checks.py},
with output
\path{agents/monodromy/outputs/supersingular_cartier_operator_checks.json}.
The same retained program performs a bounded check of constant
alternating symmetrizers for the one-factor Cartier pencils: the
space has dimension one in pole degree three and zero in the tested
odd pole degrees five, seven, nine, eleven and thirteen. This check
is not used in the proof and asserts nothing for untested degrees;
it records why the special degree-one Hessian calculation has not
been extrapolated to higher degrees.
\end{remark}


% /root/monodromy, 2026-11-19. Proved; independent root and reduction_descent audits passed.
\subsection{A natural eight-dimensional block in every supersingular degree}
\label{subsec:ss-uniform-GL8-block}
\status{Proved and checked. The parent and reduction/descent agent
independently audited both boundary constructions, the full node-cover
comparison and the extraction of an actual natural block. This determines
a uniform subgroup. The full odd-degree images are determined later in
Theorem~\ref{thm:ss-all-odd-full-groups}, and the complete all-degree
classification is Theorem~\ref{thm:ss-product-all-full-groups}.}

\begin{theorem}\label{thm:ss-uniform-GL8-block}
\status{Proved under stated hypotheses}
For $X=E_0\times E_0$, $E_0:y^2+y=x^3$, with
$L=\mathcal O(3O)\boxtimes\mathcal O(3O)$ in characteristic two,
the degree-$d$ geometric Hasse--Witt image contains, for every $d\ge1$,
a subgroup acting as the natural $\GL_8(\mathbf F_2)$ on an
eight-dimensional subspace and fixing a complementary subspace pointwise.
Consequently the whole representation admits no invariant nondegenerate
alternating bilinear form. In even degree the same conclusion holds on
the common kernel of the invariant covectors of
Theorem~\ref{thm:ss-product-exact-covectors}.
No irreducibility or full-image assertion in higher degree is made here.
\end{theorem}

\begin{proof}
We work over an algebraic closure. The degree-one case is
Theorem~\ref{thm:ss-product-degree-one-full-group}.
We use the exact ranks of Theorem~\ref{thm:DE-A3-ss-product-rank}.
The maximum stable-rank open of the coherent family includes singular
divisors; its unit Lang system restricts to the local system in the
question on the dense smooth open, with the same image.
Indeed an iterated Frobenius image of constant maximum rank is a
unit subbundle, and restriction of a finite etale cover on a normal
integral scheme to a dense open preserves its connected components.

\par\smallskip\noindent\textbf{Odd degree: an odd self-isogeny and nodal images.}
Suppose $d\ge3$ is odd. Lemma~\ref{lem:DE-A3-ss-norm} gives an
endomorphism $b$ of $E_0$ with $b^\dagger b=bb^\dagger=[d-1]$.
Put
\[
 \phi=\begin{pmatrix}1&b\\-b^\dagger&1\end{pmatrix}:X\longrightarrow X.
\]
Then $\phi^\dagger\phi=\phi\phi^\dagger=[d]I$, and $\phi$ is
an etale isogeny of degree $d^2$ with constant odd kernel $G$.
For a generic smooth degree-one member $C\in|L|$, the image
$Z=\phi(C)$ is integral and nodal, and $C$ is its normalization.
Here is why the nodal-image argument applies also to this nonproduct
isogeny and to supersingular elliptic factors.

Only translations by points of the constant kernel enter that argument.
A generic $C$ is not invariant under any nonzero translation: if a
translation preserves $L$, its projective action is nontrivial because
$L$ is very ample; otherwise it preserves no member of $|L|$.
Thus $C\to Z$ is generically one-to-one. For a nonzero kernel point
$g$, double identifications are parametrized by
$C\cap(C-g)$, of degree $L^2=18$.

The tensor first-jet calculation of
Lemma~\ref{lem:product-small-degree-jets} uses only the plane cubic
embeddings and the fact that tangents through a fixed point are finite.
For the present cubic $Y^2Z+YZ^2+X^3$, the partial derivatives are
$X^2,Z^2,Y^2$, which are linearly independent, and no nonzero quadratic
vanishes on the cubic. Thus that finiteness condition holds.
Two first jets have rank at least five; rank six fails only on the
two coordinate-fibre loci and the finite crossed-tangency locus
over either fixed first point. For pairs $P,P+g$ with both coordinates
of $g$ nonzero, the exceptional points satisfy equations
$3P_i=\text{constant}$ and are finite. Away from them, the two value
conditions and proportionality of the two nonzero gradients impose
three conditions with $P$ varying in dimension two.
At the finitely many exceptions, the two independent values already
give a proper incidence image.

If one coordinate of $g$ is zero, the joint first-jet rank is five.
After imposing the two values, its sole remaining relation involves
the two derivatives along the translated factor and is nonzero;
the two transverse derivatives remain independent. The determinant
of the two gradients is not identically zero on this three-dimensional
allowed derivative space. Tangency is again a third condition.
Thus every $C\cap(C-g)$ is transverse for generic $C$.

Three-point identifications are excluded by the same Segre value
calculation used in Lemma~\ref{lem:ordinary-product-isogeny-nodal-images}.
Three distinct rank-one Segre tensors are collinear only if one factor
is constant and the other three points are collinear. For
$P,P+g,P+h$, the exceptional locus has dimension at most one, since
it requires translations in the same factor and
$3P_i+g_i+h_i=0$. The value ranks are three off this locus and two
on it; the incidence dimensions in $|L|=\mathbf P^8$ are at most
$8-3+2=7$ and $8-2+1=7$. Etaleness of $\phi$ therefore gives only
ordinary nodes. Each $g$ contributes an etale intersection cover of
degree eighteen, with the constant identification of $g$ and $-g$.
The node permutation image is consequently a subquotient of a
product of $S_{18}$'s.

The image has exactly the required line bundle. Write
$\Theta=(\{O\}\times E_0)+(E_0\times\{O\})$, so $L=\mathcal O(3\Theta)$.
The Rosati identity gives $\phi^*[\Theta]=d[\Theta]$.
For the norm line bundle, pulling back its class gives
\[
 \phi^*[\operatorname{Nm}_\phi L]
   =\sum_{g\in G}[t_g^*L]=d^2[L]=\phi^*[L^d].
\]
Pullback by an isogeny is injective on the torsion-free Neron--Severi
group after tensoring with $\mathbf Q$, so the two line bundles
have the same numerical class and differ by a member of $\Pic^0(X)$.
Both are symmetric, and $\Pic^0(X)[2](k)=0$; hence the difference
is trivial. Thus $Z\in|L^d|$ exactly.

Its arithmetic genus is $9d^2+1$, while the normalization has
genus ten and stable rank eight. Therefore it has $9(d^2-1)$
nodes and stable rank $8+9(d^2-1)=9d^2-1$, the full generic rank.
On this boundary family its Lang representation has an exact sequence
\[
                  0\longrightarrow T\longrightarrow V_d
                    \longrightarrow Q\longrightarrow0,
 \qquad Q\simeq V_1,\quad \dim Q=8,
\]
where $T$ is the graph permutation contribution. There is no ambient
stable contribution because $X$ is supersingular.
The quotient image is full $\GL_8(\mathbf F_2)$, whereas the whole
image on $T$ has prime factors at most seventeen.

\par\smallskip\noindent\textbf{Even degree: a fixed odd member and a symmetric node cover.}
Let $d\ge2$ be even. Choose a fixed smooth maximal-rank member
$D\in|L^{d-1}|$, which exists by the odd-degree rank theorem, and
vary $C\in|L|$ through its smooth maximal-rank open, transverse to $D$.
Then $Y=C\cup D$ is nodal and
\[
 f_Y=8+\bigl(9(d-1)^2-1\bigr)+18(d-1)-1=9d^2-3.
\]
This again equals the full generic rank. The boundary Lang system has
a filtration
\[
 0\subset T\subset T_0\subset V_d,
 \qquad T_0/T\simeq H^1_{\et}(D,\mathbf F_2)\text{ constant},
 \qquad V_d/T_0=Q\simeq V_1.
\]
The graph space is $T=\mathbf F_2^N/\mathbf F_2\mathbf 1$,
where $N=18(d-1)$ is the number of intersection points.
Its permutation image is the full $S_N$.

To justify that last assertion, the one-point incidence over $D$ and
the ordered distinct-pair incidence over $D\times D\setminus\Delta$
are irreducible projective bundles: the very ample embedding separates
one and two values. Thus the node action is doubly transitive.
There exists a section with exactly one double zero on $D$ and all
other zeros simple. Choose a general point $P\in D$ whose two
elliptic coordinates are not flex points. At $P$, $L$ separates
second jets, so the quadratic coefficient along the tangent of $D$
can be made nonzero after the value and first tangent derivative
are killed. For $Q\in D\setminus\{P\}$, the two curve jets at
$P,Q$ have combined rank four except for finitely many $Q$, and
rank at least three at those exceptions, by the preceding tensor
first-jet bounds. The exceptional points are finite because $D$ is
integral and ample, hence is not a coordinate fibre. Within the
codimension-two system imposed at $P$, another double zero imposes
two conditions with $Q$ varying in dimension one, or at least one
condition at finitely many exceptions. Its incidence image is proper.
The exact double zero is thus attainable without any second collision.
The point-incidence source is smooth and the base is smooth.
These are precisely the two conditions of
\cite[Proposition~3.1, arXiv version~3, pp.~7--8]{Entin2021}, already
used in Lemma~\ref{lem:all-X-independent-nodes}; they yield a
transposition also in characteristic two and therefore the full $S_N$.
Shrinking to the smooth maximal-rank transverse $C$ open does not
change this finite permutation image.

The action of $S_N$ on $\mathbf F_2^N/\mathbf F_2\mathbf1$ is
faithful for $N\ge3$: the difference of two distinct coordinate
vectors has support two and cannot be the all-ones vector.
The whole image on $T_0$ is an extension of $S_N$ by a two-group,
since $T_0/T$ is constant. Its only nonabelian composition factor
is $A_N$. This is not $\GL_8(\mathbf F_2)$: $11$ divides $|A_N|$
for $N\ge18$, whereas it does not divide
$|\GL_8(\mathbf F_2)|=2^{28}\prod_{i=1}^8(2^i-1)$.

\par\smallskip\noindent\textbf{Extraction of an actual block.}
In either parity let $T_0$ denote the kernel of the quotient to $Q$
(in odd degree take $T_0=T$). The simple group $\GL_8(\mathbf F_2)$
has no quotient in common with the entire $T_0$ image: in odd degree
use the prime $31$, and in even degree use the just determined
composition factors. Goursat's lemma shows that the subgroup acting
trivially on all of $T_0$ still surjects onto $\GL(Q)$.
Its translation kernel in $\operatorname{Hom}(Q,T_0)$ is
$Q^*\otimes T'$ for a subspace $T'\subset T_0$, since this module
is a direct sum of copies of the absolutely simple natural dual.
Modulo that kernel, the remaining block extension is a cocycle with
values in $\operatorname{Hom}(Q,T_0/T')$.
The proved elementary vanishing
$H^1(\GL_8(\mathbf F_2),Q^*)=0$
(Lemma~\ref{lem:char2-group-inputs}) makes it a coboundary.
After conjugation it is zero. The entire translation kernel is
present, so its elements then remove the remaining $T'$-valued
columns. Hence the subgroup contains every block diagonal matrix
$\operatorname{diag}(I_{T_0},A)$ with $A\in\GL(Q)$.
This proves the actual natural-block assertion.

Finally the natural eight-dimensional $\GL_8(\mathbf F_2)$ module
has no invariant covector and no nonzero invariant bilinear form.
For the latter, the stabilizer of any nonzero vector has no fixed
covector: its arbitrary upper row first kills the covector's first
coordinate, and its $\GL_7$ block kills the rest. Thus there is
no nonzero equivariant map $Q\to Q^*$.
Any invariant bilinear form on $V_d=Q\oplus T_0$ therefore vanishes
on $Q\times Q$ and on $Q\times T_0$, so $Q$ lies in its radical.
In even degree every invariant covector vanishes on $Q$; its common
kernel is $Q\oplus(T_0\cap K)$, and the same argument applies there.
This proves the stated form exclusions.
\end{proof}


\subsection{A mixed product with a uniform exact rank formula}
\label{sec:DE-A3-mixed-generic}

The following statement uses one fixed perfect field and one fixed embedded
surface in every degree.  Specialization is used to prove the formula;
the ordinary elliptic factor is not chosen anew as the degree changes.

\begin{theorem}[The geometric generic ordinary elliptic factor]
\label{thm:DE-A3-mixed-generic}
\status{Proved; complete independent root and monodromy proof audits passed.}
Let $k=\mathbf F_2(t)^{\mathrm{perf}}$ be the perfect closure of the
rational function field, and define
\[
 E_t:\ y^2+txy+y=x^3,\qquad
 E_0:\ y^2+y=x^3,\qquad
 X=E_t\times_k E_0,
 \qquad L=\OO_{E_t}(3O)\boxtimes\OO_{E_0}(3O).
\]
Use the cubic embeddings and the Segre map to embed $X$ in
$\mathbf P^8_k$.  The first elliptic factor is ordinary and the second
is supersingular.  For every $d\ge1$, the geometric generic relative
stable rank is
\[
 \boxed{\quad
 \dim E(Y,\OO_Y)_{\ss,\mathrm{gen}}
  =\begin{cases}
       9d^2-1,& d\ \text{odd},\\
       9d^2-2,& d\ \text{even}.
    \end{cases}\quad}                                      \tag{\ddagger}
\]
The relative coherent dimension is always $9d^2-1$.
Thus its generic Frobenius is bijective exactly in odd degree.
The full generic curve defect $g(Y)-f_Y$ is one in odd degree and
two in even degree.

The same formulas hold after replacing $E_t$ by any elliptic curve
over a perfect characteristic-two field whose $j$-invariant is
transcendental over $\mathbf F_2$, and taking geometric generic ranks.
This last assertion does not include arbitrary algebraic
$j$-invariants.
\end{theorem}
\begin{proof}
Consider the smooth elliptic scheme
\[
 \mathcal E\longrightarrow
 S=\Spec\mathbf F_2[t,(t^3+1)^{-1}]
\]
given by the same Weierstrass equation.  The integral formulas with
$a_1=t,a_3=1$ and the other coefficients zero give
$c_4=t^4-24t$ and $\Delta=t^3-27$; in characteristic two these are
$t^4$ and $t^3+1$.  The Hasse coefficient is $a_1=t$.
Thus the generic fibre is ordinary, while the fibre at $t=0$ is
exactly $E_0$.  In particular this is a specialization within smooth
elliptic schemes.

Fix an odd $d\ge1$ and put $m=3d$.  On
$\mathcal X=\mathcal E\times_S(E_0\times S)$ use
$\mathcal N_m=\OO_{\mathcal E}(mO)\boxtimes\OO_{E_0}(mO)$.
Its direct image has rank $m^2$ and commutes with base change,
because both elliptic factor bundles have positive degree and no
higher cohomology.  Let $\mathcal P_m\to S$ be its projective space
of nonzero sections up to scale, and let $\mathcal D_m$ be the
universal divisor.  The relative coherent quotient is a vector
bundle $\mathcal Q_m$ of rank $m^2-1$ on all of $\mathcal P_m$.
Indeed fibrewise K\"unneth gives
\[
 H^1(\mathcal X_s,\mathcal N_{m,s}^{-1})=0,\qquad
 h^2(\mathcal X_s,\mathcal N_{m,s}^{-1})=m^2,
 \qquad h^2(\mathcal X_s,\OO)=1.
\]
Multiplication by a nonzero section maps the middle cohomology
surjectively to the last one, since its Serre-dual map sends $1$
to that nonzero section.  Its kernel is exactly the relative
quotient, including for singular divisors.  This also proves
the required base-change and local-freeness assertions.

Absolute Frobenius gives an $\OO_{\mathcal P_m}$-linear map
$F_{\mathcal P_m}^*\mathcal Q_m\to\mathcal Q_m$.
Its invertibility locus is open.  On the fibre $t=0$, the
supersingular-product theorem
\ref{thm:DE-A3-ss-product-rank} supplies a smooth section where it
is invertible, since $d$ is odd.  The determinant is therefore
a nonzero section of a line bundle on the integral scheme
$\mathcal P_m$.  It cannot vanish on its entire generic fibre:
otherwise it would vanish at the generic point of
$\mathcal P_m$, hence everywhere.  The geometric generic curve
over $\mathbf F_2(t)$ consequently has relative rank $m^2-1$.
Passing to its perfect closure preserves this rank, as does any
further algebraically closed extension.

This argument applies to every odd $d$ at the same generic point
of $S$.  It makes no assertion that an infinite intersection of
nonempty open subsets of $S$ contains a closed point.  The
elliptic curve over the fixed field $k$ is therefore the same
$E_t$ for all of these conclusions.  The smooth-member open is
nonempty in each degree because both factor line bundles have
degree at least three.

Now work over $\bar k$ and let $Y\in|L^d|$ be smooth, with $m=3d$.
The Picard embedding and connected-Prym proof of
Lemma~\ref{lem:DE-A3-ss-prym} applies word for word to
$E_t\times E_0$: it only used negative-degree K\"unneth vanishing
for the embedding, and the diagonal product polarization for
$qj=[m]$.  Thus
\[
 0\longrightarrow P\longrightarrow J_Y\xrightarrow{q}X
 \longrightarrow0
\]
is exact with $P$ an abelian variety.  Here $g(Y)=m^2+1$,
$f_X=1$, $f_Y=f_P+1$, and $\dim P=g(Y)-2$.
If $m$ is even, $j$ embeds $X[2]$ into $P[2]$.  The identity
component of $\ker V_X$ has order $2$, since $\dim X=2$ and
$f_X=1$.  Applying Frobenius twist and functorial finite-group
Verschiebung to that immersion embeds this connected group
into $\ker V_P$.  Therefore
\[
 \dim P-f_P\ge1,\qquad f_Y\le g(Y)-2.
\]
Exactness of stable parts, now subtracting the one-dimensional
ambient stable part, gives the upper bound
$\dim E(Y,\OO_Y)_{\ss}\le m^2-2$.

For even $d\ge2$, choose generic smooth members
$C\in|L|$ and $D\in|L^{d-1}|$ in the odd-degree maximal-rank
opens just proved.  Each has full genus-minus-$2$-rank defect
one, since its ambient stable rank is one.  They can be chosen
transverse, by very ampleness and Bertini.  The nodal union
$C\cup D\in|L^d|$ has full defect two: the normalization
sequence adds the graph rank bijectively for Frobenius and
adds the two component defects.  Constant coherent dimensions
and the open-rank argument in
Theorem~\ref{thm:DE-A3-ss-product-rank} deform it to a smooth
member of relative rank at least $m^2-2$.  The preceding upper
bound makes this equality.  This proves (\ddagger).
Surjectivity from ambient degree-$d$ hypersurfaces to
$H^0(X,L^d)$ follows from the two plane-cubic coordinate rings
and the Segre map exactly as in that theorem.

Finally the rational function
\[
                         j(E_t)=\frac{t^{12}}{t^3+1}
\]
is nonconstant.  If an elliptic curve $E$ has transcendental
$j(E)$, choose a root $u$ of
$u^{12}-j(E)(u^3+1)=0$ in an algebraically closed extension.
Then $u$ is transcendental over $\mathbf F_2$ and $u^3+1\ne0$.
The resulting inclusion $\mathbf F_2(t)\hookrightarrow\bar k$,
$t\mapsto u$, identifies $E$ geometrically with the base change
of $E_t$, since elliptic curves with the same $j$-invariant
are isomorphic over an algebraically closed field.  The origins
and their cubic line bundles are respected.  Geometric generic
ranks are invariant under this field extension, proving the
last assertion.
\end{proof}

\begin{remark}[The algebraic-factor refinement and its subsequent resolution]
\label{rem:DE-A3-mixed-fixed-refinement}
\status{Proved under stated hypotheses}
The specialization argument alone does not infer the formula for every
ordinary elliptic curve over $\overline{\mathbf F}_2$.
Theorem~\ref{thm:DE-A3-ppav-ranks} subsequently proves the formula
for all these factors, in every degree, as part of its universal
principal-surface result.  For the explicit algebraic
factor $y^2+xy=x^3+1$, exact retained Cartier calculations attain
the predicted bounds in degrees $1,\ldots,6$, and in every
rectangular system $(a,3)$ with $3\le a\le31$; the retained run
took $0.7093895580001117$ seconds.  The then-outstanding uniform
construction and failed sparse-section route remain recorded in
\texttt{notes/A3\_mixed\_rank\_research.md}.  These finite probes
are not substituted for the specialization proof or for the
later uniform proof.
The separate auxiliary certificate
\path{agents/obstructions_crystals/code/A3_mixed_generic_checks.py}
checks the integral Weierstrass invariants, their characteristic-two
reductions, the nonvanishing denominator in the $j$-invariant argument,
and the exact genus/rank identities for $10{,}000$ degrees.  It passed
in $0.04470579700137023$ seconds; its complete output is retained
beside the program.  These numerical identities support the displayed
construction, whereas the proof above supplies its uniform geometric
existence assertion.
\end{remark}


% Root proof, 2026-11-19. Independent audit requested before promotion.
\subsection{Two-jet geometry of a genus-two theta surface in characteristic two}
\label{subsec:ppav-two-jets}
\status{Proved. The obstruction/crystal agent independently audited the complete
jet and discriminant proofs, and separately audited the supplementary dichotomy
and ordinary-member arguments.}

\begin{lemma}\label{lem:ppav-surface-dichotomy}
\status{Proved under stated hypotheses}
Over an algebraically closed field, a principally polarized abelian surface
with an effective symmetric theta divisor is either a product of two elliptic
curves with its product theta divisor, or the Jacobian of a smooth genus-two
curve with its smooth theta divisor. This assertion includes characteristic two.
\end{lemma}
\begin{proof}
An abelian variety contains no rational curve. Indeed a nonconstant map from
$\mathbf P^1$ factors through a finite number of Frobenius maps until its
differential is nonzero, whereas invariant differentials generate the
cotangent bundle of the abelian variety and $H^0(\mathbf P^1,\Omega^1)=0$.
For any integral curve $B$ on the surface, adjunction therefore gives
$B^2=2p_a(B)-2\ge0$. Also $\Theta^2=2$ and $\Theta\cdot B>0$.
Writing $\Theta=\sum m_iB_i$ shows that a reducible theta divisor must
have exactly two components, each with multiplicity one: their positive
integer intersections with $\Theta$ sum to two. A single component of
multiplicity at least two would give $m_i^2B_i^2=2$, which is impossible.
In the reducible case, evenness of $B_i^2$ and
$B_i^2+B_1\cdot B_2=1$ force $B_i^2=0$ and $B_1\cdot B_2=1$.
Each component is consequently a smooth genus-one curve, hence a translate
of an elliptic subgroup. Addition of the two elliptic subgroups is an
isogeny of degree their intersection number, namely one. It is the
required product isomorphism with its theta divisor.

If $\Theta$ is integral, $p_a(\Theta)=2$. Its normalization cannot have
genus one: a map from an elliptic curve to an abelian variety has a smooth
elliptic translate as image, which would have self-intersection zero.
Thus its normalization has genus two and arithmetic genus two, so
$\Theta$ is smooth. Its inclusion induces an isogeny
$u:\operatorname{Jac}(\Theta)\to A$, since this ample curve generates
$A$. Let $C_J$ denote the genus-two Abel--Jacobi curve in its Jacobian.
Then $u^*\Theta\cdot C_J=2$ and $C_J^2=2$. Hodge index gives
$(u^*\Theta)^2\le2$, whereas this square is $2\deg u$.
Thus $\deg u=1$ and $u$ is an isomorphism. Its restriction maps
$C_J$ onto $\Theta$, so it identifies the principal polarizations.
\end{proof}

We work in this subsection over an algebraically closed field $k$ of
characteristic two. Let $C$ be a smooth genus-two curve, choose a theta
characteristic $\kappa$ with $\kappa^{\otimes2}=\omega_C$, and put
\[
 A=\operatorname{Pic}^0(C),\qquad
 i:C\hookrightarrow A,\quad c\longmapsto\mathcal O_C(c)\otimes\kappa^{-1},
 \qquad \Theta=i(C),\qquad L=\mathcal O_A(3\Theta).
\]
Thus $\Theta$ is smooth and symmetric and $\Theta^2=2$. These Jacobian facts, including the characteristic-two
hyperelliptic description, are recorded in
\cite[Theorem~14.23 and Corollary~14.24, pp.~230--231; Theorem~14.29, p.~233;
Example~14.34, pp.~236--237]{EdixhovenVanDerGeerMoonenAV}.
Write $\Theta_a=\Theta+a$ and
$P_a=t_a^*\mathcal O_A(\Theta)\otimes\mathcal O_A(-\Theta)$.
The principal polarization identifies $a\mapsto P_a$ with an isomorphism
$A\to\operatorname{Pic}^0(A)$. Additions and integer multiples of points
below are in this group law. Put $W=H^0(A,L)$, so $\dim W=9$.
For a point $P\in A(k)$, a first jet means the value in
$L_P\otimes\mathcal O_{A,P}/\mathfrak m_P^2$, of dimension three;
a second jet means the analogous quotient by $\mathfrak m_P^3$, of
dimension six. All jet maps in this subsection are $k$-linear evaluation maps.

\begin{lemma}\label{lem:ppav-theta-preliminaries}
\status{Proved under stated hypotheses}
Every line bundle $\mathcal O_A(2\Theta)\otimes P_a$ is globally generated
and has no positive-degree cohomology. The system $|2\Theta|$ separates
two distinct points $x,y$ whenever $y\ne -x$.
Moreover, under $\operatorname{Pic}^0(C)=A$,
\[
                     i^*P_a=-a.
\]
The Gauss map of $\Theta$ to the projective line of cotangent directions
is the canonical degree-two map of $C$.
\end{lemma}

\begin{proof}
The theorem of the square gives
$\Theta_c+\Theta_{-c}\in|2\Theta|$.
For a specified point $x$, choose $c$ outside the two curves
$\Theta+x$ and $\Theta-x$; the resulting section is nonzero at $x$.
If $y\ne x,-x$, choose $c\in\Theta+x$ outside
$(\Theta+y)\cup(\Theta-y)$. These are distinct irreducible translates:
the translation stabilizer of $\Theta$ is contained in the kernel of
its principal polarization and is trivial. Their intersections with
$\Theta+x$ are finite, so this choice exists. The resulting section
vanishes at $x$ and is nonzero at $y$. Together with a section nonzero
at both points this proves separation. Every twist by $P_a$ is a
translate of $\mathcal O(2\Theta)$, since multiplication by two on
$A(k)$ is surjective, so it is globally generated as well.

For completeness the required ample vanishing does not use
characteristic-zero Kodaira vanishing. If $n$ is odd, the finite etale
map $[n]:A\to A$ has invertible degree $n^4$, and its trace splits the
unit map of structure sheaves. Thus
$H^j(A,T)\hookrightarrow H^j(A,[n]^*T)$ for a line bundle $T$.
For $T=\mathcal O(2\Theta)$, symmetry gives
$[n]^*T\simeq\mathcal O(2n^2\Theta)$; Serre vanishing for arbitrarily
large odd $n$ proves the assertion for $j>0$. Translation proves it
for the twists. The same argument and Riemann--Roch give
$h^0(3\Theta)=(3\Theta)^2/2=9$.

Let $\iota$ be the hyperelliptic involution of $C$; then
$i(\iota c)=-i(c)$ and $c+\iota c$ is canonical. For a generic
$a=i(r)-i(s)$, the intersection $\Theta\cap\Theta_a$ consists of
$i(r)$ and $i(\iota s)$. Indeed these are intersection points, they
are distinct generically, and the intersection has length two.
The difference map $\Theta\times\Theta\to A$ is surjective, since
two ample translates intersect. Restricting $\mathcal O(\Theta_a)$
to $C$ therefore gives generically
$\mathcal O_C(r+\iota s)$. Since
$\mathcal O(\Theta_a)=\mathcal O(\Theta)\otimes P_{-a}$ and
$i^*\mathcal O(\Theta)=\omega_C=\mathcal O_C(s+\iota s)$,
we obtain $i^*P_{-a}=a$. Equality extends from this dense set to the
morphisms $A\to\operatorname{Pic}^0(C)$, proving the formula for all $a$.
Finally the restriction of invariant one-forms on the Jacobian to
$C$ is the defining isomorphism
$H^0(A,\Omega_A^1)\simeq H^0(C,\omega_C)$.
The conormal direction of $\Theta$ is consequently its canonical
map. This map has degree two and is finite and surjective onto
$\mathbf P^1$. No assertion about a characteristic-two Kummer
quotient is needed.
\end{proof}

\begin{lemma}\label{lem:ppav-three-distinct-values}
\status{Proved under stated hypotheses}
In the smooth-theta case, any three distinct points of $A$ impose three
independent value conditions on $H^0(A,3\Theta)$.
\end{lemma}
\begin{proof}
It suffices to separate the value at $P_3$ from zero values at
$P_1,P_2$. Choose $a\in\Theta+P_1$ and $b\in\Theta+P_2$.
Require $a,b\notin\Theta+P_3$ and $P_3+a+b\notin\Theta$.
The first two conditions define dense open subsets of their respective
irreducible curves, because the three points are distinct and a theta
divisor has trivial translation stabilizer. The third defines a
nonempty open subset of their product: the sum of these two translates
of $\Theta$ is all of $A$. These three opens have nonempty intersection.
The product of theta sections for $\Theta_a,\Theta_b,\Theta_{-a-b}$
then vanishes at $P_1,P_2$ and not at $P_3$. Repeating with each
point in the third position proves the assertion.
\end{proof}

\begin{theorem}\label{thm:ppav-three-theta-jets}
\status{Proved under stated hypotheses}
With the data just specified, the following assertions hold.
\begin{enumerate}
\item $L$ separates points and first jets, and it separates second jets
at every point of $A$.
\item For $P\ne Q$, the map
\[
 W\longrightarrow
 L_P\otimes\mathcal O_{A,P}/\mathfrak m_P^2\ \oplus\
 L_Q\otimes\mathcal O_{A,Q}/\mathfrak m_Q^2
\]
has rank at least four. It has rank six outside the dimension-two locus
\[
 R=\{(P,Q):2P+Q\in\Theta,\ P+2Q\in\Theta\}.
\]
Its rank is at least five outside the dimension-at-most-one locus
\[
 R_4=\bigcup_{s\in A[3](k)}
       \{(P,s-P):P+s\in\Theta\}.
\]
The rank-four locus is contained in $R_4$, whether or not the rank
actually drops to four at each listed pair.
\item If $A$ is ordinary, then for every fixed $P$ the fibre $R_P$
is finite. In particular two first jets have full rank six for all
but finitely many $Q$ on any specified integral ample curve.
\end{enumerate}
\end{theorem}

\begin{proof}
For first jets at $P$, choose a theta translate through $P$ and
multiply its equation by a section of the remaining twist of
$2\Theta$ nonzero at $P$. The possible linear terms span both
cotangent directions because the Gauss map of $\Theta$ is
surjective. A product nonzero at $P$ supplies the value. Point
separation follows by choosing a theta translate through one point
and not the other, and a nonzero section of the remaining twist
at the latter point. Thus $L$ is very ample as well.

For second jets, fix $P$ and take $p,q\in\Theta$. Put
$a=P-p$, $b=P-q$, and $c=-a-b$. The product of the three theta
sections for $\Theta_a,\Theta_b,\Theta_c$ belongs to $|3\Theta|$.
Its first two factors vanish at $P$; the third is nonzero precisely
when $3P-p-q\notin\Theta$. This condition defines a nonempty open
subset of $\Theta\times\Theta$, because the addition map to $A$
is surjective. The product of the two Gauss maps is finite and
surjective onto $\mathbf P^1\times\mathbf P^1$. The image of that
open is dense, and the associated products of two linear forms
span all of $\operatorname{Sym}^2T_P^*A$: multiplication of binary
linear forms has dense image in the projective space of binary
quadratic forms, also in characteristic two. The resulting
quadratic terms therefore span
$\mathfrak m_P^2/\mathfrak m_P^3$. Combined with first-jet
surjectivity this proves second-jet surjectivity.

Now write $s=P+Q$. The theta translates through both $P,Q$ are
parametrized by the length-two scheme
\[
                      Z=(\Theta+P)\cap(\Theta+Q).
\]
It is preserved by $a\mapsto s-a$. A fixed point of this involution
has $2a=s$; the corresponding points $P-a,Q-a$ on $\Theta$ are
hyperelliptic conjugates and have equal Gauss direction, so its
intersection multiplicity is at least two. Consequently two
distinct support points of $Z$ are exchanged, and a single support
point is fixed.

Fix one such parameter $a$, and write $p=P-a$, $q=Q-a$ as points
of $C$. The restriction exact sequence is
\[
 0\longrightarrow H^0(2\Theta\otimes P_a)
 \longrightarrow W\longrightarrow H^0(\Theta_a,L|_{\Theta_a})
 \longrightarrow0;
\]
the surjectivity follows from Lemma~\ref{lem:ppav-theta-preliminaries}.
After translation to $C$, its last line bundle is
$M_a=\omega_C^{\otimes3}\otimes(-3a)$, of degree six.
The four along-curve jet conditions have rank at least three.
Indeed their cokernel is $H^1(C,M_a(-2p-2q))$, which has dimension
at most one for this degree-two line bundle. It is one exactly
when $M_a(-2p-2q)=\omega_C$. Since
$\mathcal O_C(p+q)\otimes\omega_C^{-1}$ has class $s-2a$,
this happens exactly when
\[
                              a=2s.
\]
The sections vanishing on $\Theta_a$ supply normal derivatives of
rank at least one. They supply rank two if $a\ne -s$: choose
$b\in A(k)$ with $2b=a$. Their evaluations are those of
$2\Theta$ at $P+b,Q+b$, which are distinct and are not opposite
unless $a=-s$. Lemma~\ref{lem:ppav-theta-preliminaries} applies.
Projection of the surface jet space to the along-curve jet space
has the two normal directions as its kernel. Hence the full rank
is at least the sum of the along-curve rank and the supplied normal
rank. It is therefore at least four always, six if a parameter
$a$ avoids both $-s,2s$, and five unless those two bad parameters
coincide.

If every support point of $Z$ is among $-s,2s$, then $-s$ is a
support point. For two support points this is immediate. For one
support point equal to $2s$, its being fixed gives $4s=s$, so
$3s=0$ and $2s=-s$. Thus rank failure from six implies that
$\Theta_{-s}$ contains both $P,Q$, which is exactly the displayed
condition defining $R$. Rank below five requires $3s=0$, and then
these two containment conditions reduce to $P+s\in\Theta$.
The map $(P,Q)\mapsto(2P+Q,P+2Q)$ is an etale isogeny of
$A\times A$ (its determinant is three). Its inverse image of
$\Theta\times\Theta$ has dimension two. Since $A[3](k)$ has
$3^4=81$ elements, $R_4$ has dimension at most one.

For the final assertion suppose a fibre $R_P$ were positive
dimensional. Its first condition places it on the irreducible curve
$\Theta-2P$; the second would then force
\[
                           [2]\Theta=\Theta+3P.
\]
When $A$ is ordinary, write $[2]=V\circ F_{A/k}$, where $V$ is
etale. The restriction of $F_{A/k}$ to $\Theta$ has degree two.
The restriction of $V$ to $\Theta^{(2)}$ has generic degree one:
any larger degree would make that curve invariant under a nonzero
element of the constant kernel of $V$, whereas a principal theta
divisor has trivial translation stabilizer. Thus $[2]|_\Theta$
has degree two. The asserted image equality contradicts the
projection formula: its left interpretation gives
$([2]_*\Theta)\cdot\Theta=2\Theta^2=4$, while the projection
formula gives $\Theta\cdot[2]^*\Theta=4\Theta^2=8$.
This proves finiteness and the last assertion.
\end{proof}

\begin{corollary}\label{cor:ppav-discriminant-one-node}
\status{Proved under stated hypotheses}
The discriminant of singular divisors in $|3\Theta|$ is irreducible
of dimension seven. Its general member is integral and has exactly
one singularity, an ordinary node. These assertions hold for every
$C$ in the stated characteristic; ordinary $A$ is not needed here.
\end{corollary}

\begin{proof}
The incidence of a section singular at a specified point is a
$\mathbf P^5$-bundle over $A$, hence is irreducible of dimension
seven. Second-jet surjectivity makes an ordinary node a nonempty
open condition on this incidence: the mixed coefficient in a
binary quadratic form is nonzero. The incidence of two distinct
singular points has dimension at most six. Indeed its rank-six
stratum has dimension at most $4+(8-6)=6$, its rank-five stratum
at most $2+(8-5)=5$, and its rank-four stratum at most
$1+(8-4)=5$. A general point of the single-point incidence thus
has exactly one singular point and it is an ordinary node.
Projection to $|3\Theta|$ is generically one-to-one, proving the
dimension and irreducibility assertions.

Such a divisor is reduced, since a multiple component would give
a positive-dimensional singular locus. It cannot be reducible:
if it splits as nonzero effective divisors $B+D\sim3\Theta$,
put $t=\Theta\cdot B$, so $1\le t\le5$. Hodge index gives
$B^2\le t^2/2$, and hence
\[
 B\cdot D=3t-B^2\ge 3t-t^2/2\ge5/2.
\]
Its intersection multiplicity is therefore at least three, whereas
two components of a curve with a unique ordinary node can have
total intersection multiplicity at most one. This contradiction
proves integrality.
\end{proof}

\begin{lemma}\label{lem:ppav-ordinary-members}
\status{Proved under stated hypotheses}
If the genus-two Jacobian $A$ above is ordinary, $|3\Theta|$ contains
an ordinary connected nodal divisor with three smooth theta components.
The general discriminant member is ordinary. For every $d\ge1$, a general
smooth member of $|3d\Theta|$ is ordinary, of genus $9d^2+1$.
\end{lemma}
\begin{proof}
Take the three translates $\Theta_a,\Theta_b,\Theta_{-a-b}$.
Their sum belongs to $|3\Theta|$ by the theorem of the square.
They are distinct and pairwise transverse for general $(a,b)$.
To see transversality without a generic-smoothness assumption in positive
characteristic, the difference map $\Theta^2\to A$ is generically etale:
its differential is invertible when the two Gauss directions differ.
Its critical values thus form a proper closed subset of $A$.
Each of the differences $b-a,-2a-b,-a-2b$ varies surjectively in $A$,
so a general pair avoids all three exceptional subsets.
Triple intersections are absent for a general pair as well. A triple
intersection corresponds to $x,p,q,r$ with $p,q,r\in\Theta$ and
$3x=p+q+r$, followed by $a=x-p,b=x-q$. Since $[3]$ is finite,
this incidence has dimension three, below the four-dimensional
parameter space of $(a,b)$.

The three components are ordinary because their Jacobian is the ordinary
surface $A$. The normalization sequence for a nodal curve expresses its
$H^1(\mathcal O)$ as an extension of the component cohomologies by its
graph cohomology. Frobenius is bijective on both terms, so this union
is ordinary. Its ordinary point belongs to the irreducible discriminant
of Corollary~\ref{cor:ppav-discriminant-one-node}, whose ordinary open
is therefore dense. Ordinarity is also a nonempty open condition in
the entire linear system. Its intersection with the nonempty smooth
open supplies smooth ordinary members of $|3\Theta|$.

Finally take $d$ general such members, meeting transversely without
triple points. Their nodal union lies in $|3d\Theta|$ and is ordinary
by the same normalization sequence. The positive-degree ample vanishing
used above makes $h^1(\mathcal O_Y)=9d^2+1$ constant for this family
of effective divisors. Invertibility of its coherent Frobenius is
therefore open, and intersecting again with the smooth open gives
the asserted general ordinary curve and genus.
\end{proof}

\begin{remark}[Source scope]\label{auto:A_ppav_two_jets:0557c5a87cfe}
\status{Proved under stated hypotheses}

The source search found Bauer--Szemberg's higher-jet statements
\cite[Theorems~1 and~2, conventions p.~2]{BauerSzemberg1997}
and Alvarado--Pareschi's recent three-theta second-jet theorem
\cite[Theorem~D; Section~2 conventions]{AlvaradoPareschi2026},
but their stated conventions are in characteristic zero. We do
not import either statement in characteristic two. The proof
above uses the theta theorem of the square, adjunction on a
genus-two curve, the elementary product sections of $2\Theta$,
and the explicit incidence calculation. The fixed-point finiteness
assertion explicitly requires ordinary $A$.
\end{remark}


% Promoted after complete root and reduction/descent mathematical audits.
\subsection{Fixed-point jet finiteness without ordinarity}
\label{subsec:nonordinary-ppav-fixed-jets-candidate}
\status{Proved and independently checked.}

\begin{lemma}[The doubled image of a smooth principal theta curve]
\label{lem:all-theta-double-image-candidate}
\status{Proved under stated hypotheses}
\textup{[Proved; complete independent proof audits passed.]}
Let $(A,\Theta)$ be a principally polarized abelian surface
over an algebraically closed field of characteristic two,
with smooth theta divisor. Then the finite map
\[
                         [2]|_\Theta:\Theta\longrightarrow[2]\Theta
\]
has inseparable degree two and separable degree one. In
particular its reduced image is numerically equivalent to
$2\Theta$.
\end{lemma}
\begin{proof}
The smooth theta curve has genus two and identifies $A$
with its Jacobian by Lemma~\ref{lem:ppav-surface-dichotomy}.
Write $[2]=V\circ F_{A/k}$, with
$V:A^{(2)}\to A$ the Verschiebung.
The differential $dV$ has the rank of the Cartier operator
on the two-dimensional space of canonical forms of $\Theta$.
Here is the relevant type identification: on the Jacobian,
$V$ is pullback of line bundles by the relative Frobenius
of the curve; its tangent map is the linearized Frobenius
on $H^1(\mathcal O_\Theta)$. Serre duality identifies its
dual with relative Cartier on $H^0(\omega_\Theta)$.
This is the same Picard/Frobenius and finite-duality
identification used in the preceding rank and covector
proofs.

This rank cannot be zero. Choose a separating rational
parameter $x$ on the smooth curve. In characteristic two
every rational differential killed by Cartier has the
form $b^2\,dx$: write its coefficient in the two-basis
$1,x$ over the field of squares and apply Cartier.
If both independent canonical forms were killed, their
ratio would be a square, say $z^2$. But their ratio
defines the canonical degree-two map to $\mathbf P^1$.
The equality
\[
 [k(\Theta):k(z^2)]
       =[k(\Theta):k(z)]\,[k(z):k(z^2)]=2
\]
would give $k(\Theta)=k(z)$, making the genus-two curve
rational. This is impossible. Thus $dV$ has rank one or
two.

If its rank is two, the restriction of $V$ to
$\Theta^{(2)}$ is separable. Suppose instead its rank is
one and that this curve restriction has zero differential.
The invariant line $\ker(dV)$ is then tangent to
$\Theta^{(2)}$ everywhere. Choose a nonzero invariant
vector field $\delta$ spanning it. Locally, if $f=0$
defines that smooth Cartier divisor, tangency says
$\delta(f)\in(f)$. First-order translation by
$\varepsilon\delta$, over $k[\varepsilon]/(\varepsilon^2)$,
therefore preserves the divisor: it sends its equation to
$f+\varepsilon\delta(f)=f(1+\varepsilon a)$.
It consequently preserves its line bundle. This places
the nonzero tangent vector $\delta$ in the scheme kernel
of the principal polarization
$\lambda_{\mathcal O(\Theta^{(2)})}$, which is an
isomorphism. The contradiction proves that
$V|_{\Theta^{(2)}}$ is separable also in rank one.
No integration of a vector field or reducedness assumption
on the kernel of $V$ is used here.

Its separable degree is one. Otherwise a generic fibre
would contain two distinct points differing by a nonzero
element of the finite set $\ker(V)(k)$. One such difference
would occur on a dense subset of the irreducible curve,
forcing a nonzero translation to preserve
$\Theta^{(2)}$. Its translation stabilizer is contained
scheme-theoretically in the kernel of its principal
polarization and is trivial. This is again impossible.

Relative Frobenius restricts on the smooth theta curve
with degree two. Thus $[2]|_\Theta$ has the asserted
degrees. By the projection formula,
$[2]_*\Theta\equiv4\Theta$ on the surface, since
$[2]^*$ multiplies divisor classes numerically by four.
Its generic curve degree is two, so the reduced image
satisfies $[2]\Theta\equiv2\Theta$.
\end{proof}

\begin{corollary}[The two-jet exceptional set has finite fixed fibres]
\label{cor:all-theta-fixed-point-jets-candidate}
\status{Proved under stated hypotheses}
\textup{[Proved; complete independent proof audits passed.]}
In Theorem~\ref{thm:ppav-three-theta-jets}, the fixed-point
finiteness assertion holds for every smooth-theta abelian
surface in characteristic two, whether ordinary or not.
Explicitly, for each fixed $P\in A(k)$ the set
\[
 R_P=\{Q:2P+Q\in\Theta,\ P+2Q\in\Theta\}
\]
is finite. Consequently the first jets at $P,Q$ have full
rank six for all but finitely many $Q$ on every prescribed
integral ample curve.
\end{corollary}
\begin{proof}
If $R_P$ were infinite, its first condition would make it
dense in the irreducible curve $\Theta-2P$. Its second
condition would force
\[
                          [2]\Theta=\Theta+3P.
\]
The left-hand reduced curve has numerical class $2\Theta$
by the lemma, while the right-hand curve has class
$\Theta$, a contradiction. The rank-six assertion is
then exactly part two of the cited two-jet theorem.
\end{proof}

\begin{remark}[Scope]
\label{rem:all-theta-fixed-point-jets-scope-candidate}
\status{Proved under stated hypotheses}
\textup{[Proved; complete independent proof audits passed.]}
This removes only the ordinarity restriction from the
fixed-point exceptional-jet argument. It asserts neither
relative Hasse-divisor irreducibility nor full monodromy
for a fixed nonproduct nonordinary surface. In particular
one cannot pad an odd sharp norm curve on a nonordinary
surface with another ample component and retain its sharp
rank: the positive component defects add on a nodal union.
The elementary first-order stabilizer proof above avoids
any unsupported extension of an ordinary Hecke or Hasse
argument.
\end{remark}


% /root/monodromy, 2026-11-19. Proved; independent root and reduction_descent audits passed.
\subsection{A level-three theta symmetrizer for the Cartier pencil}
\label{subsec:ppav-theta-Cartier-symmetrizer}
\status{Proved and checked. The parent and reduction/descent agent
independently audited the complete proof, opened the cited primary sources,
and reran the exact finite Heisenberg calculation. The relative Cartier
and theta-group identifications include the Frobenius twists explicitly.}

\begin{theorem}\label{thm:ppav-level3-Cartier-symmetric}
\status{Proved under stated hypotheses}
Let $(A,\Theta)$ be a principally polarized abelian variety of dimension
$g$ over an algebraically closed field $k$ of characteristic two, and let
$L=\mathcal O_A(3\Theta)$ and $W=H^0(A,L)$.
Choose a nonzero invariant top form $\omega$ and put
$C_{\rm abs}(\omega)=c^{1/2}\omega$, with $c\in k$.
There is a theta basis of $W$, indexed by $\mathbf F_3^g$, in which
the root-removed Cartier multiplication matrix $M(s)$ is symmetric
and has diagonal
\[
                         \operatorname{diag}M(s)=c\,s.
\]
If $A$ is ordinary, normalize $\omega$ by $C_{\rm abs}(\omega)=\omega$.
Then the diagonal coordinates are independent and $\det M(s)$ is a
nonzero geometrically squarefree polynomial of degree $3^g$.
In particular, for $g=2$ the degree-one full coherent Hasse divisor of
$|3\Theta|$ is geometrically reduced and each of its components has
generic linearized Cartier corank one.
\end{theorem}

\begin{proof}
\par\smallskip\noindent\textbf{Relative Cartier and its target.}
Let $F=F_{A/k}:A\to A^{(2)}$ be relative Frobenius. Put
$L^{(2)}=L\otimes_{k,F_k}k$ on the Frobenius twist and
$W^{(2)}=H^0(A^{(2)},L^{(2)})$. The canonical map
$\iota:W\to W^{(2)}$, $w\mapsto w^{(2)}$, is square-semilinear.
The relative top Cartier trace is $k$-linear and satisfies
\[
 C_{\rm rel}(\omega)=c\,\omega^{(2)},\qquad
 C_{\rm rel}(F^*a\,\eta)=a\,C_{\rm rel}(\eta).
\]
Since $L^2=F^*L^{(2)}$, multiplication followed by this trace defines
the $k$-bilinear symmetric map
\begin{equation}\label{eq:ppav-relative-beta}
 \beta:W\otimes W\longrightarrow W^{(2)},\qquad
 C_{\rm rel}(s t\omega)=\beta(s,t)\omega^{(2)}.
\end{equation}
In particular $\beta(s,s)=c\,s^{(2)}$.
These are the relative finite-duality and coefficient-trace formulas
of \cite[Lemma~1.3.6, Proposition~1.3.7, Theorem~1.3.8, pp.~26--28]
{BrionKumar2005}. The target is the Frobenius twist; no identification
of $A$ with $A^{(2)}$ is being made.

\par\smallskip\noindent\textbf{A finite theta subgroup and its two representations.}
The theta group of $L$ is a central extension of
$K(L)=A[3]\simeq\mathbf F_3^{2g}$ by $\mathbf G_m$ with nondegenerate
commutator pairing. Choose a Lagrangian decomposition and lifts of
its generators of order three. They, together with the central
$\mu_3$, generate a finite Heisenberg subgroup $H$ of order
$3^{2g+1}$. A primitive cube root $\zeta\in k$ identifies its
Schrodinger representation on a basis $e_a$, $a\in\mathbf F_3^g$, by
\[
 T_x e_a=e_{a+x},\qquad D_y e_a=\zeta^{y\cdot a}e_a,
 \qquad Z e_a=\zeta e_a.
\]
Here $T_x,D_y,Z$ generate $H$; their commutation convention merely
chooses the sign of the standard symplectic pairing.
The natural theta action on $W$ is this representation, denoted
$\rho$. The precise primary statements are
\cite[Corollary~8.20, Lemma~8.24, Theorem~8.28, description~8.29,
pp.~119, 121--123; Corollary~9.17, pp.~133--134]
{EdixhovenVanDerGeerMoonenAV}. They give the nondegenerate theta
group, its Heisenberg form and the unique weight-one representation
on $H^0(L)$; its dimension is $3^g$.

Relative Cartier is natural under translations. A theta lift on $L$
induces its tensor square on $L^2$; under $L^2=F^*L^{(2)}$, this
tensor square is the pullback of its Frobenius-twisted lift on
$L^{(2)}$. The invariant form is unchanged by translation.
Consequently, in the bases $e_a$ and $e_a^{(2)}$, one has the
fully typed equivariance identity
\begin{equation}\label{eq:ppav-Cartier-equivariance}
 \beta(\rho(h)s,\rho(h)t)=\rho(h)^{[2]}\beta(s,t)
                         \quad(h\in H).
\end{equation}
One can also verify this directly in a local frame: if a lift acts
by a line-bundle multiplier $u$, its action on a product has multiplier
$u^2$, and the Cartier projection formula moves its Frobenius
descendant to the target. On the base the translation descends from
$x$ to $F(x)$. Thus the target action is the coefficient-square
representation, not the original weight-one action.

On the displayed generators of the finite subgroup,
\[
                \rho(h)^{[2]}=\rho(h)^{-t}.
\]
For $T_x$ both sides are the same permutation matrix; for $D_y$ and
$Z$ it is the identity $\zeta^2=\zeta^{-1}$. Since both sides are
representations, it holds for every $h\in H$.
Therefore
\[
 R:W^{(2)}\longrightarrow W^*,\qquad e_a^{(2)}\longmapsto e_a^*,
\]
is an invertible $H$-equivariant linear map. This step uses the finite
central subgroup $\mu_3$. The corresponding statement for the whole
central $\mathbf G_m$ would be false and is not used.

\par\smallskip\noindent\textbf{Cyclic symmetry is a translation in level three.}
Define the trilinear tensor
\[
                  T(s,t,u)=\langle R\beta(s,t),u\rangle.
\]
It is $H$-invariant by \eqref{eq:ppav-Cartier-equivariance} and is
symmetric in its first two arguments by multiplication.
Write $T_{abc}=T(e_a,e_b,e_c)$. Invariance under every modulation
$D_y$ forces
\[
            T_{abc}=0\quad\text{unless }a+b+c=0.
\]
Invariance under $T_x$ gives
$T_{a+x,b+x,c+x}=T_{abc}$.
If $a+b+c=0$, translation by $x=b-a$ sends the ordered triple
$(a,b,c)$ to $(b,c,a)$, because $2=-1$ in the index group
$\mathbf F_3^g$. Thus $T$ is invariant under cyclic permutation.
Together with its first-two symmetry this proves full symmetry
under all permutations of its three arguments.

Let $M(s)$ represent $\beta(s,-)$ from the $e_a$ basis to the
$e_a^{(2)}$ basis. The matrix of $R$ in these chosen bases is the
identity. Full symmetry of $T$ gives $M(s)^t=M(s)$.
Furthermore $\beta(e_a,e_a)=c e_a^{(2)}$, so
$T_{aab}=c\delta_{ab}$. By full symmetry, the $a$th diagonal
entry of $M(s)$ is $c s_a$, proving the asserted formula.

For ordinary $A$, Cartier on its invariant top form is nonzero.
If $C_{\rm abs}(\omega)=u\omega$ initially, replacing $\omega$
by $u^2\omega$ normalizes it; hence $c=1$.
Lemma~\ref{lem:symmetric-independent-diagonal-squarefree} now
applies: the determinant has the full squarefree coordinate
monomial with coefficient one, while every nontrivial matching
term has an even positive exponent. It is nonzero and has no
repeated geometric irreducible factor.
At the generic point of a reduced determinant component the
matrix has corank one: corank at least two would make every
cofactor, and hence every first derivative of the determinant,
vanish. This finishes the proof.
\end{proof}

\begin{remark}[Coordinate changes and the finite check]\label{auto:A_monodromy:bcd64c5e340f}
\status{Proved under stated hypotheses}

In an arbitrary original basis the conclusion is the existence of
a constant invertible matrix $R_0$ for which $R_0M(s)$ is symmetric
with independent diagonal coordinates. The matrix $R_0$ itself
need not be symmetric, because the change of basis on the target
is Frobenius-twisted. This is consistent with the explicit elliptic
product symmetrizers proved earlier.
The retained finite calculation
\path{agents/monodromy/code/level3_theta_tensor_checks.py}
checks the $g=2$ index action, all $243$ finite Heisenberg matrices,
and the full permutation symmetry of the five-dimensional space of
invariant first-two-symmetric tensors. It is a check of the finite
algebra, not a substitute for the geometric equivariance proof.
\end{remark}


% /root/monodromy, 2026-11-19. Proved; independent root, obstructions_crystals and reduction_descent full audits passed.
\subsection{Eventual exact groups on every ordinary principally polarized surface}
\label{subsec:ordinary-ppav-surface-full-groups}
\status{Proved and checked. The parent and both other research agents
independently audited the complete assembly, including every exceptional jet
stratum, the exact invariant quotient, the norm degeneration, actual natural
block extraction and the arithmetic image. The cited source scopes were
checked independently; the level-three finite calculation was also rerun.}

Throughout this subsection $(A,\Theta)$ is an ordinary principally polarized
abelian surface in characteristic two, $L=\mathcal O_A(3\Theta)$, and the
linear systems are complete unless ambient hypersurfaces are explicitly
mentioned. A symmetric representative of the polarization is chosen over the
algebraic closure. Translation of that representative gives an isomorphic
geometric family and does not change its geometric group.

\begin{proposition}[All invariant covectors on an ordinary surface]
\label{prop:ordinary-ppav-all-covectors}
\status{Proved under stated hypotheses}
For every $d\ge1$ the geometric generic relative rank is $r_d=9d^2-1$.
The entire geometric invariant-covector space has dimension zero for odd $d$
and three for even $d$. In even degree it is canonically the quotient
\[
 \mathbf F_2^{H}/\mathbf F_2\mathbf 1,
 \qquad H=\operatorname{Hom}(\ker F_A,\mathbf G_m),\quad |H|=4,
\]
with its natural Galois action. Equivalently its basis can be indexed by the
three nonzero characters, which Galois permutes.
\end{proposition}
\begin{proof}
The ordinary nodal degeneration of Lemma~\ref{lem:ppav-ordinary-members}
proves generic ordinarity and the rank formula in the indecomposable case.
The ordinary product degeneration proves it in the decomposable case.
The coherent relative bundle on $P=\mathbf P H^0(A,L^d)$ is $\Omega_P^1$:
indeed $H^1(A,L^{-d})=0$, $h^2(A,L^{-d})=9d^2$, and
$h^2(A,\mathcal O_A)=1$, and the multiplication map to the last group is
surjective, dual to the inclusion of the defining section.

We give the complete endomorphism reduction, to specify its scope beyond
products. Normalize an invariant two-form by $C(\omega)=\omega$ and put
$W=H^0(A,L^d)$. An invariant covector on the ordinary open extends across
all prime divisors of $P$: its equation in regular frames is
$M^t\lambda=\lambda^{[2]}$, and a largest pole of order $a>0$ would have
order $2a$ on the right and at most $a$ on the left. Normality and local
freeness extend it across codimension two. The dual Euler sequence then
represents it by a linear operator $D\in\End(W)$ modulo scalar identity.
Its Frobenius defect is a constant multiple of $s^{[2]}$; adding a scalar
identity and solving an Artin--Schreier equation removes that defect.
Thus exact lifts are the operators satisfying
\[
                 \beta(s,D(s))=D(s)^{[2]},\qquad s\in W,
\]
modulo $\mathbf F_2\operatorname{id}$, where $\beta$ is relative Cartier
multiplication and $\beta(s,s)=s^{[2]}$.

The eigenvalue and Jordan argument of
Theorem~\ref{thm:product-char2-all-parities-fixed-quotient} applies without
any product assumption: eigenvalues are zero or one, and a Jordan chain
$Dv=w,Dw=0$ would give the nonzero coefficient $\beta(w,w)=w^{[2]}$ on
substituting $v+tw$. Every exact lift is consequently an idempotent, and
its two eigenspaces are orthogonal for $\beta$.

The bundle $\mathcal F=F_{A*}L^d$ is globally generated. For an
indecomposable surface, second-jet surjectivity for $3\Theta$ from
Theorem~\ref{thm:ppav-three-theta-jets} implies surjectivity on a
Frobenius fibre: in regular parameters its algebra is
$k[x,y]/(x^2,y^2)$, a quotient of the second-jet algebra since
$(x,y)^3\subset(x^2,y^2)$. For higher $d$, multiply by a section of
$L^{d-1}$ nonzero at the point. Such a section exists since $L$ is globally
generated. On a decomposable surface the two degree-three elliptic factors
interpolate the two length-two Frobenius fibres, giving the same assertion.
Finite Frobenius duality gives the perfect symmetric Cartier pairing
$\mathcal F\otimes\mathcal F\to L^{d,(2)}$.
Orthogonality therefore makes each of the two eigenspaces preserve the
kernel of every evaluation map. Global generation and perfectness descend
$D$ to a bundle endomorphism of $\mathcal F$, exactly as in the cited
product proof; this reasoning uses only these two properties.

If $d$ is odd, apply \cite[Theorem~1.2(i), printed p.~45; proof
pp.~43--45]{Oda1971} to $F_A$ and $L^d$.
This theorem is stated for abelian varieties of arbitrary dimension and
allows inseparable isogenies. Its hypothesis is that the polarization map
of $L^d$ restricts to an isomorphism of the scheme-theoretic kernel of the
isogeny onto its image. Here this restriction is the restriction of
$3d\lambda_\Theta$ to $\ker F_A$; it is an immersion since $3d$ is odd
and $\lambda_\Theta$ is an isomorphism. Thus
$\End(F_{A*}L^d)=k$. The only exact lifts are zero and identity, and their
quotient gives no invariant covector.

If $d$ is even, write $L^d=F_A^*M$.
The ordinary splitting of $F_{A*}\mathcal O_A$, equivalently the character
decomposition for $\ker F_A\simeq\mu_2^2$, gives
\[
          F_{A*}L^d=\bigoplus_{\chi\in H}M\otimes T_\chi.
\]
The four summands are distinct ample line bundles differing by nontrivial
degree-zero two-torsion line bundles. There are no homomorphisms between
different summands, and every summand has nonzero sections by ample
vanishing and Riemann--Roch. Hence $\End(\mathcal F)=k^4$.
Their Cartier cross pairings vanish, since trace kills nonzero character
weights. Its four projectors are precisely the independent exact lifts:
each projector satisfies the displayed equation, and any idempotent in
$k^4$ is their sum. Modding out the identity gives the stated three-space.
This splitting is \cite[Lemma~4.5, p.~11]{SannaiTanaka2016}, or directly
the character algebra of the ordinary Frobenius torsor.

The construction by projectors is canonical under Galois. The zero
character is fixed, and the three nonzero characters are permuted.
Modulo the sum of all four projectors, those three projectors are a basis.
This proves the final assertion, including its arithmetic description.
\end{proof}

\begin{theorem}[Hasse geometry and the irreducible common kernel]
\label{thm:ordinary-ppav-Hasse-kernel}
\status{Proved under stated hypotheses}
For every ordinary principally polarized abelian surface, the full Hasse
divisor in $|L|$ is geometrically reduced with generic Cartier corank one
on each component. For every $d\ge2$ the Hasse divisor in $|L^d|$ is
geometrically irreducible and reduced, with generic stable loss one.
Its geometric relative group is generated by transvections.
If $K_d$ is the common kernel of its entire invariant-covector space,
then $K_d$ is irreducible, and $V_d$ has no invariant vector.
\end{theorem}
\begin{proof}
The decomposable case was proved in
Theorem~\ref{thm:ordinary-products-Hasse-irreducible} and
Corollary~\ref{cor:ordinary-products-irreducible-kernel}.
In the indecomposable case use the genus-two theta divisor and jet
statements of Theorem~\ref{thm:ppav-three-theta-jets}.
Theorem~\ref{thm:ppav-level3-Cartier-symmetric} makes the degree-one Hasse
polynomial reduced with generic corank one. Its components have smooth
generic curves, since the singular-curve discriminant is irreducible and
has ordinary generic member by Lemma~\ref{lem:ppav-ordinary-members}.

We verify all incidence bounds needed to apply the mixed-tacnode argument,
including the rank-four exceptions. Put $e=d-1\ge1$. Let $H_i$ be a
component of the degree-one Hasse divisor, and choose a general smooth
$C\in H_i$. Its Cartier kernel is one-dimensional. The nonzero tangent
functional $dh_1$ belongs to the dual first-jet space at only finitely many
points of $C$: otherwise choose one such $P$ and another outside the
finite exceptional set $R_P$ of the jet theorem. The two first-jet dual
spaces are then disjoint, a contradiction. Choose $P$ also away from the
zeros of the killed differential and from every theta translate
$\Theta-s$, $s\in A[3]$. An integral curve of class $3\Theta$ is not one
of those theta translates, so these exclusions are finite on $C$.

The affine tangent space $\ker dh_1$ to $H_i$ at $C$ surjects onto the
first jets at $P$. The incidence of $D\in|L^e|$ simply tangent to $C$ at
$P$ therefore dominates $|L^e|$: for an infinitesimal change of $D$, move
$P$ to cancel the value change and use these three freely variable first
jets of $C$ to preserve its value and the common tangent line. This is
the three-equation calculation in
Theorem~\ref{thm:product-Hasse-irred}, on the affine cone; projective
scaling imposes no additional condition.

The good contact locus is nonempty. Every $L^e$ inherits the degree-one
second jets at $P$ and the joint first-jet rank bounds at $P,Q$: multiply
by a section of $L^{e-1}$ nonzero at both points, which exists over the
infinite algebraically closed field. Our choice of $P$ eliminates $R_4$,
so all two-point ranks are at least five and are six except at finitely
many $Q$. Within the two imposed contact conditions at $P$, a singularity
at a different $Q$ imposes total rank five generically with $Q$ varying
in dimension two, and total rank at least four at the finitely many
exceptional $Q$. Both are proper loci. Another contact with $C$ imposes
rank four generically with $Q$ varying in dimension one, and rank at least
three at finitely many exceptions, again proper. Second jets exclude
higher contact and a zero normal derivative at $P$. A general member in
the incidence is thus smooth with just this simple tangency. Dominance
then allows it to be ordinary. A single general ordinary $D$ works for
all the finitely many $H_i$.

For this fixed $D$, the contact divisor $\Delta_D\subset|L|$ is
irreducible: its incidence over $P\in D$ is a projective bundle given by
two independent value/tangent conditions. We justify that its generic
point is one simple contact with smooth $C$. For pairs on $D^2$, the
rank-six locus imposes four contact conditions; its point dimension is
two. The exceptional locus $R\cap D^2$ has dimension at most one, since
$R$ is finite over each first point, and imposes at least three conditions
away from $R_4$. The locus $R_4\cap D^2$ is finite: its first point lies
on some $D\cap(\Theta-s)$. It still imposes at least two independent
value conditions. All three counts give codimension at least two in
$|L|$ for two contacts. A singularity of $C$ at a second point imposes
five conditions generically, with a three-dimensional pair space; on
$R$ it imposes at least four with pair dimension at most one, and on
$R_4$ at least three with pair dimension zero. A singularity at the
contact point and order-three contact each impose three conditions with
that point varying on $D$. Thus they too have codimension at least two.
The generic contact is simple and $C$ smooth. The joint tangency incidence
in $(P,C,D)$ is irreducible and dominates both systems by first-jet
surjectivity, so its two ordinary inverse images meet in a dense open.
For general ordinary $D$, the generic $C\in\Delta_D$ is ordinary.

On the multiplication subspace $s_D|L|\subset|L^d|$, the divisorial
nonordinary support is exactly $\bigcup_i H_i\cup\Delta_D$.
This is the union stable-rank formula
\eqref{eq:product-union-stable-rank}. Nonreduced base members have
codimension at least two by the one-node discriminant theorem. A member
containing $D$ is impossible for $e>1$ and a single projective point for
$e=1$. Each $H_i$ meets $\Delta_D$ at the good mixed contact just
constructed. The pointed Cartier lemma and the direct dualizing-form
calculation of Lemma~\ref{lem:mixed-tacnode-Hasse} apply to any ordinary
abelian surface: its only ambient input is bijective Cartier on
$H^1(\mathcal O_A)$ and $H^2(\mathcal O_A)$, and its deformation input is
first-jet surjectivity. Thus the full Hasse divisor is smooth at every
chosen mixed point.

Every global Hasse component restricts to a nonzero positive-degree
polynomial on this multiplication subspace, since the latter contains
an ordinary point. Its restriction contains one of the displayed prime
divisors. Any component containing $H_i$ or $\Delta_D$ contains their
chosen intersection point. Smoothness there forces all these assignments
to be the same component, even if restriction multiplicities occur.
This proves global irreducibility; that same smooth point proves
multiplicity one. An ordinary pair with a single simple contact has
stable loss exactly one by the union formula. Closedness of the iterated
Frobenius rank loci therefore gives generic stable loss one.
Lemma~\ref{lem:char2-defect-one-inertia} now proves transvection generation.

Finally the coherent bundle is $\Omega^1_{\mathbf P^{9d^2-1}}$.
Lemma~\ref{lem:irreducible-hasse-quotients-trivial} implies that every
nonzero proper invariant subspace has trivial quotient. By
Proposition~\ref{prop:ordinary-ppav-all-covectors}, every such subspace
contains $K_d$, making $K_d$ irreducible. Its dimension is at least $32$,
so an invariant line is impossible. This proves $V_d^G=0$ as well.
\end{proof}

\begin{lemma}[A full degree-one open in surface moduli]
\label{lem:ordinary-ppav-GL8-open}
\status{Proved under stated hypotheses}
There is a nonempty open in the ordinary moduli stack of principally
polarized abelian surfaces on which the degree-one relative geometric
image for $|3\Theta|$ is $\GL_8(\mathbf F_2)$.
\end{lemma}
\begin{proof}
Work over an algebraically closed field. Use full level five and a fixed
primitive fifth root of unity. The corresponding fine moduli variety is
smooth and geometrically irreducible and carries its universal polarized
abelian scheme, by the exact statement
\cite[Theorem~7.3.11, printed p.~76]{ConradAbelianVarietiesModuli}.
This is an opened author exposition of the moduli theorem; no claim that
its original GIT or compactification proof has been reread is required.
Near the level-five point of the certified product
$E^2$, $E:y^2+xy=x^3+1$, choose a line bundle representing the principal
polarization after an etale neighbourhood. Such representatives form a
smooth torsor under the dual abelian scheme; rigidification at zero
supplies an actual line bundle after this localization. Do not require
this family of representatives to be symmetric: that additional torsor
need not be etale in characteristic two. Each geometric fibre is a
translate of a symmetric representative, so its geometric group is the
same. Take an irreducible smooth neighbourhood $B$ containing the chosen
point.

The cube of this line bundle has a rank-nine direct image, and its
projective bundle carries a nonempty smooth ordinary open $U$ with
geometrically integral fibres. Its rank-eight full Lang frame cover
$Z\to U$ is finite etale. The fibre at the product point is connected by
the certified degree-one $\GL_8$ theorem.
The normalization argument of Lemma~\ref{lem:ordinary-pair-GL8-open}
proves that $Z$ has geometrically connected generic fibre over $B$:
a finite extension defining distinct generic components gives a finite
normalization of the excellent base; on the normal integral projective
bundle open each open-and-closed component of the finite etale cover
surjects, contradicting the connected product fibre. The geometric
generic fibre is smooth, hence normal, so it is also irreducible.
By \cite[Lemma~37.27.5]{Stacks0559}, after shrinking $B$ the frame fibres
are geometrically irreducible. Their connectedness means full $\GL_8$
monodromy. The etale map from $B$ to fine moduli is open, and forgetting
level is open; its image is the asserted nonempty open in the moduli
stack. Translating the representing theta line bundle leaves each
geometric family isomorphic, so this property is intrinsic to the
principally polarized surface.
\end{proof}

\begin{lemma}[Nodal images and an actual natural block]
\label{lem:ordinary-ppav-isogeny-block}
\status{Proved under stated hypotheses}
Let $(A',\Theta')$ be an ordinary principally polarized surface with
full degree-one relative image $\GL_8(\mathbf F_2)$, and suppose
$\phi:A'\to A$ is an odd polarized isogeny with
$\phi^*\lambda_\Theta=n\lambda_{\Theta'}$, $n$ odd.
Then $\deg\phi=n^2$, and for every $d\ge n+1$ the degree-$d$
geometric group on $A$ contains a natural $\GL_8(\mathbf F_2)$ on an
actual direct summand, fixing a complementary summand pointwise.
\end{lemma}
\begin{proof}
All work is geometric. Put $L'=\mathcal O(3\Theta')$ and choose symmetric
representatives on these two fixed surfaces. The kernel of $\phi$ is a
constant odd finite group, acting by translations.
A general smooth ordinary $C'\in|L'|$ is preserved by no nonzero kernel
translation: one preserving the line bundle has a nontrivial projective
action in its very ample embedding and hence a proper fixed locus; one
not preserving the line bundle cannot preserve a member. Thus
$C'\to\phi(C')$ is generically one-to-one.

We prove that its image has only ordinary nodes. On an indecomposable
source, for a fixed nonzero kernel element $g$, the two first jets at
$P,P+g$ have rank six except at finitely many $P$. Indeed failure is
contained in
\[
               3P+g\in\Theta',\qquad 3P+2g\in\Theta'.
\]
The two theta translates are distinct because $g\ne0$ and the principal
theta divisor has trivial translation stabilizer. Their intersection is
finite, and $[3]$ is finite. On the rank-six locus, two value conditions
and proportionality of two nonzero gradients impose three independent
conditions. Varying $P$ adds two dimensions, so tangency is generically
absent. The exceptional finite set already imposes two independent value
conditions and is avoided by a general member.

Any three distinct points on an indecomposable $A'$ impose three
independent value conditions on $|3\Theta'|$. To separate $P_3$ from
$P_1,P_2$, choose theta translates $\Theta'_a,\Theta'_b$ through $P_1,P_2$
respectively and not through $P_3$. These choices are dense opens in two
irreducible theta translates. Their addition map dominates $A'$: sums of
two theta curves cover the surface, because any two ample translates
intersect. Thus one can also require the third translate
$\Theta'_{-a-b}$ to avoid $P_3$. The product of these three equations is
a section of $3\Theta'$ vanishing at $P_1,P_2$ but not $P_3$. Cyclically
repeat the construction to prove the three-value rank assertion.
For $P,P+g,P+h$, three independent conditions with $P$ varying in dimension
two exclude triple identifications. There are only finitely many kernel
pairs. On a decomposable source the already proved cubic-product
first-jet and Segre three-value argument gives the identical nodal
conclusion, including its one-dimensional collinear exception. Therefore
$Z=\phi(C')$ is integral nodal with normalization $C'$. For each $g$,
$C'\cap(C'-g)$ has degree $L'^2=18$; its node cover is a quotient of a
union of such degree-eighteen covers. Its permutation group is consequently
a subquotient of a product of copies of $S_{18}$.

The image line bundle is $\operatorname{Nm}_\phi L'$.
Its Neron--Severi class is $3n\Theta$: push forward
$\phi^*\Theta\equiv n\Theta'$ and cancel $n$ in the torsion-free
Neron--Severi group. Both this norm and $L^n$ are symmetric, so their
quotient is a geometric two-torsion line bundle. Since $3n$ is odd,
$\lambda_{L^n}$ identifies the geometric two-torsion of $A$ with that of
its dual. Translating the whole image family by one constant two-torsion
point therefore corrects this quotient and places $Z$ in the exact
linear system $|L^n|$. Denote the corrected finite etale map again by
$\phi$; the translation does not change its kernel identifications.

Fix a general smooth ordinary $D\in|L^{d-n}|$.
Its inverse image $D'=\phi^{-1}D$ is smooth, connected and integral:
it is a smooth ample divisor on the abelian surface, whose connectedness
also follows from ample vanishing in its structure-sheaf sequence.
For varying $C'$, the intersection with $D'$ has degree
\[
                         N=18n(d-n)\ge18.
\]
Its permutation group is the full $S_N$. Here are the required incidence
checks, so that this conclusion does not assume a generic linear-series
monodromy theorem. One-point and ordered two-point incidences are
irreducible projective bundles over $D'$ and $D'^2\setminus\Delta$,
respectively, because $L'$ separates one and two values.
A simple double zero and otherwise simple zeros is attainable. On an
indecomposable source choose $P\in D'$ outside the finite intersections
of $D'$ with the theta translates occurring in the projection of $R_4$.
These intersections are finite since $D'$ is not a theta translate.
The joint surface jets at $P,Q$ have rank at least five for every
$Q\ne P$, and rank six except at finitely many $Q$; their restrictions
to the two tangent directions of $D'$ have rank at least three and
four respectively. After imposing the two conditions for a double zero
at $P$, a second double zero is therefore a proper locus: the generic
second point varies in dimension one and adds two conditions, and each
exceptional point adds at least one. Second jets at $P$ exclude order
three. On a product source choose $P$ also outside the flex fibres;
the same bounds follow from the cubic-product jet lemma, since the
exceptional fibres meet the integral ample $D'$ finitely.
The resulting length-two collision is smooth in the point incidence
and its parameter, since first jets are generated.
The single-collision criterion
\cite[Proposition~3.1, pp.~7--8 of the opened version]{Entin2021}
then gives $S_N$, including characteristic two. Equivalently its proof
uses the irreducible one- and two-point incidences and the nontrivial
transposition inertia of this one ordinary double zero.

After shrinking the coefficient space, $Z$ and $D$ meet transversely,
the image nodes of $Z$ avoid $D$, and $Y=Z\cup D$ is ordinary nodal in
$|L^d|$. The node-avoidance condition is proper: two values at $P,P+g$
with the image point restricted to $D$ impose two conditions with a
one-dimensional point parameter. The normalization sequence and the
isomorphism $H^1_{\et}(A,\mathbf F_2)\xrightarrow{\sim}
H^1_{\et}(A',\mathbf F_2)$ under an odd isogeny give
\[
 0\subset T\subset T_0\subset V_d,\qquad
 T_0/T\text{ constant},\qquad V_d/T_0=Q,
 \qquad\dim Q=8,
\]
where $Q$ is the degree-one relative system of $C'$ and $T$ is graph
cohomology. Its image on $Q$ is full $\GL_8$ even after this shrinking.

The joint node permutation image projects onto $S_N$, with kernel a
subgroup of a product of $S_{18}$'s. The latter has no prime divisor
$31$, whereas $\GL_8(\mathbf F_2)$ does. The factors of $S_N$ are
$A_N$ and cyclic groups of order two; $A_N$ is not $\GL_8(\mathbf F_2)$,
since $11$ divides $|A_N|$ for $N\ge18$ and does not divide
$|\GL_8(\mathbf F_2)|$. Thus the joint image has no composition factor
isomorphic to that simple group. The same is true of its graph image
and of the entire $T_0$ image: the latter only adds a subgroup of the
additive two-group $\Hom(T_0/T,T)$. This step does not assume that the
constant graded pieces split as local systems.

Goursat now shows that the kernel of the entire $T_0$ action still
surjects onto $\GL(Q)$. Inside its block parabolic the translation
kernel is $Q^*\otimes T'$ for a subspace $T'\subset T_0$.
The elementary $H^1(\GL_8(\mathbf F_2),Q^*)=0$ of
Lemma~\ref{lem:char2-group-inputs} splits the remaining cocycle modulo
$T'$ after conjugation. The full translation kernel then removes the
remaining $T'$ columns, producing an actual block-diagonal $\GL_8$.
The whole boundary lies in the coherent ordinary open, so its image is
a subgroup of the same group obtained on the dense smooth ordinary
open. This proves the lemma.
\end{proof}

\begin{theorem}[Eventual exact images for all ordinary principally polarized surfaces]
\label{thm:ordinary-ppav-eventual-full}
\status{Proved under stated hypotheses}
Let $(A,\Theta)$ be an ordinary principally polarized abelian surface over
any perfect field $k$ of characteristic two, with a line bundle representing
$\Theta$ over $k$, and embed it by $L=\mathcal O_A(3\Theta)$.
After a finite extension of $k$, there exists an odd polarized isogeny
$\phi:A'\to A$ with multiplier $n$ whose source has degree-one relative
geometric group $\GL_8(\mathbf F_2)$.
For the complete systems $|L^d|$, put $d_1=n+1$.
For the ambient hypersurface systems of this fixed embedding, replace it by
\[
           d_1=\max\{n+1,\operatorname{reg}(\mathcal I_A)-1\}.
\]
For every $d\ge d_1$, put $r=9d^2-1$. If $d$ is odd, the entire geometric
and arithmetic images both equal $\GL_r(\mathbf F_2)$.
If $d$ is even, put $N=r-3$ and let $\Gamma\subset S_3$ be the actual
Galois image permuting the three nonzero characters of $\ker F_A$.
After compatible choices of bases the entire images are
\[
 G_{{\rm geom},d}=
 \left\{\begin{pmatrix}B&T\\0&I_3\end{pmatrix}:
 B\in\GL_N(\mathbf F_2),\ T\in M_{N\times3}(\mathbf F_2)\right\},
\]
\[
 G_{{\rm arith},d}=
 \left\{\begin{pmatrix}B&T\\0&\gamma\end{pmatrix}:
 B\in\GL_N(\mathbf F_2),\ T\in M_{N\times3}(\mathbf F_2),\
 \gamma\in\Gamma\right\},
\]
where $\Gamma$ acts by its three-dimensional permutation matrices.
The representation is irreducible exactly in odd degree.
The integer $n$, and hence the threshold, can be found by a terminating
procedure over finite fields or effectively presented finitely generated
fields of definition; over arbitrary perfect fields the statement gives
existence without presupposing computable coefficients.
\end{theorem}
\begin{proof}
Chai's theorem \cite[Theorem~2, printed p.~477]{Chai1995} states over
any algebraically closed characteristic-$p$ field that the fixed
$\ell$-power polarized Hecke orbit of every ordinary principally
polarized abelian variety is dense in its moduli stack.
Use $g=2$, $p=2$, $\ell=3$. Its orbit meets the nonempty open of
Lemma~\ref{lem:ordinary-ppav-GL8-open}. By the defining form of this orbit
there is an odd polarized isogeny from a source in that open to $A$;
its multiplier is an odd $n$ and its degree is $n^2$. Its finitely many
defining algebraic coefficients lie in a finite extension of $k$.
This auxiliary extension is only used for the geometric subgroup
construction; the arithmetic image below is computed over the original $k$.

For $d\ge n+1\ge2$, Theorem~\ref{thm:ordinary-ppav-Hasse-kernel} gives
an irreducible common kernel $K_d$ and transvection generation.
Lemma~\ref{lem:ordinary-ppav-isogeny-block} places a natural eight-space
$Q$ inside $K_d$, since a natural $\GL_8$ has no nonzero map to a trivial
module. On $K_d$ the same subgroup fixes a complementary summand.
No invariant nondegenerate alternating form can exist: the natural
$\GL_8$ has no invariant bilinear form and no invariant covector, so
both its internal pairing and its pairing with that trivial complement
are zero. Thus $Q$ would lie in the radical. The opened transvection
classification \cite[Theorem~3.6, printed p.~22]{KantorSeress2001} now
forces the image on $K_d$ to be the full $\GL(K_d)$.
For odd $d$ this is already the entire claimed image.

For even $d$, the translation kernel in
$\Hom(V_d/K_d,K_d)$ is $K_d\otimes U$ for a subspace
$U\subset(V_d/K_d)^*$. The natural-module $H^1$ vanishing of
Lemma~\ref{lem:char2-group-inputs} conjugates the remaining extension
to its block form. Its invariant-vector space has dimension
$3-\dim U$. Theorem~\ref{thm:ordinary-ppav-Hasse-kernel} proved that
there is no invariant vector, so $U$ is the whole three-space and
all translations occur. This proves the geometric image.

Proposition~\ref{prop:ordinary-ppav-all-covectors} identifies the entire
geometric covector space together with its Galois action. Consequently
the arithmetic image is contained in the displayed preimage of $\Gamma$.
It contains the full geometric kernel just proved and projects onto
$\Gamma$, because the parameter open is geometrically connected and
its fundamental group surjects onto the absolute Galois group.
It is therefore exactly that preimage. The action on the quotient is
the dual of the covector permutation representation; permutation matrices
identify these two representations. In odd degree the ambient full
linear group leaves no room for a larger arithmetic image.

For the fixed projective embedding, regularity gives
$H^1(\mathbf P^8,\mathcal I_A(d))=0$ when
$d\ge\operatorname{reg}(\mathcal I_A)-1$.
Ambient sections then surject onto $H^0(A,L^d)$; over the nonzero-section
projective open this is an affine bundle, and the pullback Lang cover
has the same geometric and arithmetic image. This avoids a projective
normality assertion in characteristic two.

Finally enumerate all polarized three-power isogenies from the given
surface, by its finite etale three-power torsion and maximal isotropic
subgroups, and reverse them by polarized duals. Compute the degree-one
geometric Lang frame image of each source by the terminating exact
algorithms proved earlier, dovetailing the tests. Hecke density and the
nonempty full-image open ensure that some source passes. Its retained
multiplier is the required $n$. Regularity of the fixed homogeneous ideal
is computable by a finite graded resolution. All field extensions,
finite subgroup schemes and quotients here are described by finite
algebraic data over an effective field of definition. No undecided
Galois-group oracle or unproved generic-group guess enters this procedure.
The odd full module is irreducible, and in even degree its nonzero proper
common kernel of codimension three is invariant, proving the last
representation-theoretic assertion.
\end{proof}

\begin{remark}[Scope of this structural extension]\label{auto:A_monodromy:20bf13b82135}
\status{Proved under stated hypotheses}

The theorem treats every ordinary principally polarized abelian surface in
characteristic two, including indecomposable surfaces, after a computable
threshold. It does not identify degree-one groups on every such surface,
and it does not assert a structural full-image formula for arbitrary
ambient smooth projective varieties. Those are separate refinements from
the exact finite algorithms and the general residue theorem already proved.
\end{remark}

\begin{remark}[Finite arithmetic and subgroup checks]\label{auto:A_monodromy:03a3e68c4a1c}
\status{Proved under stated hypotheses}

The retained exact program
\path{agents/monodromy/code/ordinary_ppav_group_checks.py}
checks all six $\GL_2(\mathbf F_2)$ character actions on
$\mathbf F_2^4/\mathbf F_2\mathbf1$, their identification with the
three-dimensional permutation representation, and the prime-$31$ and
prime-$11$ separations used for the actual natural block. Its output is
\path{agents/monodromy/outputs/ordinary_ppav_group_checks.json}.
These finite checks support the arithmetic and group calculations; the
geometric incidence arguments and Hecke-density input are proved or cited
separately above.
\end{remark}


% /root/monodromy, 2026-11-19. Proved; independent root and reduction_descent proof/code audits passed.
\subsection{The full group in the first supersingular degree with a nilpotent part}
\label{subsec:ss-product-degree-two-full}
\status{Proved and checked. The parent and reduction/descent agent
independently audited the complete group argument, reread its group-theoretic
dependencies, and reran the separately reconstructed exact matrix verifier.
Both original discovery and independent verification outputs are retained.}

\begin{theorem}\label{thm:ss-product-degree-two-full-group}
\status{Proved under stated hypotheses}
Let $E_0:y^2+y=x^3$ in characteristic two and let
$A=E_0^2$ with $L=\mathcal O(3O)\boxtimes\mathcal O(3O)$ and its Segre
embedding. In degree two the relative coherent bundle has rank $35$,
whereas its generic stable rank is $33$. Its entire invariant-covector
space is the line generated by the covector of
Theorem~\ref{thm:ss-product-exact-covectors}.
The entire geometric and arithmetic images, over every perfect extension
of $\mathbf F_2$, coincide and are
\[
 \left\{\begin{pmatrix}B&t\\0&1\end{pmatrix}:
 B\in\GL_{32}(\mathbf F_2),\quad t\in\mathbf F_2^{32}\right\}
       =\mathbf F_2^{32}\rtimes\GL_{32}(\mathbf F_2).
\]
Thus the group is the full stabilizer of the one invariant covector,
even though two dimensions of the relative coherent bundle are generically
Frobenius-nilpotent.
\end{theorem}
\begin{proof}
\par\smallskip\noindent\textbf{The coherent and stable conventions.}
Put $N=L^2$, $W=H^0(A,N)$, and $P=\mathbf P(W)$.
Use the ordered one-factor basis
\[
             (1,x,x^2,x^3,y,xy)
\]
and its tensor basis, with index $6i+j$.
A nonnegative integer encodes a section by its binary coefficients in
these $36$ basis positions; position zero is $1\otimes1$.
With $\omega=dx_1\wedge dx_2$, let $C(s)$ denote the root-removed Cartier
multiplication matrix on $W$. Thus the full coherent Frobenius matrix
on $W^*\otimes\mathcal O_P(-1)$ is $C(s)^t$.
The relative coherent bundle is its rank-$35$ Euler kernel, and the
one-dimensional ambient quotient has zero Frobenius. Consequently the
full and relative stable parts are identical.

The rank theorem \ref{thm:DE-A3-ss-product-rank} gives generic stable rank
$33$. The three universal kernel vectors
$s,\delta_1s,\delta_2s$ for $C(s)$ are generically independent and show
that its linearized rank is at most $33$ everywhere; their independence
can be seen at $s=x_1x_2$. Here the invariant derivation on each factor is
$\delta x=1,\delta y=x^2$, and $\delta^2=0$.
On the maximal stable-rank open $U\subset P$, both $C$ and its Frobenius
iterate have rank $33$. The image of full coherent Frobenius is therefore
a rank-$33$ unit-Frobenius vector bundle there and has a finite etale Lang
frame cover. This is the same local system as the stable part of the
relative Euler kernel. Its restriction to the dense smooth-member open
has the same monodromy image. Thus rational coefficient points of $U$
can be used to obtain Frobenius elements without a separate assertion
that each retained divisor is smooth.

The unique nonzero invariant covector is
\[
                       \lambda_s=\delta_1\delta_2s\in W,
\]
restricted to the stable image. This is the projectively intrinsic
covector proved in Theorem~\ref{thm:ss-product-exact-covectors}; in a
coefficient chart its displayed coordinates satisfy
$C(s)\lambda_s=\lambda_s^{[2]}$.
Over $\mathbf F_2$ this is $C(s)\lambda_s=\lambda_s$.
Let $V$ denote its rank-$33$ Lang fibre and let
$K=\ker\lambda\subset V$, of dimension $32$.
The geometric and arithmetic groups preserve $K$ and act trivially on
$V/K$.

\par\smallskip\noindent\textbf{Two exact Frobenius certificates.}
The matrices are reconstructed from the elliptic identities
\[
 C(x^{2a})=0,\quad C(x^{2a+1})=x^a,\quad
 C(x^{2a}y)=x^{a+1},\quad C(x^{2a+1}y)=x^a y,
\]
using $y^2=y+x^3$ and tensor products. All coefficients here are binary.
At the section
\[
                         s_1=60944792081
\]
one has $\rank C(s_1)=\rank C(s_1)^2=33$ and
\[
 \det(TI-C(s_1))=T^3(T+1)f(T),
\]
where
\begin{align*}
 f(T)={}&T^{32}+T^{29}+T^{24}+T^{22}+T^{20}
       +T^{15}+T^{14}+T^{13}+T^{12}\\
       &+T^{10}+T^8+T^4+T^2+T+1.
\end{align*}
Its binary coefficient integer is $4853921047$.
The polynomial is irreducible and not reciprocal. Its roots have order
divisible by $65537$, and
$\operatorname{ord}_{65537}(2)=32$.
The retained verifier checks primality of $65537$ by trial division up
to $256$, the order by exponents $32,16$, and the root-order assertion by
\[
 T^{2^{32}-1}\equiv1\pmod f,\qquad
 T^{(2^{32}-1)/65537}\not\equiv1\pmod f.
\]
The constant term is one, so the factor $T^3$ accounts exactly for the
three full coherent nilpotent dimensions. Removing the fixed quotient
factor $T+1$ shows that $f$ is the characteristic polynomial on $K$.
The Lang arithmetic Frobenius is inverse to coherent Frobenius in a
fixed frame; its inverse is also an element of the image. Thus the
arithmetic image on $K$ contains an element with precisely this
polynomial. This convention does not require replacing an arithmetic
element by one in a different group.

At the second section
\[
                         s_2=1642027529
\]
one again has $\rank C(s_2)=\rank C(s_2)^2=33$.
Here the invariant covector and a literal image witness are
\[
              \lambda_{s_2}=20508,\qquad \mu=34911202640,
\]
and the exact binary matrix equations are
\begin{equation}\label{eq:ss-degree-two-fixed-lift-obstruction}
 (C(s_2)-I)\mu=\lambda_{s_2},\qquad
 \rank(C(s_2)-I)=35.
\end{equation}
For every $v$ fixed by $C(s_2)^t$, the first identity gives
\[
 \langle\lambda_{s_2},v\rangle
 =\langle\mu,(C(s_2)^t-I)v\rangle=0.
\]
A fixed vector automatically lies in the stable image, so this is
exactly the assertion that $\lambda$ vanishes on the fixed space of
this arithmetic Frobenius on $V$. In particular this element has no
fixed lift of the nonzero class in $V/K$.

The discovery program
\path{agents/monodromy/code/supersingular_product_degree_two.py}
retains the two complete matrices and both polynomial and image-witness
checks. The separate verifier
\path{agents/monodromy/code/verify_supersingular_degree_two_group.py}
rebuilds both matrices from the monomial rules without reading discovery
matrices or functions, computes their $36\times36$ characteristic
polynomials over the integers and then reduces modulo two, and verifies
\eqref{eq:ss-degree-two-fixed-lift-obstruction} literally. Its complete
output is
\path{agents/monodromy/outputs/verify_supersingular_degree_two_group.json}.
The original read-only discovery stdout, before the checkpoint freeze
was released, is retained separately and not overwritten by the rerun.

\par\smallskip\noindent\textbf{The full linear image on the kernel.}
Theorem~\ref{thm:ss-uniform-GL8-block} supplies an actual geometric natural
$\GL_8(\mathbf F_2)$ block. Its eight-space is contained in $K$, since
its natural representation has no map to the trivial quotient $V/K$.
Choose a nontrivial transvection in this block and let $H$ be its normal
closure under the arithmetic image on $K$.
The geometric image is normal in the arithmetic image, so $H$ is
contained in the geometric image on $K$. It is nontrivial and generated
by transvections. The first Frobenius certificate and
Lemma~\ref{lem:primitive-prime-transvection-test}, with $N=32$ and
$\ell=65537$, therefore imply
\[
                       H=\GL(K).
\]
In particular the geometric and arithmetic images on $K$ are both
$\GL_{32}(\mathbf F_2)$. The cited lemma already includes the socle and
Clifford multiplicity argument, exclusion of proper wreath products by
the primitive prime, and the opened transvection classification with
its nonreciprocity exclusion. No uniform transvection-generation claim
for the entire supersingular family is used here; $H$ was constructed
as a normal closure of specified geometric transvections.

\par\smallskip\noindent\textbf{The full translation kernel.}
Write $G$ for the geometric image in the full stabilizer of $\lambda$.
Its kernel over $\GL(K)$ is a $\GL(K)$-submodule of the natural additive
module $K$. It is either zero or all of $K$.
If it is all of $K$, then $G$ is already the asserted full stabilizer.
Suppose it were zero. Then $G$ is a complement, the graph of a
one-cocycle $\GL(K)\to K$. The elementary vanishing
$H^1(\GL_{32}(\mathbf F_2),K)=0$ from
Lemma~\ref{lem:char2-group-inputs} conjugates it to the block-diagonal
copy $\GL(K)\oplus1$.

The normalizer of this copy inside the pointwise stabilizer of $\lambda$
is itself. Indeed conjugation by a translation $t\in K$ gives the
translation column $(I-B)t$ for $B\in\GL(K)$; it vanishes for every $B$
only when $t\in K^{\GL(K)}=0$. The other diagonal factor already
normalizes the copy. Since the arithmetic image normalizes $G$, it too
would fix the complementary line, and that line maps isomorphically to
$V/K$. Every arithmetic element would then have a fixed lift of its
nonzero quotient class. This contradicts the second Frobenius witness
\eqref{eq:ss-degree-two-fixed-lift-obstruction}.
The translation kernel is therefore all of $K$.
The arithmetic image is between this full geometric stabilizer and the
same stabilizer, since $\lambda$ is defined over $\mathbf F_2$; hence it
is the identical group.

Finally geometric connectedness of the finite covers realizing this
image survives algebraically closed field extension. The invariant
covector is already defined over $\mathbf F_2$, so the equality of
geometric and arithmetic images persists over every perfect extension.
The ambient degree-two sections surject onto $H^0(A,L^2)$ by the plane
cubic coordinate rings and Segre multiplication already checked in the
rank theorem. This proves the statement for the specified embedding.
\end{proof}


% /root/monodromy, 2026-11-19. Proved; independent root and reduction_descent proof/code audits passed.
\subsection{A further certified full supersingular group}
\label{subsec:ss-product-degree-three-full}
\status{Proved and checked. The parent and reduction/descent agent
independently audited the proof and reran the complete independent
monomial-matrix, integer characteristic-polynomial and finite-field verifier.}

\begin{theorem}\label{thm:ss-product-degree-three-full-group}
\status{Proved under stated hypotheses}
For $E_0:y^2+y=x^3$, $A=E_0^2$ and
$L=\mathcal O(3O)\boxtimes\mathcal O(3O)$ in characteristic two, the
entire geometric and arithmetic relative degree-three images are
$\GL_{80}(\mathbf F_2)$ over every perfect extension of $\mathbf F_2$.
\end{theorem}
\begin{proof}
Use the one-factor basis
\[
                  (1,x,x^2,x^3,x^4,y,xy,x^2y,x^3y)
\]
and tensor bit index $9i+j$.
At the binary section
\[
                      s=1726640724722881915971767
\]
the root-removed Cartier multiplication matrix $C(s)$ has
\[
       \rank C(s)=\rank C(s)^2=80,
       \qquad \det(TI-C(s))=T f_{80}(T).
\]
The binary coefficient integer for $f_{80}$ is
\[
                         1489954141958205319760151.
\]
The polynomial is irreducible and nonreciprocal; a root has order
divisible by the prime $\ell=4278255361$, with
$\operatorname{ord}_{\ell}(2)=80$.
The independent verifier reconstructs the $81\times81$ matrix from
$y^2+y=x^3$ and the four elliptic Cartier rules displayed in the
preceding proof, computes its characteristic polynomial over the
integers before reduction modulo two, verifies irreducibility by the
Rabin exponents $16,40,80$, primality of $\ell$ by trial division through
$\lfloor\sqrt\ell\rfloor$, and the root order by polynomial powering.
All section and polynomial coefficients and the complete matrix are
retained in
\path{agents/monodromy/outputs/supersingular_product_degree_three.json}.
The separately reconstructed verifier and its output are
\path{agents/monodromy/code/verify_supersingular_degree_three_group.py}
and
\path{agents/monodromy/outputs/verify_supersingular_degree_three_group.json}.

The relative coherent bundle has rank $80$ and is generically
Frobenius-bijective by Theorem~\ref{thm:DE-A3-ss-product-rank}.
The full coherent bundle has rank $81$ and a zero-Frobenius ambient
quotient. Hence its stable part is exactly the relative system.
The retained rank equalities put this coefficient point in its maximal
open, where the finite Lang frame cover is defined. Restricting that
cover to the dense smooth-member open leaves its image unchanged.
With the transpose convention for coherent Frobenius, and replacing
an arithmetic Frobenius by its inverse if necessary, the arithmetic
image therefore contains an element with characteristic polynomial
$f_{80}$.

Theorem~\ref{thm:ss-uniform-GL8-block} supplies a geometric natural
$\GL_8(\mathbf F_2)$ block, hence a nontrivial geometric transvection.
Take its normal closure $H$ under the arithmetic image. It is nontrivial,
is generated by transvections, and remains inside the geometric image
because that image is normal. Apply
Lemma~\ref{lem:primitive-prime-transvection-test} with $N=80$ and the
specified $\ell$. It gives $H=\GL_{80}(\mathbf F_2)$, so both complete
images are this group. This argument does not assert that every local
inertia subgroup of the supersingular family is itself generated by
transvections.

As before the finite geometrically connected covers remain connected
under algebraically closed field extension. The geometric image is
already the ambient full linear group, so the arithmetic equality also
persists over every perfect extension. The cubic Segre embedding
surjects onto the complete system in every degree, giving the stated
hypersurface-section family.
\end{proof}

\begin{remark}[The remaining uniform supersingular question]\label{auto:A_monodromy:12aa8c66349f}
\status{Proved under stated hypotheses}

The exact full images are now certified in degrees one, two and three.
The proved all-degree rank, invariant-covector and natural-block theorems
do not by themselves settle the entire image in every higher degree.
In particular the nonordinary Hasse divisors can have repeated components,
and the ordinary reduced-irreducible-Hasse argument is not being applied
to them. The additional finite certificates are kept separate from that
uniform structural gate.
\end{remark}


\subsection{Exact generic ranks on every principally polarized abelian surface in characteristic two}
\label{sec:DE-A3-ppav-ranks}

The next result gives a structural formula for this entire class of
surfaces.  In particular the ordinary elliptic factor of a mixed product
may have an arbitrary algebraic $j$-invariant.  The proof does not make
an assertion about a general smooth projective ambient variety.

\begin{lemma}[The four-parameter alternating theta tensor]
\label{lem:DE-A3-theta-pfaffian}
\status{Proved; root and monodromy full independent audits passed.}
Work over an algebraically closed field of characteristic two.  Index a
nine-dimensional space $W$ by the points of $I=\mathbf F_3^2$, and let
$\mathcal D=\mathbf P^1(\mathbf F_3)$ be its four line directions.
For $i\in\mathcal D$ put
\[
 P_i(s)=\sum_{\substack{\ell\subset I\text{ affine line}\\
                         \operatorname{dir}(\ell)=i}}
                           \prod_{a\in\ell}s_a,
 \qquad P_*(s)=\sum_{a\in I}s_a^3.
\]
Let $T$ be the alternating trilinear form whose value on three distinct
basis vectors is $\alpha_i$ when their indices form a line of direction
$i$, and zero otherwise.  Let $M(s)$ be its contraction matrix:
\[
             M(s)_{ab}=T(s,e_b,e_a).
\]
Then $M(s)s=0$, and the vector of its eight-by-eight principal
Pfaffians is $q_\alpha(s)s$, where
\begin{equation}\label{eq:DE-A3-pfaffian-cubic}
 q_\alpha(s)=
   \left(\prod_{i\in\mathcal D}\alpha_i\right)P_*(s)
   +\sum_{i\in\mathcal D}
       \alpha_i\left(\sum_{j\ne i}\alpha_j^3\right)P_i(s).
\end{equation}
Consequently
\[
                \operatorname{adj}M(s)=q_\alpha(s)^2ss^t.
\]
For a nonzero coefficient vector $\alpha$, the cubic $q_\alpha$ vanishes
identically precisely in either of the following cases: exactly one
coefficient is nonzero; or exactly three are nonzero and their cubes are
equal.  These are forty points in the projective coefficient space.
In each case $W$ is a direct sum of three three-dimensional spaces on
which $T$ is the sum of three nonzero volume forms.  In particular it
has a nontrivial projector $D$ with
\[
                         T(Ds,t,u)=T(s,Dt,u).
\]
\end{lemma}
\begin{proof}
All identities may first be checked over the polynomial ring in the
$\alpha_i$ and $s_a$, and then specialized.  Alternation gives $M(s)s=0$.
The Pfaffian kernel identity says that its vector $v$ of principal
Pfaffians also lies in the kernel.  The universal matrix has rank eight:
at $s=e_0$, deleting the zero index leaves four two-by-two alternating
blocks with entries $\alpha_i$, and their Pfaffian is
$\prod_i\alpha_i$.  Thus over the fraction field $v=q s$.
Since the coordinates $s_a$ are pairwise coprime, the identities
$v_as_b=v_bs_a$ imply that each $s_a$ divides $v_a$ in the polynomial
ring.  Hence $q$ is a polynomial, cubic in $s$ and quartic in $\alpha$.

The finite level-three Heisenberg group preserves $T$.  Its generators
are translations of the index set, the diagonal modulations
$e_a\mapsto\zeta^{y\cdot a}e_a$, and scalar multiplication by
$\zeta$, where $\zeta^3=1$ and $\zeta\ne1$.  Each generator has
determinant one.  The transformation rule for principal Pfaffians of
an odd alternating matrix, or the identification
$\bigwedge^8W^*\simeq W\otimes\det(W)^*$, shows that $v$ transforms
as $s$ does.  Therefore $q$ is invariant under this group.
Modulation invariance permits only cubic monomials whose three indices,
with repetition, sum to zero.  Such a monomial is either $s_a^3$ or the
product of the three points on an affine line: if two indices coincide,
the third must coincide as well.  Translation invariance then shows
that the five polynomials $P_*,P_i$ are a basis for these invariant
cubics.

At $s=e_0$ the preceding four-block calculation gives the coefficient
$\prod_i\alpha_i$ of $P_*$.  To determine a line coefficient take the
vertical line $x=0$ and specialize its three section coordinates to
$x_0,x_1,x_2$, setting the other six to zero.  Deleting the index
$(0,0)$ leaves its two vertical companions as one block, with entry
$\alpha_\infty x_0$, and a six-by-six block with Pfaffian $\det B$, where
\[
 B=\begin{pmatrix}
 \alpha_0x_0&\alpha_1x_2&\alpha_2x_1\\
 \alpha_2x_2&\alpha_0x_1&\alpha_1x_0\\
 \alpha_1x_1&\alpha_2x_0&\alpha_0x_2
 \end{pmatrix}.
\]
Its six-term determinant is
\[
 \det B=\alpha_0\alpha_1\alpha_2(x_0^3+x_1^3+x_2^3)
       +(\alpha_0^3+\alpha_1^3+\alpha_2^3)x_0x_1x_2.
\]
Dividing the principal Pfaffian by $x_0$ proves the coefficient of
$P_\infty$.  Linear changes of the index coordinates permute the four
directions, proving every coefficient in
\eqref{eq:DE-A3-pfaffian-cubic}.  The adjugate formula is the usual
odd alternating identity $\operatorname{adj}M=vv^t$, also a polynomial
identity, so it holds at every specialization.

The five orbit sums have disjoint monomial supports.  Thus $q=0$ is
equivalent to
\[
 \prod_i\alpha_i=0,\qquad
 \alpha_i\sum_{j\ne i}\alpha_j^3=0\quad(i\in\mathcal D).
\]
If $r$ coefficients are nonzero, the last equations say that all their
cubes equal the same nonzero number $u$ and that $ru=u$ in the
characteristic-two field.  Therefore $r$ is odd, and the product equation
leaves $r=1$ or $r=3$.  Projectively the first type gives four points;
the second gives four choices of the missing direction and $3^2$ choices
of the two ratios in $\mu_3$, hence thirty-six points.

For $r=1$ the three parallel affine lines partition the basis into
three orthogonal three-spaces, and $T$ is a nonzero volume form on each.
For $r=3$, scale $T$ so that its three nonzero coefficients belong to
$\mu_3$.  A diagonal basis change
\[
                  e_a\longmapsto\zeta^{Q(a)}e_a
\]
for a homogeneous quadratic form $Q$ on $\mathbf F_3^2$ multiplies the
coefficient in direction $v$ by $\zeta^{-Q(v)}$.  Indeed
$\sum_{t\in\mathbf F_3}Q(a+tv)=-Q(v)$.
Evaluations of binary quadratics on any three distinct projective
directions are independent: a nonzero binary quadratic cannot vanish at
three projective points.  One can therefore make the three coefficients
equal to one.  A linear permutation of the index set makes the missing
direction vertical.

Here is the remaining change of basis explicitly.  Start from the
one-direction form that is the sum of volume forms on the three
horizontal lines.  Apply Fourier change of basis in the second index:
\[
              f_{x,u}=\sum_{y\in\mathbf F_3}\zeta^{uy}e_{x,y}.
\]
This matrix is invertible; its inverse has entries $\zeta^{-uy}$,
because $3=1$ in the coefficient field.  Its transformed tensor is
nonzero exactly when the three first indices are distinct and the
three second indices sum to zero.  The coefficient is
$\sum_y\zeta^{y(u_1+u_2+u_3)}$, which equals one or zero accordingly.
These triples are exactly the affine lines in the three nonvertical
directions.  Thus the normalized three-direction form is equivalent
to the one-direction form.  Its three orthogonal projectors transport
back to those asserted in the lemma.  The displayed self-adjointness
identity is immediate from orthogonality of the three summands.
\end{proof}

\begin{proposition}[Every degree-one relative Frobenius is generically bijective]
\label{prop:DE-A3-ppav-degree-one}
\status{Proved; root and monodromy full independent audits passed.}
Let $(A,\Theta)$ be any principally polarized abelian surface in
characteristic two and put $L=\OO_A(3\Theta)$.  The relative coherent
Frobenius for the geometric generic smooth member of $|L|$ is bijective
on its eight-dimensional space.  If $f_A$ is the $2$-rank of $A$, that
curve has genus ten and $2$-rank $8+f_A$.
\end{proposition}
\begin{proof}
Work over an algebraic closure and choose a symmetric theta
representative.  Any other representative is a translate, giving an
isomorphic geometric family.  The typed theta construction of
Theorem~\ref{thm:ppav-level3-Cartier-symmetric} applies to every $A$,
with relative Cartier pairing
\[
 \beta:W\otimes W\longrightarrow W^{(2)},\qquad
 W=H^0(A,L),\quad\dim W=9.
\]
Its theta basis identifies the coefficient-square representation with
the contragredient of the finite Heisenberg representation; the resulting
trilinear tensor $T$ is fully symmetric and has diagonal coefficient
given by Cartier on an invariant top form.  All Frobenius twists and
the restriction to the finite central $\mu_3$ are as in that theorem.

If $A$ is ordinary, the independent symmetric diagonal gives a nonzero
full Hasse determinant by that theorem.  Its relative rank-eight
subbundle is consequently generically bijective as well.  We now
consider the other case.  Cartier on the invariant two-form is zero,
so $T(e_a,e_a,e_b)=0$ for every pair of basis indices.  Thus $T$ is
alternating.  Modulation invariance and index translations put it
exactly in the four-parameter form of
Lemma~\ref{lem:DE-A3-theta-pfaffian}.

We show that its cubic $q_\alpha$ cannot vanish.  Relative Frobenius
$F:A\to A^{(2)}$ is finite flat of degree four, and
\[
                 \mathcal F=F_*L
\]
has a perfect symmetric pairing
$\mathcal F\otimes\mathcal F\to L^{(2)}$ whose global pairing is
$\beta$.  This is finite Frobenius duality, locally the perfect
coefficient-of-$xy$ pairing in a two-element $2$-basis, with the
invariant two-form supplying the duality trivialization; see
\cite[Lemma~1.3.6, Proposition~1.3.7 and Theorem~1.3.8,
pp.~26--28]{BrionKumar2005}.

The bundle $\mathcal F$ is globally generated.  In the smooth-theta
case, Theorem~\ref{thm:ppav-three-theta-jets} gives second jets of
$3\Theta$ at every point.  A Frobenius fibre has local algebra
$k[x,y]/(x^2,y^2)$, which is a quotient of
$k[x,y]/(x,y)^3$, so these jets give precisely the required fibre
evaluation surjection.  In the product case the two degree-three
elliptic bundles evaluate surjectively on every length-two Frobenius
fibre, since the residual line bundle has degree one and has vanishing
$H^1$; tensor their evaluation maps.  These two cases exhaust the
principal surfaces by Lemma~\ref{lem:ppav-surface-dichotomy}.
Perfectness and global generation also show that $\beta$, and hence
the coefficient vector $\alpha$, is not zero.

Suppose $q_\alpha=0$.  The lemma supplies a nontrivial orthogonal
decomposition of $W$, and let $D$ project onto one three-dimensional
summand.  It preserves the kernel of every evaluation map to a fibre
of $\mathcal F$.  Indeed write a global section vanishing in that fibre
as $a+b$, with $a$ in that summand and $b$ in its orthogonal complement.
The value of $a$ pairs to zero with every section in its complement
by orthogonality.  It pairs to zero with every section in its own
summand because $a+b$ has zero value and $b$ is orthogonal to that
summand.  Global generation and perfectness therefore force the value
of $a$ to be zero, and then the value of $b$ is zero too.
The kernel of the surjective bundle evaluation
$W\otimes\OO_{A^{(2)}}\to\mathcal F$ is locally free.  The induced
map from that kernel to $\mathcal F$ vanishes in every closed fibre,
and hence vanishes identically.  Thus $D$ descends to an endomorphism
of $\mathcal F$, acting on global sections as the original rank-three
projector.

This contradicts Oda's simplicity theorem
\cite[Theorem~1.2(i), p.~45, with its preceding proof
pp.~43--45]{Oda1971}.  That theorem concerns isogenies of abelian
varieties of arbitrary dimension, including inseparable isogenies.
Its hypothesis here is satisfied because
$\lambda_L=3\lambda_\Theta$ restricts to an immersion on
$\ker F\subset A[2]$: $\lambda_\Theta$ is an isomorphism and three
is invertible on $A[2]$.  Therefore
$\End(F_*L)=k$.  A scalar endomorphism cannot induce that nontrivial
projector.  We have proved $q_\alpha\ne0$.

For clarity, the relative determinant, not just the full linearized
rank, is determined next.  The full coherent bundle on
$\mathbf P(W)$ is $W^*\otimes\OO(-1)$, and the relative subbundle
is $\Omega^1_{\mathbf P(W)}=\operatorname{ann}(s)$ under the Euler
sequence.  Its Frobenius-domain fibre is
$\operatorname{ann}(s^{[2]})$.  The contraction matrix is symmetric
and maps that domain into $\operatorname{ann}(s)$, since $M(s)s=0$.
On the affine chart $s_0=1$, use domain basis
$e_j+s_j^2e_0$ and target basis $e_i+s_ie_0$, $1\le i,j\le8$.
The resulting matrix has entries $M_{ij}+s_j^2M_{i0}$.
Expanding its determinant by these rank-one column additions and
using $\operatorname{adj}M=q_\alpha^2ss^t$ gives
\begin{equation}\label{eq:DE-A3-relative-theta-Hasse}
                 h_{\mathrm{rel}}(s)
                    =q_\alpha(s)^2\sum_{a\in I}s_a^3.
\end{equation}
Both factors are nonzero polynomials.  This is an actual determinant
on the rank-eight relative bundle, so its nonvanishing proves
bijectivity and stable rank eight.  Intersect this nonempty open with
the smooth-member open of the very ample system $|L|$.
Finally $g(Y)=1+L^2/2=10$, and exactness of stable parts in
\[
 0\longrightarrow H^1(A,\OO_A)\longrightarrow H^1(Y,\OO_Y)
   \longrightarrow E(Y,\OO_Y)\longrightarrow0
\]
gives $f_Y=f_A+8$.
\end{proof}

\begin{lemma}[Odd polarized norms preserve the curve defect]
\label{lem:DE-A3-ppav-odd-norm}
\status{Proved; root and monodromy full independent audits passed.}
For any principally polarized abelian surface $(A,\Theta)$ in
characteristic two and every odd $d\ge1$, the system $|3d\Theta|$
contains an integral nodal curve whose normalization is a smooth
degree-one theta-section of a principally polarized surface odd-isogenous
to $A$.  It can be chosen to have arithmetic-genus-minus-$2$-rank
defect exactly $2-f_A$.
\end{lemma}
\begin{proof}
Work geometrically.  Choose a maximal isotropic direct summand
in $A[d]\simeq(\mathbf Z/d\mathbf Z)^4$.  The quotient
$\psi:A\to A'$ has degree $d^2$, and $d\Theta$ descends to a
principal polarization $\Theta'$ on $A'$.
For the precise descent and principal-degree assertions see
\cite[Corollary~8.11, p.~117, and Proposition~8.19,
p.~119]{EdixhovenVanDerGeerMoonenAV}.
The polarized dual
\[
 \phi=\lambda_\Theta^{-1}\psi^\vee\lambda_{\Theta'}:A'\to A
\]
satisfies $\phi\psi=[d]$, $\psi\phi=[d]$,
$\deg\phi=d^2$ and $\phi^*\lambda_\Theta=d\lambda_{\Theta'}$.
It is finite \'etale since $d$ is odd.  Both surfaces have the
same $2$-rank $f_A$.

Let $L'=\OO_{A'}(3\Theta')$.  Proposition~\ref{prop:DE-A3-ppav-degree-one}
gives a nonempty open of smooth $C'\in|L'|$ with genus ten and
$2$-rank $8+f_A$.  We verify that a general such curve has nodal image
under $\phi$ without an ordinarity hypothesis on either surface.
A nonzero kernel translation preserving $L'$ acts nontrivially on
its projective space of sections, because $L'$ is very ample.
Its fixed locus is proper.  A translation not preserving $L'$ cannot
preserve a member.  Thus a general $C'$ has no nontrivial kernel
stabilizer, and $C'\to\phi(C')$ is generically one-to-one.

If $\Theta'$ is smooth, for a fixed nonzero kernel point $g$ the
rank-six failure of the two first jets at $P,P+g$ is contained in
\[
                   3P+g\in\Theta',\qquad
                   3P+2g\in\Theta'.
\]
The two theta translates are distinct, so their intersection is finite;
the map $[3]$ is finite.  Consequently only finitely many $P$ are
exceptional.  At every other $P$, two values and proportionality of
two nonzero gradients impose three independent conditions on sections,
while $P$ varies in dimension two.  At each exceptional point, the two
independent value conditions already give a proper locus.  This excludes
tangential double identifications for a general member.  The three-value
Lemma~\ref{lem:ppav-three-distinct-values} excludes triple
identifications: three distinct points $P,P+g,P+h$ impose three
conditions, with $P$ varying in dimension two.  There are finitely
many kernel elements and pairs.

For a product of elliptic curves the same assertions follow directly
from the cubic Segre embedding.  When both coordinates of $g$ are
nonzero, the crossed tangent exceptions are finite, by the cubic chord
law and finiteness of $[3]$.  When one coordinate is zero, the two first
jets have rank five.  Their only relation involves the two values and
the two derivatives in that factor; the transverse derivatives are
independent.  After setting both values to zero, at least one parallel
derivative can be nonzero: otherwise a nonzero section of a degree-three
elliptic bundle would vanish twice at two distinct points.  Proportionality
of the gradients is therefore one further independent equation.
Three distinct rank-one Segre tensors are collinear only if one factor
is constant and the other three points are collinear.  For three kernel
translates the latter condition is $3P_i+g_i+h_i=0$, a
one-dimensional locus on the surface.  Its value rank is two, and
elsewhere the rank is three.  The corresponding incidences are both
proper: their dimensions in the projective eight-space are at most
$8-2+1=7$ and $8-3+2=7$.  These arguments use only degree-three elliptic
bundles and the chord law, and so include supersingular factors.

Thus $Z=\phi(C')$ is integral nodal, with normalization $C'$.
For each nonzero kernel element the intersection
$C'\cap(C'-g)$ has degree $L'^2=18$.  No triple identifications
occur, so the number of nodes is $9(d^2-1)$.

The norm line bundle has Neron--Severi class $3d\Theta$:
push forward $\phi^*\Theta\equiv d\Theta'$ and use
$\deg\phi=d^2$ and torsion-freeness of the Neron--Severi group.
Choose symmetric representatives.  Both
$\operatorname{Nm}_\phi(L')$ and $\OO_A(3d\Theta)$ are symmetric,
so their quotient is a degree-zero two-torsion line bundle.
The homomorphism $3d\lambda_\Theta$ is an isomorphism on the geometric
two-torsion points.  Translation by one constant two-torsion point
therefore makes the image lie in the exact linear system $|3d\Theta|$.
This translation changes neither the curve nor its Frobenius cohomology.

The normalization sequence of an integral nodal curve adds a
Frobenius-bijective graph term of dimension its number of nodes.
It therefore preserves genus-minus-stable-rank defect.  The chosen
normalization has defect $10-(8+f_A)=2-f_A$, which proves the lemma.
For $d=1$ the isogeny is the identity and there are no nodes.
\end{proof}

\begin{lemma}[The connected Prym gives the even-degree bound]
\label{lem:DE-A3-ppav-even-bound}
\status{Proved; root and monodromy full independent audits passed.}
Let $Y\in|3d\Theta|$ be smooth on a principally polarized abelian
surface in characteristic two.  If $d$ is even, then
\[
                         g(Y)-f_Y\ge2(2-f_A).
\]
\end{lemma}
\begin{proof}
The restriction $j:\Pic^0(A)\to J_Y$ is a closed immersion.
For a nontrivial geometric degree-zero line bundle $P$ on $A$, one has
$H^0(P)=0$ and $H^1(P(-Y))=0$, the second by duality and ample
vanishing on an abelian variety.  The restriction sequence thus gives
$H^0(Y,P|_Y)=0$, so the geometric point kernel is trivial.
The differential of $j$ is injective because $H^1(A,\OO_A(-Y))=0$.
Its kernel is consequently a finite group scheme with zero tangent
space and no nontrivial geometric points, hence is trivial.

Identify $A$ with its dual through $\lambda_\Theta$, and let
$q:J_Y\to A$ be the polarized dual of this immersion.  It is
surjective and has smooth connected kernel $P_Y$.  One can see the
connectedness without assuming that an arbitrary isogeny has smooth
kernel: quotient $J_Y$ by the reduced identity component of $\ker q$.
The remaining map to $A$ is an isogeny whose dual kernel lies in
$\ker j=0$, so it has degree one by isogeny duality
\cite[Theorem~7.5, p.~100]{EdixhovenVanDerGeerMoonenAV}.

The induced-polarization formula is
\begin{equation}\label{eq:DE-A3-general-Prym-multiplier}
                       j^*\lambda_{J_Y}=(3d)\lambda_\Theta,
                       \qquad qj=[3d].
\end{equation}
This formula includes even $3d$; it is not a prime-to-characteristic
claim.  The exact norm/polarization calculation is
\cite[Section~3, equations~(3.2)--(3.3), pp.~346--347]{Debarre1994Degrees}:
for a curve numerically equivalent to $m\Theta$ on a principal surface,
its norm endomorphism is $[m]$ and its induced Jacobian polarization is
$m\Theta$.  Equivalently, restriction of the Jacobian pairing on
prime-to-two Tate modules is cup product with $[Y]=3d[\Theta]$;
faithfulness of that Tate realization identifies the homomorphisms in
\eqref{eq:DE-A3-general-Prym-multiplier}.  The polarized-dual convention
in this display avoids changing the sign when one identifies a Picard
variety with an Albanese variety.

If $d$ is even, $j$ restricts to an immersion of finite group schemes
$A[2]\hookrightarrow P_Y[2]$.  Twist this finite-group immersion by
Frobenius and apply functoriality of Verschiebung.  The identity
component of $\ker V_A$, of order $2^{2-f_A}$, embeds in the identity
component of $\ker V_{P_Y}$.  For any abelian variety $B$ the latter
order is $2^{\dim B-f_B}$, by the degree and etale-part formulas for
Verschiebung; see \cite[Propositions~5.20--5.22, pp.~80--81]
{EdixhovenVanDerGeerMoonenAV}.  Hence
\[
              \dim P_Y-f_{P_Y}\ge2-f_A.
\]
Finally $\dim J_Y=\dim P_Y+2$ and $f_Y=f_{P_Y}+f_A$.
The last equality follows also from Poincare reducibility and
isogeny invariance of $2$-rank.  Adding the two defects proves
$g(Y)-f_Y\ge2(2-f_A)$.
\end{proof}

\begin{theorem}[The universal parity formula for principal surfaces]
\label{thm:DE-A3-ppav-ranks}
\status{Proved; root and monodromy full independent audits passed.}
Let $(A,\Theta)$ be a principally polarized abelian surface over a
perfect field of characteristic two, with a representing line bundle,
and put $L=\OO_A(3\Theta)$.  Write $f_A\in\{0,1,2\}$ for its
geometric $2$-rank.  For every $d\ge1$, the geometric generic relative
stable rank in the complete system $|L^d|$ is
\begin{equation}\label{eq:DE-A3-ppav-rank-formula}
 \boxed{\qquad
 r_d=\begin{cases}
           9d^2-1,&d\text{ odd},\\
           9d^2-3+f_A,&d\text{ even}.
       \end{cases}\qquad}
\end{equation}
The relative coherent dimension is $9d^2-1$.  The full generic curve
defect is $2-f_A$ in odd degree and $2(2-f_A)$ in even degree.
Thus relative Frobenius is generically bijective exactly when $d$ is
odd or $A$ is ordinary.

For the fixed embedding by $L$ into $\mathbf P^8$, the same formula
holds for ambient degree-$d$ hypersurfaces whenever
\[
             d\ge\max\{1,\operatorname{reg}(\mathcal I_A)-1\}.
\]
For the cubic Segre embedding of a product of elliptic curves, it holds
for every $d\ge1$, with no restriction on their geometric
$j$-invariants.
\end{theorem}
\begin{proof}
Over an algebraic closure ample vanishing gives
$H^1(A,L^{-d})=0$, $h^2(A,L^{-d})=9d^2$ and $h^2(A,\OO_A)=1$.
The multiplication map on the last two groups is surjective, Serre-dual
to the inclusion of the defining nonzero section.  Its kernel is the
relative coherent quotient, of dimension $9d^2-1$.  This also constructs
a vector bundle of that rank on the whole coefficient projective space,
including the nodal boundary.  The full $H^1(\OO_Y)$ has constant
dimension $9d^2+1$.  Frobenius and all its iterates are bundle maps after
the appropriate Frobenius base changes, so lower bounds on stable rank
are open, by minors of an iterate at least as long as the bundle rank.

For odd $d$, Lemma~\ref{lem:DE-A3-ppav-odd-norm} supplies a nodal
member of full stable dimension
$9d^2+1-(2-f_A)=9d^2-1+f_A$.
Exactness of stable parts subtracts the ambient rank $f_A$, so the
relative Frobenius at that member is already bijective on all
$9d^2-1$ dimensions.  Its nonempty invertibility open meets the
nonempty smooth-member open.  This proves the odd formula.

For even $d\ge2$, choose generic sharp smooth odd-degree members
$C\in|L|$ and $D\in|L^{d-1}|$ meeting transversely.  They exist by
the odd case, and transversality is a nonempty open condition in the
two very ample systems.  Their connected nodal union has full defect
$2(2-f_A)$, since normalization adds only a Frobenius-bijective graph
term.  Its relative stable rank is therefore $9d^2-3+f_A$.
The open rank locus meets the smooth locus, giving that lower bound
generically.  Lemma~\ref{lem:DE-A3-ppav-even-bound} gives the matching
upper bound at every smooth member.  This proves
\eqref{eq:DE-A3-ppav-rank-formula} in every degree.

Regularity gives $H^1(\mathbf P^8,\mathcal I_A(d))=0$ in the displayed
range, so ambient sections surject onto the complete linear system.
The associated parameter map over the nonzero-section locus is an
affine bundle and has dense image, preserving the geometric generic
rank.  On a product, restriction of plane polynomials to either cubic
surjects in every positive degree, by the cubic hypersurface exact
sequence.  Their tensor products are restrictions of degree-$d$ Segre
polynomials, so the ambient map is surjective in every degree.
All the statements are geometric and are therefore preserved by
extending the original perfect field to its algebraic closure.
\end{proof}

\begin{remark}[Exact certificate and scope]
\label{rem:DE-A3-ppav-rank-checks}
\status{Proved under stated hypotheses}
The retained program
\path{agents/obstructions_crystals/code/A3_ppav_rank_checks.py}
checks the symbolic Pfaffian identity in all nine coordinates, the
forty coefficient points and their orthogonal-projector description,
the Fourier and quadratic-phase identities, and the exact isogeny,
node and parity dimension formulas.  The author run passed in
$0.27075901300122496$ seconds: nine symbolic Pfaffians expanded from
$945$ perfect matchings, all forty rank-three orthogonal projectors,
$1700$ supplementary relative-determinant values, and $30000$ dimension
identities.  Its complete output is retained alongside the program.
Root independently read the entire code and reran it in
$0.29324621799969464$ seconds, preserving the author output and saving
the separate result in \path{outputs/root_A3_ppav_rank_checks.json}.
The analytic proof above supplies the classification
over every algebraically closed characteristic-two field; a finite-field
sample is not used in place of that proof.

The earlier supersingular-product and generic-mixed formulas are
special cases.  In particular the algebraic ordinary-factor rank
refinement left explicit after Theorem~\ref{thm:DE-A3-mixed-generic}
is now covered by this theorem.
No structural rank formula for arbitrary smooth projective ambient
varieties is inferred, and this theorem alone does not determine the
full monodromy groups of nonordinary principal surfaces.
\end{remark}


\subsection{The parity formula for every separable surface polarization}
\label{sec:DE-A3-separable-polarizations}

\begin{lemma}[A principal source with a prescribed norm class]
\label{lem:DE-A3-principal-norm-source}
\status{Proved; root and monodromy full independent audits passed.}
Let $A$ be an abelian surface over an algebraically closed field of
characteristic two, and let $M$ be an ample line bundle with
$D=\chi(M)$ odd.  There is a principally polarized surface
$(B,\Theta_B)$ and an \'etale isogeny $\phi:B\to A$ of degree $D$
such that
\[
                  \phi^*M\equiv D\Theta_B,
                  \qquad \operatorname{Nm}_\phi(\Theta_B)\equiv M.
\]
Here $\equiv$ denotes equality in the Neron--Severi group.
The surface $B$ has the same $2$-rank as $A$.
\end{lemma}
\begin{proof}
We give the integral construction of the dual polarization, so no
rational inverse is silently assumed to be a line bundle.
Write the separable polarization type of $M$ as $(a,b)$, with
$a\mid b$ and $ab=D$.  Then
$\ker\lambda_M\simeq(\mathbf Z/a\mathbf Z)^2\times
(\mathbf Z/b\mathbf Z)^2$; in particular its exponent divides $D$.
For this type and degree description see
\cite[Section~2, p.~345]{Debarre1994Degrees}.
The commutator pairing of $M^D$, restricted to $\ker\lambda_M$,
is the $D$-th power of that of $M$, and is therefore trivial.
The descent criterion
\cite[Corollary~8.11 and proof, p.~117]
{EdixhovenVanDerGeerMoonenAV} gives a line bundle $N$ on $A^\vee$
with
\[
                         \lambda_M^*N\simeq M^D.
\]
It is ample since its pullback under the finite surjection is ample.
The identity for polarizations of pullbacks and symmetry of
$\lambda_M$ give
\[
             \lambda_M\lambda_N\lambda_M=D\lambda_M,
             \qquad \lambda_N=D\lambda_M^{-1}.
\]
The second equality is an equality of rational homomorphisms whose
left side is now an actual integral polarization.  Intersection
degrees give
$\chi(N)=\chi(M^D)/\deg\lambda_M=D^3/D^2=D$.

Choose a maximal isotropic subgroup of $\ker\lambda_N$.
The quotient $\psi:A^\vee\to C$ has degree $D$ and a principal
polarization $\Theta_C$ with $\psi^*\Theta_C\simeq N$;
this is \cite[Proposition~8.19 and proof, p.~119]
{EdixhovenVanDerGeerMoonenAV}, together with the preceding descent
criterion.  Put $B=C^\vee$, equip it with the dual principal
polarization $\lambda_B=\lambda_C^{-1}$, and take
$\phi=\psi^\vee:B\to A$.
The equality $\phi\lambda_C\psi=D\lambda_M^{-1}$ implies
\[
                 \psi\lambda_M\phi=D\lambda_C^{-1}.
\]
Its left side is the polarization of $\phi^*M$, so
$\phi^*M\equiv D\Theta_B$.
Applying the norm and using $\deg\phi=D$ gives
$D\operatorname{Nm}_\phi(\Theta_B)\equiv DM$.
The Neron--Severi group is torsion-free, so the claimed norm class
follows.  The isogeny has odd degree, hence is \'etale and induces an
isomorphism on the $2$-divisible groups.  This proves the last assertion.
\end{proof}

\begin{lemma}[The surface polarization in the connected Prym]
\label{lem:DE-A3-general-polarization-Prym}
\status{Proved; root and monodromy full independent audits passed.}
With $M,N$ as in the preceding proof, let $Y\in|M^m|$ be smooth.
The restriction $j:A^\vee\hookrightarrow J_Y$ is an immersion of
abelian varieties, its polarized dual has smooth connected kernel,
and
\[
                   j^*\lambda_{J_Y}=m\lambda_N.
\]
If $m$ is even, then $g(Y)-f_Y\ge2(2-f_A)$.
\end{lemma}
\begin{proof}
Ample vanishing and its Serre-dual form show that
$H^1(A,P\otimes M^{-m})=0$ for every $P\in\Pic^0(A)$.
As in Lemma~\ref{lem:DE-A3-ppav-even-bound}, this proves that the
point kernel and tangent kernel of $j$ are both zero, hence $j$ is
an immersion.  Quotienting the polarized dual by its reduced identity
kernel, and applying isogeny duality, proves that its kernel is
smooth and connected.

The polarization formula is the general, nonprincipal instance of
\cite[Section~3, equations~(3.2)--(3.3), pp.~346--347]
{Debarre1994Degrees}.  Explicitly, in dimension two these equations
give
\[
 \alpha(Y,M)=mD\operatorname{id}_A,
 \qquad
 \lambda_M\bigl(j^*\lambda_{J_Y}\bigr)\lambda_M
       =mD\lambda_M.
\]
Cancelling the isogeny in the rational homomorphism space gives
$j^*\lambda_{J_Y}=mD\lambda_M^{-1}=m\lambda_N$.
Both sides are integral homomorphisms, so this is an integral equality,
including even $m$.
For even $m$, the polarized-dual composite therefore kills
$A^\vee[2]$.  This finite group scheme embeds in the connected Prym.
Apply Frobenius twist and functorial Verschiebung to that finite-group
immersion exactly as in Lemma~\ref{lem:DE-A3-ppav-even-bound}.
The connected Verschiebung kernel has order $2^{2-f_A}$, forcing
a Prym defect at least $2-f_A$.  Adding the ambient defect proves
the stated bound.
\end{proof}

\begin{theorem}[All degrees for a separable polarization]
\label{thm:DE-A3-separable-ranks}
\status{Proved; root and monodromy full independent audits passed.}
Let $A$ be an abelian surface over a perfect field of characteristic
two, and let $M$ be an ample line bundle of odd Euler characteristic
$D=\chi(M)$.  Put $L=M^3$.  For every $d\ge1$, the geometric generic
relative stable rank in the complete linear system $|L^d|$ is
\[
 r_d=\begin{cases}
           9Dd^2-1,&d\text{ odd},\\
           9Dd^2-3+f_A,&d\text{ even}.
       \end{cases}
\]
Its relative coherent dimension is $9Dd^2-1$.  The full generic
curve defect is $2-f_A$ in odd degree and $2(2-f_A)$ in even degree.
For the embedding by $L$ the same formula holds for ambient
hypersurfaces whenever
$d\ge\max\{1,\operatorname{reg}(\mathcal I_A)-1\}$.
\end{theorem}
\begin{proof}
Work over an algebraic closure.  We may replace $M$ by a translate
that is symmetric.  Indeed
$\alpha=[-1]^*M\otimes M^{-1}$ lies in $\Pic^0(A)$; divisibility
of this group gives $P$ with $P^2\simeq\alpha$, and then $M\otimes P$
is symmetric.  Surjectivity of $\lambda_M$ writes this bundle as
$t_a^*M$ for one point $a$.  That same translation identifies every
family $|M^{3d}|$ with $|(t_a^*M)^{3d}|$, so all generic ranks transfer.
We first note directly that $L$ is
very ample in the present construction.  The norm of a very ample line
bundle under a finite \'etale morphism is very ample.  Indeed, to
separate two distinct target points, choose a source section vanishing
at one selected lift of the first and at none of the finitely many
other lifts of either point.  Such a choice exists by very ampleness,
avoiding finitely many proper hyperplanes in the space of sections
vanishing at the selected lift.  Its norm separates the two target
points.  For a nonzero target tangent vector, choose a section with
nonzero derivative along its lifted tangent and nonzero values at the
other lifts.  The other local factors in its norm are units, so the
norm separates that tangent.  The same argument supplies nonvanishing
norm sections at every target point.  This proves the assertion by the
point and tangent criterion for a projective smooth variety.
Apply this to $\OO_B(3\Theta_B)$ in
Lemma~\ref{lem:DE-A3-principal-norm-source}.  Its norm is numerically
equivalent to $M^3$.  Line bundles in the same ample numerical class
on an abelian variety are translates, since their quotient is in
$\Pic^0$ and the polarization is surjective.  Thus $L$ is very ample.

For odd $d$, apply Lemma~\ref{lem:DE-A3-principal-norm-source} to
$M^d$.  The resulting principal source $B_d$ maps \'etale to $A$
with odd degree $Dd^2$, and the norm of its theta divisor has class
$dM$.  Proposition~\ref{prop:DE-A3-ppav-degree-one} supplies a
nonempty open of smooth curves in $|3\Theta_{B_d}|$ of genus ten and
$2$-rank $8+f_A$.
The nodal-image proof of Lemma~\ref{lem:DE-A3-ppav-odd-norm}
uses only that the source is principally polarized and that the
isogeny is \'etale of odd degree; it imposes no condition that the
kernel be the particular subgroup occurring there.
It therefore gives a general such curve an integral nodal image with
exactly $9(Dd^2-1)$ nodes.

Choose symmetric representatives.  The norm of the source bundle and
$M^{3d}$ differ by a degree-zero two-torsion line bundle.
The polarization $3d\lambda_M$ induces an isomorphism on the geometric
two-torsion points, because its degree is odd.  One target
two-torsion translation therefore puts the nodal image in the exact
system $|M^{3d}|$.  This image has full defect $2-f_A$, because its
normalization has that defect and the node graph contributes only a
Frobenius-bijective summand.

Riemann--Roch and ample vanishing give the constant full coherent
dimension $9Dd^2+1$ and relative dimension $9Dd^2-1$.
The odd nodal image already has the full possible relative stable
rank.  Openness of the rank of a sufficiently long Frobenius iterate
and the nonempty smooth-member open prove the odd formula.

For even $d$, smooth a transverse union of sharp smooth members of
$|L|$ and $|L^{d-1}|$.  Normalization adds their two defects, yielding
the lower bound $9Dd^2-3+f_A$.  The matching upper bound for every
smooth member is Lemma~\ref{lem:DE-A3-general-polarization-Prym}
with $m=3d$.  This proves every degree.
Finally, the stated regularity condition makes the ambient restriction
map surjective.  Its geometric parameter fibres are integral, so it
preserves the geometric generic rank.
\end{proof}

\begin{remark}[Scope of the separable extension]
\label{rem:DE-A3-separable-rank-scope}
\status{Proved under stated hypotheses}
This extension allows arbitrary odd polarization types $(a,b)$ on the
surface; a principal polarization on the surface itself is unnecessary.
It still does not assert a formula for arbitrary ambient varieties,
or for an arbitrary polarization whose Euler characteristic is even.
The line bundle defining the displayed embedding is explicitly
$M^3$.  No claim that every given embedding has that form is made.

The retained arithmetic certificate
\path{agents/obstructions_crystals/code/A3_separable_rank_checks.py}
checks the integral dual-polarization matrix and the displayed
degree, node, genus and rank identities for sixty-five odd types
and $24960$ degree/rank cases.  Its author run passed in
$0.024252046998299193$ seconds; the complete output is retained
alongside it.  The geometric existence and descent assertions are
proved above, not inferred from this finite arithmetic check.
Monodromy independently read the complete program and reran it in
$0.015273962999344803$ seconds, retaining a separate result in
\path{agents/monodromy/outputs/independent_A3_separable_rank_checks.json}.
\end{remark}


\subsection{Ordinary surfaces with arbitrary polarization degree}
\label{sec:DE-A3-ordinary-all-polarizations}

\begin{lemma}[The ordinary polarization kernel and its connected isotropic part]
\label{lem:DE-A3-ordinary-polarization-kernel}
\status{Proved; root and reduction-descent full independent proof/source audits passed.}
Let $A$ be an ordinary abelian surface over an algebraically closed
field of characteristic two, and let $M$ be ample with $D=\chi(M)$.
Then $\ker\lambda_M$ is killed by $D$.  If $D_2$ is the largest
power of two dividing $D$, the two-primary part of this kernel is
\[
             K_2=K_2^0\times K_2^{\et},
             \qquad \rank K_2^0=\rank K_2^{\et}=D_2,
\]
where $K_2^0$ is multiplicative and $K_2^{\et}$ is constant.
The commutator pairing is trivial on either factor and pairs the
factors perfectly.  In particular $K_2^0$ is maximal isotropic
in $K_2$.
\end{lemma}
\begin{proof}
We first justify the group-scheme assertions without a separability
assumption.  Over a perfect field, a finite commutative group scheme
splits into its connected and etale parts; there are no nonzero
homomorphisms between a finite local and an etale group scheme.
These are \cite[Lemma~4.44 and Proposition~4.45, p.~67]
{EdixhovenVanDerGeerMoonenAV}, including their proofs.
Ordinarity, the prime-power torsion count, and Cartier duality give
\[
              A[2^n]\simeq\mu_{2^n}^{2}
                      \times(\mathbf Z/2^n\mathbf Z)^2.
\]
For completeness, the etale part has order $2^{2n}$ by
\cite[Proposition~5.22, pp.~80--81]{EdixhovenVanDerGeerMoonenAV}.
The dual surface is isogenous to $A$, hence is ordinary by that
same proposition.  The Cartier dual of its etale torsion is a
multiplicative subgroup of $A[2^n]$ of order $2^{2n}$, by
\cite[Corollary~7.17, p.~104]{EdixhovenVanDerGeerMoonenAV}.
This is the entire connected part, since $\deg[2^n]=2^{4n}$.
Thus every finite two-primary subgroup of $A$ has multiplicative
connected part and constant etale part.

Put $K=\ker\lambda_M$.  Its commutator pairing is perfect and
alternating, and $\rank K=D^2$; see
\cite[Proposition~8.19 and Corollary~8.20, p.~119]
{EdixhovenVanDerGeerMoonenAV}.
Distinct primary components are orthogonal.  On $K_2$, the pairing
on the connected factor is trivial, because it gives a homomorphism
from a multiplicative local group to its etale Cartier dual.
The pairing on the etale factor is likewise trivial.  Perfection
therefore identifies each factor with the Cartier dual of the other.
Their orders are equal and have product $D_2^2$, proving the claimed
orders and the maximal isotropy.  Both factors are killed by $D_2$:
this follows from the elementary exponent bound for a finite abelian
group of order $D_2$, applied also to the character group of the
multiplicative factor.

For an odd prime $\ell$, the $\ell$-primary part of $K$ is a
constant finite abelian group with a perfect alternating pairing.
If $x$ has maximal order $\ell^b$, perfection supplies $y$ such
that $e_M(x,y)$ has that same order.  Pairing a relation
$ux+vy=0$ successively with $x,y$ shows that $\langle x,y\rangle$
has order $\ell^{2b}$.  Consequently
$2b\le v_\ell(\rank K)=2v_\ell(D)$.
The exponent of this primary part thus divides the $\ell$-part
of $D$.  Combining all primary parts proves $DK=0$.
\end{proof}

\begin{lemma}[An etale principal norm source in the ordinary case]
\label{lem:DE-A3-ordinary-etale-norm}
\status{Proved; root and reduction-descent full independent proof/source audits passed.}
With $A,M,D$ as above, there is an ordinary principally polarized
surface $(B,\Theta_B)$ and a finite etale isogeny $\phi:B\to A$
of degree $D$ satisfying
\[
                 \phi^*M\equiv D\Theta_B,
                 \qquad \operatorname{Nm}_\phi(\Theta_B)\equiv M.
\]
The integer $D$ may be even.
\end{lemma}
\begin{proof}
The preceding lemma makes the restriction of $e_{M^D}=e_M^D$
to $\ker\lambda_M$ trivial.  The actual group-scheme descent
criterion \cite[Corollary~8.11 and proof, p.~117]
{EdixhovenVanDerGeerMoonenAV} therefore gives an ample line bundle
$N$ on $A^\vee$ such that $\lambda_M^*N\simeq M^D$.
This criterion allows nonreduced finite isotropic subgroup schemes.
As in Lemma~\ref{lem:DE-A3-principal-norm-source}, the pullback
polarization and degree identities give
\[
                       \lambda_N=D\lambda_M^{-1},
                       \qquad \chi(N)=D.
\]
The left side is an integral polarization, constructed by descent.

In the two-primary part of $\ker\lambda_N$, choose its entire
connected factor; it is maximal isotropic by
Lemma~\ref{lem:DE-A3-ordinary-polarization-kernel} applied to
$A^\vee,N$.  In the odd-primary part choose a maximal isotropic
subgroup.  Their product $H$ is maximal isotropic of order $D$.
Its connected part is multiplicative, and its remaining part has
odd order.  Therefore its Cartier dual $H^D$ is etale.
Proposition~8.19 of the same source gives a principal quotient
\[
               \psi:A^\vee\longrightarrow C=A^\vee/H,
               \qquad \psi^*\Theta_C\simeq N.
\]
Take $B=C^\vee$ with its dual principal polarization, and put
$\phi=\psi^\vee$.  Isogeny duality identifies
$\ker\phi=H^D$ \cite[Theorem~7.5, p.~100]
{EdixhovenVanDerGeerMoonenAV}; hence $\phi$ is etale even when
$D$ is even.  The typed polarization identity is
\[
                \psi\lambda_M\phi=D\lambda_C^{-1}.
\]
It gives $\phi^*M\equiv D\Theta_B$.  Applying norm and cancelling
$D$ in the torsion-free Neron--Severi group gives the second equality.
Ordinarity is preserved under isogeny, so $B$ is ordinary.
\end{proof}

\begin{theorem}[All degrees for every polarization on an ordinary surface]
\label{thm:DE-A3-ordinary-all-ranks}
\status{Proved; root and reduction-descent full independent proof/source audits passed.}
Let $A$ be an ordinary abelian surface over a perfect field of
characteristic two.  Let $M$ be any ample line bundle, put
$D=\chi(M)$ and $L=M^3$.  Then $L$ is very ample and, for every
$d\ge1$, the geometric generic smooth member of $|L^d|$ is ordinary,
with genus $9Dd^2+1$.  In particular its relative stable rank equals
its relative coherent dimension:
\[
                              r_d=9Dd^2-1.
\]
For ambient hypersurfaces in the embedding given by $L$, the same
formula holds whenever
\[
               d\ge\max\{1,\operatorname{reg}(\mathcal I_A)-1\}.
\]
There is no restriction on the parity of $D$ or $d$.
\end{theorem}
\begin{proof}
Work over an algebraic closure.  Apply
Lemma~\ref{lem:DE-A3-ordinary-etale-norm} to $M$.
The cubic principal bundle $L_B=\OO_B(3\Theta_B)$ is very ample.
Its norm has numerical class $M^3$ and is very ample by the
point-and-tangent norm-section argument in the proof of
Theorem~\ref{thm:DE-A3-separable-ranks}.  That argument requires
only a finite etale morphism, not that its degree be invertible.
Indeed one selects a section vanishing at one chosen lift, with
nonzero derivative when needed, and avoids the finitely many value
hyperplanes at every other lift.  The remaining factors of the norm
are units.  All ample line bundles in the same numerical class on
an abelian variety are translates, by surjectivity of the polarization.
Thus $L=M^3$ is very ample.

Theorem~\ref{thm:DE-A3-ppav-ranks} for the ordinary principal source
provides a nonempty open of smooth ordinary genus-ten curves
$C\in|L_B|$.  A general such curve has integral nodal image under
$\phi$.  We spell out why even kernel order creates no exception
to the nodal-image proof already used in
Lemma~\ref{lem:DE-A3-ppav-odd-norm}.
The kernel is a finite constant group acting freely by translations.
A nonzero element that preserves $L_B$ has a nontrivial projective
action in its very ample embedding, hence has a proper locus of
invariant sections; no semisimplicity or averaging is used.
A translation not preserving $L_B$ cannot preserve a member.
Thus a general $C$ has trivial kernel stabilizer.

For an indecomposable principal source, the two-first-jet failure
locus for a fixed nonzero kernel point $g$ is contained in
\[
                   3P+g\in\Theta_B,
                   \qquad 3P+2g\in\Theta_B.
\]
These are inverse images under $[3]$ of two distinct theta translates.
Their intersection is finite, including when $g$ has order two.
Away from these finitely many points, the two vanishing values and
proportionality of two nonzero gradients impose three independent
conditions while $P$ varies in dimension two.  At an exceptional
point, the two value conditions already define a proper locus.
The three-value separation lemma for $|3\Theta_B|$ similarly
excludes triple identifications for every pair of distinct kernel
points.  For a decomposable principal source, the cubic Segre
argument in the cited nodal-image proof supplies the same conclusion:
the one-coordinate coincidence has rank-five first jets, and the
three-collinear-point exception has dimension one and value rank two.
Those incidence bounds also depend only on distinctness of points,
not on the order of a translation.  Hence a general chosen ordinary
$C$ has only transverse double identifications and no triples.

Its image $Z$ is integral nodal, with normalization $C$.  Counting
the ordered pairs gives $18(D-1)$, so $Z$ has $9(D-1)$ nodes.
This count is a count of distinct geometric points under a free
etale action; division by two here counts unordered pairs and does
not invert two in the ground field.  The normalization sequence
adds a Frobenius-bijective graph space.  Consequently $Z$ is
ordinary, of arithmetic genus $10+9(D-1)=9D+1$.

The line bundle of $Z$ is $\operatorname{Nm}_\phi L_B$, whose
quotient with $M^3$ is in $\Pic^0(A)$.  Surjectivity of
$\lambda_{M^3}:A\to A^\vee$ supplies a single target translation
placing $Z$ in the exact system $|M^3|$.  This is an arbitrary
geometric translation; no inverse on the two-torsion points is
asserted.  Translation preserves ordinarity and the node structure.

Ample vanishing gives constant coherent dimensions for every divisor
in $|L^d|$, namely $9Dd^2+1$ in degree one on the divisor, and
$9Dd^2-1$ in the relative quotient.  Frobenius and its iterates
are maps of the corresponding bundles on the coefficient space.
Invertibility is open.  The ordinary nodal member just constructed
therefore yields a nonempty ordinary open for $d=1$, which meets
the nonempty smooth-member open.
For general $d$, choose $d$ smooth ordinary members of $|L|$
meeting pairwise transversely with no triple points.  Bertini and
avoidance of finitely many points give such a choice, and their
connected nodal union lies in $|L^d|$.  Its normalization and graph
terms are ordinary, so the union is ordinary.  The same open-rank
argument gives a smooth ordinary generic member of $|L^d|$.
Its ambient stable rank is two, and subtracting that ambient part
gives the displayed relative rank.  The final regularity statement
follows from surjectivity of ambient restriction in the stated range,
whose parameter fibres are integral, as in
Theorem~\ref{thm:DE-A3-separable-ranks}.
\end{proof}

\begin{remark}[Scope of the extension]
\label{rem:DE-A3-ordinary-all-rank-scope}
\status{Proved scope and verification statement.}
This theorem removes the parity restriction on the polarization degree
for ordinary surfaces.  Theorem~\ref{thm:DE-A3-ordinary-arbitrary-embedding}
further gives the exact eventual rank in any chosen ample embedding,
without requiring its line bundle to be a cube.  It does not assert a formula for arbitrary
nonordinary even-degree polarizations, or for arbitrary ambient varieties.
The finite group-scheme argument is essential: using an arbitrary
maximal isotropic subgroup could give a nonetale dual isogeny.

The retained program
\path{agents/obstructions_crystals/code/A3_ordinary_all_rank_checks.py}
checks the integral dual-polarization matrices, connected and etale
order identities, norm-node counts and normalization-graph arithmetic
for $211$ polarization types, including $146$ of even Euler
characteristic, and $27008$ degree cases.  Its author run passed in
$0.016162138999789022$ seconds.  The complete output is retained
alongside the program.  Root independently read the complete program
and reran it in $0.014947520998248365$ seconds, preserving the author
output and retaining its result at
\path{outputs/root_A3_ordinary_all_rank_checks.json}.
This finite arithmetic certificate does not
replace the group-scheme descent or geometric existence proofs.
\end{remark}


\subsection{Ordinary surfaces in an arbitrary ample class}
\label{sec:DE-A3-ordinary-arbitrary-embedding}

\begin{lemma}[Ordinary members of every principal multiple at least three]
\label{lem:DE-A3-ordinary-principal-multiples}
\status{Proved; root, monodromy, and reduction-descent full independent proof audits passed.}
Let $(B,\Theta)$ be an ordinary principally polarized surface over an
algebraically closed field of characteristic two.  For every integer
$n\ge3$, $\OO_B(n\Theta)$ is very ample and a general smooth member
of its complete system is ordinary, of genus $n^2+1$.
\end{lemma}
\begin{proof}
First, very ampleness for $n=3$ was proved above.  For $n>3$,
write $r=n-3$.  Given one or two prescribed points, a general
translate of the effective divisor $r\Theta$ avoids them.  If $R$
is its bundle, then $\OO_B(n\Theta)\otimes R^{-1}$ has numerical
class $3\Theta$, hence is a translate of $\OO_B(3\Theta)$, by
surjectivity of its polarization.  Multiplying its sections by the
chosen section of $R$ transfers point and tangent separation at
these points.  It also supplies nonvanishing sections.  Thus the
point-and-tangent criterion proves very ampleness for every $n\ge3$.

The principal-surface dichotomy is
Lemma~\ref{lem:ppav-surface-dichotomy}.  In the indecomposable case,
$\Theta$ is an ordinary smooth genus-two curve.  For $n=3$ the
ordinary nodal theta union is Lemma~\ref{lem:ppav-ordinary-members}.
For $n\ge4$ choose general translations $a_1,\ldots,a_n$ with
$\sum_i a_i=0$.  Their theta sum belongs to $|n\Theta|$.
Projection of this parameter space onto any three chosen coordinates
is surjective, since the remaining translations can absorb their sum.
General pairs of theta translates meet transversely: the difference
map $\Theta^2\to B$ is generically etale, as shown in the proof of
Lemma~\ref{lem:ppav-ordinary-members} from the two distinct Gauss
directions.  General triples have no common point, since their
incidence has dimension five in a six-dimensional translation space.
Consequently the chosen sum is connected and nodal, with ordinary
components.

In the decomposable case both elliptic factors are ordinary.  Choose
the $a_i$ with sum zero and with their first coordinates pairwise
distinct and their second coordinates pairwise distinct.  This is a
nonempty open for $n\ge3$.  The theta sum is a grid of $n$ vertical
and $n$ horizontal ordinary elliptic curves, with $n^2$ transverse
nodes and no triple points.  It is again ordinary by normalization.

In either case the normalization sequence is an extension of ordinary
component cohomology by a Frobenius-bijective graph space.  The nodal
union is ordinary.  Ample vanishing makes $h^1(\OO_Y)=n^2+1$
constant over the complete system.  Frobenius invertibility is open,
so this ordinary open intersects the nonempty smooth-member open.
\end{proof}

\begin{lemma}[Nodal images for every principal multiple at least three]
\label{lem:DE-A3-all-multiple-nodal-norm}
\status{Proved; root, monodromy, and reduction-descent full independent proof audits passed.}
Let $\phi:B\to A$ be a finite etale isogeny of principally polarized
surface source $(B,\Theta)$ in characteristic two.  For every $n\ge3$,
a general smooth member $C\in|n\Theta|$ has integral nodal image
with normalization $C$ and exactly $n^2(\deg\phi-1)$ nodes.
No parity restriction on $\deg\phi$ is imposed.
\end{lemma}
\begin{proof}
The $n=3$ incidence proof is the proof of
Lemma~\ref{lem:DE-A3-ppav-odd-norm}, with the even-order checks in
Theorem~\ref{thm:DE-A3-ordinary-all-ranks}.  Those checks require
only etaleness, not ordinarity of the source.  We give the additional
jet arguments for $n\ge4$.

Suppose first that the principal source is indecomposable.  For any
two distinct points $P,Q$, the system $|n\Theta|$ separates their
six first-jet coordinates.  Put $r=n-3$ and consider pairs $(a,b)$
satisfying $3b=-ra$.  For
\[
 R_a=t_a^*\OO_B(r\Theta),\qquad
 \OO_B(n\Theta)\otimes R_a^{-1}=t_b^*\OO_B(3\Theta),
\]
the equality is in the Picard group, by the theorem of the square.
The parameter of pairs $(a,b)$ is finite and surjective over either
copy of $B$, and every component dominates both.  The first projection
is etale since $3$ is invertible.
The distinguished section of $R_a$ is nonzero at $P,Q$ on a nonempty
open in $a$.  For the translated cubic-theta bundle the rank-six
failure at $P,Q$ is contained in
\[
               2P+Q+3b\in\Theta,\qquad
               P+2Q+3b\in\Theta.
\]
The two theta translates are distinct, so there are only finitely
many such $b$.  Choose $(a,b)$ outside both bad loci.  Multiplication
by the section of $R_a$ then transfers the full six first jets to
$|n\Theta|$.  The same multiplication argument, using any translate
of $r\Theta$ avoiding three prescribed points, transfers the
three-value independence of Lemma~\ref{lem:ppav-three-distinct-values}
to $|n\Theta|$.

For a decomposable source, write
$\OO_B(n\Theta)=\OO_{E_1}(nO_1)\boxtimes\OO_{E_2}(nO_2)$.
On an elliptic curve a degree-$n$ bundle separates three distinct
values when $n\ge4$.  It separates the first jets at two distinct
points when $n\ge5$.  At $n=4$, those four first-jet conditions have
rank four except when the bundle equals the divisor of the two
doubled points; at that exception they have rank three.  These
assertions follow directly from the degree and Riemann--Roch of the
kernel line bundle.

Fix a nonzero kernel translation $g=(g_1,g_2)$.  If both coordinates
are nonzero, six surface first jets are supplied whenever either
factor separates its two first jets, using tensor products and a
first jet at a single point in the other factor.  At $n=4$, failure
on both factors is a finite point set: on the $i$th factor the
exception is $4P_i+2g_i=0$.  If precisely one coordinate of $g$
is nonzero, the same six-jet assertion holds away from the possible
one-dimensional exceptional locus.  On that locus the surface rank
is five.  After imposing the two values, the two derivatives in the
varying-coordinate direction form a nonzero line, while the two
transverse derivatives vary independently.  Their tangent determinant
is therefore a nonzero polynomial.  Finally any three distinct
surface points impose three independent values: separate their
first-coordinate classes and, within each class, their at most
three second coordinates.

These jet assertions give the required incidences.  At a rank-six
pair the two values and the proportional-gradient condition have
codimension three; the point moves in dimension two.  At the
rank-five locus just described there is still that third nonzero
condition, and the point moves in dimension one.  A finite exceptional
point set is avoided by the two independent value conditions.
Three-value independence excludes all triple identifications, since
the point parameter has dimension two.  The finitely many kernel
translations and pairs are handled simultaneously.

A nonzero translation preserving the bundle has nontrivial projective
action in its very ample embedding, and hence a proper fixed locus
of sections; a translation not preserving the bundle preserves no
member.  Thus a general $C$ has no kernel stabilizer and maps
birationally to its image.  The image is integral nodal because the
isogeny is etale and its remaining identifications are transverse
pairs.  Each nonzero kernel translation gives $2n^2$ ordered
intersection points.  Counting unordered pairs yields
$n^2(\deg\phi-1)$ nodes, including for even kernel order.
\end{proof}

\begin{theorem}[An eventual rank formula for every ordinary surface embedding]
\label{thm:DE-A3-ordinary-arbitrary-embedding}
\status{Proved; root, monodromy, and reduction-descent full independent proof audits passed.}
Let $A$ be an ordinary abelian surface over a perfect field of
characteristic two, and let $M$ be any ample line bundle with
$D=\chi(M)$.  For every integer $n\ge3$, a geometric generic smooth
member of $|M^n|$ is ordinary.  Its genus and relative stable rank are
\[
                         g_n=Dn^2+1,
                         \qquad r_n=Dn^2-1.
\]
If $M$ itself defines the chosen projective embedding, the same
relative rank formula holds for ambient degree-$n$ hypersurfaces
whenever
\[
               n\ge\max\{3,\operatorname{reg}(\mathcal I_A)-1\}.
\]
\end{theorem}
\begin{proof}
Work geometrically.  The ordinary principal norm source of
Lemma~\ref{lem:DE-A3-ordinary-etale-norm} is etale of degree $D$
and has $\operatorname{Nm}_\phi(\Theta_B)\equiv M$.
The first lemma supplies a nonempty open of smooth ordinary curves
in $|n\Theta_B|$, and the second supplies a nonempty open whose
images are integral nodal.  Choose a curve in both opens.  Its
normalization has genus $n^2+1$, and its image has
$n^2(D-1)$ nodes.  The normalization sequence proves that the
image is ordinary of arithmetic genus $Dn^2+1$.

Its norm line bundle has numerical class $M^n$.  One arbitrary
geometric translation, supplied by surjectivity of $\lambda_{M^n}$,
puts the image in the exact system $|M^n|$.
The norm of the very ample source bundle is very ample by the
point-and-tangent norm argument already proved; the same translation
therefore shows that $M^n$ is very ample.  The constant-dimensional
coherent family and Frobenius openness give a smooth ordinary generic
member of that exact system.  Subtracting the ordinary ambient
cohomology of rank two gives $r_n=Dn^2-1$.
The final ambient statement is the usual surjectivity furnished by
the displayed regularity range, with integral parameter fibres.
No assertion about the systems $|M|$ or $|M^2|$ is needed or inferred.
\end{proof}

\begin{remark}[Retained arithmetic verification]
\label{rem:DE-A3-ordinary-arbitrary-embedding-checks}
\status{Proved; exact certificate and independent proof audits retained.}
The program
\path{agents/obstructions_crystals/code/ordinary_arbitrary_embedding_checks.py}
checks both theta-union graph formulas, the norm-node count, genus and
relative-rank identities for $254$ values of $n$ and $32512$ pairs
$(n,D)$.  The author run passed in $0.012083560999599285$ seconds;
the complete output is retained alongside the program.  An independent
reviewer read the entire program and reran it with the author output
preserved, passing in $0.013041352001891937$ seconds.  These finite
identities supplement the geometric proofs and do not establish
existence by sample extrapolation.
\end{remark}


\subsection{Mixed elliptic products in every sufficiently large product degree}
\label{sec:DE-A3-mixed-all-products}

\begin{lemma}[A sharp mixed norm whenever the ordinary factor has odd degree]
\label{lem:DE-A3-mixed-odd-factor}
\status{Proved; root and monodromy full independent proof audits passed.}
Let $E_{\rm o}$ be an ordinary elliptic curve and $E_{\rm s}$ a
supersingular elliptic curve over an algebraically closed field of
characteristic two.  For positive integers $c,e$ with $c$ odd, a
general smooth divisor in
\[
 \bigl|\OO_{E_{\rm o}}(3cO)\boxtimes\OO_{E_{\rm s}}(3eO)\bigr|
\]
has genus $9ce+1$, full stable rank $9ce$, and relative stable rank
$9ce-1$.  No parity condition on $e$ is imposed.
\end{lemma}
\begin{proof}
There is an ordinary elliptic curve $E'_{\rm o}$ with an etale isogeny
$\phi_{\rm o}:E'_{\rm o}\to E_{\rm o}$ of any prescribed positive
degree $e$.  To see this, write $e=2^v e_0$ with $e_0$ odd.
The iterated Verschiebung
$E_{\rm o}^{(2^v)}\to E_{\rm o}$ is etale of degree $2^v$.
Choose a cyclic subgroup of order $e_0$ on $E_{\rm o}^{(2^v)}$;
the dual of its quotient isogeny gives an etale isogeny of degree
$e_0$ onto $E_{\rm o}^{(2^v)}$.  Compose the two maps.
The source stays ordinary by isogeny invariance of the $p$-rank.
The precise primary inputs are the Verschiebung degree and
factorization in \cite[Propositions~5.20--5.22 and their proofs,
pp.~80--81]{EdixhovenVanDerGeerMoonenAV}.  In dimension one and
$p$-rank one, that factorization makes Verschiebung separable of
its full degree and hence etale.  Dual isogenies and their degrees
are \cite[Theorem~7.5 and Corollary~7.17, pp.~100,104]
{EdixhovenVanDerGeerMoonenAV}.
The supersingular norm construction of
Lemma~\ref{lem:DE-A3-ss-norm}, or the same odd-torsion quotient and
dual construction, supplies an etale isogeny
$\phi_{\rm s}:E'_{\rm s}\to E_{\rm s}$ of degree $c$.
Its source is supersingular.

Give $B=E'_{\rm o}\times E'_{\rm s}$ the product principal
polarization $\Theta$.  The universal principal-surface rank theorem,
Theorem~\ref{thm:DE-A3-ppav-ranks}, gives a nonempty open of smooth
curves $C'\in|3\Theta|$ of genus ten and full stable rank nine.
Here the ambient $p$-rank is one, so the full curve defect is one.
Apply the product etale isogeny
\[
                 \phi=\phi_{\rm o}\times\phi_{\rm s}:B\to
                       E_{\rm o}\times E_{\rm s},
             \qquad \deg\phi=ce.
\]
The complete cubic nodal-image argument in
Lemma~\ref{lem:DE-A3-all-multiple-nodal-norm} applies to this
principally polarized source.  That lemma explicitly allows even
etale degree and does not require the source to be ordinary.
For clarity, on a product cubic a nonzero translation in just one
factor gives rank five on the two surface first jets; after imposing
the two values, its tangent determinant remains nonzero.  A translation
nonzero in both factors has the previously computed rank-six open,
with its smaller-rank loci of the required smaller dimensions.
Three-value failure can only occur when all points have the same
coordinate in one factor and the other three coordinates sum to zero;
that locus has dimension one because multiplication by three is finite.
These are exactly the pair and triple incidences used in the cited
lemma, and none depends on ordinarity or on inverting the isogeny degree.
Thus a general sharp source curve has integral nodal image and is
its normalization.

The norm of the product cubic has numerical class $(3c,3e)$.
The image has genus $9ce+1$ and exactly $9(ce-1)$ nodes.
Its full stable rank is consequently $9+9(ce-1)=9ce$.
Surjectivity of the target ample polarization on geometric points
allows a translation putting this image in the exact line-bundle
class displayed in the statement.
The normalization sequence makes each node contribution
Frobenius bijective.  Openness of iterated Frobenius rank, and
Bertini for the very ample target bundle, transfer this lower bound
to a general smooth member.

For the matching upper bound, the divisor sequence gives a
Frobenius-compatible injection of the ambient $H^1(\OO)$ into the
curve cohomology.  It contains a one-dimensional nilpotent subspace
coming from $E_{\rm s}$.  Hence a smooth curve in this system cannot
be ordinary, and its full stable rank is at most its genus minus one.
This is exactly $9ce$.  Subtracting the ambient stable rank one
gives the relative rank $9ce-1$.
\end{proof}

\begin{lemma}[The additional Cartier kernel in an even ordinary factor]
\label{lem:DE-A3-mixed-even-bound}
\status{Proved; root and monodromy full independent proof audits passed.}
On $A=E_{\rm o}\times E_{\rm s}$ as above, put
$N=\OO_{E_{\rm o}}(mO)\boxtimes\OO_{E_{\rm s}}(nO)$, with
$m,n\ge3$ and $m$ even.  The geometric generic relative stable
rank of a divisor in $|N|$ is at most $mn-2$.
\end{lemma}
\begin{proof}
Choose invariant differentials $\omega_{\rm o},\omega_{\rm s}$,
with $C(\omega_{\rm s})=0$, and let $\partial_{\rm o}$ be the
vector field dual to $\omega_{\rm o}$.
Because the first factor of $N$ is pulled back by Frobenius,
its canonical connection makes $\partial_{\rm o}$ a global
endomorphism of $H^0(A,N)$.  With horizontal rational
trivializations, the Cartier multiplication operator satisfies
\[
                        C_s(s)=0,\qquad
                        C_s(\partial_{\rm o}s)=0.
\]
The first identity uses $C(\omega_{\rm o}\wedge\omega_{\rm s})=0$.
For the second, locally write $s=f$.  Then
\[
 f\partial_{\rm o}f\,\omega_{\rm o}\wedge\omega_{\rm s}
                      =f\,df\wedge\omega_{\rm s}.
\]
Both $f\,df$ and $\omega_{\rm s}$ are closed.  Multiplicativity
of Cartier and $C(\omega_{\rm s})=0$ give the claimed vanishing.
Horizontal transitions in the first factor and transitions
depending only on the second ensure that this is a global identity.

The two kernel sections are generically independent.
Over the algebraic closure, an ordinary characteristic-two
Weierstrass equation can be put in the form
$y^2+xy=x^3+a_2x^2+a_6$, with $a_6\ne0$; compare
\cite[Example~5.26, p.~82]{EdixhovenVanDerGeerMoonenAV}.
Explicitly, scale to $a_1=1$, translate $x$ to remove $a_3$,
and then translate $y$ by a constant to remove $a_4$;
nonsingularity forces the resulting $a_6$ to be nonzero.
Its invariant differential is $dx/x$ and its dual vector field
satisfies $\partial_{\rm o}x=x$,
$\partial_{\rm o}y=x^2+y$, and
$\partial_{\rm o}^2=\partial_{\rm o}$.
The rational section $1+x$ and its derivative
$x$ therefore already give an independent pair in
$H^0(E_{\rm o},\OO(mO))$, since $m\ge4$ and $x$ has pole order two.
Tensor with the constant rational section of the second factor.
The corresponding rank minors vanish generically and hence
identically, so the Cartier rank is at most $mn-2$.

Finally the divisor sequence and ample vanishing identify the
stable part of $H^2(A,N^{-1})$ with the relative stable cohomology:
the ambient $H^2(\OO_A)$ has Frobenius zero, and the ambient
$H^1(\OO_A)$ has stable rank one.  Serre duality identifies the
former Frobenius with the twisted Cartier operator just used.
Linear rank bounds stable rank, proving the assertion.
\end{proof}

\begin{theorem}[All mixed product polarizations from degree two onward]
\label{thm:DE-A3-mixed-all-product-ranks}
\status{Proved; root and monodromy full independent proof audits passed.}
Let $E_{\rm o}$ and $E_{\rm s}$ be respectively ordinary and
supersingular elliptic curves over a perfect field of characteristic
two.  For arbitrary positive integers $a,b$ put
\[
 A=E_{\rm o}\times E_{\rm s},\qquad
 M=\OO_{E_{\rm o}}(aO)\boxtimes\OO_{E_{\rm s}}(bO),\qquad L=M^3.
\]
For every $d\ge2$ the geometric generic relative stable rank is
\[
                    r_d=9abd^2-1-\epsilon(ad),
 \qquad
 \epsilon(c)=
 \begin{cases}1&c\text{ even},\\0&c\text{ odd}.\end{cases}
\]
The same formula holds for $d=1$ if $a$ is odd or $b\ge2$.
The genus is $9abd^2+1$ and the relative coherent dimension is
$9abd^2-1$.  For the embedding by $L$, the formula also applies
to ambient hypersurfaces once
$d\ge\max\{2,\operatorname{reg}(\mathcal I_A)-1\}$.
\end{theorem}
\begin{proof}
Pass to an algebraic closure.  If $ad$ is odd, apply
Lemma~\ref{lem:DE-A3-mixed-odd-factor} with $c=ad$, $e=bd$.
Suppose $ad$ is even and $bd\ge2$.  Decompose the class as
\[
              (3ad,3bd)=(3,3)+(3(ad-1),3(bd-1)).
\]
Both summands have positive entries and odd ordinary-factor
coefficient.  Lemma~\ref{lem:DE-A3-mixed-odd-factor} gives
nonempty sharp smooth opens in both complete systems.
The two bundles are very ample, so a general pair in these
opens is transverse.  Their intersection number is positive.
The connected nodal union has full curve defect two:
if its components have genera $g_1,g_2$ and meet in $t>0$
points, normalization gives genus $g_1+g_2+t-1$ and stable rank
$(g_1-1)+(g_2-1)+t-1$.
It therefore has full stable rank $9abd^2-1$ and relative
stable rank $9abd^2-2$.  The rank open meets the nonempty
smooth-member open in the target system.
Lemma~\ref{lem:DE-A3-mixed-even-bound} gives the matching upper bound.
This proves all cases with $d\ge2$, as well as the stated
degree-one cases.
Adjunction and the same divisor sequence give the displayed
genus and coherent dimension.  The regularity bound makes
ambient section restriction surjective with integral parameter
fibres, so it preserves the geometric generic rank.
\end{proof}

\begin{remark}[The remaining small system and the scope]
\label{rem:DE-A3-mixed-product-scope}
\status{Proved; root and monodromy full independent proof audits passed.}
The argument deliberately leaves only the degree-one systems
$|\OO_{E_{\rm o}}(3aO)\boxtimes\OO_{E_{\rm s}}(3O)|$ with
$a$ even outside its stated range.
Corollary~\ref{cor:DE-A3-mixed-thin-vone} subsequently settles
the cases $a\equiv2\pmod4$ by a universal degree-six seed and
the correctly directed supersingular-axis norm.
At that stage only $4\mid a$ lay outside these results.
The uniform residue proof in
Corollary~\ref{cor:DE-A3-mixed-product-no-exception} subsequently
settles all of those cases, and gives the formula for every $a,b,d$.
The historical scope of the preceding theorem is retained.
The theorem determines the asymptotic
rank for every mixed product polarization, without a generic
$j$-invariant hypothesis or a restriction on the two-adic valuation.
It does not determine the generic rank for an arbitrary
nonproduct surface or arbitrary smooth projective variety.
The isogeny and nodal-image arguments used here were reread in
their full earlier proofs before applying them; there is no
appeal to a sample of elliptic curves.
The retained exact program
\path{agents/obstructions_crystals/code/mixed_product_rank_arithmetic.py}
checks $2048$ norm identities, $2016$ transverse-union identities,
and $65504$ covered degree cases, retaining the excluded small cases
separately.  Its author run passed in $0.03129147599975113$ seconds,
with full output beside the program.  The root reviewer read the entire
certificate and reran it with the author output preserved, passing in
$0.03777776699280366$ seconds.  The infinite existence assertions
are furnished by the preceding geometric proofs.
\end{remark}


\subsection{A nonproduct mixed class detected by the ordinary polarization index}
\label{sec:DE-A3-mixed-nonproduct}

\begin{lemma}[The two integral polarization indices in mixed dimension two]
\label{lem:DE-A3-mixed-polarization-blocks}
\status{Proved; root and monodromy full independent proof/source/code audits passed.}
Let $A$ be an abelian surface of $2$-rank one over an algebraically
closed field $k$ of characteristic two, and let $M$ be ample.
Put $D=\chi(M)$ and define
\[
 a(M)=\operatorname{length}_{\mathbf Z_2}
       \coker\bigl(T_2(\lambda_M):T_2(A)\longrightarrow T_2(A^\vee)\bigr).
\]
These Tate modules both have rank one.  There is an integer $b(M)\ge0$
such that
\[
                         v_2(D)=a(M)+b(M).
\]
The elementary divisors of the integral polarization pairing on the
ordinary Dieudonne summand are $(a(M),a(M))$, and those on its
height-two local-local summand are $(b(M),b(M))$.  In particular,
$\ker\lambda_M$ is killed by $D$.

There is an ample line bundle $N$ on $A^\vee$ satisfying
\[
 \lambda_M^*N\simeq M^D,\qquad
 \lambda_N=D\lambda_M^{-1},\qquad
 \chi(N)=D,\qquad
             a(N)=b(M),\quad b(N)=a(M).
\]
The inverse here is first taken in the rational homomorphism space;
the assertion includes the integrality of the displayed polarization.
\end{lemma}
\begin{proof}
We explain both the finite-flat splitting and the integral pairing.
Over a perfect field the connected-etale sequence of a finite
commutative group scheme splits canonically and functorially:
the reduced subgroup is the etale factor.  There are no nonzero
homomorphisms between a local group and an etale group.  These
statements, including the functoriality needed for transition maps,
are \cite[Lemma~4.44 and Proposition~4.45, p.~67]
{EdixhovenVanDerGeerMoonenAV}.
Apply them at every level of $A[2^\infty]$, and then apply the same
splitting to the Cartier dual of the connected factor.  The compatible
splittings give an actual product of $2$-divisible groups
\[
 A[2^\infty]=G_{\et}\times G_{\rm mult}\times G_{\rm ll},
 \qquad
 G_{\et}\simeq\mathbf Q_2/\mathbf Z_2,\quad
 G_{\rm mult}\simeq\mu_{2^\infty}.
\]
Here $G_{\rm ll}$ and its Cartier dual are connected.  It has height
two and dimension one.  Indeed the etale factor has height one by
the $2$-rank hypothesis; duality and isogeny invariance give the same
height for the multiplicative factor, leaving height two out of the
total height four.  The same reasoning applies to $A^\vee$.
Maps between distinct factors vanish.  For the multiplicative versus
local-local assertion, apply Cartier duality and the local-etale
vanishing.  Thus every homomorphism under consideration preserves
these factors, before any passage to isogeny categories.

Use the covariant Dieudonne lattice, equivalently the $W(k)$-linear
dual of Oda's contravariant lattice.  It is free, and Cartier duality
is its integral $W(k)$-linear dual, with the usual exchanged
Frobenius and Verschiebung actions.  Exactness and the length/order
dictionary are \cite[Theorem~3.14(ii)--(iii) and
Corollaries~3.15--3.16, pp.~87--88; Proposition~3.18, p.~89;
Proposition~3.22, p.~99]{Oda1969Dieudonne}.
The polarization gives a skew-symmetric integral bilinear Riemann
form on this lattice by
\cite[Propositions~1.10 and 1.12, pp.~75--76, and
Proposition~3.24(ii), pp.~100--101]{Oda1969Dieudonne}.
Since $W(k)$ has no $2$-torsion, skew symmetry implies that the
diagonal values are zero.  The form is therefore alternating even
in the present residue characteristic.  The ordinary and local-local
summands are orthogonal, by the vanishing of the cross-factor
homomorphisms just proved.

The ordinary summand pairs its rank-one etale and multiplicative
parts.  Symmetry of the polarization and integral Cartier duality
make their two indices equal; the etale one is precisely $a(M)$.
On the free rank-two local-local lattice, an alternating matrix has
the form
\[
                         \begin{pmatrix}0&c\\-c&0\end{pmatrix},
                         \qquad 0\ne c\in W(k).
\]
Set $b(M)=v_2(c)$.  Its two elementary divisors are both $b(M)$;
this integer is independent of the chosen basis.
The kernel of an isogeny has Dieudonne length equal to the valuation
of its order.  Since $\deg\lambda_M=D^2$, the two blocks give
$2v_2(D)=2a(M)+2b(M)$.  Their finite cokernels are killed by
$2^{\max\{a(M),b(M)\}}$, hence by the two-primary part of $D$.
Exact Dieudonne theory transfers this assertion to the finite
group-scheme kernel.

For an odd prime $\ell$, the kernel is a constant finite abelian
$\ell$-group with a perfect alternating commutator pairing.
If an element has maximal order $\ell^e$, perfection supplies a
partner with pairing of that same order.  The two generate a group
of order $\ell^{2e}$, by pairing any proposed relation with the two
generators.  Thus $e\le v_\ell(D)$.
Together with the two-primary calculation this proves
$D\ker\lambda_M=0$.

Consequently the commutator pairing of $M^D$ is trivial on
$\ker\lambda_M$.  The group-scheme descent criterion
\cite[Corollary~8.11 and its proof, p.~117]
{EdixhovenVanDerGeerMoonenAV}, which allows a nonreduced isotropic
subgroup, gives $N$ with $\lambda_M^*N\simeq M^D$.
It is ample because its pullback by the finite surjection is ample.
Polarizations of pullbacks give
\[
                   \lambda_M\lambda_N\lambda_M=D\lambda_M,
\]
hence the asserted integral inverse formula.
Intersection degrees give
$\chi(N)=\chi(M^D)/\deg\lambda_M=D^3/D^2=D$.
On each of the two blocks, multiplying the inverse matrix by $D$
replaces its elementary-divisor valuation by $v_2(D)$ minus that
valuation.  The two indices of $N$ are therefore $b(M)$ and $a(M)$,
as stated.
\end{proof}

\begin{lemma}[An etale principal source exactly at ordinary index zero]
\label{lem:DE-A3-mixed-index-zero-source}
\status{Proved; root and monodromy full independent proof/source/code audits passed.}
Under the hypotheses of the preceding lemma, the condition $a(M)=0$
is equivalent to the existence of a principally polarized surface
$(B,\Theta_B)$ and an etale isogeny $\phi:B\to A$ of degree $D$
such that
\[
                \phi^*M\equiv D\Theta_B,\qquad
                \operatorname{Nm}_\phi(\Theta_B)\equiv M.
\]
The surface $B$ has $2$-rank one.  No parity restriction is imposed
on $D$, and no product decomposition of $A$ is assumed.
\end{lemma}
\begin{proof}
Suppose first that $a(M)=0$.  Construct $N$ on $A^\vee$ by
Lemma~\ref{lem:DE-A3-mixed-polarization-blocks}.
Its local-local block has index zero, so the two-primary kernel of
$\lambda_N$ is entirely in the ordinary factor.  It is the product
of a multiplicative group and its etale Cartier dual, each of order
$2^{b(M)}$.  Its commutator pairing is trivial on either factor and
perfect between them: restriction to a single factor is a map from
a local multiplicative group to an etale dual, or its reverse, and
must vanish.  Thus its entire multiplicative factor is maximal
isotropic.

Choose this factor in the two-primary kernel and a maximal isotropic
subgroup in each odd-primary kernel.  Their product $H$ is isotropic
of order $D$, and $H^D$ is etale.  The descent theorem gives
\[
       \psi:A^\vee\longrightarrow C=A^\vee/H,\qquad
       \psi^*\Theta_C\simeq N
\]
with $\Theta_C$ principal
\cite[Proposition~8.19 and proof, p.~119]
{EdixhovenVanDerGeerMoonenAV}.
Put $B=C^\vee$, take its dual principal polarization, and set
$\phi=\psi^\vee$.  Isogeny duality gives
$\ker\phi=H^D$ \cite[Theorem~7.5, p.~100]
{EdixhovenVanDerGeerMoonenAV}, so $\phi$ is etale.
The typed polarization identity is
\[
                       \psi\lambda_M\phi=D\lambda_C^{-1}.
\]
It gives the first numerical equality; applying norm and cancelling
$D$ in the torsion-free Neron--Severi group gives the second.
Isogeny invariance of $2$-rank gives $f_B=1$.

Conversely, an etale isogeny induces an isomorphism on the
local-local $2$-divisible factors.  The local-local elementary
divisors of $\phi^*M$ are therefore both $b(M)$.  Since $\Theta_B$
is principal, those of $D\Theta_B$ are both $v_2(D)$.
The equality of their polarizations implies $b(M)=v_2(D)$,
and hence $a(M)=0$.
\end{proof}

\begin{lemma}[The even-degree mixed Prym bound without separability]
\label{lem:DE-A3-mixed-all-even-Prym}
\status{Proved; root and monodromy full independent proof/source/code audits passed.}
Let $A,M,D$ be as in
Lemma~\ref{lem:DE-A3-mixed-polarization-blocks}, with no restriction
on either index.  Every smooth curve $Y\in|M^m|$ with $m$ even
satisfies
\[
                              g(Y)-f_Y\ge2.
\]
\end{lemma}
\begin{proof}
The proof of Lemma~\ref{lem:DE-A3-general-polarization-Prym}
uses oddness of $D$ only to construct the integral line bundle $N$.
That construction is now supplied for every $M$ by
Lemma~\ref{lem:DE-A3-mixed-polarization-blocks}.
We give the remaining argument explicitly.
Ample vanishing and Serre duality give
$H^1(A,P\otimes M^{-m})=0$ for every $P\in\Pic^0(A)$.
For nontrivial $P$, also $H^0(A,P)=0$, so restriction to $Y$
cannot trivialize $P$.  The same vanishing for $P=\OO_A$ makes
the tangent map on Picard varieties injective.  Thus restriction
is an immersion $j:A^\vee\hookrightarrow J_Y$.
Its polarized dual $q:J_Y\to A$ has smooth connected kernel $P_Y$:
quotient by the reduced identity kernel and use isogeny duality
to exclude a remaining finite quotient, exactly as in the cited lemma.

The integral polarization identity is
\[
                 qj=j^*\lambda_{J_Y}=mD\lambda_M^{-1}=m\lambda_N.
\]
Here the first equality identifies the polarization with its
homomorphism $A^\vee\to A$; the middle identity is
\cite[Section~3, equations~(3.2)--(3.3), pp.~346--347]
{Debarre1994Degrees}.  It holds in all characteristics.
Since $m$ is even, the actual finite group scheme $A^\vee[2]$
embeds in $P_Y[2]$.
Twist this finite-group immersion by Frobenius and apply functorial
Verschiebung.  The connected part of
$\ker(V_{A^\vee}:(A^\vee)^{(2)}\to A^\vee)$ has order two and injects
into the connected part of
$\ker(V_{P_Y}:P_Y^{(2)}\to P_Y)$.
Hence $P_Y$ has $2$-rank at most $\dim P_Y-1$.
The quotient $J_Y/P_Y\simeq A$ has defect one; adding defects,
using isogeny invariance and a complementary abelian subvariety,
gives $g(Y)-f_Y\ge2$.
\end{proof}

\begin{theorem}[All degrees for mixed surfaces of ordinary index zero]
\label{thm:DE-A3-mixed-nonproduct-ranks}
\status{Proved; root and monodromy full independent proof/source/code audits passed.}
Let $A$ be an abelian surface of $2$-rank one over a perfect field
of characteristic two, and let $M$ be ample.  Compute its ordinary
index after passage to an algebraic closure, and assume $a(M)=0$.
Write $D=\chi(M)$ and $L=M^3$.  Then $L$ is very ample.
For every $d\ge1$, the geometric generic relative stable rank
in the complete system $|L^d|$ is
\[
 r_d=
 \begin{cases}
       9Dd^2-1,&d\ \text{odd},\\
       9Dd^2-2,&d\ \text{even}.
 \end{cases}
\]
The relative coherent dimension is $9Dd^2-1$, and the generic
smooth curve has genus $9Dd^2+1$.  For the embedding defined by $L$,
the same ranks hold for ambient hypersurfaces whenever
$d\ge\max\{1,\operatorname{reg}(\mathcal I_A)-1\}$.
\end{theorem}
\begin{proof}
Work over an algebraic closure.  The etale source from
Lemma~\ref{lem:DE-A3-mixed-index-zero-source}, and the norm
point-and-tangent argument in
Theorem~\ref{thm:DE-A3-ordinary-all-ranks}, show that $M^3$ is
very ample.  That argument works for any finite etale degree:
one chooses a section with the desired value or first derivative
at one chosen lift, avoids its finitely many other value hyperplanes,
and takes its norm.  Numerical equivalence is converted to a
translation because the ample polarization is surjective.

For odd $d$, the ordinary index of $M^d$ is still zero.
Apply Lemma~\ref{lem:DE-A3-mixed-index-zero-source} to $M^d$.
Its principal source $B_d$ has $2$-rank one, and the etale norm
degree is $Dd^2$.  Proposition~\ref{prop:DE-A3-ppav-degree-one}
gives a nonempty open of smooth genus-ten curves of $2$-rank nine
in $|3\Theta_{B_d}|$.
A general such curve has integral nodal image under the isogeny.
The full cubic nodal-image proof in
Theorem~\ref{thm:DE-A3-ordinary-all-ranks} applies here:
its two-first-jet and three-value incidence arguments concern
the principal source and distinct kernel points, and do not use
ordinarity of the source or invertibility of the kernel order.
Thus the image has $9(Dd^2-1)$ nodes and normalization of defect one.
An arbitrary target translation puts its line bundle in the exact
system $|M^{3d}|$, since its norm class is $3dM$.
This uses surjectivity of the polarization, not an inverse on
two-torsion points.

The normalization sequence adds a Frobenius-bijective graph space.
The nodal image therefore has full defect one and the maximal
possible relative stable rank $9Dd^2-1$.
Ample vanishing gives constant full coherent dimension
$9Dd^2+1$ and relative dimension $9Dd^2-1$ over the coefficient
space.  A sufficiently long Frobenius iterate detects the stable
rank.  Its maximal-rank open meets the nonempty smooth-member open,
proving the odd formula.

For even $d$, choose sharp smooth members of $|L|$ and
$|L^{d-1}|$ meeting transversely.  Their connected nodal union
has full defect two, because each normalization component has
defect one and the graph contribution is Frobenius-bijective.
Openness gives the lower bound $r_d\ge9Dd^2-2$.
Lemma~\ref{lem:DE-A3-mixed-all-even-Prym}, applied with $m=3d$,
gives the matching upper bound for every smooth member.
Finally the stated regularity bound makes restriction from ambient
hypersurfaces surjective, with geometrically integral parameter
fibres, and hence preserves the generic rank.
\end{proof}

\begin{remark}[What the ordinary index adds]
\label{rem:DE-A3-mixed-nonproduct-scope}
\status{Proved; root and monodromy full independent proof/source/code audits passed.}
The condition in Theorem~\ref{thm:DE-A3-mixed-nonproduct-ranks}
is intrinsic to the polarized surface and does not require an
elliptic product or an odd polarization degree.
For a product of an ordinary and a supersingular elliptic curve,
with product polarization of degrees $(u,v)$ in that order,
one has $a(M)=v_2(u)$ and $b(M)=v_2(v)$.
Thus this theorem includes the odd-ordinary-factor product cases,
while also applying to any nonproduct surface satisfying the
same integral condition.
Such even-degree classes can be constructed directly: start with
any principally polarized surface of $2$-rank one and quotient by
an etale subgroup of order $2^b$ in its rank-one etale
$2$-divisible factor.  The norm of the principal bundle is ample,
has Euler characteristic $2^b$, and its pullback has numerical
class $2^b\Theta$.  The converse in
Lemma~\ref{lem:DE-A3-mixed-index-zero-source} therefore gives
ordinary index zero.  If the chosen source is geometrically simple,
its isogenous target remains so.
The class $a(M)>0$ is not settled by this norm construction:
Lemma~\ref{lem:DE-A3-mixed-index-zero-source} shows precisely why
an etale principal source of the prescribed degree cannot exist
there.  No formula for arbitrary ambient varieties is asserted.

The retained certificate
\path{agents/obstructions_crystals/code/mixed_nonproduct_rank_checks.py}
proves the block inverse and determinant identities symbolically,
including an arbitrary unit multiplier.  It also checks $289$
index pairs, $1024$ odd norm identities and $1024$ even union identities.
The author run passed in $0.04614707099972293$ seconds, with its
complete output retained alongside the program.  The existence and
group-scheme assertions follow from the proofs above, not from the
finite arithmetic checks.
Root independently read the entire proof, original sources and program,
and reran the certificate in $0.049978896997345146$ seconds.
Monodromy independently completed the same proof/source/program audit
and reran it in $0.03037099799985299$ seconds.  Their outputs are retained
separately; the author output was preserved.
\end{remark}


\subsection{A uniform mixed-product seed of type $(6,3)$}
\label{sec:root-mixed-six-three}
\status{Proved. Complete independent obstruction/crystal and
reduction/descent proof and exact-program audits passed.}

\begin{lemma}[A universal mixed $(6,3)$ Cartier witness]
\label{lem:root-mixed-six-three-cartier}
\status{Proved under stated hypotheses}
Let $k$ be algebraically closed of characteristic two, let
$E_{\rm o}/k$ be any ordinary elliptic curve, and let
$E_{\rm s}:w^2+w=z^3$. On $A=E_{\rm o}\times E_{\rm s}$ put
\[
 N=\OO_{E_{\rm o}}(6O)\boxtimes\OO_{E_{\rm s}}(3O).
\]
A geometric generic smooth $C\in|N|$ has genus $19$, $p$-rank
$17$, and relative stable rank $16$. This includes every ordinary
$j$-invariant, not only a transcendental one.
\end{lemma}
\begin{proof}
Over $k$, an ordinary Weierstrass equation can be normalized to
\[
 E_{\rm o}:y^2+xy=x^3+c,\qquad c\in k^*.
\]
Indeed weighted scaling makes $a_1=1$. Translating $x$ by $a_3$
removes the $a_3y$ term; translating $y$ by the resulting
$x$-coefficient removes that coefficient as well. A further
change $y\mapsto y+u x$ with $u^2+u=a_2$ removes the remaining
quadratic coefficient.
Smoothness is equivalent to $c\ne0$ in this model.
Use the invariant differentials $\omega_{\rm o}=dx/x$ and
$\omega_{\rm s}=dz$, and put $t=c^{1/2}$ and $u=y+t$.
The basis of $H^0(E_{\rm o},\OO(6O))$ is
\[
                    f=(1,x,y,x^2,xy,x^3).
\]
For $v=u/x$ the curve equation gives $y=x^2+xv^2$.
Thus $f_i=A_i^2+xB_i^2$, where
\[
 A=(1,0,x,x,u,0),\qquad B=(0,1,v,0,x,x).
\]
The ordinary Cartier multiplication pairing is therefore
\[
 \frac{C(f_if_j\omega_{\rm o})}{\omega_{\rm o}}
                         =A_iA_j+xB_iB_j.
\]
In the indicated basis its entries, written as functions, are
\[
 \beta_{\rm o}=
 \begin{pmatrix}
 1&0&x&x&u&0\\
 0&x&u&0&x^2&x^2\\
 x&u&y&x^2&0&xu\\
 x&0&x^2&x^2&xu&0\\
 u&x^2&0&xu&xy&x^3\\
 0&x^2&xu&0&x^3&x^3
 \end{pmatrix}.
\]
For the supersingular basis $e=(1,z,w)$ the pairing is
\[
 \beta_{\rm s}=
 \begin{pmatrix}0&1&z\\1&0&w\\z&w&0\end{pmatrix}.
\]
Both formulas follow directly by the $p$-basis Cartier rule;
the second uses $w=w^2+z^3$.

Index the tensor basis $f_i e_j$ by $3i+j$, starting at zero.
For a section $s$ let $T_s$ be the matrix of the inverse-Frobenius
semilinear map
\[
 C_s:H^0(A,N)\longrightarrow H^0(A,N),\qquad
 \sigma\longmapsto
 \frac{C(s\sigma\,\omega_{\rm o}\wedge\omega_{\rm s})}
      {\omega_{\rm o}\wedge\omega_{\rm s}}.
\]
The tensor product of the two displayed pairings specifies every
entry of $T_s$; the coefficient of $s$ is taken to its square root.
Its Serre-dual Frobenius-linear map has matrix
$M_s=(T_s^{[2]})^t$.

Choose the same section for every $c\ne0$:
\[
 s=1+z+x(1+z+w)+xyz+x^3(z+w).
\]
Its ordinary invariant derivative is
$\partial_{
\rm o}s=x+xz+xw+x^3w$.
The two sections are independent, and the displayed pairings give
$C_s(s)=C_s(\partial_{\rm o}s)=0$.
Consequently $\rank M_s\le16$.

There is also an exact uniform lower certificate.
In $M_sM_s^{[2]}$, select the ordered rows and columns
\begin{align*}
 I&=(0,7,2,3,8,5,6,9,11,1,10,14,13,15,4,12),\\
 J&=(0,1,2,3,4,5,6,7,9,10,11,12,13,14,15,16).
\end{align*}
Their $16\times16$ determinant, computed over $\mathbf F_2[c]$, is
\[
                  \det (M_sM_s^{[2]})_{I,J}=c^5.
\]
This is a polynomial identity, not a check at finitely many
parameters. The retained search program verifies it by exact binary
polynomial fraction-free elimination. The second verifier derives
all six $p$-basis identities and all $36$ ordinary-pairing identities
from the curve equation, reconstructs the entire tensor matrix,
checks the two kernel vectors, and verifies the same determinant
using independent polynomial-domain matrix arithmetic.
Since $c\ne0$, the first and second Frobenius images both have
dimension $16$. They are therefore stable, so the stable rank is
exactly $16$.

The nonvanishing minor defines a nonempty open in $|N|$.
The line bundle $N$ is very ample, so its smooth-divisor locus is
also nonempty and open. These two opens meet in the irreducible
space $|N|=\mathbf P^{17}$.
Lemma~\ref{lem:DE-A3-mixed-even-bound} gives the matching generic
upper bound $16$; equivalently its two generic Cartier-kernel
sections give this bound on the same coefficient space.
The universal divisor sequence has $H^1(A,N^{-1})=0$, while
$H^2(A,\OO_A)$ has nilpotent Frobenius. Taking semisimple parts
shows that the stable rank of $H^1(C,\OO_C)$ is the sum of the
stable rank $16$ just computed and the ambient rank
$f(A)=1$. Adjunction gives $g(C)=1+N^2/2=19$.
Thus $f(C)=17$ and the relative stable rank is $16$, as asserted.
\end{proof}

\begin{remark}[Exact certificate and a corrected transfer direction]
\label{rem:root-mixed-six-three-certificate}
\status{Proved under stated hypotheses}
The search program \path{code/mixed_6_3_cartier_search.py} retains
all six tested coefficient vectors and their outputs; its sixth
vector gives the universal minor $c^5$ in
$0.0044763490004697815$ seconds.
The independent construction in
\path{code/mixed_6_3_cartier_verify.py} passed in
$0.18401811700459803$ seconds. Its first execution encountered a
polynomial-domain coercion error in the software's ordinary
remainder routine; the complete failed source and output are
retained. Using pseudo-remainder with the monic curve equation
resolved the implementation error without changing any identity.
The two independent complete proof/program audits reran the search
in $0.002776610002911184$ and $0.002867126000637654$ seconds,
and the symbolic verifier in $0.15830573499988532$ and
$0.16291111600003205$ seconds, respectively. Their separate
outputs are retained without changing the author's outputs.

A proposed norm transfer along an ordinary-factor isogeny of
degree $n$ would send type $(6,3)$ to $(6,3n)$, not $(6n,3)$.
We record that proposal as invalid. The latter class is obtained
by an isogeny on the supersingular factor; that isogeny is etale
only when $n$ is odd. This seed by itself therefore addresses
the previously open thin-strip classes $a\equiv2\pmod4$,
subject to the separate nodal-image argument. It makes no claim
for $a$ divisible by four.
\end{remark}


\subsection{Axis norms from a mixed degree-six by degree-three seed}
\label{sec:DE-A3-mixed-axis-norms}

\begin{lemma}[Nodal images for an axis kernel]
\label{lem:DE-A3-axis-nodal-norm}
\status{Proved; complete independent proof audits by root and
reduction-descent passed.}
Let $B=E_1\times E_2$ over an algebraically closed field of
characteristic two, and put
\[
             N=\OO_{E_1}(6O)\boxtimes\OO_{E_2}(3O).
\]
Let $\phi:B\to A$ be a finite etale isogeny of degree $n$ whose
kernel is contained in one of the two factors.
A general member $C\in|N|$ is smooth and maps birationally to
an integral curve $Z=\phi(C)$ with exactly $18(n-1)$ ordinary
nodes and no other singularities.  Thus its normalization is $C$.
Every nonempty open condition on smooth members of $|N|$ can
be imposed at the same time.

If $\phi=\phi_1\times1$, then the norm class is $(6,3n)$.
If $\phi=1\times\phi_2$, then the norm class is $(6n,3)$.
These two formulae do not interchange the two factors.
\end{lemma}
\begin{proof}
The complete system has dimension $6\cdot3-1=17$ and is very
ample.  Its smooth-member locus is therefore a nonempty open.
Write $V=H^0(E_1,\OO(6O))$ and
$W=H^0(E_2,\OO(3O))$.
Fix a nonzero kernel translation $j$.
We verify directly that the condition that $C$ and $C-j$ have
a nontransverse intersection has incidence dimension at most
sixteen in $B\times|N|$.

First suppose that $j$ acts on $E_2$.
At a point $(u,v)$ the two value conditions are independent,
because degree three separates the distinct points $v,v+j$.
Their common coefficient fibre is a projective space of
dimension fifteen.
The value and first derivative at $u$ on the degree-six factor
can be prescribed independently.  A section, to first order
in the first-factor coordinate at $u$, can consequently be
written
\[
                        w_0+t w_1,
                    \qquad w_0,w_1\in W
\]
with $w_0,w_1$ independently arbitrary.
The two value conditions require $w_0(v)=w_0(v+j)=0$.
The resulting space of $w_0$ has dimension one, generated by
the line section through these two points in the plane-cubic
embedding.  Let $b_0,b_1$ be its derivatives at the two points.
They cannot both vanish: a nonzero line section of degree three
cannot have zeros of order at least two at two distinct points.
Meanwhile the two values $a_0=w_1(v)$ and $a_1=w_1(v+j)$ are
independently arbitrary.
The intersection determinant is a nonzero scalar multiple of
\[
                              a_0b_1-a_1b_0
\]
times the free coefficient of $w_0$.
It is therefore a genuine nonzero polynomial on every coefficient
fibre.  Its vanishing cuts that fibre in dimension at most fourteen.
Adding the two-dimensional point base gives the required bound.
This argument includes a chord tangent at either of its two points.

If $j$ acts on $E_1$, evaluation of degree-six sections on the
two first jets at $u,u+j$ is surjective: the kernel bundle has
degree $6-4=2>0$ and hence has vanishing $H^1$.
Evaluation of the degree-three factor on its first jet at $v$
is likewise surjective.  Tensoring these evaluations shows that
the two values and all four first partial derivatives can be
prescribed independently.  After imposing the values, the
tangent determinant is again a nonzero polynomial on every
coefficient fibre.  The same dimension bound follows.
Since the incidence is proper, a general section has only
transverse intersections with each of its nontrivial translates.

We also exclude three preimages of a quotient point.
For distinct $0,j,k$ in the kernel, when the kernel acts on
$E_1$, the three values of a degree-six section are independent
at every three distinct points.  Thus the three-value incidence
has dimension $2+(17-3)=16$.
When the kernel acts on $E_2$, degree-three evaluation at the
three distinct points $v,v+j,v+k$ has rank three except when
\[
                              3v+j+k=0.
\]
This is a finite locus in $E_2$, since multiplication by three
is finite etale in characteristic two.  On that locus the
evaluation rank is exactly two, and otherwise it is three.
The exceptional incidence dimension is therefore
$1+(17-2)=16$, and the generic one is again sixteen.
All these proper incidences miss a nonempty open of $|N|$.
There are only finitely many differences and triples, so their
opens can be intersected.

Choose a smooth member in the resulting open.
No nontrivial translation stabilizes it, since its intersection
with every translate is finite.  Hence the restriction of the
quotient is generically injective.  It is an immersion at every
point because $\phi$ is etale.  There are at most two preimages
of every image point, and when there are two their tangent
directions are distinct.  This proves the asserted nodal
description and birationality.
The source has genus $19$ and the image has arithmetic genus
$18n+1$, so the number of nodes is $18(n-1)$.
Equivalently, summing the intersection number
$N^2=36$ over the $n-1$ nonzero kernel elements counts each
unordered pair twice.  This counting is valid also when the
etale kernel has even order; it is an equality of integers.

For completeness, norm of divisors gives the stated classes.
For $\phi_1\times1$, pullback multiplies the first class
coordinate by $n$, whereas
\[
       \phi^*\operatorname{Nm}_{\phi}(N)
           \simeq\bigotimes_{j\in\ker\phi}t_j^*N
\]
has class $(6n,3n)$.
Thus the norm class is $(6,3n)$.
For $1\times\phi_2$ the same calculation gives $(6n,3)$.
The image divisor is the norm divisor because its generic
degree is one.  Finally all conditions used are nonempty opens
in the irreducible projective space $|N|$, which proves the
last assertion about imposing further open conditions.
\end{proof}

\begin{corollary}[The mixed thin strip with ordinary degree twice an odd number]
\label{cor:DE-A3-mixed-thin-vone}
\status{Proved; complete independent proof and seed-certificate audits
by root and reduction-descent passed.}
Let $E_{\rm o}$ be ordinary and $E_{\rm s}$ supersingular over
an algebraically closed field of characteristic two.
For every odd positive integer $n$, a general smooth member of
\[
          \bigl|\OO_{E_{\rm o}}(6nO)
                     \boxtimes\OO_{E_{\rm s}}(3O)\bigr|
\]
has full stable rank $18n-1$ and relative stable rank $18n-2$.
Consequently the previously excluded degree-one product
systems with $a\equiv2\pmod4$, $b=1$, satisfy
$r_1=9a-2$.
\end{corollary}
\begin{proof}
There is an etale isogeny
$\phi_{\rm s}:E'_{\rm s}\to E_{\rm s}$ of degree $n$:
take a cyclic odd-order subgroup on $E_{\rm s}$ and the dual
of its quotient isogeny.  The source is supersingular by
isogeny invariance.  These are the same precise elliptic
isogeny inputs used in
Lemma~\ref{lem:DE-A3-mixed-odd-factor}.
The degree-six by degree-three seed
Lemma~\ref{lem:root-mixed-six-three-cartier}
gives a nonempty open of smooth curves on
$E_{\rm o}\times E'_{\rm s}$ of genus nineteen and full stable
rank seventeen.
Apply Lemma~\ref{lem:DE-A3-axis-nodal-norm} to
$1\times\phi_{\rm s}$ and intersect its nodal-image open
with that sharp open.
The resulting image has $18(n-1)$ nodes, so its full stable
rank is
\[
                         17+18(n-1)=18n-1.
\]
Here the normalization sequence adds one Frobenius-bijective
graph contribution for each node.
The norm class is $(6n,3)$.
Translation of the target puts it in the exact displayed
line-bundle class: the homomorphism associated to an ample
line bundle is an isogeny and is surjective on points over
the algebraic closure.
Openness of iterated Frobenius rank and the nonempty smooth
open in this very ample system transfer this lower bound
to a general smooth member.
The even ordinary-factor Cartier-kernel bound
Lemma~\ref{lem:DE-A3-mixed-even-bound}
gives the matching relative upper bound $18n-2$.
Subtracting the ambient stable rank one from the full stable
rank proves both equalities.
\end{proof}

\begin{remark}[The direction and remaining scope of the norm construction]
\label{rem:DE-A3-mixed-axis-scope}
\status{Proved; complete independent proof audits by root and
reduction-descent passed.}
The isogeny in the corollary acts on the supersingular factor,
and its degree is odd.
An attempted transfer by an arbitrary-degree etale isogeny on
the ordinary factor would instead produce class $(6,3n)$.
That proposed route to all classes $(6n,3)$ is invalid and is
not used here.  This particular norm argument makes no assertion
about the thin-strip cases with $4\mid a$.
Those remaining cases are subsequently proved by the uniform
residue construction in Theorem~\ref{thm:DE-A3-mixed-thin-all};
the norm statement itself retains exactly the scope proved above.
The universal seed and both retained programs were independently
read in full before use; their exact determinant certificate is
$c^5$, valid for every ordinary parameter $c\ne0$.
The dimension bounds and the two norm-class formulae above
hold uniformly, without an inference from a finite collection
of kernels or elliptic curves.
\end{remark}


\subsection{The remaining mixed thin systems by two residue obstructions}
\label{sec:DE-A3-mixed-thin-residues}

\begin{lemma}[The ordinary Frobenius algebra and its two pencils]
\label{lem:DE-A3-ordinary-two-pencils}
\status{Proved; root and reduction-descent full independent proof/source/code audits passed.}
Let $E$ be an ordinary elliptic curve over an algebraically closed
field of characteristic two, let $F:E\to E'=E^{(2)}$ be relative
Frobenius, and let $H=\OO_{E'}(qO)$, where $q\ge3$.
There is an orthogonal decomposition
\[
 F_*\OO_E=\OO_{E'}\oplus B,\qquad
 B\not\simeq\OO_{E'},\qquad
 \mu:B^{\otimes2}\xrightarrow{\sim}\OO_{E'}.
\]
Here orthogonality refers to the multiplication--Cartier pairing,
normalized to take the value one on $(1,1)$; in this decomposition
the pairing is $(a,b),(c,d)\mapsto ac+\mu bd$.
Put $V=F_*F^*H=H\oplus HB$.
One can choose $a_0\in H^0(H)$ and $b_2\in H^0(HB)$ with
reduced, disjoint zero divisors $A$ and $D$.
For every such choice the following maps are isomorphisms:
\[
\begin{split}
 H^0(H)\oplus H^0(HB)&\longrightarrow H^0(H^2),\\
 (x,y)&\longmapsto a_0x+\mu b_2y,\\
 H^0(HB)&\longrightarrow H^0(A,HB|_A),\\
 H^0(H)&\longrightarrow H^0(D,H|_D).
\end{split}
\]
\end{lemma}
\begin{proof}
Choose nonzero invariant differentials on $E,E'$.
Ordinarity says that Cartier on the invariant differential is nonzero.
After normalization it gives an $\OO_{E'}$-linear map
$\tau:F_*\OO_E\to\OO_{E'}$ with $\tau(1)=1$.
The pairing $(u,v)\mapsto\tau(uv)$ is perfect: this is also immediate
in a local $2$-basis from the Cartier formula.
For the original Cartier isomorphism and local formula see
\cite[0, Theorem~2.1.9, p.~515]{Illusie1979DRW}.
Thus $B=\ker\tau$ is a line bundle, orthogonal to the unit summand.
Every square in $F_*\OO_E$ belongs to the unit summand.
Locally a generator of $B$ therefore has square in $\OO_{E'}$,
and perfection of the pairing makes that square a unit.
This proves the stated multiplication isomorphism $\mu$.
Finally $H^0(F_*\OO_E)=k$ gives $H^0(B)=0$.
Hence $B$ is nontrivial of degree zero, and $H^1(B)=0$ as well.

The degree-$q$ line bundles $H$ and $HB$ are very ample, so general
sections have reduced disjoint zero divisors.
The first displayed map comes from the exact sequence
\[
 0\longrightarrow B
 \xrightarrow{\,t\mapsto(\mu b_2t,a_0t)\,}
 H\oplus HB
 \xrightarrow{\,(x,y)\mapsto a_0x+\mu b_2y\,}
 H^2\longrightarrow0.
\]
Its exactness follows from the absence of common zeros of $a_0,b_2$.
The vanishing of both cohomology groups of $B$ proves the first
isomorphism.  The two evaluation isomorphisms follow in the same way
from $HB(-A)=B$ and $H(-D)=B^{-1}=B$.
\end{proof}

\begin{lemma}[A nonzero stable-rank obstruction determinant]
\label{lem:DE-A3-mixed-residue-determinant}
\status{Proved; root and reduction-descent full independent proof/source/code audits passed.}
With the notation of Lemma~\ref{lem:DE-A3-ordinary-two-pencils},
there exist $a_1\in H^0(H)$ and $b_1\in H^0(HB)$ for which the
section triple
\[
 f_0=(a_0,0),\qquad f_1=(a_1,b_1),\qquad f_2=(0,b_2)
 \quad\hbox{in }H^0(V)
\]
gives multiplication--Cartier rank $6q-2$ and stable rank $6q-2$
on $E\times E_{\rm s}$, where
$E_{\rm s}:w^2+w=z^3$ and the second line bundle is $\OO(3O)$.
The corresponding section is $s=f_0+zf_1+wf_2$.
\end{lemma}
\begin{proof}
First choose $a_1,b_1$ in the nonempty open where
$(a_0,a_1)$ and $(b_1,b_2)$ are basepoint-free pencils.
These conditions say respectively that $a_1$ is nonzero on $A$
and $b_1$ is nonzero on $D$.
Write an arbitrary input triple in $H^0(V)^3$ as
$(x_i,y_i)$, with $x_i\in H^0(H)$ and $y_i\in H^0(HB)$.
On the supersingular factor the multiplication--Cartier pairing
in the basis $(1,z,w)$ is
\[
 \begin{pmatrix}0&1&z\\1&0&w\\z&w&0\end{pmatrix}.
\]
Consequently the linearized Cartier map $T_s$ has output
$g=(g_0,g_1,g_2)\in H^0(H^2)^3$ precisely when
\begin{align}
 a_0x_1+a_1x_0+\mu b_1y_0&=g_0,\label{eq:DE-residue-C0}\\
 a_0x_2+\mu b_2y_0&=g_1,\label{eq:DE-residue-C1}\\
 a_1x_2+\mu b_1y_2+\mu b_2y_1&=g_2.\label{eq:DE-residue-C2}
\end{align}
The middle equation determines $x_2,y_0$ uniquely, by the preceding
lemma.  The other two equations are the two pencil multiplication
maps, each with kernel and cokernel of dimension one.
In particular $T_s$ has kernel exactly
\[
 k\,k_a\oplus k\,k_b,\qquad
 k_a=((a_0,0),(a_1,0),(0,0)),\quad
 k_b=((0,0),(0,b_1),(0,b_2)),
\]
and its rank is $6q-2$.

Here is the precise Frobenius convention needed for stable rank.
Descent of squares gives a bijective Frobenius-semilinear map
\[
 Q:H^0(E,F^*H)\longrightarrow H^0(E',H^2).
\]
The identification $(F^*H)^{(2)}=H^2$ is canonical in the present
origin-divisor notation.  Componentwise,
$Q(a,b)=a^2+\mu b^2$.
The absolute semilinear Cartier operator is $Q^{-1}T_s$,
up to the harmless normalization of the invariant differential.
Thus its kernel meets its image trivially if and only if the map
\[
 Q(\ker T_s)\longrightarrow\coker T_s
\]
is an isomorphism.  The two output vectors to test are
\[
 Q(k_a)=(a_0^2,a_1^2,0),\qquad
 Q(k_b)=(0,\mu b_1^2,\mu b_2^2).
\]
This square-descent step prevents confusing a linear Cartier map
between Frobenius-twisted spaces with an ordinary endomorphism.

Fix a nonzero invariant differential $\omega$ on $E'$.
The two pencil cokernels, identified with $H^1(E',\OO)=k$, have
residue functionals
\[
 L_A(g)=\sum_{P\in A}\operatorname{Res}_P
              \frac{g\,\omega}{a_0a_1},\qquad
 L_D(g)=\sum_{P\in D}\operatorname{Res}_P
              \frac{g\,\omega}{\mu b_1b_2}.
\]
Indeed the two local solutions of each pencil equation differ by
the displayed fraction times its kernel generator.  This is the
two-open Cech connecting homomorphism followed by residue duality.
The original references are \cite[Theorem~2 and Theorem~3 with its
corollary, pp.~155--156; Section~4, Theorem~5 and its corollary,
pp.~157--159]{Tate1968Residues}.

In \eqref{eq:DE-residue-C1}, reduction at $A$ gives
$y_0=g_1/(\mu b_2)$, and reduction at $D$ gives $x_2=g_1/a_0$.
Subtracting these local substitutions changes the other expressions
only by functions regular at the respective residue points.
The two-obstruction matrix on $Q(k_a),Q(k_b)$ is therefore
\[
 \begin{pmatrix}t&u\\v&t\end{pmatrix},
\]
where
\[
\begin{split}
 t&=\sum_A\operatorname{Res}\frac{a_1b_1}{a_0b_2}\omega,\\
 u&=\sum_A\operatorname{Res}\frac{\mu b_1^3}{a_0a_1b_2}\omega,\\
 v&=\sum_D\operatorname{Res}\frac{a_1^3}{\mu a_0b_1b_2}\omega .
\end{split}
\]
For its lower-right entry, the same expression as $t$ is summed
over $D$.  The rational differential has poles only in $A\cup D$;
the global residue theorem and characteristic two make these sums
equal.  All displayed quotients are globally defined rational
functions: the line-bundle factors cancel using $\mu$.

Replace $b_1$ by $b_1+\lambda b_2$, with $\lambda\in k$.
The pencils remain basepoint-free, since the values on $D$ do not
change.  The entry $v$ is constant: at every simple zero of $b_2$
only the value of $b_1$ enters its residue.
The entry $t$ is constant as well, since its change is
\[
 \lambda\sum_A\operatorname{Res}\frac{a_1}{a_0}\omega=0.
\]
The entry $u$ is a polynomial of degree at most three in $\lambda$.
Its linear coefficient is
\[
 c_1=\sum_A\operatorname{Res}\frac{\mu b_1^2}{a_0a_1}\omega.
\]
For every fixed allowed $a_1$, this is a nonzero quadratic polynomial
in $b_1$.  To verify this without a genericity assumption on $E$,
use the evaluation isomorphism $H^0(HB)\simeq H^0(A,HB|_A)$.
At each reduced point of $A$ its coefficient is nonzero:
$a_0$ has a simple zero, and $a_1,\mu,\omega$ are units there.
The independent evaluation coordinates therefore exhibit $c_1$
as a nonzero sum of their squares with nonzero coefficients.

Similarly, for every fixed allowed $b_1$, the entry $v$ is a
nonzero cubic polynomial in $a_1$.  The evaluation map
$H^0(H)\simeq H^0(D,H|_D)$ is an isomorphism, and each of the
simple-pole residue coefficients on $D$ is nonzero.
Thus $v$ is a nonzero sum of cubes of independent evaluation
coordinates.  Both nonvanishing conditions define nonempty opens
in the irreducible open parameter space of allowed pairs $(a_1,b_1)$;
they have a common point.
Choose such a point.  Then $v\ne0$ and $u(\lambda)$ is nonconstant,
whereas $t$ and $v$ are constant.  Hence
\[
                  t^2+u(\lambda)v
\]
is a nonzero polynomial, and is nonzero for some $\lambda$.
The obstruction matrix is invertible for that choice.
Therefore the Cartier kernel has zero intersection with the Cartier
image.  Rank already stabilizes after the first iterate, and equals
$6q-2$, as claimed.
\end{proof}

\begin{theorem}[All mixed thin product systems]
\label{thm:DE-A3-mixed-thin-all}
\status{Proved; root and reduction-descent full independent proof/source/code audits passed.}
Let $E_{\rm o}$ be any ordinary elliptic curve and $E_{\rm s}$
any supersingular elliptic curve over an algebraically closed
field of characteristic two.  For every even $m\ge6$, a general
smooth member of
\[
 \bigl|\OO_{E_{\rm o}}(mO)\boxtimes\OO_{E_{\rm s}}(3O)\bigr|
\]
has genus $3m+1$, full stable rank $3m-1$, and relative stable
rank $3m-2$.
\end{theorem}
\begin{proof}
Write $m=2q$.  Over the algebraic closure the supersingular factor
has the displayed model used in the preceding lemma.
That lemma supplies a section for which the relative Cartier
stable rank is $3m-2$; it does not require that section itself
define a smooth curve.
The corresponding iterated-rank minor is nonzero on a nonempty
open of the complete section space.  Both line bundles have
degree at least three, so their exterior product is very ample.
Bertini supplies a nonempty smooth-member open, which intersects
the rank open.
Lemma~\ref{lem:DE-A3-mixed-even-bound} gives the matching generic
upper bound.  The divisor sequence, its stable-part exactness,
and Serre duality are precisely those proved there; the ambient
stable rank is one.  Adjunction gives genus $3m+1$.
The result is uniform in the ordinary elliptic curve and in $m$.
\end{proof}

\begin{corollary}[The mixed product formula in every degree]
\label{cor:DE-A3-mixed-product-no-exception}
\status{Proved; root and reduction-descent full independent proof/source/code audits passed.}
For every perfect field of characteristic two, every ordinary
$E_{\rm o}$ and supersingular $E_{\rm s}$, and all positive
integers $a,b,d$, the geometric generic relative stable rank in
\[
 \bigl|\OO_{E_{\rm o}}(3adO)\boxtimes
             \OO_{E_{\rm s}}(3bdO)\bigr|
\]
is
\[
            r_d=9abd^2-1-\epsilon(ad),\qquad
 \epsilon(c)=\begin{cases}1&c\ {\rm even},\\0&c\ {\rm odd}.
                         \end{cases}
\]
The relative coherent dimension is $9abd^2-1$.
\end{corollary}
\begin{proof}
Theorem~\ref{thm:DE-A3-mixed-all-product-ranks} already proves
every case except $d=b=1$ with $a$ even.
Apply Theorem~\ref{thm:DE-A3-mixed-thin-all} with $m=3a$
to the remaining cases.  Passage to the algebraic closure
computes the asserted geometric generic ranks over a perfect field.
The earlier ambient restriction bound remains applicable.
\end{proof}

\begin{remark}[Verification and scope of the residue argument]
\label{rem:DE-A3-mixed-thin-residue-checks}
\status{Proved; root and reduction-descent full independent proof/source/code audits passed.}
The entire new argument is uniform in $q$ and in the ordinary elliptic
curve.  Its two nonvanishing assertions use evaluation isomorphisms
on reduced divisors, rather than a list of elliptic parameters.
The exact symbolic program
\path{agents/obstructions_crystals/code/mixed_thin_residue_checks.py}
checks thirteen universal characteristic-two identities: the two
kernel vectors, the middle exact-sequence kernel, both regular
residue-remainder substitutions, all three terms of the cubic
variation, and the nonzero determinant coefficient $vc_1$.
The author run passed in $0.1971881189965643$ seconds, with its full
output retained beside the program.
The root and reduction-descent reviewers read the entire proof,
original source passages, and program; their author-preserving
reruns passed in $0.1764095999969868$ and
$0.19381124299980002$ seconds, respectively.
The ordinary Frobenius-algebra splitting and all global existence
steps are proved above, not inferred from that symbolic check.
This closes the thin-system exception for mixed elliptic products;
it does not supply a generic-rank formula for arbitrary ambient
varieties.
\end{remark}


\subsection{Positive ordinary index on a nonproduct mixed surface}
\label{sec:DE-A3-mixed-positive-index}

\begin{lemma}[The all-degree positive-index Prym obstruction]
\label{lem:DE-A3-positive-index-Prym}
\status{Proved; complete independent proof and original-source audits by root and reduction\_descent.}
Let $A$ be an abelian surface of $2$-rank one over an algebraically
closed field of characteristic two, let $M$ be ample, and put
$a=a(M)$ and $D=\chi(M)$ as in
Lemma~\ref{lem:DE-A3-mixed-polarization-blocks}.
If $m>0$ and $a+v_2(m)>0$, every smooth curve $Y\in|M^m|$ satisfies
\[
                         g(Y)-f_Y\ge2.
\]
In particular, if $a>0$, the generic relative stable rank for
$L=M^3$ satisfies $r_d\le9Dd^2-2$ for every $d\ge1$.
\end{lemma}
\begin{proof}
Use the integral dual polarization
$\lambda_N=D\lambda_M^{-1}:A^\vee\to A$ from
Lemma~\ref{lem:DE-A3-mixed-polarization-blocks}.
Its local-local index is $b(N)=a$.
The Picard immersion $j:A^\vee\hookrightarrow J_Y$ and its
polarized dual $q:J_Y\to A$ have the properties established in
Lemma~\ref{lem:DE-A3-mixed-all-even-Prym}: the Prym $P=\ker q$
is smooth and connected, and $qj=m\lambda_N$.
That proof, including its integral lattice and finite-flat
subgroup arguments, was reread before this application.

Let $G_{\rm ll}$ be the actual height-two local-local factor of
$A^\vee[2^\infty]$.
Both elementary divisors of $m\lambda_N$ on its Dieudonne lattice
have valuation $a+v_2(m)>0$.
Hence $m\lambda_N$ kills the finite group scheme $G_{\rm ll}[2]$,
and $j$ embeds this group into $P[2]$.
Twist that finite-group immersion by Frobenius and apply functorial
Verschiebung.  The connected group
$\ker(V:G_{\rm ll}^{(2)}\to G_{\rm ll})$, of order two, lies in
$G_{\rm ll}^{(2)}[2]$ and embeds in $\ker V_P$.
It follows that $\dim P-f_P\ge1$.
The quotient $J_Y/P\simeq A$ has defect $\dim A-f_A=1$.
Dimensions and $2$-ranks add in this exact sequence of abelian
varieties, so $g(Y)-f_Y\ge2$.
Adjunction gives $g(Y)=Dm^2+1$, and subtracting the ambient
stable rank one gives the last assertion.
\end{proof}

\begin{lemma}[Exact descent through the ordinary etale kernel]
\label{lem:DE-A3-positive-index-descent}
\status{Proved; complete independent proof and original-source audits by root and reduction\_descent.}
Under the same hypotheses, suppose $a>0$ and put $n=2^a$.
There are an etale cyclic isogeny $\pi:A\to B$ of degree $n$
and an ample line bundle $H$ on $B$ such that
\[
 M\simeq\pi^*H,\qquad a(H)=0,\qquad
 \chi(H)=D/n,\qquad b(H)=b(M).
\]
For every smooth $C\in|H^m|$, its pullback $Y=\pi^{-1}C$ is
smooth and connected, and
\[
                   g(Y)-f_Y=n\bigl(g(C)-f_C\bigr).
\]
\end{lemma}
\begin{proof}
The canonical ordinary etale factor of $\ker\lambda_M$ is the
constant cyclic group $K\simeq\mathbf Z/n\mathbf Z$.
This follows on the actual $2$-divisible group from the integral
ordinary block in
Lemma~\ref{lem:DE-A3-mixed-polarization-blocks}.
Take the quotient isogeny $\pi:A\to B=A/K$.
It is etale, since its finite-flat kernel is etale.
Existence of this quotient and duality of its kernel are
\cite[Theorem~4.38 and Example~4.40, pp.~65--66;
Theorem~7.5, p.~100]{EdixhovenVanDerGeerMoonenAV}.

Here the descent of the actual line bundle can be made explicit.
Choose a generator $P\in K(k)$ and an isomorphism
$\theta:t_P^*M\xrightarrow{\sim}M$, which exists because
$P\in\ker\lambda_M$.
Its $n$-fold composite is multiplication by a scalar $c\in k^\times$.
Rescale $\theta$ by an $n$th root of $c^{-1}$.
The new composite is the identity, so its powers give a genuine
$K$-linearization, without averaging or inverting $n$.
Faithfully flat descent gives a line bundle $H$ with
$M\simeq\pi^*H$; it is ample because its pullback is ample.
This also verifies directly the isotropic descent criterion in
\cite[Corollary~8.11 and proof, p.~117]
{EdixhovenVanDerGeerMoonenAV} in the present cyclic case.

The polarization identity is
$\lambda_M=\pi^\vee\lambda_H\pi$.
The map $T_2(\pi)$ has cokernel of length $a$.
The dual kernel is $K^D=\mu_n$, so $T_2(\pi^\vee)$ is an isomorphism.
Thus $a(M)=a(H)+a$, proving $a(H)=0$.
An etale isogeny is an isomorphism on the local-local factors,
which proves $b(H)=b(M)$.
Riemann--Roch under a finite etale isogeny gives
$\chi(M)=n\chi(H)$.

The pullback of a smooth curve is smooth because $\pi$ is etale.
It is connected: it is an ample divisor on $A$, and the divisor
sequence together with
$H^0(A,M^{-m})=H^1(A,M^{-m})=0$ gives $H^0(Y,\OO_Y)=k$.
The antiample $H^1$ vanishing is the same abelian-variety
vanishing already used in the Picard-immersion proof.
Thus $Y\to C$ is a connected etale cyclic $2$-power cover.
The precise primary Deuring--Shafarevich input is
\cite[Theorem~1.5, Corollaries~1.7--1.8, pp.~34--36,
especially equation~(1.8.2), p.~36]{Crew1984EtaleP}.
It gives $f_Y-1=n(f_C-1)$.
Etale Hurwitz gives $g(Y)-1=n(g(C)-1)$, proving the formula.

For completeness the required unramified $p$-rank formula has
a short direct proof here.
The sheaf $\pi_*\mathbf F_2$ on $C$ has geometric fibre
$\mathbf F_2[K]=\mathbf F_2[t]/(t^n)$, with its augmentation
filtration having $n$ trivial one-dimensional quotients.
The corresponding sheaf filtration has constant quotients
$\mathbf F_2$.
Since a finite etale map has no higher direct images on these
sheaves, additivity of Euler characteristic gives
$\chi(Y,\mathbf F_2)=n\chi(C,\mathbf F_2)$.
For a smooth proper connected curve, Artin--Schreier theory gives
$H^i_{\et}(-,\mathbf F_2)=0$ for $i>1$ and
$\dim H^1_{\et}(-,\mathbf F_2)=f$; thus its Euler characteristic
is $1-f$.
One may verify the vanishing directly using surjectivity of
$F-1$ on the finite-dimensional coherent $H^1(\OO)$ over the
algebraically closed field.
This is the same sheaf-filtration proof supplied in
\cite[Section~3, first alternative proof, p.~42;
Lemma~1.9.2, p.~37]{Crew1984EtaleP}.
\end{proof}

\begin{theorem}[Bounded defects and the first positive ordinary index]
\label{thm:DE-A3-positive-index-bounds}
\status{Proved; complete independent proof and original-source audits by root and reduction\_descent.}
Let $A,M,a,D$ be as above, with $a>0$, and put $L=M^3$.
For every $d\ge1$, the geometric generic relative stable rank
in $|L^d|$ satisfies
\[
 \begin{array}{ll}
 9Dd^2-2^a\le r_d\le9Dd^2-2,&d\ \text{odd},\\[2pt]
 9Dd^2-2^{a+1}\le r_d\le9Dd^2-2,&d\ \text{even}.
 \end{array}
\]
If $a=1$ and $d$ is odd, this is the exact equality
$r_d=9Dd^2-2$.
These statements hold geometrically over every perfect field
of characteristic two and do not assume a product decomposition.
\end{theorem}
\begin{proof}
Work over an algebraic closure and descend $M$ to $H$ by
Lemma~\ref{lem:DE-A3-positive-index-descent}.
Theorem~\ref{thm:DE-A3-mixed-nonproduct-ranks} applies to $(B,H)$,
because $a(H)=0$.
It supplies a smooth $C\in|H^{3d}|$ of defect one for odd $d$
and defect two for even $d$.
Its connected etale pullback has defect $2^a$ or $2^{a+1}$,
respectively.
It is already a smooth member of $|M^{3d}|$, so openness of
iterated Frobenius rank transfers its rank as a lower bound to
the geometric generic member of that complete system.
The upper bound is Lemma~\ref{lem:DE-A3-positive-index-Prym}.
When $a=1$ and $d$ is odd the bounds agree.
All constructions were made after algebraic closure, as required
for the asserted geometric generic ranks.
\end{proof}

\begin{remark}[The remaining deformation problem]
\label{rem:DE-A3-positive-index-gate}
\status{Proved; complete independent proof and original-source audits by root and reduction\_descent.}
The etale pullback construction has the exact defect stated above.
For $a\ge2$ it therefore does not prove defect two in odd degree;
even for $a=1$ it gives defect four in even degree.
These are limitations of this construction, not counterexamples
to the expected sharp upper bound.
No existence of an index-zero ample class on the original surface
has been assumed: the class $H$ lives on the quotient surface $B$.
In particular, this argument does not circumvent a possible
Neron--Severi rank-one obstruction to decomposing a polarization
into index-zero ample summands.
\end{remark}


% Post-CP16 reduction; complete independent root and reduction_descent audits passed.
\subsection{The toral first-image lattice at positive ordinary index}

\begin{lemma}[A toral lattice for the unresolved rank problem]
\label{lem:DE-A3-positive-index-image-candidate}
\status{Proved; complete independent proof and source audits passed.}
Let $A$ be a mixed abelian surface over an algebraically closed field
of characteristic two. Let $N$ be ample and suppose its ordinary
polarization index is positive. Put $b=h^0(A,N)$ and
$P=\mathbf P(H^0(A,N))$. Then $b$ is even. There is a normalized
toral invariant derivation $e$, acting on $N$, for which
\[
 W=H^0(A,N)=W_0\oplus W_1,\qquad
 e|W_i=i,\qquad\dim W_0=\dim W_1=b/2.
\]
Write $\mathcal B=W^*\otimes\mathcal O_P(-1)$ and
\[
 Z=\mathbf P(W_0)\cup\mathbf P(W_1),\qquad
 \mathcal I=\ker\bigl(\mathcal B
       \xrightarrow{(s,es)}\mathcal O_P^2\bigr).
\]
Outside $Z$, the sheaf $\mathcal I$ is locally free of rank $b-2$.
If $b/2\ge2$, its reflexive extension has
$\det\mathcal I=\mathcal O_P(-b)$ and $H^0(P,\mathcal I)=0$.
The full coherent Frobenius maps into $\mathcal I$ and preserves it.
When $N=M^{3d}$, the expected sharp relative rank $b-2$ is equivalent
to nonvanishing of the determinant of
$F_P^*\mathcal I\to\mathcal I$, a homogeneous form of degree $b$.
\end{lemma}
\begin{proof}
Positive ordinary index means that the actual finite-flat subgroup
$\mu_2$ in the ordinary multiplicative factor of $A[2]$ is contained
in $\ker\lambda_N$. Its restriction of the commutator pairing is
trivial: the associated homomorphism
$\mu_2\to\mu_2^D=\mathbf Z/2$ is zero.
Thus the isotropic descent theorem applies to
$\pi:A\to B=A/\mu_2$; compare
\cite[Corollary~8.11 and proof, p.~117]
{EdixhovenVanDerGeerMoonenAV}.
Equivalently, the restricted theta group is a commutative
multiplicative-type extension of $\mu_2$ by $\mathbf G_m$.
On character groups its extension splits because the quotient
character group is $\mathbf Z$, which is free.
The resulting splitting supplies an actual $\mu_2$-linearization.
Hence $N=\pi^*H$ for an ample $H$ on $B$.

The torsor algebra is
$\pi_*\mathcal O_A=\mathcal O_B\oplus Q$, where $Q$ is the
nontrivial two-torsion line bundle corresponding to the nonzero
character of $\mu_2$. It is nontrivial because
$H^0(A,\mathcal O_A)=k$. Its multiplication identifies $Q^2$ with
$\mathcal O_B$. The infinitesimal generator of this grading is
the unique nonzero invariant toral vector field $e$ with $e^2=e$.
It acts on the pulled-back bundle by the descended connection.
The two section spaces are $H^0(B,H)$ and $H^0(B,H\otimes Q)$.
Ample abelian-variety vanishing and Riemann--Roch make their
dimensions equal to $\chi(H)=b/2$.

Choose an invariant volume and rational two-basis so that
$e=x\partial_x$ and the volume is $(dx/x)\wedge dy$, as in the
mixed Cartier calculation already proved. Use a frame of $N$
pulled back from $B$; it is killed by $e$.
Cartier extracts the even $x$-powers and odd $y$-powers.
Writing a section coefficient as
$f=a^2+xb^2+yc^2+xyd^2$, one obtains
\[
 C(f^2\omega)=0,\qquad C(f\,e(f)\omega)=0.
\]
The second identity also follows by decomposing $f$ into its two
toral weights: the cross term has weight one and is killed by
Cartier, and the square term is killed because $C(\omega)=0$.
Thus, in the root-removed convention,
$M(s)s=M(s)es=0$. The coherent Frobenius matrix is $M(s)^t$,
so its image annihilates $s,es$, as asserted.

For $s=s_0+s_1$, the vectors $s,es$ are independent precisely
when both $s_0,s_1$ are nonzero. This proves the rank and
codimension assertions. The depth lemma proves reflexivity of
the kernel; the determinant and absence of global sections follow
from its displayed presentation. Since the full coherent model
has the same unit part as the relative model (the ambient top
coherent Frobenius is zero), its stable rank is $b-2$ exactly when
the induced Frobenius on this rank-$b-2$ lattice is generically
bijective. Its determinant lies in
$\operatorname{Hom}(\mathcal O_P(-2b),\mathcal O_P(-b))$,
giving the degree-$b$ form. Nonvanishing is the remaining issue,
not an assertion of this lemma.
\end{proof}

\begin{remark}[The exact first-variation test]
\label{rem:DE-A3-positive-index-variation-candidate}
\status{Proved reduction; determinant nonvanishing remains unproved.}
At a point outside $Z$ where the induced map on $\mathcal I$
has corank exactly two, let $K$ and $Q$ be its two-dimensional
kernel and cokernel. A tangent direction $t$ gives
\[
 D_t:K\longrightarrow Q.
\]
In a local frame the lowest possible term of the determinant
along that direction is a nonzero fixed scalar times
$\det D_t$. Let $P(s)$ be the rectangular inclusion matrix of a
local frame of $\mathcal I$ into $\mathcal B$, and choose a local
left inverse $L(s)$, so that $L(s)P(s)=1$.
The Frobenius matrix is $L(s)M(s)^tP(s)^{[2]}$. Over dual numbers the
derivative of the last factor is zero; the derivative of the
first factor vanishes on $K$. Consequently $D_t$ is obtained
by restricting $M(t)^t$ and projecting to $Q$.
This is a precise sufficient test for a deformation to the
sharp-rank open. The next lemma proves the required first-corank
statement for an etale double pullback of a sharp positive-index
source. Nonvanishing of $\det D_t$ remains unproved in general.
\end{remark}

\begin{lemma}[The restricted first corank at every sharp double pullback]
\label{lem:DE-A3-positive-index-double-corank-candidate}
\status{Proved; complete independent proof and source audits passed.}
Let $\pi:A\to B$ be an etale double isogeny of mixed abelian
surfaces. Let $H$ be ample of positive ordinary index, put
$N=\pi^*H$, and use compatible toral lifts on these bundles.
Suppose that a section $s\in H^0(B,H)$ makes Frobenius on
$I_B=\ker(s,es)\subset H^2(B,H^{-1})$ bijective.
Set $s_A=\pi^*s$ and $b=h^0(A,N)$.
Then on $I_A=\ker(s_A,es_A)$ the first Frobenius rank and stable
Frobenius rank are both $b-4$. Its kernel and cokernel have
dimension two. The full coherent first rank need not equal $b-4$.

For this point, the first-variation map $D_t$ of the preceding
remark vanishes whenever $t$ is pulled back from $H^0(B,H)$.
Consequently it factors through the surjective trace map
\[
 H^0(A,N)\longrightarrow H^0(B,H).
\]
Thus the remaining sufficient sharpness test is the nonvanishing
of a quadratic determinant of a linear map
$H^0(B,H)\to\operatorname{Hom}(k^2,k^2)$.
\end{lemma}
\begin{proof}
Write $V_B=H^2(B,H^{-1})$ and $V_A=H^2(A,N^{-1})$.
First, $s$ and $es$ are independent. Indeed, on the dense
complement of the two weight eigenspaces the full Frobenius has
rank at most $b_B-2$, where $b_B=h^0(B,H)$.
The same inequality holds everywhere by vanishing of its
$(b_B-1)$-minors. At a weight eigenpoint, $I_B$ would have
dimension at least $b_B-1$, contradicting the assumed bijectivity
of its restriction. Pullback preserves this independence, so
both section points lie outside the corresponding bad loci $Z$.
The etale double-torsor algebra has the exact sequence
\[
 0\longrightarrow\mathcal O_B
   \longrightarrow\pi_*\mathcal O_A
   \xrightarrow{\operatorname{Tr}}\mathcal O_B
   \longrightarrow0.
\]
Locally it is the Artin--Schreier algebra
$R[z]/(z^2+z-f)$; its trace sends $1$ to zero and $z$ to one.
Tensoring by $H^{-1}$ and taking cohomology gives
\[
 0\longrightarrow V_B\xrightarrow{\pi^*}V_A
   \xrightarrow{\operatorname{Tr}}V_B\longrightarrow0,
\]
because $H^1(B,H^{-1})=0$ and the surface has no third coherent
cohomology. Frobenius with multiplication by the pulled-back
section commutes with both maps: locally this uses
$\operatorname{Tr}(u^2)=\operatorname{Tr}(u)^2$.

Serre duality identifies trace with the dual of pullback on
sections. Since $es_A=\pi^*(es)$, it follows that
$I_A=\operatorname{Tr}^{-1}I_B$ and there is an exact sequence
of semilinear Frobenius spaces
\[
 0\longrightarrow V_B\xrightarrow{\pi^*}I_A
   \xrightarrow{\operatorname{Tr}}I_B\longrightarrow0.
\]
All kernel-cokernel arguments may equivalently be made after
linearizing these maps from the Frobenius pullbacks of their
domains. Since $k$ is perfect this preserves the dimensions and
exactness; in the first-variation formula the kernel is understood
in this linearized sense.
By hypothesis $F|I_B$ is bijective. Since the full Frobenius on
$V_B$ maps into $I_B$, its image is exactly $I_B$ and its kernel
$K_B$ has dimension two. Also $I_B\cap K_B=0$.
The kernel-cokernel sequence for this diagram gives
\[
 \ker(F|I_A)=\pi^*K_B,\qquad
 \operatorname{coker}(F|I_A)
       \simeq\operatorname{coker}(F|V_B).
\]
These are both two-dimensional.
If an element of $\ker(F|I_A)$ also lies in the image, write it
as $F(x)=\pi^*k$, where $k\in K_B$.
Taking trace and using bijectivity on $I_B$ gives
$\operatorname{Tr}x=0$, so $x=\pi^*v$.
Then $k=F(v)$ lies in $K_B\cap I_B=0$.
Thus the first image is already Frobenius-bijective.
Its dimension is $\dim I_A-2=b-4$.

It remains to verify the assertion on invariant tangent directions.
For $t=\pi^*u$ and $k\in K_B$, the first-variation class is the
class of $\pi^*F_u(k)$, where $F_u$ denotes the full coherent
Frobenius with multiplication by $u$.
Under the preceding cokernel identification it is the class of
$F_u(k)$ in $V_B/I_B$.
The two functionals defining that quotient are $s,es$.
Let $\beta$ be the root-removed Cartier multiplication pairing.
Symmetry and the toral identity give
\[
 \beta(s,u)=\beta(u,s),\qquad
 \beta(es,u)=\beta(s,eu).
\]
The second identity follows in a descended local frame from
$C(e(f)\omega)=0$, since toral weight-one terms have zero Cartier
image. Because $k\in\ker F_s$, it annihilates every
$\beta(s,v)$; hence both $s$ and $es$ annihilate $F_u(k)$.
The first-variation class is therefore zero.

Finally the torsor trace sequence tensored by $H$, together with
$H^1(B,H)=0$, gives
\[
 0\to H^0(B,H)\xrightarrow{\pi^*}H^0(A,N)
       \xrightarrow{\operatorname{Tr}}H^0(B,H)\to0.
\]
This proves the factorization of the linear first-variation map.
No assertion that its quadratic determinant is nonzero follows
from this factorization alone.
\end{proof}


% Promoted after complete independent root and reduction_descent proof/source audits. No determinant-nonvanishing claim.
\subsection{A rank-two bundle model for the remaining trace variation}

\begin{lemma}[The two toral bundles differ by an acyclic line bundle]
\label{lem:DE-A3-positive-two-bundles}
\status{Proved; complete independent root and reduction/descent proof and source audits passed.}
Let $B$ be a mixed abelian surface over an algebraically closed
field $k$ of characteristic two and let
$H=M^{3d}$ be ample with positive ordinary index. Put
$T=B^{(2)}$, $L=H^{(2)}$ and $b=h^0(B,H)$.
There are rank-two bundles $E_0,E_1$ on $T$ with
\[
 F_{B/k,*}H=E_0\oplus E_1,\qquad
 E_1\simeq E_0\otimes P,
\]
where $P$ is the nontrivial two-torsion point of $\Pic^0(T)(k)$.
The Cartier pairing is orthogonal for this decomposition and
restricts to perfect alternating pairings with values in $L$.
In particular $\det E_i\simeq L$. Moreover
\[
 H^q(T,E_i)=0\ (q>0),\quad h^0(T,E_i)=b/2,\quad
 c_2(E_i)=b/2,\quad H^q(T,P)=0\ \text{for all }q.
\]
Both $E_i$ are globally generated.
\end{lemma}
\begin{proof}
The previously proved integral mixed decomposition gives
$\ker F_{B/k}\simeq\mu_2\times\alpha_2$.
Factor relative Frobenius as
\[
 B\xrightarrow{\rho}C=B/\mu_2\xrightarrow{\nu}T.
\]
The second map is an $\alpha_2$-torsor.
Since $3$ is odd, positive ordinary index of $M^{3d}$ implies
that $M^d$ descends along $\rho$, by the actual theta-group
linearization already proved. Thus $H=\rho^*J$ with $J$ a
cube of an ample line bundle. The torsor algebra is
$\rho_*\mathcal O_B=\mathcal O_C\oplus Q$ for the nontrivial
two-torsion point $Q\in\Pic^0(C)(k)$, and therefore
\[
 E_0=\nu_*J,\qquad E_1=\nu_*(J\otimes Q).
\]
The dual of $\nu$ has kernel $\alpha_2$. Consequently
$\nu^*: \Pic^0(T)(k)\to\Pic^0(C)(k)$ is bijective.
The unique preimage $P$ of $Q$ is nontrivial and has order two:
$\nu^*(P^2)=\mathcal O_C$, and the kernel has only the identity
as a geometric point. Projection formula gives $E_1=E_0P$.

Finite Frobenius duality, or its elementary local two-basis
calculation, makes $C(fg\omega)$ a perfect pairing on $F_*H$
with values in $L$. Its diagonal is zero because
$C(\omega)=0$ on a mixed surface.
The toral operator is self-adjoint for this pairing:
$C(e(fg)\omega)=0$ gives $\beta(ef,g)=\beta(f,eg)$.
Its different eigenspaces are therefore orthogonal; each
restriction is perfect alternating of rank two. This proves
the determinant assertions.

Ample abelian-variety cohomology vanishing on $C$ gives the
asserted higher vanishing for $E_i$. Riemann--Roch and the
degree-two map $\rho$ give $h^0(E_i)=\chi(J)=b/2$.
On an abelian surface the tangent bundle is trivial, so
Riemann--Roch for a rank-two bundle gives
$\chi(E_i)=c_1(E_i)^2/2-c_2(E_i)$. Since $c_1(E_i)=c_1(L)$
and $L^2/2=b$, this proves $c_2(E_i)=b/2$.
The acyclicity of $P$ is
\cite[Lemma~7.19 and its full proof, p.~105]
{EdixhovenVanDerGeerMoonenAV}.
Finally $J$ and $JQ$ are cubes of ample line bundles, after
choosing a cube root of $Q$ in $\Pic^0(C)(k)$, and so are
very ample. They separate every length-two subscheme,
including every fibre of the finite flat map $\nu$.
Thus the global evaluations of $\nu_*J$ and $\nu_*(JQ)$
are surjective at every point.
\end{proof}

\begin{lemma}[An acyclic summand and a determinant curve]
\label{lem:DE-A3-positive-determinant-curve}
\status{Proved; complete independent root and reduction/descent proof and source audits passed.}
Keep the preceding notation. A general pair
$s_0\in H^0(E_0)$, $s_1\in H^0(E_1)$ has disjoint finite
local-complete-intersection zero schemes $Z_0,Z_1$, each
of length $b/2$.
Normalize the isomorphism $E_1=E_0P$ to identify its alternating
pairing with that on $E_0$, using $P^2=\mathcal O_T$.
Then
\[
 K=\ker\bigl(E_0\oplus E_1
       \xrightarrow{\beta(s_0,-)+\beta(s_1,-)}L\bigr)
\]
is locally free of rank three. It contains the line subbundle
$P\to K$, $a\mapsto(as_1,as_0)$.
The quotient $Q_s=K/P$ has rank two, determinant $LP$, and
the same cohomology as $K$.

Put $w=s_0\wedge s_1\in H^0(LP)$. For general pairs $w\ne0$.
The two sections $(s_0,0)$ and $(0,s_1)$ induce an exact sequence
\[
 0\longrightarrow\mathcal O_T^2\longrightarrow Q_s
      \longrightarrow i_*R_s\longrightarrow0,
 \qquad i:D_s=(w=0)\hookrightarrow T,
                                                        \tag{DE-TB1}
\]
where $R_s$ is an invertible sheaf on the Cartier divisor $D_s$.
The arithmetic genus of $D_s$ is $b+1$, and
$\chi(D_s,R_s)=0$, so $\deg R_s=b$.
No reducedness of the zero schemes or of $D_s$ is assumed here.

If $s=s_0+s_1$ is a sharp source in the preceding variation
lemma, then $h^0(Q_s)=h^1(Q_s)=2$ and $h^2(Q_s)=0$.
\end{lemma}
\begin{proof}
Global generation gives the usual zero incidence over $T$.
For a general section its zero scheme has dimension zero;
the two local equations are a regular sequence on the smooth
surface and its length is $c_2(E_i)=b/2$.
Requiring simultaneous vanishing of $s_0,s_1$ imposes four
independent conditions over a two-dimensional base, so the
general two zero schemes are disjoint. The contraction map
onto $L$ is then surjective. The displayed map from $P$ lies
in its kernel by symmetry of the two cross terms in
characteristic two, and is a subbundle because $s_0,s_1$
never vanish together. This proves local freeness of $Q_s$.
Determinants give $\det Q_s=LP$; acyclicity of $P$ gives the
cohomology assertion.

At a fixed point, the two section evaluations can be chosen
independently, so $w$ is not identically zero for a general
pair. The determinant of the two induced sections of $Q_s$
is $w$, with the compatible determinant identifications.
One may check this in a frame in which
$s_0=(1,0)$ and $s_1=(f,w)$. If the chosen frame $e$ of $P$
satisfies $e^2\mapsto u\in\mathcal O_T^\times$ under the
fixed trivialization of $P^2$, the two columns in $Q_s$ become
$(1,0)$ and $(u f^2,w)$. It is not necessary, and need not
be possible etale locally, to make $u$ a square.
Their determinant is still $w$. The chart with $s_1$ nonzero treats
the other points. Along $D_s$ the two columns always have
rank one, since the original sections do not vanish together.
Thus their cokernel is a line bundle on $D_s$, proving
(DE-TB1), even if the divisor is nonreduced.

Adjunction gives $p_a(D_s)=1+(LP)^2/2=b+1$.
The contraction sequence gives
$\chi(K)=\chi(E_0)+\chi(E_1)-\chi(L)=0$.
Then $\chi(Q_s)=0$ and (DE-TB1) gives $\chi(R_s)=0$.
Riemann--Roch for a line bundle on a projective Cartier
curve gives its asserted degree.
Finally sharpness says that global Cartier contraction has
kernel and cokernel of dimension two. Higher cohomology
of $E_i,L$ vanishes; the long exact sequence therefore gives
$h^0(K)=h^1(K)=2$, $h^2(K)=0$. The same holds for $Q_s$.
\end{proof}

\begin{lemma}[The actual torsor class in the contraction boundary]
\label{lem:DE-A3-positive-torsor-boundary}
\status{Proved; complete independent root and reduction/descent proof and source audits passed.}
For an etale double cover $\pi':\widetilde T\to T$, let
$\alpha\in H^1(T,\mathcal O_T)$ be the extension class of
$0\to\mathcal O_T\to\pi'_*\mathcal O_{\widetilde T}
\xrightarrow{\operatorname{Tr}}\mathcal O_T\to0$.
Choose a trace lift $q_i$ of $s_i$ in
$H^0(T,E_i\otimes\pi'_*\mathcal O_{\widetilde T})$.
The contraction $\beta(s_i,q_i)$ descends to a section
$v_i\in H^0(T,L I_{Z_i})$. Its boundary in
\[
 0\longrightarrow\mathcal O_T\xrightarrow{s_i}E_i
       \xrightarrow{\beta(s_i,-)}L I_{Z_i}
       \longrightarrow0
\]
is exactly $\alpha$.
\end{lemma}
\begin{proof}
Trace commutes with contraction, and
$\operatorname{Tr}\beta(s_i,q_i)=\beta(s_i,s_i)=0$.
Thus contraction descends. Its vanishing on $Z_i$ is local.
Choose local trace-one generators $z_a$ of the torsor algebra,
with $z_b-z_a=c_{ab}$ representing $\alpha$.
Write $q_i=z_a s_i+r_a$ in each local trivialization.
Then $\beta(s_i,q_i)=\beta(s_i,r_a)$, and
$r_b-r_a=c_{ab}s_i$ (signs agree in characteristic two).
These are precisely the Cech representatives of the
connecting homomorphism. The resulting class is $\alpha$.
Existence of the trace lift follows from $H^1(T,E_i)=0$.
\end{proof}

\begin{remark}[The remaining step]
\label{rem:DE-A3-positive-bundle-scope}
\status{Proved scope statement; residue nonvanishing remains unproved.}
These statements give actual bundles and a Cartier curve on
which a residue calculation could be made. No formula has yet
identified the quadratic variation determinant with a
nonzero residue pairing on $(D_s,R_s)$. In particular neither
acyclicity of $P$ nor the degree of $R_s$ proves sharpness.
The exact outstanding assertion remains exclusion of a common
kernel or common image for the trace-direction matrix space.
\end{remark}


% Promoted after complete independent root and monodromy audits. Sharpness remains open.
\subsection{Two zero trace directions and the remaining matrix alternative}

\begin{lemma}[The two toral trace directions vanish]
\label{lem:DE-A3-positive-trace-radical}
\status{Proved; complete independent root and monodromy proof audits passed.}
Keep the hypotheses of
Lemma~\ref{lem:DE-A3-positive-index-double-corank-candidate}.
Write $s=s_0+s_1$ for its two toral weight components and put
$R=\langle s_0,s_1\rangle=\langle s,es\rangle$.
The descended first-variation map annihilates $R$. Consequently
the quadratic determinant in that lemma is the determinant of
a linear map
\[
       H^0(B,H)/R\longrightarrow\operatorname{Hom}(k^2,k^2).
\]
This assertion does not prove that the determinant is nonzero.
\end{lemma}
\begin{proof}
We use the Cartier-dual description of the already proved
kernel--cokernel diagram. All vector spaces and maps in this
argument can be taken with their relative Frobenius twists.
Since the field is perfect, identifying those twists preserves
the kernels, cokernels, trace maps and zero identities used below.
Let $W=H^0(B,H)$, $\widetilde W=H^0(A,\pi^*H)$, and let
$C_s$ denote Cartier multiplication by $s$. The hypotheses give
\[
 \ker C_s=R,\qquad W=R\oplus\im C_s.
\]
Indeed Cartier is dual to the full coherent Frobenius, and its
induced map on $W/R$, dual to the Frobenius on $I_B$, is bijective.

The trace sequence is equivariant for the compatible toral
operators. Since the operator satisfies $e^2=e$, its two
eigenspace functors are exact. Choose $q_i\in\widetilde W_i$
with $\operatorname{Tr}q_i=s_i$, for $i=0,1$.
Cartier commutes with etale trace. The toral pairing identities
give $C_s(s_i)=0$, hence
\[
             C_{\pi^*s}(q_i)=\pi^*v_i
                 \quad\text{for some }v_i\in W.
\]
Choose $w_i\in W$ so that
$v_i+C_s(w_i)\in R$, and put $q_i'=q_i+\pi^*w_i$.
Thus $C_{\pi^*s}(q_i')\in\pi^*R$.
The classes of $q_0',q_1'$ in
$\widetilde W/\pi^*R$ form a basis of the kernel of its induced
Cartier map: they are independent because their traces are
$s_0,s_1$, and that kernel has dimension two by the preceding
lemma. They are therefore a basis of the dual of the cokernel
of the restricted coherent Frobenius.

Let $k\in K_B=\ker F_s$. By Serre duality, the matrix coefficient
of the first variation in a direction $t\in\widetilde W$,
tested against $\pi^*k$ and $q_j'$, is
\[
       \bigl\langle k,
          C_B\bigl(\operatorname{Tr}(tq_j')\,\omega_B\bigr)
       \bigr\rangle.
                                                        \tag{DE-TR1}
\]
The formula is simply the dual of restricting $F_t$ and then
passing to its cokernel; trace is dual to pullback. It also
follows by commuting relative Cartier with the etale trace.

Take $t=q_i$. The correction $\pi^*w_j$ in $q_j'$ contributes
$\beta(s_i,w_j)$ inside the pairing in (DE-TR1). This lies in
$\im C_s$: if $(w_j)_i$ is the weight-$i$ component, then
\[
                \beta(s_i,w_j)=\beta(s,(w_j)_i).
\]
It is annihilated by $k$. The remaining term is zero as well.
For $i\ne j$, the product $q_iq_j$ has toral weight one and
its Cartier image is zero. For $i=j$, the square identity
$C(q_i^2\omega_A)=q_iC(\omega_A)=0$ applies, since $A$ is mixed.
Thus (DE-TR1) vanishes for every $k$ and both $j$.
The descended variation in trace direction $s_i$ is zero.
The two directions span $R$, which proves the assertion.
\end{proof}

\begin{lemma}[The exact rank-one matrix alternative]
\label{lem:DE-A3-positive-matrix-ruling}
\status{Proved; complete independent root and monodromy proof audits passed.}
Let $k$ be an infinite field and $T:U\to\operatorname{Mat}_2(k)$
a linear map. Its determinant is identically zero if and only if
either every matrix in its image kills one fixed nonzero vector,
or all of their images lie in one fixed line.
In particular, for the preceding variation problem, nonvanishing
is equivalent to excluding both of these common annihilators.
Neither exclusion is established here.
\end{lemma}
\begin{proof}
The two stated alternatives imply the determinant vanishes.
Conversely, the zero matrix space satisfies either alternative.
Otherwise choose a nonzero image matrix. Since its determinant
is zero it has rank one, and changes of source and target bases
make it $A=\operatorname{diag}(1,0)$.
Write any other image matrix as $B=(b_{ij})$.
The identity $\det(A+tB)=0$ for all $t\in k$ gives $b_{22}=0$;
the identity $\det B=0$ then gives $b_{12}b_{21}=0$.
The image vector space is therefore the union of its two linear
subspaces defined by $b_{12}=0$ and $b_{21}=0$.
A vector space over an infinite field cannot be the union of
two proper linear subspaces. In the first case all matrices kill
the second basis vector; in the second case all images lie in
the first coordinate line. This proves the equivalence.
\end{proof}

\begin{remark}[Retained finite probe and its limitation]
\label{rem:DE-A3-positive-trace-probe-scope}
\status{Proved under stated hypotheses}
The program
\path{agents/obstructions_crystals/code/post_CP17_variation_weight_probe_v2.py}
checks one sharp characteristic-two product source. It gives
two toral weight spaces of dimension nine, variation-image
dimensions one and two on those spaces, and zero variation on
$s$ and $es$. Its predecessor and every earlier probe output are
retained. The v2 author run took $0.0163224049974815$ seconds.
These computations motivated the first lemma; its
proof is the trace and Cartier argument above. They do not
exclude either matrix alternative on an arbitrary mixed surface.
The earlier attempted factorization through the source Cartier
cokernel was refuted by the retained finite probe and is not
used: $W/R$ here is a different quotient.
\end{remark}


% Promoted after complete independent root and reduction/descent proof audits. No Petri/variation nonvanishing assertion.
\subsection{A norm pencil on the determinant curve}

\begin{lemma}[The determinant cokernel is residual to a squared pencil]
\label{lem:DE-A3-positive-pencil-candidate}
\status{Proved; complete independent root and reduction/descent proof audits passed.}
Keep the algebraically closed characteristic-two hypotheses and notation of
Lemma~\ref{lem:DE-A3-positive-determinant-curve}, and take a general pair
with disjoint zero schemes. On $D=D_s$ the images of $s_0$ and $s_1$
generate a line subbundle $R_0\subset E_0|_D$, after viewing the latter
section as a map $P^{-1}|_D\to E_0|_D$. Then
\[
 R_0\simeq\mathcal O_D(Z_0),\qquad
 R_0P\simeq\mathcal O_D(Z_1),\qquad
 P|_D\simeq\mathcal O_D(Z_1-Z_0).
\]
In particular $Z_0,Z_1$ are effective Cartier divisors on $D$, including
when $D$ is nonreduced. Put $A=R_0^{\otimes2}$. The squares of the two
sections, using the fixed trivialization of $P^2$, give a base-point-free
pencil $\langle a_0,a_1\rangle\subset H^0(D,A)$ with divisors $2Z_0,2Z_1$.
The line bundle in (DE-TB1) is
\[
                         R_s\simeq\omega_D\otimes A^{-1}.
                                                        \tag{DE-TP1}
\]
Here $\deg A=b$ and $p_a(D)=b+1$.
\end{lemma}
\begin{proof}
At every point of $D$, at least one of $s_0,s_1$ is nonzero, and their
determinant vanishes. In a frame with the nonvanishing section first,
their matrix has a unit entry and determinant zero. Its image on $D$
is therefore a line subbundle; this constructs $R_0$ even over the
nonreduced local rings of $D$. Near $Z_0$, use $s_1$ to trivialize
the common line. The two equations of $Z_0$ on the smooth surface can
then be written $(f,w)$, where $w$ cuts out $D$. They form a regular
sequence. Consequently $f$ is a nonzerodivisor on $D$ and cuts out
$Z_0$ there. The same argument treats $Z_1$. The asserted line-bundle
identifications follow from their sections. Their squares have no common
zero, since $Z_0\cap Z_1=\varnothing$.

Tensor (DE-TB1) by $\mathcal O_D$. Since $D$ is Cartier, its initial
Tor term is $R_s\otimes\mathcal O_D(-D)$, and hence
\[
 R_s(-D)=\ker\bigl(\mathcal O_D^2\longrightarrow Q_s|_D\bigr).
\]
The earlier local chart has columns $(1,0)$ and $(u f^2,w)$,
where $e^2\mapsto u$ for the chosen frame $e$ of $P$.
On $D$ their kernel is the kernel of the pencil row $(1,u f^2)$.
Globally this row is $(a_0,a_1):\mathcal O_D^2\to A$, whose kernel
is $A^{-1}$. Thus $R_s(-D)=A^{-1}$. Adjunction on the abelian surface
identifies $\mathcal O_D(D)$ with $\omega_D$, proving (DE-TP1).
The degree statements follow from $\operatorname{length}Z_i=b/2$ and
the previously computed arithmetic genus. No reducedness or smoothness
of $D$ was used.
\end{proof}

\begin{lemma}[Norm factorization and the actual ordinary differential]
\label{lem:DE-A3-positive-norm-pencil-candidate}
\status{Proved; complete independent root and reduction/descent proof audits passed.}
With the same general pair, the norm of the $\alpha_2$-quotient
$\nu:C\to T$ gives a section
\[
                   \sigma\in H^0(D,R_sP)
\]
which is a nonzerodivisor on $D$. If $S_i\in H^0(T,L)$ is the
Frobenius-twisted section corresponding to $s_i$, then
\[
                         S_i|_D=\sigma a_i.             \tag{DE-TP2}
\]
There is a canonical isomorphism
\[
 \frac{H^0(T,L)}{\langle S_0,S_1\rangle}
 \simeq
 \frac{H^0(D,\omega_DP)}{\sigma\langle a_0,a_1\rangle}.
                                                        \tag{DE-TP3}
\]
The fixed trivialization of $P^2$ also defines a nonzero invariant
differential $\eta_P\in H^0(T,\Omega_T^1)$ with
$C(\eta_P)=\eta_P$. Wherever $a_0a_1$ is invertible on $D$,
\[
                      d\log(a_1/a_0)=\eta_P|_D.         \tag{DE-TP4}
\]
This identity does not assert that its right side is nonzero on $D$.
\end{lemma}
\begin{proof}
Locally the finite flat $\alpha_2$-algebra is
$\mathcal O_T[z]/(z^2-h)$. In a trivialization of $J$, its norm is
\[
                         n(a+bz)=a^2+h b^2.
\]
Determinants of the transition matrices identify its target with
$\det\nu_*J=L$; the determinant of $\nu_*\mathcal O_C$ is trivial,
as the torsor algebra is an extension of $\mathcal O_T$ by itself.
Thus this construction gives the intrinsic quadratic norm
$n:E_0\to L$. For $E_1=E_0P$, the same norm is multiplied by the
local unit $u$ specifying $P^2=\mathcal O_T$. This is precisely the
descent of the square of the original section through relative
Frobenius, so $n(s_i)=S_i$.

The restriction of the norm to the line $R_0$ is a homomorphism
$A=R_0^2\to L|_D$. It is therefore a section of
$L|_D\otimes A^{-1}=R_sP$, using $\omega_D=(LP)|_D$.
This is $\sigma$, and homogeneity gives (DE-TP2). For a general
pair, the divisors of the sections of $J$ and $JQ$ on $C$ have no
common component: both line bundles are globally generated, and
one may choose the second section to avoid the finitely many
components of the first. Their norm divisors on $T$ likewise have
no common component, since $\nu$ is a homeomorphism on underlying
spaces. If $\sigma$ belonged to a minimal prime of $D$, (DE-TP2)
would put both norm sections in the corresponding prime. This is
impossible. The Cartier divisor $D$ is Cohen--Macaulay and has no
embedded associated primes, so $\sigma$ is a nonzerodivisor.

The restriction sequence is
\[
                 0\longrightarrow P\longrightarrow L
                    \longrightarrow L|_D\longrightarrow0.
\]
Acyclicity of $P$ identifies $H^0(T,L)$ with $H^0(D,L|_D)$.
As $L|_D=\omega_DP$, equation (DE-TP2) gives (DE-TP3).

Finally, write $e^2\mapsto u$ in any local frame of $P$.
On overlaps $u$ changes by the square of a transition unit,
so $d\log u$ glues to a global differential $\eta_P$.
The usual Cartier identity gives $C(d\log u)=d\log u$.
If this differential were zero, the Cartier kernel calculation
on the smooth perfect-base surface would make every $u$ locally
a square. The uniquely determined square roots then normalize
all transition functions to one, trivializing $P$, a contradiction.
Every global differential on an abelian variety is invariant.
In the determinant chart the pencil ratio is exactly $u f^2$,
so its logarithmic differential is $d\log u$. This proves
(DE-TP4), with no assumption that $u$ itself has a local square root.
\end{proof}

\begin{lemma}[The entire Cartier contraction is a curve extension]
\label{lem:DE-A3-positive-pencil-contraction-candidate}
\status{Proved; complete independent root and reduction/descent proof audits passed.}
The two bundles are orthogonal also in cohomology:
\[
                  \operatorname{Ext}^q_T(E_0,E_1)=0
                         \quad\text{for every }q.
                                                        \tag{DE-TP5}
\]
There is an exact sequence of bundles on $D$
\[
 0\longrightarrow A\longrightarrow G_s:=E_0|_D\otimes R_0
                  \longrightarrow L|_D\longrightarrow0,
                                                        \tag{DE-TP6}
\]
whose global-section map, under the natural identifications
\[
 H^0(D,G_s)=H^0(T,E_0\oplus E_1),\qquad
 H^0(D,L|_D)=H^0(T,L),
\]
is exactly the Cartier contraction $C_s$. Moreover $H^1(D,G_s)=0$.
Thus, for a sharp source,
\[
                  h^0(D,A)=h^1(D,A)=h^0(D,R_s)=2,
\]
and the two pencil sections form a basis of $H^0(D,A)$.

Let $\beta\in H^1(T,\mathcal O_T)$ be the class of the
$\alpha_2$-torsor $\nu$. If $\partial$ is the boundary of
(DE-TP6), then
\[
                    \partial(\sigma a)=\beta|_D\cup a
                         \quad(a\in H^0(D,A)).          \tag{DE-TP7}
\]
For a sharp source this cup-product map $H^0(D,A)\to H^1(D,A)$
is an isomorphism. The scalar Frobenius twists in identifying
the original kernel sections with their squared pencil sections
do not affect this isomorphism assertion.
\end{lemma}
\begin{proof}
The self-fibre product of the $\alpha_2$-torsor has square-zero
diagonal ideal isomorphic to $\mathcal O_C$. Consequently
$\nu^*\nu_*J$ has a two-step filtration with both quotients $J$.
Projection formula and finite flat base change identify the cohomology
computing (DE-TP5) with that of
\[
                  (\nu^*\nu_*J)^\vee\otimes JQ.
\]
This bundle has a filtration with both quotients $Q$. Since the
nontrivial degree-zero line bundle $Q$ is acyclic on $C$, the
whole cohomology vanishes. This proves (DE-TP5).

The two original sections give an exact sequence on the surface
\[
 0\longrightarrow\mathcal O_T\oplus P\longrightarrow E_0
       \longrightarrow i_*\bigl(L|_D\otimes R_0^{-1}\bigr)
       \longrightarrow0.
\]
For the cokernel identification, its Tor kernel on $D$ is
$P R_0^{-1}$, and tensoring back by $\mathcal O_D(D)=\omega_D$
gives $L|_D R_0^{-1}$. Dualizing the surface sequence gives
\[
 0\longrightarrow E_0^\vee\longrightarrow\mathcal O_T\oplus P
        \longrightarrow i_*(R_0P)\longrightarrow0.
\]
Here the last map is the row $(s_1,s_0)$, with the fixed
$P^2$ identification understood. Tensor this sequence by $E_1$.
Its left term is $E_0^\vee E_1$, acyclic by (DE-TP5), and the
other terms are $E_1\oplus E_0$ and $i_*G_s$. This gives the
claimed global-section identification and $H^1(G_s)=0$.
The inclusion of $R_0$ in $E_0|_D$, tensored by $R_0$, gives
(DE-TP6). Under the preceding map its quotient row is
$\beta(s_0,-)+\beta(s_1,-)$: this can be read in the determinant
chart, or directly by taking the wedge with the common line.
It is exactly $C_s$.

The long cohomology sequence therefore identifies $H^0(A)$
with $\ker C_s$ and $H^1(A)$ with $\operatorname{coker}C_s$.
At a sharp source these have dimension two. The two known
kernel vectors are $s_0,s_1$, and under the same identifications
they are $a_0,a_1$. Serre duality on the Gorenstein Cartier curve,
together with (DE-TP1), gives $h^0(R_s)=h^1(A)=2$.

For the last assertion choose local torsor coordinates $z_i$
with $z_j-z_i=c_{ij}$, representing $\beta$. Multiplication by
$z_i$ is a local endomorphism of $E_0=\nu_*J$; this statement
is independent of the local trivialization of $J$.
If $r$ trivializes $R_0$ on $D$, then
\[
                   r\wedge(z_i r)=n(r)=\sigma r^2.
\]
Thus, for $a=f r^2$, the section $f(z_i r)\otimes r$ of $G_s$
lifts $\sigma a$. Its transition difference is $c_{ij}a$,
which proves (DE-TP7). Finally, first rank equal to stable rank
$b-2$ is equivalent to the Frobenius-twisted kernel span
$\langle S_0,S_1\rangle$ mapping isomorphically to the cokernel
of $C_s$. By (DE-TP2) and (DE-TP7), this is precisely the
invertibility of cup product by $\beta$ on $H^0(A)$.
\end{proof}

\begin{corollary}[An exact formula for the full first rank of a pullback]
\label{cor:DE-A3-positive-first-rank-candidate}
\status{Proved; complete independent root and reduction/descent proof audits passed.}
Let $s$ be a sharp source, and let $\pi':\widetilde T\to T$ be the
etale double cover occurring in the earlier variation reduction, of
class $\alpha\in H^1(T,\mathcal O_T)$. Then the full coherent
Cartier contraction at the pulled-back section has kernel dimension
\[
        2+\dim\ker\bigl(\alpha|_D\cup-:
                     H^0(D,A)\longrightarrow H^1(D,A)\bigr).
                                                        \tag{DE-TP8}
\]
Its first rank is accordingly
$2b-2-\dim\ker(\alpha\cup-)$, which may be $2b-4$, $2b-3$,
or $2b-2$. This statement concerns the full contraction;
the previously proved first and stable ranks of its restriction
to the toral image lattice remain $2b-4$.
\end{corollary}
\begin{proof}
All determinant-curve constructions commute with etale base change.
The pullback of $E_0^\vee E_1$ remains acyclic: pushforward along
$\pi'$ is an extension of two copies of this acyclic bundle.
The same argument applies to $P$, and the higher cohomology
of the pulled-back $E_i$ vanishes. Thus the proof of (DE-TP6)
identifies the pulled-back contraction kernel with
$H^0(\widetilde D,\pi'^*A)$.
The torsor trace sequence, tensored by $A$ and restricted to $D$,
has connecting homomorphism $\alpha|_D\cup-$. Its first three
cohomology terms give
\[
 0\longrightarrow H^0(D,A)\longrightarrow
 H^0(\widetilde D,\pi'^*A)\longrightarrow
 \ker\bigl(\alpha\cup-:H^0(D,A)\to H^1(D,A)\bigr)
 \longrightarrow0.
\]
The preceding lemma gives $h^0(A)=2$, proving (DE-TP8).
The full source and target spaces each have dimension $2b$.
\end{proof}

\begin{remark}[What this model has not established]
\label{rem:DE-A3-positive-norm-pencil-scope-candidate}
\status{Proved scope statement; the quadratic variation remains unresolved.}
The two lemmas identify the quotient domain of the variation with an
explicit twisted canonical space on a Cartier curve and retain both
the norm factor and the ordinary logarithmic differential. They do
not identify the matrix-valued variation with a Petri or residue map.
The preceding two additional statements now prove $h^0(A)=2$
for a sharp source and identify the two torsor cup products exactly.
They still give no smoothness of $D$, no generic separability of its
pencil, and no lower bound of three on the variation-image dimension.
The retained finite probe in fact has total image dimension two:
the one-dimensional first-weight image is contained in the
two-dimensional second-weight image, and their displayed generators
have an invertible sum. The common-kernel/common-image exclusion
for arbitrary mixed surfaces remains the open step. The later
trace-zero theorem, Proposition~\ref{prop:DE-A3-positive-cup-trace-zero-candidate},
reduces a possible failure to one nilpotent matrix line; it does
not supply the remaining image-dimension bound.
\end{remark}


% Promoted after full independent root and reduction_descent proof/code audits; no symmetry or nonvanishing assertion.
\subsection{The variation as a derivative of the trace-lift matrix}

\begin{lemma}[A typed derivative identity]
\label{lem:DE-A3-positive-cup-derivative-candidate}
\status{Proved; complete independent root and reduction/descent proof and code audits passed.}
Keep the sharp-source and etale-double-cover hypotheses of
Lemma~\ref{lem:DE-A3-positive-index-double-corank-candidate}.
Write the source section parameters as squares, and denote their
square roots by $u$. Use compatible bases on the relative Frobenius
twists. There is a rational $2\times2$ matrix $A(u)$ on the sharp
source open, obtained by the actual etale trace lifts and projection
along the Cartier image, with the following property: the root-removed
Cartier-dual trace-direction variation is exactly
\[
                         D^{\vee,\mathrm{root}}_v=dA_u(v).
                                                        \tag{DE-CD1}
\]
The columns in (DE-CD1) are indexed by the corrected trace lifts,
which are dual to the coherent cokernel; the rows are indexed by the
dual kernel projection. Thus (DE-CD1) represents the Cartier-dual map
$Q^*\to K^*$. In conventional column-vector bases the actual coherent
map $K\to Q$ has matrix $(dA_u(v)^{[2]})^t$, and the actual section
direction is the square of $v$. In particular, trace-zero of every variation is equivalent
to $d\operatorname{tr}A=0$. Over a perfect ground field this says that
$\operatorname{tr}A$ is a square in the source parameter function field.
This derivative identity alone asserts neither condition;
Proposition~\ref{prop:DE-A3-positive-cup-trace-zero-candidate}
below subsequently proves the stronger identity $\operatorname{tr}A=0$.
Determinant nonvanishing remains open.
\end{lemma}
\begin{proof}
We retain all twists in the construction. Let $U$ be the vector space
of square-root section coordinates and write $W=U^{(2)}$ for the
ordinary section-coordinate space. Root-removed Cartier multiplication
is a symmetric alternating bilinear map
\[
             B:U\otimes U\longrightarrow W.
\]
The alternating assertion is the mixed-volume identity already proved.
Both $U$ and $W$ have compatible decompositions into the two toral
weights. The map $B$ is zero between different weights. Put
$M(u)v=B(u,v)$, write $u_i$ for the weight-$i$ part of $u$, and set
\[
 r(u)=(u_0,u_1),\qquad R(u)=(u_0^{[2]},u_1^{[2]}).
\]
Sharpness gives the direct sum $W=\im M(u)\oplus\im R(u)$.
Let $P(u):W\to k^2$ be its projection in the displayed $R$ basis.
Thus
\[
                       PM=0,\qquad PR=1.                \tag{DE-CD2}
\]
All these matrices are rational functions on the sharp open.

Choose a toral-compatible linear trace splitting
$q:U\to\widetilde U$ of the etale cover; use its Frobenius twist
$q^{(2)}:W\to\widetilde W$ on ordinary coordinates. The deck
involution sends $qv$ to $qv+\pi^*v$, and similarly on $W$.
Cartier and trace commute. There are therefore bilinear maps
$H,J,K:U\otimes U\to W$ such that
\[
 \begin{aligned}
 \widetilde B(\pi^*u,qv)&=\pi^*H(u,v)+q^{(2)}B(u,v),\\
 \widetilde B(qu,qv)&=\pi^*J(u,v)+q^{(2)}K(u,v).
 \end{aligned}
                                                        \tag{DE-CD3}
\]
Equivariance under the deck involution, followed by expansion of
both sides, gives the exact identity
\[
                  K(u,v)=H(u,v)+H(v,u)+B(u,v).            \tag{DE-CD4}
\]
The maps $H$ and $B$ vanish on different toral weights.

Define the trace-lift coefficient matrix by
\[
                          A(u)=P(u)H(u)r(u),
                                                        \tag{DE-CD5}
\]
where $H(u)v=H(u,v)$. This matrix is independent of the chosen trace
splitting. Indeed replacing $q$ by $q+\pi^*\Lambda$ changes $H$ to
\[
 H'(u,v)=H(u,v)+B(u,\Lambda v)+\Lambda^{[2]}B(u,v).
\]
On $v=u_i$, the last term is zero and the middle term lies in
$\im M(u)$, so $P H'(u)r=P H(u)r$. A toral-compatible splitting
may therefore be used for the proof even when a different linear
splitting is used to compute $A$.
Choose, locally on the sharp open, columns
$w_i\in U$ with
\[
                 H(u,u_i)+M(u)w_i=R(u)A(u)_i.            \tag{DE-CD6}
\]
Such columns exist by (DE-CD2). The square-root coordinates
$q u_i+\pi^*w_i$ are precisely the corrected trace lifts used in
Lemma~\ref{lem:DE-A3-positive-trace-radical}: their ordinary
coordinates have trace $u_i^{[2]}$, and their Cartier images belong
to $\pi^*\im R$. The root-removed Cartier-dual variation in trace direction $v$
therefore has column
\[
                  P\bigl(K(v,u_i)+B(v,w_i)\bigr).
                                                        \tag{DE-CD7}
\]
This formula is the earlier trace-dual kernel--cokernel formula,
with its semilinear types retained.

Differentiate (DE-CD2) with respect to the root parameters $u$.
In characteristic two $dR=0$, so
\[
                  (dP)R=0,\qquad(dP)M=P\,dM.
\]
Using (DE-CD6) gives
\[
 \begin{aligned}
 dA_u(v)_i
 &=P B(v,w_i)+P H(v,u_i)+P H(u,v_i).
 \end{aligned}
                                                        \tag{DE-CD8}
\]
Toral orthogonality gives $H(u,v_i)=H(u_i,v)$ and
$B(v,u_i)=B(u,v_i)\in\im M(u)$, the last term being killed by $P$.
Substitution of (DE-CD4) into (DE-CD7) now proves (DE-CD1), column by
column. The identities also show that choices of the correction
columns do not change the induced variation.

Taking traces yields the asserted scalar criterion. The sharp open
is an open subset of the section affine space, whose function field
is a rational function field over a perfect field. In such a field
the kernel of the absolute differential is exactly the subfield of
squares: expand in the basis of square-free monomials over the square
subfield and differentiate. This proves the final equivalence.
\end{proof}

\begin{remark}[Diagnostic evidence and remaining scope]
\label{rem:DE-A3-positive-cup-derivative-scope-candidate}
\status{Proved scope statement; the later trace-zero theorem supersedes the symmetry gate.}
The retained programs
\path{agents/obstructions_crystals/code/post_CP19_canonical_cup_variation_probe.py},
\path{agents/obstructions_crystals/code/post_CP19_positive_local_index_cup_probe.py}
and
\path{agents/obstructions_crystals/code/post_CP19_GF4_canonical_variation_probe.py}
use the actual Cartier-image complement, trace lifts and Frobenius
linearization. Their author runtimes were respectively
$0.022789923001255374$, $0.10455728400120279$ and
$0.16540274800354382$ seconds. The first two retain twenty examples
each, of source types $(6,3)$ and $(6,6)$, and the last retains twelve
examples over $\mathbf F_4$. All sampled canonical variation matrices
have trace zero. The second family has positive local polarization
index. These finite computations neither establish the square
condition on arbitrary mixed surfaces nor exclude a common
annihilator. The subsequent acyclic-factorization argument in
Proposition~\ref{prop:DE-A3-positive-cup-trace-zero-candidate}
proves trace zero independently of these samples. The lemma itself is a linear-algebra identity, not an
inference from those samples.
A separate direct evaluation over $\mathbf F_4[\epsilon]/(\epsilon^2)$
checks (DE-CD1) in all $216$ coordinate directions of the twelve retained
sources. It evaluates the actual Cartier products and inverts the
projection matrix over dual numbers, without using the differentiated
projection formula. The retained program
\path{agents/obstructions_crystals/code/post_CP19_GF4_cup_derivative_check.py}
passed in $1.857703956993646$ seconds.
\end{remark}


% Promoted after complete independent root and reduction_descent proof/source/code audits. No variation-image lower bound is asserted.
\subsection{An acyclic alternating factorization and trace-zero variation}

\begin{lemma}[An acyclic factorization of the two-section alternating map]
\label{lem:DE-A3-positive-acyclic-alternating-candidate}
\status{Proved; complete independent root and reduction/descent proof, source and code audits passed.}
Keep the algebraically closed characteristic-two hypotheses and the
general disjoint-zero-scheme construction of
Lemma~\ref{lem:DE-A3-positive-determinant-curve}. Write
$V=E_0$, $E_1=VP$, $Q=Q_s$, and
$a:\mathcal O_T^2\to Q$ for the two-section map in (DE-TB1).
Let
\[
 \mathfrak g=\bigl(\operatorname{End}^0(V)/\mathcal O_T\bigr)P,
 \qquad \operatorname{End}^0(V)=\ker(\operatorname{tr}).
\]
This is an acyclic rank-two bundle. The trace pairing descends to a
perfect alternating pairing on $\mathfrak g$ with values in
$\mathcal O_T$. Write $J_{\mathfrak g}:\mathfrak g^\vee\to\mathfrak g$
for its inverse, and put $J_2=\left(\begin{smallmatrix}0&1\\1&0\end{smallmatrix}\right)$.
There is a map $f:\mathfrak g\to Q$ for which
\[
                  aJ_2a^\vee=fJ_{\mathfrak g}f^\vee:
                              Q^\vee\longrightarrow Q.       \tag{DE-TZ1}
\]
In particular the map on $H^1$ induced by either side is zero.
\end{lemma}
\begin{proof}
The trace map $\operatorname{End}(V)\to\mathcal O_T$ is locally
surjective. In rank two and characteristic two its scalar subbundle
lies in its trace-zero kernel. Thus both sequences
\[
\begin{split}
0&\longrightarrow\operatorname{End}^0(V)P
 \longrightarrow\operatorname{End}(V)P\xrightarrow{\operatorname{tr}}P
 \longrightarrow0,\\
0&\longrightarrow P\longrightarrow\operatorname{End}^0(V)P
 \longrightarrow\mathfrak g\longrightarrow0
\end{split}
\]
are exact sequences of vector bundles. By (DE-TP5),
$\operatorname{End}(V)P=V^\vee E_1$ is acyclic; so is $P$.
The two long cohomology sequences prove that $\mathfrak g$ is acyclic.
Locally a trace-zero matrix has the form
$\left(\begin{smallmatrix}c&b\\d&c\end{smallmatrix}\right)$.
Modulo scalar matrices the trace pairing is
$(b,d),(b',d')\mapsto bd'+db'$. It is alternating and perfect.
The fixed trivialization of $P^2$ gives the claimed untwisted pairing.

For $\phi\in\operatorname{End}^0(V)P$, set
\[
                         \phi\longmapsto(\phi s_1,\phi s_0)
                                     \in E_0\oplus E_1.
\]
The two factors of $P$ in the first component are contracted by the
fixed $P^2$ trivialization. The rank-two alternating identity
\[
 \beta(s_0,\phi s_1)+\beta(s_1,\phi s_0)
                          =\operatorname{tr}(\phi)(s_0\wedge s_1)
\]
shows that the image lies in $K$. Scalar endomorphisms give precisely
the existing map $P\to K$. Consequently this map descends to
$f:\mathfrak g\to Q$.

We verify (DE-TZ1) retaining the nonsquare unit in the earlier chart.
Choose a frame in which $s_0=(1,0)$ and $s_1=(h,w)e$, with
$e^2\mapsto u\in\mathcal O_T^\times$ for the local frame of $P$.
The coordinates of $Q$ used in the determinant-curve proof give
\[
 a=\begin{pmatrix}1&u h^2\\0&w\end{pmatrix},\qquad
 f=\begin{pmatrix}uw&0\\0&1\end{pmatrix}.
\]
The second matrix uses the off-diagonal classes $(b,d)e$ as a frame
of $\mathfrak g$. Its trace pairing has matrix $uJ_2$, not $J_2$;
therefore $J_{\mathfrak g}=u^{-1}J_2$ in this frame. Direct multiplication gives
\[
                  aJ_2a^t=wJ_2=f(u^{-1}J_2)f^t.
\]
No square root of $u$ was chosen. The chart with $s_1$ nonzero gives
the same equality; alternatively the equality of the two intrinsic
bundle maps on the dense first chart extends over the integral
smooth surface. This proves (DE-TZ1). Its right side factors through
the acyclic bundle $\mathfrak g$, giving the final assertion.
\end{proof}

\begin{proposition}[Trace-zero for the actual torsor cup matrix]
\label{prop:DE-A3-positive-cup-trace-zero-candidate}
\status{Proved; complete independent root and reduction/descent proof, source and code audits passed.}
Assume in addition that $s$ is sharp. Let $\beta\in H^1(T,\mathcal O_T)$
be the class of the $\alpha_2$ quotient in (DE-TP7), and let
$\alpha\in H^1(T,\mathcal O_T)$ be the class of the connected
etale double cover in the variation construction. Then
\[
 T_s=(\beta\cup-)^{-1}(\alpha\cup-):
                       H^0(D,A)\longrightarrow H^0(D,A)
                         \quad\hbox{satisfies}\quad
                            \operatorname{tr}T_s=0.           \tag{DE-TZ2}
\]
Consequently the rational trace-lift matrix $A(u)$ of (DE-CD5)
has trace zero, and every actual trace-direction variation has
trace zero in the compatible kernel and cokernel identifications.
\end{proposition}
\begin{proof}
Dualizing (DE-TB1) on the smooth surface and using
$R_s=\omega_D A^{-1}$ gives
\[
 0\longrightarrow Q^\vee\xrightarrow{a^\vee}\mathcal O_T^2
                         \longrightarrow i_*A\longrightarrow0.
                                                               \tag{DE-TZ3}
\]
The last map is the squared-pencil row, up to interchanging its two
basis vectors. This harmless interchange preserves $J_2$.
At a sharp source its two sections form a basis of $H^0(D,A)$.
Thus the image
\[
 S=\operatorname{im}\bigl[H^1(Q^\vee)
                              \longrightarrow H^1(\mathcal O_T)^2\bigr]
\]
is the kernel of the ensuing map to $H^1(D,A)$. That ensuing map
sends $(\gamma_0,\gamma_1)$ to $\gamma_0\cup a_0+\gamma_1\cup a_1$.
Cup product on $H^1(\mathcal O_T)$ is perfect alternating; this is
the exterior-algebra statement
\cite[Proposition~7.27, p.~108, with Corollary~6.13 and its proof,
pp.~93--94]{EdixhovenVanDerGeerMoonenAV}.
Equip $H^1(\mathcal O_T)^2$ with the tensor pairing of that cup
product and $J_2$. By naturality of Serre duality, its restriction
to $S$ is computed by the map
$H^1(aJ_2a^\vee):H^1(Q^\vee)\to H^1(Q)$.
Lemma~\ref{lem:DE-A3-positive-acyclic-alternating-candidate} makes
this map zero. Hence $S$ is isotropic.

The two classes $\alpha,\beta$ are independent: $F\alpha=\alpha\ne0$
whereas $F\beta=0$ and $\beta\ne0$. Nontriviality follows from
connectedness of the etale isogeny and from smoothness of the
nontrivial $\alpha_2$-torsor, respectively. They form a basis of
the two-dimensional $H^1(\mathcal O_T)$. The already proved
invertibility of $\beta\cup-$ identifies $S$ with the graph
\[
                  v\longmapsto\alpha\otimes v
                                      +\beta\otimes T_s v.
\]
The scalar $\alpha\cup\beta\in H^2(\mathcal O_T)=k$ is nonzero.
Isotropy is therefore the matrix identity
\[
                         J_2T_s+T_s^tJ_2=0.
\]
For a two-by-two matrix its left side is
$(\operatorname{tr}T_s)J_2$. This proves (DE-TZ2).

Here is also the exact comparison with the trace-lift matrix, so
that no unproved pairing identification is used. The two norm sections
restrict to $\sigma a_i$, whose boundaries in (DE-TP6) are
$\beta\cup a_i$. A trace lift $q_i$ of the original section has
contraction $v_i=\beta(s_i,q_i)$. On a cover with trace-one generators
$z_j$, write $q_i=z_j s_i+r_{i,j}$. The image of $r_{i,j}s_i$ in
$G_s$ lifts $v_i|_D$; the difference on overlaps is exactly
$c_{jk}a_i$, where $c_{jk}$ represents $\alpha$.
Thus its boundary is $\alpha\cup a_i$. Projection of $v_i$ along
the Cartier image into the norm-section basis is precisely (DE-CD5).
This is the matrix of $T_s$, transported through the compatible
relative-Frobenius identifications used there. In particular its
trace vanishes; taking a coefficient Frobenius twist or its inverse
preserves this zero identity. The derivative identity (DE-CD1) now
gives trace zero for the root-removed Cartier-dual variation, and
the coefficient square and transpose give the same conclusion for
the coherent variation.
\end{proof}

\begin{corollary}[A single remaining singular matrix line]
\label{cor:DE-A3-positive-tracezero-matrix-line-candidate}
\status{Proved; complete independent root and reduction/descent proof, source and code audits passed.}
At a sharp source the trace-direction variation takes values in
$\mathfrak{sl}_2(k)$. If its determinant is identically zero, its
image is either zero or one line generated by a rank-one nilpotent
matrix. In particular an image of dimension at least two would
prove the desired determinant nonvanishing. No such uniform lower
bound is asserted here.
\end{corollary}
\begin{proof}
Suppose a linear subspace of trace-zero matrices consists entirely
of singular matrices and contains a nonzero element $N$. Cayley--Hamilton
gives $N^2=0$, so conjugation puts it in the form
$\left(\begin{smallmatrix}0&1\\0&0\end{smallmatrix}\right)$.
Write another element as
$B=\left(\begin{smallmatrix}a&b\\c&a\end{smallmatrix}\right)$.
The identity
\[
                 \det(N+tB)=tc+t^2(a^2+bc)=0
                              \quad\text{for every }t\in k
\]
first gives $c=0$ and then $a=0$. Therefore $B$ is a multiple of
$N$. This proves the assertion. A one-dimensional image generated
by an invertible trace-zero matrix can of course also have a nonzero
determinant; the dimension-at-least-two condition is sufficient,
not necessary.
\end{proof}

\begin{remark}[Scope and exact local certificate]
\label{rem:DE-A3-positive-tracezero-scope-candidate}
\status{Proved scope statement; the variation-image lower bound remains open.}
The trace identity is derived from an actual acyclic bundle and Serre
duality, independently of the retained finite product probes. It
does not prove that the variation image has dimension at least two,
or exclude a single common nilpotent line. Thus uniform sharpness
at arbitrary positive ordinary polarization index remains open.
The local matrix identities and their characteristic-two trace
calculation are checked symbolically by the retained program
\path{agents/obstructions_crystals/code/post_CP20_trace_zero_factorization.py}.
Its author run passed in $0.05211363000125857$ seconds, checking
twenty universal scalar identities and ninety independent pairs of
nonzero nilpotent trace-zero matrices over $\mathbf F_4$.
\end{remark}


% Equivalent reduction promoted after complete root and reduction_descent audits. Original candidate retained unchanged.
\subsection{A determinant-variation expression for the remaining rank map}

\begin{lemma}[The actual symplectic cohomology sequence]
\label{lem:DE-A3-positive-monad-note}
\status{Proved equivalent reduction; complete independent root and reduction/descent proof audits passed.}
Keep the notation of the acyclic-factorization lemma and put
$\mathcal M=\mathcal O_T^2\oplus\mathfrak g^\vee$. There is an exact
sequence of bundles
\[
0\longrightarrow Q^\vee\xrightarrow{i=(a^\vee,f^\vee)}
\mathcal M\xrightarrow{p=(aJ_2,fJ_{\mathfrak g})}Q
\longrightarrow0.                                      \tag{DE-MN1}
\]
The composite is zero by (DE-TZ1). Its middle bundle is fixed as
the sections $s_0,s_1$ vary. Its rank-two image is Lagrangian for
the alternating pairing $\mathcal M\to\mathcal M^\vee$
with matrix $\operatorname{diag}(J_2,J_{\mathfrak g})$.
At a sharp source it gives
\[
0\longrightarrow H^1(Q^\vee)\longrightarrow
H^1(\mathcal O_T)^2\longrightarrow H^1(Q)\longrightarrow0.
                                                               \tag{DE-MN2}
\]
Also $\operatorname{coker}(f)=i_*A$, canonically up to the fixed
determinant identifications.
\end{lemma}
\begin{proof}
In the earlier chart, $a$ has columns $(1,0),(uh^2,w)$ and $f$
has columns $(uw,0),(0,1)$. Their images together generate $Q$
even on $D$. Since $J_2,J_{\mathfrak g}$ are invertible, $p$ is
surjective. Its dual, composed with the perfect alternating form on
$\mathcal M$, identifies its rank-two kernel with the displayed
image of $i$; the composite is already zero, and this identification
can equally be checked in the same chart. This proves exactness
and the Lagrangian assertion.

The cokernel of $f$ is locally the first-coordinate line on $D$.
The two sections induced by $a$ are the row $(1,uh^2)$, precisely
the squared pencil. This identifies that line with $A$ on overlaps.
As $\mathfrak g$ is acyclic, $H^j(Q)=H^j(D,A)$.
At a sharp source these spaces have dimension two for $j=0,1$
and vanish for $j=2$. The map $H^0(\mathcal M)=k^2\to H^0(Q)$
is an isomorphism: multiplication by $J_2$ only interchanges the
two known basis sections. Its Serre dual makes
$H^2(Q^\vee)\to H^2(\mathcal M)$ an isomorphism.
The long cohomology sequence then gives (DE-MN2).
\end{proof}

\begin{lemma}[Determinant variation gives the cohomological tangent map]
\label{lem:DE-A3-positive-monad-variation-note}
\status{Proved equivalent reduction; complete independent root and reduction/descent proof audits passed.}
On the sharp disjoint-zero-scheme parameter open, let
$S_s\subset H^1(\mathcal O_T)^2$ be the Lagrangian subspace in
(DE-MN2). For a section direction $t=(t_0,t_1)$ put
\[
 \dot w=s_0\wedge t_1+t_0\wedge s_1\in H^0(T,LP).
\]
Under (DE-MN2), the derivative of the subspace $S_s$ is the map
\[
 H^1(Q^\vee)\xrightarrow{H^1(\dot w J_Q)}H^1(Q),        \tag{DE-MN3}
\]
where $J_Q:Q^\vee\to Q\otimes(\det Q)^{-1}$ is the canonical
inverse alternating identification. All bundle determinants here
use $\det Q=LP$.
Equivalently, after the already specified graph identifications
using the torsor classes $\alpha,\beta$, this is the derivative
of the actual cup matrix $T_s$.
\end{lemma}
\begin{proof}
For an infinitesimal deformation of a Lagrangian subbundle $i_s$
of the fixed bundle $\mathcal M$, its derivative is represented
by $p\dot i:Q^\vee\to Q$. This map is independent of a local
choice of frame for the deformed $Q^\vee$: changing such a frame
adds a term $iX$ to $\dot i$, killed by $p$.

Choose a local determinant-one frame of $V=E_0$ normalizing the
nonvanishing section $s_0$ to $(1,0)$, also over the dual numbers.
The induced first-order change of frame on
$\operatorname{End}^0(V)/\mathcal O_T$ is zero. Indeed the
infinitesimal change is the commutator $[X,Y]$ of two trace-zero
two-by-two matrices, and in characteristic two this commutator
is scalar. Equivalently the rank-two quotient representation
of $\operatorname{SL}_2$ is its Frobenius twist. The frame of $P$
is fixed, so the frame of $\mathfrak g$ is stationary to first
order. The chosen determinant-one frame also leaves the local
coordinate of the section $w$ of $LP$ unchanged.

In this frame the two matrices retain their earlier form.
Because $d(h^2)=0$ and the local unit $u$ is parameter-independent,
\[
 \dot a=\begin{pmatrix}0&0\\0&\dot w\end{pmatrix},
 \qquad
 \dot f=\begin{pmatrix}u\dot w&0\\0&0\end{pmatrix}.
\]
Consequently
\[
 p\dot i=aJ_2\dot a^t+f(u^{-1}J_2)\dot f^t
        =\begin{pmatrix}0&\dot w\\\dot w&0\end{pmatrix}.
\]
The expression is intrinsic, so the same computation on the other
nonvanishing chart proves globally that $p\dot i=\dot wJ_Q$.

Finally differentiate a Cech representative of a section of the
cohomology subspace $S_s$. Terms caused by differentiating its
representative inside $H^1(Q^\vee)$ still lie in $S_s$ and
disappear in the quotient. The remaining term is the cohomology
map induced by $p\dot i$. The constant dimensions on the sharp
open justify cohomology base change; (DE-MN2) identifies the
quotient with $H^1(Q)$. This proves (DE-MN3).
In graph coordinates $S_s=\{\alpha v+\beta T_s v\}$, its
derivative is $v\mapsto\beta(dT_s)v$ modulo $S_s$, giving the
last statement with the same basis interchange as in (DE-TZ3).
\end{proof}

\begin{remark}[Precise unproved lower bound]
\label{rem:DE-A3-positive-monad-gate-note}
\status{Proved scope statement; the required variation-image lower bound remains unproved.}
Put
\[
\begin{split}
\Delta_s&=\operatorname{im}\bigl[H^0(E_0)\oplus H^0(E_1)
\xrightarrow{(t_0,t_1)\mapsto s_0\wedge t_1+t_0\wedge s_1}
H^0(LP)\bigr],\\
\Psi_s:\Delta_s&\longrightarrow
\operatorname{Hom}_k\bigl(H^1(Q^\vee),H^1(Q)\bigr),
\qquad \dot w\longmapsto H^1(\dot wJ_Q).
\end{split}
\]
The uniform sharpness problem would follow from
$\dim_k\operatorname{im}\Psi_s\ge2$ at one suitable sharp source.
This is a condition on the span of the entire family of variation
maps, not a rank assertion about one fixed member of that family.
The acyclicity of $\mathfrak g$,
the exactness of (DE-MN1), and the trace-zero theorem do not by
themselves imply that lower bound. No Petri multiplication
identification or generic dimension claim is made in this note.
\end{remark}


% Promoted after complete root and reduction_descent proof/source audits; candidate preserved.
\subsection{Removing a common two-power from the polarization}

\begin{lemma}[The actual maximal two-power root]
\label{lem:DE-A3-mixed-common-power-root}
\status{Proved; complete independent proof and original-source audits passed.}
Let $A$ be an abelian surface of $2$-rank one over an algebraically
closed field of characteristic two, and let $M$ be ample.  Use the
integral polarization indices $a=a(M)$ and $b=b(M)$ of
Lemma~\ref{lem:DE-A3-mixed-polarization-blocks}, and put
\[
             r=\min\{a,b\},\qquad n=2^r,\qquad D=\chi(M).
\]
There is an ample line bundle $H$ on the same surface such that
\[
 M\simeq H^n,\qquad
 a(H)=a-r,\qquad b(H)=b-r,\qquad \chi(H)=D/n^2.
                                                        \tag{DE-CP1}
\]
Moreover $r$ is the largest nonnegative integer for which $M$ has
an actual $2^r$th root in $\Pic(A)$.  The root is not asserted to be
defined over a smaller, nonalgebraically closed ground field.
\end{lemma}
\begin{proof}
We use the actual integral blocks, not merely the rational slope
decomposition.  On each of the ordinary and local-local Dieudonne
summands, the matrix of $\lambda_M$ has both elementary divisors
divisible by $n$.  Equivalently it annihilates the entire finite
group scheme $A[n]$.  The integral splitting and exact Dieudonne
dictionary proving this equivalence were established in
Lemma~\ref{lem:DE-A3-mixed-polarization-blocks}; all its hypotheses
hold here.

The finite faithfully flat quotient $[n]:A\to A$ is the quotient
by $A[n]$.  Consequently there is a homomorphism
$\varphi:A\to A^\vee$ with
\[
                   \lambda_M=\varphi\circ[n]=n\varphi.
\]
Dualizing gives $n\varphi^\vee=n\varphi$.
The group $\Hom(A,A^\vee)$ has no integer torsion: if $n f=0$,
the morphism $f$ factors through the finite group scheme
$A^\vee[n]$, and a morphism from the proper geometrically
connected reduced variety $A$ to a finite affine scheme is constant;
as a group homomorphism it is zero.  Thus $\varphi$ is symmetric.

Every symmetric homomorphism of an abelian variety over an
algebraically closed field is induced by a line bundle, in arbitrary
characteristic, including two.  We use the full proof of
\cite[Proposition~11.2 and Corollary~11.3, pp.~159--160]
{EdixhovenVanDerGeerMoonenAV}.  That proof obtains the line bundle
by finite-flat isotropic descent from a Poincare pullback and does
not divide a theta-group argument by two.
Choose $H_0$ with $\lambda_{H_0}=\varphi$.
Then $P=M\otimes H_0^{-n}$ lies in $\Pic^0(A)$, since its associated
homomorphism is zero.  Multiplication by $n$ on the dual abelian
variety is surjective, hence surjective on algebraically closed
points.  There is therefore $Q\in\Pic^0(A)$ with $Q^n\simeq P$.
Set $H=H_0\otimes Q$.  This gives the actual isomorphism
$H^n\simeq M$.  A line bundle is ample if a positive tensor power
is ample, so $H$ is ample.

The equality of polarizations $\lambda_M=n\lambda_H$ subtracts
$r$ from each integral elementary-divisor valuation, proving both
index formulas.  Riemann--Roch on a surface gives
$\chi(H^n)=n^2\chi(H)$, proving the Euler-characteristic formula.
Conversely, if $M=H'^{\,2^s}$ for a line bundle $H'$, it is ample
and its integral polarization blocks have nonnegative indices.
The equality $\lambda_M=2^s\lambda_{H'}$ then gives $a,b\ge s$.
Thus $s\le r$, proving maximality.
\end{proof}

\begin{theorem}[All degrees when the ordinary index is no larger]
\label{thm:DE-A3-mixed-index-ordered-ranks}
\status{Proved; complete independent proof and original-source audits passed.}
Let $A$ be an abelian surface of $2$-rank one over a perfect field
of characteristic two, and let $M$ be ample.  Compute $a(M),b(M)$
after algebraic closure, and assume
\[
                         0<a(M)\le b(M).
\]
Put $D=\chi(M)$ and $L=M^3$.  For every $d\ge1$, the geometric
generic relative stable rank in the complete system $|L^d|$ is
\[
                             r_d=9Dd^2-2.              \tag{DE-CP2}
\]
The relative coherent dimension is $9Dd^2-1$, and the generic
smooth curve has genus $9Dd^2+1$ and $2$-rank $9Dd^2-1$.
For the embedding by the very ample bundle $L$, the same relative
stable rank holds for ambient degree-$d$ hypersurfaces whenever
restriction to the complete system is surjective; in particular
the earlier bound
$d\ge\max\{1,\operatorname{reg}(\mathcal I_A)-1\}$ suffices.
\end{theorem}
\begin{proof}
Work over an algebraic closure, as permitted by the geometric
generic assertion.  Apply the preceding lemma with $r=a(M)>0$.
It gives an actual ample root $H$ on $A$ with $a(H)=0$ and
$M\simeq H^{2^a}$.  The complete systems are therefore identical:
\[
             |M^{3d}|=|H^{3d'}|,\qquad d'=2^a d.
\]
The integer $d'$ is even.  The full proof of
Theorem~\ref{thm:DE-A3-mixed-nonproduct-ranks} applies to $(A,H)$:
it supplies sharp odd-degree nodal norms from an etale principal
source, obtains every even degree by transverse unions, and uses
the integral Prym obstruction for the matching upper bound.
In particular its even formula gives
\[
 r_d=9\chi(H)(d')^2-2
     =9(D/2^{2a})(2^a d)^2-2
     =9Dd^2-2.
\]
All systems here are the full systems of the same actual line
bundle; no claim of genericity under a restricted pullback map is
needed.  Adjunction, ample vanishing, and subtraction of the ambient
stable rank one give the remaining dimensions, as in that theorem.
The bundle $M^3$ is very ample.  Surjectivity from ambient
hypersurfaces, with geometrically integral fibres on coefficient
spaces, transfers the same generic rank, exactly as in the already
proved regularity argument.
\end{proof}

\begin{corollary}[A primitive form of the remaining mixed problem]
\label{cor:DE-A3-mixed-primitive-rank-gate}
\status{Proved exact reduction; primitive-case nonvanishing remains open.}
For every ample $M$ on a mixed abelian surface, the same-surface
root $H$ of (DE-CP1) identifies the degree-$d$ rank problem for
$M^3$ with the degree-$2^r d$ problem for $H^3$.
Thus, together with the proved ordinary-index-zero theorem and
(DE-CP2), the unresolved nonproduct positive-index case reduces to
ample bundles satisfying
\[
                         a(H)>0,\qquad b(H)=0,
                                                        \tag{DE-CP3}
\]
in all positive degrees.  This reduction does not prove the
restricted cup-product radical is zero, and does not assert that
an arbitrary mixed surface has an index-zero ample class.
\end{corollary}
\begin{proof}
The actual line-bundle equality identifies the systems and their
families, hence their geometric generic Frobenius ranks.  If
$a(M)\le b(M)$, its root has ordinary index zero and the relevant
degree has the parity already used above.  If $a(M)>b(M)$, take
$r=b(M)$.  Its root has the indices in (DE-CP3), and the exponent
changes by the displayed factor.  The root stays on the original
surface; only its line bundle and the exponent change.
\end{proof}


% Promoted exact reduction after complete root and reduction_descent audits; candidate preserved.
\subsection{The precise restricted cup-product pairing}

\begin{lemma}[The actual determinant directions and their dual pairing]
\label{lem:DE-A3-positive-restricted-Petri-note}
\status{Proved equivalent reduction; complete independent audits passed.}
Keep a sharp source in the preceding monad construction, put
$\mathcal L=\det Q=LP$, and write $U=H^1(T,Q^\vee)$.
The determinant-direction map has the exact sequence
\[
0\longrightarrow H^0(KP)\longrightarrow
H^0(E_0)\oplus H^0(E_1)\xrightarrow{\dot w}
H^0(\mathcal L)\longrightarrow H^1(KP)\longrightarrow0.
\tag{DE-PT1}
\]
In particular, its image $\Delta_s$ satisfies
\[
\operatorname{codim}\Delta_s=h^1(KP)=h^0(KP).
\tag{DE-PT2}
\]
The symmetric bilinear cup-wedge map
\[
\mu:\operatorname{Sym}^2 U\longrightarrow
H^2(T,\mathcal L^{-1}),\qquad
x\,y\longmapsto \wedge_*(x\mathbin{\smile}y)
\tag{DE-PT3}
\]
is well defined. Choose the fixed invariant volume form to identify
$H^2(T,\mathcal L^{-1})=H^0(T,\mathcal L)^*$ and let
$\mu_{\Delta}$ be its composite with restriction to $\Delta_s^*$.
Then $\mu_{\Delta}$ is the linear dual of the family
\[
\Psi_s:\Delta_s\longrightarrow
\operatorname{Hom}_k(U,H^1(Q)),\qquad
v\longmapsto H^1(vJ_Q),
\]
after Serre duality identifies $H^1(Q)$ with $U^*$.
Consequently
\[
\dim\operatorname{im}\Psi_s=\dim\operatorname{im}\mu_{\Delta}.
\tag{DE-PT4}
\]
No smoothness or reducedness of the determinant curve is required.
\end{lemma}
\begin{proof}
Twist the defining sequence for $K$ by $P$ and interchange its two
summands using $E_1=E_0P$ and the fixed $P^2$ trivialization.
Its map to $LP$ becomes
$(t_0,t_1)\mapsto s_0\wedge t_1+t_0\wedge s_1$.
The absence of a common zero of the two vector-bundle sections makes
this sheaf map surjective. Since $H^j(E_i)=0$ for $j>0$,
its cohomology gives (DE-PT1). Both the middle section space and
$H^0(LP)$ have dimension $b$, proving (DE-PT2).

Interchanging $x,y$ in degree-one cup product contributes its graded
sign, and interchanging the two factors of the alternating bundle
pairing contributes the same sign. Their product is positive.
Thus the bilinear map in (DE-PT3) is symmetric, including in
characteristic two, and factors through the stated symmetric-power
quotient. No identification of symmetric and divided powers is used.

For $v\in\Delta_s$ and $x,y\in U$, functoriality of cup product and
Serre duality gives
\[
\bigl\langle H^1(vJ_Q)x,y\bigr\rangle
=\operatorname{Tr}_T\bigl(v\,\wedge_*(x\mathbin{\smile}y)\bigr).
\tag{DE-PT5}
\]
This can also be checked in a frame of $Q$, where $J_Q$ is the
inverse alternating matrix and contraction with $y$ is the wedge
pairing on $Q^\vee$. Equation (DE-PT5) is precisely the duality
assertion. The image of $\Psi_s$ consists of symmetric bilinear
forms, which is the space $(\operatorname{Sym}^2U)^*$; taking the
dual map therefore proves (DE-PT4). All operations are on vector
bundles and their cohomology on the smooth surface $T$.
\end{proof}

\begin{corollary}[An exact common-radical test]
\label{cor:DE-A3-positive-Petri-radical-note}
\status{Proved equivalent reduction; no radical vanishing asserted.}
At the fixed sharp source, a trace-direction variation with nonzero
determinant exists if and only if
\[
\{x\in U:\mu_{\Delta}(x\,y)=0\text{ for every }y\in U\}=0.
\tag{DE-PT6}
\]
Equivalently, no nonzero $x\in U$ satisfies
\[
\mu(xU)\subset\Delta_s^\perp
=\operatorname{im}\bigl(H^1(KP)^*\longrightarrow
H^0(LP)^*\bigr).
\tag{DE-PT7}
\]
\end{corollary}
\begin{proof}
By (DE-PT5) the displayed radical is exactly the common kernel of
all the maps $\Psi_s(v)$. An invertible member forces that common
kernel to be zero. Conversely, the proved trace-zero theorem says
that, in the compatible graph bases, a linear space consisting only
of singular variations is zero or a single nilpotent matrix line.
Such a space has a nonzero common kernel. Changing to the
Serre-duality bases preserves the common-kernel condition, so a zero
radical forces an invertible member. The final equality follows
by dualizing (DE-PT1).
\end{proof}

\begin{remark}[What is still missing]
\label{rem:DE-A3-positive-Petri-radical-scope-note}
\status{Proved scope statement; the geometric nonvanishing remains open.}
Neither (DE-PT2) nor (DE-PT3) bounds the restricted cup-product
image from below. In particular the determinant directions cannot
be replaced by all of $H^0(LP)$: their actual annihilator is the
dual of $H^1(KP)$, with the twist by $P$ essential.
The curve pencil trick, when its additional hypotheses hold, controls
a different multiplication map on the determinant curve; no
identification with (DE-PT3) is assumed. The remaining uniform
problem is to prove (DE-PT6), or a sufficient two-dimensional image
bound in (DE-PT4), at one suitable sharp source on each surface.
\end{remark}


% Promoted after complete root and reduction/descent audits; candidate preserved.
% Separate geometric input; no restricted-Petri nonvanishing is claimed.
\subsection{Smooth determinant curves in a sufficiently positive range}

\begin{lemma}[First jets of the two finite pushforwards]
\label{lem:DE-A3-positive-pushforward-jets}
\status{Proved; independent proof audits passed.}
Keep the mixed-surface, $\alpha_2$-torsor and line-bundle notation
of Lemma~\ref{lem:DE-A3-positive-two-bundles}. Assume that each of
$J$ and $JQ$ is the third power of a very ample line bundle on $C$.
Then $E_0=\nu_*J$ and $E_1=\nu_*(JQ)$ separate first jets at every
point of $T$:
\[
 H^0(T,E_i)\longrightarrow
 E_i\otimes_{\mathcal O_T}\mathcal O_T/\mathfrak m_t^2
 \quad\hbox{is surjective for every }t\in T(k).
                                                        \tag{DE-SJ1}
\]
In particular the hypothesis holds if $J=J_0^{3d}$, with $J_0$
ample and $d\ge3$. In the original polarization notation this
applies whenever $M$ itself descends through $\rho$, with descended
bundle $J_0$, and the section bundle is $M^{3d}$ with $d\ge3$.
\end{lemma}
\begin{proof}
The finite flat isogeny $\nu:C\to T$ is a torsor under the smooth
surface's nonzero invariant $\alpha_2$ direction. At its unique
geometric point over $t$, suitable completed regular parameters
identify the map with
\[
                  k[[x,z^2]]\longrightarrow k[[x,z]].
\]
This follows directly from the local torsor algebra
$\mathcal O_T[z]/(z^2-h)$: smoothness of $C$ forces $dh\ne0$
at the point, so $h$ is a regular parameter on $T$.
The inverse image of the first infinitesimal neighbourhood of $t$
therefore has ideal
\[
                   (x^2,xz^2,z^4)\subset k[[x,z]].
\]
Its quotient has basis $1,z,z^2,z^3,x,xz$, and this ideal contains
$(x,z)^4$. Thus it suffices that $J$ and $JQ$ separate third jets
at the corresponding point of $C$.

If $V$ is very ample, choose a section $v_0$ nonzero at the point
and sections whose ratios with $v_0$ give a basis of
$\mathfrak m/\mathfrak m^2$. Products of three sections of $V$,
using $v_0$ for the remaining factors, give every monomial of total
degree at most three modulo $\mathfrak m^4$. These monomials span
that local Artin algebra, independently of characteristic.
Consequently $V^3$ separates third jets. The completed calculation
computes the same finite local quotient before completion, proving
(DE-SJ1) by finite pushforward.

For the stated sufficient condition, $J_0^d$ is very ample for
$d\ge3$. Choose a cube root of $Q$ in $\Pic^0(C)(k)$; every
$\Pic^0$ twist of $J_0^d$ is a translate of it, since its
polarization isogeny is surjective on algebraically closed points.
Thus $JQ$ is also a cube of a very ample line bundle. Descent
preserves ampleness under the finite surjection $\rho$ and gives
$J=J_0^{3d}$ in the final situation.
\end{proof}

\begin{proposition}[A smooth determinant curve and a separable pencil]
\label{prop:DE-A3-positive-smooth-separable-pencil}
\status{Proved; independent proof audits passed.}
Under the preceding hypotheses, a general pair $(s_0,s_1)$ has
reduced disjoint zero schemes and a smooth integral determinant
curve $D=(s_0\wedge s_1=0)$. The logarithmic differential
$\eta_P|_D$ is nonzero. Hence the base-point-free pencil
$\langle a_0,a_1\rangle$ on $A=R_0^2$ is separable.

If a sharp-source open is nonempty for these bundles, the general
pair may also be chosen in that open. At such a pair the ordinary
curve multiplication map
\[
 H^0(D,A)\otimes H^0(D,\omega_D A^{-1})
       \longrightarrow H^0(D,\omega_D)
                                                        \tag{DE-SJ2}
\]
has rank at least three, and has rank four unless
$\omega_D A^{-2}$ is trivial.
\end{proposition}
\begin{proof}
First-jet separation implies that the value and its first
derivatives of each vector-bundle section can be chosen
independently at every point. The incidence where a rank-two
section vanishes and its derivative has rank less than two has
codimension at least three in the section space, over a surface
of dimension two. A general zero scheme is therefore finite and
reduced. Simultaneous vanishing of both sections imposes four
independent value conditions, so the two zero schemes are disjoint.

On this disjoint-zero open, a point of the determinant divisor has
matrix rank exactly one. Normalize one nonzero section value to
the first coordinate vector. The determinant and its two first
derivatives are then three independent conditions on the first
jet of the other section. Their incidence has codimension three
over the two-dimensional base. Thus a general determinant divisor
is smooth. Its line bundle $LP$ is ample, so it is nonempty and
connected: the restriction sequence and ample vanishing give
$H^0(D,\mathcal O_D)=k$. A smooth connected curve over the
algebraically closed field is integral.

Suppose $\eta_P$ restricted to $\Omega_D^1$ were zero. The conormal
sequence would lift its restriction as a section of
$\mathcal O_D(-D)$. This line bundle has negative degree, so has
no nonzero section. But the nonzero invariant differential
$\eta_P$ is nowhere zero as a section of $\Omega_T^1$, and hence
its restriction to $\Omega_T^1|_D$ is nonzero, a contradiction.
The already proved identity
$d\log(a_1/a_0)=\eta_P|_D$ now proves that the pencil is separable.

All conditions just established are nonempty opens in the
irreducible space of section pairs; they meet any nonempty
sharp-source open. There $h^0(A)=h^0(\omega_D A^{-1})=2$ by the
proved contraction theorem, and the two pencil sections form a
basis of $H^0(A)$. The base-point-free pencil sequence, tensored
by $\omega_D A^{-1}$, identifies the kernel of (DE-SJ2) with
\[
                          H^0(D,\omega_D A^{-2}).
\]
This line bundle has degree zero on the smooth integral curve.
Its section space has dimension zero unless it is trivial, in
which case the dimension is one. The domain of (DE-SJ2) has
dimension four, proving both assertions.
\end{proof}

\begin{remark}[The distinction from the restricted variation]
\label{rem:DE-A3-positive-smooth-pencil-scope}
\status{Proved scope; no restricted-Petri conclusion.}
The preceding multiplication map is a curve-pencil map. No
identification of it with the actual restricted surface cup-wedge
map in (DE-PT3)--(DE-PT7) has been proved. Thus its rank bound does
not establish the required common-radical vanishing. The power
hypothesis is explicit: positivity of the ordinary index of
$M^d$ alone does not imply that $M$ itself descends along $\rho$.
In particular no unproved low-degree jet separation or determinant
curve smoothness is used in earlier arguments.
\end{remark}


% /root/monodromy, 2026-11-19. Root and RD complete proof audits passed; no uniform image claim.
\subsection{A local self-tacnode criterion for a nonordinary surface}
\label{subsec:ss-self-tacnode}
\status{Proved under the stated hypotheses. The parent and reduction/descent agent independently audited both complete local proofs. This subsection isolates
one local geometric step. It does not assert irreducibility of a Hasse divisor
or determine any new full monodromy image.}

\begin{lemma}[Two-point Cartier obstruction]
\label{lem:ss-two-point-Cartier}
\status{Proved under stated hypotheses}
Let $C$ be a smooth projective curve over an algebraically closed field
of characteristic two. Let $P\ne Q$ be points, with chosen uniformizers
$x_P,x_Q$, and put $A_C=\ker(C:H^0(\omega_C)\to H^0(\omega_C))$.
Let $U\subset H^0(\omega_C(2P+2Q))$ be the subspace of forms with
zero residues at both points. Taking the double-pole coefficients gives
\[
 0\longrightarrow H^0(\omega_C)\longrightarrow U
   \xrightarrow{\mathrm{pp}}k^2\longrightarrow0.
\]
Cartier maps $U$ to $H^0(\omega_C)$ and induces a semilinear map
\[
 \delta:k^2\longrightarrow\operatorname{coker}
       (C:H^0(\omega_C)\to H^0(\omega_C)).
\]
Under Serre and finite-Frobenius duality its dual, with the appropriate
scalar Frobenius twists, is the evaluation map
\[
          A_C\longrightarrow k^2,
          \qquad\eta\longmapsto(\eta(P),\eta(Q)).
\]
In particular a Cartier-killed form with equal, nonzero double-pole
coefficients exists if and only if
$\eta(P)+\eta(Q)=0$ for every $\eta\in A_C$.
\end{lemma}
\begin{proof}
Riemann--Roch gives $\dim H^0(\omega_C(2P+2Q))=g+3$.
The residue theorem imposes the one relation between the two residues,
and the residue map has rank one. Thus $\dim U=g+2$.
The only forms in $U$ with both double-pole coefficients zero are
holomorphic, proving the exact sequence. The local rule for Cartier
kills the double-pole term and takes the square root of the residue;
hence it has the asserted target.

For clarity we identify the dual boundary, rather than inferring it
from a dimension count. Finite Frobenius duality identifies the full
meromorphic Cartier map with
\[
 H^1(\mathcal O_C(-P-Q))\xrightarrow{F}
 H^1(\mathcal O_C(-2P-2Q)).
\]
The respective kernels of the maps to $H^1(\mathcal O_C)$ are the
constant-difference line $k^2/k$ and
$H^0(\mathcal O_{2P+2Q})/k$. Frobenius is bijective on the
constant-difference line, so it can be eliminated before taking the
remaining two first-jet coordinates.

Explicitly, represent a Frobenius-kernel class on $H^1(\mathcal O_C)$
by a Cech cocycle $c_{ij}$, with $c_{ij}^2=b_i-b_j$.
Then $db_i$ glue to a global Cartier-killed form $\eta$.
The sequence
$0\to\mathcal O_C\xrightarrow F F_*\mathcal O_C
 \xrightarrow d B_C^1\to0$
and bijectivity of Frobenius on constants identify all such classes
with $A_C$, with the scalar twist just indicated. Choose disjoint
small affine neighbourhoods of $P,Q$ in the cover. Lifting the class
to $H^1(\mathcal O_C(-P-Q))$ is ambiguous by the
constant-difference line. There is a unique adjustment of that line
which kills the constant difference of its Frobenius image.
Subtracting one common constant then makes both zeroth-order terms
of $b$ zero. The remaining boundary has first-jet coordinates
\[
                   db(P)=\eta(P),\qquad db(Q)=\eta(Q).
\]
Serre duality pairs the coefficients of $x^{-2}dx$ with these
first-jet coordinates, by the residue of their product. This gives
the stated evaluation map. This is the two-point version of the
fully written one-point calculation in Lemma~\ref{lem:pointed-Cartier};
the finite-duality identification is
\cite[Lemma~1.3.6 and Proposition~1.3.7, pp.~26--27]{BrionKumar2005}.
Pairing the vector $(1,1)$ with that evaluation map proves the last
assertion. All relevant assertions are vanishing assertions and are
unchanged by the scalar square roots in Cartier.
\end{proof}

\begin{lemma}[Smoothness at a mixed self-tacnode]
\label{lem:ss-mixed-self-tacnode}
\status{Proved under stated hypotheses}
Let $A',A$ be superspecial abelian surfaces in characteristic two,
and let $\phi:A'\to A$ be an odd-degree isogeny. Fix invariant volume
forms and identify the invariant one-forms using $\phi^*$.
Let $C'\subset A'$ be a smooth ample curve and suppose its image
$Z=\phi(C')$ is birationally normalized by $C'$. Assume that $Z$ has
ordinary nodes and exactly one simple tacnode, the latter obtained
by identifying $P,Q=P+g$ for a nonzero $g\in\ker\phi$.
Their tangent lines agree after translation. Choose the same tangent
coordinate on these two branches.

Assume that
\[
 \ker C_{C'}=
       H^0(A',\Omega^1_{A'})|_{C'}\oplus k\eta,
       \qquad \eta(P)\ne\eta(Q),
\]
and that $\mathcal N=\mathcal O_A(Z)$ has a separable polarization.
Suppose also that $\mathcal N$ is very ample. Then the relative
coherent Frobenius on $|\mathcal N|$, on the quotient by ambient
$H^1(A,\mathcal O_A)$, has corank one at $Z$, and its determinant
has a nonzero first derivative there. Thus the relative Hasse divisor
is smooth at $Z$.
\end{lemma}
\begin{proof}
Invariant one-forms take equal values on the two translated tangent
vectors. The hypothesis on $\eta$ and
Lemma~\ref{lem:ss-two-point-Cartier} show that no
Cartier-killed meromorphic form on $C'$ has equal nonzero
principal coefficients at $P,Q$.

Here are the dualizing conditions at the singularities. At the
tacnode, choose completed coordinates in which the two branches
are $y=0$ and $y=\varphi(x)$, where
$\varphi(x)=c x^2+O(x^3)$ and $c\ne0$.
The completed local ring is the subring of
$k[[x]]\oplus k[[x]]$ consisting of pairs congruent modulo $x^2$.
The adjunction generator restricts on both branches to
$dx/\varphi(x)$. Thus a dualizing form has at most double poles,
with equal double-pole coefficients and equal residues.
At a node the two poles are simple, with opposite residues.
Cartier acts on the normalization componentwise by the local trace
rule and takes square roots of residues. Consequently a
Cartier-killed dualizing form has zero residues at every branch.
The two-point lemma now excludes a nonzero double-pole coefficient
at the tacnode. It follows that
\[
             \ker C_Z=\ker C_{C'},\qquad \dim\ker C_Z=3.
\]

We identify the ambient boundary precisely. For every effective
Cartier divisor $Z$ on an abelian surface the adjunction sequence is
\[
 0\longrightarrow\mathcal O_A\xrightarrow{s}\mathcal N
 \longrightarrow\omega_Z\longrightarrow0,
\]
using the chosen invariant volume form. Its boundary on the
restriction of invariant one-forms is, after contraction with that
volume form, the differential of the polarization
$\lambda_{\mathcal N}:A\to A^\vee$.
Indeed, on trivializing opens write $s_i=g_{ij}s_j$. For an invariant
vector field $v$, the local lift of its contracted one-form is
$v(s_i)$. The boundary of these lifts is
$v(g_{ij})/g_{ij}$, which is exactly the differential of
$a\mapsto t_a^*\mathcal N\otimes\mathcal N^{-1}$.
Separability of the polarization makes this two-dimensional map
an isomorphism. In particular the restrictions of the two invariant
forms are independent, as required in the displayed direct sum.

Ample vanishing gives $H^1(A,\mathcal N)=0$. Thus the image of
$H^0(\mathcal N)/ks$ in $H^0(\omega_Z)$ is the kernel of that
boundary. It is Cartier-stable, and its Cartier kernel is
one-dimensional: subtract a unique invariant form from $\eta$
to put it in the boundary kernel. Denote the resulting form by
$\eta_0$. Its two values still differ, since invariant forms have
equal values. Lift $\eta_0$ to $u\in H^0(\mathcal N)$, uniquely
modulo $ks$.

We finish with the first derivative, keeping track of the quotient.
In a local frame of $\mathcal N$, holomorphicity on both normalized
branches implies
\[
        u(P)=0,\qquad
        u=a y+b x^2+\text{terms irrelevant to the branch values}.
\]
The two branch values of $\eta_0$ are $b/c$ and $a+b/c$.
Thus $a\ne0$. Since $s=y(y-\varphi(x))$ has quadratic term $y^2$
and no $xy$ term, the Cartier trace of $sv$ at the singular point
vanishes for every regular $v$.

In root-removed coordinates, Cartier is the bilinear map
$M(s):W\to W^{(2)}$, with $W=H^0(\mathcal N)$, and the relative
map is
\[
                   W/ks\longrightarrow W^{(2)}/ks^{(2)}.
\]
The Frobenius-twisted evaluation at the singular point annihilates
its image and is nonzero on its target quotient; it spans the
one-dimensional left kernel. Very ampleness supplies a global
$\dot s$ with zero value and first jet $x$ at that point. The
$xy$ coefficient of $\dot s\,u$ is $a\ne0$, and hence
\[
             \operatorname{ev}^{(2)}_P(M(\dot s)u)\ne0.
\]
This is the intrinsic derivative of the quotient map: a derivative
of the chosen lift $u$ contributes $M(s)\dot u$, annihilated by
evaluation; the target kernel line $ks^{(2)}$ has zero first
variation in characteristic two, and its value at the singular
point is zero. The rank-one adjugate formula therefore gives a
nonzero derivative of the relative determinant. The calculation
is on the affine coefficient cone; scaling contributes only the
radial derivative, which vanishes at a determinant-zero point, so
the conclusion descends to the projective system.
\end{proof}

\begin{remark}[The remaining global gate]
\label{rem:ss-self-tacnode-global-gate}
\status{Proved under stated hypotheses}
The lemma requires a specified self-tacnode with
$\eta(P)\ne\eta(Q)$. It does not prove that such points occur on
every translation-collision divisor of a norm family, and it does
not address its triple-identification divisors. Those incidence
statements remain unproved. In particular this local criterion
alone does not imply irreducibility of any higher-degree Hasse
divisor or a uniform supersingular full-group theorem.
\end{remark}


% /root/monodromy, 2026-11-19. Both local proofs independently audited by root and RD; no global extension asserted.
\subsection{The corresponding local criterion at an ordinary triple point}
\label{subsec:ss-triple-point}
\status{Proved under the stated hypotheses. The parent and reduction/descent agent independently audited both complete local proofs. No global incidence or full-monodromy assertion is made.}

\begin{lemma}[Smoothness at a mixed ordinary triple point]
\label{lem:ss-mixed-triple-point}
\status{Proved under stated hypotheses}
Let $A',A$ be superspecial abelian surfaces in characteristic two,
and let $\phi:A'\to A$ be an odd-degree isogeny. Let $C'\subset A'$
be smooth and suppose $\phi|_{C'}$ is birational onto its image $Z$.
Assume $Z$ has ordinary nodes and exactly one ordinary triple point,
with preimages $P_1,P_2,P_3$. Choose tangent parameters on these
three branches such that their tangent vectors, identified by
$\phi$, satisfy $v_3=v_1+v_2$.
Suppose
\[
 \ker C_{C'}=H^0(A',\Omega^1_{A'})|_{C'}\oplus k\eta,
 \qquad \eta(P_1)+\eta(P_2)+\eta(P_3)\ne0,
\]
where values use the chosen tangent vectors. Suppose
$\mathcal N=\mathcal O_A(Z)$ is globally generated and its
polarization is separable. Then the relative coherent Cartier map
has corank one at $Z$, and its determinant has a nonzero first
derivative there. The relative Hasse divisor is therefore smooth
at $Z$.
\end{lemma}
\begin{proof}
The three-point version of Lemma~\ref{lem:ss-two-point-Cartier}
is obtained by replacing $P+Q$ with $P_1+P_2+P_3$ in its proof.
For completeness, the zero-residue subspace of
$H^0(\omega_{C'}(2P_1+2P_2+2P_3))$ has dimension $g(C')+3$,
and its quotient by holomorphic forms is the space $k^3$ of
double-pole coefficients. Its Cartier map to the holomorphic
Cartier cokernel is dual, up to the scalar Frobenius twists, to
\[
     \ker C_{C'}\longrightarrow k^3,
     \quad\xi\longmapsto(\xi(P_1),\xi(P_2),\xi(P_3)).
\]
Indeed the constant-difference kernel in the twisted Frobenius
calculation is now $k^3/k$. Frobenius is bijective on it; after
its unique adjustment, subtracting a common constant kills all
three zeroth-order terms. The three first-jet terms of the Cech
primitives are $\xi(P_i)$, and the local residue pairing gives
the displayed map. This repeats the complete two-point calculation
without any change to its duality hypotheses.

We spell out the ordinary-triple-point matching. Choose completed
coordinates so that the first two branches are $y=0,x=0$, and
the third is $y=\varphi(x)$, where
$\varphi(x)=x+O(x^2)$. A local equation is
\[
                         f=xy(y-\varphi(x)).
\]
The three adjunction generators have leading terms
$dx/x^2$, $dy/y^2$, and $dx/x^2$ respectively.
Thus all three double-pole coefficients of a dualizing form are
equal. Their residues sum to zero. These are also sufficient:
multiplication of the adjunction generator by the constant and
the two linear terms of a numerator produces the full
three-dimensional space of allowed polar parts. Equivalently,
the conductor is the square of the maximal ideal on each
normalized branch; its polar conditions are equality of the
three leading coefficients and the single residue relation.

Cartier-killed forms have zero residues. The only possible polar
vector at the triple point is therefore a multiple of $(1,1,1)$.
Invariant ambient forms have value sum zero since $v_1+v_2+v_3=0$.
The hypothesis on $\eta$ and the three-point calculation exclude
a nonzero such polar vector. Ordinary nodes likewise contribute
no pole to a Cartier-killed form. Consequently
\[
               \ker C_Z=\ker C_{C'},\qquad\dim\ker C_Z=3.
\]
The adjunction boundary on the two invariant forms is the
differential of $\lambda_{\mathcal N}$, after contraction with
an invariant volume form. Its explicit Cech calculation is in
Lemma~\ref{lem:ss-mixed-self-tacnode}. Separability makes it an
isomorphism. Thus a unique ambient correction of $\eta$ lies in
the relative Cartier kernel; denote it by $\eta_0$. Its three-value
sum is still nonzero, and that relative kernel has dimension one.

Lift $\eta_0$ to $u\in W=H^0(A,\mathcal N)$ modulo $ks$, where
$s$ defines $Z$. Holomorphicity on all branches forces the
constant and both linear terms of $u$ to vanish. Write its
quadratic part as
\[
                          u=a x^2+bxy+c y^2.
\]
After normalizing the invariant volume form at the singular point,
the branch values are $a,c,a+b+c$. Their sum is $b$, hence $b\ne0$.
The equation $s$ has order three. Therefore the Cartier trace of
$sv$ has zero value at the singular point for every regular $v$,
since that value extracts the $xy$ coefficient.

In root-removed coordinates the relative map is
\[
                  M(s):W/ks\longrightarrow W^{(2)}/ks^{(2)}.
\]
Frobenius-twisted evaluation at the singular point spans its
one-dimensional left kernel. Global generation supplies a
section $\dot s$ whose value there is one. The $xy$ coefficient
of $\dot s\,u$ is $b\ne0$, so
\[
                \operatorname{ev}^{(2)}(M(\dot s)u)\ne0.
\]
As in the self-tacnode lemma, derivatives of the input lift are
killed by the left kernel, and the target line $ks^{(2)}$ has no
first variation in characteristic two. The adjugate derivative
formula proves that the determinant has a nonzero derivative.
All quantities are intrinsic on the projective coefficient space:
the radial derivative vanishes at the determinant-zero point.
\end{proof}

\begin{remark}[A remaining global, rather than local, question]\label{auto:A_monodromy:9292c61fd3ca}
\status{Proved under stated hypotheses}

The hypothesis on the three values has not been proved on every
triple-identification divisor in a norm family. In particular,
triples lying on a single elliptic fibre can have a different
incidence dimension from three general points. This lemma does
not discard those components and makes no assertion that the
needed value sum is generically nonzero on them.
\end{remark}

\begin{lemma}[A tangent triple joining two collision types]
\label{lem:ss-tangent-triple-bridge}
\status{Proved under stated hypotheses}
In the same superspecial-isogeny setting, suppose the normalization
$C'$ is sharp in the sense that
\[
             \ker C_{C'}=H^0(A',\Omega^1_{A'})|_{C'}.
\]
Suppose its image $Z$ has ordinary nodes and one further singular
point with three smooth branches: the first two have simple contact,
and the third is transverse to their common tangent. Assume
$\mathcal N=\mathcal O_A(Z)$ is very ample and its polarization is
separable. Then the relative Hasse divisor is smooth at $Z$.
\end{lemma}
\begin{proof}
Let $P,Q,R$ be the preimages of the exceptional singularity, with
$P,Q$ the tangent pair. Choose completed coordinates so that the
branches are $y=0$, $y=\varphi(x)$, and $x=0$, with
$\varphi(x)=c x^2+O(x^3)$ and $c\ne0$. A local equation is
\[
                         f=xy(y-\varphi(x)).
\]
The adjunction generators have pole orders three on the first
two branches and order two on the third. A Cartier-killed
meromorphic differential has no term $x^{-3}dx$, because Cartier
maps it to a nonzero multiple of $x^{-2}dx$; it also has zero
residue. Thus the only possible nonzero polar terms are double
poles at $P,Q,R$.

The dualizing conditions for those zero-residue terms reduce to
\[
                     (a_P,a_Q,a_R)=(a,a,0).
\]
One can read this directly from adjunction: multiplication of
$dx/(x\varphi(x))$ by a numerator with leading term $x$ gives
matching double-pole coefficients at $P,Q$ and none on the third
branch. A quadratic correction to that numerator removes the
matching residues. Conversely multiplication of a dualizing form
by either local coordinate must have residues summing to zero.
The coordinate $x$ gives $a_P+a_Q=0$, and the coordinate $y$
gives $a_R=0$. These are all conditions on zero-residue double
terms: their one-dimensional solution space is realized by the
numerator just described. Holomorphic terms are unrestricted.

By the three-point Cartier boundary calculation, a killed form
with polar vector $(1,1,0)$ exists precisely when
$\xi(P)+\xi(Q)=0$ for every $\xi\in\ker C_{C'}$.
This holds because the kernel consists of invariant ambient
forms and the two tangents agree under translation. The killed
meromorphic forms with zero polar vector are exactly the two
ambient holomorphic forms. Hence
\[
                        \dim\ker C_Z=3.
\]
The separable-polarization boundary is an isomorphism on those
two ambient forms. After their unique subtraction the relative
Cartier kernel is therefore one-dimensional, and its generator
still has nonzero double-pole coefficients on the tangent pair.

Let $u\in H^0(\mathcal N)$ be a numerator lifting that generator.
The preceding local matching shows that $u$ has nonzero linear
$x$ term and zero constant and linear $y$ terms. The defining
section $s$ has order three, so Frobenius-twisted evaluation at
the singular point annihilates the entire Cartier image and
spans its relative left kernel. Very ampleness supplies a global
$\dot s$ with zero value and first jet $y$. Then the $xy$
coefficient of $\dot s\,u$ is nonzero. The same quotient
adjugate derivative calculation as above proves that the relative
Hasse determinant has nonzero derivative at $Z$.
\end{proof}

\begin{remark}[What the tangent-triple bridge would supply]\label{auto:A_monodromy:2b2a821f869e}
\status{Proved under stated hypotheses}

A norm-family parameter with this singularity belongs both to a
pair-tangency locus and to a triple-identification locus. If it is
a smooth point of the relevant target Hasse divisor, any target
Hasse components containing the images of these loci coincide.
The lemma proves that local smoothness under its stated hypotheses.
It does not prove the existence or dominance of such parameters,
or the irreducibility of the source collision loci. Those remain
separate global incidence questions.
\end{remark}


% /root/monodromy, 2026-11-19. Root discovery/certificate; independent complete audit and rerun passed.
\subsection{An exact global contact in the degree-three norm family}
\label{subsec:ss-global-contact-witness}
\status{Proved and checked. The parent constructed and checked the exact
certificate. The monodromy agent independently audited its complete code,
recomputed the local data by truncated elliptic-curve addition, and reran
all four source-chart and all 144 collision checks, preserving the author's
outputs. This is one geometric witness, not a uniform full-group theorem.}

\begin{lemma}[A mixed contact with no other nonnodal collisions]
\label{lem:ss-global-contact03}
\status{Proved under stated hypotheses}
Let $E:y^2+y=x^3$, $X=E^2$, and
$L=\mathcal O_E(3O)\boxtimes\mathcal O_E(3O)$.
Over $\mathbf F_{64}=\mathbf F_2[z]/(z^6+z+1)$ encode
$\sum_{i=0}^5a_i z^i$ by the integer $\sum_{i=0}^5a_i2^i$.
Set $\omega=58$, so $\omega^2+\omega+1=0$, and let
\[
 i(x,y)=(x+1,y+x+\omega),\qquad b=1+i,\qquad
 \phi=\begin{pmatrix}1&b\\-b^\dagger&1\end{pmatrix}:X\to X.
\]
Then $\phi^\dagger\phi=3$ and $\ker\phi=\{(-bT,T):T\in E[3]\}$.
In the tensor basis $(1,x,y)\otimes(1,x,y)$, take
\[
 s=(38,14,10,47,47,43,30,4,8).
\]
The curve $C=(s=0)$ is geometrically smooth. Its norm
$Z=\phi(C)\in|L^3|$ has exactly one simple self-tacnode and otherwise
only ordinary nodes. The two preimages of the tacnode are
\[
 P=((17,52),(8,53)),\qquad Q=((18,35),(19,10)),
 \qquad Q-P=G=((58,59),(0,0)).
\]
The source Cartier kernel is the two-dimensional ambient kernel plus
one additional form whose values at these two translated tangent
vectors differ. Consequently the degree-three relative Hasse divisor
is smooth at $Z$.
\end{lemma}
\begin{proof}
The quaternion identity $i^\dagger=-i$, $i^2=-1$ gives
$b^\dagger b=2$, and hence the stated isogeny identities and kernel.
The exact norm line bundle is $L^3$, as proved in
Theorem~\ref{thm:ss-uniform-GL8-block}.
Write $S=(s_{ij})$ and $J=J_3\otimes J_3$, with $J_3$ reversing the
three basis vectors. The coefficient certificate gives
\[
 \det S=43\ne0,\qquad s^{[2]t}Js=0,\qquad\rank M(s)=8,
\]
and the literal identity
\[
 M(s)u=s^{[2]},\qquad
             u=(3,5,56,20,4,38,38,12,0).
\]
Here $M(s)$ is the root-removed Cartier multiplication matrix,
whose one-factor matrix is
\[
 M_E(a)=\begin{pmatrix}a_1&a_0&0\\a_2&0&a_0\\0&a_2&a_1\end{pmatrix}.
\]
Thus the induced Cartier map $W/ks\to W^{(2)}/ks^{(2)}$ has
one-dimensional kernel, represented by $u$.
Indeed $\ker M(s)=ks$ and $s^{[2]}$ lies in its image; quotienting
by that image line increases the kernel dimension by one before
quotienting the source by $ks$.

On the smooth source curve, the two invariant ambient one-forms
are Cartier-killed. Their adjunction boundary is the differential
of $\lambda_L=3\lambda_\Theta$, and is an isomorphism.
Therefore its full Cartier kernel has dimension three, and the
additional kernel form is represented by
$\operatorname{Res}(u\,dx_1\wedge dx_2/s)$, up to ambient correction.
Direct evaluation gives
\[
 s(P)=s(Q)=0,\qquad ds(P)=ds(Q)=(34,53),\qquad
                         u(P)+u(Q)=60\ne0.
\]
The derivatives use the invariant derivations
$\delta x=1$, $\delta y=x^2$ on each factor. The common tangent
vector $(53,34)$ makes the displayed difference precisely the
value difference of the additional differential; ambient corrections
have equal values on translated tangent vectors.
The independent truncated-series calculation, using the rational
elliptic addition law, finds an actual source-curve branch through
$P$ to order two and obtains
\[
                 s(\gamma(t))=O(t^3),\qquad
                 s(\gamma(t)+G)=55t^2+O(t^3).
\]
The nonzero coefficient proves simple contact, not higher tangency.

We describe the global exclusion, including its projective coverage.
The two charts $Z=1$ and $Y=1$ cover the cubic
$Y^2Z+YZ^2=X^3$. Both have equation $y^2+y=x^3$; the section
basis is $(1,x,y)$ in the first and $(y,x,1)$ in the second.
On all four product charts the ideal generated by the two elliptic
equations, $s$, and its two invariant derivatives has Groebner basis
$\{1\}$. Hence $C$ has no geometric singular point.

For every nonzero kernel translation $g$, an invertible projective
matrix for translation on each cubic was first computed. It was
then certified by exact rational addition identities modulo
$y^2+y+x^3$; agreement on the nine torsion points alone was not
used as a proof. For each of the eight $g$ and each of the four
product charts, the ideal
\[
       (F_1,F_2,s,t_g^*s,\det(ds,d(t_g^*s)))
\]
was checked. Its support is empty except at $P$ for $g=G$ and
at $Q$ for $g=-G$. For these two exceptions, reduction of
$(x_j-a_j)^{32}$ to zero in the Groebner basis proves that the
support over the algebraic closure is just the indicated point.
Frobenius is injective on an algebraically closed field, so this
power test does not merely exclude rational extra points.

Finally, for all twenty-eight unordered pairs $g\ne h$ of nonzero
kernel elements and on all four charts, the ideal
$(F_1,F_2,s,t_g^*s,t_h^*s)$ has Groebner basis $\{1\}$.
Thus no three distinct source points have the same image.
The eight translation matrices are genuine global translations,
so these $32+112=144$ checks exhaust all possible nonnodal
identifications. They also exclude invariance of $C$ by a nonzero
kernel translation. Since $C$ is a smooth connected ample divisor,
it is integral; $\phi|_C$ is therefore birational to its image.
Etaleness of $\phi$ gives no other singularities. The only nontransverse
identification is the certified simple contact; all other
identifications are transverse pairs. This proves the assertion
about the norm curve. Lemma~\ref{lem:ss-mixed-self-tacnode} now
proves smoothness of the relative Hasse divisor at its image.
\end{proof}

\begin{remark}[Retained discovery, failures and independent checks]\label{auto:A_monodromy:0638c557453d}
\status{Proved under stated hypotheses}

The parent's discovery and exact global verifiers are
\path{code/ss_contact_witness.py},
\path{code/ss_contact_witness_screened.py},
\path{code/verify_ss_contact_geometry.py}, and
\path{code/verify_ss_contact_collisions.py}; the successful input is
\path{outputs/ss_contact_witness_03.json}. The full successful outputs are
\path{outputs/verify_ss_contact_witness_03_geometry.json} and
\path{outputs/verify_ss_contact_witness_03_collisions.json}.
The third discovery used seed $202611191119$, and succeeded at its
first base pair and twenty-ninth section trial; its complete run
record is retained. The first two candidates failed the global test
because they had unwanted triple collisions. Their discovery and
failed verification outputs were retained; neither was promoted.

The independent local program
\path{agents/monodromy/code/verify_level3_contact03_local.py}
reconstructs the coefficient-field arithmetic, tensor matrix and
actual translated power series; it passed in
$0.0036840360007772688$ seconds.
The complete independent rerun of the four source charts passed in
$3.4523235300002852$ seconds and that of all $144$ collision ideals in
$72.10777428700021$ seconds. The parent's corresponding successful
runs took $3.2001894870008982$ and $74.52918819499973$ seconds.
The independent outputs are in the monodromy agent's output directory,
with the same verification basenames; the author's outputs were
preserved. The rerun driver and log are
\path{agents/monodromy/code/rerun_contact03_global_audit.py} and
\path{agents/monodromy/outputs/independent_contact03_global_run.log}.

This witness joins one pair-tangency locus to the source cubic $Q$.
It does not by itself establish irreducibility of the target Hasse
divisor: the two tangent-triple incidence orbits in the proposed
finite reduction still require their own global witnesses.
\end{remark}


% /root/monodromy, 2026-11-19. Root and reduction-descent complete audits passed.
\subsection{Two exact tangent triples in the norm family}
\label{subsec:ss-global-tangent-triples}
\status{Proved and checked. The parent and reduction-descent agents independently
audited the full proof and exact verification programs. The reduction-descent
agent reran both complete global certificates and the separate local verifier.
These are two geometric witnesses, not a new full-image theorem.}

\begin{lemma}[Global witnesses for both affine triple types]
\label{lem:ss-global-tangent-triples}
\status{Proved under stated hypotheses}
Use the field, isogeny, monomial tensor basis and kernel elements $G,H$
of Lemma~\ref{lem:ss-global-contact03} and
Lemma~\ref{lem:ss-level3-collision-orbits}.
The following two coefficient vectors define geometrically smooth
source curves:
\[
 \begin{split}
 s_{\mathrm{line}}&=(24,2,49,5,40,45,54,16,46),\\
 s_{\mathrm{nonline}}&=(37,55,55,1,1,2,61,1,5).
 \end{split}
\]
Both are sharp: their Cartier kernels consist exactly of the two
ambient invariant forms. Their norm curves each have one three-branch
singularity consisting of a simply tangent pair and a transverse third
branch, and otherwise only ordinary nodes.
For the first curve the three source points are
\[
 \begin{split}
 P&=((24,23),(56,26)),\\
 Q=P+G&=((51,23),(23,34)),\\
 R=P-G&=((57,25),(55,21)).
 \end{split}
\]
For the second they are
\[
 \begin{split}
 P&=((3,44),(14,44)),\\
 Q=P+G&=((18,34),(50,30)),\\
 R=P+H&=((12,26),(25,53)).
 \end{split}
\]
In both cases $P,Q$ are the tangent pair. In particular these are
witnesses for the affine-line and noncollinear triple orbits,
respectively, and the degree-three relative Hasse divisor is smooth
at both norm curves.
\end{lemma}
\begin{proof}
All coefficients use the integer encoding of
$\mathbf F_2[z]/(z^6+z+1)$ in
Lemma~\ref{lem:ss-global-contact03}. The following exact values
are retained in the discovery and verification outputs:
\[
\begin{array}{c|ccccc}
 &D(s)&Q(s)&c\text{ in }ds(Q)=c\,ds(P)
 &\det(ds(P),ds(R))&\text{contact coefficient}\\\hline
 \mathrm{line}&55&32&31&41&3\\
 \mathrm{nonline}&62&1&35&49&39
\end{array}
\]
In particular $D(s)^2Q(s)\ne0$, so the relative source Cartier map
is bijective. The ambient invariant forms are killed by Cartier and
map isomorphically under the adjunction boundary. Thus the full source
Cartier kernel consists exactly of these two forms.

The local checks are independent of the discovery formula.
The program \path{agents/monodromy/code/verify_tangent_triple_local.py}
reconstructs the finite field by binary polynomial reduction,
checks the translated points and all displayed values, and computes
actual order-two elliptic-curve power series. It chooses a branch
$\gamma(t)$ of the source curve through $P$ to order two and verifies
\[
 s(\gamma(t))=O(t^3),\qquad
 s(\gamma(t)+G)=a t^2+O(t^3),
 \qquad a=3\ \text{or }39.
\]
The nonzero determinant in the table proves that the third branch
is transverse. Hence the intended singularity is the precise
one in Lemma~\ref{lem:ss-tangent-triple-bridge}.

For global smoothness, the same four projective product charts used
in Lemma~\ref{lem:ss-global-contact03} give the unit ideal for the
source singular-point equations, separately for both coefficient
vectors. The norm-collision verifier uses the same eight globally
certified translation matrices. It checks all eight pair-tangency
ideals and all twenty-eight triple-identification ideals on each of
the four charts. Pair-tangency support is exactly $P$ for $G$ and
$Q$ for $-G$, and is otherwise empty.

For the triple ideals, the allowed support is determined before the
Groebner computation: for given $g,h$, retain exactly those
$T\in\{P,Q,R\}$ such that $T+g,T+h\in\{P,Q,R\}$ and that lie
in the chosen chart. If this list is empty, the ideal must be the
unit ideal. For a nonempty list $T_1,\ldots,T_m$, $m\le3$, choose
a linear form $\ell$ separating them. Write $\alpha_i=\ell(T_i)$
and let $p_j$ be the degree-at-most-$m-1$ interpolation polynomial
with $p_j(\alpha_i)=(T_i)_j$. The program reduces
\[
  \prod_i(\ell-\alpha_i)^{32},\qquad
          (x_j-p_j(\ell))^{32}\quad(1\le j\le4)
\]
to zero in the computed Groebner basis. These identities prove
support containment over the algebraic closure: at a common zero,
$\ell=\alpha_i$ for some $i$, and injectivity of the thirty-second
power map then forces every coordinate to equal that of $T_i$.
The prescribed points themselves satisfy the equations. This is
an exact finite-support test, not a search for rational points.

All $144$ collision checks pass for each curve. Thus no additional
triple, quadruple or nontransverse double identification is possible.
The source curves are smooth connected ample divisors and hence
integral. The finite pair-intersection checks exclude translation
invariance, so each source is the normalization of its norm image.
Etaleness of the isogeny excludes any further branch singularity.
This proves both global assertions. Finally
Lemma~\ref{lem:ss-tangent-triple-bridge} applies and proves Hasse
smoothness at both norm curves.
\end{proof}

\begin{remark}[Exact execution record and independent audits]\label{auto:A_monodromy:ac724428800d}
\status{Proved under stated hypotheses}

The discovery program is
\path{agents/monodromy/code/ss_tangent_triple_search.py}.
It uses seeds $202611191140$ and $202611191141$, and accepts the
second and fourth sections at the first base point, respectively.
Its rational orbit and tangency screens only reject candidates.
The complete discovery records are
\path{agents/monodromy/outputs/ss_tangent_triple_collinear.json} and
\path{agents/monodromy/outputs/ss_tangent_triple_noncollinear.json}.
The discovery runtimes were $0.8046595620035077$ and
$0.7408189790003235$ seconds. The first noncollinear attempt had an
extra point in the intended fibre, and the next two had other rational
triple fibres; these rejections are retained in that output.

The global collision program is
\path{agents/monodromy/code/verify_ss_tangent_triple_collisions.py}.
Its full outputs are
\path{agents/monodromy/outputs/verify_ss_tangent_triple_collinear_collisions.json}
and
\path{agents/monodromy/outputs/verify_ss_tangent_triple_noncollinear_collisions.json}.
They passed in $72.36891810200177$ and $69.24018129799515$ seconds.
The separate four-chart source smoothness checks, using the retained
root verifier, passed in $4.423437987999932$ and $3.9984959799985518$
seconds; their complete outputs have the corresponding
\texttt{geometry.json} suffixes in the same directory.
The independent local-series verifier passed in
$0.005125144998601172$ seconds, with output
\path{agents/monodromy/outputs/verify_tangent_triple_local.json}.
The parent independently read the complete proof and passed it. The
reduction-descent agent independently read all programs and the proof,
then reran both $144$-check collision certificates in
$76.3690687260023$ and $74.86146398200071$ seconds. Its four-chart
source checks passed in $3.7411868540002615$ and
$3.3895406299998285$ seconds, and its separate local-series rerun
passed in $0.007957189001899678$ seconds, again producing $3$ and $39$.
The independent outputs are retained in the reduction-descent agent's
output directory with stems
\path{independent_ss_tangent_triple_collinear},
\path{independent_ss_tangent_triple_noncollinear}, and
\path{independent_tangent_triple_local}. The author outputs were
preserved. Both complete independent audits passed.
\end{remark}


% /root/monodromy, 2026-11-19. Complete proof; independent parent audit passed.
\subsection{A finite collision reduction for the degree-three norm family}
\label{subsec:ss-level3-collision-reduction}
\status{Proved and independently audited. The parent and reduction-descent agents read and checked the
complete reduction, including the conductor stable-rank formula, the
local-ring Frobenius-unit graph comparison and every incidence and
component-assignment step. All three global witnesses have complete
independent proof and program audits. The later
Theorem~\ref{thm:ss-all-odd-full-groups} gives the full image in
every odd degree, using the subsequent uniform collision proofs.}

Work over an algebraically closed field of characteristic two. Let
$E:y^2+y=x^3$, $X=E^2$, $L=\mathcal O(3O)\boxtimes\mathcal O(3O)$,
and choose $\omega\in\mathbf F_4$ with $\omega^2+\omega+1=0$.
Put $i(x,y)=(x+1,y+x+\omega)$, $b=1+i$, and
\[
 \phi=\begin{pmatrix}1&b\\-b^\dagger&1\end{pmatrix}:X\to X.
\]
Then $b^\dagger b=2$, $\phi^\dagger\phi=3$, and
$K=\ker\phi=\{(-bP,P):P\in E[3]\}$ has order nine. The norm gives
\[
 \nu:\mathbf P(H^0(L))=\mathbf P^8\longrightarrow |L^3|,
 \qquad\nu^*\mathcal O(1)=\mathcal O(9).
\]
There are no base points in this map: the norm of a nonzero section
is nonzero. Its exact target line bundle was proved in
Theorem~\ref{thm:ss-uniform-GL8-block}.
In the monomial tensor basis put
$D(s)=\det(s_{ij})$ and $Q(s)=s^{[2]t}(J_3\otimes J_3)s$,
where $J_3$ reverses the three basis vectors. The source Hasse
polynomial is $D^2Q$.

\begin{lemma}[The finite kernel has three collision types]
\label{lem:ss-level3-collision-orbits}
\status{Proved under stated hypotheses}
A group of product automorphisms preserving $K$, together with
translations by $K$, acts transitively on nonzero kernel differences.
Its unordered triples have precisely two orbits: a three-point affine
line in $K\simeq\mathbf F_3^2$, and a noncollinear triple.
One may use the representatives
\[
 G=((\omega,\omega^2),(0,0)),\qquad
 H=((\omega^2,\omega^2),(1,\omega)),
\]
with pair difference $G$ and triples $\{0,G,-G\}$ and $\{0,G,H\}$.
Every nonzero kernel difference has both elliptic coordinates nonzero.
\end{lemma}
\begin{proof}
The automorphisms $i$, $j(x,y)=(x+\omega,y+\omega^2x+\omega)$,
and $z(x,y)=(\omega x,y)$ preserve $O$ and the cubic. Their action
on $E[3]$ generates $\mathrm{SL}_2(\mathbf F_3)$.
This is also checked without any group identification in the retained
finite certificate: the nine points of $E(\mathbf F_4)$ are all of
$E[3]$; the displayed generators give twenty-four distinct permutations,
all preserving addition, and their action on nonzero points is transitive.
The action is faithful, since an automorphism of the cubic fixing all
nine flexes is the identity; it preserves the Weil pairing, so its
order identifies it with $\mathrm{SL}_2(\mathbf F_3)$.

The map $-b$ on $E[3]$ is invertible and has determinant two. For
$v\in\mathrm{SL}_2(\mathbf F_3)$ the conjugate
$u=(-b)v(-b)^{-1}$ again lies in $\mathrm{SL}_2(\mathbf F_3)$.
The product automorphism $(u,v)$ preserves the graph $K$.
Its action on the graph is the action of $v$ on the second coordinate.
The resulting affine special linear action has order $9\cdot24=216$.
It is transitive on nonzero differences. On unordered triples, the
collinear triples form one orbit of size twelve. A noncollinear triple
can be translated to contain zero, and a reordering changes the sign
of its nonzero determinant, so all such triples form the other orbit,
of size seventy-two. Since $b$ is invertible on $E[3]$, neither
coordinate of a nonzero graph point is zero. Direct addition gives
the stated $G,H$.

The retained exact program
\path{agents/monodromy/code/supersingular_level3_kernel_orbits.py}
checks the full addition table, the twenty-four automorphisms, conjugation
closure, all graph points and both affine orbits. Its output is
\path{agents/monodromy/outputs/supersingular_level3_kernel_orbits.json};
the author run passed in $0.007294220002222573$ seconds.
\end{proof}

For $g\ne0$, let $R_g$ denote the closure of the pair-tangency incidence
image: $C_s$ contains $P,P+g$ and the two translated tangents coincide.
For distinct $0,g,h$, let $T_{g,h}$ denote the closure of the locus
where $C_s$ contains $P,P+g,P+h$.

\begin{lemma}[Irreducibility and separation from the two source factors]
\label{lem:ss-level3-collision-incidence}
\status{Proved under stated hypotheses}
Each $R_g$ has one possible divisorial component arising from its
main incidence, and no additional divisorial component from the
exceptional jet locus. Each $T_{g,h}$ is an irreducible incidence
image. If these images contain a smooth noninvariant curve, they
are irreducible divisors. None of these divisors equals $D=0$ or
$Q=0$.
\end{lemma}
\begin{proof}
For fixed $g$, parametrize a common tangent by $(P,\ell)$ with
$P\in X$ and $\ell\in\mathbf P(T_PX)$. Containment and vanishing
of the derivative along $\ell$ at both translated points give four
linear conditions on $H^0(L)$. Since both coordinates of $g$ are
nonzero, the tensor jet calculation in
Theorem~\ref{thm:ss-uniform-GL8-block} gives rank six for the joint
first jets except at finitely many $P$, and rank at least five there.
Thus the four imposed conditions have rank four on the main open
and at least three at those exceptions. The main incidence is a
$\mathbf P^4$-bundle over an irreducible three-dimensional base,
so is irreducible of dimension seven. The exceptional part has
dimension at most $1+(8-3)=6$.

For triples, the three value tensors are independent everywhere.
Indeed three distinct rank-one Segre tensors can be collinear only
if one factor is constant, whereas both projections of $K$ are
injective. The incidence is therefore a $\mathbf P^5$-bundle over
$X$, irreducible of dimension seven.
On a smooth curve not invariant under a nonzero kernel translation,
all translated intersections are finite. Hence the incidence maps
have finite fibres there. The invariant-curve locus cannot dominate:
a nonzero three-torsion translation acts with three distinct eigenvalues
on $H^0(E,\mathcal O(3O))$, since a repeated projective eigenspace
would be a fixed line intersecting the cubic. Its tensor action has
eigenspaces of dimension three, so invariant sections form only a
union of $\mathbf P^2$'s. This proves the divisor assertions under
the stated nonemptiness hypothesis.

No vector subspace of dimension greater than four lies in $Q=0$.
If a subspace $U$ did, polarization of the cubic would imply
$U^{[2]t}JU=0$, so $U\subset(JU^{[2]})^\perp$ and
$2\dim U\le9$. The generic coefficient fibres in the two incidences
have vector dimensions five and six. Thus neither incidence image
can be the cubic $Q$.

For $D$, normalize the two point-value directions in each factor to
$e_0,e_1$. At a generic translated pair the first-factor tangent at
either point has nonzero $e_2$ coefficient. Choose a common tangent
with nonzero first-factor component. The two value equations are
$S_{00}=S_{11}=0$. Set $S_{01}=S_{10}=1$; the two tangent equations
solve for $S_{20},S_{21}$ after the other entries are chosen.
Neither tangent equation involves $S_{22}$, and its coefficient
in $\det S$ is one. Thus a generic pair-tangent coefficient fibre
contains an invertible matrix.

For a triple, its first-factor value directions are generically
independent: dependence would say $3P_1+g_1+h_1=0$ on the elliptic
curve, a finite condition on that coordinate. Normalize them to a
basis. The three rows of $S$ must lie respectively in the three
distinct hyperplanes orthogonal to the second-factor values.
One can choose three independent row vectors in these planes:
otherwise their sum would have dimension at most two and all
three planes would coincide. This gives an invertible $S$ in the
coefficient fibre. Hence neither incidence image equals $D$.
\end{proof}

\begin{proposition}[A reduction to three global geometric witnesses]
\label{prop:ss-level3-three-witness-reduction}
\status{Proved under stated hypotheses}
Suppose the following three witnesses exist.
\begin{enumerate}
\item A smooth source curve on $Q\setminus D$ whose norm has one
simple self-tacnode of difference $G$ and otherwise only ordinary
nodes, with unequal values of its extra Cartier-killed form on
the two tangent branches.
\item A smooth source curve with $D Q\ne0$ whose norm has one
three-branch singularity of the tangent-triple type in
Lemma~\ref{lem:ss-tangent-triple-bridge}, with source
triple in the affine-line orbit, and otherwise only ordinary nodes.
\item The same kind of sharp-source tangent-triple witness for the
noncollinear orbit.
\end{enumerate}
Then the degree-three relative Hasse divisor on $|L^3|$ is
geometrically irreducible and reduced.
\end{proposition}
\begin{proof}
The witnesses give the nonemptiness required by the incidence lemma.
More precisely, the six first-jet rows at their respective tangent
pairs have ranks $6,6,6$, so all three witnesses lie in the main
rank-six pair incidence. This is checked directly in the retained
program \path{agents/monodromy/code/verify_ss_collision_main_incidence.py},
whose output includes the rows and their reduced echelon forms;
it passed in $0.0005378700006986037$ seconds. The reduction-descent
auditor also computed these three ranks independently.
They also ensure that the indicated source curves are smooth at
points of each collision divisor. The source singular-curve
discriminant is irreducible: the singular-point incidence is a
projective bundle over $X$. Its generic member has a unique ordinary
node, by the first/second-jet bounds (the incidence of two singular
points has dimension at most six). Its degree is
$c_2(J^1L)=3L^2=54$: on a generic coefficient pencil its singular
members are the zeros of the rank-three first-jet section on
$X\times\mathbf P^1$, and its top Chern number is $54$.
The ordinary-node Hessian is invertible, so the incidence has generic
separable degree one. In particular this discriminant differs from
both cubic source factors.

Outside the source discriminant and the collision divisors, the
norm curve is nodal with normalization $C_s$. The normalization
sequence has a Frobenius-unit graph term, so the relative determinant
vanishes exactly when $D^2Q$ does. At the generic point of the source
discriminant, no collision divisor occurs, by the smooth-source
witnesses. The source node maps isomorphically to a target node;
all additional gluing terms are again Frobenius-unit. Thus this
discriminant is not an additional factor of $\nu^*h_3$.
There are no other codimension-one exceptions: noninvariance and
smoothness exclude all failures of the finite normalization map,
and nonordinary identifications require either tangency or a triple.
Consequently every irreducible factor of $\nu^*h_3$ is one of
$D,Q,R_g,T_{g,h}$.

We also need the converse inclusion for the collision divisors.
For a reduced local curve singularity with $r$ branches, write
$A\subset\widetilde A$ for its completed local ring and normalization.
The branch constants in $\widetilde A/A$ form $k^r/k$, on which
Frobenius is bijective. Modulo these constants every element has
positive order on each branch; a sufficiently high Frobenius power
lies in the conductor and is therefore zero in the quotient.
Thus the Frobenius-stable rank of $\widetilde A/A$ is exactly $r-1$.
The global normalization sequence, and exactness of the stable-part
functor over the perfect field, now give for a connected reduced
curve with normalization components $C_i$
\[
 g(Z)-f(Z)=\sum_i\bigl(g(C_i)-f(C_i)\bigr)
              +\sum_x(\delta_x-r_x+1).
\]
Here $g=h^1(\mathcal O)$ and $f$ is its Frobenius-stable rank.
At a generic point of any $R_g$ or $T_{g,h}$ the source is smooth,
by the checked witnesses and irreducibility, and it lies outside
$D\cup Q$ by the preceding lemma. Hence its defect is exactly two.
A tangency of distinct smooth branches has
$\delta-r+1\ge1$, and so does an identification of at least three
smooth branches in the smooth target surface. The target curve
therefore has defect at least three. Its relative coherent dimension
subtracts the two-dimensional ambient cohomology, whose stable rank
is zero. Thus its relative Frobenius is not bijective. This proves
that every generic collision image lies in the target Hasse divisor;
no unproved ordinarity assertion about a triple image is used.

Finally consider the local family at the generic point of $D$.
The source curve is smooth and the norm has only transverse double
identifications. After a finite etale base change ordering the nodes
and their branches, the normalization quotient is a direct sum of
copies of the base ring, with Frobenius the ordinary square map.
The resulting graph term in $R^1\mathcal O$ is a finite free
Frobenius-unit subbundle. The ambient $H^1(\mathcal O)$ subbundles
on source and target identify under the odd isogeny. Passing to
relative coherent cohomology therefore gives an exact sequence of
vector bundles whose graph term has invertible Frobenius determinant
and whose quotient is the source relative bundle. Determinants in
this exact sequence multiply. The equality is valid over the local
ring, not just on its generic fibre, so the vanishing order along
$D$ in $\nu^*h_3$ is exactly the order two in $D^2Q$.

The first witness and Lemma~\ref{lem:ss-mixed-self-tacnode} identify
the global target Hasse component containing the image of $Q$ with
the one containing the image of $R_G$. The target divisor is smooth
at that witness, so both assignments are unique. Every source
normalizer automorphism descends under $\phi$ and preserves its
norm line bundle. It preserves $Q$, the unique multiplicity-one
component of the source divisor $D^2Q$, and hence fixes the unique
target component just identified. Transitivity on kernel differences
therefore assigns every $R_g$ to that same target component.
The other two witnesses and the tangent-triple bridge identify the
two triple orbits with it as well. Smoothness at each witness makes
these assignments valid even if a pullback has multiplicity.

Any other target Hasse component would therefore pull back to a
nonzero divisor supported only on $D$, with multiplicity at most two.
Its pullback degree would be at most six. But every nonzero target
hypersurface has pullback degree $9\deg(H)\ge9$, since
$\nu^*\mathcal O(1)=\mathcal O(9)$. This contradiction proves
irreducibility. Smoothness at the first witness proves that the
multiplicity of the single component is one.
\end{proof}

\begin{corollary}[The degree-three Hasse divisor]
\label{cor:ss-degree3-hasse-irreducible}
\status{Proved under stated hypotheses}
The relative Hasse divisor on $|L^3|$ is geometrically irreducible
and reduced.
\end{corollary}
\begin{proof}
Lemma~\ref{lem:ss-global-contact03} supplies the first witness, and
Lemma~\ref{lem:ss-global-tangent-triples} supplies the second and
third. All three sources, every global collision, the local
Cartier hypotheses and the nonzero transverse derivatives have
been independently checked. Apply the preceding proposition.
\end{proof}

\begin{remark}[Retained failures and scope]\label{auto:A_monodromy:1e009f1c7aa8}
\status{Proved under stated hypotheses}

The first local contact candidate had coefficient vector
$(8,3,45,21,31,9,33,17,4)$, Cartier-value difference $38$ and
contact coefficient $29$. Its later global check found an unwanted
triple. A second contact candidate also failed a global triple
check. Their discovery and verification records were retained;
neither failed candidate is used in the proof. The third contact
candidate and the two tangent-triple candidates are the successful
witnesses cited above.

This corollary determines the geometry of the degree-three Hasse
divisor. The full degree-three group $\mathrm{GL}_{80}(\mathbf F_2)$
was already proved by a separate certificate. The subsequent uniform collision witnesses extend this argument to
every odd degree in Theorems~\ref{thm:ss-all-odd-Hasse} and
\ref{thm:ss-all-odd-full-groups}. The remaining even-degree images are determined in
Theorem~\ref{thm:ss-product-all-full-groups}.
\end{remark}


% /root/monodromy, 2026-11-19. Complete independent parent proof and code audit passed.
\subsection{A symbolic mixed contact for an arbitrary translation}
\label{subsec:ss-uniform-contact}
\status{Proved local calculation with an exact universal polynomial
certificate; the parent independently read the complete proof and program.
The following subsection proves source smoothness. Other collisions
for odd kernels with injective coordinate projections are excluded in
Proposition~\ref{prop:ss-uniform-graph-mixed-contact}. Uniform
tangent-triple bridges are proved in
Proposition~\ref{prop:ss-uniform-tangent-triple}, and every odd-degree
full image is determined in Theorem~\ref{thm:ss-all-odd-full-groups}.}

Let $E:y^2+y=x^3$ over an algebraically closed field of characteristic
two, let $X=E^2$, and let $L=\mathcal O(3O)\boxtimes\mathcal O(3O)$.
Take a translation
\[
 g=((a,b),(c,d)),\qquad b^2+b=a^3,\quad d^2+d=c^3,
\]
whose two coordinates differ from $O$, and suppose first that
$(c,d)\notin E[3]$. Put $P=(O,O)$ and $Q=g$.

\begin{lemma}[A uniform local mixed-contact slice]
\label{lem:ss-uniform-local-contact}
\status{Proved under stated hypotheses}
There is a nonempty irreducible three-dimensional parameter family
of coefficient matrices $S$, on a dense open of which the following
hold. The section $s$ vanishes at $P,Q$ and is smooth at both points;
the two translated tangent directions coincide, with contact order
exactly two. One has $\det S\ne0$ and $Q_H(S)=0$, where
$Q_H(S)=\sum S_{ij}^2S_{2-i,2-j}$ is the source Hasse cubic.
There is an explicitly given section $u$ with
\[
                         M(s)u=s^{[2]},
\]
and the local residue differential represented by $u/s$ has unequal
values on the two translated tangent vectors. The assertions concern
the indicated two points and the global coefficient operator only;
the source is not asserted to be smooth elsewhere.
\end{lemma}
\begin{proof}
Use three independent parameters $F,T,U$, all nonzero, and adjoin $v$
satisfying
\begin{equation}\label{eq:ss-uniform-contact-AS}
 v^2+v+U^3+U(c^4+c)+F^2(T+a^4+a)+FT^2=0.
\end{equation}
This Artin--Schreier equation is irreducible over $k(F,T,U)$:
its right side has a pole of odd order three at $U=\infty$,
whereas a nonzero pole of $h^2+h$ has even order. Hence it defines
an irreducible three-dimensional family. Set
\[
\begin{aligned}
 C&=T+a,& E_0&=U+c^2,& D_0&=v+Uc+d^2+Fa^2,\\
 B&=Tc^2+E_0a,& A&=Td^2+D_0a+Fb,\\
 S&=\begin{pmatrix}A&B&C\\D_0&E_0&1\\F&0&0\end{pmatrix}.
\end{aligned}
\]
Here the section is expressed in the tensor basis $(1,x,y)$ on
both factors. Direct substitution gives
\[
 s(P)=s(Q)=0,\qquad ds(P)=(1,0),\qquad ds(Q)=(v,0),
 \qquad\det S=FTU.
\]
The first derivative at $P$ uses the $Y=1$ chart at $O$ and the
invariant local tangent coordinate. At $Q$ it uses the invariant
derivations $\delta x=1,\delta y=x^2$.
Moreover
\[
 Q_H(S)=D_0^2+D_0+E_0^3+C^2F+F^2C,
\]
which is precisely the left side of
\eqref{eq:ss-uniform-contact-AS} after the elliptic equations
are used. The parameter $v$ is generically nonzero, so both indicated
points are smooth.

Let $J$ reverse the three coordinates, put $R=JS^{[2]}J$, and define
\[
                       u=\frac{S R^t S}{FTU}.
\]
We regard the right side as the matrix of its nine coefficients.
The root-removed Cartier matrix, after multiplication on the output
by $J\otimes J$, is the Hessian of $\det S$.
The inverse-Hessian identity gives
$M(s)u=s^{[2]}$ when $Q_H(S)=0$; equivalently, for a universal matrix
$S$, every coefficient of
\[
 M(s)\operatorname{vec}(S R^t S)
                   -(\det S)s^{[2]}
\]
is divisible by $Q_H(S)$. This identity is also verified coefficient
by coefficient by the exact program below.
The values of the residue form on the common translated tangent
are $u(P)$ and $u(Q)/v$. They satisfy
\[
 u(P)=\frac{F}{FTU},\qquad
 \frac{u(Q)}{v}=\frac{T\bigl(v^2+F^2(a^4+a)\bigr)}{FTU}.
\]
Their difference is therefore
\begin{equation}\label{eq:ss-uniform-contact-value}
       \frac{F+T\bigl(v^2+F^2(a^4+a)\bigr)}{FTU}.
\end{equation}
It is generically nonzero. Indeed if its numerator vanished, the
Artin--Schreier equation would give
\[
 v=F/T+U^3+U(c^4+c)+F^2T+FT^2.
\]
Squaring this and substituting back into
$v^2+F^2(a^4+a)=F/T$ produces a polynomial in $U$ with nonzero
leading term $U^6$, which is impossible for the independent parameters.

We finally check the contact order using actual elliptic addition,
not merely tangent vectors. At $O$ write the $Y=1$ coordinates as
$(t,z)$ with $z+z^2=t^3$, hence $z=t^3+O(t^6)$. Since the section
has derivative $(1,0)$ at $P$ and its restriction to the second
elliptic fibre is $Fz$, its local branch through $P$ has first
coordinate of order at least three. Translation by $g$ therefore
fixes the first elliptic coordinate through order two. On the second
factor the group law gives
\[
 x((c,d)+(t,z))=c+t+c^2t^2+O(t^3),\qquad
 y((c,d)+(t,z))=d+c^2t+ct^2+O(t^3).
\]
On the first-coordinate fibre at $(a,b)$ the section equals
$T(d^2+c^2x+y)$. Its value along the translated source branch is
thus
\[
                         T(c^4+c)t^2+O(t^3).
\]
The coefficient is nonzero because $(c,d)\notin E[3]$: on this cubic
its finite flexes have $c\in\mathbf F_4$, and the corresponding
nine points including $O$ are exactly $E[3]$. This proves contact
order two and completes all stated local assertions.
\end{proof}

\begin{remark}[The other coordinates and the global consequences]\label{auto:A_monodromy:b303eb6cd268}
\status{Proved under stated hypotheses}

If the second coordinate is a nonzero three-torsion point but the
first is not, interchange the two factors. If both coordinates are
nonzero three-torsion points, product automorphisms preserving $O$
act transitively on these pairs, so the already checked contact03
witness supplies local mixed contacts after an automorphism.
These facts remove a local polynomial-identity obstruction for any
translation with both coordinates nonzero. Source smoothness is proved in
Lemma~\ref{lem:ss-uniform-source-smoothness}. Proposition~\ref{prop:ss-uniform-graph-mixed-contact} excludes all other
nonordinary identifications for every odd kernel with injective
coordinate projections. Proposition~\ref{prop:ss-uniform-tangent-triple}
supplies the remaining triple bridges, giving the odd-degree theorem
\ref{thm:ss-all-odd-full-groups}.
\end{remark}

\begin{remark}[Exact certificate]\label{auto:A_monodromy:ee94dcb570ad}
\status{Proved under stated hypotheses}

The retained program
\path{agents/monodromy/code/supersingular_uniform_contact_identities.py}
works in the polynomial ring over $\mathbf F_2$, reducing only by the
two displayed elliptic equations. It checks all containment, gradient,
determinant and cubic identities; all nine universal inverse-Hessian
coefficients; both residue evaluations; the nonzero degree-six
obstruction to their equality; and the rational elliptic-addition
expansion through order two. Its output is
\path{agents/monodromy/outputs/supersingular_uniform_contact_identities.json}.
The completed run passed in $0.34270756500336574$ seconds. No geometric
existence assertion beyond the local lemma is inferred from this code.
\end{remark}


% /root/monodromy, 2026-11-19. Complete independent parent proof audit passed.
\subsection{Global source smoothness in the symbolic contact slice}
\label{subsec:ss-uniform-source-smoothness}
\status{Proved and independently audited. The parent checked all chart
incidences, the four exact discriminant multiplicities and the degree
argument. This addresses source smoothness. Other collisions for odd graph
kernels are excluded in Proposition~\ref{prop:ss-uniform-graph-mixed-contact}.
The later uniform tangent-triple and odd-degree full-image results are
Proposition~\ref{prop:ss-uniform-tangent-triple} and
Theorem~\ref{thm:ss-all-odd-full-groups}.}

Keep the notation of Lemma~\ref{lem:ss-uniform-local-contact},
including $g_1,g_2\ne O$ and $g_2\notin E[3]$. Homogenize its slice
by a fifth parameter $h$ and write
\[
\begin{aligned}
 C&=T+ha,&E_0&=U+hc^2,&D_0&=v+Uc+hd^2+Fa^2,\\
 B&=Tc^2+E_0a,&A&=Td^2+D_0a+Fb,\\
 S&=\begin{pmatrix}A&B&C\\D_0&E_0&h\\F&0&0\end{pmatrix}.
\end{aligned}
\]
These matrices form a projective linear space $\Lambda\simeq\mathbf P^4$
with independent homogeneous coordinates $[F:T:U:v:h]$.
On the affine product put
\[
 X=x+a,\quad Y=y+b,\quad Z=z+c,\quad W=w+d,
 \qquad \ell=w+c^2z+d^2.
\]
Its sections have the equivalent expression
\begin{equation}\label{eq:ss-contact-linear-system}
 s=T\ell+D_0X+E_0Xz+hXw+FY.
\end{equation}
Thus its function basis is $X,Y,\ell,Xz,X\ell$.
The source Hasse cubic restricts to the irreducible cubic
\begin{equation}\label{eq:ss-contact-cubic-homogeneous}
 q=h v^2+h^2v+U^3+h^2U(c^4+c)
                       +F^2T+hF^2(a^4+a)+FT^2.
\end{equation}
Its irreducibility follows by setting $h=1$ and using the odd-pole
Artin--Schreier argument of the preceding lemma; $h$ is not a factor.

\begin{lemma}[A smooth source exists in every nonflex contact slice]
\label{lem:ss-uniform-source-smoothness}
\status{Proved under stated hypotheses}
The restriction to $\Lambda$ of the reduced source discriminant has
exactly four forced hyperplane factors, with multiplicities
\[
                            F^3T^3h^3v^2,
\]
and one further irreducible factor, with some positive multiplicity.
The total degree of that remaining factor with multiplicity is $43$.
Consequently the cubic $q=0$ is not contained in the source
discriminant. A geometric generic member of the local family in
Lemma~\ref{lem:ss-uniform-local-contact} is globally smooth.
\end{lemma}
\begin{proof}
We first classify all possible divisorial singularity incidences of
this linear system. Its base points are
\[
              P=(O,O),\qquad Q=(g_1,g_2),
                    \qquad R=(g_1,-2g_2).
\]
The last two are distinct because $g_2$ is not a flex. This follows
on the affine chart from $X=Y=\ell=0$, and on the infinity charts
from the displayed coefficient matrix.

On the affine locus $X\ne0$, ratios of the five basis functions give
\[
                 Y/X,\qquad\ell/X,\qquad z,\qquad\ell.
\]
The second tangent direction is detected by $z$. The first is lost
only when both
\[
                    A_1:=Y+Xx^2=0,\qquad\ell=0.
\]
The first equation has only finitely many points on the first elliptic
curve: it says that the fixed point $g_1$ lies on the tangent at
$(x,y)$, and its displayed nonzero polar polynomial is not identically
zero on the cubic. The second equation has only two geometric points
on the second factor. Thus these exceptions are finite. At each,
first-jet rank is still at least two, since the value and the derivative
of the ratio $z$ are independent. They cannot give a divisorial
singularity incidence in $\mathbf P^4$.

The locus $X=0$ consists of the first-factor fibres through $g_1$
and $-g_1$. On the former the section is $T\ell$. Outside its two
base points, any singular section therefore lies in $T=0$.
At $Q$ the only first derivative is $v$, giving the hyperplane $v=0$.
At $R$ the second derivative direction is $T\partial\ell(R)$,
which is nonzero when $T\ne0$; the two first derivatives at this
base point are independent linear conditions, so its singularity
incidence has codimension two. On the fibre through $-g_1$ one has
$Y=1$, and the value together with the two derivatives have rank
three except when $z=c$, a finite set of points. Their ranks there
are at least two.

On the first-factor fibre through $O$, the section restricts to
$F$ times the section $1$ of $\mathcal O_E(3O)$. Consequently it
is avoided for $F\ne0$ except at $P$; there the only first derivative
is $h$. On the second-factor fibre through $O$, at a finite first
point the value is $T+hX$ and the first derivative is $h$.
The full first-jet rank is three when $X\ne0$; when $X=0$ there
are only two finite exceptional points and rank is at least two.
These calculations cover all four product charts. They show that,
apart from $F,T,h,v$, every divisorial singularity must arise from
the main rank-three jet incidence. That incidence is a
$\mathbf P^1$-bundle over an irreducible open of the surface, and
is therefore irreducible of dimension three.

For clarity its image really is a divisor, and is distinct from all
four hyperplanes. At a moving point satisfying
$X A_1\ell(z+c)\ne0$, solve the three singularity equations as
\[
 F=\frac{T\ell}{A_1},\qquad
 E_0=\frac{T(z^2+c^2)}X+hz^2,\qquad
 D_0=E_0z+hw+Fx^2.
\]
The two remaining parameters are $T,h$. The mixed Hessian is
\[
       \partial_1\partial_2s=E_0+hz^2
                       =\frac{T(z^2+c^2)}X\ne0
\]
when $T\ne0$. Thus the moving point is an ordinary node. The
implicit function theorem for its two first derivatives makes this
node persist only along one local discriminant equation, so the
incidence has divisorial image. Its generic coefficients have
$F,T,h\ne0$. Moreover
\[
 v=\frac{T(z+c)^3}{X}
           +h\bigl(z^2(z+c)+w+d\bigr)+FX^2.
\]
The coefficient of $h$ is a nonzero polar function of the second
elliptic point, so it is generically nonzero. Hence $v$ is also
nonzero at a generic point of this incidence. Its irreducible image,
call it $\Delta_{\mathrm{mov}}$, is distinct from $F,T,h,v$.
The same incidence dimension bounds prove that a generic section
of $\Lambda$ is smooth. In particular restriction of the full
source discriminant is a nonzero homogeneous polynomial.

We compute its four forced multiplicities. At a generic point of
$F=0$, the first-factor fibre through $O$ is a component. Its
residual curve has numerical type $(2,3)$ and meets it at the three
zeros of
\[
                         D_0+E_0z+hw.
\]
These coefficients range freely, so the three zeros are simple,
finite, and avoid the two zeros of $\ell$. The residual curve is
smooth: away from these zeros its first invariant derivative is
nonzero; at a zero the derivative in the second factor is nonzero.
There are therefore exactly three transverse component nodes.
Varying $F$ smooths each to first order, since its restriction at
those finite points is a nonzero local unit. The discriminant order
along $F=0$ is three.

At a generic point of $T=0$, the first-factor fibre through $g_1$
is a component. On its complement the residual equation is
\[
       D_0+E_0z+hw+F\rho=0,\qquad
                     \rho=\frac{y+b}{x+a}.
\]
The rational map $\rho:E\to\mathbf P^1$ has degree two and is
separable, as its derivative $A_1/X^2$ is not identically zero.
It has only finitely many ramification points. For generic
$D_0,E_0,h,F$, no critical value on the second factor coincides
with the image of such a ramification point: after choosing the
other coefficients, only finitely many values of $D_0$ are forbidden.
At the poles the two simple poles of $\rho$ give smooth local
branches. Thus the residual curve is smooth. Its intersections
with the component are the three simple zeros of
\[
                         D_0+E_0z+hw+Fa^2.
\]
They avoid $Q,R$ for a generic choice. Variation of $T$ evaluates
to $\ell$, which is nonzero at all three intersection points.
Hence these nodes are each smoothed to first order and the
multiplicity along $T=0$ is also three.

At a generic point of $h=0$, the only singular point is $P$:
all other possible divisorial incidences have already been shown
to be the distinct hypersurfaces $F,T,v,\Delta_{\mathrm{mov}}$.
In the two $Y=1$ elliptic parameters at $P$ the local equation has
leading terms
\[
          h x+E_0xy+Fy^3+Cx^3+\text{higher terms},
\]
with $E_0,F\ne0$. The Hessian has nonzero mixed coefficient
$E_0$, so the critical point of this equation is unique as a formal
function of $h$. Its critical value has leading term
$Fh^3/E_0^3$: the second coordinate is $h/E_0+O(h^2)$ and the
first is of order at least two. Thus the discriminant has order
exactly three along $h=0$.

At a generic point of $v=0$, its only singularity is similarly the
fixed point $Q$. Choose translated invariant parameters $x,z$ there.
Its quadratic part and transverse parameter have the form
\[
            vx+\alpha x^2+Uxz+T(c^4+c)z^2,
\]
where $U,T(c^4+c)\ne0$. Its critical value therefore has leading
term $T(c^4+c)v^2/U^2$. The discriminant order along $v=0$ is
exactly two. The higher coefficients and the dependence of
$E_0,D_0$ on these transverse parameters do not change the displayed
leading terms, since the critical point has the indicated orders.

The reduced discriminant of the full Segre cubic system has degree
$c_2(J^1L)=3L^2=54$, as proved in the source-discriminant calculation
of Proposition~\ref{prop:ss-level3-three-witness-reduction}.
Its nonzero restriction to $\Lambda$ thus has total degree $54$.
We have exhausted its divisorial support, and found the four orders
$3,3,3,2$. Consequently its remaining irreducible factor
$\Delta_{\mathrm{mov}}$, with its positive multiplicity, has total
degree $54-3-3-3-2=43$.
If the irreducible cubic $q$ were contained in the source discriminant,
it would have to equal $\Delta_{\mathrm{mov}}$, since it is not a
hyperplane. That would make $43$ divisible by three, a contradiction.
Hence a dense open of $q$ parametrizes globally smooth curves.
Intersect it with the nonzero-parameter, simple-contact and unequal
Cartier-value opens of the preceding lemma. Irreducibility of $q$
makes this intersection nonempty, proving the final assertion.
\end{proof}

\begin{remark}[The norm-collision consequences]\label{auto:A_monodromy:9f1f96604bfd}
\status{Proved under stated hypotheses}

This result supplies globally smooth source curves for arbitrary
translations with both coordinates nonzero (interchanging factors
when necessary, and using the checked three-torsion witness in the
finite flex case). Proposition~\ref{prop:ss-uniform-graph-mixed-contact}
now excludes every further nontransverse or triple identification
for odd kernels with injective coordinate projections. The uniform
tangent-triple bridges of Proposition~\ref{prop:ss-uniform-tangent-triple}
then complete the odd-degree Hasse-component argument in
Theorem~\ref{thm:ss-all-odd-Hasse}.
\end{remark}


% /root/monodromy, 2026-11-19. Complete independent root, RD and OC audits passed.
\subsection{A point-incidence criterion for excluding all other kernel collisions}
\label{subsec:ss-uniform-collision-exclusion}
\status{Proved. Three independent complete proof audits and exact
program audits passed. This removes the other-collision gate for a
prescribed mixed contact under every finite odd kernel with injective
coordinate projections. Uniform tangent-triple bridges are established
in Proposition~\ref{prop:ss-uniform-tangent-triple}; the full image
in every odd degree follows in Theorem~\ref{thm:ss-all-odd-full-groups}.}

Work over an algebraically closed field $k$ of characteristic two.
Let $E:y^2+y=x^3$, $X=E^2$ and
$L=\mathcal O_E(3O)\boxtimes\mathcal O_E(3O)$.
Let $K\subset X(k)$ be a finite subgroup of odd order such that both
projections $K\to E(k)$ are injective. Fix $0\ne g\in K$, written
$g=((a,b),(c,d))$. In particular both coordinates are nonzero.
Use the homogeneous contact slice $\Lambda=\mathbf P^4$ with coordinates
$[F:T:U:v:h]$ from
\eqref{eq:ss-contact-linear-system}--\eqref{eq:ss-contact-cubic-homogeneous}.
We also use that slice when $g$ has order three. Write
$\mathcal Q=(q=0)\subset\Lambda$ and $P=(O,O)$, $Q=g$.

\begin{lemma}[The contact cubic and its evaluation map]
\label{lem:ss-contact-cubic-evaluation}
\status{Proved under stated hypotheses}
The cubic $\mathcal Q$ is smooth. It contains no projective plane.
Its intersection with every projective hyperplane is an integral
cubic surface. For each $0\ne j\in K$, the locus where the two value
functionals at $R,R+j$ on $\Lambda$ have rank at most one is finite.
Their common rank is zero only when $j=\pm g$, at the corresponding
ordered pair of base points $P,Q$. No three distinct points in a
$K$-orbit are all base points of $\Lambda$.
\end{lemma}
\begin{proof}
Set $A=a^4+a$ and $B=c^4+c$. In the ordered variables $(F,T,U,v,h)$,
the cubic $q=z^{[2]t}A_qz$ has coefficient matrix
\[
 A_q=\begin{pmatrix}
 0&1&0&0&A\\
 1&0&0&0&0\\
 0&0&1&0&0\\
 0&0&0&0&1\\
 0&0&B&1&0
 \end{pmatrix},\qquad \det A_q=1.
\]
Its gradient is $A_q^t z^{[2]}$, so it has no projective singular
point. If a vector subspace $W$ is contained in its affine cone,
expanding $q(u+t w)$ gives
$u^{[2]t}A_qw=0$ for all $u,w\in W$.
Consequently $2\dim W\le5$. This excludes projective planes.
A reducible cubic in a projective hyperplane $\mathbf P^3$ would
have a linear factor and hence contain a projective plane. Therefore
every such section is integral.

For the value-map calculation put
\[
 X_1=x+a,\quad Y_1=y+b,\quad Z_1=z+c,\quad
 \ell=w+c^2z+d^2,\quad P_1=Y_1+a^2X_1.
\]
The section can be rewritten as
\begin{equation}\label{eq:ss-contact-five-functions}
       s=(T+hX_1)\ell+FP_1+vX_1+UX_1Z_1.
\end{equation}
Equivalently the five functions defining the rational value map are
$X_1,Y_1,\ell,X_1z,X_1\ell$.
Where both factors are affine and $X_1\ell\ne0$, their ratios
recover $z,\ell,X_1,Y_1$, and therefore the point of $X$.
The images of points outside this open lie outside its image: in
these coordinates their first or fifth homogeneous coordinate is zero.

It remains to examine the exceptional fibres. The three vertical
fibres through $O,g_1,-g_1$ have respectively constant image $F$,
constant image $T$, and nonconstant image on the line spanned by
$F,T$. The horizontal fibre through $O$ maps nonconstantly to the
line spanned by $T,h$. The zero divisor of $\ell$ has support
$g_2,-2g_2$; its two horizontal fibres have nonconstant images on
the lines spanned respectively by $F,v$ and
$F,v+\bigl(x(-2g_2)+c\bigr)U$.
For a nonflex $g_2$ these are distinct lines. For a flex the two
horizontal fibres coincide, and there is just the line spanned by
$F,v$. These assertions follow directly from
\eqref{eq:ss-contact-five-functions}, including the infinity
charts, where the section is divided by the local trivializing
monomial of $\mathcal O(3O)$.

A nonzero element of $K$ translates both elliptic coordinates
nontrivially. Thus an exceptional vertical or horizontal fibre cannot
be translated to itself. Distinct fibres of the same type above
either have distinct constant images, or a constant image and a
nonconstant image, or nonconstant images on distinct lines.
In every case equality of the two values occurs at only finitely
many points. This proves the rank-one assertion; any positive
dimensional equality locus would have had to be such an exceptional
fibre.

The base locus is $P,Q,(g_1,-2g_2)$, with the last two points
coinciding when $3g=0$. Injectivity of the first projection of $K$
shows that the third point has no $K$-difference with either of the
first two unless it is $Q$. This proves both remaining assertions.
\end{proof}

\begin{lemma}[The moving two-point covers are generically separable]
\label{lem:ss-contact-pair-covers}
\status{Proved under stated hypotheses}
For a geometric generic section of $\mathcal Q$, the intersections
of its divisor with all its nontrivial $K$-translates are finite.
For $j\ne\pm g$, every intersection with its $j$-translate is
transverse. For $j=g$, there is the prescribed intersection at $P$;
its length is two if $3g\ne0$ and three if $3g=0$. All other
intersections with the $g$-translate are transverse. The assertion
for $-g$ is the translated version of this one.
\end{lemma}
\begin{proof}
Consider the proper incidence
\[
 I_j=\{(R,s)\in X\times\mathcal Q:s(R)=s(R+j)=0\}.
\]
On its generic rank-two value locus over $X$, the coefficient fibre
is a plane cubic, never the whole plane by the preceding lemma.
We first prove that this generic plane cubic is smooth.

Suppose $j\ne\pm g$. At $R=P$ the first value is zero and the
second is a nonzero functional $l_j$. Since $j\ne\pm g$,
$X_j=x(j_1)+a\ne0$. Put
\[
 r=\ell(j)/X_j,\qquad p=P_1(j)/X_j,\qquad Z_j=x(j_2)+c.
\]
The equation $l_j=0$ is
$v=rT+\ell(j)h+pF+Z_jU$.
In the two $Y=1$ coordinates at $P$, approach $P$ along
$(t_1,t_2)=(\kappa t^3,t)$. The first value, divided by $t^3$,
tends to $F+\kappa h$. Hence the limiting coefficient plane is
defined by $l_j=F+\kappa h=0$. Substitution into $q$ shows that
its Frobenius-form coefficient matrix, in variables $(T,U,h)$,
has determinant
\begin{equation}\label{eq:ss-pair-plane-determinant}
                        (\kappa+r^2)(\kappa^2+r).
\end{equation}
Choose $\kappa$ outside the finite zero set of this expression.
The limiting plane cubic is smooth, so the geometric generic
plane cubic is smooth and integral.

For $j=g$, instead approach $P$ along $(t,0)$. The two values,
divided by $t$, tend to $h,v$. The limiting plane cubic is
\[
                           U^3+F^2T+FT^2=0,
\]
again smooth. This proves the generic integrality in this case too.

Consequently each $I_j$ has exactly one three-dimensional component
dominating $X$, which is generically reduced. All rank-one strata
are finite over $X$ and have two-dimensional coefficient fibres.
For $j\ne\pm g$ there is no rank-zero stratum, hence no other
component can dominate $\mathcal Q$. For $j=g$ the only additional
three-dimensional support is the constant section $R=P$.

For $j\ne\pm g$ we exhibit an etale point of $I_j\to\mathcal Q$.
Choose $s$ in the integral cubic surface $\mathcal Q\cap(l_j=0)$
such that $h\ne0$ and $\partial_2s(j)\ne0$.
Both are genuine nonzero linear conditions on $\ker l_j$:
$l_j$ has nonzero $v$ coefficient $X_j$, whereas
\[
 \partial_2s(j)=(T+hX_j)\bigl(x(j_2)^2+c^2\bigr)+UX_j
\]
has no $v$ coefficient and has nonzero $U$ coefficient.
The Jacobian of the two value equations with respect to $R$ at
$R=P$ has determinant $h\partial_2s(j)\ne0$.
This point therefore belongs to the unique main component, whose
projection is dominant and generically etale. After removing a
proper closed subset of $\mathcal Q$, its entire finite cover is
etale; all lower-dimensional components have also disappeared.

For $j=g$ and $3g\ne0$, the same argument at $R=Q$ applies to
the other value at $2g$. Here $2g\ne\pm g$ and
$x(2g_2)+c\ne0$. The integral surface
$\mathcal Q\cap(s(2g)=0)$ is not contained in $v=0$ or in
$\partial_2s(2g)=0$. The Jacobian determinant is
$v\partial_2s(2g)$, so the main component is again generically
etale. The constant intersection at $P$ has length two by
Lemma~\ref{lem:ss-uniform-local-contact}; its leading coefficient
is $T(c^4+c)$.

If $3g=0$, the constant intersection has length exactly three.
Indeed the branch through $P$ has first coordinate
$(F/h)t^3+O(t^4)$, and translation on the second elliptic factor
gives $\ell(g_2+T(t))=t^3+O(t^4)$. The other value on that branch
therefore has leading coefficient
\begin{equation}\label{eq:ss-flex-contact-length}
                             T+vF/h.
\end{equation}
It is generically nonzero on the integral cubic $\mathcal Q$.
Every component of $I_g$ has dimension at most three, so its
generic fibre over $\mathcal Q$ is finite: a positive-dimensional
fibre over the generic point would require a component of dimension
at least four. Intersection theory then gives its total length
$L^2=18$. Subtracting the length-three constant component leaves
degree $15$ on the unique integral main component. Its inseparable
degree is a power of two dividing $15$, so it is separable.
After shrinking $\mathcal Q$ this moving cover is etale.
The same dimension argument and intersection number apply in the
preceding cases. This proves every assertion.
\end{proof}

\begin{lemma}[There are generically no three kernel translates]
\label{lem:ss-contact-no-triples}
\status{Proved under stated hypotheses}
A geometric generic section of $\mathcal Q$ contains no three
distinct points of a $K$-orbit.
\end{lemma}
\begin{proof}
For distinct nonzero $j,k\in K$, consider the three-value incidence
over $X$. We first show that its generic three value functionals
are independent and that $q$ does not vanish identically on their
common vector kernel of dimension two.

One may reorder and translate the unordered triple. Unless it is
a coset of $\{0,g,-g\}$ with $3g=0$, choose its initial point so
that neither $j$ nor $k$ equals $g$. This is possible because a
three-element set in which every point has its $g$-successor in the
set is a three-element orbit of translation by $g$.
At $R=P$, the other two value rows $l_j,l_k$ are independent modulo
$F,h$. For $j\ne\pm g$ their rows in $(T,v,U)$ are
\[
                        (\ell(j),X_j,X_jZ_j).
\]
Different second-factor $x$-coordinates give different normalized
last entries. Equal second-factor $x$-coordinates imply $k=-j$,
by injectivity of the second projection of $K$; then $X_k=X_j$
and $\ell(k)=\ell(j)+1$, so the rows are still independent.
The exceptional allowed index $j=-g$ has row $(1,0,0)$, also
independent from every other permitted row. Thus
$l_j,l_k,F,h$ have rank four.

Approach $P$ along $(\kappa t^3,t)$ as before. The limiting three
rows are $F+\kappa h,l_j,l_k$, of rank three. If the cubic vanished
on the generic common kernel, it would vanish on every limiting
kernel
\[
                    \ker(l_j,l_k,F+\kappa h).
\]
As $\kappa$ varies, these two-dimensional vector spaces have dense
union in the three-dimensional space $\ker(l_j,l_k)$, because
$F,h$ are independent on that space. This would put a projective
plane inside $\mathcal Q$, contrary to
Lemma~\ref{lem:ss-contact-cubic-evaluation}.

For the exceptional triple take $(j,k)=(g,-g)$ and $3g=0$.
Then $a^4+a=c^4+c=0$. Along the same path, after division of the
first two rows by $t^3$, the limiting rows are
\[
                      F+\kappa h,\quad T+\kappa v,\quad F+T.
\]
The coefficient one in the middle row is the actual elliptic
translation coefficient
$\ell(g_2+T(t))=t^3+O(t^4)$.
For $\kappa\ne0$ the three rows have rank three, and their kernel
has $F=T=\kappa h$, $v=h$. The cubic restricts to $U^3$, which
is not zero. This proves the required generic assertion in all cases.

The incidence dimension now proves the lemma. On the generic
rank-three stratum the cubic cuts the coefficient line in finitely
many points. Where it contains that line, the base has dimension at
most one, so this stratum still has total dimension at most two.
The rank-two locus has dimension at most one, and its coefficient
plane is never contained in $\mathcal Q$, again giving total
dimension at most two. The rank-one locus is finite by the pair
calculation, with coefficient fibres of dimension two. There is
no rank-zero triple. Thus the full proper three-value incidence
has dimension at most two. Its image cannot fill the
three-dimensional cubic $\mathcal Q$. There are only finitely many
ordered pairs $j,k$, so they can all be excluded simultaneously.
\end{proof}

\begin{proposition}[A mixed contact with no other collisions for every graph kernel]
\label{prop:ss-uniform-graph-mixed-contact}
\status{Proved under stated hypotheses}
For every finite odd subgroup $K\subset E^2(k)$ with injective
coordinate projections, and every $0\ne g\in K$, there is a smooth
source $C\in|L|$ on the source cubic $Q_H=0$ and outside $\det S=0$
whose image in $X/K$ has exactly one simple self-tacnode, with
preimages differing by $g$, and otherwise only ordinary nodes.
The extra source Cartier-kernel differential has unequal values on
the two identified tangent vectors. In particular the local
mixed-contact criterion applies to this norm image.
\end{proposition}
\begin{proof}
We first record the irreducibility needed to combine separate open
conditions. Let $\mathcal J_g$ parametrize triples $(R,\tau,s)$,
where $\tau$ is a common translated tangent direction at $R,R+g$,
and the section vanishes at both points and has derivative zero
along $\tau$ at both points. The base $(R,\tau)$ has dimension
three. Off finitely many $R$, the six first jets at $R,R+g$ are
independent: the tensor tangent-space calculation reduces a failure
to the crossed conditions
$3R_i+g_i=0$ or $3R_i+2g_i=0$ on both elliptic factors.
Thus on this open base the four contact conditions are independent,
and its coefficient spaces form a rank-five vector bundle.
The full source cubic $Q_H$ is a nondegenerate Frobenius form in
nine variables. It contains no vector space of dimension five,
by the same polarization argument as above. Consequently every
coefficient fibre is a genuine cubic threefold.
At $(R,\tau)=(P,\text{second-factor tangent})$ the fibre is the
smooth cubic $\mathcal Q$. The six jets here are indeed independent:
the origin rows select coefficients $S_{22},S_{12},S_{21}$, and
the three translated rows have a unit triangular minor on
$S_{00},S_{10},S_{01}$.
The relative cubic is nonzero in every fibre, so it is a flat
relative hypersurface in the projective bundle. Its constant fibre
dimension three excludes vertical components. The smooth integral
fibre just exhibited therefore gives a geometrically integral
generic fibre, and the total family on the good base is integral.
There is a unique six-dimensional component of
$\mathcal J_g\cap(Q_H=0)$.

There is no further component of that dimension over the finitely
many exceptional $R$. Indeed the two values and the derivative in
any nonzero tangent direction at the first point have rank at
least three. To see this, the Segre point representing $R+g$
does not lie in the embedded tangent space at $R$: both its factor
coordinates differ from those of $R$, and its image in the tensor
product of the two one-dimensional-point quotients is nonzero.
Thus the contact rank is at least three. The coefficient vector
space then has dimension at most six; intersecting with $Q_H$
reduces its projective dimension to at most four. Together with
the one-dimensional choice of tangent, this gives dimension at
most five. We work on the unique main component from now on.

The preceding two lemmas give, on its particular vertical slice,
a section with no triple and with no unwanted nontransverse
identification. All pair intersections are finite. These conditions
give a nonempty open on the main contact component. The only
openness point needing care is $3g=0$, when the prescribed
intersection in the slice has length three, not two. Near a chosen
point the proper intersection family is finite flat of degree
$18$: the two Cartier equations form a regular sequence in every
fibre of the smooth relative surface, so their zero-dimensional
complete intersection is flat over the base. Its length-three cluster at the marked point
contains the prescribed intersection of length at least two.
A further tangency in that cluster would contribute another
length at least two, requiring total length at least four.
It is therefore impossible. The finite-cluster decomposition is
made etale locally on the coefficient base; the resulting open
condition descends to the base. The other clusters are etale and
remain so after shrinking. This proves the stated nonempty open also in
the flex case. The triple-avoidance conditions are open directly
by their proper incidences. For $j\ne\pm g$, transversality is
likewise an open condition by the proper two-value incidence.

There is also a nonempty open of this same main component on
which the source is globally smooth, the prescribed contact is
simple, $\det S\ne0$, and the two Cartier values are unequal.
If $3g\ne0$, these assertions follow from
Lemmas~\ref{lem:ss-uniform-local-contact} and
\ref{lem:ss-uniform-source-smoothness}; injectivity of the two
projections makes both coordinates have the same non-three order.
If $3g=0$, product automorphisms of $E$ preserving $O$ take the
already certified contact of Lemma~\ref{lem:ss-global-contact03}
to the prescribed nonzero pair of three-torsion coordinates.
These automorphisms preserve the linear system and the two source
Hasse factors. All asserted conditions are open on the main
incidence. At the actual two points of contact03 the six jets have
rank six, as independently certified by
\path{agents/monodromy/code/verify_ss_collision_main_incidence.py},
with the six rows and their row reduction retained in its output.
Thus this witness, as well as the entire vertical slice at $P$,
belongs to the unique main incidence.

The two nonempty opens intersect by irreducibility. For a section
in their intersection, all translates meet in finitely many
points, so the map of its smooth source to $X/K$ is birational
onto its image. The absence of triples and the transverse
intersection conclusions give precisely ordinary nodes away from
the one simple tacnode. The unequal-value assertion is exactly
the hypothesis of the already proved local mixed-contact Cartier
criterion. This completes the proof.
\end{proof}

\begin{remark}[Scope and exact algebra certificate]\label{auto:A_monodromy:c25cfd1730a5}
\status{Proved under stated hypotheses}

The only group restriction above is injectivity of both coordinate
projections. In particular it includes balanced odd quaternionic
isogenies with kernel $\{(-bP,P):P\in E[d]\}$ and
$b^\dagger b=d-1$. The exact program
\path{agents/monodromy/code/ss_uniform_collision_identities.py}
checks the homogeneous coefficient matrix, its determinant, the
plane-cubic determinant, the two flex limiting equations and the
order-three elliptic translation coefficient. It does not replace
the incidence, flatness or openness arguments. The subsequent uniform
tangent-triple theorem merges every possible triple-collision component
of these norm families; its odd-degree full-group consequence is
Theorem~\ref{thm:ss-all-odd-full-groups}. The even-degree conclusion is proved in
Theorem~\ref{thm:ss-product-all-full-groups}.
\end{remark}


\subsection{Smooth sources after imposing a third kernel value}
\label{sec:DE-A4-contact-third-value}

\begin{lemma}[The moving discriminant spans the contact coefficient space]
\label{lem:DE-A4-contact-discriminant-span}
\status{Proved; root and monodromy full independent proof/code audits passed.}
Keep the homogeneous contact space $\Lambda=\mathbf P^4_{[F:T:U:v:h]}$
of \eqref{eq:ss-contact-linear-system}, for any two nonzero points
$g_1=(a,b)$ and $g_2=(c,d)$ of $E:y^2+y=x^3$.
The second point may be a flex.
There is an irreducible hypersurface
$\Delta_{\rm mov}\subset\Lambda$ such that the locus of singular
source curves is contained in
\[
              (FThv=0)\ \cup\ \Delta_{\rm mov}.
\]
The hypersurface $\Delta_{\rm mov}$ is not contained in any
projective hyperplane.  In particular, every hyperplane of
$\Lambda$ distinct from $F=0,T=0,h=0,v=0$ has a nonempty open
parametrizing globally smooth source curves.
\end{lemma}
\begin{proof}
We first check that the source-singularity support calculation
extends to the flex case; none of its multiplicity calculations
will be needed.  Put
\[
 X=x+a,\quad Y=y+b,\quad Z=z+c,\quad
 \ell=w+c^2z+d^2,\quad A_1=Y+Xx^2.
\]
The five functions of the contact system are
$X,Y,\ell,Xz,X\ell$.  Its base points are
$P=(O,O)$, $Q=(g_1,g_2)$ and $(g_1,-2g_2)$, with the last two
coinciding when $g_2$ is a flex.
On the affine open $X\ne0$, the ratios
$Y/X,\ell/X,z,\ell$ detect both tangent directions except where
$A_1=\ell=0$.  This is a finite set: $A_1$ is a nonconstant
function of the first elliptic point and $\ell$ of the second.
At those exceptions the value and the derivative of $z$ still
give first-jet rank at least two.

On the fibre through $g_1$ the section is $T\ell$.
Away from its base points, a singularity therefore lies in $T=0$.
At $Q$ the only first derivative is $v$.  If the additional base
point $(g_1,-2g_2)$ is distinct, its two first derivatives impose
independent conditions; if it coincides with $Q$, it is already
accounted for by $v=0$.  On the first-factor fibre through $-g_1$
the value and first derivatives have rank three except at finitely
many second points, and rank at least two at those exceptions.
On the first-factor fibre through $O$, the value is $F$, except
at $P$, where the first derivative is $h$.  On the second-factor
fibre through $O$, the value and first derivative are $T+hX,h$;
the full first-jet rank is three except for the two finite points
with $X=0$, where it is at least two.
These are exactly the four-chart calculations in the proof of
Lemma~\ref{lem:ss-uniform-source-smoothness}; they remain valid
when the two affine base points merge.
Every finite rank-two exceptional set contributes a singularity
incidence of dimension at most two, and hence no additional divisor
in the four-dimensional coefficient space.

The remaining singularity incidence is a $\mathbf P^1$-bundle over
the irreducible rank-three jet open of the surface.  Its image
closure is irreducible.  It is a hypersurface because the following
explicit singular sections have an ordinary node at the selected
moving point.  On the dense open $XA_1\ell Z\ne0$, solve the
three singularity equations with free parameters $T,h$:
\[
 \begin{aligned}
 F&=\frac{T\ell}{A_1},\\
 U&=Z^2\left(\frac TX+h\right),\\
 v&=\frac{TZ^3}{X}
       +h\bigl(z^2Z+w+d\bigr)+FX^2.
 \end{aligned}
\]
These are the earlier moving-node formulas, after replacing
$E_0$ by $U+hc^2$.  The mixed Hessian is $TZ^2/X$ and is nonzero
when $T\ne0$.  The two critical-point equations then determine
the critical point formally, and its value is a nonzero first-order
linear condition on the coefficients.  Thus the image has codimension
one.  Call it $\Delta_{\rm mov}$.  The remaining incidence strata
have smaller images, and their closures are already contained in
the source discriminant.  The same dimension bounds show that a
generic member of $\Lambda$ is smooth.  Thus the full-system
discriminant hypersurface of
Proposition~\ref{prop:ss-level3-three-witness-reduction}
restricts to a nonzero equation on $\Lambda$.
Since its zero locus is a hypersurface,
the preceding divisorial-support calculation gives the asserted
containment.  Equivalently, any additional component of the
discriminant would have a divisorial singularity incidence, which
the chart calculation has exhausted.

Suppose a linear form
$\alpha F+\beta T+\gamma U+\delta v+\epsilon h$ vanished on
$\Delta_{\rm mov}$.  Substitute the displayed moving-node formulas.
Since $T,h$ vary freely, its coefficient of $h$ is identically zero:
\[
                  \gamma Z^2+\delta(z^2Z+w+d)+\epsilon=0.
\]
On the second elliptic curve the three functions
$z^2Z+w+d$, $Z^2$, and $1$ have respective pole orders $6,4,0$
at $O$.  Their leading terms cannot cancel, so
$\delta=\gamma=\epsilon=0$.
The coefficient of $T$ now gives
$\alpha\ell/A_1+\beta=0$ on a dense open of the product.
Fix a second point with $\ell\ne0$ and vary the first.
The function
$A_1=y^2+b+ax^2$ has pole order six and is nonconstant.
It follows that $\alpha=\beta=0$.
Thus no nonzero linear form vanishes on $\Delta_{\rm mov}$.

Finally a projective hyperplane contained in the singular-source
locus would, by irreducibility and the displayed support, equal one
of the four forced hyperplanes or equal $\Delta_{\rm mov}$.
The latter is impossible by the span assertion.  The complement
of the singular-source locus in every other hyperplane is therefore
a nonempty open.
\end{proof}

\begin{remark}[Exact symbolic checks]
\label{rem:DE-A4-third-value-checks}
\status{Proved under stated hypotheses}
The retained program
\path{agents/obstructions_crystals/code/contact_third_value_checks.py}
checks the universal moving-node equation, both first derivatives,
the nonzero mixed-Hessian formula, the coefficient of $h$ in an
arbitrary linear relation, the alternate formula for $A_1$,
$\det S=FTU$, and $l_{-g}=F+T$.
Its author run passed in $0.12054040400107624$ seconds; the complete
output is retained alongside the program.
Root and monodromy independently audited the entire proof and program,
with author-output-preserving reruns in $0.15617483800451737$ and
$0.1453385850036284$ seconds, respectively.
The pole-order independence and geometric incidence arguments are
proved in the text and are not inferred from the computation.
\end{remark}

\begin{corollary}[Every third-value hyperplane has smooth sharp sources]
\label{cor:DE-A4-third-value-smooth}
\status{Proved; root and monodromy full independent proof/code audits passed.}
Let $K\subset E^2(k)$ be a finite odd subgroup with both coordinate
projections injective, fix $0\ne g\in K$, and construct $\Lambda$
using $g_1,g_2$ as above.
For every $k\in K\setminus\{0,g\}$, the hyperplane
\[
                         H_k=\{s(k)=0\}\subset\Lambda
\]
has a nonempty open of globally smooth source curves with
$\det S\ne0$ and $Q_H(S)\ne0$.
Every such curve contains the three distinct points $P,Q,k$.
The assertion concerns smoothness and the source Hasse factors;
it does not claim that the three tangent directions have the
required tangent-triple type or that other norm collisions are absent.
\end{corollary}
\begin{proof}
In the contact coordinates the value functional at an affine point is
\[
      l_k=P_1(k)F+\ell(k)T+X(k)Z(k)U
                     +X(k)v+X(k)\ell(k)h,
      \qquad P_1=Y+a^2X.
\]
If $k=-g$, this is $F+T$.
Otherwise injectivity of both coordinate projections gives
$X(k)\ne0$ and $Z(k)\ne0$: equality of either corresponding
$x$-coordinate would force $k=g$ or $k=-g$.
Thus both the $U$ and $v$ coefficients of $l_k$ are nonzero.
In either case $H_k$ is distinct from the four forced hyperplanes
in Lemma~\ref{lem:DE-A4-contact-discriminant-span}, and from $U=0$.
That lemma gives a dense open of smooth sources in $H_k$.

The homogeneous determinant on the contact slice is
$\det S=FTU$, so its nonvanishing is another nonempty open on $H_k$.
The restricted source cubic $q=Q_H|_\Lambda$ is a nondegenerate
Frobenius cubic in five variables.  Its coefficient matrix has
determinant one, as in
Lemma~\ref{lem:ss-contact-cubic-evaluation}.
Polarization of this cubic excludes a projective plane in it,
and in particular excludes the hyperplane $H_k$.
Hence $q\ne0$ is also a nonempty open on $H_k$.
The irreducibility of $H_k\simeq\mathbf P^3$ lets these three opens
intersect.  The stated point containments are the defining contact
and third-value equations.
\end{proof}


\subsection{The simple tangent-triple type for every translated triple}
\label{sec:DE-A4-uniform-tangent-triple-type}

\begin{lemma}[A uniform local type on the main triple-contact incidence]
\label{lem:DE-A4-uniform-triple-local-type}
\status{Proved. Complete independent proof and program audits by root
and monodromy passed.}
Let $E:y^2+y=x^3$ over an algebraically closed field of
characteristic two, $X=E^2$, and
$L=\OO_E(3O)\boxtimes\OO_E(3O)$.
Let $K\subset X(k)$ be a finite odd subgroup with injective
coordinate projections.  Fix $0\ne g\in K$ and
$\eta\in K\setminus\{0,g\}$.
Consider sections containing $R,R+g,R+\eta$, with a common translated
tangent at the first two points.

The locus where these five linear conditions are independent is
an irreducible nonempty open in the base of marked points and
tangent directions.  Over that open the coefficient incidence is
a $\mathbf P^3$-bundle.  It contains a nonempty open of sections
whose source curve is globally smooth, satisfies
$\det S\ne0$ and $Q_H(S)\ne0$, has simple contact of order exactly
two at the first two translated branches, and has a transverse
third branch.
There is no exception when $g$ has order three or when
$\{0,g,\eta\}$ is a cyclic three-element subgroup.
\end{lemma}
\begin{proof}
Write $h_0=0,h_1=g,h_2=\eta$, and let
\[
                    p_i=R_1+(h_i)_1,\qquad
                    q_i=R_2+(h_i)_2
                    \quad(0\le i\le2).
\]
All three $p_i$ are distinct, as are all three $q_i$.
Choose $R$ in a dense open where each triple of value vectors
in the plane-cubic embedding is a basis.
Such a choice is possible: three points on this cubic are
collinear exactly when their sum is zero, so the failures are
the finite conditions $3R_j+g_j+\eta_j=0$ on each elliptic factor.
Normalize the value vectors to $e_0,e_1,e_2$ separately in the two
factors.  A section is then represented by a matrix $S$.
Its containment of the three points says
\[
                             S_{00}=S_{11}=S_{22}=0.
\]

Use the invariant formal-group parameters on both factors.
They identify the tangent directions at all translated points.
Write the first-factor tangent vectors in the normalized frames as
$a_i=(a_{i0},a_{i1},a_{i2})$, and the second-factor tangent vectors
as $b_i=(b_{i0},b_{i1},b_{i2})$.
We may also require
\[
             a_{02}a_{12}b_{02}b_{12}b_{20}b_{21}\ne0,
             \qquad q_0,q_1\notin E[3].
\]
Every excluded condition is finite in the corresponding $R_j$.
For example, $a_{02}=0$ says that the tangent at $p_0$ contains
$p_1$, equivalently $3R_1+g_1=0$.
The condition $a_{12}=0$ is $3R_1+2g_1=0$.
The analogous conditions for $b_{02},b_{12}$ use the second
factor, and those for $b_{20},b_{21}$ are
$3R_2+2\eta_2+g_2=0$ and $3R_2+2\eta_2=0$.
The nonflex conditions are again finite.  Multiplication by three
is etale and finite in characteristic two, so these exclusions
cannot exhaust either elliptic curve.

Take the common tangent direction to be the second-factor direction.
After the three value conditions, the two tangent conditions
are independent equations on rows zero and one.
All sections satisfying these five conditions therefore have the
form
\[
 S=
 \begin{pmatrix}
       0&\alpha b_{02}&\alpha b_{01}\\
       \beta b_{12}&0&\beta b_{10}\\
       r&s&0
 \end{pmatrix}.
\]
Here $\alpha,\beta,r,s$ are free.  Fix $\alpha\beta\ne0$.
The row-zero and row-one sections on the second cubic are
nonzero multiples of its tangent-line sections at $q_0,q_1$.
Because these points are not flexes, their local expansions have
nonzero quadratic coefficients $\alpha\kappa_0$ and
$\beta\kappa_1$, with $\kappa_0\kappa_1\ne0$.

The first-factor derivatives of the full source equation at the
first two points are
\[
        l_0=a_{01}\beta b_{12}+a_{02}r,\qquad
        l_1=a_{10}\alpha b_{02}+a_{12}s.
\]
They are independent affine functions of $r,s$.
Choose them nonzero, so the source is smooth at both points and
has the indicated common tangent.
In the common invariant formal parameters $(u,t)$, the two
branches then have expansions
\[
       u=\frac{\alpha\kappa_0}{l_0}t^2+O(t^3),
       \qquad
       u=\frac{\beta\kappa_1}{l_1}t^2+O(t^3).
\]
The displayed parameters are identified by translation, so this is
the actual contact comparison.  It has order exactly two provided
\[
                  \alpha\kappa_0l_1+\beta\kappa_1l_0\ne0.
\]
This excludes one genuine affine linear equation in $r,s$:
its $r$ and $s$ coefficients are nonzero.
At the third point the derivative in the second-factor direction is
\[
                               rb_{20}+sb_{21}.
\]
Its nonvanishing excludes one more genuine linear equation and
ensures both smoothness at the third point and transversality to
the common tangent.  Since $k$ is infinite, all these finitely
many proper affine lines in the $(r,s)$-plane can be avoided.
This proves the existence of the required local type with five
independent linear conditions for every fixed $g,\eta$.

It remains to combine this local type with global smoothness and
the two nonvanishing source Hasse factors.  Let
$B=X\times\mathbf P^1$ parametrize $(R,\tau)$, where $\tau$ is
a common translated tangent direction.  The two value conditions
at $R,R+g$, their two directional-derivative conditions, and the
value condition at $R+\eta$ are algebraic linear conditions on
$H^0(X,L)$.  Their rank-five locus $\Omega\subset B$ is nonempty
by the construction above.  It is open and irreducible, since
$B$ is irreducible.  Their vector kernels form a rank-four bundle
on $\Omega$, whose projectivization $\mathcal T_{g,\eta}$ is therefore
an irreducible $\mathbf P^3$-bundle.

The point $(R,\tau)=((O,O),\text{second-factor tangent})$ also
belongs to $\Omega$: the first four conditions define the contact
$\mathbf P^4$ considered earlier, and the third-value functional
$l_\eta$ is nonzero there by
Corollary~\ref{cor:DE-A4-third-value-smooth}.
That corollary gives a point of this fibre whose source is globally
smooth and satisfies $\det S\,Q_H(S)\ne0$.
These conditions define a nonempty open of $\mathcal T_{g,\eta}$.
The local type constructed above likewise defines a nonempty open:
it is specified by nonvanishing first derivatives, a nonzero
contact-order-two coefficient, and a nonzero third derivative
in the marked tangent direction.
The two opens meet because $\mathcal T_{g,\eta}$ is irreducible.
Their intersection proves the lemma in full.
\end{proof}

\begin{remark}[Scope of the local-type construction]
\label{rem:DE-A4-uniform-triple-local-scope}
\status{Proved under stated hypotheses}
The lemma supplies a globally smooth sharp source with the required
three marked branches on the main rank-five incidence.
Exclusion of other identifications under the full quotient by $K$
is a further geometric condition.  It can be combined with this
lemma once it gives a nonempty open on the same incidence.
The proof makes no assertion that a finite list of small-kernel
witnesses controls arbitrary kernels.

The retained universal program
\path{agents/obstructions_crystals/code/uniform_triple_local_type_checks.py}
checks the rank-five minor, both tangent equations, the independent
first gradients, both nonzero coefficients of the simple-contact
condition, the third vertical derivative, and the standard cubic
tangent-line coefficient $x^4+x$.
Its author run passed in $0.0512757739998051$ seconds; its complete
output is retained alongside the source.
Independent complete program reads and author-output-preserving reruns
passed in $0.04299991099833278$ seconds (root) and
$0.05460318699624622$ seconds (monodromy).
The finite-exception and irreducible-incidence arguments are proved
above and are not inferred from this algebraic check.
\end{remark}


% /root/monodromy, 2026-11-19. Complete root, reduction-descent and obstructions-crystals audits passed.
\subsection{Uniform tangent triples for odd graph kernels}
\label{subsec:ss-uniform-tangent-triples}
\status{Proved and independently checked. The parent, reduction/descent and
obstructions/crystals agents independently audited the complete geometric
proof and both dependency files. The actual rank-five witness certificate
was independently read and rerun. This proves uniform tangent-triple
witnesses for all odd graph kernels; its full-group consequence is stated
separately.}

Continue with $E:y^2+y=x^3$, $X=E^2$, $L=\mathcal O(3O)\boxtimes
\mathcal O(3O)$ and a finite odd subgroup $K\subset X(k)$ whose
two coordinate projections are injective. Given a triple of distinct
elements, translate and reorder it as $\{0,g,\eta\}$.
We may and do arrange
\[
 \eta\ne-g,\qquad\hbox{or}\qquad 3g=0\ \hbox{and}\ \eta=-g.
\]
Indeed if the original triple is $\{0,g,-g\}$ and $3g\ne0$, use
the translated triple $\{0,2g,g\}$. Take the contact space
$\Lambda$ defined by $g$, and its third-value hyperplane
$H_\eta=(l_\eta=0)\simeq\mathbf P^3$, where $l_j(s)=s(j)$.
The coordinates $F,T,U,v,h$, functions $X_1,Y_1,Z_1,\ell,P_1$,
and map $\Phi$ are those of
Lemma~\ref{lem:ss-contact-cubic-evaluation}.

\begin{lemma}[The projected value map and its three-value ranks]
\label{lem:ss-third-value-projection}
\status{Proved under stated hypotheses}
If $\eta\ne-g$, the base locus of $H_\eta$ consists of
$P=(O,O),Q=g,\eta$, and the possible additional point
$(g_1,-2g_2)$ already in the contact base locus. The additional
point has no $K$-difference with the other three.
For every $0\ne j\in K$, the two values at $R,R+j$ on $H_\eta$
have rank one only at finitely many $R$, and generically have rank
two. For every three distinct translates their three value
functionals are generically independent.

If $3g=0$ and $\eta=-g$, the base locus consists of $P,Q,-g$.
The pair-rank assertion still holds for $j\notin\{0,g,-g\}$.
For $j=\pm g$, any additional positive-dimensional rank-one
stratum contributes only sections with $F=0$; it can be discarded
on $F\ne0$. The generic three-value independence holds in this
case as well.
\end{lemma}
\begin{proof}
The new value map is projection of $\Phi$ from the point
$\Phi(\eta)$. Away from the exceptional fibres in
Lemma~\ref{lem:ss-contact-cubic-evaluation}, $\Phi$ is injective.
For $\eta\ne\pm g$, its only possible exceptional horizontal
fibre is $\eta=-2g$, since $\ell(\eta)=0$ exactly there.
On that fibre the map is the degree-two pencil
\[
                         \rho(x,y)=\frac{y+b}{x+a}.
\]
The fibre through $-2g_1$ has other point $g_1$, which is already
the base point $(g_1,-2g_2)$. Hence no further base point appears.
If $\eta=-g$ and $3g=0$, the value is $F+T$; its fibre on the
vertical exceptional line is given by $\ell=1$, whose support
is just $-g_2$ because this is a flex. These observations give
the stated base loci. The assertion about $K$-differences follows
from injectivity of the first projection, as before.

We make the pair-rank assertion precise. Where the first coordinates
of $R,R+j$ avoid $O,\pm g_1$, divide the five value coordinates
by $X_1$. They become
\begin{equation}\label{eq:ss-projected-value-coordinates}
              (b_0,\rho b_0,b_2/X_1,b_1,b_2),
\end{equation}
where $(b_0,b_1,b_2)$ is the homogeneous value vector of the
second elliptic point in the basis $(1,z,\ell)$.
Three value rows at $R,R+j,\eta$ can be dependent only if their
projections to $(b_0,b_1,b_2)$ are dependent. If
$\eta\ne-g$, its projection is nonzero. The three second-factor
points must therefore be collinear, unless two coincide. The
noncoincident case gives
\[
                            2R_2+j_2+\eta_2=0,
\]
so there are only finitely many possibilities for $R_2$.
Coincidences also give only finitely many possibilities.

Fix one of these second-factor values. In a dependence relation,
the second coordinate in
\eqref{eq:ss-projected-value-coordinates} would give a
nontrivial relation
\[
                  A\rho(R_1+j_1)+B\rho(R_1)+C=0.
\]
There is no such identity in $R_1$. The pole divisor of $\rho$
is exactly $(O)+(-g_1)$, with two simple poles. If both $A,B$
are nonzero, an identity would make this divisor invariant under
translation by $j_1$. Its stabilizer has order dividing two,
whereas $j_1$ has nontrivial odd order. If one coefficient is zero,
the nonconstant remaining function gives the contradiction.
The case in which a second-factor value is $O$ uses the same
homogeneous formula: the other moving value has $b_0\ne0$, and
its coefficient cannot be zero unless the fixed second point
were $O$, which it is not. Thus the relation remains nontrivial.
We conclude that the rank-one points on these finitely many
horizontal fibres are finite.

The omitted first-factor exceptional fibres cause no new curve
when $\eta\ne-g$. Their projections to $(b_0,b_1,b_2)$ are zero.
If only one moving point lies on such a fibre, dependence forces
the other second point to equal $\eta_2$, giving finitely many
points. If both do, dependence reduces to equality of their
original $\Phi$-values, already proved to occur at only finitely
many points. For $\eta=-g$, the fixed projection is zero.
On the nonexceptional open dependence would require the two
distinct second-factor points to have proportional value vectors,
which is impossible. Two exceptional vertical fibres can give
a curve only when $j=\pm g$ (since $3g=0$). Their values lie on
the line spanned by $F,T$; on $H_{-g}$ one has $T=F$, so a
generic point of every such incidence requires $F=0$.
All remaining exceptional points are finite. This proves the
pair assertions, including their stated restriction in the cyclic case.

We finally prove generic independence for three values on
$H_\eta$. If they were dependent, the fixed value
$\Phi(\eta)$ would lie in the span of the three original values
at $R,R+j,R+k$ for a geometric generic $R$.
The second-factor value vectors of these three points are generically
a basis: dependence would impose $3R_2+j_2+k_2=0$.
If $\eta=-g$, its projection is zero, forcing all three coefficients
to vanish, which is impossible.
Otherwise normalize the first coordinate of each second-factor
value to one. The coefficients $\alpha_i\in k(R_2)$ expressing
the fixed second-factor point have sum one. Their other value
equation would force
\begin{equation}\label{eq:ss-three-pole-relation}
       \alpha_0\rho(R_1)+\alpha_1\rho(R_1+j_1)
             +\alpha_2\rho(R_1+k_1)=\rho(\eta_1).
\end{equation}
The poles of these three functions are the three pairs
$\{-j_i,-g_1-j_i\}$, with $j_i=0,j_1,k_1$.
Every nonzero coefficient must have all its poles cancelled.
Viewing the pairs as edges in orbits of translation by $g_1$,
a set of at most three distinct edges can have no endpoint of
degree one only if it is an entire three-cycle. Thus the only
possible case is $3g=0$ and $\{0,j,k\}=\{0,g,-g\}$.
The two residues of $\rho$ against an invariant differential are
equal and nonzero (by the residue theorem in characteristic two).
Cancellation in this cycle therefore forces all $\alpha_i$ equal,
and their sum one forces each to be one. But the sum of the three
second-factor $x$-coordinates is not constant: it has a pole at
each of the three distinct translated origins, with no cancellation.
It cannot equal the fixed $x(\eta_2)$ for generic $R_2$.
This contradicts \eqref{eq:ss-three-pole-relation}, and
completes the proof.
\end{proof}

\begin{lemma}[Two linear independence tests]
\label{lem:ss-third-value-linear-tests}
\status{Proved under stated hypotheses}
Let $i,\eta\in K\setminus\{0,g\}$ be distinct, with
$i\ne-g$. On $\ker(l_i,l_\eta)\subset\Lambda$, neither $h$ nor
$\partial_2s(i)$ is identically zero.
If additionally $i=2g$, $3g\ne0$ and $\eta\ne\pm g,2g$,
then $v$ is not identically zero there either.
\end{lemma}
\begin{proof}
For a point $j\ne\pm g$ write
\[
 X_j=x(j_1)+a,\quad Z_j=x(j_2)+c,\quad
 p_j=\rho(j_1)+a^2,\quad r_j=\ell(j)/X_j.
\]
The normalized value row is
\[
 l_j/X_j=p_jF+r_jT+Z_jU+v+\ell(j)h,
\]
whereas
\[
 \partial_2s(j)/X_j=(Z_j^2/X_j)T+U+Z_j^2h.
\]
The rows of $l_i,l_\eta$ modulo $F,h$ are independent, by the
same argument with $(\ell(j),X_j,X_jZ_j)$ used in
Lemma~\ref{lem:ss-contact-no-triples}. This proves the assertion
for $h$. If $\eta=-g$, then $l_\eta=F+T$, and independence of
the derivative follows at once from its nonzero $U$ coefficient
and the nonzero $v$ coefficient of $l_i$.

Assume $\eta\ne\pm g$. A relation expressing the derivative in
the span of the two value rows must cancel their $v$ coefficients.
Its $F$ coefficient then requires $p_i=p_\eta$.
The degree-two pencil $\rho$ has fibre divisors summing to $-g_1$;
as $i\ne\eta$, this would imply $i+\eta=-g$ by injectivity of
the first projection. Hence the three distinct second-factor points
$i_2,\eta_2,g_2$ lie on one affine line, of slope $\beta$.
Put $L_0=\beta+c^2$. Then
\[
 \ell(i)=L_0Z_i,\qquad \ell(\eta)=L_0Z_\eta,
 \qquad Z_i+Z_\eta\ne0.
\]
Comparing the $U$ and $h$ coefficients in the alleged relation
gives $L_0=Z_i^2$, or $\beta=x(i_2)^2$.
Thus that line would be tangent at $i_2$, despite already containing
the three distinct points $i_2,\eta_2,g_2$. This is impossible.
The derivative assertion follows.

For the final assertion, $i=2g$ has $\ell(i)\ne0$: the zeros of
$\ell$ are $g_2,-2g_2$, neither equal to $2g_2$ for nonzero odd
$g$ with $3g\ne0$. If $v$ lay in the span of $l_i,l_\eta$, the
four coordinates other than $v$ in the normalized rows would be
proportional. Their $T$ and $h$ entries force $X_i=X_\eta$.
Hence $\eta=-i$ (the other possibility $\eta=i$ is excluded).
Their nonzero $U$ entries then force proportionality scalar one,
but $\ell(-i)=\ell(i)+1$, a contradiction. This proves the last
claim.
\end{proof}

\begin{lemma}[No unwanted collisions on the third-value slice]
\label{lem:ss-third-value-other-collisions}
\status{Proved under stated hypotheses}
Suppose first that $\eta\ne-g$. A geometric generic section of
$H_\eta$ has exactly its prescribed three-point orbit and no
other triple identifications under $K$. All pair identifications
are transverse except the pair $P,Q$. Its contact length is two
if $3g\ne0$ and three if $3g=0$.

If $3g=0$ and $\eta=-g$, a generic section of $H_\eta$ has no
nontransverse pair of difference outside $\langle g\rangle$,
and no triple with at least one pairwise difference outside that subgroup.
No assertion about its prescribed three tangent directions is
made in this second case.
\end{lemma}
\begin{proof}
For each $j\ne0$ use the proper two-value incidence over $H_\eta$.
Its generic coefficient fibre over the moving point $R$ is a
projective line. The projected-value lemma shows that there is
one integral three-dimensional main component, and that all other
components dominating $H_\eta$ are the constant base-point pairs.
In the cyclic case the additional curve strata lie over $F=0$,
which we remove. The intersection family is generically finite,
and its total length is $L^2=18$.

First assume $\eta\ne-g$ and $j\ne\pm g$.
If $j\ne\eta$, set $R=P$ and impose $l_j=0$ as well as
$l_\eta=0$. By the linear tests choose $h\partial_2s(j)\ne0$.
The two value equations have invertible Jacobian in $R$, so this
is an etale point of the main component. If $j=\eta$, instead
take $R=-\eta$, impose $l_{-\eta}=0$, and use the same test
with the constant second point $P$. In both cases the moving
component is dominant and generically etale.

For $j=g$ with $3g\ne0$ and $\eta\ne2g$, take $R=Q$.
The two value equations are the prescribed zero at $Q$ and the
additional condition $l_{2g}=0$; the last linear test permits
$v\partial_2s(2g)\ne0$. This again gives an etale point of the
main component. If instead $\eta=2g$, the constant pairs have
lengths two at $P$ and one at $Q$: the third branch has nonzero
second-factor derivative, as its derivative functional has nonzero $U$
coefficient whereas $l_\eta$ has nonzero $v$ coefficient.
The moving component therefore has degree $18-3=15$, which is
odd and hence separable in characteristic two.
If $3g=0$ (still $\eta\ne-g$), the sole constant $g$-pair is
$P,Q$ with length three. Its coefficient $T+vF/h$ is not identically
zero on $H_\eta$. Indeed its homogeneous numerator $hT+vF$ is an
irreducible quadratic in the independent variables $F,T,v,h$,
so no hyperplane divides it. The moving degree is again $15$.
The contact length is two in the nonflex case by
Lemma~\ref{lem:ss-uniform-local-contact}. The assertion for $-g$
is its translate. All other constant pairs involve the third point
and are transverse on a dense open, by the nonzero second-factor derivative
there and the nonzero first derivatives $h,v$ at $P,Q$.
Thus after shrinking the coefficient base all moving intersections
are etale and the only nontransverse prescribed pair is $P,Q$.

In the cyclic case $H_{-g}$ has $T=F$.
For every $j\notin\{0,g,-g\}$, take $R=P$ and impose $l_j=0$.
The linear tests give $h\partial_2s(j)\ne0$, proving that the
unique moving component is generically etale. There is no constant
base pair of such a difference.

For triples, the projected-value lemma gives generic rank three;
the coefficient fibre is then a point. The rank-two stratum has
base dimension at most one and coefficient fibres of dimension
one. Rank-one strata are finite with coefficient fibres of dimension
two; in the cyclic exceptional case their curve parts require
$F=0$ and have been removed. Thus every nonconstant triple
incidence has dimension at most two, and cannot dominate
$H_\eta\simeq\mathbf P^3$. The only rank-zero triples are the
prescribed three base points. This proves precisely the assertions
about triples in both cases. There are finitely many differences
and triples, so all stated open restrictions can be imposed at once.
\end{proof}

\begin{proposition}[A tangent-triple bridge for every kernel triple]
\label{prop:ss-uniform-tangent-triple}
\status{Proved under stated hypotheses}
Let $K\subset E^2(k)$ be a finite odd subgroup with injective
coordinate projections. For every triple of distinct elements of
$K$, there exists a smooth sharp source $C\in|L|$ whose image
in $X/K$ has one singularity with those three preimages: two
branches have simple contact and the third is transverse.
Every other singularity is an ordinary node.
Thus the already proved local tangent-triple Cartier criterion
applies to every triple-collision incidence.
\end{proposition}
\begin{proof}
Use the ordering $\{0,g,\eta\}$ fixed above.
Lemma~\ref{lem:DE-A4-uniform-triple-local-type} constructs the
irreducible main rank-five triple-contact incidence, a
$\mathbf P^3$-bundle over an open of the moving point and tangent
base. It contains our vertical slice $H_\eta$.
That lemma and Corollary~\ref{cor:DE-A4-third-value-smooth} give
a nonempty open of globally smooth sharp sources with two simply
tangent branches and a transverse third branch.

If $\eta\ne-g$, the preceding lemma supplies another nonempty
open, excluding every unwanted collision. We justify openness
at its prescribed points. The third branch is transverse already
on the vertical slice, so the three-value incidence is reduced
at the prescribed triple: two of its value derivatives are
independent. Thus an extra triple cannot specialize into it.
All constant transverse pairs are likewise etale. The only remaining
cluster is the prescribed tangent pair, of length two or three.
In the full contact incidence it always has length at least two;
a second tangency in that cluster would require length at least
four. The finite flat intersection argument of
Proposition~\ref{prop:ss-uniform-graph-mixed-contact} therefore
gives the required open even when $3g=0$.
Intersect the two nonempty opens in the irreducible main incidence.

For $3g=0$ and $\eta=-g$, apply the preceding lemma only to
differences outside $\langle g\rangle$ and triple patterns with at least
one pairwise difference outside $\langle g\rangle$.
They have no prescribed base pairs or triples, so their avoidance
is open directly by the proper incidence maps.
For the collisions inside $\langle g\rangle$, use the already
certified collinear tangent-triple witness of
Lemma~\ref{lem:ss-global-tangent-triples}.
Product automorphisms preserving $O$ carry its nonzero three-torsion
generator to $g$ in both coordinates. The resulting witness has
exactly one transverse-third tangent triple and no other
$\langle g\rangle$-collision. This supplies a nonempty open of
the same main incidence. Its four contact rows and its third value
have rank five, as checked at the actual witness in the retained
program
\path{agents/monodromy/code/verify_ss_tangent_incidence_rank.py}.
All source, contact and internal-collision conditions are open at
this witness; the transverse third makes the prescribed triple
reduced. Irreducibility lets this open intersect all of the
outside-collision opens obtained from the vertical slice.

In either case translated intersections are finite, so the smooth
source maps birationally to its image. The collision conditions
give exactly the asserted singularities. Sharpness is
$\det S\,Q_H(S)\ne0$, equivalently the only Cartier-killed source
forms are the two invariant ambient forms. The local conclusion
is Lemma~\ref{lem:ss-tangent-triple-bridge}. No classification of
the resulting full monodromy group is included in this proposition.
\end{proof}


% /root/monodromy, 2026-11-19. Complete root, reduction-descent and obstructions-crystals proof/source audits passed.
\subsection{Every odd degree on the supersingular product}
\label{subsec:ss-all-odd-groups}
\status{Proved and independently checked. The parent, reduction/descent and
obstructions/crystals agents independently audited all three complete
proofs and their prerequisites, and each actually reread the original
transvection-group classification. The conclusion covers every odd
positive degree. The even-degree conclusion is established later in
Theorem~\ref{thm:ss-product-all-full-groups}.}

Put $E:y^2+y=x^3$, $X=E\times E$, and
$L=\mathcal O_E(3O)\boxtimes\mathcal O_E(3O)$ in characteristic two.
We first work over an algebraically closed field.
For $d\ge1$ write $P_d=|L^d|$ and let $h_d$ be the determinant of
relative coherent Frobenius on $P_d$.

\begin{lemma}[All odd norm families and their divisorial exceptions]
\label{lem:ss-all-odd-norm-factors}
\status{Proved under stated hypotheses}
For every odd $d\ge3$ there is an etale self-isogeny
$\phi_d:X\to X$ of degree $d^2$, with kernel having injective
coordinate projections, and a norm morphism
\[
 \nu_d:P_1=\mathbf P^8\longrightarrow P_d,
 \qquad \nu_d^*\mathcal O_{P_d}(1)=\mathcal O_{\mathbf P^8}(d^2).
\]
In the tensor monomial basis write
$D(s)=\det(s_{ij})$ and
$Q(s)=s^{[2]t}(J_3\otimes J_3)s$.
Every irreducible factor of $\nu_d^*h_d$ is one of $D,Q$, a
pair-tangency divisor $R_g$, or a triple-identification divisor
$T_{g,\eta}$ for distinct $0,g,\eta\in\ker\phi_d$.
Every $R_g$ and $T_{g,\eta}$ is an irreducible divisor distinct
from $D,Q$, and its generic image lies in $h_d=0$.
The multiplicity of $D$ in $\nu_d^*h_d$ is exactly two.
\end{lemma}

\begin{proof}
Choose $b\in\End(E)$ with $b^\dagger b=bb^\dagger=[d-1]$,
using the explicit quaternion norm construction in
Lemma~\ref{lem:DE-A3-ss-norm}, and put
\[
 \phi_d=\begin{pmatrix}1&b\\-b^\dagger&1\end{pmatrix}.
\]
Then $\phi_d^\dagger\phi_d=\phi_d\phi_d^\dagger=[d]I$.
Its degree is $d^2$, and its kernel is the graph
\[
 K_d=\{(-bP,P):P\in E[d]\}.
\]
The restriction of $b$ to $E[d]$ is invertible, since
$d-1$ is relatively prime to $d$; hence both kernel projections
are injective. All these finite groups are etale since $d$ is odd.
The exact norm-line calculation of
Theorem~\ref{thm:ss-uniform-GL8-block} gives
$\operatorname{Nm}_{\phi_d}L\simeq L^d$:
the numerical classes agree, and their symmetric difference is
geometric two-torsion in $\Pic^0(X)$, which is zero on this
supersingular product. Norms of sections are homogeneous of degree
$d^2$, and the norm of a nonzero section is nonzero. They therefore
give the asserted morphism and pullback line bundle.

We give the incidence details for arbitrary $K_d$, rather than
using the finite orbit calculation for degree three.
For $g\ne0$ the pair-contact conditions at $P,P+g$, with common
tangent $\tau$, impose four linear conditions on $H^0(L)$.
Away from finitely many $P$ the two full first jets have rank six,
and these four conditions have rank four. At the exceptions their
rank is at least three. The main incidence is thus an irreducible
$\mathbf P^4$-bundle over a three-dimensional open of
$X\times\mathbf P^1$, and its exceptional part has dimension at
most six. For triples, both kernel projections are injective, so
the three value tensors are independent everywhere: three distinct
rank-one Segre tensors can be collinear only if one tensor factor
is constant. Hence the triple incidence is an irreducible
$\mathbf P^5$-bundle over $X$, of dimension seven.

Propositions~\ref{prop:ss-uniform-graph-mixed-contact} and
\ref{prop:ss-uniform-tangent-triple} supply smooth sources in
these main incidences with finite translated intersections.
Consequently their incidence maps have finite fibres there and
their images are irreducible divisors. A section invariant under
a nonzero kernel translation cannot be a generic exception: such
a translation must preserve $L$, hence must be three-torsion.
On each elliptic degree-three space a nonzero three-torsion
translation has three distinct projective eigenlines; its tensor
action has eigenspaces of dimension three. Thus invariant curves
lie in a finite union of $\mathbf P^2$'s. If the translation does
not preserve $L$, no divisor in $|L|$ is invariant under it.

Neither collision divisor equals $Q$. Indeed the nondegenerate
Frobenius cubic $Q$ contains no vector subspace of dimension
greater than four: an identically zero restriction to $U$ gives
$U^{[2]t}(J_3\otimes J_3)U=0$, and so $2\dim U\le9$.
The pair and triple coefficient fibres have vector dimensions
five and six. Neither equals $D$, either. At a generic pair,
normalize the two values in each factor to $e_0,e_1$, choose the
common tangent with nonzero first component, and impose
$S_{00}=S_{11}=0$, $S_{01}=S_{10}=1$.
The tangent equations solve for $S_{20},S_{21}$, while the
coefficient of the still free $S_{22}$ in $\det S$ is one.
For a generic triple the first-factor values form a basis, and
the three rows of $S$ may be chosen independently in the three
distinct planes orthogonal to the second-factor values.
Those planes admit three independent row vectors. Thus both
generic coefficient fibres contain invertible matrices.

The source singular-curve discriminant $\Delta$ is irreducible
and has degree $c_2(J^1L)=3L^2=54$, with a generic single ordinary
node. The singular-point incidence is a projective bundle; the
two-point jet bounds give dimension at most six for two singular
points, and the second-jet calculation supplies the ordinary node
and generic separable degree one. These are the same first-jet
and Chern-class calculations written out in
Proposition~\ref{prop:ss-level3-three-witness-reduction}; they
involve only the fixed source linear system and are independent
of the chosen odd isogeny. In particular $\Delta$ is distinct
from the two cubic source factors. It is also distinct from
every collision divisor, since each has a smooth-source witness.
At the generic point of $\Delta$ the source node is not identified
with another source point: imposing a singularity at $P$ and
containment at $P+g$ imposes four conditions, except at finitely
many $P$ where at least three remain, so this extra incidence has
dimension at most six. All other identifications there are nodes.

The source relative determinant is $D^2Q$ by
Proposition~\ref{prop:ss-product-relative-Hasse}.
Outside $\Delta$ and the collision divisors, the norm curve is
nodal with smooth normalization $C_s$. The normalization sequence
adds only a Frobenius-unit graph space. At the generic point of
$\Delta$ the same comparison holds, with its one source node
mapping isomorphically to a target node. Thus $\Delta$ adds no
factor to the pullback determinant. The noninvariant integral
source maps birationally onto its image under the etale isogeny;
all failures of ordinary double identification require a tangency
or a triple. The remaining invariant locus has codimension at
least two. This proves the asserted list of all possible factors.

We verify that every collision divisor does occur in the target
Hasse divisor. For a reduced curve singularity with $r$ branches,
the Frobenius-stable rank of the normalization quotient is $r-1$.
Its branch-constant subspace is $k^r/k$, with bijective Frobenius;
modulo these constants, every branch expansion has positive order,
so a sufficiently high Frobenius power lies in the conductor.
The global normalization sequence therefore gives
\[
 g(Z)-f(Z)=\sum_i\bigl(g(C_i)-f(C_i)\bigr)
                  +\sum_x(\delta_x-r_x+1).
\]
The generic source on each collision divisor is smooth and outside
$D\cup Q$, so its defect is two. A tangency, or at least three
branches in a smooth surface, contributes
$\delta-r+1\ge1$. The target defect is consequently at least
three. The relative coherent dimension subtracts the two ambient
dimensions, whose stable rank is zero, so its Frobenius is not
bijective. This proves the converse inclusion just required.

Finally, at the generic point of $D$ the source is smooth and all
norm identifications are transverse double points. Order the nodes
and their branches by a finite etale base change of this local
family. Its graph term in coherent cohomology is then a finite
free Frobenius-unit subbundle. The odd isogeny identifies the two
ambient coherent cohomology subbundles. Hence the relative target
bundle is an extension of the relative source bundle by that unit
graph bundle, over the local ring itself. Multiplication of
Frobenius determinants in this exact sequence gives vanishing
order exactly two along $D$, as in the source polynomial $D^2Q$.
\end{proof}

\begin{theorem}[Irreducible reduced Hasse divisors in all odd degrees]
\label{thm:ss-all-odd-Hasse}
\status{Proved under stated hypotheses}
For every odd $d\ge3$, the relative Hasse divisor on $P_d$ is
geometrically irreducible and reduced. At its geometric generic
point the relative Frobenius has stable rank $9d^2-2$,
one less than its rank $9d^2-1$ off the divisor.
\end{theorem}

\begin{proof}
Use the norm family of the preceding lemma. For each $g\ne0$,
the uniform mixed-contact witness lies on $Q\setminus D$ and
on the main $R_g$. Its target Hasse divisor is smooth there by
Lemma~\ref{lem:ss-mixed-self-tacnode}; the line bundle $L^d$ is
very ample and its polarization is separable because $d$ is odd.
Since $Q$ and $R_g$ are irreducible and their generic images lie
in the target Hasse divisor, they are assigned to the same unique
global component. All $R_g$ are thus assigned to the component
containing the image of $Q$, without requiring any transitivity
of automorphisms on the arbitrary kernel.

For every triple, Proposition~\ref{prop:ss-uniform-tangent-triple}
gives a smooth sharp source with exactly one tangent triple and
otherwise nodes. Lemma~\ref{lem:ss-tangent-triple-bridge} makes
the target Hasse divisor smooth at that point. That point belongs
to its triple divisor and to the tangency divisor of its tangent
pair. It therefore assigns every $T_{g,\eta}$ to the same global
component as $Q$ and all $R_g$. The uniqueness at these smooth
points justifies the assignments even when a pullback factor has
multiplicity.

Any additional target Hasse component would pull back to a
nonzero divisor supported only on $D$, with multiplicity at most
two. Its pullback degree would be at most six. But a nonzero
target hypersurface has pullback degree at least $d^2\ge9$.
The pullback cannot vanish identically, since generic sharp source
curves have nodal norm images and thus maximal relative rank.
This contradiction proves irreducibility. Smoothness at any one
of the mixed-contact witnesses proves multiplicity one.

For the stable-rank assertion, fix a nonzero $g$.
The locus in $R_g$ with globally smooth source, one simple target
tacnode, and otherwise only nodes is nonempty and open, by the
mixed-contact witness and its finite-flat collision analysis.
Since $R_g$ is irreducible and distinct from $D,Q$, this open
contains a sharp source. Its normalization defect is two, and
the single simple tacnode contributes
$\delta-r+1=2-2+1=1$; all nodes contribute zero.
Thus the target has defect exactly three, giving relative stable
rank $9d^2-2$. The rank of a sufficiently high Frobenius iterate
is lower semicontinuous. On the now irreducible Hasse divisor it
is generically at least this value and at most one less than the
full rank. Equality follows.
\end{proof}

\begin{theorem}[Full images in every odd degree on the supersingular product]
\label{thm:ss-all-odd-full-groups}
\status{Proved under stated hypotheses}
Let $k$ be any perfect field of characteristic two, and give
$X=E\times E$ the embedding by
$L=\mathcal O_E(3O)\boxtimes\mathcal O_E(3O)$.
For every odd $d\ge1$, the geometric and arithmetic images of
the relative mod-two Hasse--Witt local system on its maximal-rank
smooth degree-$d$ hypersurface-section locus are
\[
 G_{\mathrm{geom},d}=G_{\mathrm{arith},d}
          =\GL_{9d^2-1}(\mathbf F_2).
\]
In particular this local system is geometrically irreducible in
every odd degree.
\end{theorem}

\begin{proof}
Degree one is Theorem~\ref{thm:ss-product-degree-one-full-group}.
Let $d\ge3$ be odd, first over an algebraically closed field,
and put $W_d=H^0(X,L^d)$, $b=\dim W_d=9d^2$.
The universal divisor sequence, Kunneth vanishing
$H^1(X,L^{-d})=0$, and Serre duality identify the relative
coherent bundle on $P_d=\mathbf P(W_d)$ with the kernel of
\[
 W_d^*\otimes\mathcal O_{P_d}(-1)
                  \longrightarrow\mathcal O_{P_d},
\]
namely $\Omega^1_{P_d}$ by the Euler sequence.
It has rank $b-1$, and it carries regular Frobenius on all of
$P_d$, including its singular-divisor parameters. The ambient
$H^1(X,\mathcal O_X)$ has zero Frobenius, so its unit Lang system
is the relative system in question. It is generically bijective
by the exact odd-degree rank theorem.

The preceding Hasse theorem and determinant allocation
(Lemma~\ref{lem:irreducible-hasse-quotients-trivial}) imply that
every nonzero proper invariant subspace of this Lang
representation would have a nonzero trivial quotient.
But Theorem~\ref{thm:ss-product-exact-covectors} gives no nonzero
invariant covectors in odd degree. Hence the geometric
representation is irreducible.

At the generic point of its single noninvertibility divisor,
stable rank drops by exactly one. Lemma~\ref{lem:char2-defect-one-inertia}
therefore shows that the geometric group is generated by
transvections. Explicitly, the strict henselian Lang finite part
fixes a codimension-one subspace pointwise, and the one-dimensional
quotient over $\mathbf F_2$ is also fixed. Each nontrivial inertia
element has rank-one square-zero difference from the identity.
Purity extends the quotient by all such inertia to a finite etale
cover of projective space, which is trivial.

The group contains an actual natural $\GL_8(\mathbf F_2)$ block
fixing a complementary summand, by
Theorem~\ref{thm:ss-uniform-GL8-block}. Thus it preserves no
nondegenerate alternating bilinear form: the natural block has
neither an invariant bilinear form nor an invariant covector,
so it lies in the radical of any invariant form on the whole
representation.
Apply the irreducible transvection-group classification
\cite[Theorem~3.6, printed p.~22]{KantorSeress2001}, whose exact
primary statement was reopened for this application. Over
$\mathbf F_2$ its alternatives are the full special linear
group, a symplectic group, an orthogonal group in its natural
even-dimensional representation, or the indicated symmetric
groups on their deleted permutation modules. All alternatives
except the special linear group preserve a nondegenerate
alternating form: in the orthogonal case use the polar form,
and in the permutation cases use the dot product on the
sum-zero space, taking its constant-line quotient when necessary.
They have just been excluded. Since
$\SL_{b-1}(\mathbf F_2)=\GL_{b-1}(\mathbf F_2)$, the claimed
geometric group follows.

Restriction from the Frobenius-invertibility open to its dense
smooth-divisor open preserves connected components of every finite
etale frame cover, so this is the group on the requested smooth
family. The embedding also gives the full complete linear system
in every degree. Indeed each plane cubic has its section ring
generated in degree one (restrict polynomials on $\mathbf P^2$;
$H^1(\mathbf P^2,\mathcal O(d-3))=0$). The degree-$d$ product
sections are tensors of products of $d$ degree-one sections in
the two factors, hence products of $d$ tensor degree-one sections.
Ambient degree-$d$ forms therefore surject onto $H^0(L^d)$.
The induced ambient parameter map has affine-space fibres after
removing its kernel, so pullback preserves the geometric image.

Finally all geometric frame-cover connectedness assertions persist
under algebraically closed extension. Over an arbitrary perfect
base field the arithmetic image contains the full geometric
general linear group and is itself a subgroup of that same
general linear group. It is therefore equal to it.
\end{proof}

\begin{remark}[Scope of the uniform conclusion]\label{auto:A_monodromy:0c2f1cb265b6}
\status{Proved under stated hypotheses}

The theorem closes the entire odd-degree full-image problem for
this fixed supersingular product and its indicated embedding.
Even degrees have generic rank $9d^2-3$ and exactly one invariant
covector. The following first-image sheaf and boundary degree comparison
turn these restrictions into the exact uniform even image in
Theorem~\ref{thm:ss-product-all-full-groups}.
\end{remark}


% /root/monodromy, 2026-11-19. Complete root, reduction-descent and obstructions-crystals proof audits passed.
\subsection{The first-image sheaf and the remaining even-degree gate}
\label{subsec:ss-even-image-sheaf}
\status{Proved and independently checked. The parent, reduction/descent and obstructions/crystals agents audited all three complete proofs. The conditional geometric hypothesis is discharged for every even degree at least four in Theorem~\ref{thm:ss-even-unit-Hasse-irreducible}.}

Continue with $E:y^2+y=x^3$, $X=E^2$ and
$L=\mathcal O_E(3O)\boxtimes\mathcal O_E(3O)$ over an
algebraically closed field of characteristic two. Let $d\ge2$ be
even, $m=3d$, $N=L^d$, $W=H^0(N)$ and $b=m^2=9d^2$.
Write $P=\mathbf P(W)$, $\mathcal B=W^*\otimes\mathcal O_P(-1)$,
and let $\delta_1,\delta_2$ be the canonical lifted invariant
derivations on $N=F_X^*N_{m/2}$.

\begin{lemma}[A negative reflexive extension of the unit image]
\label{lem:ss-even-image-sheaf}
\status{Proved under stated hypotheses}
The space $W$ is a free module of rank $b/4$ over
\[
 \mathcal A=k[\delta_1,\delta_2]/(\delta_1^2,\delta_2^2).
\]
Define reflexive sheaves on $P$ by
\[
 \mathcal I=\ker\bigl(\mathcal B\xrightarrow{(s,\delta_1s,\delta_2s)}
                                  \mathcal O_P^3\bigr),\qquad
 \mathcal K=\ker\bigl(\mathcal B\xrightarrow{
 (s,\delta_1s,\delta_2s,\delta_1\delta_2s)}\mathcal O_P^4\bigr).
\]
Let $Z_0\subset P$ be the linear locus $\delta_1\delta_2s=0$.
It has codimension $b/4$. On $P^\circ=P\setminus Z_0$ these are
vector bundles of ranks $b-3,b-4$ and fit into a Frobenius-equivariant
exact sequence
\[
 0\longrightarrow\mathcal K\longrightarrow\mathcal I
        \xrightarrow{\lambda}\mathcal O_{P^\circ}\longrightarrow0,
 \qquad \lambda=\delta_1\delta_2s,
\]
whose quotient has the constant unit Frobenius. Moreover
\[
 \det\mathcal I=\det\mathcal K=\mathcal O_P(-b),\qquad
 H^0(P,\mathcal I)=H^0(P,\mathcal K)=0.
\]
The full coherent Frobenius maps into $\mathcal I$; its restriction
there is generically bijective and has the required relative unit
Lang system. Its determinant, and the determinant on $\mathcal K$,
give the same nonzero homogeneous polynomial
$h_d^{\mathrm u}\in H^0(P,\mathcal O_P(b))$, up to a common
nonzero constant.
\end{lemma}

\begin{proof}
On either elliptic factor the proof of
Theorem~\ref{thm:ss-product-exact-covectors} gives
\[
 F_*\mathcal O_E(mO)=U\otimes\mathcal O_{E^{(2)}}((m/2)O),
 \qquad 0\to\mathcal O\to U\to\mathcal O\to0
\]
with a nonsplit extension, and $\delta$ is a nonzero multiple of
its quotient-followed-by-inclusion map. Twist this sequence by
the positive-degree line bundle $\mathcal O((m/2)O)$.
Its first cohomology vanishes, so on global sections the kernel
and image of $\delta$ agree and both have dimension $m/2$.
Thus $H^0(\mathcal O_E(mO))$ is a free module of rank $m/2$
over $k[\delta]/\delta^2$. Tensoring the two factors proves
the asserted free $\mathcal A$-module structure.

The operator $\delta_1\delta_2$ has rank $b/4$ on this free
module, proving the codimension of $Z_0$.
If $\delta_1\delta_2s\ne0$, the four vectors
\[
                  s,\quad\delta_1s,\quad\delta_2s,
                       \quad\delta_1\delta_2s
\]
are independent. Apply $\delta_1\delta_2$, then $\delta_1$,
then $\delta_2$ to any relation to kill its first three
coefficients; its last coefficient is then zero as well.
This proves all bundle and surjectivity assertions on $P^\circ$.
More precisely the first three vectors lose rank only when
$(a\delta_1+b'\delta_2)s=0$ for some $[a:b']\in\mathbf P^1$.
Every such nonzero nilpotent operator has rank $b/2$.
Its projective kernel incidence has dimension $b/2$, with
generically unique $[a:b']$, so this smaller rank-drop locus
has codimension $b/2-1$ and lies inside $Z_0$.

Each displayed kernel sheaf is reflexive: it is the kernel of
a vector bundle mapping to a torsion-free image sheaf, so the
depth lemma gives the depth-two condition. Since the evaluation
maps are surjective outside a locus of codimension at least two,
their determinant lines on that open extend as
$\mathcal O_P(-b)$ over all of $P$.
Their inclusions in $\mathcal B$ give the claimed vanishing of
global sections.

Let $M(s)$ denote the root-removed Cartier multiplication matrix,
so full Frobenius has matrix $M(s)^t$ with the usual Frobenius
twists. The previously proved identities are
\[
 M(s)s=M(s)\delta_1s=M(s)\delta_2s=0,
 \qquad M(s)\delta_1\delta_2s=(\delta_1\delta_2s)^{[2]}.
\]
They may also be read off in a pullback frame from
$s=a^2+x_1b_1^2+x_2b_2^2+x_1x_2c^2$ by taking the
$x_1x_2$ coefficient. The first three identities put the full
Frobenius image in $\mathcal I$; the fourth says that $\lambda$
is a morphism to the constant unit-Frobenius line and that
$\mathcal K$ is Frobenius-stable. The exact even-degree stable
rank is $b-3$, so Frobenius on $\mathcal I$ is generically
bijective. Every unit vector of the full coherent bundle lies
in its Frobenius image, hence in $\mathcal I$; consequently the
two unit Lang systems agree on their common dense open.

The determinant on $\mathcal I$ is a regular section of
$\det\mathcal I\otimes(F^*\det\mathcal I)^{-1}
=\mathcal O_P(b)$ away from $Z_0$.
Normality extends it across this codimension-at-least-nine locus.
The quotient in the displayed exact sequence is Frobenius-unit,
so the determinant on $\mathcal K$ agrees with it. Generic
bijectivity proves that this homogeneous polynomial is nonzero.
\end{proof}

\begin{lemma}[An even-degree point with stable loss exactly one]
\label{lem:ss-even-one-loss-point}
\status{Proved under stated hypotheses}
For every even $d\ge2$, the divisor $h_d^{\mathrm u}=0$ has a
point in $P^\circ$ at which the Frobenius-stable rank of
$\mathcal I$ is $b-4$.
\end{lemma}

\begin{proof}
Choose a smooth sharp odd-degree curve $D\in|L^{d-1}|$.
At a general point $P_0\in D$, both coordinate tangent components
are nonzero: each projection $D\to E$ has odd degree $3(d-1)$
and is therefore separable. Choose $P_0$ also away from the two
elliptic flex sets. The codimension-two linear system of
$C\in|L|$ containing $P_0$ and tangent to $D$ there is a
projective six-space.

A generic member of this six-space is globally smooth. At $P_0$
the remaining normal derivative is free. At another point $Q$,
its singularity conditions together with the two imposed
conditions have rank five off the exceptional joint-jet locus,
and at least four on that locus. The latter has dimension at
most one, by the product-cubic first-jet calculation. Thus the
moving singular-member incidence has dimension at most five
inside the indicated six-space. Second jets at $P_0$ are
separated, so the contact can be required to have exactly order
two. The two curve-jet conditions at another point of $D$ have
combined rank four except at finitely many points and at least
three there. The same incidence count excludes every second
double intersection, as in the node-cover proof of
Theorem~\ref{thm:ss-uniform-GL8-block}.

The sharp-source open also meets this six-space. It cannot lie
in the cubic $Q$, which contains no projective linear space of
dimension greater than three. It does not lie in $\det S=0$,
either: normalize the two factor values at $P_0$ to $e_0$ and
their tangent directions to $e_1$. Then the conditions are
$S_{00}=0$ and $aS_{10}+b'S_{01}=0$, with $a,b'\ne0$.
Choose nonzero $S_{10},S_{01}$ satisfying that equation,
set $S_{22}=1$, and all other entries zero. This is an
invertible matrix. All the nonempty opens just described
therefore intersect in the irreducible six-space.

For the resulting $Y=C\cup D$, the sole nonnodal singularity
is a simple tacnode between the two components.
Each normalized component has defect two, and the tacnode
contributes $\delta-r+1=1$. The normalization formula gives
\[
                 g(Y)-f(Y)=2+2+1=5,
 \qquad f(Y)=b-4.
\]
There are also $18(d-1)-2>0$ ordinary intersection nodes.
At such a node the local equation is a product of two transverse
parameters, and $\delta_1\delta_2s$ has nonzero value: it is the
mixed second coefficient in the invariant tangent coordinates.
Thus $Y\notin Z_0$. On $P^\circ$ the full Frobenius image lies
in $\mathcal I$, so its stable rank equals the stable rank on
$\mathcal I$. The full and relative stable ranks agree, since
the ambient supersingular cohomology is nilpotent. This gives
the asserted point and rank.
\end{proof}

\begin{theorem}[A precise conditional even-degree full-image reduction]
\label{thm:ss-even-full-group-if-unit-Hasse-irreducible}
\status{Proved under stated hypotheses}
Let $k$ be any perfect field of characteristic two, and fix even
$d\ge2$ on the indicated surface and embedding. If the divisor
$h_d^{\mathrm u}=0$ is
geometrically irreducible and reduced, then the entire geometric
and arithmetic images are the full pointwise stabilizer of the
unique nonzero invariant covector:
\[
 G_{\mathrm{geom},d}=G_{\mathrm{arith},d}
   =\mathbf F_2^{\,9d^2-4}\rtimes\GL_{9d^2-4}(\mathbf F_2)
       \subset\GL_{9d^2-3}(\mathbf F_2).
\]
\end{theorem}

\begin{proof}
First work over an algebraic closure. The determinant-allocation
proof of Lemma~\ref{lem:irreducible-hasse-quotients-trivial}
applies to $\mathcal I$. We give the only changes required by
its reflexive extension. Saturation of an invariant subspace in
$\mathcal I$ is reflexive by the depth lemma, and its quotient
is torsion-free. In codimension one all terms are vector bundles,
so their effective Frobenius determinant divisors add to the
reduced irreducible $h_d^{\mathrm u}$ divisor. A zero determinant
divisor makes the corresponding reflexive Frobenius map an
isomorphism everywhere, by reflexive extension across codimension
two. Lemma~\ref{lem:periodic-reflexive-projective-trivial} then
makes that sheaf a trivial unit-Frobenius bundle.
Allocating no divisor to a nonzero subobject would contradict
$H^0(P,\mathcal I)=0$. Thus every proper nonzero invariant
subspace has trivial quotient.

The full invariant-covector space has dimension one, by
Theorem~\ref{thm:ss-product-exact-covectors}.
Consequently the common kernel $K$ of that covector is the
only proper nonzero invariant subspace: any such subspace has
constant quotient of dimension at most one, and the quotient
covector is the unique one. In particular $K$ is irreducible
and the whole representation has no invariant line.

The preceding lemma and irreducibility of the determinant
divisor imply generic stable loss exactly one there, by lower
semicontinuity of the rank of a high Frobenius iterate.
The Lang finite-part and purity argument of
Lemma~\ref{lem:char2-defect-one-inertia} therefore makes the
geometric group generated by transvections. Removing $Z_0$
causes no change to that purity argument: it has codimension
at least nine, so the quotient cover is still unramified at
every codimension-one point of the regular projective space.
Restriction to $K$ is thus an irreducible group generated by
transvections (discard any generators restricting to identity).
The actual natural $\GL_8(\mathbf F_2)$ block lies in $K$ and
excludes invariant alternating forms on it, by
Theorem~\ref{thm:ss-uniform-GL8-block}.
The exact classification
\cite[Theorem~3.6, printed p.~22]{KantorSeress2001}
therefore gives the full image $\GL(K)$.

The kernel of the map from the whole group to $\GL(K)$ lies
in the translation module $K$ and is either zero or all of $K$.
If zero, the subgroup is the graph of a one-cocycle on $\GL(K)$.
The proved elementary vanishing
$H^1(\GL_n(\mathbf F_2),\mathbf F_2^n)=0$ for $n\ge4$
(Lemma~\ref{lem:char2-group-inputs}) makes it conjugate to a
linear complement. That complement has a fixed line outside
$K$, contradicting the preceding determinant-allocation result.
Thus every translation occurs, proving the stated geometric group.
The covector is defined over $\mathbf F_2$, so the arithmetic
image is contained in the same stabilizer and contains the
geometric image. Equality follows, with the same dense-open and
ambient-linear-system passage used in
Theorem~\ref{thm:ss-all-odd-full-groups}.
\end{proof}

\begin{remark}[The geometric hypothesis and retained experiments]\label{auto:A_monodromy:feea204b71c7}
\status{Proved under stated hypotheses}

The homogeneous polynomial $h_d^{\mathrm u}$ is now specified
intrinsically, with its degree, vector bundle on a complement of
known codimension, and one genuine stable-loss-one point proved
in the preceding lemmas. The next subsection proves its geometric irreducibility and
reducedness for every even $d\ge4$ by boundary degree comparisons,
thereby discharging the conditional hypothesis in those degrees.
The retained exploratory programs
\path{agents/monodromy/code/ss_even_unit_hasse_pencils.py} and
\path{agents/monodromy/code/ss_even_generic_hasse_pencils.py}
compute exact degree-two quotient determinants on six constant-minor
and four moving-minor pencils. Their outputs retain both repeated
and squarefree pencil factorizations. They do not establish a
uniform geometric assertion; the following proof supplies the needed
geometric hypothesis independently of these experiments.
\end{remark}


% /root/monodromy, 2026-11-19. Complete root, reduction-descent and obstructions-crystals proof audits passed.
\subsection{Two boundary degree comparisons for the even images}
\label{subsec:ss-all-even-groups}
\status{Proved and independently checked. The parent, reduction/descent and obstructions/crystals agents audited all four complete proofs, including the inseparable double-point incidence, its reduced Frobenius determinant, every degree comparison and the all-degree group conclusion. No exploratory pencil computation is used.}

Use $X=E^2$, $E:y^2+y=x^3$, $L=\mathcal O(3O)\boxtimes\mathcal O(3O)$,
and the polynomial $h_d^{\mathrm u}$ of
Lemma~\ref{lem:ss-even-image-sheaf}.
Work over an algebraic closure until the final group statement.

\begin{lemma}[The intersection-point Hasse divisor]
\label{lem:ss-intersection-Hasse-degree}
\status{Proved under stated hypotheses}
Let $a,e$ be positive odd integers and let $D\in|L^e|$ be a fixed
smooth sharp curve. On $P_a=|L^a|$, the divisor $R_D$ parametrizing
nonreduced intersections $C\cap D$ is irreducible and reduced,
and has degree
\[
                       \deg R_D=18ae+9e^2.
\]
Its generic member has exactly one intersection of multiplicity
two and all other intersections simple. Its Frobenius equation
is the determinant on the finite intersection algebra.
\end{lemma}

\begin{proof}
Remove the locus of sections vanishing identically on $D$.
It is empty if $a<e$. If $a=e$ it is a point in a projective
space of dimension $9a^2-1\ge8$. If $a>e$ its codimension is
\[
                  9a^2-9(a-e)^2=9e(2a-e)>1.
\]
On the remaining open $B$, the intersection family
$Z\subset D\times B$ is finite flat of degree $n=18ae$.
Indeed it is a relative effective Cartier divisor, defined by a
section nonzero on every curve fibre; it is flat and has
zero-dimensional fibres, so properness makes it finite.
Write $\mathcal A=\pi_*\mathcal O_Z$. Cohomology of the divisor
sequence gives an exact sequence of vector bundles
\[
 0\to\mathcal O_B\to\mathcal A\to
 H^1(D,L^{-a}|_D)\otimes\mathcal O_B(-1)
       \to H^1(D,\mathcal O_D)\otimes\mathcal O_B\to0.
\]
Here $g(D)=9e^2+1$ and Riemann--Roch gives
$h^1(D,L^{-a}|_D)=n+g(D)-1=n+9e^2$.
Thus $\det\mathcal A=\mathcal O_B(-n-9e^2)$.
Its Frobenius determinant is a section of
$\mathcal O_B(n+9e^2)$ and extends over the removed
codimension-at-least-two locus. It vanishes precisely when the
finite geometric intersection scheme is nonreduced.

The double-zero incidence over $D$ is a projective bundle:
the two conditions on the value and tangent derivative at any
point are independent because $L^a$ is very ample.
It is irreducible of dimension $\dim P_a-1$.
We check the generic contact and its multiplicity.
For $a\ge2$, the two factors of degree $3a\ge6$ interpolate
first jets at two points and second jets at one point. Tensoring
these interpolations gives full joint surface first jets and
full surface second jets. This follows directly from
$H^1(E,\mathcal O_E(3aO)(-T))=0$ for every effective
subscheme of length at most four. Hence a prescribed double
zero can be made exactly double, and another double zero
imposes two further conditions with its point varying in a
curve. The latter incidence is proper inside the double-zero
incidence. If $a=1$, choose the point of $D$ away from the two
elliptic flex sets. The joint curve jets at this point and a
second point have rank four except at finitely many points and
at least three there: the exceptional surface-jet locus is the
union of two coordinate fibres and finitely many crossed
tangencies, and $D$ is integral and ample. The same dimension
count applies, and the product-cubic second jets give exact
contact two. This is the complete double-zero calculation also
used in Lemma~\ref{lem:ss-even-one-loss-point}.

Thus the generic member of the incidence has exactly one
double zero. The local length-two factor of its finite algebra
has a Weierstrass equation
\[
                             z^2+b'z+c.
\]
The value and first-jet interpolation makes $b',c$ independent
local deformation parameters. Frobenius in the basis $1,z$
has matrix $\left(\begin{smallmatrix}1&c\\0&b'\end{smallmatrix}\right)$,
so its determinant is $b'$, with vanishing order one.
All other points contribute unit determinants.
The marked double-point incidence may map generically inseparably
of degree two to its image in characteristic two; no separability
of that incidence map is used. The displayed determinant directly
computes the reduced equation on the coefficient base itself.
Consequently the extended Frobenius determinant has the reduced
irreducible double-zero divisor as its entire divisor. Its
degree is exactly $n+9e^2$, as asserted.
\end{proof}

\begin{lemma}[An exact factorization on an odd-component boundary]
\label{lem:ss-even-boundary-factorization}
\status{Proved under stated hypotheses}
Let $a\ge3$ and $e\ge1$ be odd, put $d=a+e$, and fix a smooth
sharp $D\in|L^e|$. The linear embedding
\[
 j_D:P_a\hookrightarrow P_d,\qquad C\longmapsto C\cup D
\]
satisfies, up to a nonzero scalar,
\[
                         j_D^*h_d^{\mathrm u}=h_a\,r_D,
\]
where $h_a$ is the irreducible reduced odd-degree relative
Hasse polynomial and $r_D$ is the irreducible equation of
$R_D$ from the preceding lemma.
For $a=1$, the same restriction is nonzero and is divisible
by $r_D$; no complete factorization in this case is needed below.
\end{lemma}

\begin{proof}
A generic smooth sharp $C$ meets $D$ transversely. Its nodal
union has normalization defect $2+2=4$, so its stable rank is
$9d^2-3$, the generic even rank. At each transverse node the
mixed derivative $\delta_1\delta_2(s_Cs_D)$ is nonzero, as in
Lemma~\ref{lem:ss-even-one-loss-point}. Thus the union
lies in the vector-bundle domain $P^\circ$ and the restricted
polynomial is not zero.

For $a\ge3$, the odd Hasse divisor is irreducible and reduced
of degree $9a^2$, by Theorem~\ref{thm:ss-all-odd-Hasse}.
Its generic curve is smooth. To check this explicitly, the
singular-curve discriminant for $|L^a|$ is irreducible of
degree $c_2(J^1L^a)=54a^2$: joint first jets and second jets
give a generic unique ordinary node, with separable degree one
for the singular-point incidence. Its degree differs from
$9a^2$, so it cannot equal the Hasse divisor.
Also $R_D$ has degree $18ae+9e^2\ne9a^2$; equality would give
$(a+e)^2=2a^2$, impossible for positive integers.
Thus the generic curve on $h_a=0$ is smooth and transverse to
$D$. Its defect is at least three, so the union defect is at
least five and its unit-image determinant vanishes.
There remain transverse intersection nodes, so these points
are in $P^\circ$. Hence $h_a$ divides $j_D^*h_d^{\mathrm u}$.

The generic member of $R_D$ is smooth and sharp as a source
curve. For $a\ge3$ this follows from irreducibility and the
degree comparisons with $h_a$ and the singular-curve
discriminant. The latter degree cannot equal $18ae+9e^2$,
since equality would give $(a+e)^2=7a^2$.
For $a=1$, the tangent-six-space construction in
Lemma~\ref{lem:ss-even-one-loss-point} supplies a
smooth sharp member directly, so the same assertion follows
by irreducibility of $R_D$.
Its union with $D$ has one simple cross tacnode and otherwise
nodes. Its defect is $2+2+1=5$, so the restricted unit-image
Hasse polynomial vanishes on $R_D$. The remaining ordinary
nodes again place the generic point in $P^\circ$.
This proves the claimed divisibility by $r_D$ in all cases.

For $a\ge3$ the two irreducible factors are distinct, and
their degrees add to the full degree of the restriction:
\[
             9a^2+(18ae+9e^2)=9(a+e)^2=9d^2.
\]
The nonzero restriction therefore equals their product, with
both multiplicities exactly one. No unproved comparison of
coherent first-image lattices is used: divisibility follows
from stable ranks, and the equality follows from this degree
identity.
\end{proof}

\begin{theorem}[The even unit-image Hasse divisor]
\label{thm:ss-even-unit-Hasse-irreducible}
\status{Proved under stated hypotheses}
For every even $d\ge4$, the divisor $h_d^{\mathrm u}=0$ is
geometrically irreducible and reduced.
\end{theorem}

\begin{proof}
Use the boundary with $(a,e)=(d-1,1)$. Its restricted
polynomial is a product of two distinct irreducible polynomials
of degrees
\[
                        9(d-1)^2,\qquad18d-9.
\]
Every nonzero homogeneous global factor restricts to a nonzero
homogeneous polynomial of the same positive degree: the whole
restriction is nonzero, and $j_D$ is a linear projective
embedding. The squarefree restriction thus proves that
$h_d^{\mathrm u}$ is reduced and has at most two irreducible
components. If there are two, their degrees must be exactly
the displayed pair.

For $d\ge6$, use instead a fixed sharp degree-three component,
so $(a,e)=(d-3,3)$. The same argument would force the degrees
of the same two global components to be
\[
                       9(d-3)^2,\qquad54d-81.
\]
The two unordered pairs are different. Equality in the same
order would give $(d-1)^2=(d-3)^2$, impossible for $d\ge6$.
Equality in the opposite order would give
\[
                  d^2-8d+10=0,
                  \quad\text{or}\quad(d-4)^2=6,
\]
also impossible for an integer. This excludes two components.

For $d=4$, the first pair is $81,63$.
Fix a sharp degree-three curve and vary a degree-one curve.
The preceding lemma shows that the restriction of $h_4^{\mathrm u}$
is nonzero and divisible by the irreducible point-intersection
factor of degree $18\cdot1\cdot3+9\cdot3^2=135$.
If the global divisor had two components of degrees $81,63$,
their nonzero restrictions would have those same degrees.
Their product cannot be divisible by an irreducible polynomial
of degree $135$, since such a polynomial must divide one of
the two factors. This is the remaining contradiction.
\end{proof}

\begin{theorem}[All degrees on the fixed supersingular product]
\label{thm:ss-product-all-full-groups}
\status{Proved under stated hypotheses}
Let $k$ be any perfect field of characteristic two, and embed
$X=E\times E$, $E:y^2+y=x^3$, by
$L=\mathcal O_E(3O)\boxtimes\mathcal O_E(3O)$.
For every $d\ge1$, the entire geometric and arithmetic images
of the relative mod-two Hasse--Witt local system are
\[
 G_{\mathrm{geom},d}=G_{\mathrm{arith},d}=
 \begin{cases}
 \GL_{9d^2-1}(\mathbf F_2),&d\text{ odd},\\[2pt]
 \mathbf F_2^{\,9d^2-4}\rtimes\GL_{9d^2-4}(\mathbf F_2),
                                      &d\text{ even}.
 \end{cases}
\]
The even image is the full pointwise stabilizer of the unique
nonzero invariant covector in rank $9d^2-3$. The odd system
is geometrically irreducible; the even system has its unique
proper nonzero invariant subspace of codimension one.
\end{theorem}

\begin{proof}
The odd degrees are Theorem~\ref{thm:ss-all-odd-full-groups}.
For even $d\ge4$, apply the preceding Hasse theorem to the
conditional reduction in
Theorem~\ref{thm:ss-even-full-group-if-unit-Hasse-irreducible}.
Degree two is the separately proved exact certificate
Theorem~\ref{thm:ss-product-degree-two-full-group}.
The natural translation module and general linear block of its
full stabilizer show that its common covector kernel is also
the unique proper nonzero invariant subspace.
All geometric and arithmetic identifications, and the passage
to ambient hypersurfaces, are those already checked in these
theorems. No claim for an arbitrary ambient variety follows
from this fixed-surface classification.
\end{proof}


% /root/monodromy, 2026-11-19. Root, reduction-descent and obstructions-crystals complete independent proof/source/code audits passed.
\subsection{All degrees for a generic mixed elliptic product}
\label{subsec:generic-mixed-product-groups-candidate}
\status{Proved and independently checked. The parent, reduction/descent and obstructions/crystals agents audited the complete proof and independently reran the coefficient reconstruction. The ordinary elliptic modulus is assumed transcendental over the prime field; every fixed algebraic ordinary modulus remains outside this theorem.}

Let $E_0:y^2+y=x^3$ and let $E$ be an ordinary elliptic curve
in characteristic two. Put
\[
 X=E\times E_0,\qquad
 L=\mathcal O_E(3O)\boxtimes\mathcal O_{E_0}(3O),\qquad
 b_d=9d^2.
\]
The rank and entire invariant-covector results already proved give
relative ranks $b_d-1$ in odd degree and $b_d-2$ in even degree,
with respectively zero and two geometric invariant covectors.
Here and below a sharp curve means a smooth curve attaining the
indicated maximal rank; on a mixed product its full curve defect
$g-f$ is one in odd degree and two in even degree.

\begin{lemma}[Generic odd images and odd Hasse divisors]
\label{lem:generic-mixed-odd-groups-candidate}
\status{Proved under stated hypotheses}
If $j(E)$ is transcendental over $\mathbf F_2$, the geometric
relative degree-$d$ image is $\GL_{b_d-1}(\mathbf F_2)$ for every
odd $d\ge1$. For odd $d\ge3$, the determinant on the relative
coherent bundle is a geometrically irreducible reduced homogeneous
polynomial of degree $b_d$, with generic stable loss one on its
divisor.
\end{lemma}

\begin{proof}
Consider the one-parameter family
\[
 E_t:\quad y^2+txy+(1+t)y=x^3+1,
 \qquad \Delta(t)=1+t^3+t^5,
 \qquad j(t)=\frac{t^{12}}{1+t^3+t^5}.
\]
Work over $B=\Spec\mathbf F_2[t,\Delta(t)^{-1}]$.
Its fibres at $t=0,1$ are smooth; the first is geometrically
isomorphic to $E_0$ and the second is
$y^2+xy=x^3+1$. The displayed nonconstant $j$ makes the generic
ordinary modulus transcendental. The standard pole-order bases
of $H^0(E_t,\mathcal O(3dO))$ trivialize its coefficient projective
bundle over $B$.

For odd $d$, the relative coherent bundle on each coefficient
projective space is $\Omega^1$, and it is generically Frobenius
bijective by Theorem~\ref{thm:DE-A3-ppav-ranks}. This construction
and its rank-$b_d-1$ Lang frame cover extend over the open of
Frobenius invertibility in the relative projective bundle over
$B$. The geometric $t=0$ frame fibre is connected, by
Theorem~\ref{thm:ss-all-odd-full-groups}. The geometric generic
frame fibre is therefore connected as well. Here no properness
of the invertibility open is assumed: normalize $B$ in a finite
separable extension defining the hypothetical generic components.
The normalization is finite, since $B$ is excellent. Its pulled-back
coefficient open is normal and integral. Components of a finite
etale cover of this open are disjoint, open and closed, and each
surjects onto the open. All would consequently meet a geometric
fibre over $t=0$, contradicting that fibre's connectedness.
This is precisely the nonproper argument of
Lemma~\ref{lem:ordinary-pair-GL8-open}.
Connectedness of the full frame cover proves the asserted generic
general linear group.

The determinant on $\Omega^1$ has degree $b_d$ and varies
algebraically with $t$. For $d\ge3$ its $t=0$ polynomial is
geometrically irreducible and reduced, by
Theorem~\ref{thm:ss-all-odd-Hasse}. Geometric reducibility of
a nonzero homogeneous form of fixed degree is a closed condition:
for each possible pair of positive factor degrees it is the image
of the proper multiplication map between the corresponding
projective spaces of forms. Their finite union includes every
nonreduced form. Thus the geometric generic polynomial is
irreducible and reduced. A point of stable loss one on the
special divisor belongs to the open where an appropriate
$(b_d-2)$-minor of a sufficiently high Frobenius iterate is
nonzero. The total determinant divisor is flat over a
neighbourhood of that point: its defining form is nonzero on
every fibre. Therefore this nonempty open meets the generic
divisor, giving generic stable loss one. Equivalently one may
use the first nonzero minor and the absence of a vertical
component of this flat hypersurface.

Finally any elliptic curve with transcendental $j$ becomes
isomorphic over an algebraic closure to a base change of this
geometric generic curve, since $j(t)$ is nonconstant.
Finite etale connectedness is preserved under algebraically
closed field extension. This proves each assertion for that
curve. The same generic point works for all $d$; no intersection
of countably many closed-point open conditions is being taken.
\end{proof}

\begin{lemma}[The mixed first-image sheaf and a natural block]
\label{lem:generic-mixed-image-block-candidate}
\status{Proved under stated hypotheses}
Work over an algebraically closed field, and suppose the odd
degree-one image on $X$ is $\GL_8(\mathbf F_2)$.
For even $d\ge2$ put $N=L^d$, $W=H^0(N)$, $b=b_d$,
$P=\mathbf P(W)$ and $\mathcal B=W^*\otimes\mathcal O_P(-1)$.
Normalize invariant derivations $e,n$ so that
$e^2=e$, $n^2=0$ and $en=ne$, with $e$ from $E$ and $n$ from
$E_0$. They lift canonically to $N$, since $N$ is a Frobenius
pullback. Set
\[
 \mathcal I=\ker\bigl(\mathcal B\xrightarrow{(s,es)}\mathcal O_P^2\bigr),
 \qquad
 \mathcal K=\ker\bigl(\mathcal B\xrightarrow{(s,es,ns,ens)}\mathcal O_P^4\bigr).
\]
There is a closed subset $Z\subset P$, of codimension $b/4$,
outside which these are vector bundles and
\[
 0\longrightarrow\mathcal K\longrightarrow\mathcal I
       \xrightarrow{(ns,ens)}\mathcal O_{P\setminus Z}^2
       \longrightarrow0
\]
is Frobenius-equivariant with constant unit quotient. Their
reflexive extensions satisfy
\[
 \rank\mathcal I=b-2,\quad\rank\mathcal K=b-4,\quad
 \det\mathcal I=\det\mathcal K=\mathcal O_P(-b),\quad
 H^0(P,\mathcal I)=H^0(P,\mathcal K)=0.
\]
Frobenius on $\mathcal I$ is generically bijective, its Lang
system is the required relative unit system, and its determinant
is a nonzero degree-$b$ form $h_d^{\mathrm u}$.
The geometric group contains an actual natural
$\GL_8(\mathbf F_2)$ block fixing a complementary subspace.
Its natural eight-space lies in the common covector kernel.
\end{lemma}

\begin{proof}
On the ordinary factor,
$F_*\mathcal O_E=\mathcal O\oplus B_E$ with $B_E$ the nontrivial
geometric two-torsion line bundle. The two summands are the
weight-zero and weight-one spaces for $e$. On the supersingular
factor the positive twist of the nonsplit extension
$0\to\mathcal O\to F_*\mathcal O_{E_0}\to\mathcal O\to0$
has vanishing first cohomology; hence its global sections are
free over $k[n]/n^2$. Consequently
\[
 W=W_0\oplus W_1,\qquad e|W_i=i,
\]
and each $W_i$ is free of rank $b/4$ over $k[n]/n^2$.
Write $s=s_0+s_1$. The locus
\[
 Z=\{ns_0=0\}\cup\{ns_1=0\}
\]
is the union of two linear subspaces of codimension $b/4$.
Outside it, $s_0,ns_0$ and $s_1,ns_1$ are independent in their
respective weight spaces. This proves independence of
$s,es,ns,ens$ and all the claimed bundle ranks and surjectivity.
The kernel sheaves are reflexive by the depth lemma, and
codimension at least two permits their determinant calculation
from the displayed maps. Their inclusions into $\mathcal B$
give the vanishing of global sections.

The mixed row of \eqref{eq:nonordinary-Cartier-table} says
\[
 M(s)s=M(s)es=0,\qquad
 M(s)ns=(ns)^{[2]},\qquad M(s)ens=(ens)^{[2]},
\]
in the root-removed Cartier convention. Thus full Frobenius
maps into $\mathcal I$ and the two quotient covectors are
Frobenius-fixed. The exact generic relative stable rank is
$b-2$, so its restriction to $\mathcal I$ is bijective on a
dense open. The full coherent model has the same unit part
as the relative system: the ambient top coherent cohomology
is Frobenius-nilpotent. Its determinant extends across $Z$
as a nonzero section of $\mathcal O_P(b)$, and agrees with
the determinant on $\mathcal K$ because the quotient is unit.

For the block, fix a smooth sharp odd-degree curve
$D\in|L^{d-1}|$ and vary a sharp $C\in|L|$ meeting it
transversely. The full curve defect of their nodal union is
$1+1=2$, precisely the generic even defect. At an ordinary
intersection node of its equation $s$, the first derivatives
vanish and the mixed derivative $ens$ is nonzero. Thus both
$ns_1=ens$ and $ns_0=ns+ens$ are nonzero there, placing the
union outside $Z$. The boundary Lang system therefore embeds
with the full generic rank, as in
Theorem~\ref{thm:ss-uniform-GL8-block}.

There are $q=18(d-1)$ intersection points. Their monodromy
is $S_q$: the one- and two-point incidences on the smooth
integral curve $D$ are projective bundles, using respectively
one and two independent values. The double-zero incidence
has a general unique double zero with independent value and
first-jet deformation parameters; its quadratic local algebra
$z^2+b'z+c$ supplies a single transposition, including in
characteristic two. The joint curve-jet rank computation for
two plane cubics excludes a second double zero, exactly as in
Lemma~\ref{lem:ss-intersection-Hasse-degree}. Its hypotheses
also hold for an ordinary plane cubic: its quadratic partials
are independent, so the tangents through a fixed point form a
finite set. Removing the sharp-source complement does not
change the image of a finite cover of a normal integral base.

The normalization filtration of the boundary relative unit
system $V$ has graph term $T$ and a quotient
\[
 V/T=(H^1_{\et}(C,\mathbf F_2)\oplus
          H^1_{\et}(D,\mathbf F_2))/H^1_{\et}(X,\mathbf F_2).
\]
The ambient diagonal line maps injectively to either curve.
Projection to the relative degree-one system on $C$ has
kernel an extension $T_0$ of the graph term by a constant
space. The image on $T_0$ has a normal two-group with quotient
a subgroup of $S_q$; on its faithful graph quotient it is the
full $S_q$. Hence its only possible nonabelian simple factor
is $A_q$. It has no factor $\GL_8(\mathbf F_2)$: the latter
is simple, has no prime factor eleven, whereas $A_q$ does
since $q\ge18$.

Goursat's lemma therefore gives a subgroup still surjecting
onto $\GL_8$ and acting trivially on $T_0$. In the resulting
block parabolic its translation kernel is a sum of natural-dual
copies. The elementary vanishing
$H^1(\GL_8(\mathbf F_2),(\mathbf F_2^8)^*)=0$
splits its quotient complement after conjugation; removing
the remaining translation columns gives an actual natural
$\GL_8$ block fixing a complement, as proved in the cited
block theorem. The two invariant covectors vanish on this
nontrivial simple eight-dimensional block, so it lies in their
common kernel.
\end{proof}

\begin{lemma}[The mixed even Hasse divisor in degrees at least four]
\label{lem:generic-mixed-even-Hasse-candidate}
\status{Proved under stated hypotheses}
Suppose $j(E)$ is transcendental over $\mathbf F_2$.
For every even $d\ge4$, the polynomial $h_d^{\mathrm u}$ of
the preceding lemma is geometrically irreducible and reduced,
and its generic divisor has stable loss exactly one.
\end{lemma}

\begin{proof}
We spell out the changes to the two boundary degree comparison
proof, so no assertion for an arbitrary fixed ordinary modulus
is implicit. Fix odd $e$ and a smooth sharp $D\in|L^e|$.
For odd $a$, the finite intersection algebra on $P_a=|L^a|$
has determinant line
\[
 \det\pi_*\mathcal O_{C\cap D}
       =\mathcal O_{P_a}(-18ae-9e^2).
\]
The cohomology sequence, removal of the codimension-two
containment locus, and characteristic-two local matrix
$\left(\begin{smallmatrix}1&c\\0&b'\end{smallmatrix}\right)$
are unchanged from Lemma~\ref{lem:ss-intersection-Hasse-degree}.
Thus the reduced irreducible point-intersection divisor
$R_D$ has degree $18ae+9e^2$. For $a\ge3$ the two factor
line bundles separate joint first jets and second jets.
For $a=1$, the product-cubic calculation still gives full
joint curve jets away from finitely many points, rank at
least three at those points, and full second jets at a
nonflex point. These statements follow by tensoring the two
elliptic evaluation spaces; the ordinary cubic is no more
exceptional than the supersingular one. The generic point
of $R_D$ therefore has precisely one contact of order two.
The marked incidence can be generically inseparable; the
local finite-algebra determinant proves the reduced equation
on the coefficient space itself.

For $a\ge3$ the preceding odd lemma gives an irreducible
reduced source Hasse polynomial $h_a$ of degree $9a^2$.
The singular-curve discriminant is irreducible of degree
$c_2(J^1L^a)=54a^2$, by the same first/second-jet calculation
and a generic unique ordinary node. The three degrees
\[
 9a^2,\quad54a^2,\quad18ae+9e^2
\]
are pairwise distinct: either equality involving the last
would require $(a+e)^2=2a^2$ or $7a^2$.
Thus the generic curve on $h_a$ is smooth and transverse
to $D$, and the generic curve on $R_D$ is smooth and sharp.
The respective union defects are at least $2+1=3$ and
exactly $1+1+1=3$. These exceed the generic even defect two.
There are also transverse nodes, which place these unions
outside $Z$ by the mixed-Hessian calculation above.
Consequently both $h_a$ and $r_D$ divide the nonzero restriction
of $h_{a+e}^{\mathrm u}$ under multiplication by $s_D$.
Their degrees add to $9(a+e)^2$, so
\[
                  j_D^*h_{a+e}^{\mathrm u}=c_D h_a r_D,
                  \qquad c_D\ne0.
\]
This uses stable-rank divisibility and total degree, not an
unproved exactness assertion for coherent first-image lattices.

For $a=1,e=3$, the restriction is still nonzero and divisible
by the irreducible $r_D$ of degree $135$. Its generic source
curve is sharp because its Hasse polynomial has degree only
nine, and smooth because the full singular-curve discriminant
has degree $54$. Hence this divisibility needs no irreducibility
of the degree-one source Hasse divisor.

The boundary pairs $(a,e)=(d-1,1)$ and $(d-3,3)$ for $d\ge6$
would force any two global components to have both unordered
degree pairs
\[
 \{9(d-1)^2,18d-9\},\qquad
 \{9(d-3)^2,54d-81\}.
\]
They cannot agree: the crossed equality would give
$(d-4)^2=6$. The first squarefree restriction already forces
the global form to be reduced with at most two components.
For $d=4$, their hypothetical degrees would be $81,63$,
contradicted by the irreducible degree-$135$ restriction
factor just obtained. This proves reduced irreducibility in
every claimed degree. On the generic $R_D$ boundary with
$a=d-1,e=1$ the stable loss is exactly one and the point lies
outside $Z$. Irreducibility and lower semicontinuity of the
rank of a high Frobenius iterate give generic loss one on
the global determinant divisor.
\end{proof}

\begin{lemma}[Two exact mixed degree-two coefficient points]
\label{lem:mixed-degree-two-certificate-candidate}
\status{Proved under stated hypotheses}
On the coefficient space for
\[
 E_o:y^2+xy=x^3+1,\qquad E_0:y^2+y=x^3,
 \qquad N=\mathcal O(6O)\boxtimes\mathcal O(6O)
\]
over $\mathbf F_2$, there are two points of the rank-$34$
unit-image open with the following properties. At the first,
the Frobenius on the common kernel of the two invariant
covectors has irreducible nonreciprocal characteristic
polynomial of degree $32$. At the second, every fixed vector
lies in that common kernel.
\end{lemma}

\begin{proof}
Order the factor basis as $(1,x,x^2,x^3,y,xy)$ and its tensor
basis by $6i+j$, with $0\le i,j<6$. A nonnegative integer
encodes its coefficient vector by binary digits in this order.
The section integers are
\[
                 s_A=34346068584,\qquad s_B=68001692233.
\]
Let $C_s$ be the root-removed Cartier multiplication matrix.
The exact independent reconstruction gives
$\rank C_s=\rank C_s^2=34$ at both points. The two fixed
covectors $ns,ens$ have integer vectors
\[
 \begin{array}{c|cc}
  &ns&ens\\\hline
 A&1494603548&420583168\\
 B&26911823620&74729216.
 \end{array}
\]
The full degree-$36$ characteristic polynomial at $A$ is
$T^2(T+1)^2 f(T)$, where
\[
 \begin{split}
 f(T)={}&T^{32}+T^{30}+T^{29}+T^{28}+T^{27}+T^{25}
       +T^{24}+T^{23}+T^{22}\\
       &+T^{20}+T^{18}+T^{15}+T^{14}+T^{11}+T^9+T^8
       +T^5+T+1.
 \end{split}
\]
Its binary polynomial integer is $6372510499$.
The certificate checks $T^{2^{32}}=T\pmod f$ and
$\gcd(T^{2^{16}}-T,f)=1$, proving irreducibility, and checks
that its reciprocal differs from $f$. The prime $65537$
is verified by trial division and has $\operatorname{ord}_{65537}(2)=32$.

At $B$ the literal vectors
\[
                         \mu_0=8601123076,\qquad
                         \mu_1=41642123014
\]
satisfy $(C_{s_B}-I)\mu_0=ns_B$ and
$(C_{s_B}-I)\mu_1=ens_B$.
The coherent Frobenius matrix is $C_s^t$, so these identities
annihilate both covectors on every Frobenius-fixed vector.
Nilpotent directions contain no nonzero fixed vectors;
therefore the assertion holds on the required unit image.
The two covectors are independent, and the rank checks put
both points outside $Z$ and on the unit-image open.
Their curves need not be smooth: restriction from this normal
coefficient open to its dense smooth-member open preserves
the image of every finite cover.

The discovery program and its original output are retained as
\path{agents/monodromy/code/mixed_degree_two_group_discovery.py}
and \path{agents/monodromy/outputs/mixed_degree_two_group_discovery.json}
(306 trials, 51 exact characteristic polynomials,
$4.062721072004933$ seconds).
The independent program
\path{agents/monodromy/code/verify_mixed_degree_two_group.py}
imports no discovery code or stored matrix. It multiplies in
the two coordinate rings and reconstructs Cartier by inverting
the square map on the six-dimensional section space, using
\[
 C_o(f\,dx/x)^2=(1+e)f,\qquad C_0(f\,dx)^2=n f.
\]
It then forms both tensor matrices, computes their integer
characteristic polynomials and reduces modulo two, and checks
all the displayed ranks and vector identities. Its retained
output \path{agents/monodromy/outputs/verify_mixed_degree_two_group.json}
is PASS in $0.4645676829968579$ seconds. These are finite
coefficient checks, not a claim that a single finite-field
fibre has the desired geometric group.
\end{proof}

\begin{lemma}[The geometric generic mixed degree-two image]
\label{lem:generic-mixed-degree-two-group-candidate}
\status{Proved under stated hypotheses}
If $j(E)$ is transcendental over $\mathbf F_2$, the degree-two
geometric image is the entire pointwise stabilizer of its
two independent invariant covectors:
\[
                 \operatorname{Hom}(\mathbf F_2^2,\mathbf F_2^{32})
                          \rtimes\GL_{32}(\mathbf F_2).
\]
\end{lemma}

\begin{proof}
Use the family $E_t$ from the first lemma, now restricted to
$t\ne0$, and the maximal-rank even first-image open over it.
It is a normal integral open of a projective bundle. Let $G$
be its full finite arithmetic Lang image, and $H$ the image
of its geometric generic coefficient fibre.
Then $H$ is normal in $G$. Indeed every connected finite
etale cover of the normal total open is integral, and its
generic fibre over the integral curve is a localization of
that integral scheme, hence integral. Thus the generic-fibre
fundamental group surjects onto the total fundamental group.
The geometric generic-fibre fundamental group is normal in
the former, by the exact fundamental-group sequence over
the function field. Passing to finite images proves the claim.

The two covectors are defined over this ordinary base and
fixed arithmetically, including over its imperfect function
field. Explicitly the invariant derivation $D$ on $E_t$ is
\[
 D(x)=tx+1+t,\qquad D(y)=x^2+ty,\qquad D^2=tD.
\]
On $t\ne0$ put $e=D/t$, so $e^2=e$, and use the standard
nilpotent derivation $n=\partial_{x_0}$ on $E_0$.
The invariant volume
\[
             \left(\frac{t\,dx}{tx+1+t}\right)\wedge dx_0
\]
has contraction one with $(e,n)$. The direct local
coefficient-extraction identities of
\eqref{eq:nonordinary-Cartier-table} therefore give the two
fixed covectors $ns,ens$ over this base itself; this does not
require a perfect-field descent argument over $\mathbf F_2(t)$.
At $t=1$ the derivations and volume specialize to exactly
those in the coefficient certificate.
Consequently $G$ fixes the quotient $V/K\simeq\mathbf F_2^2$
pointwise, where $\dim K=32$.
The generic-fibre block lemma gives an actual natural
$\GL_8$ block in $H$, fixing a complement and lying in $K$.
Let $N$ be the normal closure in $G$ of its transvections.
Since $H\triangleleft G$, one has $N\subset H$.
On $K$, the group $N_K$ is nontrivial, normal in $G_K$,
and generated by transvections.

The first certified coefficient point is a rational point of
the total open. Its arithmetic Frobenius, or its inverse if
the opposite convention is used, has the irreducible
nonreciprocal polynomial $f$ or its reciprocal on $K$.
Lemma~\ref{lem:product-frobenius-group-test}, in dimension
$32$, now gives $N_K=G_K=\GL(K)$. In particular $H_K$ is full.

The translation kernel of $G\to\GL(K)$ is a submodule of
$\operatorname{Hom}(\mathbf F_2^2,K)=K\otimes(\mathbf F_2^2)^*$.
Because the natural module has endomorphism ring $\mathbf F_2$,
it has the form $K\otimes U$ for a subspace
$U\subset(\mathbf F_2^2)^*$. The vanishing of the first
cohomology of $\GL(K)$ on sums of its natural module makes
$G$ conjugate to $(K\otimes U)\rtimes\GL(K)$.
If $U$ were proper, this group would have a nonzero fixed
vector outside $K$. The second certified Frobenius element
has no fixed vector outside $K$, excluding that possibility.
Thus $G$ is the entire pointwise stabilizer.

Finally $H$ is normal in this entire stabilizer and surjects
onto $\GL(K)$. Its commutators with all translations include
every element of $K\otimes(\mathbf F_2^2)^*$: for each quotient
coordinate the vectors $(A-I)v$, with $A\in\GL(K)$, span $K$.
Normality puts these commutators in $H$. Hence $H$ contains
all translations and is the entire stabilizer as well.
This step is why the two finite-field coefficient points
give a geometric generic-fibre conclusion, rather than only
a total-family arithmetic conclusion.
\end{proof}

\begin{theorem}[All generic mixed-product images]
\label{thm:generic-mixed-product-all-groups-candidate}
\status{Proved under stated hypotheses}
Let $k$ be a perfect field of characteristic two, let $E/k$
be ordinary with $j(E)$ transcendental over $\mathbf F_2$,
and let $X=E\times E_0$ have the Segre embedding defined
by $L$. For every $d\ge1$, the entire geometric and arithmetic
relative Hasse--Witt images are
\[
 G_{\mathrm{geom},d}=G_{\mathrm{arith},d}=
 \begin{cases}
 \GL_{9d^2-1}(\mathbf F_2),&d\text{ odd},\\[2pt]
 \operatorname{Hom}(\mathbf F_2^2,\mathbf F_2^{9d^2-4})
       \rtimes\GL_{9d^2-4}(\mathbf F_2),&d\text{ even}.
 \end{cases}
\]
The even image is the full pointwise stabilizer of its entire
two-dimensional invariant-covector space in rank $9d^2-2$.
Its common covector kernel is irreducible and is contained in
every nonzero invariant subspace; the quotient is the trivial
two-dimensional representation. In particular the odd systems
are geometrically irreducible and the even systems are not.
\end{theorem}

\begin{proof}
Odd degrees and degree two are the preceding lemmas. For
even $d\ge4$, use the reduced irreducible unit-image Hasse
divisor. The determinant-allocation argument of
Lemma~\ref{lem:irreducible-hasse-quotients-trivial} applies
to the negative reflexive sheaf $\mathcal I$ exactly as in
Theorem~\ref{thm:ss-even-full-group-if-unit-Hasse-irreducible}.
In codimension one, determinants of a saturated subobject
and its torsion-free quotient add to the reduced irreducible
determinant divisor. Assigning no divisor to a nonzero
subobject would make its reflexive extension a trivial
Frobenius-unit bundle, contrary to $H^0(\mathcal I)=0$.
Thus every proper nonzero invariant subspace has trivial
quotient. The entire covector space has dimension two, by
Theorem~\ref{thm:nonordinary-ppav-all-covectors}; hence every
such subspace contains its common kernel $K$.
It follows that $K$ is irreducible and that $V$ has no
nonzero fixed vector.

Generic stable loss one gives generation by transvections,
using the Lang finite part and purity argument of
Lemma~\ref{lem:char2-defect-one-inertia}. Removing $Z$ causes
no exception, because its codimension is at least nine.
The restriction to $K$ is thus irreducible and generated
by transvections. Its actual natural $\GL_8$ block excludes
a nondegenerate invariant alternating form: the eight-space
has no invariant form and pairs trivially with its fixed
complement. The exact alternatives of
\cite[Theorem~3.6, printed p.~22]{KantorSeress2001}
therefore give the full $\GL(K)$ image.

As in the degree-two lemma, the translation kernel has the
form $K\otimes U$, and the vanishing of natural-module
$H^1$ gives a linear complement modulo that kernel.
If $U$ were proper, a nonzero quotient direction would
give a fixed vector of $V$. There are no such vectors,
so $U=(\mathbf F_2^2)^*$ and the entire stabilizer occurs.
Both invariant covectors descend over every perfect base
field by the arithmetic part of the cited covector theorem.
Thus the arithmetic image is bounded by the same stabilizer
and contains the geometric image, proving equality.

The plane cubic coordinate rings are generated in degree
one, and the diagonal tensor product has the same property.
Ambient degree-$d$ forms therefore surject onto $H^0(X,L^d)$
for every $d$. The corresponding linear projection has
geometrically integral fibres, and restriction to the dense
smooth maximal-rank open preserves the finite image.
This proves the statement for the requested ambient
hypersurface sections as well as the complete linear systems.
\end{proof}

\begin{remark}[The exact remaining modulus restriction]
\label{rem:generic-mixed-modulus-gate}
\status{Proved under stated hypotheses}
The theorem treats every degree at the geometric generic
ordinary modulus, and hence every ordinary elliptic factor
with transcendental $j$. It does not prove these images for
every fixed ordinary elliptic curve over an algebraic extension
of $\mathbf F_2$. Specialization gives the generic conclusion
used above; reversing that direction would be invalid.
The independent exact-rank theorem for every fixed ordinary
factor does not remove this monodromy restriction.
\end{remark}


% /root/monodromy, 19 November 2026. Complete root and reduction/descent audits passed.
\subsection{Uniform degree-one image for every fixed mixed product}
\label{subsec:mixed-uniform-degree-one-inertia-candidate}
\status{Proved and independently checked. The root and reduction/descent
agents audited all six complete assertions, the source dependencies and
the full-image assembly. Both independently read and reran the universal
Pfaffian and stable-loss certificates; the rank-compatibility certificate
was independently read and rerun by the author and reduction/descent agent.
The result covers every ordinary modulus.}

Work over an algebraically closed field $k$ of characteristic two. For
$a\in k^\times$ put
\[
 E_a:y^2+xy=x^3+a,\qquad E_0:z^2+z=w^3,
 \qquad X_a=E_a\times E_0,
 \qquad L=\mathcal O_{E_a}(3O)\boxtimes\mathcal O_{E_0}(3O).
\]
Write $W=H^0(X_a,L)$, $P=\mathbf P(W)$ and $s=(s_{ij})$ in
the tensor basis $(1,x,y)\otimes(1,w,z)$; put $D(s)=\det(s_{ij})$.
The chosen invariant volume is $(dx/x)\wedge dw$.

\begin{proposition}[Two cubic components at every ordinary modulus]
\label{prop:mixed-uniform-degree-one-Hasse-candidate}
\status{Proved under stated hypotheses}
Let
\[
 R_o=\begin{pmatrix}1&0&a\\0&0&a\\0&a&a\end{pmatrix},
 \qquad J=(\delta_{i,2-j})_{0\le i,j\le2},
 \qquad R=R_o\otimes J.
\]
The relative coherent Frobenius bundle is $\Omega_P^1$, and its
determinant is, up to a nonzero constant,
\[
                   h(s)=D(s)^2Q(s),\qquad Q(s)=s^tRs^{[2]}.
\]
The cubics $D,Q$ are geometrically irreducible and distinct, and
$Q$ is smooth. The generic relative stable rank is eight. At the
geometric generic point of $Q$, it is seven, and the linearized
relative Frobenius has a one-dimensional kernel.
\end{proposition}

\begin{proof}
The root-removed Cartier multiplication matrices on the two factors are
\[
 M_o(u)=\begin{pmatrix}
 u_0&a u_2&a u_1\\u_2&u_1&u_0\\0&u_2&u_1+u_2
 \end{pmatrix},\qquad
 M_0(u)=\begin{pmatrix}u_1&u_0&0\\u_2&0&u_0\\0&u_2&u_1\end{pmatrix}.
\]
For the first factor, the underlying Cartier identities are
$C(1)=1$, $C(x)=0$, $C(y)=x$, and $C(xy)=y+\sqrt a$,
together with $C(f^2)=f$ for an invariant Cartier-fixed
differential. The second-factor identities were proved in
Proposition~\ref{prop:ss-product-relative-Hasse}.
If $M_{o,i},M_{0,j}$ are the coefficient matrices, put
$M(s)=\sum s_{ij}M_{o,i}\otimes M_{0,j}$.
The symmetric multiplication identity gives $M(s)s=0$.
More explicitly, $A(s)=RM(s)$ is alternating, and its coefficient
tensor is symmetric in all three arguments. One may check the latter
directly from
\[
 R_oM_o(u)=\begin{pmatrix}
 u_0&0&a u_2\\
 0&a u_2&a(u_1+u_2)\\
 a u_2&a(u_1+u_2)&a(u_0+u_1+u_2)
 \end{pmatrix}
\]
and the alternating volume tensor on the second factor.
The nine size-eight Pfaffians satisfy
\[
                    \operatorname{Pf}_8(A(s))=a^3D(s)s,
                    \qquad \det R=a^6.
\]
These are polynomial identities over $\mathbf F_2[a,s_0,\ldots,s_8]$:
the retained exact certificate expands every Pfaffian recursively and
checks all nine identities. The alternating adjugate identity consequently
gives
\[
                  \operatorname{adj}M(s)=D(s)^2s\,s^tR.
\]
The universal hypersurface sequence is the Euler sequence, since
$H^1(X_a,L^{-1})=0$ and $h^2(X_a,\mathcal O_{X_a})=1$.
Frobenius on the last one-dimensional quotient is zero, because the
surface is nonordinary. Thus the relative map is the restriction of
$M(s)^t$ from $(s^{[2]})^\perp$ to $s^\perp$. Its cofactor formula
is exactly $D(s)^2s^tRs^{[2]}$, as in the preceding supersingular
calculation. Its degree is nine, agreeing with
$\det\Omega_P^1=\mathcal O_P(-9)$.

The determinant cubic $D$ is irreducible by its matrix-factorization
parametrization. Since
$\nabla Q=Rs^{[2]}$ and $R$ is invertible, $Q$ is smooth and hence
irreducible. The coefficient vector with $s_{01}=1$ and all other
entries zero has $D=0,Q=1$, so the two factors are distinct. Their
complement is a nonempty full-rank open, proving generic rank eight.

We verify stable loss one on $Q$ uniformly in $a$, not merely at a
single ordinary modulus. Let $z$ be an auxiliary nonzero parameter
and take, in lexicographic order,
\[
             s=(z^3/a,z,0,0,1,0,0,0,1).
\]
Then $D=z^3/a\ne0$ and $Q=0$. The vector
\[
 u=\left(0,\frac{z^3}{a},\frac a z,0,
          \frac{a+z^3}{z^3},\frac{a^2+z^5}{az^2},0,0,\frac a{z^3}\right)^t
\]
satisfies $M(s)u=s^{[2]}$, and
\[
             s^tR u^{[2]}=\frac{a^4+a^3z^4+z^{10}}{a^2z^3}.
\]
Both identities are also checked over $\mathbf F_2(a,z)$ by the
retained program. The last expression is not identically zero as
a function of $z$ for any $a\ne0$.
On $Q\setminus D$ the full matrix has kernel $ks$, and the
relative Cartier kernel is the line generated by $u$ modulo $ks$.
Its squared representative does not lie in $\operatorname{im}M(s)$
precisely when the displayed pairing is nonzero, since $s^tR$
spans the left kernel and $s^{[2]}$ already lies in the image.
Thus the first relative image has rank seven and is already
Frobenius-stable at such a point. Openness of the corresponding
iterate minor proves the generic assertion on the irreducible $Q$.
\end{proof}

\begin{lemma}[Full elementary inertia from a moving kernel line]
\label{lem:moving-kernel-full-root-inertia-candidate}
\status{Proved under stated hypotheses}
Let $Y$ be a smooth variety over a perfect field of characteristic two,
and let $D\subset Y$ be an integral divisor. Let a vector bundle
$\mathcal E$ of rank $r$ carry regular semilinear Frobenius. Suppose
its geometric generic rank is $r$, its stable rank at the generic point
of $D$ is $r-1$, its linearized kernel there is a line, and its
determinant has order one along $D$.
Suppose that on a neighbourhood of the generic point of $D$ there is
an inclusion $\mathcal E\subset V\otimes\mathcal O_Y$ as a subbundle of a constant
bundle, and the linearized kernel line in $F^*\mathcal E$ defines
a rational map
\[
                      \kappa:D\dashrightarrow\mathbf P(V^{(2)})
\]
whose differential has rank $r-1$ at the generic point.
Then the divisorial inertia image of the generic Lang local system
contains the entire elementary group
\[
           \operatorname{Hom}(\mathbf F_2,\mathbf F_2^{r-1}),
\]
acting trivially on the regular $(r-1)$-dimensional unit subspace
and on the one-dimensional quotient. In particular it contains
every transvection with that fixed hyperplane.
\end{lemma}

\begin{proof}
Pass to the strict henselization at the generic point of $D$.
Write $R$ for this discrete valuation ring, $t$ for a uniformizer,
and $K$ for its separably closed residue field. The finite part of
the Lang scheme supplies $r-1$ regular Frobenius-fixed sections
$w_1,\ldots,w_{r-1}$, with independent special values spanning
the stable part. This is the finite-part construction already
proved using Stacks~\cite{Stacks0BTN,Stacks09ZS}. Complete them
to a frame by $e$. In that frame
\[
                   F=\begin{pmatrix}I_{r-1}&b\\0&a\end{pmatrix},
                  \qquad a=t c,\qquad c\in R^\times.
\]
No perfectness of $K$ is assumed. In particular the constant
column $b_0=b\bmod t$ need not be removable by a regular change
of frame. The linearized kernel is the line generated by
\[
             e^{[2]}+\sum_{i=1}^{r-1}b_{0,i}w_i^{[2]}.
\]
Use the given inclusion in the constant bundle. In Kähler
differentials over the perfect ground field, every entry of
$e^{[2]}$ and $w_i^{[2]}$ has zero differential. The differential
rank of the projective kernel line is therefore the rank of
$db_{0,1},\ldots,db_{0,r-1}$. The assumed rank $r-1$ makes
these differentials linearly independent over $K$. This calculation
is unchanged by the separable residue-field extension: separable
algebraic extension preserves the differential rank.

The nonzero fixed vector of the quotient has coordinate $a^{-1}$.
Lifting it requires the independent Artin--Schreier equations
\[
                      y_i^2+y_i=b_i/a^2.
\]
For any nonzero $(\epsilon_i)\in\mathbf F_2^{r-1}$, the leading
coefficient of $\sum\epsilon_i b_i/a^2$ at order $t^{-2}$ is
$c_0^{-2}\sum\epsilon_i b_{0,i}$. Its differential is nonzero,
so it is not a square in $K$. If this sum were $y^2+y$ in the
fraction field of $R$, its valuation would force $v_t(y)=-1$,
whose leading coefficient would make that coefficient a square.
This is impossible. The same argument applies in the henselian
field itself and does not require completion.

The $r-1$ classes are consequently linearly independent in
$\operatorname{Frac}(R)/\{y^2+y\}$. Successive quadratic
Artin--Schreier extensions therefore have compositum Galois group
$\mathbf F_2^{r-1}$. This standard elementary assertion follows
also by taking the annihilator of the image in
$\mathbf F_2^{r-1}$: a proper image would give exactly one of the
excluded nonzero linear combinations. Over a strictly henselian
trait this is its inertia image, including possible purely
inseparable residue extensions. Its action adds the arbitrary
constant combination $\sum\epsilon_iw_i$ to one chosen lift of
the quotient vector, and fixes all the $w_i$. This is exactly
the asserted full elementary group.
\end{proof}

\begin{corollary}[Seven independent root directions for every mixed modulus]
\label{cor:mixed-every-modulus-full-root-inertia-candidate}
\status{Proved under stated hypotheses}
For every $a\ne0$, the geometric degree-one relative image for
$X_a,L$ contains the group of order $2^7$ consisting of all
transvections with a fixed seven-dimensional hyperplane, together
with the identity.
\end{corollary}
\begin{proof}
Apply the preceding lemma to the cubic $Q$ on $P=\mathbf P^8$.
The stable rank and simple determinant order were just proved.
On an affine coefficient chart, the relative Euler bundle embeds
in the constant nine-dimensional bundle. The kernel of the
linearization on $Q\setminus D$ is the line
$kR^ts\subset (s^{[2]})^\perp$. Thus its projective map is
the restriction of the projective linear isomorphism
$[s]\mapsto[R^ts]$. Its differential on the smooth cubic has
rank seven. All hypotheses hold. The inertia is geometric,
since the construction was over an algebraically closed field.
Restricting the coefficient invertibility open to its dense
smooth-curve open preserves the finite image, as before.
\end{proof}

\begin{lemma}[The exceptional ranks of the mixed Cartier tensor]
\label{lem:mixed-degree-one-exceptional-ranks-candidate}
\status{Proved under stated hypotheses}
For the alternating matrix $A(w)$ above put
\[
       \mathcal Z=\{(s,v):v\ne0,\ A(R^{-t}v^{[2]})s=v\}.
\]
The part of $\mathcal Z$ where the $3\times3$ matrix $v$ has
rank at most two has dimension at most eight.
\end{lemma}
\begin{proof}
Write $B(u)=R_oM_o(u)$. Directly,
\[
 \det B(u)=a^2\bigl(u_0^2u_2+u_0u_1u_2+u_0u_1^2+a u_2^3\bigr).
\]
The cubic in parentheses is smooth for every $a\ne0$.
Its three partial derivatives are
$u_1^2+u_1u_2$, $u_0u_2$, and
$u_0^2+u_0u_1+a u_2^2$.
If $u_2=0$, their simultaneous vanishing forces $u_1=u_0=0$;
if $u_2\ne0$, the second forces $u_0=0$ and the third is
nonzero. Thus the projective cubic is smooth and contains no
line. In particular $B(u)$ has rank three off that cubic
and rank two on it: rank at most one would make all cofactors,
and therefore all three partial derivatives, zero.

Suppose first that $v$ has matrix rank two. Write
\[
 v=t_1\otimes q_1+t_2\otimes q_2,
 \qquad T=(t_1,t_2),\qquad Q_0=(q_1,q_2),
\]
where both displayed $3\times2$ frames have full rank.
All such representations form a $\GL_2$-torsor of dimension
four over the rank-two matrix locus. Put
\[
 u_i=R_o^{-t}t_i^{[2]},\qquad r_i=Jq_i^{[2]},\qquad
 w=R^{-t}v^{[2]}=u_1\otimes r_1+u_2\otimes r_2.
\]
The plane spanned by $u_1,u_2$ cannot lie in the smooth cubic
$\det B=0$. A change of the two-element frame therefore makes
$B(u_1)$ invertible; square roots needed for that change exist
over the algebraically closed field. Thus the open of such
frames meets every rank-two orbit.

Contract the second tensor factor with $r_1,r_2$, and complete
them to a basis. The alternating volume tensor on that factor
shows that a vector lies in $\operatorname{im}A(w)$ precisely
when its first two contractions can be written
$B(u_2)z,B(u_1)z$ for one vector $z$. The remaining contraction
is unrestricted, because $B(u_1)$ is invertible. In particular
$A(w)$ has rank six. Define
\[
 H=(h_{ij})_{1\le i,j\le2},\qquad
 h_{ij}=q_i^tJq_j^{[2]},\qquad C=B(u_2)B(u_1)^{-1}.
\]
The required compatibility is exactly
\[
                         T H_{\bullet1}=CT H_{\bullet2}.
\]
For fixed $T$, these are $k_T=\dim\langle T,CT\rangle$
independent linear equations on the four entries of $H$,
where $k_T$ is two or three.

The Gram map $Q_0\mapsto Q_0^tJQ_0^{[2]}$ on full-rank
$3\times2$ frames is smooth of relative dimension two.
Indeed its derivative is
\[
                     dH=(dQ_0)^tJQ_0^{[2]},
\]
which is surjective onto all $2\times2$ matrices: each row
of $(dQ_0)^t$ may be chosen independently and
$JQ_0^{[2]}$ has rank two. Hence the compatible $Q_0$ have
dimension at most $6-k_T$.

The exceptional condition $k_T=2$ is a proper closed condition
on the six-dimensional open of $T$-frames under consideration.
To check properness uniformly, take $T=(e_0,e_1)$. Then
\[
 u_1=(1,1,0)^t,\quad u_2=(0,a^{-1},a^{-1})^t,\qquad
 B(u_1)=\begin{pmatrix}1&0&0\\0&0&a\\0&a&0\end{pmatrix},
 \quad B(u_2)=J,
\]
and $Ce_0=e_2\notin\langle e_0,e_1\rangle$.
The locus with $k_T=3$ therefore has compatible-frame
dimension at most $6+(6-3)=9$; the locus with $k_T=2$ has
dimension at most $5+(6-2)=9$. Subtracting the four-dimensional
frame fibre gives a compatible $v$-locus of dimension at most
five. The affine space of solutions $s$ has dimension
$9-\rank A(w)=3$, giving the required bound eight.

For rank one write $v=t\otimes q$, with $t,q\ne0$, and put
$u=R_o^{-t}t^{[2]}$, $r=Jq^{[2]}$.
Then $A(w)$ is the tensor product of $B(u)$ and the
rank-two alternating volume matrix determined by $r$.
If $B(u)$ is invertible, the rank is six; the entire rank-one
$v$-locus already has dimension five, so its incidence has
dimension at most $5+(9-6)=8$.
If $B(u)$ has rank two, the rank of $A(w)$ is four.
Compatibility necessarily requires
\[
                             q^tJq^{[2]}=0,
\]
since the image on the second factor is the plane annihilating
$r$. The condition on $u$ cuts the nonzero $t$-space to
dimension two, and the displayed nonzero cubic cuts the
nonzero $q$-space to dimension two. Quotienting their
one-dimensional scaling ambiguity gives dimension at most
three for compatible $v$. Adding the solution dimension five
again gives eight. These cases exhaust all nonzero $v$ of
matrix rank at most two.
\end{proof}

\begin{lemma}[Transitivity on all nonzero relative vectors]
\label{lem:mixed-degree-one-Lang-transitive-candidate}
\status{Proved under stated hypotheses}
For every $a\ne0$, the geometric degree-one relative image is
transitive on the $255$ nonzero vectors of its eight-dimensional
natural representation. In particular it is irreducible.
\end{lemma}
\begin{proof}
Work on the affine coefficient cone where $D(s)Q(s)\ne0$.
The nonzero Lang scheme is finite etale of degree $255$ there.
In the ambient nine-dimensional frame its equations are
\[
                    v=M(s)^tv^{[2]},\qquad s^tv=0,\qquad v\ne0.
\]
Put $w=R^{-t}v^{[2]}$. Full symmetry of the coefficient tensor
gives $A(s)w=A(w)s$, so the first equation is equivalent to
$A(w)s=v$. The second equation then holds automatically,
because $s^tA(w)s=0$.

When $v$ has matrix rank three, $D(w)\ne0$ and
$\ker A(w)=kw$. The equation is solvable precisely when
\[
                         v^tR^{-t}v^{[2]}=0.
\]
This is a smooth projective cubic, since its gradient is
$R^{-t}v^{[2]}$; its punctured affine cone is irreducible.
Its full-matrix-rank open is nonempty: its monomials have a
squared variable, while every determinant monomial has three
distinct variables, so it is not the determinant cubic.
Over that open the solution space in $s$ is an affine-line
bundle, hence irreducible of dimension nine.

The previous lemma bounds the remaining incidence by dimension
eight. Every component of a finite etale cover of the integral
nine-dimensional coefficient cone has dimension nine and
surjects onto that cone. Thus none can lie in the exceptional
incidence. All meet the one irreducible affine-line incidence,
so the nonzero Lang cover is irreducible. Its degree is $255$,
which proves the asserted transitivity. Pullback to the affine
cone can only decrease the monodromy image, so the same holds
on the projective parameter open and then on its dense
smooth-curve open. A nonzero proper invariant subspace would
have a proper nonempty invariant set of nonzero vectors,
contradicting transitivity.
\end{proof}

\begin{theorem}[Full degree-one image for every mixed product]
\label{thm:mixed-every-fixed-modulus-degree-one-candidate}
\status{Proved under stated hypotheses}
Let $k$ be any perfect field of characteristic two, let $E/k$
be an ordinary elliptic curve and let $E'/k$ be a supersingular
elliptic curve. On
\[
 X=E\times E',\qquad
 L=\mathcal O_E(3O)\boxtimes\mathcal O_{E'}(3O),
\]
the relative Hasse--Witt system on the maximal-rank smooth
open of $|L|$ has rank eight and
\[
                         G_{\mathrm{geom}}=G_{\mathrm{arith}}
                              =\GL_8(\mathbf F_2).
\]
This holds for every ordinary modulus, including every algebraic
one. Its geometric representation is irreducible.
\end{theorem}
\begin{proof}
Over an algebraic closure the pair has the displayed form
$E_a\times E_0$ for some $a\ne0$. Let $G$ be its geometric
image on $V=\mathbf F_2^8$. The preceding lemma proves
irreducibility. The inertia corollary supplies, for a nonzero
$\phi\in V^*$, the entire group
\[
                  U_\phi=\{1+u\otimes\phi:\phi(u)=0\}.
\]
The span of the $G$-orbit of $\phi$ is all of $V^*$, by dual
irreducibility. Choose a basis $\phi_1,\ldots,\phi_8$ from
that orbit and let $e_1,\ldots,e_8$ be its dual basis.
Conjugation gives $U_{\phi_i}\subset G$ for every $i$.
Consequently $1+e_j\otimes\phi_i\in G$ whenever $i\ne j$.
These are all elementary matrices in that basis and generate
$\SL_8(\mathbf F_2)=\GL_8(\mathbf F_2)$. This closure uses
only elementary matrices, without a finite-group classification.
The arithmetic group contains the geometric group and is
contained in the same general linear group, proving equality.
Geometric base extension preserves connectedness of the finite
frame cover, so the conclusion holds over every stated perfect
field and for all twists of the elliptic factors.
\end{proof}

\begin{remark}[Retained computations and higher-degree scope]
\label{rem:mixed-uniform-inertia-scope-candidate}
\status{Proved under stated hypotheses}
The exact programs
\path{agents/monodromy/code/verify_mixed_degree_one_pfaffians.py}
and
\path{agents/monodromy/code/explore_mixed_universal_Q_pencil.py}
retain all universal polynomial and rational identities used above.
Earlier failed point choices are preserved separately and are not
inputs to this proof. The parent independently supplied and checked
the exceptional-rank dimension argument. Its exact auxiliary certificate is
\path{code/mixed_lang_rank_compatibility.py}, with retained output
\path{outputs/mixed_lang_rank_compatibility.json}; it checks the
cubic determinant and partials, the uniform exceptional-frame witness,
the Gram differential and its block-minor identities. These finite
identities support the full dimension proof above.
The theorem closes the fixed-modulus degree-one case. Higher degrees
at an arbitrary fixed algebraic ordinary modulus require a further
uniform argument; the already proved generic-modulus all-degree
theorem remains a separate statement.
\end{remark}


% /root/monodromy, 19 November 2026. Complete root and reduction/descent audits passed.
\subsection{A wreath and grid argument in degree two on every mixed product}
\label{subsec:mixed-degree-two-wreath-grid-candidate}
\status{Proved and independently checked.}
The root and the two companion reviewers each audited all five complete
proofs, their primary-source dependencies and the exact graph certificate,
including independent reruns. The result covers every fixed ordinary modulus.

Throughout this subsection $k$ is algebraically closed of characteristic
two, $X=E\times E_0$ with $E$ ordinary and $E_0$ supersingular, and
$L=\mathcal O_E(3O)\boxtimes\mathcal O_{E_0}(3O)$.
The degree-two relative unit system has rank $34$. Its entire invariant
covector space has dimension two; write $K$ for their common kernel,
of rank $32$. These ranks and covectors are the previously proved
ppav rank and nonordinary covector theorems.

\begin{lemma}[A labelled-transposition group]
\label{lem:labelled-transposition-wreath-candidate}
\status{Proved under stated hypotheses}
Let $G$ be a finite elementary abelian two-group and $N\ge3$.
Suppose $H\subset G^N\rtimes S_N$ acts transitively on
$G\times\{1,\ldots,N\}$ and is generated by labelled transpositions:
each generator interchanges two fibres, with a translation $c$ in
one direction and $-c$ in the other, and fixes all other fibres
pointwise. After changing the origins of the $N$ fibres,
\[
                   H=\ker(\Sigma:G^N\to G)\rtimes S_N.
\]
\end{lemma}
\begin{proof}
The underlying transpositions form a connected graph, by transitivity.
Choose a spanning tree from them. Successively changing the origins
along the tree makes all tree labels zero. These tree transpositions
generate the ordinary $S_N$ subgroup. If another generator exchanges
$i,j$ with label $c$, multiplying it by the unlabelled transposition
gives the translation with $c,-c$ at $i,j$. Conjugating by $S_N$
gives the same difference at any pair of indices.
Let $G_0\subset G$ be generated by all the remaining labels. The
group generated is therefore exactly
$\ker(\Sigma:G_0^N\to G_0)\rtimes S_N$.
Its orbits on $G\times\{1,\ldots,N\}$ are indexed by $G/G_0$.
Transitivity forces $G_0=G$.
\end{proof}

\begin{lemma}[Lifted intersection-point monodromy]
\label{lem:mixed-lifted-intersection-wreath-candidate}
\status{Proved under stated hypotheses}
Fix $\alpha\in\Pic^0(X)$ and a smooth sharp curve
$C\in|L\otimes\alpha|$, where sharp means two-rank nine.
Let $\widetilde C\to C$ be a connected finite etale torsor
with group $G$, a finite elementary abelian two-group.
As $D$ varies through a dense open of $|L\otimes\alpha^{-1}|$
with transverse intersection with $C$, the lifted cover of its
$N=18$ intersection points has image
\[
                     \ker(\Sigma:G^{18}\to G)\rtimes S_{18},
\]
up to choices of origins in the point fibres. The same assertion
holds after requiring $D$ to be smooth and sharp.
\end{lemma}
\begin{proof}
Use the image of $H^0(X,L\otimes\alpha^{-1})$ in
$H^0(C,(L\otimes\alpha^{-1})|_C)$ as a linear system on $C$.
It is very ample, because the line bundle on $X$ is very ample.
Its projective dimension is eight, or seven when the restriction
has its one-dimensional kernel. The original parameter space maps
onto it with integral affine fibres away from the possible centre;
this pullback does not change a finite geometric monodromy image.

We spell out the generic tangency assertion in characteristic two.
Each projection of $C$ to an elliptic factor has degree three,
so it is separable and its ramification set is finite. A degree-three
line bundle on an elliptic curve separates second jets except at
the finite set of points $P$ with that line bundle equal to
$\mathcal O(3P)$. Consequently the product line bundle separates
second jets on the surface at a general point of $C$, and hence
on $C$ itself.

For a fixed surface point the joint first-jet rank at it and a
second point is six away from the two factor fibres and the
finitely many crossed tangent incidences of
Lemma~\ref{lem:product-small-degree-jets}. The same proof
applies to the present two cubics and to their translates.
Its polar-finiteness input holds uniformly: the three partials
of $Y^2Z+XYZ+X^3+aZ^3$ are linearly independent, as seen
from the respective coefficients of $X^2,Y^2,XZ$; the
partials of $Y^2Z+YZ^2+X^3$ are $X^2,Z^2,Y^2$.
A nonzero polar quadratic therefore cannot vanish on either
entire cubic, so only finitely many tangents pass through
any prescribed point. The exceptional intersection with $C$
is finite, since both projections of $C$ are finite. Thus on
$C\times C$ the four curve first jets have rank four away from
a subset of dimension at most one. Rank two means that the two
embedded curve tangent lines coincide. At a point whose two
projection derivatives are nonzero, its tangent line contains
no other point of the Segre surface: a nonzero point on that
line has tensor rank two. The remaining points of $C$ are
finite, and each of their tangent lines meets $C$ in finitely
many points. Hence the rank-two pairs form a finite set.

The incidence of two distinct tangencies consequently has
codimension at least two in the hyperplane parameter space:
its three strata have point dimensions at most $2,1,0$ and
impose at least $4,3,2$ conditions. Higher contact at one
point also has codimension at least two, by generic second-jet
separation; at the finitely many exceptional points, the two
first-jet conditions alone give the same bound.
The tangent incidence is an irreducible projective
bundle over $C$, and its generic fibre over its image is finite.
Its image is therefore an irreducible divisor $\Delta$, whose
geometric generic member has precisely one contact of order
two and otherwise simple intersections.

The resulting unlabelled $S_{18}$ conclusion is a known
instance of the quasireflexive hyperplane-monodromy theorem
\cite[Theorem~2, p.~3; definition on p.~6; Proposition~3.1
and the proof of Theorem~2, pp.~7--8]{Entin2021}:
the jet argument above verifies its geometric hypothesis.
The labelled sum-zero wreath extension and its later gluing
application are the additional arguments proved here.

The local degree-two root cluster can be written
$z^2+bz+c=0$: after henselization the simple roots split off,
and $1,z$ form a basis of its remaining length-two algebra.
The independent constant and linear jet parameters make
$db,dc$ independent at a simple tangency, so $b,c$ are part
of a local parameter system in the coefficient base.
In particular $\Delta$ is locally $b=0$, and $c$ is not a
square at its separably closed generic residue field. The
equivalent equation $y^2+y=c/b^2$ proves that inertia interchanges
the two roots. This handles the possible purely inseparable
generic tangent-point field: no separability of the marked
tangency incidence is assumed. Since $\widetilde C\to C$ is
etale, the lifted inertia is a labelled transposition, fixing
all other point fibres.

The incidence of a point of $\widetilde C$ and a hyperplane
through its image is an irreducible projective bundle over
$\widetilde C$. Its map to the hyperplane space is finite:
no member of this embedding contains all of $C$.
It follows that the lifted-point image is transitive.
The global deck transformations commute with monodromy, so
the image is a subgroup of $G^{18}\rtimes S_{18}$.
The Galois closure has no branch divisor except $\Delta$.
Its quotient by the normal subgroup generated by divisorial
inertia is unramified in codimension one and therefore etale
over projective space by purity \cite{Stacks0BMB}; it is trivial
because geometric projective space is simply connected, as proved
in Lemma~\ref{lem:char2-defect-one-inertia}.
The image is thus generated
by these labelled transpositions. Apply the preceding lemma.

Finally smooth sharp $D$ form a nonempty dense open by the
rank theorem and the degree-one result. Restricting a finite
etale cover of a normal integral open to a dense open preserves
its connected components and monodromy image.
\end{proof}

\begin{lemma}[The nodal boundary contains all gluing translations]
\label{lem:mixed-degree-two-boundary-group-candidate}
\status{Proved under stated hypotheses}
Fix $\alpha\in\Pic^0(X)$, and vary smooth sharp transverse
$C\in|L\otimes\alpha|$, $D\in|L\otimes\alpha^{-1}|$.
Their nodal unions belong to the rank-$34$ relative unit open
inside $|L^2|$. On this boundary the relative representation
has an exact sequence
\[
 \begin{gathered}
 0\longrightarrow T\longrightarrow V\longrightarrow R\longrightarrow0,
 \qquad T=\mathbf F_2^{18}/\mathbf F_2\mathbf1,\\
 R=(H^1_{\et}(C,\mathbf F_2)\oplus
             H^1_{\et}(D,\mathbf F_2))/H^1_{\et}(X,\mathbf F_2).
 \end{gathered}
\]
Put $T^0=\{(x_i):\sum x_i=0\}/\mathbf F_2\mathbf1$,
of dimension sixteen. The boundary image contains all
translations $\operatorname{Hom}(R,T^0)$ and an actual
$\GL_8(\mathbf F_2)\times\GL_8(\mathbf F_2)$ acting naturally
on two eight-dimensional direct summands and fixing a
complement. Moreover every nonzero invariant subspace of
this boundary representation contains $T^0$.
\end{lemma}
\begin{proof}
Each normalized curve has genus ten and two-rank nine.
Their eighteen nodes add seventeen Frobenius-unit graph
directions. Thus the union has two-rank $9+9+17=35$;
the ambient rank is one, giving relative rank $34$.
The normalization sequence gives the displayed exact sequence,
with $\dim R=17$.
Write $A$ for the one-dimensional ambient unit line.
Both relative curve images are $\GL_8$ by the preceding
uniform degree-one theorem. This applies to every Picard
twist: the polarization of $L$ is an isogeny, so over $k$
every $L\otimes\alpha$ is a translate of $L$.
Since $C,D$ vary independently in a dense open of a product,
their two relative images are independent.

We give the needed group-extraction details. The image on
$H^1_{\et}(D,\mathbf F_2)$ fixes $A$ pointwise and maps
onto $\GL_8$; its unipotent kernel is either zero or the
entire natural-dual module of dimension eight. The previously
proved $H^1(\GL_8,\mathrm{nat}^*)=0$ splits the possible
complement. In either case this image is perfect and its
only nonabelian simple composition factor is $\GL_8$.
The analogous statement holds for $C$.
On $R$ the image is similarly an extension of
$L_0=\GL_8\times\GL_8$ by a submodule of
$(\mathrm{nat}_C^*\oplus\mathrm{nat}_D^*)\otimes A$.
These two simple modules are nonisomorphic; the same
first-cohomology vanishing makes the image a semidirect
product and therefore perfect.

The point-permutation image on $T$ is $S_{18}$ by the
preceding lemma. It and the image on $R$ have no common
nontrivial quotient. Indeed the latter is perfect with
nonabelian simple factors $\GL_8$, whereas the former has
only $A_{18}$ and cyclic factors. The two simple groups
are different (for example $11$ divides $|A_{18}|$ but
does not divide $|\GL_8(\mathbf F_2)|$).
Goursat's lemma gives their product as joint image.
Take the inverse image of $1\times L_0$ after choosing
the semidirect complement on $R$. Its matrix translation
kernel is an $L_0$-submodule of $\operatorname{Hom}(R,T)$.
Under $L_0$ this is a direct sum of copies of the two
natural-dual modules and the trivial module. Every quotient
has zero $H^1$: this uses the proved natural-dual vanishing
and perfectness of $L_0$ for the trivial summands.
The usual graph-of-a-cocycle argument therefore supplies
an actual block-diagonal $L_0$ in the boundary image,
fixing $T$ and a complementary ambient line in $R$.

Let $J$ be the kernel of the boundary action on $T\oplus R$.
Then $J\subset\operatorname{Hom}(R,T)$ is additive and is
an $L_0$-submodule. Fix $C$ and choose its connected maximal
elementary two-cover, whose group is
$G=H^1_{\et}(C,\mathbf F_2)^*$ of dimension nine.
The preceding lemma gives the lifted-node group
$G^{17}\rtimes S_{18}$. Its only nonabelian simple factor
is $A_{18}$. Its joint image with the moving $D$ cohomology
is their product, by the same Goursat argument: a common
quotient of the perfect $D$ image would have to have a
nonabelian simple factor, and none is shared.

Choose a pure lifted-point translation $(g_i)$ with
$\sum g_i=0$ and trivial moving-$D$ cohomology action.
It gives an element of $J$. On the injected
$H^1_{\et}(C,\mathbf F_2)\subset R$ its translation is
\[
                       \xi\longmapsto(\xi(g_i))_i
                                     \bmod\mathbf F_2\mathbf1.
\]
For precision, a class $\xi$ is a fixed etale two-torsor
on $C$. Extend it to the union by the trivial torsor on
$D$ and a chosen gluing at every node. Continuing those
gluing choices around the loop changes them by exactly
the displayed node evaluations. The two constant changes
of trivialization remove their diagonal class. This proves
the formula directly, rather than identifying an unspecified
extension class with a point-cover action.

As $(g_i)$ varies, these maps are all of
$\operatorname{Hom}(H^1_{\et}(C,\mathbf F_2),T^0)$.
The $L_0$-module $\operatorname{Hom}(R,T)$ is semisimple,
with three isotypic parts corresponding to the two
relative eight-dimensional spaces and to $A$.
An $L_0$-submodule is the direct sum of its intersections
with these parts. The preceding surjectivity therefore
puts all maps from the $C$ relative space and from $A$
to $T^0$ in $J$, with the other columns zero.
Interchanging $C,D$ supplies the remaining relative
columns. Hence $\operatorname{Hom}(R,T^0)\subset J$.

If an invariant subspace has nonzero image in $R$, these
translations force it to contain $T^0$. If it lies in
$T$, take a nonconstant representative of a nonzero
vector and apply a transposition exchanging unequal
coordinates. Their difference is the class of $e_i+e_j$.
Its $S_{18}$-orbit spans $T^0$. This also proves the last
assertion for subspaces of $T$.
\end{proof}

\begin{lemma}[Comparison at a common grid through the finite Lang part]
\label{lem:mixed-degree-two-common-grid-candidate}
\status{Proved under stated hypotheses}
Let $S$ be a nonzero sub-local system of the degree-two
relative system, and suppose $S\subset K$. Then $S=K$.
In particular the geometric representation on $K$ is irreducible.
\end{lemma}
\begin{proof}
Choose six distinct points $P_0,\ldots,P_5$ on $E$ and
six distinct points $Q_0,\ldots,Q_5$ on $E_0$, each six-point
divisor linearly equivalent to $6O$. Put
\[
 Z=\bigcup_i(\{P_i\}\times E_0)\ \cup\
   \bigcup_j(E\times\{Q_j\})\in|L^2|.
\]
Its dual graph is the complete bipartite graph $K_{6,6}$.
The graph has first Betti number $36-12+1=25$; the six
ordinary components contribute six further unit directions,
and the ambient unit line removes one. Thus the relative
stable rank at $Z$ is $30$, below the generic rank $34$.
We use only its finite Lang part, not a specialization
isomorphism of full fibres.

For any row and column three-element subsets $I,J$, group
the corresponding three vertical and three horizontal
components into $C_{I,J}$, and the complementary components
into $D_{I,J}$. Their line bundles are
$L\otimes\alpha_{I,J}$ and $L\otimes\alpha_{I,J}^{-1}$
for a suitable Picard class. This opposite twist is essential:
we do not require distinct partitions of six points to
give two triples each linearly equivalent to $3O$.
The two curves meet at the eighteen cross-edges, all simple,
and are smooth along those intersection points.
Their complete parameter systems give a boundary family
through $Z$; smooth sharp transverse pairs form a dense
open in each such family.

Here is the precise compatibility of the rank-dropping
comparison. Take the strict henselization $R_Z$ of the
local coefficient ring of $|L^2|$ at $Z$. It is a normal
local domain. The full coherent Frobenius bundle is
regular there, and its affine Lang scheme is separated,
quasi-finite and etale. Its finite part is a constant
finite etale group with special fibre the thirty-dimensional
relative unit space $W_Z$. This is the already proved
henselian finite-part construction. Its restriction to
the nonempty generic rank-$34$ open injects into $V$.

Extend $V$ and $S$ to the full coherent maximal-rank open,
including nodal unions. This does not change their images:
the original smooth open is dense in a normal integral
parameter space. For each constant section of the finite
part, membership in $S$ on that open is an open-and-closed
condition. The open is integral, so membership is either
everywhere or nowhere. These sections define one fixed
subspace $S_Z\subset W_Z$.

Pull back to the strict local parameter scheme of any
boundary family through $Z$. Its eighteen cross intersection
points remain finite etale, so its graph space $T$ and
heart $T^0$ extend as constant unit systems even at $Z$.
Their special injections into $W_Z$ agree with the
restrictions of the ambient finite Lang sections: two
sections of a separated etale scheme over a henselian
local scheme with the same closed value coincide.
The preceding boundary lemma puts $T^0$ inside $S$ at
the generic boundary point. It follows that each of these
special cross-hearts lies in the same $S_Z$.

We now compare two partitions. First take
$I_1=J_1=\{0,1,2\}$, and then
$I_2=J_2=\{0,1,3\}$.
Represent graph cohomology by edge cochains modulo vertex
coboundaries, and denote the edge $(i,j)$ by $e_{ij}$.
The second cross-heart contains
\[
                       e_{02}+e_{12},\qquad e_{34}+e_{35}.
\]
Both have even sum on its cross-edges. The first is supported
on the internal graph of $C_{I_1,J_1}$, and is nonzero
there: its evaluation on the four-cycle with rows $0,2$
and columns $1,2$ is one. The second is supported on the
internal graph of $D_{I_1,J_1}$, and its evaluation on
the four-cycle with rows $3,4$ and columns $3,4$ is one.

The two corresponding constant finite-Lang sections belong
to $S$ on the first boundary family. Projection from union
cohomology to the two component relative cohomologies is
regular and Frobenius-equivariant even at the grid. Indeed
$H^1(X,(L\otimes\alpha)^{-1})=0$ by Kunneth, so ambient
$H^1(\mathcal O_X)$ injects for every member. Connectedness,
constant arithmetic genus and cohomology-and-base-change give
the relevant coherent vector bundles. Restriction commutes
with Frobenius, as do the quotients by ambient cohomology.
The displayed graph classes survive that quotient: on each
$3\times3$ grid, ambient coherent cohomology maps injectively
to the sum of the six elliptic-component cohomologies, while
graph classes map to zero there. Hence the projections are nonzero at
$Z$ and consequently nonzero generically in that boundary
family.

The image of $S$ on the two relative component systems
therefore has nonzero projection to each of the independent
natural $\GL_8$ representations. It contains their direct
sum, of dimension sixteen. Its kernel already contains
the first cross-heart, of dimension sixteen. Thus
$\dim S\ge32=\dim K$, proving $S=K$.
\end{proof}

\begin{theorem}[Full degree-two image for every mixed product]
\label{thm:mixed-every-fixed-modulus-degree-two-candidate}
\status{Proved under stated hypotheses}
For every perfect field of characteristic two, every ordinary
elliptic curve $E$ and every supersingular elliptic curve $E_0$,
the relative system for $|L^2|$ on $E\times E_0$ has
\[
 G_{\mathrm{geom}}=G_{\mathrm{arith}}
   =\operatorname{Hom}(\mathbf F_2^2,\mathbf F_2^{32})
                         \rtimes\GL_{32}(\mathbf F_2),
\]
the full pointwise stabilizer of its two invariant covectors.
Its common covector kernel is irreducible and is contained
in every nonzero invariant subspace.
\end{theorem}
\begin{proof}
Work geometrically first. The boundary lemma shows that every
nonzero invariant subspace of $V$ contains a cross-heart
$T^0\subset K$. The common-grid lemma makes $K$ irreducible,
so every such subspace contains $K$. In particular $V$ has
no invariant vector.

In a boundary frame, the two actual natural eight-dimensional
blocks and $T^0$ lie in $K$: an invariant covector vanishes
on each of these nontrivial simple boundary constituents.
Their dimensions sum to $8+8+16=32$, so they are all of $K$.
The boundary translations include every rank-one map from
either natural block to $T^0$, acting trivially on the other
summands.

Let $N$ be the normal subgroup of the geometric image on $K$
generated by its transvections. Since that image is irreducible,
the restriction to $N$ is semisimple: the nonzero $N$-socle
is invariant under the whole group. Each isotypic multiplicity
is one, because a rank-one transvection acts nontrivially on
only one simple constituent with multiplicity one; an
isotypic component fixed by all such generators would be
$N$-fixed and hence globally invariant. Thus the simple
$N$-summands form a direct sum. The first $\GL_8$ natural
block lies in a single summand. Each rank-one map from it
to an arbitrary vector of $T^0$ forces that vector into
the same summand. The analogous maps from the second
natural block put that block in the same summand as well.
Consequently this summand is all of $K$, so $N$ is irreducible.

The transvection classification
\cite[Theorem~3.6, printed p.~22]{KantorSeress2001}, in the
already used $\mathbf F_2$ form, now gives $N=\GL(K)$ or
a group preserving a nondegenerate alternating form.
The latter alternative is impossible. The actual natural
$\GL_8$ block has no invariant bilinear form, and its
pairing with a complement fixed by that block is zero
because the natural representation has no invariant
linear functional. Such a block would lie in the radical.
Thus the image on $K$ is $\GL_{32}(\mathbf F_2)$.

The quotient $V/K$ is the constant space $\mathbf F_2^2$.
The remaining translation kernel is $K\otimes U$ for a
subspace $U\subset(\mathbf F_2^2)^*$. Modulo that kernel,
the subgroup is a graph of a cocycle. The proved vanishing
$H^1(\GL_{32},K)=0$ makes it a linear complement. If $U$
were proper, a nonzero quotient direction annihilated by
$U$ would have a fixed lift, contradicting absence of
invariant vectors. Therefore $U=(\mathbf F_2^2)^*$ and
the full displayed stabilizer occurs. The two entire
covectors descend arithmetically over every perfect base
field by the nonordinary covector theorem. Arithmetic
equality and descent of the geometric statement follow.
\end{proof}

\begin{remark}[Scope and retained checks]
\label{rem:mixed-degree-two-wreath-grid-scope-candidate}
\status{Proved under stated hypotheses}
This argument uses opposite Picard twists at the common grid
and only its finite Lang part of rank thirty. It does not
identify that special fibre with the generic rank-$34$ space.
The new uniform result concerns degree two; the fixed-modulus
all-higher-degree problem remains separate. The supporting graph
certificate is
\path{agents/monodromy/code/mixed_degree_two_grid_checks.py},
with its retained output. It checks the actual graph quotients,
cross-hearts and two nonzero projections; it does not replace
the finite-Lang comparison or geometric monodromy proof.
The author execution took $0.000195655$ seconds. Independent
root, reduction/descent and obstructions/crystals reruns took
$0.000138148$, $0.000105329$ and $0.000102624$ seconds,
respectively, preserving the original output. All geometric and
group steps have passed the complete
independent proof audits described above.
\end{remark}


% /root/monodromy, post-CP18. Root, reduction/descent and obstructions/crystals full proof audits passed.
\subsection{A uniform strip of mixed product polarizations}
\label{subsec:mixed-three-by-b-strip-candidate}
\status{Proved and independently checked. Complete proof audits passed;
the supporting finite certificate was independently rerun.}

Write $X=E\times E_0$, where $E$ is ordinary and $E_0$ is
supersingular in characteristic two. For an integer $b\ge3$ put
\[
 N_b=\mathcal O_E(3O)\boxtimes\mathcal O_{E_0}(bO),\qquad m=b-3.
\]
The geometric arguments below are over an algebraically closed field.
The final theorem states the descent to an arbitrary perfect field.
All local systems are the relative unit systems, with the ambient
one-dimensional etale cohomology divided out.

\begin{lemma}[Lifted degree-three branch monodromy]
\label{lem:mixed-strip-lifted-triple-candidate}
\status{Proved; complete independent proof audits passed.}
There is a dense open of smooth sharp curves $C\in|N_3|$ such
that the projection $f:C\to E_0$ has degree three and all its
ramified fibres have one double point and one simple point.
For such $C$, and every connected etale torsor $\widetilde C\to C$
with finite elementary abelian two-group $G$, the monodromy of
the lifted three-point fibre of $f$ is
\[
                  \ker(\Sigma:G^3\to G)\rtimes S_3.
\]
\end{lemma}
\begin{proof}
Choose plane cubic embeddings of the two factors. An invertible
coefficient matrix for a section of $N_3$ identifies the second
cubic with a projective transform of a cubic in the dual plane
of the first. Its point $Q$ specifies the line cutting out the
three points of $f^{-1}(Q)$ on $E$. The flex lines of the first
cubic form a finite set. Indeed a contact of order at least
three at $P$ means $\mathcal O_E(3O)=\mathcal O_E(3P)$,
and multiplication by three is finite etale. A general projective
transform of the second cubic avoids these finitely many points
of the dual plane. Simultaneously its incidence curve is smooth
and sharp, by the smooth open and the proved uniform degree-one
rank theorem. Thus every nonsimple line intersection is exactly
double plus simple. The projection has degree three, hence is
separable; at a double point its local ramification index is two.
Its Galois inertia is consequently a transposition. This reasoning
does not assume that the Gauss map or marked tangency incidence
is separable. Branching is nonempty, for an etale cover of an
elliptic curve has genus one, whereas $C$ has genus ten.

The lifted monodromy $H$ is transitive on $G\times\{1,2,3\}$
and lies in $G^3\rtimes S_3$. Since the torsor is etale on $C$,
each inertia generator is a labelled transposition, fixing the
remaining fibre pointwise. Let $J$ be their normal closure in
$H$. The projection of $H$ to $S_3$ is transitive and contains a
transposition, so it is $S_3$. The projection of $J$ is also
$S_3$. The spanning-tree proof of
Lemma~\ref{lem:labelled-transposition-wreath-candidate}, without
its final transitivity conclusion, therefore gives, after changing
origins,
\[
                 J=\ker(\Sigma:G_0^3\to G_0)\rtimes S_3
\]
for a subgroup $G_0\subset G$.

We record the normalizer calculation, so no residual cover is
silently discarded. A translation $(g_1,g_2,g_3)$ normalizes this
group precisely when its three entries have the same image in
$G/G_0$: conjugating an unlabelled transposition gives the label
$g_i-g_j$. Modulo the sum-zero translations in $G_0^3$, the
normalizing translation group is elementary abelian of dimension
$\dim G$. The permutation group already lies in $J$. Hence
$H/J$ is a two-group. The corresponding quotient of the Galois
closure of $\widetilde C\to E_0$ is unramified everywhere and
therefore etale over $E_0$. A supersingular elliptic curve has no
nontrivial etale two-cover, and hence no nontrivial etale
two-group quotient. Thus $H=J$. Transitivity on the lifted
fibre forces $G_0=G$, proving the formula.
\end{proof}

\begin{lemma}[A sharp boundary in a fixed strip system]
\label{lem:mixed-strip-boundary-candidate}
\status{Proved; complete independent proof audits passed.}
For $b\ge4$, the system $|N_b|$ has a boundary family of nodal
curves consisting of a translated sharp curve of type $(3,3)$
and $m$ distinct horizontal ordinary fibres. These curves have
the maximal relative rank $r=3b-1$.
On this boundary the relative unit system has an exact sequence
\[
 0\longrightarrow T\longrightarrow V\longrightarrow
 R=U\oplus A^m\longrightarrow0,
 \qquad T=\bigoplus_{j=1}^m T_j,
 \quad T_j=\mathbf F_2^3/\mathbf F_2\mathbf1,
\]
where $\dim U=8$, $A=H^1_{\et}(E,\mathbf F_2)$ has dimension
one, and the symmetric group on the fibres permutes the $m$
summands. The image contains an actual natural $\GL_8$ block
fixing a complement, and all translations
\[
            \operatorname{Hom}(U\oplus A_j,T_j)
                    \quad(1\le j\le m),                 \tag{Strip-1}
\]
with every other column zero. Every nonzero invariant subspace
of the boundary representation contains all of $T$.
\end{lemma}
\begin{proof}
Begin with $C\in|N_3|$ and ordered points $Q_1,\ldots,Q_m$
of $E_0$, with distinct unramified fibres on $C$. Choose $t\in E_0$
with
\[
                         [b]t=\sum_j Q_j.
\]
Translate the union $C\cup\bigcup_j(E\times\{Q_j\})$ by
$-t$ in the second factor. Its line bundle is exactly $N_b$.
The parameter space with this chosen $t$ is smooth and integral:
on the ordered space solve
$Q_m=[b]t-\sum_{j<m}Q_j$, leaving coordinates
$(t,Q_1,\ldots,Q_{m-1})$ on an open of $E_0^m$.
The translates are actual isomorphisms, and therefore identify
the component cohomology with that of $C$ and the fixed curve $E$.

There are three transverse nodes for each ordinary fibre.
The graph rank is $2m$. The normalization has two-rank $9+m$,
so the union has two-rank $9+3m$, and relative rank
$8+3m=3b-1$. This is the full coherent relative dimension:
$H^1(X,N_b^{-1})=0$ by Kunneth, the genus is $3b+1$, and the
ambient coherent dimension is two. It follows both that the
boundary lies in the maximal-rank open and that this open meets
the smooth locus. Thus the asserted rank is also generic.

The degree-three projection $C\to E$ supplies a canonical
splitting
\[
 H^1_{\et}(C,\mathbf F_2)=A\oplus U.
\]
Pullback followed by trace is multiplication by three, hence the
identity. The image on $U$ is $\GL_8$ by the proved degree-one
theorem. If the ambient coordinate of a class on $C$ is $a$,
and those on the $m$ ordinary fibres are $a_j$, identify the
relative constant body coordinates with
\[
                              r_j=a_j+a.                 \tag{Strip-2}
\]
This gives the displayed exact sequence.

Fix a curve $C$ in the dense open of the preceding lemma and
take its maximal elementary two-torsor, with group
$G=H^1_{\et}(C,\mathbf F_2)^*$. Independently moving the ordered
$Q_j$ gives the product of $m$ copies of
$H_0=G^2\rtimes S_3$; allowing their permutation gives
$H_0\wr S_m$. The cover obtained by choosing $t$ does not shrink
these images. Indeed $[b]$ on the supersingular elliptic curve
has a radicial two-primary part and an etale abelian part of odd
degree. Radicial pullback preserves etale fundamental groups.
The abelianization of $H_0$ is $\mathbf Z/2$: the coinvariants
of the sum-zero $G^3$ under $S_3$ vanish. The abelianization of
its wreath product, as well as of its ordered product, is a
two-group. Therefore these images have no common nontrivial
quotient with the odd abelian deck group. Their restriction to
the $t$-cover remains the full image.

For a fixed $j$, a pure lifted-point translation with three
labels summing to zero acts on the gluing of etale torsors by
node evaluation. The $C$-ambient torsor and the ordinary-fibre
torsor are both pulled back from $E$, so their node evaluations
coincide. After quotienting the global ambient class, their
coefficient is precisely $r_j$ in (Strip-2). The other eight
coefficients are those of $U$. Thus these labels induce every
map in (Strip-1), while fixing all other component and graph
coordinates. This is the same explicit gluing calculation as
in Lemma~\ref{lem:mixed-degree-two-boundary-group-candidate};
here three is odd, so the sum-zero node space maps isomorphically
to $\mathbf F_2^3/\mathbf F_2\mathbf1$.

For completeness an actual $\GL_8$ block, rather than just a
quotient, can be extracted. The graph and fibre-permutation
image is a quotient of $S_3\wr S_m$, whose nonabelian simple
factors are at most $A_m$. None is $\GL_8(\mathbf F_2)$:
if $m\le10$ its order has no factor $17$, whereas if $m\ge11$
the factor $A_m$ has a factor $11$, absent from $|\GL_8|$.
Goursat's lemma therefore allows the graph and fibre permutations
to be killed while retaining the full $\GL_8$ image on $U$.
The remaining translation module is a direct sum of natural-dual
and trivial $\GL_8$ modules. Its every quotient has zero $H^1$,
by the previously proved natural-dual vanishing and perfectness.
The cocycle argument gives an actual natural block fixing a
complement.

Finally, if a nonzero invariant subspace has a nonzero $U$
coordinate, (Strip-1) forces every $T_j$ into it. If its body
coordinate is nonzero only in some $A_j$, it contains $T_j$;
permutation of the fibres then gives all $T$. If it is contained
in $T$, the independent $S_3$ factors and the simplicity of their
two-dimensional natural modules give one $T_j$, and again all
of them by permutation. This proves the last assertion.
\end{proof}

\begin{lemma}[The strip grid and the regular trace projections]
\label{lem:mixed-strip-grid-candidate}
\status{Proved; complete independent proof audits passed.}
For odd $b\ge3$ the geometric relative representation for
$|N_b|$ is irreducible. For even $b\ge4$, trace along the first
projection defines a nonzero invariant covector $\lambda$;
its kernel $K$ has dimension $3b-2$, is irreducible, and is
contained in every nonzero invariant subspace of $V$.
\end{lemma}
\begin{proof}
The case $b=3$ is the proved degree-one theorem, so assume $b\ge4$.
Choose three distinct points $P_i$ on $E$ whose sum is zero and
$b$ distinct points $Q_j$ on $E_0$ whose sum is zero. The union
of the corresponding three vertical and $b$ horizontal fibres
is a divisor $Z\in|N_b|$ with graph $K_{3,b}$. Its relative
unit space $W_Z$ has dimension
\[
            (2b-2)+b-1=3b-3.                            \tag{Strip-3}
\]
In particular it is two dimensions smaller than the generic
space; no full-fibre specialization isomorphism is used.

Let $S$ be a nonzero sub-local system of $V$. In the even case,
first suppose $S\subset K$. Work over the strict henselization
of the coefficient space at $Z$. The henselian finite part of
the coherent Lang scheme is constant with fibre $W_Z$ and
injects into the generic unit system. Membership of one of its
sections in $S$ is clopen on the integral maximal-rank open,
and therefore gives a fixed subspace $S_Z\subset W_Z$. These
are exactly the finite-part and membership constructions used
and proved in Lemma~\ref{lem:mixed-degree-two-common-grid-candidate}.

For every three-element subset $J$ of the horizontal components,
group the three vertical components and the three horizontals
in $J$ into a curve $C_J$ of type $(3,3)$. The other $m$
components remain ordinary fibres. This is a boundary family
of the preceding lemma through the same $Z$, with its appropriate
Picard translate. More explicitly choose $\tau$ with
$3\tau=-\sum_{j\in J}Q_j$, so that $C_J+\tau\in|N_3|$.
The remaining fibre points become $Q_k+\tau$, and their sum
is $b\tau$. The preceding normalization therefore permits
$t=\tau$ and returns the original grid.

For this partition the three nodes on each outside fibre remain
finite etale near $Z$. Their two-dimensional graph space extends
as a constant unit system, and agrees with the ambient finite
Lang part by uniqueness of henselian etale lifts. The boundary
lemma puts it in $S_Z$. Vary $J$. For every column of $K_{3,b}$,
some choice excludes that column, and its three edge cochains
modulo the column-star coboundary give the corresponding
two-dimensional space. These column spaces span all graph
cohomology. Hence
\[
                         H^1(K_{3,b},\mathbf F_2)\subset S_Z.
                                                               \tag{Strip-4}
\]
For a fixed $J$, its restriction to the internal $K_{3,3}$ graph
is nonzero. This graph survives quotienting by ambient cohomology,
because ambient classes inject into the normalization cohomology.
The coherent restriction to the relative cohomology of $C_J$
is regular near $Z$, by Kunneth vanishing, connectedness, and
cohomology and base change. Thus the projection of $S$ to its
generic relative eight-dimensional system is nonzero. The actual
$\GL_8$ block in the boundary lemma forces its image to contain
the entire $U$. Since $S$ already contains $T$, it contains the
full inverse image of $U\subset U\oplus A^m$.

We explain precisely how this assertion supplies component
directions at the rank-dropping grid. Trace to $E$ remains a
regular coherent map for the family of curves $C_J$, even when
vertical components are contracted. Define it by dualizing the
pullback map
\[
               H^0(E,\Omega_E^1)\longrightarrow H^0(C_J,\omega_{C_J}).
\]
The relative dualizing sheaf gives the same construction in the
whole family. On the smooth dense open this is the ordinary
degree-three trace, commutes with Frobenius, and its composite
with pullback is three times the identity. These equalities
extend between coherent vector bundles across the grid. At the
grid, normalization and Serre duality identify it with the sum
of the three ordinary-component classes; graph and vertical
component classes contribute zero.

Consequently all the constant-body projections $r_k$ in
(Strip-2) are regular Frobenius-equivariant maps even at $Z$.
A finite Lang section whose special values under these maps
are zero has those values identically zero: their targets are
the constant unit line $A$. It therefore lies generically in
the full inverse image of $U$, already contained in $S$.
It belongs to $S_Z$. This uses kernels of regular constant-body
maps, not a presumed surjectivity of a source kernel under a
rank drop.

Represent the ordinary-component quotient of $W_Z$ by
$\mathbf F_2^b/\mathbf F_2\mathbf1$, with coordinates $a_j$.
For the partition $J$, trace on $C_J$ is
$a=\sum_{j\in J}a_j$, and the outside body coordinates are
$r_k=a_k+a$. In particular every component vector supported
on $J$ with sum zero has all $r_k=0$. The preceding paragraph,
together with (Strip-4), places all these classes in $S_Z$.
As $J$ varies, they span the image of
\[
              \{(a_j)\in\mathbf F_2^b:\sum_j a_j=0\}
                  \quad\text{in }\mathbf F_2^b/\mathbf F_2\mathbf1.
                                                               \tag{Strip-5}
\]
Indeed each pair vector $e_i+e_j$ is supported on a triple.

Here is the final dimension calculation. Fix one partition $J$.
On the space in (Strip-5), put $a=\sum_{j\in J}a_j$ and
$r_k=a_k+a$ for its $m$ outside indices. The sum-zero condition
gives
\[
                         \sum_{k\notin J}r_k=(m+1)a.       \tag{Strip-6}
\]
If $b$ is odd, then $m$ is even, and the image is all of $A^m$:
given the $r_k$, take $a=\sum r_k$, choose the three inside
coordinates with sum $a$, and the outside coordinates $r_k+a$.
If $b$ is even, then $m$ is odd, and the image is precisely
the hyperplane $\sum r_k=0$, of dimension $m-1$; one can choose
$a=0$ to realize any vector in that hyperplane.

For odd $b$, $S$ thus has body image containing $U$ and all
$m$ constant directions, and contains $T$. Its dimension is
at least $8+m+2m=3b-1$, so $S=V$.
For even $b$, trace of the full curve to $E$ annihilates the
ambient line, since its degree is $b$. It gives the covector
$\lambda$, whose boundary expression is $\sum r_k$.
It is nonzero and has kernel of dimension $3b-2$. The same
count for $S\subset K$ gives
\[
                    \dim S\ge 8+(m-1)+2m=3b-2,
\]
so $S=K$. Finally any nonzero invariant subspace of $V$ contains
$T\subset K$ by the boundary lemma; irreducibility of $K$
then forces it to contain $K$. This completes the proof.
\end{proof}

\begin{theorem}[Full images for all mixed three-by-$b$ strips]
\label{thm:mixed-three-by-b-full-groups-candidate}
\status{Proved; complete independent proof audits passed.}
Let $k$ be any perfect field of characteristic two, let $E/k$
be ordinary and $E_0/k$ supersingular, and let $b\ge3$.
The generic relative rank for $|N_b|$ is $3b-1$. Its full
geometric and arithmetic images are equal, and are
\[
 \begin{cases}
  \GL_{3b-1}(\mathbf F_2),&b\text{ odd},\\[2pt]
  \operatorname{Hom}(\mathbf F_2,\mathbf F_2^{3b-2})
          \rtimes\GL_{3b-2}(\mathbf F_2),&b\text{ even}.
 \end{cases}
\]
In the even case the displayed group is the full pointwise
stabilizer of the trace covector from the first projection,
and that is its entire invariant covector space.
\end{theorem}
\begin{proof}
For $b=3$ use the uniform degree-one theorem. Otherwise let $K=V$
if $b$ is odd, and let $K=\ker\lambda$ if $b$ is even. The
preceding lemma proves that the geometric image on $K$ is
irreducible. Let $J$ be its normal subgroup generated by
transvections. Its representation on $K$ is semisimple: the
nonzero $J$-socle is invariant under the entire irreducible group.
Every nontrivial isotypic multiplicity is one, because some
rank-one generator acts on it. The $J$-fixed part is zero,
since it is invariant under the whole group and the actual
$\GL_8$ block acts nontrivially.

The natural eight-dimensional block belongs to one simple
$J$-summand. The translations in (Strip-1) from that block to
arbitrary vectors of every $T_j$ force all of $T$ into the
same summand. Every other $J$-summand is killed by these
translations and by the natural block. It is also killed by
the translations from the individual constant-body coordinates
to $T_j$. These coordinate functionals span the dual of the
constant body (or its sum-zero hyperplane in the even case).
Such a summand must consequently lie in $T$, already in the
first summand. Hence there are no others and $J$ is irreducible.

Apply the transvection classification
\cite[Theorem~3.6, printed p.~22]{KantorSeress2001} in its
previously used $\mathbf F_2$ form. The alternatives other
than the full general linear group preserve a nondegenerate
alternating form. The actual natural $\GL_8$ block has no
invariant bilinear form or functional, while fixing its
complement. It would therefore lie in the radical of such
a form, a contradiction. Thus the image on $K$ is $\GL(K)$.

There is nothing further in the odd case. In the even case,
the translation kernel is either zero or all of $K$.
If it were zero, the proved vanishing $H^1(\GL(K),K)=0$
would give a fixed lift of the one-dimensional quotient.
This contradicts the fact that every nonzero invariant
subspace contains $K$, whose dimension exceeds one.
Thus the full pointwise stabilizer occurs.

All the constructions are geometric and invariant under
extension of algebraically closed fields. For arithmetic
descent in the even case, the trace map to
$H^1_{\et}(E_{\bar k},\mathbf F_2)$ is defined over $k$.
This one-dimensional $\mathbf F_2$ space has trivial
arithmetic action, since $\GL_1(\mathbf F_2)=1$.
Hence the arithmetic image fixes the same nonzero covector.
The geometric image already fills its entire stabilizer,
so arithmetic equality follows. The odd case is immediate.
\end{proof}

\begin{remark}[Scope of the strip result]
\label{rem:mixed-strip-scope-candidate}
\status{Proved under stated hypotheses}
The first factor degree in this theorem is exactly three.
It is not an assertion about all balanced systems
$\mathcal O_E(3dO)\boxtimes\mathcal O_{E_0}(3dO)$.
The common-grid finite part has dimension $3b-3$, not
$3b-1$. The two missing directions are recovered in the
dimension argument by the full relative degree-one block
and the regular constant-body projections, rather than by
an identification of special and generic fibres.

The retained supporting program is
\path{agents/monodromy/code/mixed_strip_grid_checks.py}, with output
\path{agents/monodromy/outputs/mixed_strip_grid_checks.json}.
It checks the graph quotients, triple-supported component spans,
parity-dependent body images and recovered dimensions for
$4\le b\le32$, and the zero coinvariants of the three-point
label heart. It is finite supporting evidence; the universal
conclusions follow from the proofs above. The author run took
$0.029003258$ seconds; independent root and reduction/descent
runs took $0.030066450$ and $0.032449207$ seconds, preserving
that author output. All four complete geometric and group proofs
passed separate root, reduction/descent and obstructions/crystals
audits.
\end{remark}


% Promoted after complete root and reduction/descent proof audits; candidate retained.
\subsection{A cyclic norm and common partial normalizations in mixed degree three}
\label{subsec:mixed-degree-three-norm-candidate}
\status{Proved and independently checked. Complete proof audits and
independent exact graph reruns passed.}

Work geometrically in characteristic two. Let $E$ be ordinary,
let $\psi:E'_0\to E_0$ be a cyclic etale isogeny of degree three
of supersingular elliptic curves, and put $\phi=1_E\times\psi$.
Its kernel is $K=\{0,g,-g\}$. Norms of source divisors of type
$(3,9)$ have target type $(9,9)$. Picard twists with trivial
norm are permitted throughout. The target relative unit rank
is $80$ and its invariant covector space is zero, by the proved
rank and nonordinary-covector theorems. The strip theorem gives
the full source relative image $\GL_{26}(\mathbf F_2)$.

\begin{lemma}[A cubic source with separated translated branch sets]
\label{lem:mixed-d3-cubic-norm-source-candidate}
\status{Proved under stated hypotheses}
\textup{[Proved; complete independent proof audits passed.]}
A general smooth sharp source curve $C_1$ of type $(3,3)$ has
the following properties. Its norm $D$ has exactly eighteen
ordinary nodes and no other singularities. Its degree-three
projection $f:C_1\to E'_0$ has only double-plus-simple ramified
fibres, and its branch divisor and its two translates by $g$
and $-g$ have disjoint supports.
\end{lemma}
\begin{proof}
The double-plus-simple assertion is
Lemma~\ref{lem:mixed-strip-lifted-triple-candidate}.
The simultaneous ramification conditions at $(P,Q)$ and
$(R,Q+g)$ are four independent linear conditions on the
$(3,3)$ section: its value and its first-factor derivative
vanish at each point. Independence follows by first using the
two independent evaluations at $Q,Q+g$ on the second cubic,
and then the two first jets on the first cubic at $P$ and $R$
separately. The three parameters $P,R,Q$ therefore give a locus
of dimension at most seven in the eight-dimensional section
space. The same holds with $-g$. Thus the three translated
branch supports are disjoint for general $C_1$.

Here is a direct check of norm nodality in this small system.
In the plane cubic model of $E'_0$, send $Q$ to the chord
joining $Q$ and $Q+g$. This defines a morphism to the dual plane,
with pullback of its line bundle of degree six: its coordinates
are the cross products of the two degree-three coordinate vectors,
and they have no common zero because $g\ne0$.
The chord morphism is generically injective. If a line gave two
different pairs of this kind, its three intersections with the
cubic would be the complete orbit $Q+K$. Such an orbit is
collinear exactly when $3Q=0$, so these exceptional chords are
finite. Separability requires a separate check in characteristic
 two. If the differential of the chord map vanished at a general
 $Q$, differentiate both incidences $\ell(Q)(Q)=0$ and
 $\ell(Q)(Q+g)=0$. The derivative of its line coefficients is
 proportional to those coefficients, so the two equations force
 the chord to contain both tangent directions. Translation by
 $g$ has invertible differential. The same line would therefore
 be tangent at the two distinct cubic points, giving total
 intersection multiplicity at least four, which is impossible
 for a line and a cubic. Thus the chord map is separable.
 Together with generic injectivity this gives degree one onto
 its image. That image is consequently an integral sextic.

An invertible coefficient matrix for $C_1$ moves the first
elliptic cubic by an arbitrary projective transformation in
this same dual plane. A general transform meets the sextic
transversely in eighteen distinct smooth points and avoids its
finitely many exceptional chords. This transversality uses only
the separable transitive action of $\mathrm{PGL}_3$ on the plane:
the incidence of a smooth intersection point and coincident
tangent lines imposes one extra condition on a projective
transformation. No separability of a Gauss map is asserted.
The resulting intersections are precisely the pairs identified
by $\phi$, and transversality says that their image branches
have distinct tangents. Full three-point orbits were excluded.
The source is smooth on a dense open, and $\phi$ is etale, so
there are no other singularities. Intersecting these dense
opens with the sharp open proves all assertions simultaneously.
\end{proof}

\begin{lemma}[The self-node cover of the norm family]
\label{lem:mixed-d3-self-node-symmetric-candidate}
\status{Proved under stated hypotheses}
\textup{[Proved; complete independent proof audits passed.]}
For a general smooth sharp source $C$ of type $(3,9)$, including
each fixed Picard twist, its norm is nodal with $54$ nodes and
normalization $C$. The permutation image on these nodes, as
$C$ varies, is the full symmetric group $S_{54}$.
\end{lemma}
\begin{proof}
The additional double points are parametrized by
\[
                   Z_C=C\cap(C+g).
\]
At a nodal norm image each unordered pair has exactly one such
orientation, because $g$ has order three. Its length is
$2\cdot3\cdot9=54$.
Nonemptiness of the nodal open can already be seen from the
preceding lemma. Add six distinct ordinary horizontal fibres
to $C_1$, avoiding its branch points, its translated branch
points and its norm nodes, and with distinct target fibres.
The source has only its eighteen ordinary component nodes;
its norm has these nodes and fifty-four further ordinary nodes,
all disjoint. The additional nodes and the corresponding
intersection scheme $Z_C$ are etale in a neighbourhood of this
coefficient point. A general deformation there smooths the
source nodes and their images, and is sharp by the strip rank
theorem. It gives the required smooth nodal norm open. For a
specified source Picard twist the sum of the six fibre points
can be prescribed; its ordered parameter space is an open of
$(E'_0)^5$, so the same finitely many avoidance conditions can
be imposed. A first-factor three-torsion twist is placed on
$C_1$ itself; its degree-three plane embedding is simply a
translate, and its norm twist is trivial. Thus this construction
also treats every first-factor twist permitted in the statement.

The incidence of $(C,P)$ with $P,P-g\in C$ is an irreducible
projective bundle over the source surface: the two value
conditions are independent. The incidence of two different
nodes is likewise irreducible. It imposes four independent
value conditions at $P,P-g,R,R-g$, with all four points
distinct. If the two first coordinates differ, separate
those two first-coordinate values and then the two second
values within each pair. If the first coordinates coincide,
the four second points are distinct and the degree-nine line
bundle separates their values. Removing the diagonals, triple
orbits and the bad coefficient loci leaves dense opens of
these irreducible bundles. The node image is therefore doubly
transitive.

We verify an isolated transposition, including the small
first-factor degree. At a general first-factor point the
degree-three bundle separates second jets; its exceptional
flex set is finite. The degree-nine second-factor bundle
separates second jets at two prescribed distinct points
(total length six). Thus one can prescribe at $P,P-g$
two nonzero equal tangent directions and distinct quadratic
terms along that direction. The translated curves have
contact exactly two. The marked-contact incidence, with its
common tangent direction, is a projective bundle: the two
values and the two directional derivatives are four
independent linear conditions. Its point-and-direction base
has dimension three, so its image is a divisor. To justify
generic finiteness here, a translation-invariant source section
is an eigenvector for the order-three action. Its three
eigenspaces each have dimension nine, by descent of the
degree-nine second-factor bundle to degree three. A contact
linear system has projective dimension twenty-two, and cannot
be contained in this union of three eight-dimensional spaces.
For a smooth irreducible noninvariant source the translated
intersection is finite; the generic source has these properties
by the contact-system smoothness argument below.

For clarity, the following counts ensure that its generic
member has only this one contact. For two disjoint translated
pairs whose four second coordinates are distinct, the
degree-nine bundle separates their eight first jets. Hence
the two contacts impose eight independent conditions, against
six point-and-direction parameters. If the second coordinates
overlap, the point base has dimension at most three. With three
distinct second coordinates the eight conditions have rank
at least seven. At the shared coordinate, the two values
leave a one-dimensional first-factor section space; its two
directional derivatives lose at most one further condition.
With just two second coordinates one has a rectangle. If the
two directions have nonzero second components, the vertical
derivative values supply both conditions and the rank is eight.
If both directions are horizontal, the rank is six, but both
directions are fixed, leaving only three point parameters.
If exactly one is horizontal, rank can fall to six only when
the residual degree-three section has a double zero at that
first point. This is the relation
$\mathcal O_E(3O)=\mathcal O_E(2P+R)$ (with the appropriate
Picard twist); it removes one first-point parameter, and the
horizontal direction removes another. Thus this stratum also
has codimension at least two in the coefficient space.
Overlapping source pairs are a three-point orbit. Two contacts
there impose the three values and three directional derivatives,
of rank six, against three parameters. These counts show that
two contacts do not occur at the generic contact divisor.

The same jet argument excludes a singular source or a further
triple orbit at its generic point. Here are the potentially
exceptional cases. Away from the two prescribed second fibres,
sections vanishing to first order at the prescribed pair still
separate surface first jets, since the second-factor conditions
have total length at most six. On either prescribed second
fibre, work on the dense open where the prescribed common
tangent has nonzero second component. Its derivative equation
then determines the vertical-derivative value from the
first derivative; the first-factor value section is exactly in
$H^0(\mathcal O_E(3O)(-P))$. Its first jets at another point
have rank two except at the finite points where this bundle
equals $\mathcal O_E(2R)$; the vertical derivative value is
independent. The rank is therefore three off a finite set and
at least two on it. Singular-point incidence is consequently
proper in the contact system, and the prescribed two gradients
are nonzero on an open.
An additional triple orbit imposes three value conditions.
If its second orbit is disjoint from the prescribed one, these
and the four contact conditions are independent (second-factor
length at most seven). If the second orbits agree, the rank is
at least six against at most four parameters. The only drop
to five requires a horizontal common tangent and
$\mathcal O_E(3O)=\mathcal O_E(2P+R)$; it leaves only two
parameters. A triple containing the contact pair itself has
five independent conditions against three parameters.
Thus none of these loci dominates the contact divisor.

At its generic point the length-two intersection cluster has
the form $z^2+uz+v$. The independent constant and linear jets
make $u,v$ independent coefficient parameters. As in
Lemma~\ref{lem:mixed-lifted-intersection-wreath-candidate},
the equation $w^2+w=v/u^2$ gives nontrivial inertia exchanging
these two points, while all the other points are etale.
This is one transposition in the doubly transitive image.
A doubly transitive permutation group containing a
transposition is $S_{54}$: its conjugates give all transpositions.
The dense sharp and smooth restrictions do not alter the image.
\end{proof}

\begin{lemma}[The normalization block and its gluing heart]
\label{lem:mixed-d3-norm-block-shears-candidate}
\status{Proved under stated hypotheses}
\textup{[Proved; complete independent proof audits passed.]}
On the nodal norm family there is an exact sequence
\[
       0\longrightarrow T\longrightarrow V\longrightarrow U
          \longrightarrow0,
       \qquad \dim T=54,\quad\dim U=26.
\]
The image contains an actual natural $\GL_{26}$ block fixing a
complement, has full $S_{54}$ permutation image on $T$, and
contains all $\operatorname{Hom}(U,T^0)$ translations, where
$T^0\subset\mathbf F_2^{54}$ is the sum-zero heart. These
assertions hold for every source Picard twist with trivial norm.
\end{lemma}
\begin{proof}
The normalization sequence gives the graph kernel $T$, one
unit direction for each self-node. The odd isogeny induces an
isomorphism on ambient etale two-cohomology in degree one,
so quotienting by it gives the displayed relative sequence.
The source image on $U$ is $\GL_{26}$ by the strip theorem,
which applies to every Picard twist by translation.
The joint image on $T\oplus U$ is $S_{54}\times\GL_{26}$.
Indeed $\GL_{26}(\mathbf F_2)$ is simple and is not $A_{54}$
(its order is at least $2^{650}$ whereas
$54!<64^{54}=2^{324}$), so Goursat gives the product.
Kill the permutation action. The remaining additive kernel
is a submodule of $\operatorname{Hom}(U,T)$, a direct sum
of natural-dual modules. The proved vanishing of their $H^1$
and of the $H^1$ of every quotient supplies an actual
block-diagonal $\GL_{26}$ by the cocycle argument.

It remains to prove the gluing assertion for this self-node
family; it is not inferred from independent intersection
systems. Fix a curve $C_1$ from the first lemma and use its
maximal elementary two-torsor, with group
$G=H^1_{\et}(C_1,\mathbf F_2)^*$.
Each of the three lifted covers above $q,q+g,q-g$ has image
$H_0=G^2\rtimes S_3$ by
Lemma~\ref{lem:mixed-strip-lifted-triple-candidate}.
Their branch supports are disjoint. Inertia at a branch point
of one factor is a labelled transposition in that factor alone.
Its conjugates generate the entire factor, by the same lemma.
Hence their joint image is the full product $H_0^3$.

Now take a source boundary
$C=C_1\cup\bigcup_{j=1}^6(E\times\{q_j\})$ with the sum
of the $q_j$ prescribed to give the chosen source twist.
The same independence holds for all eighteen shifted covers.
On the ordered sum fibre one may vary one $q_j$ through one
branch point and vary a second point to maintain the sum,
keeping that second point and every other shifted point
unramified. Each projection to a chosen $q_j$ has geometrically
connected fibres, so its individual lifted-cover image remains
$H_0$. Thus single-factor inertia is still present for
each factor. Its normal conjugates generate that full factor,
and the joint image remains $H_0^{18}$.

Choose a pure label translation in a shifted factor
$q_j+g$ and the identity in every other factor, in particular
in the unshifted factor at $q_j$. The source cohomology action
is the identity: $C_1$ is fixed, the ordinary components are
constant, and their source-node gluings depend on the unshifted
factors. All norm-node permutations are also the identity.
The norm gluing, however, changes at the three nodes in the
chosen shifted block by their torsor evaluations. The three
labels sum to zero. They may be chosen in the annihilator of
the ambient one-dimensional class in $H^1(C_1)$; this eliminates
ambient-body columns and still permits all relative evaluations.
On the relative eight-dimensional part
of $H^1(C_1)$ these evaluations give every map to the
two-dimensional sum-zero space of that triple. This is
nonzero, while the other norm-node coordinates are fixed.
The equality follows directly by continuing the chosen node
gluings of the etale torsor, as in the earlier boundary proof.

Let $J$ denote the translation kernel on the entire norm
family. As a $\GL_{26}$ submodule of $U^*\otimes T$, it is
$U^*\otimes T'$ for a subspace $T'\subset T$.
The preceding actual translation puts a nonzero pair vector
of one triple in $T'$. The full symmetric permutation group
on $T$ preserves $T'$, so its orbit spans all $T^0$.
Consequently $\operatorname{Hom}(U,T^0)\subset J$.
The boundary used here lies in the full source and target
unit opens; its additional norm nodes are disjoint and etale.
Thus these monodromy elements belong to the same extended
norm-family local system as the smooth-source elements.
\end{proof}

\begin{lemma}[Common partial normalizations force irreducibility]
\label{lem:mixed-d3-common-partial-normalizations-candidate}
\status{Proved under stated hypotheses}
\textup{[Proved; complete independent proof audits passed.]}
The geometric relative representation of rank $80$ in the
balanced system of type $(9,9)$ is irreducible.
\end{lemma}
\begin{proof}
Let $S$ be a proper nonzero invariant subspace. On any norm
family, a nonzero projection of $S$ to $U$ would give the whole
natural $26$-dimensional quotient and, by the gluing translations,
the $53$-dimensional heart $T^0$. It would have dimension at
least $79$. If proper it would be a hyperplane, giving a
nonzero invariant covector, contrary to the proved odd-degree
covector theorem. Thus $S$ is contained in $T$ on every such
norm family, including all the allowed source twists.

Fix $C_1$ and its norm $D$ as in the first lemma. Choose six
distinct target ordinary fibres, with sum zero, whose inverse
images on $C_1$ are unramified and avoid all self-nodes of $D$.
Let $Y$ be their union with $D$. For each of the six fibres
choose independently one of its three source lifts, and call
the resulting source union $C_{\mathbf r}$. All $3^6$ unions
have norm exactly the same $Y$. Their line bundles have type
$(3,9)$, with a possible Picard twist in the kernel of norm,
so the preceding lemma applies to each.

The source two-rank is $9+6+12=27$. The target two-rank is
$9+6+18+48=81$, hence relative rank $80$. In particular $Y$
is in the maximal-rank open. This comparison involves no
rank drop at the target, and $S_Y$ has its generic dimension.
Each partial normalization $C_{\mathbf r}\to Y$ separates
the eighteen self-nodes of $D$ and six of the nine cross-nodes
on each ordinary fibre. Its graph kernel $T_{\mathbf r}$ has
dimension $18+6\cdot6=54$. The additional nodes remain
ordinary and disjoint locally, so these kernels extend as
unit systems in the corresponding norm families. Thus
$S_Y\subset T_{\mathbf r}$ for every choice.

Use graph cochains with one normalization vertex for $C_1$,
six ordinary-fibre vertices, eighteen loops, and nine parallel
edges to each other vertex. For one fibre, the three source
choices select three disjoint triples of its nine edges.
Membership in the corresponding partial-normalization kernel
says that the cochain is constant on the selected triple,
modulo the full nine-edge star coboundary. Membership for all
three choices says it is constant separately on all three
triples, a two-dimensional space modulo that star. Hence
\[
                \dim\bigcap_{\mathbf r}T_{\mathbf r}
                         =18+6\cdot2=30.
\]
An invariant subspace of the permutation module
$\mathbf F_2^{54}$ containing a nonconstant vector contains
the sum-zero heart: subtract a transposition of two unequal
coordinates, and then conjugate the resulting pair vector.
It therefore has dimension at least $53$. Since $S_Y$ has
dimension at most thirty, it must be the diagonal line in
each $T_{\mathbf r}$ if it is nonzero.

These lines differ. Modulo the nine-edge star, the sum of
the six unselected edges is the indicator of the selected
triple. Changing one source lift changes the diagonal by
the indicators of two distinct triples. This is neither
zero nor the full nine-edge star, and hence is nonzero in
graph cohomology. The two diagonal lines over $\mathbf F_2$
are distinct and have zero intersection. This contradicts
$S_Y\ne0$, proving irreducibility.
\end{proof}

\begin{theorem}[Full fixed-modulus mixed image in degree three]
\label{thm:mixed-every-fixed-modulus-degree-three-candidate}
\status{Proved under stated hypotheses}
\textup{[Proved; complete independent proof audits passed.]}
For every perfect field of characteristic two, every ordinary
elliptic curve $E$ and every supersingular elliptic curve $E_0$,
the full geometric and arithmetic relative images for
$|\mathcal O_E(9O)\boxtimes\mathcal O_{E_0}(9O)|$ are
\[
                         \GL_{80}(\mathbf F_2).
\]
\end{theorem}
\begin{proof}
Work geometrically and let $J$ be the normal subgroup generated
by transvections. Irreducibility of the entire image makes
the restriction to $J$ semisimple, by the invariant-socle
argument already proved. All nontrivial isotypic multiplicities
are one, and there is no fixed part. The actual natural
$26$-dimensional block lies in one simple summand. Every
rank-one map from that block to $T^0$ is present, so all of
$T^0$ lies in the same summand. Its dimension is at least
$26+53=79$. Any other simple summand would have dimension at
most one and would be fixed by all transvections, impossible.
Thus $J$ is irreducible.

The already used transvection classification
\cite[Theorem~3.6, printed p.~22]{KantorSeress2001} gives the
full general linear group or a group preserving a nondegenerate
alternating form. The actual natural $\GL_{26}$ block fixes
a complement and has no invariant bilinear form or functional,
so the latter alternative is impossible. This proves the
geometric statement. Every supersingular elliptic curve has
a cyclic etale three-isogeny after algebraic closure, and
the argument is unchanged by geometric base extension or
Picard translation. The arithmetic image contains the full
geometric general linear group and is therefore equal to it.
\end{proof}

\begin{remark}[Scope and exact graph supplement]
\label{rem:mixed-d3-norm-scope-candidate}
\status{Proved under stated hypotheses}
\textup{[Proved; complete independent proof audits passed.]}
This is a degree-three claim only; higher odd balanced degrees
remain separate. The independent root graph program
\path{code/mixed_degree_three_partial_normalizations.py}
checks all $729$ partial-normalization kernels, their
$53$-dimensional hearts, their $30$-dimensional common
intersection and the $729$ distinct diagonal classes.
Its author execution took $0.060238889$ seconds. Independent
reduction/descent and monodromy executions took $0.068377761$
and $0.062999599$ seconds, respectively, preserving the author
output. The independent wrappers and outputs are retained.
It does not establish norm nodality, the self-node symmetric
image or the gluing translations; those are the new geometric
proofs above. The full mathematical proofs were independently
rederived by the root and reduction/descent agents before promotion;
the earlier complete strip and transvection-classification audits
are reused with their scope recorded.
\end{remark}


% Promoted after complete root and reduction/descent proof, source and code audits.
\subsection{All odd balanced degrees on a fixed mixed elliptic product}
\label{subsec:mixed-all-odd-norm-candidate}
\status{Proved and independently checked.}

Work first over an algebraically closed field of characteristic two.
Let $E$ be ordinary and let $E_0$ be supersingular. Fix an odd integer
$d\geq3$ and a cyclic etale isogeny
$\psi:E'_0\longrightarrow E_0$ of degree $d$, with kernel $K$.
Put $\phi=1_E\times\psi$ and
\[
 \mathcal H=(K\setminus\{0\})/\{h\sim-h\},\qquad
 t=\frac{d-1}{2},\quad m=3d-3,\quad N=18d.
\]
The source systems have type $(3,3d)$, with any fixed Picard twist
whose norm is trivial. Their norms have target type $(3d,3d)$.
Write
\[
 u=9d-1,\qquad \delta=9d(d-1)=tN,\qquad r=u+\delta=9d^2-1.
\]
The proved strip theorem gives source relative image $\GL_u(\mathbf F_2)$.
The proved rank theorem gives target generic relative rank $r$.

\begin{lemma}[A small source with all translated pairs separated]
\label{lem:mixed-odd-all-translated-small-source-candidate}
\status{Proved under stated hypotheses}
\textup{[Proved; complete independent proof audits passed.]}
A general smooth sharp curve $C_1$ of type $(3,3)$ has the following
properties. Its norm $D$ has eighteen ordinary nodes of each type
$[h]\in\mathcal H$, with no triple points or other singularities.
Its degree-three projection to $E'_0$ has double-plus-simple ramified
fibres, and the branch supports of all its translates by $K$ are
pairwise disjoint.
\end{lemma}
\begin{proof}
The double-plus-simple assertion is
Lemma~\ref{lem:mixed-strip-lifted-triple-candidate}.
For any nonzero $h\in K$, simultaneous ramification at
$(P,Q)$ and $(R,Q+h)$ imposes four independent linear conditions
on a $(3,3)$ section: its value and its first-factor derivative
vanish at both points. Separate the two second-factor values and
then the two first jets at the respective first-factor points.
There are only three point parameters. The simultaneous-ramification
locus therefore has dimension at most seven in the section space
$\mathbf P^8$. Applying this to the finitely many nonzero differences
in $K$ proves the branch-support assertion on a dense open.

For each nonzero $h$, the chord map
\[
 Q\longmapsto \overline{Q\,(Q+h)}
\]
from the plane cubic $E'_0$ to its dual plane has line-bundle degree
six. Its cross-product coordinates have no common zero, since
$Q\ne Q+h$. Two distinct translated pairs on the same line must
share a point; otherwise that line meets a cubic in four distinct
points. The overlapping pairs have three consecutive points
$Q,Q+h,Q+2h$, and their collinearity requires $3(Q+h)=0$.
There are only finitely many such points. This also treats order-three
$h$, when the three points form an entire orbit. Thus the chord map
is generically injective. It is separable: a vanishing differential
would, by differentiating both incidence equations, make the chord
tangent at both distinct endpoints, giving line-intersection
multiplicity at least four. Its image is consequently an integral
sextic, and it has degree one onto that image.

An invertible coefficient matrix for $C_1$ moves the first elliptic
cubic by an arbitrary projective transformation in this dual plane.
A general transform meets each of the finitely many chord sextics
transversely in eighteen distinct smooth points, avoiding their
finitely many exceptional chords. For completeness, transversality
here follows from the separable transitive action of
$\mathrm{PGL}_3$: matching two smooth points is a codimension-two
condition, and matching their tangent lines imposes one further
condition. Varying the two points leaves a proper tangency locus.
No Gauss-map separability is used.

A triple of distinct points $Q,Q+h,Q+k$ on the second cubic can
lie on one line only if $3Q+h+k=0$. There are finitely many such
triples as $h,k$ vary in $K$. The same general projective transform
avoids their lines. Hence no norm point has three source preimages.
The transverse chord intersections give precisely the eighteen
ordinary double points of each type. The source is smooth on a
dense open and $\phi$ is etale, so there are no other singularities.
Intersecting the finitely many dense opens with the nonempty sharp
open completes the proof.
\end{proof}

\begin{lemma}[Independent permutation groups for all node types]
\label{lem:mixed-odd-node-product-candidate}
\status{Proved under stated hypotheses}
\textup{[Proved; complete independent proof audits passed.]}
For a general smooth sharp source $C$ of type $(3,3d)$, with any
fixed norm-trivial Picard twist, its norm is nodal with $N$ nodes
of each type $[h]\in\mathcal H$. The joint permutation image on
these nodes is
\[
                         \prod_{[h]\in\mathcal H}S_N.
\]
\end{lemma}
\begin{proof}
Choose one representative $h$ for each type. Its nodes are
parametrized by $Z_h=C\cap(C+h)$, of length $2\cdot3\cdot3d=N$.
Each unordered source pair has exactly one such orientation because
$h\ne-h$. Start with $C_1$ from the preceding lemma and add $m$
ordinary horizontal fibres. Choose their points with any prescribed
sum, with distinct images under $\psi$, and away from all translated
branch points and norm nodes. The ordered sum fibre is an open of
$(E'_0)^{m-1}$; the finitely many forbidden conditions are proper.
The source has $3m$ ordinary component nodes. Its norm has these
nodes and, in addition, eighteen nodes of each type from $D$ and
six per ordinary fibre of each type. Thus each additional type has
$18+6m=N$ nodes. All these points are disjoint and ordinary.
They remain etale locally in the coefficient space. A general
deformation smooths the source component nodes and their images,
leaving the additional norm nodes. Sharpness is open and holds
at the boundary by the strip rank calculation. This proves
nonemptiness of the smooth nodal norm open. First-factor twists
may be placed on $C_1$; every degree-three line bundle on an
elliptic curve is a translate of the chosen one.

For one fixed type, the incidence of $(C,P)$ with $P,P-h\in C$
is an irreducible projective bundle. The incidence of two different
nodes is also irreducible: the four values at
$P,P-h,R,R-h$ are independent. If the first coordinates differ,
separate their two values first and then the two second values
within each pair. If they coincide, the four second points are
distinct and the degree-$3d$ bundle separates them. Thus the
permutation group of every type is doubly transitive.

We check that a transposition can occur in just one type. Mark
$P,P-h$ and a common nonzero tangent direction. The two values
and two directional derivatives are four independent linear
conditions. At a nonflex first-coordinate point, the first cubic
separates second jets, and the second bundle separates second
jets at the two points (length six). Consequently two smooth
branches with contact exactly two can be prescribed.

Here are the incidence bounds excluding every additional contact.
For pairs $(P,P-h)$ and $(R,R-k)$ with four distinct second
coordinates, the degree-$3d$ bundle separates their eight first
jets. There are eight independent conditions and six
point-and-direction parameters. With three second coordinates,
the point base has dimension at most three, and the eight
conditions have rank at least seven: at the shared coordinate
the two values leave one first-factor section, and its two
directional derivatives lose at most one condition. Two second
coordinates can occur only when $[h]=[k]$. This is exactly the
rectangle calculation of
Lemma~\ref{lem:mixed-d3-self-node-symmetric-candidate}:
the rank is eight when both directions have nonzero second
components; it is six when both are horizontal, but then there
are only three point parameters; with exactly one horizontal
direction, rank six additionally requires the relation
$A=\mathcal O_E(2P+R)$ for the degree-three first-factor bundle
$A$, removing another parameter. If the two translated pairs
share an actual source point, smoothness forces their directions
to agree there. Their three values and three directional
derivatives have rank six, against only three parameters.
Thus a second contact never dominates the first contact divisor.
For distinct types the two-coordinate rectangle does not arise.

The source-smoothness calculation in that same lemma uses only
second-factor interpolation of length at most six, so it holds
unchanged for degree $3d\geq9$. More explicitly, restrict to the
dense open where the common direction has nonzero second
component. On a prescribed second fibre the remaining value
section is $H^0(A(-P))$, with rank-two first jets except at the
finite points where $A(-P)=\mathcal O_E(2R)$. Its independent
second-factor derivative gives rank three off that finite set
and at least two on it. Away from the two prescribed fibres,
the remaining sections separate surface first jets. All the
singular-point incidences are therefore proper.

For an additional triple with differences in $K$, the three
second points are distinct. If disjoint from the contact's second
points, its three values and the four contact conditions are
independent, by interpolation of total length at most seven.
If just one second coordinate overlaps, the union has four
second fibres. The three-value calculation has rank seven
generically and at least six everywhere. Coincident first
points leave only three parameters; a horizontal residual-jet
exception also lowers the parameter count. Thus this case has
rank at least six against at most four parameters.
With two overlapping second coordinates the rank is at least
six against at most four parameters. A rank-five exception
requires a horizontal direction and $A=\mathcal O_E(2P+R)$,
leaving only two parameters. A triple containing the contact
pair itself imposes five independent conditions against three
parameters. These are the bounds in the cited lemma; their proof
uses distinct second coordinates and not the order of $h$ or
the third difference. There are only finitely many difference
choices. No further triple dominates the contact incidence.

The generic source in this incidence is therefore smooth.
It is connected, since Kunneth gives
$H^1(E\times E'_0,L_{\mathrm{source}}^{-1})=0$; the divisor
sequence then gives $H^0(\mathcal O_C)=k$. Smoothness and
connectedness make it irreducible. It is not invariant under
translation by $h$.
Indeed, if $e$ is the odd order of $h$, its translation
eigenspaces have dimension $9d/e\leq3d$, whereas a contact
linear system has projective dimension $9d-5>3d-1$.
For a smooth irreducible noninvariant source, intersection with
its translate is finite. The marked-contact image is therefore
a divisor. At its generic point just one type has one length-two
cluster, of the form $z^2+uz+v$ with independent constant and
linear deformation parameters. The Artin--Schreier equation
$w^2+w=v/u^2$ gives nontrivial inertia interchanging that pair;
every other point of every type is etale. This is a transposition
in that factor alone. Its conjugates under the doubly transitive
coordinate image generate the entire symmetric factor. Repeating
for every type proves the full direct product assertion.
\end{proof}

\begin{lemma}[The large normalization block and every gluing heart]
\label{lem:mixed-odd-normalization-block-hearts-candidate}
\status{Proved under stated hypotheses}
\textup{[Proved; complete independent proof audits passed.]}
On the nodal norm family the relative unit system has an exact
sequence
\[
 0\longrightarrow T\longrightarrow V\longrightarrow U
 \longrightarrow0,
 \qquad T=\bigoplus_{[h]\in\mathcal H}T_h,\quad
 T_h=\mathbf F_2^N,\quad\dim U=u.
\]
Its image contains an actual natural $\GL_u$ block fixing $T$,
has the full product of symmetric groups on $T$, and contains
\[
             \operatorname{Hom}\left(U,\bigoplus_hT_h^0\right),
 \qquad T_h^0=\ker\left(\sum:\mathbf F_2^N\to\mathbf F_2\right).
\]
These assertions hold for every fixed norm-trivial source twist.
\end{lemma}
\begin{proof}
The normalization sequence gives $T$. The odd isogeny induces
an isomorphism on ambient $H^1_{\et}(-,\mathbf F_2)$, so the
source relative quotient is the strip system of dimension $u$.
Its image is $\GL_u$. This simple group is not $A_N$:
\[
 |\GL_u(\mathbf F_2)|>2^{u(u-1)}>2^{54d^2}
       \geq(18d)^{18d}>N!.
\]
Here $18d\leq2^{3d}$ for $d\geq3$, and
$(9d-1)(9d-2)>54d^2$. The nonabelian simple factors of the
node product are precisely copies of $A_N$. Goursat therefore
gives the full joint quotient
$\GL_u\times\prod_hS_N$. Kill the permutations. The remaining
translation kernel is a submodule of $U^*\otimes T$; the proved
vanishing of natural-dual $H^1$ and of its quotient modules
gives an actual block section by the earlier cocycle argument.

For the translations, fix $C_1$ from the first lemma and its
maximal elementary two-torsor, with group
$G=H^1_{\et}(C_1,\mathbf F_2)^*$. Each lifted three-point cover
over $q+k$, $k\in K$, has image $H_0=G^2\rtimes S_3$ by the
strip theorem. Their branch supports are pairwise disjoint.
Single-factor branch inertia is a labelled transposition, and
its normal conjugates generate that entire $H_0$ factor.
Thus the joint image is $H_0^d$.

Now vary the $m$ source ordinary fibres subject to their prescribed
sum. Each projection to $q_j$ has geometrically connected fibres,
so each individual lifted cover retains its image. Vary one
$q_j$ through a chosen shifted branch point, compensating with
a second coordinate while keeping every other shifted point
unramified. The same normal-closure argument gives the full
joint image $H_0^{md}$ on all shifted covers.

Choose a pure label translation in a shifted factor $q_j+h$,
with identity in all other factors, especially the unshifted
source-node factor at $q_j$. Choose its three sum-zero labels
in the annihilator of the ambient class in $H^1(C_1)$.
The entire source relative cohomology action and all node
permutations are identity. The norm gluing changes only at
the corresponding three type-$[h]$ nodes, and its evaluations
give every map from the relative eight-dimensional part of
$H^1(C_1)$ to that triple's two-dimensional sum-zero space.
This is the explicit torsor-gluing calculation of
Lemma~\ref{lem:mixed-d3-norm-block-shears-candidate}; all other
shifted factors remain fixed, so no other type is changed.

The full translation kernel $J$ is a $\GL_u$ submodule of
$U^*\otimes T$, hence $J=U^*\otimes T'$ for a subspace $T'$.
The exhibited translation puts a pair vector of this triple
in $T'$. The $S_N$ action in that type spans $T_h^0$ from that
vector. This works for every $h$, giving the stated translations.
The boundary source and target have maximal relative rank;
all additional norm nodes are disjoint and etale there.
The finite unit systems therefore extend to this boundary, so
these are actual elements of the same norm-family image.
\end{proof}

\begin{lemma}[Two common-graph assertions for any odd kernel]
\label{lem:mixed-odd-common-graph-candidate}
\status{Proved under stated hypotheses}
\textup{[Proved; complete independent proof audits passed.]}
Let $K$ be a finite abelian group of odd order $d>1$, and
$\mathcal H=(K\setminus0)/\{\pm1\}$. Form a graph with one
central vertex, $m\geq1$ further vertices, eighteen central
loops of each type $h\in\mathcal H$, and $3d$ parallel edges
to each further vertex, labelled $(k,a)$ with $k\in K$ and
$1\leq a\leq3$. Put $\mathcal G=H^1(\Gamma,\mathbf F_2)$.
For $\mathbf r=(r_1,\ldots,r_m)\in K^m$, let $T_{\mathbf r,h}$
be spanned by the eighteen type-$h$ loops and the six edges
$k=r_j\pm h$ at every vertex. Write $T_{\mathbf r,h}^0$ for
its sum-zero hyperplane and $T_{\mathbf r}=\bigoplus_hT_{\mathbf r,h}$.
Then:
\begin{enumerate}
\item the spaces $T_{\mathbf r,h}^0$, for all $\mathbf r,h$,
span all of $\mathcal G$;
\item if $S\subset\bigcap_{\mathbf r}T_{\mathbf r}$ is invariant
under the independent full permutation groups on the bases of
each $T_{\mathbf r,h}$, then $S=0$.
\end{enumerate}
\end{lemma}
\begin{proof}
All the indicated edge vectors are independent in their
$T_{\mathbf r}$: a nonzero star coboundary uses the selected
three edges $k=r_j$, which do not occur in that space.
A cochain belongs to $T_{\mathbf r}$ exactly when its three
values on every selected triple are constant; a star removes
that constant.

For the first assertion, let $\lambda\in\mathcal G^*$ kill every
heart. Describe it by coefficients on loops and edges. Its
edge coefficients sum to zero at every noncentral vertex,
because it annihilates that vertex's star coboundary.
On the basis of $T_{\mathbf r,h}$ its coefficients must all
equal one constant $c_{\mathbf r,h}$. The fixed central loops
anchor this constant independently of $\mathbf r$, so denote
it by $c_h$. Fix one vertex and one edge $(k,a)$. Varying its
selected lift $r\ne k$ gives
\[
                    \lambda_{k,a}=c_{[k-r]}.
\]
Every nonzero difference occurs; consequently all $c_h$ have
one value $c$, and every edge coefficient is $c$. The vertex
star relation is $3d\,c=c=0$. Thus every loop and edge
coefficient is zero. Duality proves the spanning assertion.

For the second assertion, suppose some $v\in S$ has a
nonconstant component in a type block for a choice $\mathbf r$.
Subtracting a transposition and then conjugating gives the
entire heart $T_{\mathbf r,h}^0\subset S$. Take its pair vector
on two of the three edges at one unselected lift $k=r_j+h$.
Change that vertex's selected lift to $k$. The pair vector
has values $(1,1,0)$ on this newly selected triple, and hence
does not belong to the new $T_{\mathbf r'}$, a contradiction.
Therefore every element of $S$ is in the span of the type
diagonals for every $\mathbf r$.

Fix such an element. Its loop values give constants $c_h$
independent of $\mathbf r$. On one vertex's edge triples it
has constant values $v_k$ modulo the full star. Membership
in every diagonal span gives
\[
                   v_k+v_r=c_{[k-r]}\qquad(k\ne r).
\]
Set $f(k)=v_k+v_0$. Then $f(0)=0$, $f(h)=c_{[h]}$, and
$f(k)+f(r)=f(k-r)$. Taking $k=x+y$, $r=y$ shows that $f$
is a homomorphism $K\to\mathbf F_2$. It is zero because $K$
has odd order. All $c_h$ are zero, and the edge values at
every vertex are a star coboundary. Thus the original element
is zero in $\mathcal G$, proving $S=0$.
\end{proof}

\begin{proposition}[Irreducibility in every odd balanced degree]
\label{prop:mixed-odd-balanced-irreducibility-candidate}
\status{Proved under stated hypotheses}
\textup{[Proved; complete independent proof audits passed.]}
The geometric relative representation of rank $r=9d^2-1$
in the balanced system of type $(3d,3d)$ is irreducible.
\end{proposition}
\begin{proof}
Let $S$ be a nonzero proper global invariant subspace. On any
norm family its projection to $U$ is either zero or all of $U$.
In the latter case the translations force every $T_h^0$ into
$S$, so
\[
                         \dim S\geq u+\delta-t=r-t.
\]
In the former case $\dim S\leq\delta$. Since $u>t$, these
ranges are disjoint. Therefore either the projection is zero
on every norm family and every allowed source twist, or it
is surjective on every one of them.

Fix $C_1$ from the first lemma, its nodal norm $D$, and $m$
distinct target ordinary fibres whose points sum to zero,
avoiding the branch and norm-node loci. Let $Y$ be their union
with $D$. Independently choose any one of the $d$ source lifts
of each fibre. The resulting $C_{\mathbf r}$ all have norm
exactly $Y$; their source twists have trivial norm. Their
two-rank is
\[
                 9+m+2m=9d,
\]
whereas that of $Y$ is
\[
       9+m+9(d-1)+m(3d-1)=9d^2.
\]
After removing the ambient one-dimensional unit class these
are precisely the maximal ranks $u$ and $r$. The target point
$Y$ is therefore in the full-rank coefficient open, and $S_Y$
has the generic dimension. No rank-drop specialization is used.

The graph of $Y$ is exactly the graph in the preceding lemma:
the normalization is $C_1$ together with $m$ ordinary curves;
each node type contributes eighteen loops, and each ordinary
curve meets $D$ at $3d$ points, in $d$ triples. The graph kernel
of the partial normalization $C_{\mathbf r}\to Y$ is
$T_{\mathbf r}$. These kernels extend in the corresponding
norm families because the additional nodes are ordinary and
etale and both unit ranks are maximal. The global subsystem
$S$ restricts to the same fibre $S_Y$ in all these families.

If every projection is zero, then $S_Y\subset T_{\mathbf r}$
for all $\mathbf r$, with invariance under each product of
symmetric groups. Assertion two of the graph lemma gives
$S_Y=0$, a contradiction.

If every projection is surjective, $S_Y$ contains every
$T_{\mathbf r,h}^0$. Assertion one gives the whole graph
space $\mathcal G\subset S_Y$. For any choice $\mathbf r$,
the projection $S_Y\to U_{\mathbf r}$ is surjective, and its
kernel $T_{\mathbf r}$ is contained in $\mathcal G\subset S_Y$.
Exactness therefore gives $S_Y=V_Y$, again a contradiction. The geometric
representation is irreducible.
\end{proof}

\begin{theorem}[Full images in all odd balanced mixed degrees]
\label{thm:mixed-every-fixed-modulus-all-odd-candidate}
\status{Proved under stated hypotheses}
\textup{[Proved; complete independent proof audits passed.]}
Let $k$ be any perfect field of characteristic two, let $E/k$
be ordinary and $E_0/k$ supersingular. For every odd integer
$d\geq1$, the full geometric and arithmetic relative images
for
\[
       |\mathcal O_E(3dO)\boxtimes\mathcal O_{E_0}(3dO)|
\]
are both $\GL_{9d^2-1}(\mathbf F_2)$.
\end{theorem}
\begin{proof}
Degree one is the proved uniform fixed-modulus theorem.
For $d\geq3$, work geometrically and let $J$ be the normal
subgroup generated by transvections. Irreducibility of the
whole image makes its restriction to $J$ semisimple: the
nonzero $J$-socle is invariant under the full group and hence
is the entire representation.

The actual natural $\GL_u$ block is generated by transvections
and is contained in $J$. Its natural constituent occurs with
multiplicity one in $V=U\oplus T$, so $U$ lies in a single
simple $J$-summand. The rank-one gluing maps from $U$ to all
$T_h^0$ place all those hearts in the same summand. Its dimension
is at least $r-t>r/2$. Every conjugate is again a simple
$J$-summand of that dimension; two distinct such summands
would be disjoint, which is impossible. This summand is thus
invariant under the full irreducible image, and equals $V$.
Consequently $J$ is irreducible.

Apply the already audited characteristic-two transvection
classification \cite[Theorem~3.6, printed p.~22]{KantorSeress2001}.
The alternatives other than the full general linear group
preserve a nondegenerate alternating form. The actual natural
$\GL_u$ block fixes a complement and has no invariant
bilinear form or functional, so it would be in the radical
of such a form. This contradiction gives $J=\GL_r$.

Every odd cyclic isogeny used above exists after algebraic
closure of the given perfect field. Picard translations
identify the relevant source systems, and geometric base
extension preserves the image. The arithmetic image contains
the full geometric general linear group and is therefore
equal to it.
\end{proof}

\begin{remark}[Scope of this extension]
\label{rem:mixed-all-odd-norm-scope-candidate}
\status{Proved under stated hypotheses}
\textup{[Proved; complete independent proof audits passed.]}
This is an all-odd theorem for fixed ordinary times
supersingular elliptic products. It is not a general ambient
monodromy classification. The odd-kernel graph lemma is a
universal linear proof; any finite graph certificate is only
a supplement. No new source theorem is imported by analogy:
the input strip monodromy, natural-module cohomology and
transvection classification were already independently read
and audited, while the multiple-type geometric incidence and
common-graph steps have now passed complete independent root and
reduction/descent audits. The exact supporting program
\path{agents/monodromy/code/mixed_odd_common_graph_checks.py}
checks nine actual cochain cases, including cyclic kernels of orders
$3,5,7,9,15,21$ and the noncyclic group $(\mathbf Z/3)^2$.
Its author run took $0.036200414$ seconds; independent root and
reduction/descent runs took $0.050087592$ and $0.040483062$
seconds, preserving the author output. These finite checks do not
replace the universal proof or any geometric incidence argument.
\end{remark}


% Promoted after complete root and reduction/descent proof, source and code audits.
\subsection{All even balanced degrees on a fixed mixed elliptic product}
\label{subsec:mixed-all-even-wreath-candidate}
\status{Proved and independently checked.}

Work geometrically in characteristic two and write
$X=E\times E_0$, with $E$ ordinary and $E_0$ supersingular, and
$L=\mathcal O_E(3O)\boxtimes\mathcal O_{E_0}(3O)$.
Let $d\geq2$ be even and put
\[
 a=d-1,\qquad u=9a^2-1,\qquad N=18a,\qquad
 r=9d^2-2,\qquad \kappa=9d^2-4.
\]
The proved rank and exact-covector theorems give relative rank $r$
and exactly two invariant covectors on $|L^d|$. Their common
kernel is denoted $K_d$ and has dimension $\kappa$.
The preceding all-odd theorem supplies the full relative images
$\GL_u$ and $\GL_8$ in the component systems used below.

\begin{lemma}[Lifted intersections for unequal odd components]
\label{lem:mixed-even-unequal-lifted-wreath-candidate}
\status{Proved under stated hypotheses}
\textup{[Proved; complete independent proof audits passed.]}
Fix $\alpha\in\Pic^0(X)$. Let
\[
 C\in|L^a\otimes\alpha|,\qquad
 D\in|L\otimes\alpha^{-1}|
\]
be smooth sharp curves. Fix either component and a connected
elementary abelian two-torsor on it, with group $G$, and vary
the other component in its complete system. On the transverse
intersection open the lifted cover of the $N$ points has image
\[
                     \ker(\Sigma:G^N\to G)\rtimes S_N.
\]
The conclusion persists after restricting the moving component
to its smooth sharp open.
\end{lemma}
\begin{proof}
We verify the only new geometric issue in extending
Lemma~\ref{lem:mixed-lifted-intersection-wreath-candidate}:
the generic hyperplane section on the fixed component has
one contact of order two and no other multiple point.

First fix $C$ and move $D$ in the cubic product system.
Both projections of $C$ have degree $3a$, which is odd;
they are therefore separable, with finite ramification sets.
The two degree-three factor line bundles separate second
jets except at their finitely many flex points. Their tensor
product consequently separates surface second jets at a
general point of $C$, and hence curve second jets there.

The proof of Lemma~\ref{lem:product-small-degree-jets} applies
to any translates of the two cubic embeddings. Its
polar-finiteness assertion is uniform in the ordinary modulus:
for the ordinary model the partials of
$Y^2Z+XYZ+X^3+cZ^3$ are independent by their $X^2,Y^2,XZ$
coefficients; for the supersingular model they are
$X^2,Z^2,Y^2$. Thus only finitely many tangents pass through
a fixed point of either cubic. For a fixed surface point,
the exceptional second points for joint surface first jets
lie in the two factor fibres and finitely many crossed
tangent incidences. Their intersection with the ample
irreducible curve $C$ is finite. On $C\times C$ the four
curve first jets therefore have rank four off a subset of
dimension at most one.

Rank two means that the two embedded curve tangent lines
coincide. At a point where both projection derivatives are
nonzero, its tangent line meets the Segre surface only in
that point: every other point on the line has tensor rank
two. The remaining first points are finite, and their
tangent lines meet $C$ in finitely many points. Hence the
rank-two pairs are finite. The three pair strata have
dimensions at most $2,1,0$ and impose at least $4,3,2$
conditions. Two tangencies have codimension at least two
in the hyperplane space. Higher contact has codimension
at least two by generic second-jet separation; the finitely
many exceptional points already impose two first-jet
conditions. The tangent incidence is an irreducible
projective bundle over $C$, generically finite over its
image, and that image is a divisor with precisely one
generic contact of order two.

Conversely fix $D$ and move $C$. If $a=1$, the preceding
argument applies with the two curves interchanged. If
$a\geq3$, the factor bundles of degree $3a\geq9$ separate
length-four jets at every two points and second jets at
every point. Their product therefore separates all joint
surface first jets at two distinct points and surface
second jets at one point. The same tangent-incidence
argument is then immediate, with no exceptional pair
strata. These assertions hold for every Picard twist.

In either direction use the image of the surface section
space on the fixed curve. Its linear system is very ample.
Projection from the possible restriction kernel has
geometrically integral affine fibres, so passing from the
projective image system back to the original parameter
space does not change finite geometric monodromy.
The unlabelled symmetric image is also a direct instance
of the quasireflexive hyperplane theorem
\cite[Theorem~2 and Proposition~3.1]{Entin2021}; the jet
argument above verifies its hypothesis.

For clarity the labelled part uses the whole earlier
proof, not just its unlabelled conclusion. The incidence
of a point on the connected torsor and a hyperplane through
its image is an irreducible projective bundle, finite over
the hyperplane space. It makes the lifted point image
transitive. At the generic simple contact the length-two
root cluster is $z^2+bz+c$ with independent constant and
linear parameters. The equation $y^2+y=c/b^2$ gives
nontrivial inertia even at an imperfect separably closed
generic residue field. Etaleness of the torsor makes this
a labelled transposition, fixing the other point fibres.
The only branch divisor is the tangent divisor. Quotient
the Galois image by the normal subgroup generated by these
inertias. Purity \cite{Stacks0BMB} extends the quotient
etale over projective space, where it is trivial by the
previously proved geometric simple connectedness.
Lemma~\ref{lem:labelled-transposition-wreath-candidate}
therefore gives the displayed sum-zero wreath image.
Finally, the moving smooth sharp locus is nonempty and
dense by the all-odd theorem and the rank theorem.
Restriction to it preserves every finite-cover component.
\end{proof}

\begin{lemma}[The unequal-component boundary group]
\label{lem:mixed-even-unequal-boundary-candidate}
\status{Proved under stated hypotheses}
\textup{[Proved; complete independent proof audits passed.]}
Vary smooth sharp transverse pairs
$C\in|L^a\otimes\alpha|$ and
$D\in|L\otimes\alpha^{-1}|$. Their nodal unions lie in
the maximal-rank open of $|L^d|$. Their relative system
has an exact sequence
\[
 \begin{gathered}
 0\longrightarrow T\longrightarrow V\longrightarrow R\longrightarrow0,
 \qquad T=\mathbf F_2^N/\mathbf F_2\mathbf1,\\
 R=(H^1_{\et}(C,\mathbf F_2)\oplus H^1_{\et}(D,\mathbf F_2))
                      /H^1_{\et}(X,\mathbf F_2).
 \end{gathered}
\]
Put $T^0=\ker(\sum:\mathbf F_2^N\to\mathbf F_2)/\mathbf F_2\mathbf1$,
of dimension $N-2$. The boundary image contains
$\operatorname{Hom}(R,T^0)$ and an actual block-diagonal
$\GL_u\times\GL_8$, acting naturally on two direct summands
and fixing a complement. Every nonzero invariant subspace
of the boundary representation contains $T^0$.
\end{lemma}
\begin{proof}
The normalized two-ranks are $9a^2$ and $9$. The $N$ nodes
add $N-1$ unit graph directions. Removing the ambient unit
line gives $9a^2+9+N-2=9d^2-2=r$, the maximal relative rank.
The normalization sequence gives the displayed extension,
with $\dim R=u+9$.

Let $A$ denote the one-dimensional ambient unit space.
The component relative images are $\GL_u$ and $\GL_8$ by
the all-odd theorem and degree one. They are independent
because the components vary independently on a dense open
of a product. Every Picard twist is a translate of the
untwisted system over the algebraic closure: its polarization
homomorphism is an isogeny and is surjective on geometric
points, including when inseparable.

We include the group-extraction details. The full cohomology
image of either component fixes $A$, has its indicated
general linear quotient and has additive kernel either
zero or the entire corresponding natural-dual module.
The previously proved vanishing of its natural-dual $H^1$
splits the complement. Such an image is perfect, and its
only nonabelian simple factor is that general linear group.
On $R$, the image is similarly an extension of
$L_0=\GL_u\times\GL_8$ by a submodule of the two
nonisomorphic natural-dual modules. Their $H^1$ vanishes,
so it is a split semidirect product and is perfect.

The node permutation image is $S_N$ by the preceding lemma.
It has no common nontrivial quotient with the perfect image
on $R$. The latter has only $\GL_u,\GL_8$ as nonabelian
simple factors, whereas $S_N$ has $A_N$. For $\GL_8$ the
prime $11$ distinguishes them, since $N\geq18$ and $11$
does not divide $|\GL_8(\mathbf F_2)|$. For $a\geq3$,
the order bound in the odd norm proof applies with
$u\geq9a-1$ and $N=18a$, and gives $|\GL_u|>N!$.
For $a=1$ both factors are $\GL_8$, already treated.
Goursat therefore supplies the joint product on $T\oplus R$.

Choose the complement $L_0$ on $R$ and kill the node
permutations. The remaining matrix kernel is a submodule
of $\operatorname{Hom}(R,T)$, which is a direct sum of
the two natural-dual types and trivial modules. Every
quotient has zero $H^1$, by natural-dual vanishing and
perfectness for the trivial part. The graph-of-a-cocycle
argument gives an actual block-diagonal $L_0$, fixing
$T$ and a complementary ambient line in $R$.

Fix $C$ and its maximal elementary two-torsor, whose group
$G=H^1_{\et}(C,\mathbf F_2)^*$ has dimension $9a^2$.
The lifted node image is $G^{N-1}\rtimes S_N$.
Its only nonabelian simple factor is $A_N$, whereas the
moving $D$ cohomology image is perfect with factor $\GL_8$.
Their joint image is a product. A pure lifted translation
$(g_i)$ with $\sum g_i=0$ can therefore be chosen with
trivial moving-$D$ cohomology action. It is an actual
element of the kernel $J\subset\operatorname{Hom}(R,T)$,
and on $H^1_{\et}(C)$ its gluing action is
\[
                   \xi\longmapsto(\xi(g_i))_i\pmod{\mathbf F_2\mathbf1}.
\]
Indeed extend the fixed torsor on $C$ trivially over $D$
with chosen node gluings. Continuation changes these gluings
by exactly the displayed evaluations, while constant
trivializations remove the diagonal. This gives every
map from $H^1_{\et}(C)$ to $T^0$.

Under $L_0$, the module $\operatorname{Hom}(R,T)$ is
semisimple with three isotypic parts, from the two natural
blocks and the ambient line. Projecting the $L_0$-submodule
$J$ to these parts puts all maps from the $C$ block and
from $A$ to $T^0$ into $J$, with other columns zero.
Fixing $D$ and moving $C$ gives the remaining block's
columns: its lifted wreath has factor $A_N$ and the moving
cohomology image has factor $\GL_u$, still with no common
quotient. Thus $\operatorname{Hom}(R,T^0)\subset J$.

If an invariant subspace has nonzero image in $R$, these
translations put all of $T^0$ in it. If it is nonzero and
lies in $T$, choose a nonconstant representative; subtract
a transposition of two unequal coordinates. The resulting
class of $e_i+e_j$ and its $S_N$ orbit span $T^0$.
This proves the last assertion as well.
\end{proof}

\begin{lemma}[The common even grid determines the whole kernel]
\label{lem:mixed-even-uniform-common-grid-candidate}
\status{Proved under stated hypotheses}
\textup{[Proved; complete independent proof audits passed.]}
Every nonzero global sub-local system $S\subset K_d$
equals $K_d$. In particular $K_d$ is geometrically irreducible.
\end{lemma}
\begin{proof}
Set $n=3d$. Choose $n$ distinct points on each elliptic
factor, with each total divisor linearly equivalent to
$nO$. Their vertical and horizontal fibres give a grid
$Z\in|L^d|$ with graph $K_{n,n}$. Its graph rank is
$(n-1)^2$, and its $n$ ordinary components contribute
$n$ units. After removing the ambient line, its relative
unit rank is
\[
                (n-1)^2+n-1=9d^2-3d,
\]
below $r=9d^2-2$. We use the finite Lang part, never a
full-fibre specialization isomorphism.

For any three row and column indices, put the corresponding
six components into a small curve $D$ and the complementary
components into a large curve $C$. Their line bundles are
$L\otimes\alpha^{-1}$ and $L^a\otimes\alpha$ for an
appropriate Picard class. The twists are opposite because
the total grid lies in $|L^d|$; no individual triple is
required to sum to zero. Their $N=18a$ cross-edges are
simple intersections and both curves are smooth there.
They lie in a complete boundary family whose dense open
consists of smooth sharp transverse pairs.

We recall the exact finite-part comparison so the deeper
rank drop causes no unstated surjectivity claim. At the
strict henselization of the ambient coefficient space at
$Z$, the coherent Frobenius bundle has a separated
quasi-finite etale Lang scheme. Its finite part is constant,
with closed fibre the relative unit space $W_Z$ of the
rank just computed, and injects into the generic relative
system. The maximal-rank open is nonempty and integral.
Extend $S$ there from the dense smooth open. Membership
of each constant finite-Lang section in $S$ is open and
closed, hence constant. This defines one subspace
$S_Z\subset W_Z$.

On the strict local parameter scheme of any boundary
family through $Z$, its cross-points remain etale, so
its graph heart extends as a constant unit space. Its
closed injection agrees with the ambient finite Lang
part: separated etale sections over a henselian local
scheme with the same closed value coincide. The boundary
lemma places this heart in $S$ at its generic full-rank
point. Thus every special cross-heart lies in the same
$S_Z$. This is the finite-part construction proved in
Lemma~\ref{lem:mixed-degree-two-common-grid-candidate},
with no change to its normality or membership hypotheses.

Here is the uniform graph comparison. First choose small
row and column sets $I_0=J_0=\{0,1,2\}$, and then
$I_1=J_1=\{0,1,3\}$. Each cross set has $N$ edges.
A vertex coboundary supported on the cross set must be
constant on each of the two internal connected blocks.
Modulo the global constant it is a multiple of the sum
of all cross-edges. Since $N$ is even, that coboundary
lies in the sum-zero cochain hyperplane. The cross-heart
therefore has dimension $(N-1)-1=N-2$ in graph cohomology.

Writing $e_{ij}$ for an edge cochain, the second heart
contains
\[
                       v_D=e_{02}+e_{12},\qquad
                       v_C=e_{34}+e_{35}.
\]
The first restricts only to the small internal block of
the first partition, the second only to its large block.
Their evaluations on the respective four-cycles
\[
 e_{01}+e_{02}+e_{21}+e_{22},\qquad
 e_{33}+e_{34}+e_{43}+e_{44}
\]
are both one. Thus both internal restrictions are nonzero.
All these indices exist for $n\geq6$.

The associated constant finite-Lang sections belong to
$S$ on the first boundary family. Projection to the two
component relative coherent cohomology bundles is regular
and Frobenius-equivariant even at $Z$. Indeed Kunneth gives
$H^1(X,(L^b\otimes\beta)^{-1})=0$ for each $b>0$ and
Picard twist $\beta$. Connectedness, constant arithmetic
genus and cohomology-and-base-change give the bundles and
the injection of ambient $H^1(\mathcal O_X)$ throughout.
At an internal grid, ambient cohomology injects into the
sum of its elliptic-component cohomologies, whereas graph
classes map to zero there. The two graph classes therefore
survive the relative quotient. Nonzero coherent projections
at $Z$ remain nonzero generically on the integral boundary
family.

Consequently the image of $S$ has nonzero projection to
both independent natural component representations. It
contains their direct sum of dimension $u+8$. Its kernel
already contains the first heart, of dimension $N-2$.
Hence
\[
                 \dim S\geq u+8+N-2=9d^2-4=\dim K_d.
\]
This proves $S=K_d$.
\end{proof}

\begin{theorem}[Full images in all even balanced mixed degrees]
\label{thm:mixed-every-fixed-modulus-all-even-candidate}
\status{Proved under stated hypotheses}
\textup{[Proved; complete independent proof audits passed.]}
Let $k$ be any perfect field of characteristic two, $E/k$
ordinary and $E_0/k$ supersingular. For every even $d\geq2$
the full geometric and arithmetic relative images for
$|\mathcal O_E(3dO)\boxtimes\mathcal O_{E_0}(3dO)|$ are
\[
 \operatorname{Hom}(\mathbf F_2^2,\mathbf F_2^{9d^2-4})
                       \rtimes\GL_{9d^2-4}(\mathbf F_2).
\]
This is the full pointwise stabilizer of the two invariant
covectors. Their common kernel is irreducible and is contained
in every nonzero invariant subspace of the full relative system.
\end{theorem}
\begin{proof}
Work geometrically. Every nonzero invariant subspace contains
a boundary heart. That heart lies in $K_d$: invariant
covectors vanish on its nontrivial simple symmetric-group
representation. The preceding lemma makes $K_d$ irreducible,
so every such subspace contains $K_d$. In particular the
full relative system has no invariant vector.

The two actual natural boundary blocks also lie in $K_d$.
Their dimensions and the heart dimension sum to
$u+8+N-2=\kappa$, so they give all of $K_d$.
Let $J$ be the normal subgroup of the image on $K_d$
generated by transvections. Irreducibility makes its
restriction semisimple by the invariant-socle argument.
The actual eight-dimensional general linear block is
contained in $J$, with its natural constituent occurring
once; it lies in a single simple $J$-summand. Rank-one
maps from it to every heart vector put the whole heart
in that summand. The large natural block likewise lies
in one simple $J$-summand, since its natural constituent
for its own general linear factor occurs with multiplicity one.
The analogous maps from that block to the heart then put it in the
same summand: otherwise they would carry a vector of
a different simple summand into their zero intersection.
The summand is all of $K_d$, so $J$ is irreducible.

The previously audited transvection classification
\cite[Theorem~3.6, printed p.~22]{KantorSeress2001}
gives $\GL(K_d)$ or a group preserving a nondegenerate
alternating form. The latter is impossible: an actual
natural $\GL_8$ block fixing a complement has no invariant
bilinear form or functional and would be in the radical.
Thus the image on $K_d$ is the whole $\GL(K_d)$.

The quotient by $K_d$ is the constant space $\mathbf F_2^2$.
The remaining translation kernel is $K_d\otimes U$ for
a subspace $U\subset(\mathbf F_2^2)^*$. Modulo it, the
group is the graph of a cocycle. The proved natural-module
$H^1$ vanishing makes this a linear complement. If $U$
were proper, a nonzero quotient direction annihilated
by $U$ would have a fixed lift, contradicting absence
of invariant vectors. The full translation module occurs.

The exact-covector theorem proves that both geometric
covectors are defined over every perfect base field:
the unique toral invariant derivation and the normalized
nilpotent derivation give the two exact operators.
Thus arithmetic monodromy is contained in the same
pointwise stabilizer and contains its full geometric
subgroup. Equality follows.
\end{proof}

\begin{corollary}[Every balanced degree for every fixed mixed product]
\label{cor:mixed-every-fixed-modulus-all-degrees-candidate}
\status{Proved under stated hypotheses}
\textup{[Proved; complete independent proof audits passed.]}
Under the preceding hypotheses, for every $d\geq1$ the
geometric and arithmetic images coincide and are
\[
 \begin{cases}
 \GL_{9d^2-1}(\mathbf F_2),&d\text{ odd},\\
 \operatorname{Hom}(\mathbf F_2^2,\mathbf F_2^{9d^2-4})
       \rtimes\GL_{9d^2-4}(\mathbf F_2),&d\text{ even}.
 \end{cases}
\]
The odd representation is irreducible. In even degree the
common kernel of its two invariant covectors is irreducible
and is contained in every nonzero invariant subspace.
\end{corollary}
\begin{proof}
Combine the all-odd theorem with the preceding all-even theorem.
\end{proof}

\begin{remark}[Scope and exact graph support]
\label{rem:mixed-all-even-wreath-scope-candidate}
\status{Proved under stated hypotheses}
\textup{[Proved; complete independent proof audits passed.]}
The geometric extension concerns unequal odd components and
the finite Lang comparison at a deeper grid degeneration.
The special rank is $9d^2-3d$, never silently identified
with the generic rank $9d^2-2$. The opposite Picard twists
are part of every grid comparison. The separate root
program \path{code/mixed_even_degree_common_grid.py}
checks actual cross-hearts and the two projection witnesses
for $d=2,4,\ldots,20$; its author run took
$0.015011040$ seconds. Independent monodromy and
reduction/descent reruns took $0.013861630$ and
$0.013912451$ seconds, preserving the author output.
That finite calculation supports
only the explicitly proved uniform graph argument, not
the monodromy or specialization statements. The complete
unequal-component geometry, finite-part comparison and group
assembly passed independent root and reduction/descent audits.
This settles every balanced degree on each fixed mixed elliptic
product, while making no claim for arbitrary nonproduct surfaces.
\end{remark}


\subsection{A full elementary inertia group for every nonordinary principal surface}
\label{sec:nonordinary-uniform-Q-inertia-candidate}
\status{Proved and independently checked.}

Throughout this subsection the ground field is algebraically closed of
characteristic two. Use the nine theta coordinates and the four direction
coefficients $\alpha_i$ of Lemma~\ref{lem:DE-A3-theta-pfaffian}. Write
$M(s)$ for the alternating contraction matrix, $q=q_\alpha$ for its
Pfaffian cubic, and $Q(s)=\sum_i s_i^3$. The actual nonordinary principal
surface tensor has $q\ne0$, by
Proposition~\ref{prop:DE-A3-ppav-degree-one}. The relative Frobenius
bundle is the Euler bundle $\Omega^1_{\mathbf P^8}$ and its determinant
is $q^2Q$, by \eqref{eq:DE-A3-relative-theta-Hasse}.

\begin{lemma}[Stable loss one on the Hermitian cubic]
\label{lem:nonordinary-uniform-Q-stable-loss-candidate}
\status{Proved under stated hypotheses}
For every four-parameter tensor with $q\ne0$, the relative Frobenius
has linear rank and stable rank seven at the generic point of $Q=0$.
Its determinant has order one there, and the linearized kernel map
from that divisor to the constant projective eight-space is $[s]\mapsto[s]$.
\end{lemma}
\begin{proof}
The cubic $Q$ is smooth, since its partial derivatives are $s_i^2$;
it is therefore geometrically integral. Also $q$ is not proportional
to $Q$. Indeed, if all its line coefficients vanished while its pure
cube coefficient were nonzero, all four $\alpha_i$ would be nonzero
and would satisfy $\sum_{j\ne i}\alpha_j^3=0$. Subtracting any two
of these equations makes their cubes equal, but then each sum of the
other three equals that common nonzero cube, a contradiction.
Consequently $Q\setminus\{q=0\}$ is dense and the determinant has
order one along $Q$.

At such a point $\ker M=ks$ and $\operatorname{im}M=s^\perp$.
Since $Q=0$, there is a vector $u$ with $Mu=s^{[2]}$. On the dual
relative Cartier quotient its linearized kernel is $k\bar u$ in
$W/ks$. The rank-seven first image is already stable precisely when
\begin{equation}\label{eq:nonordinary-Q-stable-test-candidate}
                          s^t u^{[2]}\ne0.
\end{equation}
To check this criterion directly, the kernel of the next squared
linearization is represented by $u^{[2]}$ modulo $ks^{[2]}$.
The latter line already belongs to $\operatorname{im}M$.
Thus this next kernel meets the first image if and only if
$u^{[2]}\in\operatorname{im}M=s^\perp$. Equality of the first two
linearized image ranks implies stabilization of all subsequent ranks.
Changing $u$ by a multiple of $s$ does not change the test, because
$Q=0$.

First suppose that for some direction, renamed vertical, its line
coefficient
\[
          c_0=\alpha_0(\alpha_1^3+\alpha_2^3+\alpha_3^3)
\]
is nonzero and two of the other cubes differ. The stabilizer of a
direction in $\GL_2(\mathbf F_3)$ permutes its other three directions
as the full symmetric group. We may therefore arrange
$\alpha_1^3+\alpha_3^3\ne0$. Order coordinates lexicographically
by $(x,y)\in\mathbf F_3^2$, with the vertical direction first, followed
by slopes zero, one, two. Take the pencil
\[
                 s=(1,1,z,z,0,0,0,0,0).
\]
Here $Q=0$ identically and $q=c_0z$. Delete row and column zero of
$M$ to obtain the invertible alternating eight-by-eight matrix $B$.
Its Pfaffian is $q$. The inverse alternating pairing between
$e_0+e_1$ and $e_2+e_3$, with coordinate zero deleted, has numerator
\[
 P_6(0,1,2)+P_6(0,1,3)
                    =(\alpha_1^3+\alpha_3^3)z,
\]
where $P_6(i,j,k)$ denotes the Pfaffian after deleting those three
indices. For completeness the two Pfaffians are
\[
 \begin{split}
 P_6(0,1,2)&=(\alpha_1^3+\alpha_2^3+\alpha_3^3)z
                         +\alpha_1\alpha_2\alpha_3z^3,\\
 P_6(0,1,3)&=\alpha_2^3z+\alpha_1\alpha_2\alpha_3z^3.
 \end{split}
\]
These identities follow by the fifteen matching terms for each
six-by-six Pfaffian. The inverse-matrix Pfaffian formula, or its
entrywise product with $B$, now gives a solution $u_0=0$ of
$Mu=s^{[2]}$ satisfying
\[
 s^t u^{[2]}=(z^4+z)
       \left(\frac{(\alpha_1^3+\alpha_3^3)z}{c_0z}\right)^2.
\]
One can also derive this last formula without expanding the solution:
write $s=v+zw$ with binary vectors $v=e_0+e_1$, $w=e_2+e_3$.
The inverse of $B$ is alternating; its self-pairings on $v,w$ vanish,
while its two cross-pairings agree. Expanding
$s^t(B^{-1}s^{[2]})^{[2]}$ gives exactly the displayed expression.
It is nonzero in $k(z)$, proving the required generic stable rank.

Here is the complete coefficient reduction. Put $x_i=\alpha_i^3$.
If exactly two $x_i$ are nonzero, either nonzero direction has
$c_i\ne0$ and two unequal remaining cubes. If exactly three are
nonzero and not equal, at least one nonzero $x_i$ differs from
$S=\sum_jx_j$; that direction has $c_i\ne0$, and its remaining
cubes include zero and a nonzero value. Three equal nonzero cubes,
or just one nonzero cube, are precisely excluded cases $q=0$.
If all four are nonzero, some $x_i\ne S$, since otherwise $S=4S=0$.
Choose such an $i$. If its other three cubes are not equal we are
done. If they equal $v$ and $x_i\ne v$, choose one of these three
instead: its complementary sum is $x_i\ne0$ and its complementary
cubes are unequal. The only remaining case is therefore
\[
                         x_0=x_1=x_2=x_3\ne0.
\]

In this last case a scalar multiple of the tensor and a unitary
diagonal theta change reduce all coefficients to one. To see this
explicitly, divide by one coefficient and write the others as powers
of a primitive cube root $\zeta$. A diagonal change
$e_a\mapsto\zeta^{R(a)}e_a$, with $R$ a binary quadratic form over
$\mathbf F_3$, multiplies the direction coefficient by
$\zeta^{-R(v_i)}$. The evaluations at the four directions form exactly
the hyperplane of four-tuples with sum zero: the formula for
$R=Ax^2+Bxy+Cy^2$ gives the four values
$C,A,A+B+C,A-B+C$. Adding a common exponent supplies its complementary
one-dimensional space. Thus such a diagonal change, together with a
common tensor scalar, makes every coefficient one. The diagonal change
satisfies $g^{[2]}=g^{-t}$, so preserves $Q$ and the typed relative
Frobenius system. Multiplying its Frobenius matrix by a nonzero
constant preserves every iterated rank.

For this single normal form an exact point over $\mathbf F_8$ suffices.
Let $\theta^3+\theta+1=0$, and encode
$a_0+a_1\theta+a_2\theta^2$ by $a_0+2a_1+4a_2$.
Take
\[
 \begin{split}
 s&=(0,5,3,3,6,0,4,7,5),\\
 u&=(3,3,3,1,5,7,6,0,0).
 \end{split}
\]
The contraction rule gives, in that field,
\[
        Q(s)=0,\qquad q(s)=7,\qquad Mu=s^{[2]},
                         \qquad s^tu^{[2]}=7.
\]
These are exact finite-field equalities, not a sampling inference.
They establish \eqref{eq:nonordinary-Q-stable-test-candidate} in the
remaining normal form. Openness of the second-iterate rank-seven
condition on the integral $Q$ proves the general assertion.
Finally on $Q\setminus\{q=0\}$ the kernel of the coherent
linearization, whose domain is $\operatorname{ann}(s^{[2]})$,
is exactly $ks$. Its projective map is the stated identity.
\end{proof}

\begin{corollary}[Full seven-dimensional elementary geometric inertia]
\label{cor:nonordinary-ppav-uniform-root-inertia-candidate}
\status{Proved under stated hypotheses}
For every nonordinary principally polarized abelian surface in
characteristic two, the degree-one relative geometric image on the
maximal-rank smooth-member open of $|3\Theta|$ contains the entire
group of order $2^7$ fixing a seven-dimensional hyperplane pointwise
and acting trivially on its one-dimensional quotient.
\end{corollary}
\begin{proof}
Apply Lemma~\ref{lem:moving-kernel-full-root-inertia-candidate}
to $Q\subset\mathbf P^8$ and the relative Euler bundle, which embeds
in a constant nine-dimensional bundle on an affine coefficient chart.
The preceding lemma verifies stable rank seven, a one-dimensional
linearized kernel, and determinant order one. Its kernel map is the
identity on the smooth seven-dimensional cubic $Q$, so its differential
has rank seven. The lemma gives all $2^7$ elementary inertia elements.
Restricting the coefficient invertibility open to its dense smooth-curve
open preserves the finite geometric image. All assertions are geometric
and therefore remain valid after algebraically closed extension.
\end{proof}

\begin{remark}[Certificate and remaining global gate]
\label{rem:nonordinary-ppav-uniform-root-scope-candidate}
\status{Proved under stated hypotheses}
The exact symbolic certificate is
\path{agents/monodromy/code/ppav_theta_Q_pencil_pfaffians_complete.py};
it verifies the two six-Pfaffians, every coordinate of $Mu=q s^{[2]}$
for the numerator vector, and its pairing, in
$\mathbf F_2[\alpha_0,\ldots,\alpha_3,z]$.
The exact normal-form witness is independently reconstructed by
\path{agents/monodromy/code/ppav_theta_equal_cubes_F8_witness.py}.
Their initial author executions took respectively
$0.1280531329975929$ and $0.0005276640004012734$ seconds.
The earlier bounded finite-field exploration and both failed binary
isotropic-pencil searches are retained separately; neither is a proof
input. This corollary does not assert full monodromy: connectivity of
the nonzero Lang cover, or another global irreducibility argument,
remains to be proved for an arbitrary fixed nonproduct principal surface.
\end{remark}


% Root result promoted after two independent complete proof/source audits; original candidate preserved.
\subsection{The actual Cartier cubic and its rank strata}

\begin{lemma}[Identifying the Cartier cubic on a nonproduct surface]
\label{lem:root-nonordinary-Coble-identification-candidate}
\status{Proved under the stated nonproduct, nonordinary, characteristic-two hypotheses.}
Let $(A,\Theta)$ be a nonproduct principally polarized abelian surface
over an algebraically closed field of characteristic two, and suppose
that $A$ is nonordinary. Set $L=\mathcal O_A(3\Theta)$ and
$W=H^0(A,L)$. Use the typed theta basis and the alternating tensor
$T\in\bigwedge^3W^*$ of
Theorem~\ref{thm:ppav-level3-Cartier-symmetric}.
Let $M(v)$ be its contraction matrix and $q=q_\alpha$ the nonzero
Pfaffian cubic of Lemma~\ref{lem:DE-A3-theta-pfaffian}.
There is a $3\Theta^{(2)}$ embedding
\[
                 j:A^{(2)}\hookrightarrow\mathbf P(W)
\]
such that $q$ is, up to scalar, the Heisenberg-invariant Coble cubic
of this embedded surface. Its reduced singular locus is exactly
$j(A^{(2)})$. The rank-at-most-four locus of $M$ has this same
underlying set, and $M$ has rank exactly four at every point of it.
\end{lemma}
\begin{proof}
Write $F:A\to A^{(2)}$ for relative Frobenius and put
$E=F_*L$, a bundle of rank four. The established finite-duality
pairing is perfect and has global map
\[
 \beta:W\otimes W\longrightarrow W^{(2)}.
\]
The bundle $E$ is globally generated, as proved in
Proposition~\ref{prop:DE-A3-ppav-degree-one}. The map
$R:W^{(2)}\to W^*$ sends $e_a^{(2)}$ to $e_a^*$.
An evaluation functional of $L^{(2)}$ at $P\in A^{(2)}$, up to
its fibre trivialization, is an element of $(W^{(2)})^*$.
Define $j(P)$ by applying $(R^*)^{-1}$ to this functional.
Thus $j$ is precisely the complete $|L^{(2)}|$ embedding followed
by a fixed projective linear isomorphism; it does not identify
$A$ with $A^{(2)}$.

For any representative $v_P$ of $j(P)$, contraction $M(v_P)$ is
the perfect fibre pairing on $E_P$ pulled back by the surjection
$W\to E_P$. Its rank is therefore exactly four. In particular all
six-by-six Pfaffians vanish on the reduced embedded surface, so
the inclusion in the rank-at-most-four scheme holds scheme-theoretically.
On a chart $v_i\ne0$, the principal eight-by-eight Pfaffian is
$v_iq(v)$. Its differential is a linear combination of six-by-six
Pfaffians. At rank at most four both that Pfaffian and its
differential vanish; division by $v_i$ gives $q=dq=0$.
Hence $q$ is singular along $j(A^{(2)})$.

The finite Heisenberg action on the target of $\beta$ is the
coefficient-square representation, and $R$ identifies it with the
contragredient representation. Consequently the just-defined
embedding is equivariant for the same level-three Heisenberg group
acting on $\mathbf P(W)$. The cubic $q$ is invariant under this
group by the Pfaffian calculation already proved.

For clarity, the uniqueness used here is invariant uniqueness,
which is valid in characteristic two. Put $B=A^{(2)}$ and
$L_B=\mathcal O_B(3\Theta^{(2)})$. By
\citet[Theorem~2.4, p.~723]{Sekiguchi1977PNII}, applied with
the ample bundle $\mathcal O_B(\Theta^{(2)})$ and exponents $3,3$,
the multiplication map
\[
 \operatorname{Sym}^2H^0(B,L_B)\longrightarrow H^0(B,L_B^2)
\]
is surjective. Its dimensions are $45$ and $36$, respectively,
so $h^0(\mathcal I_{j(B)}(2))=9$.
\citet[Proposition~2.1, Remark~(3), Corollary~2.2, pp.~191--192]
{Beauville2003Coble} then gives a one-dimensional space of
Heisenberg-invariant cubics singular along $j(B)$.
The characteristic restriction there is only characteristic
different from three for this assertion. Its separate
Proposition~3.1 about unrestricted uniqueness is not needed.

The nonproduct principal surface is the Jacobian of its smooth
genus-two theta divisor, by
Lemma~\ref{lem:ppav-surface-dichotomy}. Choose a Weierstrass point
over the algebraically closed field. Translating the Weierstrass-point
representative of the principal theta divisor identifies its complete
$3\Theta$ embedding with the fixed complete embedding above; the induced
linear isomorphism transports the theta group and the tensor together.
In this identified embedding, the construction of
\citet[Section~4, Lemmas~4.1--4.4 and Proposition~4.5(a)]
{RainsSam2018Trivectors} gives a stable alternating tensor whose
rank-four surface is $j(B)$; the identification follows also from
its Propositions~5.1--5.2. Lemma~3.3 of that paper says its Pfaffian
cubic has singular locus $j(B)$ as a set. This cubic is invariant
under the same finite theta group. The preceding uniqueness makes
it a scalar multiple of $q$.

Every rank-at-most-four point of our matrix is singular on $q$ by
the local Pfaffian argument above. Conversely every point of
$j(B)$ has rank exactly four. Since the reduced singular locus
of $q$ is $j(B)$, this proves all the set-theoretic and exact-rank
claims. No equality of Jacobian schemes is inferred from the
quoted singular-locus statement.
\end{proof}

\begin{lemma}[A plane cubic rules out instability]
\label{lem:root-nonordinary-Cartier-stability-candidate}
\status{Proved under the stated nonproduct, nonordinary, characteristic-two hypotheses.}
The actual alternating Cartier tensor $T$ in the preceding lemma
is stable for $\operatorname{SL}(W^*)/\mu_3$.
Its rank-at-most-four scheme is smooth of dimension two and equals
$j(A^{(2)})$ scheme-theoretically. There are no nonzero points of
rank at most two, and $q=0$ is smooth at every rank-six point.
\end{lemma}
\begin{proof}
Suppose that $T$ is not stable. Apply
\citet[Proposition~2.12 and the proof of Lemma~3.2]
{RainsSam2018Trivectors}, with its $V_9=W^*$.
There is a six-dimensional subspace $U\subset W^*$ with
\[
 T\in\bigwedge^3U+\bigwedge^2U\otimes(W^*/U).
\]
On the plane $\Pi=\mathbf P(U^\perp)\subset\mathbf P(W)$,
the contraction matrix is supported on $U$. Its rank-at-most-four
locus is therefore cut out on $\Pi$ by one six-by-six Pfaffian,
a cubic polynomial. If that polynomial is identically zero, the
whole plane lies in the rank-at-most-four set, hence in $j(A^{(2)})$,
which is impossible: an abelian variety contains no rational curve.
Otherwise its zero scheme is a plane cubic supported in this
abelian surface.

Each irreducible component of the reduced support has degree at
most three. A line, a conic, or a singular integral plane cubic
has rational normalization, and would give a nonconstant map
from $\mathbf P^1$ to an abelian variety. Thus the only possibility
is a smooth integral plane cubic $E$, with multiplicity one; its
degree already exhausts the total degree three. After translating
an origin, its inclusion in the abelian surface is a homomorphism
from an elliptic curve, by the rigidity property of morphisms
between abelian varieties. Since the embedding is $3\Theta^{(2)}$,
\[
                         \Theta^{(2)}\cdot E=1.
\]
This forces the principal polarization to decompose, as follows.
If $i:E\hookrightarrow B=A^{(2)}$ is the inclusion, then
$i^*\lambda_\Theta i:E\to E^*$ is the degree-one polarization
of $E$ and is an isomorphism. Therefore
\[
 r=(i^*\lambda_\Theta i)^{-1}i^*\lambda_\Theta:B\to E
\]
is a retraction. The homomorphism
$E\times\ker r\to B$, $(e,b)\mapsto i(e)+b$, has the explicit
inverse $b\mapsto(r(b),b-i(r(b)))$ and is an isomorphism.
The kernel is consequently a connected elliptic curve. The two
factors are orthogonal for $\lambda_\Theta$ by construction, and
the degree-one principal polarization restricts principally to
both. This contradicts the nonproduct hypothesis.
Thus $T$ is stable.

Lemma~3.1 of the same source now gives smoothness and dimension two
for the actual rank-at-most-four scheme and excludes rank at most
two. The preceding lemma identifies its underlying set with the
closed smooth surface $j(B)$. Their scheme-theoretic inclusion
is an equality, since both schemes are reduced with that support.
Lemma~3.3 gives smoothness of $q=0$ away from that surface, as
claimed. This deduction uses smoothness of the Pfaffian scheme
only after stability has been established.
\end{proof}

\begin{corollary}[Only the rank-six Lang incidence remains]
\label{cor:root-nonordinary-Lang-rank-four-excluded-candidate}
\status{Proved under the stated nonproduct, nonordinary, characteristic-two hypotheses.}
In the affine coefficient model for the degree-one relative Lang
cover, write its nonzero fixed-vector equation as
\[
                    M(s)x^{[2]}=x,\qquad x\ne0.
\]
The part where $M(x^{[2]})$ has rank at most four cannot dominate
the nine-dimensional section-parameter space.
The part where that matrix has rank eight is irreducible and dominates,
with
an affine-line fibre over a nonempty open of the Fermat cone
$\sum x_a^3=0$.
Thus any further dominating component must come from the
rank-six compatibility locus. Global connectivity and full
monodromy are not asserted by this corollary.
\end{corollary}
\begin{proof}
Full symmetry of $T$ rewrites the equation as
$M(x^{[2]})s=x$. At rank four the allowed nonzero $x$ lie in
the Frobenius preimage of the affine cone over $j(A^{(2)})$.
This set has dimension three, since Frobenius is finite and
the projective surface has dimension two. For any such $x$,
the solution set in $s$ is empty or an affine space of dimension
$9-4=5$. The total incidence therefore has dimension at most
eight. Lower ranks do not occur for nonzero $x$ by the preceding
lemma. This proves nondominance of this entire part.

At rank eight, the kernel is $kx^{[2]}$, so compatibility is
exactly $x^t x^{[2]}=0$. The Fermat cubic in $\mathbf P^8$ is
smooth: its gradient is $(x_a^2)$ and has no projective zero.
It is integral, since distinct positive-degree hypersurface
components in projective space would meet and make their product
singular. Its nonzero affine cone is therefore irreducible of
dimension eight. The rank-eight open is nonempty because its
complement is the cubic $q(x^{[2]})=0$, and the proved
nonproportionality of $q$ with the Fermat cubic persists under
the coefficient-square-root twist: explicitly,
$q_\alpha(x^{[2]})=q_{\alpha^{[1/2]}}(x)^2$.
The affine-linear solutions in
$s$ form an affine-line torsor on this open, hence an irreducible
incidence of dimension nine.
The Jacobian of $M(s)x^{[2]}-x$ in the $x$ variables is the identity.
Thus its projection to the $s$ space is etale. On the nonempty
rank-eight open its image is open and hence dense, proving dominance.
The only other possible nonzero rank is six. Removing the
rank-at-most-four part does not bound that remaining incidence,
so the last sentence is a limitation, not a connectivity claim.
\end{proof}


% Post-CP20 conditional argument. No assertion that its extra hypothesis
% holds on the actual nonordinary Cartier locus.
\subsection{A conditional polar-tensor test for the remaining Lang locus}
\status{Proved under the explicit extra stability hypothesis; independently checked.}

\begin{lemma}[A cubic containing a Frobenius-defined smooth sixfold]
\label{lem:polar-third-derivative-Lang-candidate}
\status{Proved under stated hypotheses}
Let $k$ be algebraically closed of characteristic two, let $V=k^9$, and
let $f$ be a homogeneous cubic on $V$. Suppose that on an open subset
of $V$ there are three rational vector-valued functions $k_i(x^{[2]})$
which are everywhere linearly independent. Set
\[
                 g_i(x)=x^t k_i(x^{[2]}),\qquad i=1,2,3.
\]
If an irreducible six-dimensional component of $g_1=g_2=g_3=0$
is contained in $f=0$, the alternating polar three-form of $f$
is not stable for $\operatorname{SL}(V)$.
\end{lemma}
\begin{proof}
Shrink to an open on which all functions are regular. Ordinary
coordinate differentiation gives
\[
 dg_i=k_i(x^{[2]})^t dx,\qquad D_vD_wg_i=0
\]
for all constant vector fields $v,w$, because all coefficients of
$k_i(x^{[2]})$ have zero differential, including rational denominators.
The three independent differentials make their zero scheme smooth
of dimension six. At a general point $z$ of the indicated component,
its local ideal is therefore generated by $g_1,g_2,g_3$; after shrinking
again one can write $f=\sum_i a_i g_i$ with regular $a_i$.
Put $U=T_z\{g_1=g_2=g_3=0\}$, so $\dim U=6$.
For any $u,v,w\in U$, the product rule gives
\[
                    (D_uD_vD_w f)(z)=0.
\]
Indeed every summand contains either $g_i(z)$, a first derivative
$D_u g_i(z)$ with $u\in U$, or a second or higher derivative of $g_i$.
All of these vanish. The left side is the constant alternating
polar three-form of the cubic evaluated on $u,v,w$. Hence this form
vanishes on $\bigwedge^3 U$.

Choose a complement to $U$ in $V$. On $V^*$ give its three-dimensional
subspace $U^\perp$ weight $2$ and a complementary six-space weight $-1$.
These weights sum to zero. Every monomial of the polar form contains
at least one factor from $U^\perp$, so it has nonnegative weight.
The resulting nontrivial one-parameter subgroup has a finite limit
on that form. The Hilbert--Mumford stable-orbit criterion therefore
excludes stability. Equivalently, if the limit has only weight-zero
terms it has a positive-dimensional stabilizer, and if there are
positive-weight terms the original orbit is not closed with finite
stabilizer. This includes the case of zero polar form.
\end{proof}

\begin{theorem}[Full degree-one image under polar stability]
\label{thm:nonordinary-ppav-polar-stable-conditional-candidate}
\status{Proved under stated hypotheses}
Let $(A,\Theta)$ be a nonproduct nonordinary principally polarized
abelian surface over a perfect field $k$ of characteristic two.
Use its typed theta tensor and cubic $q=q_\alpha$ from
Lemma~\ref{lem:root-nonordinary-Coble-identification-candidate}.
Over $\overline k$ put
\[
 f(x)=q_{\alpha^{[1/2]}}(x),\qquad
 b_i=\left(\alpha_i\sum_{j\ne i}\alpha_j^3\right)^{1/2}.
\]
Suppose, as an additional hypothesis, that the alternating polar
three-form of $f$ is stable. In the affine-line theta basis its
four Cartan coefficients are precisely $b_i$.
Then the degree-one relative unit local system has rank eight and
both its geometric and arithmetic images equal $\GL_8(\mathbf F_2)$.
\end{theorem}
\begin{proof}
The rank is already established. For completeness we verify the
remaining exceptional-incidence implication. Write the nonzero
Lang equation as
\[
                 M(s)x^{[2]}=x
       \quad\Longleftrightarrow\quad M(x^{[2]})s=x.
\]
At a rank-six point choose an invertible principal alternating
six-by-six block. Solving for the first six coordinates with each
of the other three coordinates free gives three rational vectors
$k_i(w)$, independent throughout this block chart. On the rank-six
locus they form the kernel of $M(w)$. Since the matrix is alternating,
compatibility of the displayed linear equation in $s$ is exactly
\[
                    x^t k_i(x^{[2]})=0,\qquad i=1,2,3.
\]
The three functions have independent differentials. Their zero
scheme consequently has dimension six. Also the rank-six condition
implies $q(x^{[2]})=f(x)^2=0$.
If the rank-six compatible locus had a six-dimensional component,
its closure in this block chart would be a component of those three
smooth equations contained in $f=0$. The preceding lemma would make
the polar form of $f$ nonstable, contrary to the hypothesis. Thus
its dimension is at most five. Its affine solution spaces in $s$
have dimension three, so this entire exceptional incidence has
dimension at most eight.

Corollary~\ref{cor:root-nonordinary-Lang-rank-four-excluded-candidate}
already bounds the rank-four part by eight, excludes smaller nonzero
ranks, and identifies the rank-eight part as one irreducible
nine-dimensional affine-line torsor over an open of
$\sum x_i^3=0$. The projection to the section-parameter space is
\'etale, since the derivative in $x$ of $M(s)x^{[2]}-x$ is $-I$.
On the maximal-rank open the nonzero Lang scheme is finite \etale;
every component dominates that open and has dimension nine.
The bounds just proved therefore force it to have exactly one
component. Geometric monodromy is transitive on the 255 nonzero
vectors, in particular irreducible.

By Corollary~\ref{cor:nonordinary-ppav-uniform-root-inertia-candidate},
this group contains all transvections
\[
             1+u\otimes\varphi,\qquad u\in\ker\varphi,
\]
for some nonzero covector $\varphi$. The orbit of $\varphi$ spans
the whole dual space, since its annihilator is an invariant subspace
and the representation is irreducible. Choose an orbit basis
$\varphi_1,\ldots,\varphi_8$ and its dual basis $e_1,\ldots,e_8$.
Conjugating the displayed subgroup gives
$1+e_j\otimes\varphi_i$ for every $j\ne i$.
These elementary matrices generate $\GL_8(\mathbf F_2)$.
The arithmetic group contains the geometric group and is a subgroup
of the same general linear group, so equality holds for it as well.
\end{proof}

\begin{remark}[The new hypothesis does not follow from tensor stability]
\label{rem:nonordinary-polar-stability-gap-candidate}
\status{Proved under stated hypotheses}
This theorem is conditional. The extra hypothesis is verified for every
nonordinary principal surface in
Theorem~\ref{thm:every-nonordinary-ppav-degree-one-full-candidate} below;
the present argument alone does not verify it on the actual Cartier locus. In particular, stability of the original theta tensor does
not suffice by a formal invariant-theory implication.
In the four-dimensional Cartan coefficient space, consider
\[
                    \alpha_1^3+\alpha_2^3+\alpha_3^3=0.
\]
This is an irreducible cubic hypersurface: its plane base is a smooth
Fermat cubic in characteristic two. It is not contained in the finite
union of reflection hyperplanes. By
\citet[Proposition~2.10]{RainsSam2018Trivectors}, its general point
therefore defines a stable original tensor, with all $\alpha_i\ne0$.
But its polar coefficient $b_0$ is zero, so the polar tensor lies on
a coordinate reflection hyperplane and is nonstable by
\citet[Proposition~2.9]{RainsSam2018Trivectors}.
This proves the failure of the proposed implication on the entire
stable Cartan family. It does not identify that hypersurface with
any locus of actual Cartier tensors. The subsequent normal-form
calculations show that the actual nonordinary Cartier locus satisfies
polar stability, and thereby prove the unconditional degree-one
fixed-surface theorem without the false general implication.
\end{remark}


% Post-CP21 candidate; exact universal polynomial certificate retained.
\subsection{Degree-one full image on every principal surface of two-rank zero}
\status{Proved and independently checked.}

\begin{lemma}[The polar tensor for every genus-two curve of two-rank zero]
\label{lem:p-rank-zero-genus-two-polar-stable-candidate}
\status{Proved under stated hypotheses}
Let $C$ be a smooth genus-two curve over an algebraically closed field
of characteristic two, and suppose its Jacobian has two-rank zero.
In the complete three-theta embedding of its Jacobian, the alternating
polar three-form of its Coble cubic is stable for $\operatorname{SL}_9$.
\end{lemma}
\begin{proof}
We first make the range of normal forms explicit. A genus-two curve
is hyperelliptic, and its degree-two canonical map is separable:
a purely inseparable degree-two extension of $k(x)$ would be
$k(x^{1/2})$, and hence rational. Choose a geometric Weierstrass
point as infinity. Its equation has the form
\[
       y^2+h(x)y=f_5(x),\qquad
       h=h_0+h_1x+h_2x^2,
\]
where $f_5$ has degree five and $h\ne0$. This is also the curve
normal form in \citet[Proposition~2.1]{RainsSam2018Trivectors}.
A basis of regular differentials is $\omega=dx/h$ and $x\omega$.
The rational Cartier rule gives
\[
 C\omega=\sqrt{h_1}\,\omega,
 \qquad C(x\omega)=(\sqrt{h_0}+\sqrt{h_2}x)\omega.
\]
Indeed one writes $dx/h=h\,dx/h^2$ and similarly for $x\,dx/h$;
Cartier retains exactly the odd powers of $x$. This triangular
semilinear matrix has stable rank zero precisely when
$h_1=h_2=0$. Thus $h=h_0\ne0$. Scaling $y$, then $x$, normalizes
the equation to $y^2+y=x^5+$ lower terms. Changes $y\mapsto y+u(x)$
remove the even powers and constant term over the algebraically
closed field. In particular every curve under consideration occurs
in the slightly larger family
\begin{equation}\label{eq:p-rank-zero-normal-form-candidate}
                  y^2+y=x^5+a x^3+b x^2+c x.
\end{equation}
Every member of this family is smooth, including its unique point
at infinity: its affine $y$-derivative is one, and its odd pole of
order five gives the smooth genus-two Artin--Schreier model at infinity.

Use the normal form of the alternating tensor in the same cited
Proposition~2.1, with $c_{15}=1,c_{12}=a,c_{18}=b,c_{24}=c$ and
all other coefficients zero. With $[ijk]=e_i\wedge e_j\wedge e_k$,
it is
\[
 \gamma=[267]+[258]+[348]+[169]+[357]+[249]+[178]+[456]
                 +[235]+a[147]+b[145]+c[134].
\]
\citet[Proposition~2.13 and Proposition~3.8, including their proofs]
{RainsSam2018Trivectors} show that this tensor is stable, that its
rank-four surface is the Jacobian of \eqref{eq:p-rank-zero-normal-form-candidate}
in its complete three-theta embedding, and that its Pfaffian cubic
is the corresponding Coble cubic. Let $B$ be the alternating polar
three-form of that cubic. The following explicit expression and
minor identity are over the polynomial ring $\mathbf F_2[a,b,c]$:
\begin{align*}
 B={}&[129]+[138]+[156]+[179]+[236]+[259]+[268]
                 +[349]+[367]+[457]\\
    &+a\bigl([134]+[148]+[157]\bigr)
                 +b[146]+c\bigl([126]+[135]\bigr).
\end{align*}
To specify the minor without any numerical specialization, form
\[
 \mathcal D_B:\operatorname{End}(k^9)\longrightarrow\bigwedge^3 k^9,
                   \qquad X\longmapsto X\cdot B.
\]
Its 81 columns are $E_{ij}$ in lexicographic order and its 84 rows
are the lexicographically ordered triples. Delete the three rows
$[467],[478],[578]$. The resulting square matrix satisfies
\begin{equation}\label{eq:p-rank-zero-polar-unit-minor-candidate}
                         \det\mathcal D'_B=1
                   \quad\text{in }\mathbf F_2[a,b,c].
\end{equation}
This is a finite polynomial identity, with a retained exact
unit-pivot elimination certificate described below. It proves that
the infinitesimal stabilizer of $B$ in $\operatorname{GL}_9$ is zero
for every value of $a,b,c$, and hence also in $\operatorname{SL}_9$.

It remains to justify why this infinitesimal test proves stability
in this setting. Since $\gamma$ is stable, it is conjugate into the
standard Cartan subspace, by
\citet[Corollary~5.6]{RainsSam2018Trivectors}.
The Pfaffian cubic construction and ordinary cubic polarization are
covariant under change of basis, up to a nonzero determinant scalar.
In a Cartan basis the cubic is the sum of pure cubes and affine-line
products, so its alternating polar tensor is again in the same
four-dimensional Cartan subspace. Therefore $B$ is conjugate, up
to scalar, to a Cartan tensor. By the same source's
Propositions~2.9--2.10, a nonstable Cartan tensor lies on a reflection
hyperplane and has a positive-dimensional stabilizer. Such a
stabilizer has nonzero Lie algebra. This contradicts
\eqref{eq:p-rank-zero-polar-unit-minor-candidate}. Thus $B$ is stable.
\end{proof}

\begin{remark}[A universal polynomial certificate]
\label{rem:p-rank-zero-polar-unit-certificate-candidate}
\status{Proved under stated hypotheses}
The executed exact program is
\path{agents/monodromy/code/ppav_supersingular_polar_uniform_lie.py},
with full output
\path{agents/monodromy/outputs/ppav_supersingular_polar_uniform_lie.json}.
It builds the alternating contraction matrix directly from the twelve
terms of $\gamma$, computes all nine principal eight-Pfaffians by
recursive matching, and verifies their common cubic after division
by the omitted coordinate. It then obtains $B$ from the squarefree
cubic monomials and constructs the infinitesimal action matrix from
$E_{ij}\cdot[e_r\wedge e_s\wedge e_t]$ by the exterior-product rule.
Polynomial Gaussian elimination uses only pivots equal to the constant
one; all 81 pivot pairs and all coefficients are retained.
The three unused exterior rows are exactly $[467],[478],[578]$;
the added trace row is unused as well. Thus the elimination proves
the displayed square-minor identity over the entire polynomial ring,
not just at a list of field values. The author execution at
18:29:35.876074 UTC on 19 November 2026 took
0.019472674997814465 seconds. The parent, reduction/descent, and obstruction/crystal agents independently
audited every proof and reread the cited original sources. The parent
independently reconstructed the displayed tensor action and the fixed minor
in 0.06436865900468547 seconds; reduction/descent independently rebuilt
all nine Pfaffians and that minor in 0.10397645600460237 seconds;
obstruction/crystals independently reconstructed the displayed minor
in 0.06390757599729113 seconds. Their author-output-preserving executions
of the complete retained program took respectively
0.02088291999825742, 0.02283291499770712, and
0.02654484500089893 seconds. All checks passed.
The earlier finite normal-form and Lie probes remain retained as
discovery evidence; they are not the proof of the uniform identity.
\end{remark}

\begin{theorem}[Every fixed principal surface of two-rank zero]
\label{thm:p-rank-zero-ppav-degree-one-full-candidate}
\status{Proved under stated hypotheses}
Let $(A,\Theta)$ be a principally polarized abelian surface of
geometric two-rank zero over any perfect field $k$ of characteristic
two. For the smooth maximal-rank locus in $|3\Theta|$, the relative
unit local system has rank eight, and both its geometric and
arithmetic images are $\GL_8(\mathbf F_2)$.
In particular the representation is geometrically irreducible.
\end{theorem}
\begin{proof}
Work first over an algebraic closure. If the principal surface is
a product, both elliptic factors are supersingular. In characteristic
two they are geometrically isomorphic to $y^2+y=x^3$, and the result
is Theorem~\ref{thm:ss-product-degree-one-full-group}.
Otherwise $A$ is the Jacobian of its smooth genus-two theta divisor.
The same is true of $A^{(2)}$, whose two-rank remains zero.
The preceding lemma proves stability of the alternating polar of
the Coble cubic of $A^{(2)}$.

Lemma~\ref{lem:root-nonordinary-Coble-identification-candidate}
identifies the actual Cartier Pfaffian cubic with this Coble cubic
in the exact complete three-theta embedding. Consequently their
polar tensors are identified under the same projective linear
change of basis and a nonzero scalar. No identification of the
actual Cartier trivector with a particular marked-curve trivector
is needed for this step: equality of the cubics is sufficient.
If different principal line-bundle representatives are used,
translation on the abelian surface identifies the complete systems
and the associated embeddings. Finally stability is unchanged by
the coefficient-square-root automorphism of the algebraically
closed field. The extra hypothesis in
Theorem~\ref{thm:nonordinary-ppav-polar-stable-conditional-candidate}
is therefore satisfied, and that theorem gives the full geometric
image. The arithmetic image contains it and lies in the same
$\GL_8(\mathbf F_2)$, giving equality over the original perfect field.
\end{proof}

\begin{remark}[The remaining degree and two-rank-one gates]
\label{rem:p-rank-zero-ppav-degree-one-scope-candidate}
\status{Proved under stated hypotheses}
This theorem concerns degree one for the complete three-theta
embedding. The two-rank-one extension is proved in
Theorem~\ref{thm:every-nonordinary-ppav-degree-one-full-candidate} below.
Higher-degree full images on a fixed nonproduct nonordinary surface
remain separate structural gates. The earlier conditional polar
theorem is nonvacuous on the actual Cartier locus by the uniform
normal-form results.
\end{remark}


% Post-CP21 candidate. Exact universal b^4 minor retained.
\subsection{Degree-one full image on every nonordinary principal surface}
\status{Proved and independently checked.}

\begin{lemma}[The two-rank-one polar tensor]
\label{lem:p-rank-one-genus-two-polar-stable-candidate}
\status{Proved under stated hypotheses}
Let $C$ be a smooth genus-two curve over an algebraically closed field
of characteristic two. If its Jacobian has two-rank one, the alternating
polar three-form of its Coble cubic is stable.
\end{lemma}
\begin{proof}
Begin with the degree-five Weierstrass equation and the triangular
Cartier matrix calculated in
Lemma~\ref{lem:p-rank-zero-genus-two-polar-stable-candidate}.
Stable rank one means exactly one of $h_1,h_2$ is nonzero. If $h_2=0$,
$h$ is already linear. If $h_1=0$ and $h_2\ne0$, translate the double
root of $h$ to zero, so that $h=h_2x^2$. Translate $y$ by a constant
to make $f_5(0)=0$. Smoothness at $(0,0)$ implies $f'_5(0)\ne0$,
because both the $y$-derivative and $h'(0)$ vanish there. Now set
\[
                         z=x^{-1},\qquad v=yx^{-3}.
\]
The transformed equation is
\[
                  v^2+h_2zv=z^6f_5(z^{-1}).
\]
Its right side has degree exactly five: the original constant
coefficient is zero and its linear coefficient is nonzero. Thus
its $h$ is linear, and the new infinity is the other Weierstrass point.
This explicit transformation needs no assertion about the number of
points in an inseparable branch scheme.

In the linear case, translate its root to zero, and translate $y$
by a constant to arrange $f_5(0)=0$. Nonzero rescalings of $x,y$
normalize $h=x$ and the leading coefficient of $f_5$ to one:
if the initial coefficients are $h_1x$ and $f_5^{\rm lead}x^5$,
one chooses the $x$ scale with cube $h_1^2/f_5^{\rm lead}$ and the
$y$ scale equal to $h_1$ times that $x$ scale.
The substitution $y\mapsto y+ux^2$ adds $u^2x^4+ux^3$, so removes
the fourth-degree term. The substitution $y\mapsto y+vx$ adds
$(v^2+v)x^2$, so removes the quadratic term. Both equations for
$u,v$ have solutions over the algebraically closed field.
Every curve in question consequently has the form
\begin{equation}\label{eq:p-rank-one-normal-form-candidate}
                   y^2+xy=x^5+a x^3+b x,\qquad b\ne0.
\end{equation}
Here $b\ne0$ is precisely smoothness at $(0,0)$: it is the only
possible affine singular point since the $y$-derivative is $x$.
The standard degree-five model is smooth at infinity. Conversely
the triangular Cartier rule gives stable rank one for every $b\ne0$.

In the explicit normal form of
\citet[Proposition~2.1]{RainsSam2018Trivectors}, take
$c_9=1,c_{12}=a,c_{24}=b$, with all other coefficients zero. The tensor is
\[
 \gamma=[267]+[258]+[348]+[169]+[357]+[249]+[178]+[456]
                         +[148]+a[147]+b[134].
\]
Propositions~2.13 and 3.8 of that source identify its rank-four
surface and Coble cubic with the Jacobian of
\eqref{eq:p-rank-one-normal-form-candidate}, and show $\gamma$ is stable.
The alternating polar of its Pfaffian cubic is
\begin{align*}
 B={}&[158]+[167]+[179]+[247]+[259]+[268]+[349]+[367]+[457]\\
    &+a\bigl([148]+[157]\bigr)+b\bigl([126]+[135]\bigr).
\end{align*}
Form the infinitesimal action matrix
$\mathcal D_B:\operatorname{End}(k^9)\to\bigwedge^3k^9$ as in
Lemma~\ref{lem:p-rank-zero-genus-two-polar-stable-candidate}.
Deleting rows $[123],[345],[356]$ gives the exact polynomial identity
\begin{equation}\label{eq:p-rank-one-polar-minor-candidate}
                         \det\mathcal D'_B=b^4
                        \quad\text{in }\mathbf F_2[a,b].
\end{equation}
The certificate below specifies and checks the complete computation.
Since $b\ne0$, the Lie stabilizer of $B$ is zero. As in the preceding
lemma, stability of $\gamma$ and Corollary~5.6 of the cited source
place $B$, up to change of basis and scalar, in the Cartan subspace.
Propositions~2.9--2.10 then turn zero Lie stabilizer into stability:
every nonstable Cartan tensor has a positive-dimensional stabilizer.
\end{proof}

\begin{remark}[The uniform nonzero minor]
\label{rem:p-rank-one-polar-unit-certificate-candidate}
\status{Proved under stated hypotheses}
The complete exact program and output are
\path{agents/monodromy/code/ppav_two_rank_one_polar_uniform_minor.py}
and
\path{agents/monodromy/outputs/ppav_two_rank_one_polar_uniform_minor.json}.
They independently build all nine Pfaffians from $\gamma$, extract
$B$, and build its infinitesimal action by the exterior-product rule.
Over $\mathbf F_2[a,b]$, 77 unit pivots leave eight rows and four
columns. Four of those rows, in their retained order, have the form
\[
 \begin{pmatrix}
 b&ab&A&b^2\\
 b&ab+1&A&b^2\\
 ab+1&a^2b&D&ab^2\\
 0&0&b^2&0
 \end{pmatrix},
\quad
 \begin{aligned}
 A&=a^2b^4+b^5,\\
 D&=a^3b^4+a^2b^3+ab^5+ab^2+b.
 \end{aligned}
\]
Adding its first two rows and expanding gives determinant $b^4$.
The 77 pivot rows together with these four rows are precisely all
exterior rows except $[123],[345],[356]$; the added trace row is
unused. Hence this is a certificate of
\eqref{eq:p-rank-one-polar-minor-candidate} over the polynomial ring,
without specialization. The successful author run at
18:36:54.364729 UTC on 19 November 2026 took
0.01662725299684098 seconds. A first execution failed before any
calculation because a generated assertion omitted a space;
that exact failed source and diagnostic are retained in the history
and outputs directories. The earlier exact exploratory residual-minor
calculation and its original executed source/output are also retained.
No finite sampling is used in the proof. The parent, reduction/descent,
and obstruction/crystal agents independently audited the complete proof
and original source passages. Independent polynomial reconstructions of
the displayed tensor and fixed minor took respectively
0.04347756900097011, 0.08703309600241482, and
0.0531565969940857 seconds; reduction/descent also independently
reconstructed all nine Pfaffians. Their author-output-preserving
executions of the complete retained program took
0.019327093999891076, 0.019681259997014422, and
0.019665626001369674 seconds. All checks passed.
\end{remark}

\begin{theorem}[Every fixed nonordinary principal surface in degree one]
\label{thm:every-nonordinary-ppav-degree-one-full-candidate}
\status{Proved under stated hypotheses}
Let $(A,\Theta)$ be any nonordinary principally polarized abelian
surface over a perfect field $k$ of characteristic two.
For the smooth maximal-rank locus in $|3\Theta|$, the relative
unit local system has rank eight and both its geometric and
arithmetic images are $\GL_8(\mathbf F_2)$.
\end{theorem}
\begin{proof}
The geometric two-rank-zero case is
Theorem~\ref{thm:p-rank-zero-ppav-degree-one-full-candidate}.
If the geometric two-rank is one and the principally polarized
surface decomposes geometrically, it is the product of an ordinary
and a supersingular elliptic curve, so
Theorem~\ref{thm:mixed-every-fixed-modulus-degree-one-candidate}
applies after identifying the principal polarization factors.

Otherwise $A^{(2)}$ is the Jacobian of its smooth genus-two theta
curve, whose two-rank is one. The preceding lemma proves stability
of the polar tensor of its Coble cubic. The exact Cartier--Coble
identification in
Lemma~\ref{lem:root-nonordinary-Coble-identification-candidate}
transfers this statement to the actual Cartier Pfaffian cubic under
the complete three-theta embedding. Taking coefficient square roots
preserves stability. Thus all hypotheses of
Theorem~\ref{thm:nonordinary-ppav-polar-stable-conditional-candidate}
are satisfied, and it gives the full geometric image. The arithmetic
image is then the same full general linear group. All identifications
of principal line-bundle representatives are by translations on the
geometric abelian surface, so they identify the complete linear-system
relative local systems as well.
\end{proof}

\begin{remark}[The next fixed-surface gate]
\label{rem:nonordinary-ppav-degree-one-next-gate-candidate}
\status{Proved under stated hypotheses}
This closes degree one on every fixed nonordinary principally
polarized surface. Higher-degree structural images on a fixed
nonproduct nonordinary surface remain separate. In particular this
result does not extend the product all-degree theorems to arbitrary
principal surfaces by specialization or by padding a nodal norm
with a positive-defect ample component.
\end{remark}


% Post-CP21 higher-degree intermediate; no irreducibility assertion.
\subsection{Actual natural blocks in every odd degree on a fixed nonordinary surface}
\status{Proved and independently checked.}

\begin{lemma}[Sharp pure norms in every odd principal degree]
\label{lem:nonordinary-ppav-odd-sharp-norm-candidate}
\status{Proved under stated hypotheses}
Let $(A,\Theta)$ be a principally polarized abelian surface over an
algebraically closed field of characteristic two, with two-rank $f<2$.
For every odd positive integer $d$, there is an odd polarized isogeny
$\phi:(B,\Theta_B)\to(A,\Theta)$ of degree $d^2$ and multiplier $d$.
After one constant target translation, the norms of a nonempty open
of smooth maximal-rank curves $C\in|3\Theta_B|$ form integral nodal
curves $D\in|3d\Theta|$ with normalization $C$ and
\[
             \delta(D)=9(d^2-1),\qquad
             f(D)-f=9d^2-1.
\]
Thus this entire open norm family lies in the maximal-relative-unit-rank
locus on the target. Its graph representation has image a subgroup
of a product of symmetric groups $S_{18}$.
\end{lemma}
\begin{proof}
Choose a maximal isotropic subgroup of $A[d]$, take its principal
polarized quotient, and use the polarized dual isogeny. The standard
finite theta-group descent gives the stated $B,\phi$ and
$\phi^*\lambda_\Theta=d\lambda_{\Theta_B}$; this is the same
odd principal-isogeny construction used in
Theorem~\ref{thm:DE-A3-ppav-ranks}. Its degree is $d^2$ and it is
\etale. In particular $B$ has the same two-rank $f$ as $A$.
Put $K=\ker\phi$, a constant group of odd order, and $L_B=3\Theta_B$.
The smooth maximal-rank open in $|L_B|$ is nonempty and, by
Theorem~\ref{thm:every-nonordinary-ppav-degree-one-full-candidate},
has relative image $\GL_8(\mathbf F_2)$.

We detail the nodality input, which uses no ordinarity of the ambient
surface. A general member is invariant under no nonzero $g\in K$:
if $g$ does not preserve $L_B$, invariance is impossible; otherwise
it acts nontrivially on the complete very ample projective system,
so its invariant members lie in a proper union of eigenspaces.
On a nonproduct source, for each fixed $g\ne0$ the two first jets
at $P,P+g$ have rank six except at finitely many $P$, by
Corollary~\ref{cor:all-theta-fixed-point-jets-candidate}.
Indeed its exceptional set is contained in
\[
                     3P+g\in\Theta_B,\qquad
                     3P+2g\in\Theta_B.
\]
The two theta translates are distinct and intersect finitely, and
$[3]$ is finite. On the rank-six open, the two value equations and
proportionality of the nonzero gradients impose three conditions;
$P$ varies in dimension two, so tangency is absent for a general
member. The finite exceptional set already imposes two independent
value conditions and is avoided. Three distinct points impose three
independent values on $3\Theta_B$, by the theta-product separation
argument in Lemma~\ref{lem:ordinary-ppav-isogeny-block}. Applying this
to $P,P+g,P+h$ excludes triple identifications, since their point
parameter has dimension two.

For completeness the product case uses only the cubic-product jet
calculation, not ordinary factors. If the two translated points differ
in both coordinates, Lemma~\ref{lem:product-small-degree-jets} says
that failure of rank six requires the cross-tangency equations
$3P_1+2g_1=0$, $3P_2+g_2=0$, or the equations with $g$ and $2g$
interchanged. This gives only finitely many points, since $[3]$
is finite, and two independent values avoid them. If one coordinate
agrees, the joint jets have rank exactly five everywhere: two
distinct points of a smooth plane cubic have distinct tangent lines,
as otherwise their line would have intersection multiplicity at least
four with the cubic. In the tensor first-jet calculation the two
three-dimensional jet spaces therefore meet in exactly one dimension.
After imposing the two values, the remaining three-dimensional space
of pairs of gradients cannot consist entirely of proportional pairs:
a linear subspace of rank-at-most-one $2$-by-$2$ matrices has dimension
at most two. Thus tangency again adds one condition, with no additional
axis exception. These statements hold for both ordinary and
supersingular plane cubics. Three value vectors of the Segre cubic
product are independent unless one coordinate is shared and the
three vectors in the other plane cubic are collinear. For a fixed
triple of odd translations this collinearity requires
$3P_1+g_1+h_1=0$ (or its second-factor version), so that exceptional
point locus has dimension one; two independent values still exclude
it generically. This proves the same nodality statement on a product.

There are finitely many kernel elements and pairs, so all these
nonempty conditions hold simultaneously. Thus $C\to D=\phi(C)$ is
birational, $C$ is the normalization, and $D$ is integral nodal.
For each unordered type $\{g,-g\}$, the intersections
$C\cap(C+g)$ give 18 distinct nodes, since $L_B^2=18$ and triple
identifications are excluded. Hence the number of nodes is
$18(|K|-1)/2=9(d^2-1)$. Their permutation image is a subgroup of
$\prod_{\{g,-g\}}S_{18}$.

Choose symmetric representatives of the two principal polarizations.
The norm of $L_B$ has N\'eron--Severi class $3d\Theta$ and is symmetric;
its difference from that precise target bundle is geometric
$2$-torsion in $\Pic^0(A)$. Since $3d$ is odd, translation by a
constant geometric $2$-torsion point corrects this difference,
as in Lemma~\ref{lem:ordinary-ppav-isogeny-block}.
This identifies the norm family with a family in the exact system
$|3d\Theta|$.

The normalization has genus ten and two-rank $8+f$.
Graph cohomology of an integral nodal curve adds one unit direction
per node, so
\[
             f(D)=8+f+9(d^2-1),\qquad
             f(D)-f=9d^2-1.
\]
This equals the full relative coherent dimension in the odd target
degree, and therefore proves maximal rank there. Notice that no
additional ample component has been attached to the norm.
\end{proof}

\begin{theorem}[An actual eight-dimensional block in every odd degree]
\label{thm:nonordinary-ppav-odd-natural-block-candidate}
\status{Proved under stated hypotheses}
Let $(A,\Theta)$ be any nonordinary principally polarized abelian
surface over a perfect field of characteristic two. For every odd
$d\ge1$, the degree-$d$ geometric relative monodromy contains a
subgroup conjugate to
\[
                      \operatorname{diag}(\GL_8(\mathbf F_2),I)
                  \subset\GL_{9d^2-1}(\mathbf F_2).
\]
In particular this is an actual natural block fixing a complementary
summand pointwise, rather than merely a general-linear quotient.
\end{theorem}
\begin{proof}
All constructions may be made after passing to an algebraic closure,
which is sufficient for the geometric subgroup assertion. For $d=1$
the preceding degree-one theorem gives the claim. For odd $d>1$,
use the norm family of the lemma and shrink its source open as needed.
Such shrinking does not change the full degree-one geometric image.
The normalization sequence, together with the odd-isogeny isomorphism
\[
             H^1_{\et}(A,\mathbf F_2)
                      \xrightarrow{\sim}H^1_{\et}(B,\mathbf F_2),
\]
gives the exact sequence of relative local systems
\[
 0\longrightarrow T\longrightarrow V_D\longrightarrow U_C
       \longrightarrow0,\qquad
 \dim T=9(d^2-1),\quad \dim U_C=8.
\]
Here $T$ is the graph permutation representation and the image on
$U_C$ is $\GL_8(\mathbf F_2)$. The image on $T$ has order prime to31,
since it embeds in a product of $S_{18}$, whereas31 divides
$|\GL_8(\mathbf F_2)|$. The latter group is simple.
Goursat's lemma therefore implies that the kernel of the action on
$T$ still surjects onto $\GL(U_C)$: any common nontrivial quotient
would have to be that simple group, which cannot be a quotient of
a group of order prime to31.

Choose a vector-space splitting of the displayed sequence and restrict
to this kernel. Its image lies in
$\Hom(U_C,T)\rtimes\GL(U_C)$. Its translation kernel is a
$\GL(U_C)$-submodule of $U_C^*\otimes T$, hence equals
$U_C^*\otimes T'$ for a subspace $T'\subset T$. This assertion follows
from the full matrix algebra generated by the natural general-linear
representation; it does not require a semisimplicity assertion for
arbitrary modules. Modulo $T'$ the extension is the graph of a
one-cocycle. The elementary vanishing
\[
                 H^1(\GL_8(\mathbf F_2),U_C^*)=0
\]
of Lemma~\ref{lem:char2-group-inputs} makes that cocycle a coboundary.
After conjugating the splitting, every residual upper-right column
lies in the known translation kernel, so it can be removed by an
actual kernel element. Thus the image contains the claimed
block-diagonal subgroup.

Finally the nodal norm family is inside the target maximal-rank
open, by the preceding lemma. The relative Lang system extends
\etale over this open, including its nodal curves. Its monodromy
on the dense smooth-member open is the same, because the ambient
open in the projective parameter space is normal and a dense-open
inclusion induces a surjection on geometric fundamental groups.
Consequently the subgroup just obtained belongs to the geometric
monodromy requested in the theorem.
\end{proof}

\begin{remark}[The remaining odd-degree irreducibility gate]
\label{rem:nonordinary-odd-block-scope-candidate}
\status{Proved under stated hypotheses}
The actual block above gives transvections but does not prove that
the whole degree-$d$ representation is irreducible or that its image
is the full general linear group. In particular the product common-grid
argument has not been imported for a simple surface. The remaining
graph-kernel and global invariant-subspace questions are separate.
\end{remark}


% Post-CP21 work. All geometric hypotheses below are explicit.
\subsection{A rank-drop-compatible three-theta graph reduction}
\status{Proved under the stated geometric hypotheses; independently checked.}

\begin{lemma}[Clopen closure of an invariant Lang subgroup]
\label{lem:normal-base-clopen-Lang-closure-candidate}
\status{Proved under stated hypotheses}
\textup{[Proved under the stated hypotheses; complete independent audits passed.]}
Let $B$ be normal, integral and Noetherian, let $E$ be a vector bundle
with Frobenius-semilinear map, and let $\mathcal L=\{Fv=v\}$ be its
Lang scheme. Suppose $U\subset B$ is a nonempty open where the stable
rank is constant, and let $\mathcal S_U\subset\mathcal L_U$ be a
finite \etale subgroup local system arising from an invariant
$\mathbf F_2$-subspace. Its closure $\overline{\mathcal S}$ in
$\mathcal L$ is open and closed and is a subgroup scheme.
If $C$ is integral and connected, $C\to B$ meets $U$ generically,
and $T\hookrightarrow\mathcal L_C$ is a finite \etale subgroup,
then $T\cap\overline{\mathcal S}_C$ is finite \etale of constant
rank. If $\mathcal S$ is generically contained in $T$, its entire
closure on $C$ is contained in $T$ and retains its generic rank.
If $T$ is generically contained in $\mathcal S$, it is contained
in the closure on every fibre.
\end{lemma}
\begin{proof}
The Lang equations have derivative $-I$, so $\mathcal L\to B$ is
separated, quasi-finite and \etale. It is normal and Noetherian,
and its finitely many irreducible components are disjoint and open
and closed. Each component dominates $B$: an \etale map is open,
and every nonempty open of the integral base contains its generic
point. Its intersection with $U$ is therefore a nonempty irreducible
open. Since $\mathcal S_U$ is open and closed in $\mathcal L_U$,
its closure is a union of entire components of $\mathcal L$.
The subgroup property follows by checking addition and inverse on
the dense inverse image of $U$; every component of
$\mathcal L\times_B\mathcal L$ also dominates $B$.

A finite \etale monomorphism $T\to\mathcal L_C$ is an open immersion,
and is also closed because it is finite and the target is separated
over $C$. Its intersection with the clopen closure is consequently
finite \etale, with locally constant rank on connected $C$.
If the closure has a point outside $T$, that open complement is
\etale over $C$, so its image meets the generic point, contradicting
the assumed generic containment. The reverse containment follows
from the same argument on the complement inside finite \etale $T$.
These arguments do not identify an entire rank-dropping Lang fibre
with the generic fibre.
\end{proof}

\begin{lemma}[The three-component graph with odd labels]
\label{lem:three-theta-odd-labelled-graph-candidate}
\status{Proved under stated hypotheses}
\textup{[Proved under the stated hypotheses; complete independent audits passed.]}
Let $K$ be a finite abelian group of odd order $k>1$, and put
$H=(K\setminus\{0\})/\{h\sim-h\}$. Form a graph with three vertices,
two loops $\ell_{i,h,\epsilon}$ for each vertex $i$ and type $h\in H$,
and two edges $e_{ij,g,\epsilon}$ for each $i<j$, $g\in K$,
$\epsilon\in\{0,1\}$. Let $G$ be its cochain space modulo vertex
coboundaries, so $\dim G=9k-5$.
For each $r=(r_1,r_2,r_3)$ with $r_1+r_2+r_3=0$, retain only the
six edges with $g=r_j-r_i$, and let $T_r\subset G$ be the kernel
of restriction to this doubled triangle. Then $\dim T_r=9k-9$,
and it has a direct sum decomposition into eighteen-dimensional
coordinate spaces $T_{r,h}$, consisting of the six loops of type
$h$ and the twelve cross edges with $g=r_j-r_i\pm h$.
Write $T_{r,h}^0$ for the coordinate-sum-zero hyperplane.
Then:
\begin{enumerate}
\item The span of all $T_{r,h}^0$, for all permitted $r,h$, is the
codimension-one subspace $G^0\subset G$ on which the sum of all edge
coordinates is zero.
\item If $S\subset G$ lies in every $T_r$ and is invariant under
$\prod_h S_{18}$ acting on the coordinate spaces of $T_0$, then $S=0$.
\end{enumerate}
\end{lemma}
\begin{proof}
There are $9k-3$ edges and two independent vertex coboundaries.
The selected doubled triangle has first cohomology dimension four.
The unselected edge cochains are independent modulo coboundaries:
a coboundary vanishing on its connected selected triangle is zero.
They therefore give precisely the stated coordinates of $T_r$.

Pair cochains modulo coboundaries with cycles. A cycle annihilating
$T_{r,h}^0$ has the same coefficient $c_{r,h}$ on all its eighteen
edges. The six loop edges of type $h$ are independent of $r$, so
$c_{r,h}=c_h$. Each individual difference $r_j-r_i$ ranges over all
of $K$ even with $\sum r_i=0$: assign these two entries $0,g$ and
the third $-g$. Thus every cross-edge coefficient equals $c_h$ for
every type $h$, forcing a single common coefficient $c$ on every
edge of the graph. The all-ones edge vector is a cycle, since the
cross degree at every vertex is $4k$, which is even. Hence this
annihilator is exactly one-dimensional, proving the first assertion.

For the second assertion, regard $S$ inside $T_0$. If any vector has
nonconstant coordinates in one type, subtract its image under a
transposition in that factor. This produces a pair vector in $S$;
conjugating by the whole factor gives every pair vector and hence
$T_{0,h}^0\subset S$. Choose such a pair vector supported on one
cross edge of type $h$ and one loop of the same type. Choose a
permitted $r$ which selects that cross edge. Its restriction to the
doubled triangle is nonzero, since precisely one of the two parallel
edges has value one. Thus this vector is not in $T_r$, a contradiction.
Consequently all vectors of $S$ are type-constant. Write one as
$\sum_h a_h\mathbf1_{T_{0,h}}$, and extend the coefficients to an
even function $a:K\to\mathbf F_2$ with $a(0)=0$.
For $r=(-u,0,u)$, restriction to the selected triangle vanishes only if
\[
                         a(u)+a(u)+a(2u)=0.
\]
Hence $a(2u)=0$ for every $u$. Multiplication by two is an
isomorphism of $K$, so all coefficients are zero.
\end{proof}

\begin{theorem}[Conditional odd-degree irreducibility and full image]
\label{thm:three-theta-conditional-full-image-candidate}
\status{Proved under stated hypotheses}
\textup{[Proved under the stated hypotheses; complete independent audits passed.]}
Let $A$ be a nonordinary principally polarized surface in
characteristic two and let $d>1$ be odd. Choose the polarized
isogeny $\phi:B\to A$ of degree $d^2$ used in
Lemma~\ref{lem:nonordinary-ppav-odd-sharp-norm-candidate}, and put
$K=\ker\phi$. Assume the following two additional geometric inputs.
\begin{enumerate}
\item A principal theta curve on $B$ has integral nodal norm, with
exactly two distinct transverse self-identifications for every
unordered nonzero kernel type and no triple identifications.
Three general theta translates with sum of translating points zero
have norms meeting transversely, with all their mutual intersections
and self-nodes distinct. Their independently $K$-translated source
unions with $r_1+r_2+r_3=0$ all lie in a connected coefficient open
where the additional norm identifications remain distinct ordinary
nodes, even though the source union itself is nodal.
\item On the smooth sharp norm family in the single source system
$|3\Theta_B|$, the node permutation image is the full product
$\prod_hS_{18}$, and, after a vector-space splitting, the boundary
image contains the actual $\GL_8$ block and every translation in
$\Hom(U,T_h^0)$ for each type $h$, where $\dim U=8$.
\end{enumerate}
Then the target odd-degree geometric and arithmetic images equal
$\GL_{9d^2-1}(\mathbf F_2)$.
\end{theorem}
\begin{proof}
Work geometrically and write $R=9d^2-1$, $\delta=9d^2-9$,
$t=(d^2-1)/2$. The pure norm sequence is
$0\to T\to V\to U\to0$, with $\dim U=8$.
Let $S$ be a nonzero proper global invariant subspace. Its projection
to $U$ on this norm family is either zero or all of $U$.
This choice is uniform on the connected dense maximal-rank part of
the source coefficient open:
the rank of a morphism of finite \etale local systems is locally
constant. The clopen closure argument below carries the resulting
alternative to the three-theta specializations; no constancy of
the rank-eight quotient is asserted at those rank-dropping points.
Their translating points have sum zero, so their source divisors
belong to this same linear system.
No inequality between $8$ and the number of node types is used.

Let $Y$ be the common target union of three normed theta translates.
The hypotheses identify its graph with the preceding lemma, and
identify the partial-normalization kernels with its $T_r$.
The graph cohomology embeds in the relative Lang fibre: ambient
$H^1(B,\mathbf F_2)$ restricts injectively to each principal theta
normalization, and therefore has zero intersection with the graph
part of the normalization sequence.
The relative coherent Frobenius bundle on the target projective
coefficient base is locally free of rank $R$, including $Y$:
for every ample divisor the ambient coherent $H^1$ injects and the
relative quotient has constant dimension, by the ample-vanishing
structure-sheaf sequence. Apply the clopen Lang closure lemma to
$S$ on the maximal-rank open and denote its special fibre by
$\overline S_Y$. It is the same subspace of the relative Lang fibre
at $Y$ for every source specialization. The graph subgroups
$T_r$ and their hearts are finite \etale on the source neighbourhoods
by the distinct-node hypothesis.

If the generic projection of $S$ to $U$ is zero, the closure lemma
puts $\overline S_Y$ inside every $T_r$, preserves its full dimension,
and makes it invariant under the product node group on $T_0$.
The second graph assertion gives $\overline S_Y=0$, a contradiction.

If the projection is surjective, the stated translations put every
$T_h^0$ into $S$ on the generic norm family. The closure lemma then
puts every $T_{r,h}^0$ into $\overline S_Y$, so its first graph
assertion gives $G^0\subset\overline S_Y$.
For a fixed $r$, the intersection $T_r\cap G^0$ has dimension
$\delta-1$: the total-edge-sum functional is nonzero on each
unselected-edge basis vector of $T_r$.
The finite \etale subgroup $T_r\cap\overline S$ has constant rank
on its connected norm-family base. Therefore
\[
            \dim(S\cap T)\ge\delta-1,
        \qquad \dim S\ge(\delta-1)+8=R-1.
\]
A proper $S$ must be a hyperplane. Over $\mathbf F_2$ its quotient
is a trivial one-dimensional representation, contradicting the
proved absence of invariant covectors in every odd nonordinary
principal degree. Thus the global representation is irreducible.

Let $J$ be its normal subgroup generated by transvections.
The nonzero $J$-socle is invariant under the full irreducible group,
so the representation restricted to $J$ is semisimple.
The actual natural $\GL_8$ block lies in $J$, and its natural
constituent occurs with multiplicity one. It therefore belongs to
one simple $J$-summand. All rank-one gluing maps into the hearts
place those hearts in the same summand. Its dimension is at least
\[
                       8+17t>\frac{8+18t}{2}=R/2.
\]
Its conjugates cannot be distinct disjoint summands, so it is
invariant under the full group and equals the whole representation.
Hence $J$ is irreducible. The characteristic-two transvection
classification, as applied in
Theorem~\ref{thm:mixed-every-fixed-modulus-all-odd-candidate}, leaves
only the full general linear group or a group preserving a
nondegenerate alternating form. The actual natural $\GL_8$ block
fixing a complement has no invariant bilinear form or functional,
so would lie in the radical of such a form. This excludes the latter
case, proving the full geometric image; the arithmetic image follows.
\end{proof}

\begin{remark}[The exact geometric gates retained]
\label{rem:three-theta-conditional-geometric-gates-candidate}
\status{Proved under stated hypotheses}
\textup{[Proved under the stated hypotheses; complete independent audits passed.]}
Neither additional hypothesis of the theorem is asserted uniformly
here. In particular, theta translates do not move their intrinsic
self-intersections with a fixed kernel translate. For higher odd
kernel order these intersections can be tangent, and theta norms
can also have triple identifications. Pairwise transversality alone
does not exclude the latter. Likewise an actual $\GL_8$ block does
not by itself supply the full product node groups or all gluing
hearts. The unconditional result currently available is the actual
block in every odd degree. This reduction isolates the remaining
geometric inputs while treating the common boundary's stable-rank
drop rigorously.
\end{remark}


% Post-CP22 candidate. No assertion is imported from characteristic-zero sources.
\subsection{Three total ramifications force ordinarity in genus two}
\status{Proved and independently checked.}

\begin{lemma}[Three totally ramified fibres of a trigonal map]
\label{lem:char2-genus-two-three-total-ramifications-candidate}
\status{Proved under stated hypotheses}
\textup{[Proved under the stated hypotheses; complete independent audits passed.]}
Let $C$ be a smooth projective curve of genus two over an algebraically
closed field of characteristic two. If a degree-three morphism
$C\to\mathbf P^1$ has at least three distinct totally ramified
points, then $C$ is ordinary.
\end{lemma}
\begin{proof}
The morphism is separable, since its degree is odd. The different has
degree eight by Riemann--Hurwitz. Each of the stated tame ramification
points has different exponent two. The remaining different degree is
two. Consequently there are exactly two possibilities: a fourth tame
point of ramification index three, or one point of ramification index
two and different exponent two. A wild quadratic point has different
exponent at least two, so there are no other possibilities.
These different formulas are the characteristic-independent statements
in \citet[Section~53.12, especially Lemmas~53.12.2 and~53.12.4]{Stacks0C1B}.

Let $Y$ be the smooth Galois closure. Its transitive permutation group
is either $C_3$ or $S_3$. If all four ramification points are tame of
index three, the sign quotient has no ramification. A nontrivial
unramified double cover of $\mathbf P^1$ is impossible by
Riemann--Hurwitz. Thus the group is $C_3$ in this case.
Choose infinity away from the four branch points. Kummer theory
\citep[Lemma~9.24.1, Tag~09DX]{Stacks09DX}, which applies because
three is prime to the characteristic, gives
\[
                   y^3=F(z)/G(z),
\]
where $F$ and $G$ are monic separable quadratics with no common root.
Indeed the four nonzero branch exponents are two $+1$ and two $-1$
modulo three, since their sum is zero. Put
$\eta=yG$, so that $\eta^3=FG^2$.
The two forms
\[
             \omega=dz/\eta,\qquad \omega'=G\,dz/\eta^2
\]
are regular and independent. At a branch point their orders are zero
or one; at infinity both have order zero. They therefore form a
basis of the genus-two space of regular differentials.
Write $F=z^2+f_1z+f_0$ and $G=z^2+g_1z+g_0$.
The Cartier rules in characteristic two give
\[
              C\omega=\sqrt{f_1}\,\omega',\qquad
              C\omega'=\sqrt{g_1}\,\omega.
\]
For example $C(dz/\eta)=\eta^{-2}C(FG^2\,dz)
=\eta^{-2}G C(F\,dz)$. Both $f_1$ and $g_1$ are nonzero,
because the respective quadratics are separable. Cartier is
therefore bijective, proving ordinarity in the cyclic case.

In the remaining case the Galois group is $S_3$. Let
$Z=Y/C_3$ be its quadratic sign quotient. It is unramified above the
three tame branch values. At the remaining branch value its completed
extension is exactly the wild quadratic extension of different
exponent two: the third sheet of the original trigonal cover is
unramified, and the ramified local degree-two extension is already
Galois. Thus $Z\to\mathbf P^1$ has total different degree two,
and $Z$ is rational. Choose a coordinate $z$ on $Z$ in which the
nontrivial involution is $z\mapsto z+1$, whose only fixed point is
infinity.

The cyclic cubic cover $Y\to Z$ is unramified above infinity and
ramified at six distinct finite points, in three pairs interchanged
by $z\mapsto z+1$. The two exponents in each pair are opposite,
since this involution acts on $C_3$ by inversion. Choosing the
positive exponent in each pair gives a monic separable cubic
\[
 F(z)=z^3+a z^2+b z+c,\qquad
 F^\sigma(z)=F(z+1),\qquad
                    y^3=F(z)/F^\sigma(z).
\]
Here $F$ and $F^\sigma$ have disjoint roots. Multiplying a Kummer
coordinate by a cube root of unity makes the chosen involution of
$Y$ exactly
\[
                         \sigma(z,y)=(z+1,y^{-1}).
\]
The quotient $Y/\langle\sigma\rangle$ is the original curve $C$,
up to the choice of a conjugate order-two subgroup of $S_3$.

The curve $Y$ has genus four. A basis of its regular differentials is
\[
 \omega_i=\frac{z^i\,dz}{yF^\sigma},\qquad
 \omega'_i=\frac{z^i\,dz}{y^2F^\sigma},\qquad i=0,1.
\]
To verify regularity directly, at a root of $F$ the respective orders
are one and zero, and at a root of $F^\sigma$ they are zero and one.
At each of the three points over infinity the orders are $1-i$.
The four forms are independent by their two distinct $C_3$
characters and the independence of $1,z$.
The involution sends
\[
 \omega_0\longleftrightarrow\omega'_0,\qquad
 \sigma\omega_1=\omega'_1+\omega'_0,\qquad
 \sigma\omega'_1=\omega_1+\omega_0.
\]
Its invariant subspace is thus two-dimensional, with basis
\[
       \Omega_0=\omega_0+\omega'_0,\qquad
       \Omega_1=\omega_1+\omega'_1+\omega'_0.
\]
Pullback along the separable map $Y\to C$ injects the two-dimensional
space $H^0(C,\omega_C)$ into this invariant subspace, hence identifies
it with the whole subspace. Cartier commutes with this pullback.
This also follows directly from its defining square-semilinearity
and the identities $C(dh)=0$ and $C(h\,dh)=dh$: a separating
parameter on $C$ remains a separating parameter in the separable
function-field extension.

Set $d=a+b+1$. Then
\[
 F^\sigma=z^3+(a+1)z^2+(b+1)z+(c+d).
\]
The four Cartier computations are
\begin{align*}
 C\omega_0&=\sqrt b\,\omega'_0+\omega'_1,&
 C\omega_1&=\sqrt c\,\omega'_0+\sqrt a\,\omega'_1,\\
 C\omega'_0&=\sqrt{b+1}\,\omega_0+\omega_1,&
 C\omega'_1&=\sqrt{c+d}\,\omega_0+\sqrt{a+1}\,\omega_1.
\end{align*}
In the invariant basis $(\Omega_0,\Omega_1)$ the Cartier matrix is
\[
           H=\begin{pmatrix}
                \sqrt{b+1}&\sqrt{a+c}\\
                1&\sqrt a
              \end{pmatrix},\qquad
                         (\det H)^2=ab+c.
\]
The last scalar is nonzero. Indeed $F'(z)=z^2+b$ has its unique
geometric root $\sqrt b$, and $F(\sqrt b)=ab+c$; its vanishing
would contradict separability of $F$. Thus Cartier on $C$ is
bijective in the $S_3$ case as well.
\end{proof}

\begin{corollary}[No triple three-torsion packet on a nonordinary theta curve]
\label{cor:nonordinary-genus-two-no-three-torsion-triple-candidate}
\status{Proved under stated hypotheses}
\textup{[Proved under the stated hypotheses; complete independent audits passed.]}
Let $C$ be nonordinary of genus two in characteristic two. There
are no three distinct points $P,Q,R\in C$ for which
\[
                         3P\sim3Q\sim3R.
\]
Consequently, for any subgroup $K\subset\operatorname{Jac}(C)[3]$,
each coset of $K$ meets an Abel--Jacobi theta curve in at most two
geometric points. This conclusion does not require isotropy of $K$.
\end{corollary}
\begin{proof}
Riemann--Roch gives $h^0(C,\mathcal O_C(3Q))=2$. The three stated
effective divisors define three distinct members of this pencil.
It is basepoint-free because two of these divisors have disjoint
supports. It therefore gives a degree-three map totally ramified
at all three points, contradicting the lemma. Points in a common
$K$-coset have differences annihilated by three, which gives the
same linear equivalences.
\end{proof}

\begin{corollary}[Every three-torsion theta norm is nodal]
\label{cor:nonordinary-theta-three-torsion-nodal-norm-candidate}
\status{Proved under stated hypotheses}
\textup{[Proved under the stated hypotheses; complete independent audits passed.]}
Let $B$ be a nonordinary indecomposable principally polarized
abelian surface in characteristic two, with a smooth symmetric
theta curve $C$, and let $K\subset B[3]$. The norm of $C$ under
$B\to B/K$ is integral and nodal, its normalization is $C$,
and it has exactly $|K|-1$ nodes. For each unordered nonzero type
$\{g,-g\}$ there are exactly two nodes.
\end{corollary}
\begin{proof}
Fix a Weierstrass point $W$ to identify $C$ with its Abel--Jacobi
image. The restriction of invariant differentials on $B$ to $C$
is an isomorphism onto $H^0(C,\omega_C)$, so equality of its
translated tangent lines is equivalent to equality of canonical
images. Thus a tangency between $C$ and $C+g$, with $g\ne0$,
would give $g=2(P-W)$ for a point $P$ with $P\ne\iota P$.
Since $3g=0$, the class $3(P-W)$ is geometric two-torsion.

If the two-rank is zero, this class is zero, so $3P\sim3W$.
But $|3W|=W+|K_C|$ by Riemann--Roch and $2W\sim K_C$;
thus $P=W$, a contradiction. If the two-rank is one, the
characteristic-two normal form in
Lemma~\ref{lem:p-rank-one-genus-two-polar-stable-candidate}
has exactly two Weierstrass points $W,W'$. Their nonzero
two-torsion difference generates the geometric two-torsion.
The remaining possibility is $3(P-W)=W'-W$, hence
$3P\sim K_C+W'$. Again $|K_C+W'|=W'+|K_C|$ by
Riemann--Roch, so $P=W'$, which is fixed by $\iota$ and gives
$g=0$. Thus every nonzero kernel translate intersects $C$
transversely, in exactly two points since $C^2=2$.

The preceding corollary excludes all three-point fibres of the
norm map. It is generically one-to-one: a nonzero translation
cannot stabilize the principal divisor, since the kernel of its
polarization is trivial. The quotient is \etale, so the local
image of the smooth curve has precisely ordinary nodes, each
arising from one pair of transverse branches. The two points of
$C\cap(C+g)$ give two such nodes for each type $\{g,-g\}$;
no identifications among these nodes remain because a fibre has
at most two points. The total is therefore $|K|-1$.
\end{proof}

\begin{remark}[Exact polynomial verification]
\label{rem:three-total-ramification-Cartier-certificate-candidate}
\status{Proved under stated hypotheses}
\textup{[Proved under the stated hypotheses; complete independent audits passed.]}
The executed program
\path{agents/monodromy/code/genus_two_three_total_ramification_checks.py},
with full output
\path{agents/monodromy/outputs/genus_two_three_total_ramification_checks.json},
reconstructs all four Cartier columns from odd polynomial coefficients,
checks the involution and its two-dimensional fixed basis, verifies
the restricted matrix, and computes the cubic separability resultant
$(ab+c)^2$. The author run at 19:25:11.722810 UTC on 19 November
2026 passed in 0.06777819900162285 seconds. These are universal
polynomial identities; the Galois-closure and ramification arguments
are proved above and are not inferred from finite examples.
\end{remark}

\begin{remark}[Scope]
\label{rem:nonordinary-three-torsion-triple-scope-candidate}
\status{Proved under stated hypotheses}
\textup{[Proved under the stated hypotheses; complete independent audits passed.]}
This establishes nodality of individual theta norms for kernels killed
by three, not for general odd kernels. It does not by itself prove
independence of node permutations or the gluing-translation input of the three-theta
conditional monodromy theorem. The calculation is intrinsic to
characteristic two; no characteristic-zero nodality assertion is
being specialized.
\end{remark}


% Post-CP22 candidate. All characteristic-two local steps are explicit.
\subsection{Three-torsion node permutations and gluing hearts}
\status{Proved and independently checked.}

\begin{lemma}[Finite eigenmaps for a three-torsion translation]
\label{lem:nonproduct-three-theta-finite-eigenmaps-candidate}
\status{Proved under stated hypotheses}
\textup{[Proved under the stated hypotheses; complete independent audits passed.]}
Let $(B,\Theta)$ be an indecomposable principally polarized abelian
surface over an algebraically closed field of characteristic two,
and let $0\ne g\in B[3]$. Each of the three eigenspaces of a lift
of $g$ acting on $H^0(B,3\Theta)$ has dimension three, is
basepoint-free, and defines a finite morphism $B\to\mathbf P^2$.
\end{lemma}
\begin{proof}
Put $E=B/\langle g\rangle$ and let $\pi:B\to E$ be the quotient.
The cyclic subgroup is isotropic for the theta pairing, so
$L=\mathcal O_B(3\Theta)$ descends to an ample line bundle $M$
on $E$. Its three possible descents differ by the three characters
of the cyclic cover; write them $M_i$. The direct-image character
decomposition identifies the three eigenspaces with $H^0(E,M_i)$.
Riemann--Roch gives $M_i^2=6$ and $\chi(M_i)=3$; ample vanishing
gives $h^0(M_i)=3$. These standard cohomology formulas and their
characteristic-independent proofs are
\citet[Theorems~9.11 and~9.14, Proposition~9.18]{EdixhovenVanDerGeerMoonenAV}.

We give the basepoint argument, including fixed divisors. First no
elliptic curve $F\subset E$ has $M_i\cdot F=1$. Translate $F$ to
an elliptic subgroup and let $F'$ be the identity component of its
inverse image in $B$. The restriction $F'\to F$ has degree
$e\in\{1,3\}$. The projection formula gives
$3\Theta\cdot F'=e$, hence $e=3$ and $\Theta\cdot F'=1$.
An elliptic curve of principal degree one splits a principally
polarized surface as a product, contrary to the hypothesis.

Suppose the complete system $|M_i|$ has a nonzero fixed divisor
$D$, and write $M_i\sim D+N$, with $N$ the moving part.
Every effective divisor on an abelian variety is nef: translate
it away from any prescribed integral curve to compute a
nonnegative intersection. Thus both $D,N$ are nef.
If $N^2>0$, then $N$ is ample on the surface. Indeed the Hodge
index theorem shows that a curve with $N$-intersection zero
would have negative square, impossible for an effective divisor
on an abelian surface; Nakai--Moishezon then applies.
Ample vanishing and $h^0(N)=h^0(M_i)=3$ give $N^2=6$.
But $(M_i-N)\cdot(M_i+N)>0$, since $D\ne0$, $M_i$ is ample,
and $N$ is nef, contradicting $M_i^2=N^2=6$.

If $N^2=0$, each irreducible component has square zero and any
two have intersection zero. By adjunction its normalization has
genus at most one; an abelian variety contains no rational curve,
so the components are translates of elliptic subgroups. Distinct
elliptic directions have positive intersection, hence they are
all fibres of one elliptic quotient. Thus $N$ is the pullback of
a degree-$m$ divisor on the elliptic quotient and $h^0(N)=m$.
It follows that $m=3$. Writing its fibre class as $F$, we obtain
\[
              6=M_i^2=D^2+6D\cdot F.
\]
Ampleness makes $D\cdot F\ge1$, while $D^2\ge0$, so
$D\cdot F=1$ and $M_i\cdot F=1$, again impossible.
Consequently $|M_i|$ has no fixed divisor.

Its remaining base locus is finite. The theta group translates
it by $K(M_i)$, an \etale group of order $\chi(M_i)^2=9$.
Translations act freely on the surface, so a base point would
give at least nine distinct base points. Two general independent
members, having no common curve component, intersect with total
length $M_i^2=6$. This contradiction proves global generation.
The associated morphism $E\to\mathbf P^2$ is finite, since its
pullback of $\mathcal O(1)$ is the ample bundle $M_i$; its
composition with $\pi$ is finite as well.
\end{proof}

\begin{lemma}[The theta difference cover and its independent shifted branches]
\label{lem:theta-difference-cover-char2-candidate}
\status{Proved under stated hypotheses}
\textup{[Proved under the stated hypotheses; complete independent audits passed.]}
Let $C=\Theta\subset B$ be a smooth symmetric principal theta
curve in characteristic two. Away from zero the difference map
\[
            \beta:C\times C\longrightarrow B,\qquad(P,Q)\mapsto P-Q
\]
is finite flat of degree two and generically \etale. Its only
divisorial branch support is the reduced curve
$\Delta=[2]C$, of numerical class $2\Theta$.
The inertia permutation at its generic point is the nontrivial
transposition. No nonzero odd torsion translation preserves
$\Delta$.

For three theta translates with translating points
$a_1+a_2+a_3=0$, the two-point intersection covers for any one
pair and any prescribed shift by an odd torsion point have an
inertia transposition which fixes all covers belonging to other
pair/shift branch divisors.
\end{lemma}
\begin{proof}
For a nonzero difference, the fibre is the intersection of two
distinct theta translates and has length two. Properness gives
finiteness away from zero; the source is Cohen--Macaulay and the
target is smooth of the same dimension, giving flatness.
The differential fails to be invertible exactly when the two
translated tangent lines agree. These are the canonical tangent
directions of the genus-two curve. Apart from the diagonal, whose
image is zero, equality means $Q=\iota P$. Thus the only
divisorial branch support is $[2]C$.
The map is generically \etale since the canonical map is
nonconstant. Its quadratic generic extension is Galois, with
involution $(P,Q)\mapsto(\iota Q,\iota P)$.

By Lemma~\ref{lem:all-theta-double-image-candidate}, the ramification
curve $Q=\iota P$ maps to its reduced image with inseparable
degree two and separable degree one, and that image has class
$2\Theta$. In particular the quadratic cover is not \etale at
the generic point of this divisor. Over its strict henselian
local base its nontrivial separable generic extension therefore
has inertia of order two. This includes the possible purely
inseparable residue extension; no tame-ramification assertion
is made in characteristic two.

If translation by $h$ preserves $\Delta$, it preserves its line
bundle, hence $h\in K(\mathcal O(\Delta))=B[2]$; numerical
equivalence suffices here because the polarization homomorphism
depends only on the N\'eron--Severi class. An odd torsion $h$ is
therefore zero. Distinct odd shifts of $\Delta$ are distinct
irreducible divisors.

Finally the map from $(a_1,a_2)$, with $a_3=-a_1-a_2$, to any
two independent pair differences is an isogeny with determinant
three, hence \etale in characteristic two. Branch divisors
for a fixed pair are distinct by the preceding translation
argument. Those for different pairs are distinct because the
two corresponding differences vary independently. At a generic
point of any chosen branch divisor all the other covers are
\etale. Its inertia consequently acts as the one specified
transposition and as the identity on the other covers.
\end{proof}

\begin{proposition}[The full product of eighteen-node permutation groups]
\label{prop:nonordinary-three-torsion-node-product-candidate}
\status{Proved under stated hypotheses}
\textup{[Proved under the stated hypotheses; complete independent audits passed.]}
Let $(B,\Theta)$ be nonordinary and indecomposable, and let
$K\subset B[3]$ be a subgroup of order nine defining a polarized
quotient $\phi:B\to A$. On the generic norm family of curves in
$|3\Theta|$, the additional nodes have permutation image
\[
                      \prod_{h\in(K\setminus0)/\pm} S_{18}.
\]
\end{proposition}
\begin{proof}
Fix $0\ne g\in K$. Its eighteen oriented node points are
$P$ satisfying $s(P)=s(P+g)=0$; this is in bijection with the
eighteen nodes of unordered type $\{g,-g\}$ when the norm is
nodal. The one-point incidence is a projective-six bundle over
$B$, hence irreducible. We prove two-transitivity by examining
two distinct such points $P,Q$.

Choose the eigenbasis in the first lemma. An evaluation vector
has three nonzero components $v_0(P),v_1(P),v_2(P)$, and the
orbit plane is their direct sum. The chord
$\langle v(P),v(P+g)\rangle$ contains no eigen-coordinate vertex:
canceling two different eigenvalues in a combination of its
two spanning vectors would cancel the whole combination.
If the two chords for $P,Q$ meet, their orbit planes must
therefore meet in at least two coordinate directions. Thus at
least two of the eigen-projections of $P,Q$ coincide. Each
eigenmap is finite, so the locus of such pairs has dimension
at most two. Outside it the four value conditions are
independent. The rank-four open in $B\times B$ has dimension
four and is irreducible; its coefficient incidence is a
projective-four bundle of total dimension eight.

At four distinct points the rank is at least three, since any
three values of $3\Theta$ are independent. The exceptional
pair locus just bounded therefore contributes dimension at
most $2+(8-3)=7$. If $Q=P+g$ or $P-g$, only three distinct
points occur, with three independent values and a two-dimensional
point locus, giving the same bound. Hence these loci cannot
contribute a component of the finite \etale ordered-two-point
cover over the generic coefficient open. That cover is
irreducible, proving two-transitivity.

The needed transpositions can be obtained on a three-theta
boundary without imposing an unproved structured higher-jet
condition. By
Corollary~\ref{cor:nonordinary-theta-three-torsion-nodal-norm-candidate},
each individual theta norm has only its eight distinct nodes.
Choose three general translates with translating points of sum
zero. Their norms have distinct self-nodes, mutually transverse
intersections, and no triple points. To check the last assertion
directly, the map
\[
 (P,a_1,a_2)\longmapsto(P-\phi a_1,\ P-\phi a_2,\ P+\phi a_1+\phi a_2)
\]
is a finite \etale isogeny of three surfaces (its scalar
determinant is three). The inverse image of the product of
the three norm curves has dimension three, while the translating
parameter space has dimension four. Pair tangencies and incidence
with the finitely many self-nodes are likewise proper conditions,
by generic separability of the difference map.

In a fixed node type, six points are intrinsic self-nodes of the
three translated theta curves and are constant. The remaining
twelve form six two-point covers: the two orientations of each
of the three component pairs. Their branch supports are exactly
the shifted difference divisors of the preceding lemma. A
generic point of one such divisor supplies a transposition in
this type, fixing every other type and every other node.
These inertia elements belong to the generic norm-node image:
the finite \etale additional-node cover extends over the normal
coefficient open containing the generic three-theta boundary;
its smooth-source dense open surjects on the relevant geometric
fundamental group. Two-transitivity and one transposition give
$S_{18}$ in each projection. Normal conjugation of an isolated
factor transposition by this surjective projection yields the
whole $S_{18}$ supported in that factor. The full product follows.
\end{proof}

\begin{proposition}[All type hearts occur as gluing translations]
\label{prop:nonordinary-three-torsion-gluing-hearts-candidate}
\status{Proved under stated hypotheses}
\textup{[Proved under the stated hypotheses; complete independent audits passed.]}
In the setting of the preceding proposition, write the relative
norm sequence on its smooth maximal-rank open as
\[
                    0\longrightarrow T\longrightarrow V
                       \longrightarrow U\longrightarrow0,
       \qquad T=\bigoplus_h T_h,\quad \dim T_h=18,\quad\dim U=8.
\]
Its monodromy contains every translation in
$\Hom(U,T_h^0)$, where $T_h^0$ is the coordinate-sum-zero
hyperplane, for every type $h$.
\end{proposition}
\begin{proof}
Let $N=\prod_hS_{18}$. The joint image on $T$ and $U$ is
$N\times\GL_8(\mathbf F_2)$: each projection is surjective,
and the simple group $\GL_8(\mathbf F_2)$ has the prime divisor
31, which does not divide $|N|$, so Goursat's lemma excludes
a common nontrivial quotient. The kernel of this joint image
inside $\Hom(U,T)$ is $U^*\otimes T'$ for an $N$-stable
subspace $T'\subset T$. This is the elementary full-matrix-algebra
consequence for the natural $\GL_8$ module used in the actual
block proof. Modulo $T'$ the translation kernel is zero.
The vanishing $H^1(\GL_8,U^*)=0$ conjugates the restriction
of this extension to $\GL_8$ to a block-diagonal action.
The remaining $N$ factor commutes with $\GL_8$; its possible
upper-right column is therefore a $\GL_8$-invariant element
of $U^*\otimes(T/T')$, and is zero. We obtain an equivariant
splitting
\[
                         V/T'=(T/T')\oplus U.             \tag{NTG1}
\]

We justify using this splitting at the smaller-rank three-theta
boundary. Let $\mathcal B$ be the normal open of the source
projective coefficient space where all additional norm
identifications are distinct transverse pairs of smooth source
points. It contains the generic three-theta members; source
internal nodes and their selected-pair tangencies are allowed.
The partial-normalization structure-sheaf sequence gives on
$\mathcal B$ an exact sequence of coherent Frobenius bundles
\[
                         0\to T\otimes\mathcal O
                              \to\mathcal V\to\mathcal U\to0.
\]
The ambient $H^1(\mathcal O)$ maps are identified by the odd
isogeny. Both coherent relative bundles have constant ranks,
although their unit parts may drop. The graph term is a
Frobenius-unit bundle associated with the finite \etale
permutation system $T$.

The subspace $T'$ extends to a finite \etale subgroup of $T$
on $\mathcal B$: the dense maximal-rank open surjects on
$\pi_1(\mathcal B)$, and its action on $T$ already factors
through that group. On Lang schemes the projection
\[
              \mathcal L_{\mathcal V}/T'
                         \longrightarrow\mathcal L_{\mathcal U}
\]
is a finite \etale torsor under $T/T'$. This can be checked
after trivializing the graph bundle, when the equations for
a lift are $t_i^2-t_i=b_i(u)$.
The generic splitting (NTG1) is a section of this torsor on
the dense maximal-rank open. The full Lang scheme
$\mathcal L_{\mathcal U}$ is normal, being \etale over
the normal base. Every component dominates the base. A
section of a finite \etale morphism over a dense open of a
normal integral component extends uniquely: its closure is
finite and birational and hence is that normal component.
Thus (NTG1) extends to every finite Lang part at the boundary,
compatibly with addition and all monodromy. No full-rank
specialization isomorphism has been asserted.

Now work in the three-theta translating family and take inertia
at the branch divisor of a selected, unshifted component-pair
intersection. The difference-cover lemma makes this inertia
swap exactly its two selected nodes and fix all additional
norm nodes. Its action on cohomology of the nodal union has
rank one, with image
\[
                     q_{ij}=e_{ij,1}+e_{ij,2}
\]
in graph cohomology. Here is the local gluing justification in
characteristic two. The component curves are constant after
translation, so their finite \etale torsors are constant. Near
the coalescing pair, each such torsor is \etale and is trivial
over the strict henselian local branch neighbourhood. Its two
gluing labels can therefore be chosen in the same local
trivialization. Inertia exchanges these two labels and fixes
all other labels. Hence its image is exactly the displayed
one-dimensional graph line; a possible change of splitting
does not alter that image. This argument applies also when
the residue extension of the quadratic collision is purely
inseparable.

On the source doubled-triangle graph all three classes
$q_{12},q_{13},q_{23}$ coincide with one nonzero class $q$,
because their pairwise sums are vertex coboundaries. This
graph class remains nonzero in relative cohomology, since
ambient $H^1(B,\mathbf F_2)$ injects into each theta
normalization and meets the graph part trivially.
In (NTG1) the additional graph factor is fixed by each of
these inertia elements. The common source image $q$ must
therefore have the same lift in $V/T'$ for all three pairs.
Consequently
\[
                              q_{12}+q_{13}\in T'.       \tag{NTG2}
\]
In the target graph, this sum is the sum of all unselected
edges incident with vertex one, since the full vertex star
is a coboundary. In each type its eighteen coordinates have
eight ones on incident cross edges and ten zeros on the
other cross edges and the six loops. It is nonconstant in
every type. A transposition in any one $S_{18}$ factor
subtracts two unequal coordinates in that factor only;
normal conjugation then gives its entire sum-zero heart.
Thus $T_h^0\subset T'$ for every $h$, which is precisely
the claimed gluing-translation subgroup.
\end{proof}

\begin{theorem}[The degree-three image from a nonproduct three-neighbour]
\label{thm:nonordinary-degree-three-from-neighbour-candidate}
\status{Proved under stated hypotheses}
\textup{[Proved under the stated hypotheses; complete independent audits passed.]}
Let $\phi:(B,\Theta_B)\to(A,\Theta_A)$ be an \etale polarized
isogeny of multiplier three and degree nine in characteristic two,
with $B$ nonordinary and indecomposable. Then the relative
degree-three geometric monodromy on the smooth maximal-rank
open of $|9\Theta_A|$ is $\GL_{80}(\mathbf F_2)$.
The arithmetic image over any perfect field of definition is
the same full group.
\end{theorem}
\begin{proof}
The individual theta norms are nodal by
Corollary~\ref{cor:nonordinary-theta-three-torsion-nodal-norm-candidate}.
The three general translates and all their kernel-translated
partial normalizations satisfy the first geometric hypothesis
of Theorem~\ref{thm:three-theta-conditional-full-image-candidate},
by the incidence argument in the node-product proposition.
All the source unions whose translating points have sum zero
belong to the same $|3\Theta_B|$ and to the same connected
coefficient open. The preceding two propositions establish
the full type permutation groups and every gluing heart.
The actual natural $\GL_8$ block is already supplied by
Theorem~\ref{thm:nonordinary-ppav-odd-natural-block-candidate}.
Thus both geometric hypotheses of the conditional theorem
hold. Applying it with $d=3$ proves the geometric image
$\GL_{9\cdot3^2-1}=\GL_{80}$. All constructions may be
made after geometric base change; the arithmetic image
contains this full geometric group in the same general
linear group, so equality follows over the stated base.
\end{proof}

\begin{remark}[Remaining scope]
\label{rem:nonordinary-degree-three-neighbour-scope-candidate}
\status{Proved under stated hypotheses}
\textup{[Proved under the stated hypotheses; complete independent audits passed.]}
The theorem as stated assumes a nonproduct polarized three-neighbour.
Existence of such a neighbour for every fixed nonproduct nonordinary
surface is a separate assertion. The eigenmap proof uses that the
kernel is killed by three; neither it nor the theta nodality proof
has been extended here to arbitrary odd multiplier.

The finite graph supplement
\path{agents/monodromy/code/degree_three_theta_graph_checks.py}, with
full output
\path{agents/monodromy/outputs/degree_three_theta_graph_checks.json},
checks all 81 sum-zero translating triples in $K=\mathbf F_3^2$.
It reproduces graph dimension 76, partial-normalization kernel
dimension 72, the four-heart dimension 68, the total heart span
dimension 75, the rank-four constraints on type-constant vectors,
and the eight-one/ten-zero vertex star in each type. The author
run at 19:41:14.217748 UTC on 19 November 2026 passed in
0.012680367006396409 seconds. The parent and reduction/descent agent
independently read every proof and the complete program, and reran
the latter without changing the author output in respectively
0.014382761997694615 and 0.013090728003589902 seconds; both passed.
This finite supplement does not
replace any geometric incidence, inertia, or gluing argument.
\end{remark}


% Root proof promoted after full monodromy and obstruction/crystal independent audits.
\subsection{A direct exclusion of three total trigonal ramifications}

\begin{lemma}[The trigonal pencil on a nonordinary genus-two curve]
\label{lem:root-nonordinary-trigonal-three-excluded-candidate}
\status{Proved; both independent proof and exact-certificate audits passed.}
Let $C$ be a smooth projective genus-two curve over an algebraically
closed field $k$ of characteristic two, with two-rank zero or one.
No degree-three morphism $C\to\mathbf P^1_k$ has three distinct
totally ramified points.
\end{lemma}
\begin{proof}
Suppose such a morphism exists, and choose one totally ramified point
$Q$. It cannot be Weierstrass: the gap sequence at a genus-two
Weierstrass point is $1,3$, so no rational function has its sole pole
of exact order three there. Choose a Weierstrass point $W$ over
infinity for the hyperelliptic coordinate. The already proved
nonordinary normal forms allow the equation to have $h$ constant
or linear and a degree-five right side. For two-rank one we choose
the Weierstrass point with the order-three Artin--Schreier pole over
infinity, as in the preceding normal-form proof. Since $Q$ is not
Weierstrass, it is a finite point with $h(Q)\ne0$.

Translate its hyperelliptic coordinate to zero, translate its
$y$-coordinate to zero, and scale $y$ so that $h(0)=1$.
An additional change $y\mapsto y+v x^2$, with $v$ chosen to cancel
the degree-four coefficient, is possible over $k$ and does not move
$Q$. Thus we obtain
\[
 y^2+(1+t x)y=u x^5+a x^3+b x^2+c x,
 \qquad Q=(0,0),\quad u\ne0.                         \tag{RTG1}
\]
Here $t=0$ in two-rank zero and $t\ne0$ in two-rank one. These
changes are invertible on the function field and preserve the
smooth projective curve and its specified ramification question.
Put $e=b+c^2+ct$ and define
\[
                 z=\frac{y+1+(c+t)x+e x^2}{x^3}.       \tag{RTG2}
\]
At $Q$ the numerator is a unit, so $z$ has a pole of order three.
At its hyperelliptic conjugate $Q'=(0,1)$, the coefficient expansion
of (RTG1) gives
\[
        y=1+(c+t)x+(b+c^2+ct)x^2+O(x^3),
\]
so there is no pole there. There are no other finite possible poles.
At $W$, the pole orders of $x,y$ are two and five; hence $z$ has
a zero of order exactly one. Thus $z$ defines a separable degree-three
map with pole divisor $3Q$, and is unramified at $W$.
Riemann--Roch gives $h^0(3Q)=2$, so every degree-three map totally
ramified at $Q$ differs from $z$ by a target fractional linear
transformation. After moving that ramified value to infinity, the
change is affine. We may therefore count all the other totally
ramified fibres using $z$.

Substituting (RTG2) into (RTG1) and cancelling $x^3$ gives
\[
 z^2x^3+u x^2+(e^2+t z)x+z+a+te=0.                    \tag{RTG3}
\]
The value $z=0$ is not totally ramified, since its fibre contains
the unramified point $W$. If a nonzero finite value $z_0$ is totally
ramified, the cubic in $x$ in (RTG3) must be
$z_0^2(x-r)^3$. Its quadratic coefficient forces $r=u/z_0^2$.
The remaining coefficients give
\[
 \begin{aligned}
 P(z_0)&=0,&P(z)&=t z^3+e^2z^2+u^2,\\
 Q_5(z_0)&=0,&Q_5(z)&=z^5+(a+te)z^4+u^3.
 \end{aligned}                                       \tag{RTG4}
\]
If $t=0$, then $P(z)=(ez+u)^2$ has at most one root. There
cannot be two further totally ramified values, proving the assertion
in this case.

Suppose $t\ne0$. Put
\[
 \begin{aligned}
 D&=e^4+ua+ute,\qquad A=te^2+u,\\
 R(z)&=t^2z^2+uz+D,\\
 R_1&=t^2D+uA,\qquad R_0=u^2t^3+AD,\qquad R_2=A^2+ut^5.
 \end{aligned}                                       \tag{RTG5}
\]
The identity $P(z)^2+uQ_5(z)=z^4R(z)$ shows that every
nonzero common root in (RTG4) is a root of $R$.
Two further totally ramified points would give two distinct such
values. The quadratic $R$ is separable, since its derivative is
$u\ne0$, so it would divide $P$. The exact polynomial division is
\[
              t^3P(z)+(t^2z+A)R(z)=R_1z+R_0.
                                                               \tag{RTG6}
\]
Consequently $R_1=R_0=0$. The identity
$t^2R_0+AR_1=uR_2$ then gives $R_2=0$.

These conditions contradict smoothness of (RTG1). Indeed define
\[
                  S=u^2+a^2t^4+ut^5+at^7+et^8.
\]
There is the universal polynomial identity over $\mathbf F_2$
\[
                  u^2S=R_1^2+ut^5R_1+e^4t^2R_2.       \tag{RTG7}
\]
Thus $S=0$. If $f(x)$ denotes the right side of (RTG1), direct
substitution gives
\[
          S=t^8\bigl(f'(1/t)^2+t^2 f(1/t)\bigr).       \tag{RTG8}
\]
At $x=1/t$ the coefficient of $y$ is zero. The unique point with
$y^2=f(1/t)$ is singular precisely when
$ty+f'(1/t)=0$, which is equivalent to the vanishing in (RTG8).
This is a singular point of the affine model, contradicting its
smoothness. The transformations to (RTG1) were invertible, so this
is a contradiction to the original hypothesis.

The exact program
\path{code/nonordinary_trigonal_direct_certificate.py}
checks the elimination, the two triple-root conditions, both
remainder identities and the singularity identity as polynomial
identities over $\mathbf F_2$, including the separate $t=0$ case.
It does not replace the Riemann--Roch or ramification arguments.
\end{proof}

\begin{corollary}[No theta triples under a group killed by three]
\label{cor:root-nonordinary-theta-no-three-triples-candidate}
\status{Proved and independently checked.}
Let $(B,\Theta)$ be a nonproduct nonordinary principally polarized
abelian surface over $k$ as above, and let $K\subset B[3](k)$ be
any subgroup. No three distinct points of its smooth principal
theta curve have the same image in $B/K$. No isotropy assumption
on $K$ is needed for this assertion.
\end{corollary}
\begin{proof}
Identify $B$ with the Jacobian of its genus-two theta curve $C$,
using the already proved principal-surface dichotomy. Three such
points $P,Q,R$ would satisfy
$3P\sim3Q\sim3R$ as divisors on $C$. A rational function with
divisor $3P-3Q$ gives a degree-three map totally ramified at
$P,Q$. Since $h^0(3Q)=2$, the function with divisor $3R-3Q$
is an affine linear combination of that function and the constant
function. The same map is therefore totally ramified at $R$.
This contradicts the lemma.
\end{proof}

\begin{remark}[Scope and independent route]
\label{rem:root-trigonal-direct-scope-candidate}
\status{Proved scope statement.}
This is a direct argument on the actual nonordinary characteristic-two
normal forms. It neither imports a generic complex nodality theorem
nor asserts a condition for higher odd kernel exponent. A separate
Galois-closure and Cartier argument has also been proved independently;
the direct proof above has no dependence on that argument or on its
ramification classification. The existence of full node permutation
groups and gluing translations remains a separate monodromy question.
\end{remark}


% Root proof promoted after full monodromy and obstruction/crystal independent audits.
\subsection{A nonproduct three-neighbour and its clean theta boundary}

\begin{lemma}[A nonproduct principally polarized three-neighbour]
\label{lem:root-nonordinary-nonproduct-three-neighbor-candidate}
\status{Proved; both complete independent mathematical audits passed.}
Let $(A,\Theta_A)$ be a nonproduct nonordinary principally polarized
abelian surface over an algebraically closed field of characteristic two.
There is a nonproduct principally polarized surface $(B,\Theta_B)$
and a polarized isogeny $\phi:B\to A$ such that
\[
 \deg\phi=9,\qquad \ker\phi\subset B[3],\qquad
 \phi^*\lambda_{\Theta_A}=3\lambda_{\Theta_B}.
\]
The surface $B$ has the same two-rank as $A$. In two-rank zero,
every quotient of $A$ by a maximal isotropic subgroup of $A[3]$
is nonproduct. In two-rank one, at most one such subgroup has a
product principally polarized quotient.
\end{lemma}
\begin{proof}
Use the principal alternating Weil pairing on the etale group
$A[3]\simeq\mathbf F_3^4$. A maximal isotropic subgroup $H$ has
order nine. The polarized quotient construction already proved in
the odd-isogeny argument gives
\[
 \lambda:A\longrightarrow B=A/H,\qquad
 \lambda^*\lambda_{\Theta_B}=3\lambda_{\Theta_A}.
\]
Its polarized dual $\phi:B\to A$ satisfies
$\phi\lambda=[3]_A$, $\lambda\phi=[3]_B$ and the displayed
pullback formula. Both maps are etale. In particular they induce
isomorphisms on the two-divisible groups and preserve two-rank.

Suppose first that this rank is zero. On a nonproduct principal
surface the smooth theta curve has genus two. Its previously
proved characteristic-two normal form has equation $y^2+y=f_5(x)$;
the regular forms $dx,x\,dx$ have Cartier images $0,dx$.
Thus Cartier on its two-dimensional differential space has rank one.
The restriction of invariant differentials from its Jacobian is
an isomorphism. A rank-zero product principal surface is a product
of two supersingular elliptic curves, and its Cartier operator is
zero. Pullback by the etale isogeny $\lambda$ is an isomorphism on
invariant differentials and commutes with Cartier. Its target cannot
therefore be such a product.

Now suppose the rank is one and a quotient $B$ is a product
$E_o'\times E_s'$ with its product principal polarization. Exactly
one elliptic factor is ordinary and the other supersingular.
Let $E_o,E_s\subset A$ be the connected inverse-image elliptic
subvarieties of these axes. The restrictions of $\lambda$ are
isogenies onto the axes, so their respective types are ordinary
and supersingular. There is at most one elliptic subvariety of
each type in $A$: two distinct elliptic subvarieties span $A$ and
their addition map is an isogeny; two ordinary ones would give
two-rank two and two supersingular ones two-rank zero.

For either $E=E_o,E_s$, the projection formula for the restriction
of the polarized quotient gives
\[
             \deg(\lambda|_E)=3(\Theta_A\cdot E).
\]
Its kernel is a subgroup of $E[3]$, so this degree is one, three,
or nine. The intersection $\Theta_A\cdot E$ is a positive integer.
It cannot be one on a nonproduct principal surface. Indeed, in that
case the restricted polarization $\lambda_E:E\to E^\vee$ is an
isomorphism, and for the inclusion $i:E\hookrightarrow A$ the map
\[
 r=\lambda_E^{-1}\circ i^\vee\circ\lambda_{\Theta_A}:A\longrightarrow E
\]
satisfies $ri=1_E$. The maps $(e,z)\mapsto i(e)+z$ and
$a\mapsto(r(a),a-i(r(a)))$ are inverse group-scheme isomorphisms
$E\times\ker r\simeq A$. Hence $\ker r$ is a connected smooth
elliptic curve. The kernel is orthogonal to $E$ for the principal
polarization, which therefore becomes the product of the two
principal polarizations. This contradicts the nonproduct hypothesis.

It follows that $\Theta_A\cdot E=3$ and
$\deg(\lambda|_E)=9$. Thus $H=E[3]$ for either axis.
The uniqueness of the ordinary and supersingular elliptic
subvarieties shows that this possible bad subgroup $H$ is unique.
There are at least two maximal isotropic subgroups: a symplectic
basis supplies the spans of its two $e$ vectors and its two $f$
vectors. Choose one different from the possible bad subgroup.
Its quotient and polarized dual give the required nonproduct $B$.
\end{proof}

\begin{lemma}[A clean common three-theta boundary]
\label{lem:root-nonordinary-three-theta-clean-boundary-candidate}
\status{Proved; both complete independent mathematical audits passed.}
Let $(B,\Theta)$ be a nonproduct nonordinary principally polarized
abelian surface in characteristic two, and let $K\subset B[3]$
have order nine. Put $q:B\to B/K$. There are theta translates
$C_i=\Theta+a_i$, $1\le i\le3$, with $a_1+a_2+a_3=0$, such that:
each $q(C_i)$ is integral with eight ordinary nodes; the three
images meet pairwise transversely; and all self-nodes and mutual
intersections are distinct. For every $r_i\in K$ with
$r_1+r_2+r_3=0$, the source union
\[
                       C_r=\bigcup_{i=1}^3(C_i+r_i)
\]
belongs to the same complete system $|3\Theta|$, has six ordinary
intersection nodes, and has the same target image $Y$.
All these finitely many source unions lie in a connected coefficient
open in $|3\Theta|$ on which the additional identifications under
$q$ form finite etale node schemes, with eighteen nodes per
unordered nonzero type. Thus the first geometric hypothesis of
Theorem~\ref{thm:three-theta-conditional-full-image-candidate}
holds for this kernel.
\end{lemma}
\begin{proof}
Individual theta norms have the asserted eight nodes by
Corollary~\ref{cor:nonordinary-theta-three-torsion-nodal-norm-candidate}.
This also excludes triple identifications within one theta curve.
We choose the translates in the irreducible four-dimensional
parameter space $B^2$, with $a_3=-a_1-a_2$.

The locus of differences giving a tangency between two distinct
theta translates is contained in the image of the one-dimensional
theta curve under multiplication by two, together with the zero
difference. To see this, identify the theta curve with its
Abel--Jacobi image based at a Weierstrass point. Equality of tangent
lines is equality of canonical images, so the two points coincide
or are hyperelliptic conjugates. Their differences are respectively
zero or $2(P-W)$. Each pair-difference map from $B^2$ to $B$ is
surjective, including pairs involving $a_3$; one of its two
coefficients is $1$ or $-1$. Avoiding the finitely many translates
of this proper locus by $K$ makes all cross intersections transverse
and keeps distinct components distinct.

For three distinct components, a triple intersection after any
specified kernel translations is parametrized by the inverse image
of $\Theta^3$ under the affine translate of the homomorphism
\[
 B^3\longrightarrow B^3,\qquad
 (P,a_1,a_2)\longmapsto(P-a_1,P-a_2,P+a_1+a_2).
\]
Its integral coefficient matrix has determinant three, so it is
a finite etale isogeny in characteristic two. This incidence has
dimension three, and its image in the four-dimensional parameter
space is proper, since the projection is proper. There are only
finitely many kernel choices. For a triple involving two branches
from one component and one from another, the first two branches
come from one of finitely many fixed self-identifications on the
theta curve. Requiring a translate of another theta curve to contain
that point is again a proper condition: its pair difference lies
on a fixed translate of a curve. These exclusions also keep every
self-node off the other components and prevent coincidences among
the different mutual intersections. A finite union of proper closed
subsets cannot cover $B^2$. Choose a point in the remaining open.

The theorem of the square and $\sum a_i=\sum r_i=0$ give the actual
line-bundle identity $\mathcal O_B(C_r)\simeq\mathcal O_B(3\Theta)$.
The selected three pairs of theta translates each meet transversely
in two points, and no triples remain, so $C_r$ has exactly six nodes.
Translations in $K$ do not change its image under $q$.

Here is the required family statement near these reducible sources.
For fixed $g\in K\setminus\{0\}$, the two equations defining a
point of $C_s\cap(C_s+g)$ have independent differentials at every
point of each chosen $C_r$: both source branches are smooth there,
and their translated tangent lines are transverse. A source
intersection node cannot occur in such a pair, since it would give
at least three target branches, already excluded. The incidence
over the coefficient space is therefore etale at these points;
properness and quasi-finiteness make it finite etale after shrinking
to an open containing the finitely many chosen sources. Each fibre
has eighteen points, since $(3\Theta)^2=18$.

Impose the same open conditions for all nonzero $g$, exclude triple
identifications and component stabilizers, and identify $g$ with
$-g$. A node of unordered type $\{g,-g\}$ is represented exactly
once in the incidence for one chosen representative $g$; exponent
three and absence of triples ensure this assertion. The resulting
node scheme for each type has degree eighteen. The intersection of
these opens contains every chosen $C_r$ and is a nonempty open of
the irreducible coefficient space, hence connected. The usual
partial-normalization exact sequence applies there: the finite map
$C_s\to q(C_s)$ is a local isomorphism away from these pairs, and
at each pair it identifies two smooth transverse branches. This
gives the finite etale graph subgroups even when $C_s$ itself is
nodal. It proves precisely the stated geometric hypothesis.
\end{proof}

\begin{remark}[The remaining monodromy input]
\label{rem:root-three-neighbor-boundary-scope-candidate}
\status{Proved scope statement.}
These two lemmas supply a nonproduct source and a clean common
boundary for every fixed nonproduct nonordinary principal target
in degree three. They do not supply the product node-permutation
group or the gluing translations. Those are the second hypothesis
of the conditional graph theorem and require a separate proof.
No existence statement for clean theta norms at higher odd kernel
exponent is inferred.
\end{remark}


\begin{theorem}[Degree three on every fixed nonordinary principal surface]
\label{thm:root-every-nonordinary-ppav-degree-three-full-candidate}
\status{Proved; complete consequence and independent dependency audits passed.}
Let $(A,\Theta)$ be a principally polarized abelian surface over a
perfect field of characteristic two, and assume its geometric
two-rank is zero or one. On the smooth maximal-rank open of the
complete linear system $|9\Theta|$, the relative unit local system
has rank eighty and both its geometric and arithmetic monodromy
groups are
\[
                            \GL_{80}(\mathbf F_2).
\]
Equivalently this is degree three for the embedding by $3\Theta$,
whenever the ambient hypersurface family surjects onto this complete
system, with the same restriction qualification as before.
\end{theorem}
\begin{proof}
We first work over an algebraic closure. If the principal surface
is a product, the previously proved complete supersingular-product
and mixed-product odd-degree theorems give exactly this group at
$d=3$. If it is nonproduct,
Lemma~\ref{lem:root-nonordinary-nonproduct-three-neighbor-candidate}
supplies a nonproduct nonordinary principal surface $B$ with a
polarized etale isogeny $B\to A$ of multiplier three. Apply
Theorem~\ref{thm:nonordinary-degree-three-from-neighbour-candidate}.
Its clean-boundary input is also checked independently in
Lemma~\ref{lem:root-nonordinary-three-theta-clean-boundary-candidate}.
The known all-degree principal-surface rank formula gives
$9\cdot3^2-1=80$, and the theorem supplies the full geometric
group in that dimension. Thus the conclusion holds for every
geometric principal surface, with no generic-modulus assumption.

The arithmetic image over the original perfect field contains
the geometric image in the same general linear group. Since that
geometric image is already the whole group, equality follows.
The chosen neighbour need only exist geometrically; no descent
of that auxiliary isogeny is required. If restriction from ambient
hypersurfaces to the complete system is surjective, its parameter
fibres are geometrically integral affine spaces on coefficient
charts, so the established restriction-transfer argument gives
the same monodromy there.
\end{proof}


\subsection{An induction through powers of three for every nonordinary principal surface}

\begin{lemma}[Uniform group inputs for the norm step]
\label{lem:RD-A4-three-power-groups}
\status{Proved; full root and monodromy proof/code audits passed.}
Let $r\ge8$, put $s=2r+2$, and let $t\ge1$.
The simple group $\GL_r(\mathbf F_2)$ is not a composition factor
of $S_s^t$.  Moreover
\[
 H^1(\GL_r(\mathbf F_2),\mathbf F_2^r)
 =H^1(\GL_r(\mathbf F_2),(\mathbf F_2^r)^*)=0.
\]
Consequently a subgroup of $S_s^t\times\GL_r(\mathbf F_2)$
surjecting onto both factors is the whole product.
\end{lemma}
\begin{proof}
The only simple composition factors of $S_s^t$ are $C_2$ and
$A_s$.  The general linear group is simple because over $\mathbf F_2$
it is $\mathrm{PSL}_r$ and $r\ge3$.  It cannot be $A_s$, since
\[
 v_2|\GL_r(\mathbf F_2)|=r(r-1)/2>2r
       \ge v_2(s!)-1=v_2|A_s|.
\]
Here $v_2(s!)=s-s_2(s)$, with $s_2(s)\ge1$ the binary digit sum.
Goursat's lemma now proves the last assertion.

For completeness, the cohomology proof is uniform in this range.
For a cocycle $z$ on the natural module, put $t_{ij}=1+E_{ij}$.
The relation $t_{ij}^2=1$ kills the $j$th coordinate of $z(t_{ij})$.
For $l\ne i,j$, choose $k$ outside $\{i,j,l\}$; the commuting
relation for $t_{ij},t_{kl}$ then kills its $l$th coordinate.
Thus $z(t_{ij})=a_{ij}e_i$.  Applying the cocycle to
$[t_{ij},t_{jk}]=t_{ik}$ gives $a_{ik}=a_{jk}$, so $a_{ij}=b_j$
for a single vector $b$.  Hence $z(t_{ij})=(t_{ij}-1)b$.
These transvections generate the group.  The cocycle is a
coboundary, and inverse transpose gives the dual assertion.
This is the full argument of Lemma~\ref{lem:char2-group-inputs},
with its dimension requirement verified rather than a fixed-prime
divisibility argument.
\end{proof}

\begin{lemma}[The many-theta graph]
\label{lem:RD-A4-many-theta-graph}
\status{Proved; full root and monodromy proof/code audits passed.}
Let $m\ge3$ and let $K$ be a finite abelian group of odd order $k>1$.
Put $\mathcal H=(K\setminus0)/\pm$.  Form the graph with $m$ vertices,
two loops at every vertex for every $h\in\mathcal H$, and two cross
edges $e_{ij,g,\epsilon}$ for every $i<j$, $g\in K$,
$\epsilon\in\{0,1\}$.  Let $\mathcal G$ be its $\mathbf F_2$
cochain space modulo vertex coboundaries.  For
$r=(r_1,\ldots,r_m)\in K^m$ with $\sum_i r_i=0$, retain the two
edges $g=r_j-r_i$ for every pair, and let $T_r$ be the kernel of
restriction to that selected doubled complete graph.  Then
\[
 \begin{gathered}
 \dim\mathcal G=m^2k-2m+1,\qquad
 \dim H^1(\text{selected graph})=(m-1)^2,\\
 \dim T_r=m^2(k-1),\qquad
 T_r=\bigoplus_{h\in\mathcal H}T_{r,h},\quad
 \dim T_{r,h}=2m^2.
 \end{gathered}                                         \tag{RD-A4-3P.1}
\]
The type coordinates consist of its $2m$ loops and the four
cross edges $g=r_j-r_i\pm h$ for each pair.  Write $T_{r,h}^0$
for their sum-zero hyperplane.
The span of all these hearts, as $r,h$ vary, is the codimension-one
total-edge-sum-zero subspace of $\mathcal G$.  If a subspace
$S\subset\bigcap_rT_r$ is invariant under
$\prod_hS_{2m^2}$ on $T_0$, then $S=0$.
\end{lemma}
\begin{proof}
There are $m(k-1)+m(m-1)k=m^2k-m$ edges and $m-1$ independent
vertex coboundaries.  The selected graph is connected with
$m(m-1)$ edges.  Its unselected-edge cochains give independent
coordinates on $T_r$: a coboundary zero on the selected graph
is zero.  This proves all dimensions and the type decomposition.

Pair cochains with cycles.  A cycle annihilating a type heart is
constant on all its $2m^2$ coordinates.  The fixed loops of type
$h$ make that constant $c_h$ independent of $r$.  Every difference
$r_j-r_i$ ranges over $K$ under the sum-zero constraint: set those
two entries to $0,g$, a third entry to $-g$, and all others to zero.
It follows that every cross-edge coefficient equals every $c_h$.
Thus the annihilator consists precisely of the all-ones cycle.
It is indeed a cycle, since each cross degree is $2k(m-1)$.
This proves the codimension-one assertion.

For the last assertion, a nonconstant type component of a vector
in $S$ yields, by a transposition and then conjugation, every pair
vector in that type.  Choose one supported on a cross edge and a
loop, and choose $r$ selecting that cross edge.  Its restriction
has exactly one of a parallel pair nonzero, so cannot be a
coboundary.  This contradicts $S\subset T_r$.
Hence every vector of $S$ is type-constant.  Extend its constants
to an even function $a:K\to\mathbf F_2$ with $a(0)=0$.
Take $r=(-u,0,u,0,\ldots,0)$.  The triangle on its first three
vertices gives the necessary coboundary relation
$a(u)+a(u)+a(2u)=0$.  Therefore $a(2u)=0$ for every $u$.
Multiplication by two is bijective on $K$, so the vector is zero.
\end{proof}

\begin{proposition}[Node products and gluing hearts for the uniform three-step]
\label{prop:RD-A4-three-step-boundary}
\status{Proved; full root and monodromy proof/code audits passed.}
Let $\phi:(B,\Theta_B)\to(A,\Theta_A)$ be a polarized etale isogeny
of multiplier three and degree nine over an algebraically closed
field of characteristic two.  Suppose $B$ is nonordinary and
indecomposable.  Let $m=3e$ be odd, and suppose the degree-$e$
relative system on $|m\Theta_B|$ has rank $r=m^2-1$ and full
geometric image $\GL_r(\mathbf F_2)$.
For its generic pure norm family put
\[
 K=\ker\phi,\quad \delta=8m^2,\quad s=2m^2,
 \qquad 0\to T\to V\to U\to0.
\]
Then $\dim T=\delta$, $\dim U=r$, and $\dim V=R=9m^2-1$.
The node permutation image is $\prod_{h\in(K\setminus0)/\pm}S_s$.
The boundary image contains an actual natural $\GL_r$ block fixing
a complement, and every translation in $\Hom(U,T_h^0)$.
All kernel-shifted unions of $m$ general theta translates with
sum zero belong to one connected source coefficient open and
give the common graph of Lemma~\ref{lem:RD-A4-many-theta-graph}.
\end{proposition}
\begin{proof}
The norm line bundle has numerical class $3m\Theta_A$.
Its fixed $\Pic^0$ discrepancy is removed by a single fixed target
translation, since the polarization map of $3m\Theta_A$ is
surjective.  This identifies the entire norm family with a subfamily
of the actual complete system $|3m\Theta_A|$, rather than only
its numerical class.

We verify the geometric changes from the proved degree-three
node and gluing argument, rather than assuming its dimensions
are unchanged.  For $0\ne g\in K$, descend $3\Theta_B$ to
$E=B/\langle g\rangle$ as the three globally generated ample
bundles $M_i$ of square six in
Lemma~\ref{lem:nonproduct-three-theta-finite-eigenmaps-candidate}.
The three descents of $m\Theta_B$ are
$M_0^{e-1}\otimes M_i$.  Tensor products of globally generated
bundles are globally generated, and these bundles are ample.
They therefore define finite eigenmaps, now to
$\mathbf P^{3e^2-1}$; their section spaces have dimension $3e^2$.

The value chord $\langle v(P),v(P+g)\rangle$ contains no vertex
of the three-dimensional orbit plane.  Meeting two such chords
therefore requires coincidence of at least two finite eigenmap
projections.  The exceptional pair locus has dimension at most
two.  Four independent values give an incidence of dimension
$4+(m^2-5)=m^2-1$.  On the exceptional locus the rank is at
least three, so its dimension is at most
$2+(m^2-4)=m^2-2$.  The same estimate handles overlapping
orbits.  The three-value assertion follows from
Lemma~\ref{lem:ppav-three-distinct-values} by multiplying by a
section of $(m-3)\Theta_B$ nonzero at the three specified points;
this bundle is globally generated when $m>3$, and is trivial
when $m=3$.  Hence the ordered-two-point node cover is irreducible.

The generic norm is nodal.  Three-point identifications are
excluded by these three independent values and their two-dimensional
point incidence.  For a fixed nonzero translate, first jets at
$P,P+g$ have rank six except at finitely many $P$ for $3\Theta_B$:
the exceptional equations are $3P+g,3P+2g\in\Theta_B$, and the
two distinct theta translates intersect finitely.  Multiplication
by a section of $(m-3)\Theta_B$ nonzero at the two points preserves
that rank.  On this open, a tangency is one extra equation after
the two value equations; the resulting incidence has codimension
at least one in the coefficient space.  The finitely many exceptional
points cannot occur generically, already by their value equations.
Bertini gives smooth generic source members.  Thus each type has
$s=(m\Theta_B)^2=2m^2$ nodes, and there are four types.
The normalization and its $\delta$ graph directions give full
target relative rank $R$, so the relative Lang system extends
over this nodal norm family.

For the common boundary, choose $m$ theta translates with sum
zero.  Each theta norm has its eight distinct ordinary nodes by
Corollary~\ref{cor:nonordinary-theta-three-torsion-nodal-norm-candidate}.
Choose the translating points generally so all component norms
meet transversely, avoid the self-nodes and have no triple points.
For $m=3$ this is the complete incidence argument in
Lemma~\ref{lem:root-nonordinary-three-theta-clean-boundary-candidate}.
For $m\ge5$, any two or three selected translating points vary
independently: an additional component absorbs their sum.
The three-membership incidence therefore has dimension one less
than the parameter space, and the pair difference cover is
generically etale.  The same proper exclusions apply to the
finitely many intrinsic self-nodes.  Impose these open conditions
also for every one of the finitely many $r\in K^m$ with sum zero.
Their intersection is nonempty in the irreducible translation
parameter space.  All source unions then lie in $|m\Theta_B|$,
and their common target has exactly the graph in the preceding lemma.

The shifted difference covers supply isolated transpositions.
For a fixed pair their branch supports are the distinct shifts
of $[2]\Theta_B$.  For distinct pairs the two difference variables
are independent: when $m=3$ the relevant matrix has determinant
three; when $m\ge5$, an unused translating point absorbs the sum
and the two difference maps have a unit rank-two minor.
Thus generic inertia on any one branch divisor exchanges its
two points and fixes every other node cover.  This is precisely
the characteristic-two difference-cover proof in
Lemma~\ref{lem:theta-difference-cover-char2-candidate}; its possible
purely inseparable residue extension causes no tame-inertia assumption.
The normal coefficient-open and dense smooth-source argument in
Proposition~\ref{prop:nonordinary-three-torsion-node-product-candidate}
places these permutations in the generic norm image.
Two-transitivity plus an isolated transposition gives each $S_s$,
and normal conjugation gives the full product supported factorwise.

By Lemma~\ref{lem:RD-A4-three-power-groups}, the joint image on
$T,U$ is $N\times\GL_r$, where $N=S_s^4$.
Its translation kernel is $U^*\otimes T'$ for an $N$-stable
$T'\subset T$, by the full matrix algebra generated on $U^*$.
Modulo $T'$, the uniform $H^1$ vanishing splits the $\GL_r$ factor.
The commuting $N$ factor has no remaining column, since
$(U^*)^{\GL_r}=0$.  Thus
\[
                         V/T'=(T/T')\oplus U.            \tag{RD-A4-3P.2}
\]
This also supplies an actual natural block before quotienting:
all residual columns lie in the known translation kernel and can
be removed by its actual elements.

The splitting extends to the smaller-rank common boundary exactly
as a section of the finite etale $T/T'$ torsor on the full source
Lang scheme.  To specify every hypothesis, take the normal source
coefficient open where the additional identifications are distinct
transverse pairs of smooth source points, allowing internal source
nodes and selected-pair tangencies.  The partial-normalization
sequence gives an exact sequence of coherent Frobenius bundles
with graph-unit kernel.  The source Lang scheme is etale over this
normal base and hence normal; all its components dominate the base.
The generic section of the finite etale torsor extends uniquely
over each component, by finite birational closure and normality.
Addition and descent compatibility extend by density.  These are
the full steps of
Proposition~\ref{prop:nonordinary-three-torsion-gluing-hearts-candidate},
and do not identify a rank-dropping Lang fibre with a generic fibre.

For each selected unshifted component pair, local inertia has
rank-one image $q_{ij}$, the sum of its two edge coordinates.
This is the same strict-henselian gluing calculation: component
torsors are constant after translation, their local evaluations
can be trivialized simultaneously at the coalescing pair, and
inertia exchanges only those two labels.  It fixes every additional
node.  In the extended splitting, $q_{ij}$ in the target must be
the unique lift of its class in the source graph.  For each vertex,
the latter classes satisfy $\sum_{j\ne i}q_{ij}=0$.
Consequently the target sum lies in $T'$.  Modulo the full vertex
star it is the unselected-edge star at that vertex.  In each type
it has
\[
                  4(m-1)\ \text{ones and}\quad
                  2m^2-4(m-1)\ \text{zeros}.            \tag{RD-A4-3P.3}
\]
Both numbers are positive.  A transposition in any one factor
therefore produces a pair vector of that type in $T'$; conjugating
produces its whole heart.  Hence $T_h^0\subset T'$ for every type,
which gives all the asserted gluing translations.
\end{proof}

\begin{theorem}[The uniform multiplier-three monodromy step]
\label{thm:RD-A4-three-step}
\status{Proved; full root and monodromy proof/code audits passed.}
Under the hypotheses of Proposition~\ref{prop:RD-A4-three-step-boundary},
the geometric and arithmetic relative images on the smooth
maximal-rank part of $|3m\Theta_A|$ are
$\GL_{9m^2-1}(\mathbf F_2)$.
\end{theorem}
\begin{proof}
Work geometrically.  Let $S$ be a nonzero proper invariant
subspace of $V$.  Its projection to the simple source module $U$
on the connected dense maximal-rank source open is zero or all
of $U$.  Use Lemma~\ref{lem:normal-base-clopen-Lang-closure-candidate}
to extend this alternative to every common theta specialization.
The full graph injects into relative cohomology because ambient
$H^1(B,\mathbf F_2)$ injects into each theta normalization and
therefore misses the graph.

In the first case, the closure has the full original dimension,
lies in every $T_r$, and is invariant under the node product on
$T_0$.  Lemma~\ref{lem:RD-A4-many-theta-graph} makes it zero, a
contradiction.  In the second case, the gluing translations put
every type heart in $S$; their common-boundary span is the
total-edge-sum-zero graph hyperplane.  Its intersection with
each $T_r$ has dimension $\delta-1$.  The intersection with the
clopen Lang closure is finite etale and has constant rank.
It follows that $\dim S\ge\delta-1+r=R-1$.
A proper $S$ is a hyperplane, giving a fixed quotient covector
over $\mathbf F_2$.  This contradicts the proved odd-degree
nonordinary covector theorem
\ref{thm:nonordinary-ppav-all-covectors}.  The full image is irreducible.

Let $J$ be the normal subgroup generated by transvections.
Its nonzero socle is invariant under the whole image, hence is
all of $V$; the $J$-module is semisimple.  The actual natural
$\GL_r$ block is generated by transvections and occurs once.
It therefore lies in one simple $J$-summand.  The rank-one
gluing maps put all four hearts in that same summand.  Its
dimension is at least
\[
                         r+4(s-1)=R-4>R/2.
\]
Two distinct conjugate simple summands would be disjoint, which
this inequality forbids.  The summand is invariant under the
whole group and equals $V$.  Thus $J$ is irreducible.
The already audited transvection classification
\cite[Theorem~3.6, printed p.~22]{KantorSeress2001} now leaves
the full general linear group or a group preserving a
nondegenerate alternating form.  The actual natural $\GL_r$
block fixing a complement has neither an invariant bilinear
form nor an invariant functional; it would lie in the radical
of such a form.  The latter alternative is impossible.
The geometric image is full, and the arithmetic image contains
it in that same general linear group, so is full as well.
\end{proof}

\begin{corollary}[Every power-of-three degree on every fixed nonordinary surface]
\label{cor:RD-A4-powers-three}
\status{Proved; full root and monodromy proof/code audits passed.}
Let $(A,\Theta)$ be any principally polarized abelian surface over
a perfect field of characteristic two, of geometric two-rank zero
or one.  For every integer $j\ge0$, put $d=3^j$.
On the smooth maximal-rank open of the complete system
$|3d\Theta|$, both geometric and arithmetic relative images are
\[
                         \GL_{9d^2-1}(\mathbf F_2).
\]
The same holds for ambient hypersurface families whenever the
restriction map onto this complete system is surjective, with
the previously stated embedding qualification.
\end{corollary}
\begin{proof}
Degree one is the proved fixed-modulus nonordinary principal
theorem.  Suppose the assertion holds in degree $e=3^{j-1}$
on every such surface.  Products already have the full image in
all odd degrees by the complete supersingular-product and
mixed-product theorems.  For a nonproduct $A$, the proved
Lemma~\ref{lem:root-nonordinary-nonproduct-three-neighbor-candidate}
supplies an indecomposable nonordinary principal $B$ with an
etale polarized isogeny of multiplier three to $A$.
Apply the induction hypothesis on $B$ to $|3e\Theta_B|$, and
then Theorem~\ref{thm:RD-A4-three-step} with $m=3e$.
Its target is $|9e\Theta_A|=|3d\Theta_A|$ with $d=3e$,
and its rank is $9d^2-1$.  The construction is geometric;
the auxiliary neighbour need not descend to the original field.
Arithmetic equality and surjective-restriction transfer follow
as in the degree-three theorem.
\end{proof}

\begin{remark}[Scope of the induction]
\label{rem:RD-A4-powers-three-scope}
\status{Proved; full root and monodromy proof/code audits passed.}
The induction uses only kernels killed by three.  It proves an
infinite sequence of degrees for every fixed nonordinary principal
surface, and does not assert the other odd degrees or any even
degree on a nonproduct surface.  The step from $e$ to $3e$ is
uniform in $e$, including its group-theoretic and rank-drop
arguments; it does not reuse the prime $31$ after the source
rank has grown.

The retained exact supplement
\path{code/reduction_descent/verify_nonordinary_powers_three_graphs.py}
checks the graph dimensions, type partitions, vertex-star relations
and type-constant constraints for $m=3,5,9,27$.  For $m=3,5,9$
it also computes the heart span from an explicitly retained finite
set of sum-zero shifts, obtaining dimensions $75,215,711$; it
does not enumerate all $9^{m-1}$ shifts.  At $m=27$ that span
calculation is not performed.  Eight separate exact order and
dimension checks cover $m=3,5,9,15,27,81,243,729$.  The author run
passed in $0.1809211840009084$ seconds at 19:54:40 UTC.
These finite computations supplement the universal graph proof
and establish no geometric genericity.
\end{remark}


% Promoted after complete root and obstructions/crystals proof/source/code audits.
\subsection{Degree two on every fixed nonordinary principal surface}
\label{subsec:nonordinary-ppav-degree-two-wreath-theta-candidate}
\status{Proved and independently checked.}

Until the final theorem let $k$ be algebraically closed of characteristic
two, let $(A,\Theta)$ be an indecomposable principally polarized abelian
surface, and put $f=f(A)\in\{0,1\}$ and $L=\mathcal O_A(3\Theta)$.
We may take $\Theta$ symmetric. A different representative of the
principal polarization is a translate and gives an isomorphic family.
The known rank and covector theorems give, in degree two,
\[
 \dim V=33+f,\qquad \dim(V^*)^{\pi_1}=1+f,
 \qquad \dim K=32,
\]
where $K$ is the common kernel of all invariant covectors. We use the
already proved full degree-one group $\GL_8(\mathbf F_2)$ on every
fixed nonordinary principal surface and every Picard twist of $L$.
For a degree-one curve, ``sharp'' means that its two-rank is $8+f$.

\begin{lemma}[Lifted intersections on a fixed cubic-theta curve]
\label{lem:nonordinary-ppav-two-lifted-intersections-candidate}
\status{Proved under stated hypotheses}
\textup{[Proved; independent audits passed.]}
Fix $\alpha\in\Pic^0(A)$ and a smooth sharp curve
$C\in|L\otimes\alpha|$. If $\widetilde C\to C$ is a connected
finite etale torsor with finite elementary abelian two-group $G$,
then the monodromy of its fibres over the eighteen points $C\cap D$,
as $D$ varies through smooth sharp transverse members of
$|L\otimes\alpha^{-1}|$, is
\[
                 \ker(\Sigma:G^{18}\to G)\rtimes S_{18},
\]
up to choices of origins in the eighteen torsor fibres.
\end{lemma}
\begin{proof}
The image of $H^0(A,L\otimes\alpha^{-1})$ on $C$ is a very ample
linear system. Its projective dimension is eight, except that it
is seven when the restriction has its one-dimensional kernel
$H^0(A,\alpha^{-2})$. The original coefficient space maps onto
the image system with geometrically integral affine fibres off
the possible centre. It suffices to calculate the monodromy for
the image system on $C$.

Here are the characteristic-two tangency details. Every translate
of $L$ separates second jets on the surface, by
Theorem~\ref{thm:ppav-three-theta-jets}; consequently its restricted
system separates second jets on $C$ at every point. On $C\times C$
off the diagonal, the joint curve first-jet rank is four outside
a set of dimension at most one, is at least three outside a finite
set, and is at least two everywhere. Indeed the corresponding
surface ranks are six, at least five, and at least four. Passing
to curve jets discards only the two normal derivative coordinates.
The rank-six exceptional locus has finite fibres over either
fixed surface point by
Corollary~\ref{cor:all-theta-fixed-point-jets-candidate}, so its
intersection with $C^2$ has dimension at most one. The possible
surface rank-four locus is contained in a finite union of curves
\[
                 (P,Q)=(P,s-P),\qquad P+s\in\Theta,
                 \qquad 3s=0,
\]
after a translation for the chosen twist. Its intersection with
$C^2$ is finite: the first coordinate lies on a theta translate,
which cannot equal the smooth integral curve $C$ of numerical
class $3\Theta$.

Thus the incidence of two distinct tangencies has codimension at
least two in the coefficient space. Its three strata have point
dimensions at most $2,1,0$ and impose at least $4,3,2$ independent
linear conditions, respectively. Contact of order at least three
at one point also has codimension at least two, because three
curve jets are independent everywhere. The single-tangency
incidence is an irreducible projective bundle over $C$; a section
has finitely many zeros, so its map to coefficient space has finite
fibres. Its image is an irreducible divisor, whose geometric generic
member has exactly one double zero and sixteen other simple zeros.
In particular the unlabelled $S_{18}$ assertion is an instance of
the quasireflexive hyperplane theorem
\cite[Theorem~2 and Proposition~3.1]{Entin2021} with its hypothesis
verified here; we also give the labelled argument.

At the generic double zero the degree-two root cluster is
$z^2+bz+c=0$. The independent constant and linear jets give
independent coefficient parameters $b,c$. Along the reduced
discriminant $b=0$ the residual $c$ is not a square, including
after separable closure of the residue field. The equation
$y^2+y=c/b^2$ therefore has nontrivial inertia of order two.
This accounts for the possibly purely inseparable tangent-point
residue extension. The etale torsor is locally trivial on this
cluster, so inertia is a labelled transposition and fixes the
other sixteen torsor fibres pointwise.

The incidence of a point of $\widetilde C$ and a hyperplane
through its image is an irreducible projective bundle over
$\widetilde C$. Its map to coefficient space is finite, since a
nonzero member of the image system cannot vanish identically
on $C$. Hence lifted-point monodromy is transitive. It commutes
with the global deck transformations, and therefore lies in
$G^{18}\rtimes S_{18}$. The Galois closure is branched only over
the discriminant. Quotienting by the normal subgroup generated
by its divisorial inertia gives a cover unramified in codimension
one. Purity \cite{Stacks0BMB} and geometric simple connectedness
of projective space make that quotient trivial, exactly as in
Lemma~\ref{lem:mixed-lifted-intersection-wreath-candidate}.
Thus the image is generated by labelled transpositions.
Lemma~\ref{lem:labelled-transposition-wreath-candidate} gives
the displayed group. Finally the smooth sharp members form a
nonempty dense open by the rank theorem. Restricting to that open
does not change the connected components of any finite etale
cover of the normal coefficient open.
\end{proof}

\begin{lemma}[All boundary gluing columns]
\label{lem:nonordinary-ppav-two-boundary-gluing-candidate}
\status{Proved under stated hypotheses}
\textup{[Proved; independent audits passed.]}
For fixed $\alpha$, vary independent smooth sharp transverse
$C\in|L\otimes\alpha|$ and $D\in|L\otimes\alpha^{-1}|$.
Their union is in the maximal relative-rank open of $|L^2|$.
Write $A_0=H^1_{\et}(A,\mathbf F_2)$ and $H_C,H_D$ for the two
component cohomologies. The boundary representation has
\[
 \begin{gathered}
 0\longrightarrow T\longrightarrow V\longrightarrow R\longrightarrow0,
 \qquad T=\mathbf F_2^{18}/\mathbf F_2\mathbf1,\\
 R=(H_C\oplus H_D)/A_0,
 \qquad \dim R=16+f,\qquad
 T^0=\{(x_i):\sum x_i=0\}/\mathbf F_2\mathbf1.
 \end{gathered}
\]
It contains an actual $\GL_8\times\GL_8$ acting on two natural
eight-dimensional summands and fixing a complement, and all
translations $\operatorname{Hom}(R,T^0)$. Every nonzero invariant
subspace of this boundary representation contains $T^0$.
In its block frame the span of the two natural summands and $T^0$
is precisely the global common covector kernel $K$.
\end{lemma}
\begin{proof}
Each component has two-rank $8+f$, and the eighteen nodes add
seventeen unit graph directions. Subtracting the ambient rank
$f$ gives $2(8+f)+17-f=33+f$, the maximal relative rank. The
normalization sequence gives the displayed exact sequence.
Restriction from $A$ injects into either component cohomology;
on coherent cohomology this follows from
$H^1(A,(L\otimes\alpha)^{-1})=0$, and hence holds on unit parts.

Put $U_C=H_C/A_0$, $U_D=H_D/A_0$. Their independent relative
images are $\GL_8\times\GL_8$: both variables range over a dense
open of a product, and every Picard twist is a translate of $L$.
If $f=0$ the full component images are already $\GL_8$.
If $f=1$, a component image fixes $A_0$ and surjects onto $\GL_8$;
its unipotent kernel is zero or the entire natural-dual module.
The proved vanishing $H^1(\GL_8,\mathrm{nat}^*)=0$ splits its
complement. Thus in either case the component image is perfect
and has only $\GL_8$ as a nonabelian simple composition factor.
The image on $R$ is likewise a split extension of
$L_0=\GL_8\times\GL_8$ by a submodule of
$(U_C^*\oplus U_D^*)\otimes A_0$; it is perfect. The two
nontrivial simple modules here are nonisomorphic, and the same
first-cohomology vanishing proves the assertion about splitting.

The preceding lemma gives $S_{18}$ on $T$. This group and the
image on $R$ have no common nontrivial quotient: the latter is
perfect with only $\GL_8$ nonabelian simple factors, while
$S_{18}$ has factors $A_{18}$ and $C_2$. For example $31$ divides
$|\GL_8(\mathbf F_2)|$ but not $|S_{18}|$.
Goursat's lemma therefore gives the product joint image.
Take the inverse image of $1\times L_0$. Its translation kernel
lies in $\operatorname{Hom}(R,T)$, which under $L_0$ is a direct
sum of natural-dual and trivial modules. All its quotients have
zero first cohomology: use natural-dual vanishing and perfectness
for the trivial modules. The graph-of-a-cocycle argument gives
an actual block-diagonal $L_0$, fixing $T$ and an ambient
$f$-dimensional complement in $R$.

Let $J\subset\operatorname{Hom}(R,T)$ be the kernel of the action
on $T\oplus R$. Fix $C$ and its maximal connected elementary
two-cover, with group $G=H_C^*$ of dimension $8+f$. Its lifted
intersection image is $G^{17}\rtimes S_{18}$. Its only nonabelian
simple factor is $A_{18}$, so its joint image with the moving
$D$ cohomology is their product by Goursat and perfectness of
the latter. Hence a pure lifted-node translation $(g_i)$ with
$\sum g_i=0$ can be chosen with trivial moving-$D$ cohomology
action. It is an element of $J$, and on the injected $H_C\subset R$
its gluing column is
\[
                  \xi\longmapsto (\xi(g_i))_i
                                    \bmod\mathbf F_2\mathbf1.
\]
Indeed extend the fixed etale two-torsor represented by $\xi$
on $C$ using the trivial torsor on $D$ and chosen node gluing
values. Continuing the choices changes them by these node
evaluations; changing either component trivialization removes
the diagonal vector. This directly identifies the gluing action.

These columns exhaust $\operatorname{Hom}(H_C,T^0)$. Under $L_0$,
$\operatorname{Hom}(R,T)$ is semisimple, with isotypic pieces
for $U_C^*$, $U_D^*$ and, if $f=1$, the trivial ambient column.
Consequently their surjectivity on $H_C$ isolates all maps from
$U_C$ and $A_0$ to $T^0$, with the other columns zero. Reversing
$C,D$ supplies all maps from $U_D$. This proves
$\operatorname{Hom}(R,T^0)\subset J$.

If an invariant subspace has a nonzero projection to $R$, these
translations put all of $T^0$ in it. If it is contained in $T$,
choose a nonconstant representative of a nonzero vector. A
transposition exchanging unequal coordinates gives its difference
$e_i+e_j$, whose $S_{18}$-orbit spans $T^0$.
Finally every global invariant covector kills both natural
blocks and, by the translations, $T^0$. Their dimensions sum
to $8+8+16=32=\dim K$, proving the final assertion.
\end{proof}

\begin{lemma}[Two partitions of a six-theta curve]
\label{lem:nonordinary-ppav-two-six-theta-comparison-candidate}
\status{Proved under stated hypotheses}
\textup{[Proved; independent audits passed.]}
The geometric representation on $K$ is irreducible. Every
nonzero invariant subspace of $V$ contains $K$.
\end{lemma}
\begin{proof}
Choose general $a_1,\ldots,a_6\in A$ with $\sum a_i=0$.
Their theta translates are distinct, all pairs meet transversely
in two points, and there are no triple intersections. To see
nonemptiness under the sum constraint, pair differences remain
surjective parameters, and the theta difference map is generically
etale. For three chosen indices, vary their three translations
freely and compensate the sum using a fourth index. The incidence
of a common point has dimension at most five in their
six-dimensional translation space, so a generic triple avoids
it. There are only finitely many pairs and triples. The theorem
of the square gives a nodal curve
\[
                    Y=\bigcup_{i=1}^6(\Theta+a_i)\in|L^2|.
\]
Its graph has six vertices and two edges $e_{ij,0},e_{ij,1}$
for every unordered pair $i,j$. Thus it has thirty edges and
graph cohomology of dimension $30-6+1=25$. Its relative stable
rank is $25+6f-f=25+5f$, below the generic rank $33+f$.

For any partition into triples $I\sqcup J=\{1,\ldots,6\}$,
put $C_I=\bigcup_{i\in I}(\Theta+a_i)$ and similarly $D_J$.
They belong to $|L\otimes\alpha_I|$ and
$|L\otimes\alpha_I^{-1}|$. The two twists need not be trivial
and may differ between partitions. The two triple curves are
smooth at their eighteen cross intersection points, so these
points and their graph subgroup $T$ remain finite etale on a
neighbourhood in the complete product of the two coefficient
systems, including at this nodal pair.

Let $0\ne S\subset K$ be globally invariant. Extend it from
the smooth maximal-rank open to the entire coherent maximal-rank
open. This is legitimate because the full constant-rank Lang
scheme is finite etale there and a dense open of a normal
integral base has the same connected components of such a
cover. The full coherent relative Frobenius bundle is regular
on coefficient space: all these ample divisors are connected,
their arithmetic genus is constant, and the negative ample
cohomology vanishing injects ambient $H^1(\mathcal O_A)$.
Take the closure $\overline S$ in its full relative Lang scheme.
Lemma~\ref{lem:normal-base-clopen-Lang-closure-candidate} says
that it is an open-and-closed subgroup, including over $Y$.

On the generic maximal-rank locus of each boundary product,
the preceding lemma puts its entire cross-heart $T^0$ in $S$.
Since $T^0$ is finite etale and extends over the nodal pair,
the clopen closure lemma puts its special fibre in the same
space $\overline S_Y$, independently of the partition. This
uses only those extending finite Lang sections; there is no
identification of the full rank-dropping special fibre with $V$.
The graph space injects into relative cohomology: ambient
$H^1(\mathcal O_A)$ restricts isomorphically to the cohomology
of each principal theta normalization, whereas graph classes
restrict to zero on every normalization. The same holds on
unit parts.

Compare the partitions
\[
        \{1,2,3\}\mid\{4,5,6\},\qquad
        \{1,2,4\}\mid\{3,5,6\}.
\]
The second cross-heart contains
\[
             v_C=e_{13,0}+e_{23,0},\qquad
             v_D=e_{45,0}+e_{46,0}.
\]
Both vectors have even cross-edge sum. On the first partition,
$v_C$ has a nonzero projection to the internal graph of the
first triple and zero projection to the other; $v_D$ has the
opposite property. Nonvanishing can be read on a parallel-edge
two-cycle: the coordinate on $e_{13,0}$ differs from that on
$e_{13,1}$, and likewise for $e_{45,0},e_{45,1}$. In particular
neither internal projection is a vertex coboundary.

Both vectors belong to $\overline S_Y$. Pull back to the strict
henselization of the first boundary parameter space at this
pair. Each point of the etale scheme $\overline S$ has a unique
section with its specified closed value. Hence these two
sections belong to $S$ on the nonempty generic maximal-rank
part of that family. Restriction to each triple curve, followed
by the ambient coherent quotient, is a regular Frobenius map
of vector bundles. Its two nonzero graph projections at $Y$
remain nonzero generically. This proves that $S$ has nonzero
projection to both relative natural systems $U_C,U_D$ of the
first boundary. It does not assert surjectivity of a kernel
under arbitrary rank-dropping specialization.

The two actual independent $\GL_8$ blocks now force the image
of $S$ in $R$ to contain $U_C\oplus U_D$. Its kernel already
contains $T^0$. Thus $\dim S\ge16+16=32=\dim K$, and $S=K$.
For a general nonzero invariant subspace of $V$, the boundary
lemma first supplies $T^0\subset K$ in its intersection with
$K$. The irreducibility just proved then puts all of $K$ in it.
\end{proof}

\begin{theorem}[The full degree-two pointwise stabilizer]
\label{thm:every-nonordinary-ppav-degree-two-full-candidate}
\status{Proved under stated hypotheses}
\textup{[Proved; independent audits passed.]}
Let $(A,M)$ be a principally polarized abelian surface over any
perfect field of characteristic two, with principal polarization
represented by $M$ over that field and geometric two-rank
$f\in\{0,1\}$. Put $L=M^3$. On the smooth maximal-rank open of
$|L^2|$ the relative unit system has rank $33+f$ and
\[
 G_{\mathrm{geom}}=G_{\mathrm{arith}}
   =\operatorname{Hom}(\mathbf F_2^{\,1+f},\mathbf F_2^{32})
                      \rtimes\GL_{32}(\mathbf F_2).
\]
This is the full pointwise stabilizer of its entire invariant
covector space. Their common kernel is irreducible and is
contained in every nonzero invariant subspace.
\end{theorem}
\begin{proof}
First work over an algebraic closure and suppose the principal
surface is indecomposable. The comparison lemma gives the
irreducibility of $K$ and absence of invariant vectors in $V$:
a fixed line could not contain the sixteen-dimensional $T^0$.
The two actual natural blocks and $T^0$ form all of $K$, and
the boundary contains every rank-one map from either natural
block to $T^0$, zero on the other summands.

Let $N$ be the normal subgroup of the image on $K$ generated
by its transvections. Its restriction is semisimple: the
nonzero $N$-socle is invariant under the whole irreducible
group. Every isotypic multiplicity is one. Indeed a nontrivial
transvection has rank one and therefore can act nontrivially
on an isotypic component only with multiplicity one; a component
fixed by all generators would give a global invariant
$N$-fixed space. The first natural $\GL_8$ block belongs to
one simple $N$-summand, because its natural constituent has
multiplicity one in $K$ restricted to that actual block and
the block lies in $N$. The rank-one maps from it to all vectors
of $T^0$ put $T^0$ in that same summand. The analogous maps
from the second natural block put the second block there as
well. Thus $N$ is irreducible on all of $K$.

The previously used transvection classification
\cite[Theorem~3.6, printed p.~22]{KantorSeress2001} gives
$N=\GL(K)$ or an alternative preserving a nondegenerate
alternating form. The actual natural $\GL_8$ block has no
invariant bilinear form, and its pairing with its fixed
complement is zero since it has no invariant linear functional.
Such a block would lie in the radical of that alternating
form. Therefore the image on $K$ is $\GL_{32}(\mathbf F_2)$.

Put $Q=V/K$, the constant space of dimension $1+f$. The
translation kernel in $\operatorname{Hom}(Q,K)$ is
$K\otimes Q'$ for a subspace $Q'\subset Q^*$. Modulo this
kernel the subgroup is the graph of a cocycle, and the proved
vanishing $H^1(\GL_{32},K)=0$ makes it a linear complement.
If $Q'\ne Q^*$, a nonzero direction of $Q$ annihilated by
$Q'$ would then have a fixed lift in $V$, a contradiction.
Thus the translation kernel is all of $\operatorname{Hom}(Q,K)$.

If the principal surface is geometrically a product, the
already proved all-degree supersingular-product and mixed-product
theorems give the same degree-two formula. These cases exhaust
the two possible nonordinary product types. Finally
Theorem~\ref{thm:nonordinary-ppav-all-covectors} proves that the
entire geometric covector space descends pointwise over every
perfect base field. The arithmetic image is consequently
contained in the displayed stabilizer and contains its full
geometric subgroup, so equality follows.
\end{proof}

\begin{remark}[Scope and exact graph supplement]
\label{rem:nonordinary-ppav-degree-two-scope-candidate}
\status{Proved under stated hypotheses}
\textup{[Proved; independent audits passed.]}
The six-theta comparison takes place at relative rank $25+5f$,
not at generic rank $33+f$. Its use of clopen Lang closure and
regular component projections is essential. Opposite Picard
twists are allowed separately for the two partitions.
For an ambient hypersurface family the conclusion transfers
whenever restriction to the complete linear system is surjective,
by the previously proved geometrically integral coefficient-fibre
argument. No assertion about all higher degrees on a fixed
nonproduct surface is inferred here. The exact finite supplement
\path{agents/monodromy/code/nonordinary_degree_two_six_theta_graph.py}
checks the doubled six-vertex graph, both cross-hearts and the
two internal projections. It is a finite verification of the
displayed graph calculation, not a substitute for the geometric
or finite-Lang arguments. The author execution took
$0.000139070$ seconds. Independent root and obstructions/crystals
runs took $0.000148985$ and $0.000128674$ seconds, respectively;
both preserved the author output. Both reviewers read all proofs,
their dependencies and the entire graph program.
\end{remark}


% Promoted after complete root and reduction/descent mathematical audits.
\subsection{Adding two odd degrees on a fixed nonordinary principal surface}
\label{subsec:nonordinary-even-addition-candidate}
\status{Proved and independently checked.}

\begin{lemma}[No alternating composition factor collision]
\label{lem:nonordinary-addition-simple-group-separation-candidate}
\status{Proved under stated hypotheses}
\textup{[Proved; independent audits passed.]}
For every $r\ge8$ and $n\ge5$, the simple group
$\GL_r(\mathbf F_2)$ is not isomorphic to $A_n$.
Consequently no product of symmetric groups has a nonabelian
simple composition factor isomorphic to $\GL_r(\mathbf F_2)$.
\end{lemma}
\begin{proof}
Put $u=2^{\lceil\log_2 r\rceil}$. Every element of two-power
order in $\GL_r(\mathbf F_2)$ is unipotent, and its order is at
most $u$, by its Jordan blocks. If $n\ge2u+2$, the product of
a $2u$-cycle and a disjoint transposition belongs to $A_n$ and
has order $2u$. Hence the two groups are not isomorphic.
Otherwise
\[
       v_2|A_n|=n-s_2(n)-1\le n-2\le2u-1
                           < r(r-1)/2=v_2|\GL_r(\mathbf F_2)|.
\]
The strict inequality holds directly for $r=8,9$. For $r\ge10$,
use $u<2r$ and $r(r-1)/2>4r-1$. The last displayed inequality
at $r=10$ is $45>39$, and it increases thereafter. The same
conclusion follows if $r$ is a power of two, when $u=r$.
The general linear group is simple in this range, while the
nonabelian simple factors of a symmetric group are alternating
groups of degree at least five. This proves the consequence.
\end{proof}

\begin{proposition}[A two-component boundary with unequal odd degrees]
\label{prop:nonordinary-even-addition-boundary-candidate}
\status{Proved under stated hypotheses}
\textup{[Proved; independent audits passed.]}
Let $(A,\Theta)$ be an indecomposable nonordinary principal
surface over an algebraically closed field of characteristic
two. Put $L=\mathcal O_A(3\Theta)$ and $f=f(A)\in\{0,1\}$.
Let $a,b$ be positive odd integers for which the relative
geometric images in degrees $a,b$ are respectively
$\GL_{r_a}(\mathbf F_2)$ and $\GL_{r_b}(\mathbf F_2)$, where
$r_a=9a^2-1$, $r_b=9b^2-1$. For fixed $\alpha\in\Pic^0(A)$,
vary sharp smooth transverse curves
\[
       C\in|L^a\otimes\alpha|,\qquad
       D\in|L^b\otimes\alpha^{-1}|.
\]
Set $n=18ab$. Their unions lie in the maximal-rank degree-$a+b$
open. The boundary representation has graph kernel
$T=\mathbf F_2^n/\mathbf F_2\mathbf1$, graph heart
$T^0=\ker(\Sigma)/\mathbf F_2\mathbf1$ of dimension $n-2$,
and body
\[
       R=(H^1_{\et}(C,\mathbf F_2)\oplus
             H^1_{\et}(D,\mathbf F_2))/H^1_{\et}(A,\mathbf F_2).
\]
It contains actual natural blocks $\GL_{r_a}\times\GL_{r_b}$,
fixing a complement, and all translations $\Hom(R,T^0)$.
Every nonzero invariant subspace contains $T^0$.
\end{proposition}
\begin{proof}
The hypotheses apply to every Picard twist: multiplication
by $3a$ or $3b$ on $A$ is surjective, so each twisted ample
bundle is a translate of its untwisted bundle. The full
component ranks are $r_a+f$ and $r_b+f$. The $n$ nodes add
$n-1$ graph directions. Subtracting ambient rank $f$ gives
\[
 (r_a+f)+(r_b+f)+(n-1)-f=9(a+b)^2-3+f,
\]
the known maximal even-degree relative rank.

The lifted intersection proof in
Lemma~\ref{lem:nonordinary-ppav-two-lifted-intersections-candidate}
applies in both directions. Here are all changes needed for
unequal degrees. To cut a fixed $C$ by $L^b\otimes\alpha^{-1}$,
use the image of its sections on $C$; its kernel is now
$H^0(A,L^{b-a}\otimes\alpha^{-2})$, possibly of dimension
greater than one. The projection onto the image system still
has geometrically integral affine fibres off its centre.
The line bundle is very ample and separates all second jets:
write it as a translate of $3\Theta$ times a globally generated
bundle numerically equivalent to $3(b-1)\Theta$. At any two
prescribed points choose the latter section nonzero at both.
Multiplication then preserves the joint ranks of the first
factor on those jets. Thus the curve joint first-jet ranks
are four off a locus of dimension at most one, at least three
off a finite set, and at least two everywhere. Indeed the
rank-four surface exceptional curves are theta translates,
none equal to the integral ample $C$ of class $3a\Theta$.
The fixed-point exceptional fibres are finite by the proved
uniform theta jet theorem. At a single point the same local-unit
multiplication preserves second jets.

The two-tangency and higher-contact dimension counts are therefore
unchanged. The generic discriminant has one double zero, with
the characteristic-two quadratic cluster and labelled inertia
proved in the cited lemma. The incidence above a connected
elementary two-torsor on $C$ is an irreducible projective bundle.
Purity and geometric simple connectedness of the image coefficient
space show that its monodromy is generated by those labelled
transpositions. Hence it is $G^{n-1}\rtimes S_n$. The same proof
works with $C,D$ interchanged. Restricting to the nonempty dense
smooth sharp loci preserves these images.

The group extraction is also uniform. Each full component image
is a split extension of $\GL_{r_i}$ by either zero or a natural-dual
ambient column; it is perfect. This uses the natural/dual first
cohomology vanishing in
Lemma~\ref{lem:RD-A4-three-power-groups}, whose proof holds for
every $r_i\ge8$. The body image is a split perfect extension
of $L_0=\GL_{r_a}\times\GL_{r_b}$ by its two distinct natural-dual
isotypic modules tensored with the ambient space of dimension $f$.
The modules remain distinct if $a=b$, because the two independent
factors act on different summands.

Lemma~\ref{lem:nonordinary-addition-simple-group-separation-candidate}
shows that the node group $S_n$ has no nontrivial common quotient
with this perfect body group. Goursat gives their product. The
same vanishing for natural-dual and trivial summands in
$\Hom(R,T)$ then supplies an actual $L_0$ block. Fixing $C$,
the maximal elementary cover has group $H^1_{\et}(C,\mathbf F_2)^*$.
Its lifted node group has only alternating and cyclic simple
factors and is independent of the moving $D$ cohomology by the
same Goursat argument. Its pure translations supply all columns
from $H^1_{\et}(C,\mathbf F_2)$ to $T^0$, through the explicit
torsor-gluing evaluation formula of
Lemma~\ref{lem:nonordinary-ppav-two-boundary-gluing-candidate}.
The semisimple $L_0$-isotypic decomposition isolates the $C$ and
ambient columns. Reversing the roles supplies the $D$ columns.
Thus all of $\Hom(R,T^0)$ occurs. Finally a subspace projecting
nontrivially to $R$ contains $T^0$ by these translations, and
a nonzero subspace inside $T$ contains $T^0$ by taking a
transposition difference and its $S_n$-orbit.
\end{proof}

\begin{theorem}[Adding two known odd degrees]
\label{thm:nonordinary-even-addition-candidate}
\status{Proved under stated hypotheses}
\textup{[Proved; independent audits passed.]}
Under the geometric hypotheses of the preceding proposition,
put $d=a+b$ and $k_d=9d^2-4$. The relative image in degree $d$ is
\[
       \Hom(\mathbf F_2^{1+f},\mathbf F_2^{k_d})
                               \rtimes\GL_{k_d}(\mathbf F_2).
\]
It is the entire pointwise stabilizer of the known invariant
covectors. Their common kernel is irreducible and is contained
in every nonzero invariant subspace.
\end{theorem}
\begin{proof}
The exact even-degree covector theorem gives $1+f$ invariant
covectors and common kernel $K$ of dimension $k_d$. The actual
two natural blocks and $T^0$ in the boundary lie in $K$ and
have total dimension
\[
       r_a+r_b+(18ab-2)=9(a+b)^2-4=k_d.
\]
They consequently span $K$.

Use $m=3a+3b$ general theta translates with sum zero. Pair
intersections are transverse of length two, and there are no
triples, by the independent-translation incidence argument in
Lemma~\ref{lem:nonordinary-ppav-two-six-theta-comparison-candidate}.
This remains valid for every $m\ge6$: an unused translating
point absorbs the sum. The total graph is the doubled complete
graph on $m$ vertices. Its graph rank is $(m-1)^2$, and its
relative stable rank is $(m-1)^2+(m-1)f$, with no assertion that
this equals the generic rank.

Partition the vertices into $I$ of size $3a$ and $J$ of size
$3b$. The two component unions belong to opposite Picard
twists of $L^a,L^b$, and their $18ab$ cross nodes are distinct
smooth points of both components. Therefore their cross-heart
extends as a finite etale graph subgroup near this common
boundary. If $0\ne S\subset K$ is globally invariant, the
preceding proposition puts this heart in $S$ generically for
every partition. The clopen Lang closure puts every such
special heart in the same $\overline S_Y$.

Choose $1,2,3\in I$ and $4,5,6\in J$, and exchange vertices
$3$ and $4$ to obtain the second partition. Its heart contains
\[
                  e_{13,0}+e_{23,0},\qquad
                  e_{45,0}+e_{46,0}.
\]
On the original partition, the first vector has nonzero
projection only to the first internal graph, and the second
only to the other graph. Parallel-edge cycles detect both
nonzero projections. Ambient cohomology injects into each
theta normalization, so these graph projections survive the
relative quotient. The same strict-henselian finite-Lang
sections and regular component restriction maps used in the
six-theta proof show that $S$ has nonzero generic projection
to each natural block of the original boundary. No surjectivity
of a source kernel at a rank drop is assumed. Independence
of the two blocks puts their full sum in the image of $S$;
its kernel contains $T^0$. The displayed dimension equality
forces $S=K$.

Every nonzero global invariant subspace meets $K$ in $T^0$,
so contains $K$. In particular $V$ has no fixed vector. The
normal closure of its transvections is irreducible on $K$:
the first actual natural block has multiplicity one in one
simple summand, the rank-one gluing columns put all of $T^0$
in that summand, and the columns from the second block put
it there as well. These spaces span $K$. The exact socle,
classification and no-form argument in
Theorem~\ref{thm:every-nonordinary-ppav-degree-two-full-candidate}
therefore gives $\GL(K)$. The quotient has dimension $1+f$;
the natural-module $H^1$ vanishing and absence of fixed lifts
force the full translation kernel, as in that theorem.
\end{proof}

\begin{corollary}[Sums of two powers of three]
\label{cor:nonordinary-even-three-power-sums-candidate}
\status{Proved under stated hypotheses}
\textup{[Proved; independent audits passed.]}
Let $(A,M)$ be a principally polarized abelian surface over a
perfect field of characteristic two with $f(A)\in\{0,1\}$,
and put $L=M^3$. For every pair $i,j\ge0$ and
$d=3^i+3^j$, the geometric and arithmetic relative groups on
the smooth maximal-rank part of $|L^d|$ are
\[
 \Hom(\mathbf F_2^{1+f(A)},\mathbf F_2^{9d^2-4})
                               \rtimes\GL_{9d^2-4}(\mathbf F_2).
\]
\end{corollary}
\begin{proof}
Geometrically indecomposable surfaces satisfy the two odd-degree
hypotheses by Corollary~\ref{cor:RD-A4-powers-three}; apply the
addition theorem. Geometric products already have the same
full even-degree stabilizer by the proved product theorems.
The entire invariant covector space descends pointwise over
every perfect field, so the arithmetic image equals the full
geometric stabilizer. Surjective ambient restriction transfers
the statement exactly as before.
\end{proof}

\begin{remark}[Scope]
\label{rem:nonordinary-even-addition-scope-candidate}
\status{Proved under stated hypotheses}
\textup{[Proved; independent audits passed.]}
The addition theorem is conditional only on the two stated
odd-degree full images on the same fixed surface. The displayed
infinite family follows from the proved powers-of-three theorem.
It does not assert all even degrees: an arbitrary even integer
need not be a sum of two powers of three. The boundary is already
of maximal rank because it is a union of two sharp odd-degree
components; no ordinary auxiliary curve is introduced. The root and
reduction/descent reviewers independently read and rederived every
proof and its stated dependencies. This symbolic extension uses the
previous graph and geometric proofs; no new numerical experiment
or original-source reread is claimed by these two audits.
\end{remark}


% Post-CP23 candidate. This is a fixed-source geometric certificate, not a uniform image theorem.
\subsection{Twenty nodal five-self-isogenies of a supersingular Jacobian}
\label{subsec:five-theta-self-isogenies-candidate}
\status{Proved and independently checked.}

Let $k=\overline{\mathbf F}_2$, let $C$ be the smooth projective curve
\[
                         y^2+y=x^5,
\]
let $W$ be its unique point at infinity, and embed $C$ as the symmetric
principal theta divisor of $J=\operatorname{Pic}^0(C)$ by $P\mapsto[P-W]$.
Write $\mathbf F_{16}=\mathbf F_2[t]/(t^4+t+1)$, with field elements
encoded by their four binary coefficients. The Cartier matrix on
$\langle dx,x\,dx\rangle$ is nilpotent, since $\mathcal C(dx)=0$
and $\mathcal C(x\,dx)=dx$. Thus $f(J)=0$ and $J[2](k)=0$.

\begin{lemma}[Explicit fifth-order divisor functions and pairings]
\label{lem:five-theta-explicit-Weil-candidate}
\status{Proved under stated hypotheses}
\textup{[Proved; independent audits passed.]}
The curve has thirty-three $\mathbf F_{16}$-points. Its thirty-two
finite points lie in $J[5]$, and they generate the entire four-dimensional
symplectic $\mathbf F_5$-space $J[5](k)$. In particular every geometric
fifth-torsion point of $J$ is $\mathbf F_{16}$-rational.
For distinct finite points $P=(a,b)$ and $Q$, put
\[
             f_P=y+a^8x^2+a^4x+b^4.
\]
Then
\[
 \operatorname{div}(f_P)=5P-5W,\qquad
 e_5(P-W,Q-W)=\frac{f_P(Q)}{f_Q(P)}.
\]
The inverse convention for $e_5$ gives the same isotropic subspaces.
\end{lemma}
\begin{proof}
For $a\in\mathbf F_{16}$, $a^5\in\mathbf F_4$, whose absolute trace
from $\mathbf F_{16}$ is zero. Hence $b^2+b=a^5$ has two solutions
in $\mathbf F_{16}$, giving the thirty-two finite points.
Since $b^4=b+a^{10}+a^5$, direct substitution gives
\[
                      f_P^2+f_P=(x+a)^5.
\]
At $P$ the function vanishes, at the other point above $a$ it is one,
and its only pole is $W$, of order five. This proves its divisor.
These are the explicit functions underlying
\cite[Example~1.2(ii), pp.~30--31]{BoxallGrantLeprevost2001}.

Here is the common-pole normalization in the pairing formula.
The parameter $z=x^2/y$ is a uniformizer at $W$ and
$z^5 f_P\to1$, independently of $P$. In the divisor definition of
the Weil pairing, move $Q-W$ to a divisor disjoint from $P-W$.
For completeness choose a rational function $h$ having a simple zero
at $W$ and no zero or pole at $P,Q$; weak approximation supplies one.
If $h\sim d z$ and $f_P\sim c_Pz^{-5}$ at $W$, the tame-symbol
product formula gives
\[
           f_P(\operatorname{div}(h)-W)
                   =\frac{h(P)^5}{c_Pd^5}.
\]
Indeed the symbols at $P,W$ are respectively $h(P)^{-5}$ and
$c_Pd^5$, and all other symbols give the evaluation on the left.
Using the disjoint divisor $Q-W+\operatorname{div}(h)$ and its
function $f_Qh^5$ therefore gives
$f_P(Q)c_Q/(f_Q(P)c_P)$. Here $c_P=c_Q=1$; the usual sign is also
one in characteristic two. This proves the displayed formula.

The retained exact program enumerates the thirty-three points,
evaluates these functions, and expresses all pairings as powers of
$\zeta=t^3$. Four finite points, numbered $1,3,5,9$ in its displayed
point list, have pairing matrix
\[
 J_5=\begin{pmatrix}
 0&0&4&3\\0&0&4&2\\1&1&0&3\\2&3&2&0
 \end{pmatrix}\quad\text{over }\mathbf F_5.
\]
Its determinant is nonzero, so their classes are independent.
The group $J[5](k)$ has dimension four; these rational classes are
therefore a basis. The program reconstructs the coordinates of
every curve point from its four pairings and verifies all pairings
again in those coordinates. In particular this computation does
not assume a Jacobian group-law formula.
\end{proof}

\begin{lemma}[A complete geometric theta-packet census]
\label{lem:five-theta-geometric-census-candidate}
\status{Proved under stated hypotheses}
\textup{[Proved; independent audits passed.]}
Among the $156$ maximal isotropic subgroups $K\subset J[5]$, exactly
twenty give a geometrically integral nodal image of $C$ in $J/K$.
Each such image has twenty-four ordinary nodes. For the other
subgroups, the largest rational packet has size four for eighty
subgroups and size five for fifty-six subgroups.
\end{lemma}
\begin{proof}
For a nonzero $g\in J(k)$, $C\cap(C-g)$ has length two: theta has
self-intersection two and no nontrivial translation stabilizer.
If $P,P+g\in C$, symmetry also gives $-P-g,-P\in C$.
These describe the two intersection points, with their possible
coincidence. Suppose a geometric packet contains three distinct
points $P,P+g,P+h$, where $g,h\in J[5]$ are distinct and nonzero.
Both intersection schemes are defined over $\mathbf F_{16}$.
If $P$ were not rational, its Frobenius image would have to be
the other point of each length-two scheme. Thus
\[
             \operatorname{Frob}_{16}(P)=-P-g=-P-h,
\]
which contradicts $g\ne h$. Consequently every geometric packet
of size at least three is already in the rational point list.

The tangent directions of this theta embedding are the fibres of
the canonical degree-two map of $C$. Hence distinct identified
points can have equal tangents only when they are $P$ and $-P$.
If $2P\in K\subset J[5]$, then $[5]P\in J[2](k)=0$; therefore
$P\in J[5](\mathbf F_{16})$. Moreover multiplication by two is
invertible on $K$, so $P\in K$. The packet containing $W$ then
contains the three distinct points $W,P,-P$. Thus a rational
census with every packet of size at most two also excludes every
geometric tangency.

The program enumerates the $156$ projective lines in
$\mathbf F_5^4$, and the $156=(5+1)(5^2+1)$ isotropic planes by
spanning every pair of independent orthogonal lines. For each
plane it partitions the thirty-three curve vectors into cosets.
It retains every plane basis and every packet, and finds
\[
 \begin{array}{c|ccc}
 \text{largest packet}&2&4&5\\ \hline
 \text{number of planes}&20&80&56.
 \end{array}
\]
The preceding geometric argument proves that the twenty first
planes have only transverse double identifications, and that the
other planes cannot have nodal images.

The quotient is etale. Its restriction to $C$ is birational onto
its image: a positive-degree generic identification would make
$C$ invariant under a nonzero element of $K$, contrary to its
principal polarization. Thus the image is integral with
normalization $C$, and the identified transverse pairs give exactly
ordinary nodes. Its numerical class on the principally polarized
quotient is $5\Theta'$, so its arithmetic genus is $26$ and the
node number is $26-2=24$.
\end{proof}

\begin{proposition}[Explicit nodal polarized five-self-isogenies]
\label{prop:five-theta-twenty-self-isogenies-candidate}
\status{Proved under stated hypotheses}
\textup{[Proved; independent audits passed.]}
Let $\alpha,\beta\in\mathbf F_{16}$ satisfy
\[
        \alpha^5\in\mathbf F_4\setminus\mathbf F_2,
        \qquad \beta^2+\beta=\alpha^5.
\]
The automorphism
\[
 g_{\alpha,\beta}(x,y)
   =(x+\alpha,\ y+\alpha^8x^2+\alpha^4x+\beta)
\]
induces an automorphism $g$ of the principal Jacobian with
$g^2=[-1]$. The isogeny $\phi=[2]+g$ has degree twenty-five and
$\phi^*\Theta\equiv5\Theta$. Its kernel is one of the twenty good
planes of the preceding lemma. As $(\alpha,\beta)$ varies, these
twenty kernels are distinct and exhaust the good planes.
Consequently every such $\phi$ gives a nodal theta image on this
same fixed Jacobian, defined over $\mathbf F_{16}$.
\end{proposition}
\begin{proof}
Substitution in the curve equation verifies the automorphism.
Its square fixes $x$ and adds
$\alpha^{10}+\alpha^5=1$ to $y$, so it is the hyperelliptic
involution. It fixes $W$ and preserves the principal polarization;
its Rosati adjoint is $g^{-1}=-g$. Therefore
\[
                 \phi^\dagger\phi=(2-g)(2+g)=[5].
\]
It is a separable polarized isogeny of degree twenty-five, and its
kernel is maximal isotropic in $J[5]$.

There are ten possible $\alpha$ and two possible $\beta$ for each.
The second retained program computes their actions on the four
explicit fifth-torsion basis points of the first lemma. It verifies
the entire action on all thirty-three points, $g^2=-I$, and
$g^{\mathsf t}J_5g=J_5$. It then enumerates
$\ker(2I+g)\subset\mathbf F_5^4$ and verifies that the twenty
resulting order-twenty-five subgroups are distinct and exactly
the twenty good planes of the complete census.

The norm divisor is symmetric. Its line bundle differs from
$\mathcal O_J(5\Theta)$ by a symmetric element of $\operatorname{Pic}^0(J)$,
which is a geometric two-torsion point and hence zero. Thus these
are actual nodal divisors in the indicated complete system,
not merely numerical classes.
\end{proof}

\begin{remark}[Certificate and remaining monodromy scope]
\label{rem:five-theta-self-isogeny-scope-candidate}
\status{Proved under stated hypotheses}
The executed programs are
\path{code/monodromy/nonordinary_five_theta_packet_probe.py} and
\path{code/monodromy/nonordinary_five_theta_self_isogenies.py}, with complete
JSON outputs of the same stems in \path{agents/monodromy/outputs/}. The first passed
at 20:08:39 UTC in $0.31384029000037117$ seconds; an independent
output-preserving rerun passed in $0.2980953239966766$ seconds.
The second passed at 20:24:45 UTC in $0.051931495996541344$ seconds.
The parent completed the full mathematical and source audit at 20:31:15 UTC and independently reran both complete programs, preserving the author outputs, in $0.334064817$ and $0.058760470$ seconds. The independent reduction/descent audit also checked the Weil functions and finite census; no broader geometric scope is attributed to that bounded audit.
This establishes good degree-five sources and targets only for the
specified Jacobian. It does not establish a suitable five-neighbour
for every fixed nonordinary surface, nor the required degree-five
node permutation and gluing image. The order-three eigenmap
argument cannot be imported for a fifth-torsion translation,
which need not preserve $\mathcal O(3\Theta)$.
\end{remark}


% Post-CP23 candidate; no degree-five full-image conclusion is asserted.
\subsection{Uniform supersingular five-theta geometry and four evaluations}
\label{subsec:uniform-five-theta-geometry-candidate}
\status{Proved and independently checked.}

\begin{lemma}[A cyclic order-four presentation of every supersingular genus-two curve]
\label{lem:uniform-f0-order-four-presentation-candidate}
\status{Proved under stated hypotheses}
\textup{[Proved; independent audits passed.]}
Over an algebraically closed field of characteristic two, every genus-two
curve of two-rank zero is isomorphic to
\[
 C_z:\quad y^2+y=a x^5+b x^3,\qquad
 a=z^2\ne0,\quad b=z^2+z+1.
\]
Choose $\beta$ with $\beta^2+\beta=z+1$. The formula
\[
       g(x,y)=(x+1,\ y+z x^2+(z+1)x+\beta)
\]
defines an automorphism fixing the unique point $W$ at infinity,
and $g^2=\iota$ is the hyperelliptic involution.
\end{lemma}
\begin{proof}
The two-rank-zero hyperelliptic normal form has one branch point;
after placing it at infinity, Artin--Schreier reduction gives
$y^2+y=x^5+\mu x^3+\nu x$ after scaling and removing the constant.
This is also the normal-form input in
Lemma~\ref{lem:p-rank-zero-genus-two-polar-stable-candidate}.
For clarity the linear term can be removed directly. Choose a root
$t$ of
\[
                 t^{16}+\mu^4t^8+\mu^2t^2+t+\nu^4=0.
\]
After $x\mapsto x+t$, remove the fourth and second powers by the
Artin--Schreier shift $y\mapsto y+A x^2+B x$, where
$A^2=t$ and $B^2=\mu t+A$. The chosen equation for $t$ implies
$B=t^4+\mu t^2+\nu$, by injectivity of squaring, so the remaining
linear coefficient is zero. Remove the new constant by a constant
Artin--Schreier shift, possible over the algebraically closed field.
Thus the model is
$y^2+y=x^5+\mu x^3$; the same normal form is explicitly recorded in
\cite[pp.~30--31]{BoxallGrantLeprevost2001}.

Choose a nonzero root $z$ of
\[
                       (z^2+z+1)^5=\mu^5z^6.
\]
This polynomial has degree ten and constant coefficient one, so such
a root exists. If $\mu=0$, both cubic coefficients are zero and any
fifth-root scaling with $\rho^5=z^2$ gives the required model.
If $\mu\ne0$, choose $\rho^5=z^2$ and adjust it by a fifth root of
unity so that $\mu\rho^3=z^2+z+1$. The displayed fifth-power
identity permits this adjustment, since cubing permutes the fifth
roots of unity. Replacing $x$ by $\rho x$ gives $C_z$.

Put $q(x)=z x^2+(z+1)x+\beta$ and $F(x)=a x^5+b x^3$.
Directly
\[
       q(x)^2+q(x)=F(x+1)+F(x),\qquad q(x+1)+q(x)=1.
\]
These identities prove that $g$ is an automorphism and its square
adds one to $y$. Its polynomial formula extends over the unique
point at infinity, which it fixes.
\end{proof}

\begin{theorem}[A nodal polarized five-self-isogeny on every supersingular Jacobian]
\label{thm:uniform-f0-five-self-isogeny-nodal-candidate}
\status{Proved under stated hypotheses}
\textup{[Proved; independent audits passed.]}
Let $C$ be any genus-two curve of two-rank zero over an algebraically
closed field of characteristic two. Embed $C=\Theta\subset J(C)$
using its unique Weierstrass point $W$. There is a separable isogeny
$\phi:J(C)\to J(C)$ with
\[
                \deg\phi=25,\qquad \phi^*\Theta\equiv5\Theta,
\]
whose theta image is an integral curve in $|5\Theta|$ with exactly
twenty-four ordinary nodes and no other singularities.
\end{theorem}
\begin{proof}
Use the preceding presentation and let $g$ also denote its action
on the principal Jacobian. It preserves the polarization and has
$g^2=[-1]$, so $g^\dagger=-g$. Set $\phi=[2]+g$; then
\[
                   \phi^\dagger\phi=(2-g)(2+g)=[5].
\]
It is a separable polarized isogeny of degree twenty-five.

First $\ker\phi\cap\Theta=\{W\}$. For a finite point $P$, membership
would give $2P+gP\sim3W$. But $H^0(C,\mathcal O_C(3W))$
has basis $1,x$, so $W$ is a base point of $|3W|$. The effective
divisor $2P+gP$ does not contain $W$, a contradiction.

We next bound all fibres of $\phi|_\Theta$ using a direct reduced
divisor calculation. Write $P=(p,y_0)$, set $X=x+p$, and put
\[
 \begin{split}
 c_1&=a p^4+b p^2,\\
 c_2&=c_1+z p^2+(z+1)p+\beta,\\
 v_P(x)&=y_0+c_1X+c_2X^2,\\
 \ell_P&=a p+c_2^2+a.
 \end{split}
\]
The function $v_P$ interpolates the value and first derivative of
$y$ at $P$, and its value at $p+1$ is the $y$-coordinate of $gP$.
Using $y_0^2+y_0=F(p)$ gives the exact identity
\[
 v_P(x)^2+v_P(x)+F(x)
       =X^2(X+1)\bigl(aX^2+\ell_PX+b+\ell_P\bigr).
\]
The function $y+v_P(x)$ has its sole pole of order five at $W$.
Its zero divisor is therefore $2P+gP+E_P$, where $E_P$ is the
effective finite degree-two residual divisor whose $x$-polynomial,
with leading coefficient $a$, is
\[
 U_P(x)=a x^2+\ell_Px+
                        a p^2+\ell_Pp+b+\ell_P.
\]
This statement permits multiplicities and accidental overlap with
$P$ or $gP$: the displayed factorization subtracts their prescribed
effective divisors and leaves the residual length two.

The residual divisor cannot contain a hyperelliptic pair, since
$y+v_P(x)$ cannot vanish at both points above the same finite $x$:
their $y$-coordinates differ by one. Thus $E_P$ is nonspecial.
Equivalently, its degree-two linear system has dimension zero.
The divisor equation says
\[
                  \phi(P-W)=-[E_P-2W].
\]
Consequently equal images of finite $P,Q$ imply $E_P=E_Q$ and
$U_P=U_Q$. Writing $Q=(q,y_1)$, comparison of the linear and
constant coefficients gives
\[
 \ell_Q=\ell_P=\ell,\qquad
                 (p+q)\bigl(a(p+q)+\ell\bigr)=0.
\]
For a fixed image there are at most two possible $x$-coordinates:
$p$ and $p+\ell/a$.

There cannot be two distinct points with the same $x$-coordinate
in one fibre. They would be $P$ and $\iota P=-P$ in the Jacobian,
so $2\phi(P-W)=0$. Two-rank zero gives $J[2](k)=0$, hence
$\phi(P-W)=0$, contrary to the first paragraph. The point $W$
has no other point in its fibre by that paragraph. Every geometric
fibre of $\phi|_\Theta$ therefore contains at most two points.

The restriction is birational onto its image, since otherwise the
theta divisor would be invariant under a nontrivial kernel translation.
It is an immersion because $\phi$ is etale. Its two branches at an
identified pair have equal tangent lines only if they are a
hyperelliptic pair: the Gauss map of a principal theta curve is its
canonical degree-two map. Such a pair was just excluded. Thus its
only singularities are transverse ordinary double points.
The image is symmetric and has numerical class $5\Theta$.
Its difference from $\mathcal O_J(5\Theta)$ is a symmetric
degree-zero line bundle, hence a geometric two-torsion point;
it is trivial because $f(J)=0$. Its arithmetic genus is $26$,
so normalization genus two gives exactly twenty-four nodes.
\end{proof}

\begin{lemma}[Four evaluations and two translated pairs]
\label{lem:nonproduct-three-theta-four-values-candidate}
\status{Proved under stated hypotheses}
\textup{[Proved; independent audits passed.]}
Let $(A,\Theta)$ be an indecomposable principally polarized surface
over an algebraically closed field of characteristic two, with
$\Theta$ symmetric. If four distinct points $P_0,P_1,P_2,P_3$
fail to impose four independent conditions on
$H^0(A,\mathcal O_A(3\Theta))$, then
\[
                    \sum_{j\ne i}P_j\in\Theta
                           \quad(0\le i\le3).
\]
For fixed nonzero $h,k\in A$, the locus of pairs $(P,Q)$ for which
the four distinct points $P,P+h,Q,Q+k$ fail to impose four
conditions is finite. The analogous assertions hold for every
Picard twist of $\mathcal O_A(3\Theta)$ after translating coordinates.
\end{lemma}
\begin{proof}
First any three distinct points impose three independent conditions.
To separate $P_2$ from $P_0,P_1$, choose
$a\in P_0-\Theta$ outside $P_2-\Theta$ and the single point
$-P_1-P_2$. Next choose $b\in P_1-\Theta$ outside
$P_2-\Theta$ and $-a-P_2+\Theta$. The latter two intersections
are finite: equality of the second pair of theta translates would
force the excluded value $a=-P_1-P_2$. With $c=-a-b$, the product
of the three theta sections vanishes at $P_0,P_1$ and not at $P_2$.
Repeating for each excluded point proves the assertion. This also
recovers Lemma~\ref{lem:ppav-three-distinct-values} directly.

We exhibit a section vanishing at three chosen points and not at
the fourth whenever their sum is outside $\Theta$. Put
$t=P_0+P_1+P_2$ and suppose $t\notin\Theta$. In the surface
$\Theta\times\Theta$ consider the nonempty effective divisor
\[
        Z_t=\{(u,v):u+v\in\Theta+t\}.
\]
The addition map $\Theta^2\to A$ is dominant, so this is a proper
divisor. None of its curve components has constant first projection,
constant second projection, or constant sum. Indeed such a component
would respectively have $u=t$, $v=t$, or $u+v=0$: equality of the
corresponding theta translates follows from containment of integral
curves and the trivial translation stabilizer of a principal theta
divisor. Each alternative forces $t\in\Theta$.

Set $a=P_0-u$, $b=P_1-v$, and $c=-a-b$. The product of the three
theta sections has divisor
$\Theta+a+\Theta+b+\Theta+c$ and lies in $|3\Theta|$, by the
theorem of the square. It vanishes at $P_0,P_1,P_2$.
The conditions that it also vanish at $P_3$ are, respectively,
\[
 \begin{split}
 u&\in\Theta\cap(\Theta+P_0-P_3),\\
 v&\in\Theta\cap(\Theta+P_1-P_3),\\
 u+v&\in(\Theta+t)\cap(\Theta+P_0+P_1+P_3).
 \end{split}
\]
These are finite sets because the points are distinct. On any
component of $Z_t$ the three relevant maps are nonconstant, so
one can avoid all three conditions. Thus $P_3$ can be separated
from the other three. Repeat with each excluded index to prove
the necessary four-sum condition.

For the indicated two pairs put
\[
             U=P+2Q+h+k,\qquad V=2P+Q+h+k.
\]
The necessary conditions say
$U,U-h\in\Theta$ and $V,V-k\in\Theta$.
Each allowed set is finite. The homomorphism
$(P,Q)\mapsto(P+2Q,2P+Q)$ has determinant $-3$; it is a finite
etale isogeny in characteristic two. Hence its inverse image of
those finite sets is finite. Finally every Picard twist of
$3\Theta$ is a translate of that bundle, since $[3]$ is surjective.
\end{proof}

\begin{remark}[Universal certificate and remaining image gate]
\label{rem:uniform-five-theta-geometry-scope-candidate}
\status{Proved under stated hypotheses}
The complete program
\path{code/monodromy/nonordinary_uniform_five_self_isogeny_identities.py}
checks nine universal characteristic-two polynomial identities,
including the automorphism, the residual quadratic, its global
coefficients and the equal-fibre factorization. It passed at
20:34:32 UTC in $0.08886713300307747$ seconds; the complete result is
\path{agents/monodromy/outputs/nonordinary_uniform_five_self_isogeny_identities.json}.
An initial coefficient-domain setup failure in its final leading-term
check is retained with the exact failed program in \path{history/};
no mathematical identity failed. The parent completed the full proof and code audit at 20:40:04 UTC and independently reran all identities in $0.101561034$ seconds. Reduction/descent completed the full proof and code audit at 20:39:50 UTC, with an output-preserving rerun in $0.36796294799569296$ seconds. The theorem proves a nodal theta
five-self-isogeny for every fixed supersingular Jacobian. It does
not yet prove the full degree-five monodromy: isolation of the
node-cover inertia and all gluing columns remains a separate gate.
\end{remark}


% /root, 19 November 2026. Independently audited by both other Part-A agents.
\subsection{All invariant covectors on nonordinary principally polarized surfaces}
\label{subsec:nonordinary-ppav-covectors}
\status{Proved and checked. Both the monodromy and obstruction/crystal agents
independently rederived the complete proof, checked the source scope, and
reran the retained exact local certificate without changing the root output.}

\begin{lemma}[Removing an arbitrary nilpotent length]
\label{lem:linear-covectors-any-nil-length}
\status{Proved under stated hypotheses}
Let $k$ be algebraically closed of characteristic two, let $W$ be a
finite-dimensional $k$-space with coordinates $s_1,\ldots,s_b$, and put
$P=\mathbf P(W)$ and $\mathcal B=W^*\otimes\mathcal O_P(-1)$.
Suppose a semilinear Frobenius on $\mathcal B$ has matrix
$M(s)^t$, where $M(s)=\sum_i s_iM_i$ and
$\bigcap_i\ker M_i=0$. On the open where the iterated linearized image
has maximal stable rank, let $\mathcal I$ be that unit-Frobenius image.
Every invariant $\mathbf F_2$-covector on its Lang local system is the
restriction of exactly one linear expression
\[
             \lambda_s=D(s),\qquad D\in\operatorname{End}_k(W),
             \qquad M(s)D(s)=D(s)^{[2]}.
\]
No assertion that the first Frobenius image is already stable is needed.
\end{lemma}
\begin{proof}
We first work at the generic point with field $K=k(P)$.
Choose $e$ such that the image of the $e$-fold linearized Frobenius is
the stable image $I$, and put $q=2^e$. We may increase $e$ so $e\geq1$.
An invariant covector is a $K$-linear map $f:I\to K$ satisfying
$f(Fx)=f(x)^2$. Extend it, provisionally over $K^{1/q}$, by
\[
                    \lambda(x)^q=f(F^e x),\qquad x\in\mathcal B_K.
\]
The unique $q$-th root is additive and $K$-linear in $x$.
It agrees with $f$ on $I$, and
$\lambda(Fx)=\lambda(x)^2$. Its coefficient vector raised to the
$q$-th power is a rational section
$w\in F_P^{e*}\mathcal B^*=W^{[q]}\otimes\mathcal O_P(q)$, with
\begin{equation}\label{eq:any-nil-lift-polynomial}
                    M(s)^{[q]}w=w^{[2]}.
\end{equation}
Here the brackets mean coefficientwise powers, and the equation is a
section equation in $W^{[2q]}\otimes\mathcal O_P(2q)$.
The formulas in different coefficient charts agree because a covector
on $\mathcal O(-1)$ has transition weight one.

If $w$ had a pole of largest order $h>0$ along a prime divisor, the
right side would have a coordinate of order $-2h$, whereas the left
side, whose matrix is regular, has every coordinate of order at least
$-h$. This is impossible. Normality of $P$ extends $w$ across the
remaining codimension-two set. Thus $w$ is a homogeneous polynomial
vector of degree $q$.

Write $w=\sum_{|\alpha|=q}w_\alpha s^\alpha$.
We prove inductively that every exponent occurring is divisible by
$2^j$, for $j=0,\ldots,e$. Suppose this holds for $j<e$, and choose
$\alpha$ not divisible coordinatewise by $2^{j+1}$.
It is not a pure exponent $qe_i$, so every $\alpha_i<q$.
For each $i$, the coefficient of $s_i^q s^\alpha$ on the left of
\eqref{eq:any-nil-lift-polynomial} is uniquely
$M_i^{[q]}w_\alpha$: no other pair $(i',\alpha')$ gives this monomial.
On the right, every exponent is divisible by $2^{j+1}$, so that
coefficient is zero. The common-kernel hypothesis, preserved under
Frobenius over perfect $k$, gives $w_\alpha=0$.
At $j=e$ only pure exponents remain. Taking their unique $q$-th roots
in $k$ gives $w=D(s)^{[q]}$, and then $\lambda=D(s)$ and the stated
matrix identity. Conversely that identity defines an invariant
covector. If its restriction to $I$ is zero, iteration gives
$\lambda(x)^q=\lambda(F^ex)=0$ for every $x$, so $D=0$.
\end{proof}

\begin{lemma}[The four-dimensional endomorphism algebra]
\label{lem:surface-frobenius-distributions}
\status{Proved under stated hypotheses}
Let $A/k$ be an abelian surface over an algebraically closed field of
characteristic two, and let $F=F_{A/k}:A\to A^{(2)}$.
Writing $K=\ker F$, there are canonical algebra identifications
\[
 \operatorname{End}_{A^{(2)}}(F_*\mathcal O_A)
       =\mathcal O(K)^*
       =k[\delta_1,\delta_2]/
          (\delta_i^2-\delta_i^{[2]}),
\]
where $\delta_1,\delta_2$ form a basis of invariant vector fields,
the multiplication on $\mathcal O(K)^*$ is convolution, and
$\delta_i^{[2]}$ is the invariant vector field obtained by composing
$\delta_i$ with itself. The vector-space basis is
$1,\delta_1,\delta_2,\delta_1\delta_2$.
If $f=f(A)$ is the two-rank, a basis of the invariant vector fields
can be chosen with the following relations:
\[
\begin{array}{c|cc}
 f\text{ and nilpotent type}&\delta_1^2&\delta_2^2\\\hline
 2&\delta_1&\delta_2\\
 1&\delta_1&0\\
 0,\ \text{zero map}&0&0\\
 0,\ \text{one chain}&\delta_2&0.
\end{array}
\]
\end{lemma}
\begin{proof}
The difference map gives
$A\times_{A^{(2)}}A\simeq A\times K$, since $F$ is a $K$-torsor.
Finite flat base change and adjunction therefore identify the
endomorphism space with
\[
 \operatorname{Hom}_A(F^*F_*\mathcal O_A,\mathcal O_A)
 =\operatorname{Hom}_A(\mathcal O_A\otimes_k\mathcal O(K),\mathcal O_A)
 =\mathcal O(K)^*,
\]
using $H^0(A,\mathcal O_A)=k$. Composition is convolution by the
torsor action. At the identity, smooth local parameters identify the
coordinate algebra of $K$ with $k[u,v]/(u^2,v^2)$ as an algebra.
The invariant differential operators $1,\delta_1,\delta_2,
\delta_1\delta_2$ have independent symbols dual to
$1,u,v,uv$. They therefore form a basis of its distribution algebra.
They commute, and their squares are primitive invariant derivations;
the displayed relations follow, with no additional relation by the
dimension count.

The square map on the two-dimensional space of invariant derivations
is semilinear. Its stable and nilpotent parts split over perfect $k$.
On the stable part Lang's equation supplies a basis fixed by squaring.
On a nilpotent part of dimension at most two there is either a zero
map or one chain; scaling the chain gives the displayed forms.
If the stable dimension is $a$, the distribution algebra has $2^a$
geometric points: each toral generator has values zero or one, and
nilpotent generators have value zero. But this algebra is the
coordinate algebra of the Cartier dual $K^D$.
The dual-isogeny identity gives
$K^D=\ker(V_{A^\vee/k})$; as relative Frobenius is bijective on
geometric points, this group has $2^{f(A^\vee)}=2^{f(A)}$ points.
Hence $a=f$.
The exact sources for these last identifications are
\citet[Theorem~7.5, p.~100; Proposition~7.34 and proof, p.~111;
Proposition~5.22 and proof, pp.~80--81]{EdixhovenVanDerGeerMoonenAV}.
In particular the argument concerns the whole scheme-theoretic
Frobenius kernel, not only its one geometric point.
\end{proof}

\begin{theorem}[Entire invariant quotient on every nonordinary surface]
\label{thm:nonordinary-ppav-all-covectors}
\status{Proved under stated hypotheses}
Let $(A,\Theta)$ be a principally polarized abelian surface over a
perfect field $k$ of characteristic two, with the polarization
represented by the line bundle $\mathcal O_A(\Theta)$ over $k$.
Assume $f(A)<2$ and put $L=\mathcal O_A(3\Theta)$.
For the relative Hasse--Witt local system $V_d$ on the smooth
maximal-rank open in the complete linear system $|L^d|$, for every
$d\geq1$ one has
\begin{equation}\label{eq:nonordinary-ppav-entire-covectors}
 \dim_{\mathbf F_2}(V_d^*)^{G_{\rm geom}}
 =\dim_{\mathbf F_2}(V_d^*)^{G_{\rm arith}}
 =\begin{cases}0,&d\text{ odd},\\2^{f(A)},&d\text{ even}.
 \end{cases}
\end{equation}
In particular the entire even-degree invariant quotient has dimension
one in both two-rank-zero types, and dimension two in two-rank one.
For hypersurfaces in the fixed $3\Theta$ embedding the same statements
hold whenever restriction onto $H^0(A,L^d)$ is surjective, in particular
for $d\geq\max(1,\operatorname{reg}(\mathcal I_A)-1)$.
This determines the invariant quotient, not the whole monodromy group.
\end{theorem}
\begin{proof}
We first work over $\bar k$ and choose an invariant top differential
$\omega$. Set $N=L^d$, $W=H^0(A,N)$, and $P=\mathbf P(W)$.
The full coherent bundle is
$\mathcal B=H^2(A,N^{-1})\otimes\mathcal O_P(-1)$, identified with
$W^*\otimes\mathcal O_P(-1)$ by Serre duality and $\omega$.
Its relative Euler subbundle is $\mathcal E=\Omega_P^1$.
The quotient $\mathcal B/\mathcal E$ is $H^2(A,\mathcal O_A)\otimes
\mathcal O_P$ with Frobenius zero, because $A$ is nonordinary.
Thus $\mathcal B$ and $\mathcal E$ have the same stable local system.
In the notation of the previous Cartier sections, the relative,
root-removed pairing is the $k$-bilinear map
\[
 \beta:W\otimes_kW\longrightarrow W^{(2)},\qquad
           (s,t)\longmapsto C_{\rm rel}(st\omega)/\omega^{(2)}.
\]
Its multiplication matrix $M(s)$ is linear in $s$, and the Frobenius
matrix on $\mathcal B$ is $M(s)^t$.

We check the common-kernel and descent hypotheses, including small
degrees. The vector bundle $F_*N$ is globally generated. In the
indecomposable theta case, Theorem~\ref{thm:ppav-three-theta-jets}
separates second jets for $3\Theta$; a Frobenius fibre has local
algebra $k[x,y]/(x^2,y^2)$, a quotient of the second-jet algebra.
In the product case the two degree-three elliptic bundles separate
length-two fibres, and their tensor product separates the full
Frobenius fibre. For higher $d$, multiply by a section of $L^{d-1}$
nonzero at the point: it is a unit on that finite local fibre.
The relative Cartier pairing
\[
                    F_*N\otimes F_*N\longrightarrow N^{(2)}
\]
is perfect, by the finite-Frobenius duality already checked in
\citet[Proposition~1.3.7 and Theorem~1.3.8, pp.~26--28]{BrionKumar2005}.
If $M(s)t=0$ for all $s$, fibrewise global generation and this
perfection force $t=0$. This proves the common-kernel hypothesis.

Lemma~\ref{lem:linear-covectors-any-nil-length} now represents every
invariant covector uniquely by a $D\in\operatorname{End}(W)$ such that
\begin{equation}\label{eq:nonordinary-exact-D}
                       \beta(s,D(s))=D(s)^{[2]}.
\end{equation}
Polarization in $s,t$ gives $\beta(s,Dt)=\beta(t,Ds)$.
If $s$ vanishes in one Frobenius fibre, this identity and the perfect
pairing show that $Ds$ vanishes there too. The kernel of the globally
surjective evaluation map is locally free. Thus the identity on all
closed fibres shows that $D$ descends to an
$\mathcal O_{A^{(2)}}$-linear endomorphism of $F_*N$.
No assumption on the nilpotent length was used.

If $d$ is odd, the polarization of $N$ is $[3d]\lambda_\Theta$.
Its restriction to the scheme $\ker F$ is injective, since $3d$ is odd.
\citet[Theorem~1.2(i), p.~45, with proof on pp.~43--45]{Oda1971}
therefore gives $\operatorname{End}(F_*N)=k$.
For a scalar identity the left side of
\eqref{eq:nonordinary-exact-D} is zero, since $C(\omega)=0$.
The only exact scalar lift is zero. This proves the odd-degree case.

For even $d$, write $N=F^*M$ with
$M=(L^{(2)})^{d/2}$ on $A^{(2)}$.
Projection formula gives
$\operatorname{End}(F_*N)=\operatorname{End}(F_*\mathcal O_A)$.
The operators of Lemma~\ref{lem:surface-frobenius-distributions}
act on $N$ by its Frobenius-pullback connection: in a pullback frame
they act on its coefficient function. We calculate all of them,
not merely a selected subspace.

Here are rational coordinates giving every required normal form.
All derivations annihilate the field of squares. In the zero-map
case choose $z$ with $a=\delta_1\delta_2z\ne0$; this exists because
the top distribution acts nontrivially on the generic torsor.
It is a square-field constant. Put
$x=\delta_2z/a$, $y=\delta_1z/a$. Then
$\delta_1=\partial_x$ and $\delta_2=\partial_y$.
In the one-chain case choose $z$ with $a=\delta_1^3z\ne0$ and put
$x=\delta_1^2z/a$, $y=\delta_1z/a$.
Then $\delta_1=\partial_x+x\partial_y$ and
$\delta_2=\partial_y$. In both cases $x,y$ are a two-basis and
the invariant volume can be normalized to $dx\wedge dy$:
its contraction with the invariant frame is a nonzero global constant.

In the mixed case the eigenspaces of $\delta_1^2=\delta_1$ on the
four-dimensional square-field extension have dimensions two and two;
multiplication by a nonzero weight-one element identifies them.
The nilpotent $\delta_2$ acts nontrivially on the weight-zero space.
Otherwise it acts there by zero, and on a weight-one generator $u$
as $\delta_2u=bu$; squaring gives $b^2=0$, hence $\delta_2=0$,
a contradiction. Choose weight-zero $z$ with $a=\delta_2z\ne0$.
Then $a$ is killed by both derivations, so belongs to the square field;
$y=z/a$ has weight zero and $\delta_2y=1$.
There is a nonzero weight-one $x$ killed by $\delta_2$: use any such
element if it is killed, and otherwise its nonzero $\delta_2$-image.
Then $\delta_1=x\partial_x$, $\delta_2=\partial_y$, and the volume
can be normalized to $(dx/x)\wedge dy$.

Write a general local section coefficient as
$s=a^2+xb^2+yc^2+xye^2$. On the ordered operator basis
$1,\delta_1,\delta_2,\delta_1\delta_2$, define $Q$ by
$C(sD(s)\omega)=Q(D)(s)\omega$.
Direct extraction of the odd exponents in $dx\wedge dy$ gives
\begin{equation}\label{eq:nonordinary-Cartier-table}
\begin{array}{c|cccc}
 &Q(1)&Q(\delta_1)&Q(\delta_2)&Q(\delta_1\delta_2)\\\hline
 \text{mixed}&0&0&\delta_2&\delta_1\delta_2\\
 \text{zero map}&0&0&0&\delta_1\delta_2\\
 \text{one chain}&0&\delta_2&0&\delta_1\delta_2.
\end{array}
\end{equation}
The map $Q$ is inverse-Frobenius-semilinear in scalar coefficients.
For example in the one-chain case
$\delta_1s=b^2+xc^2+(y+x^2)e^2$ and
$C(s\delta_1s\,dx\wedge dy)=(c^2+xe^2)dx\wedge dy$,
whereas $\delta_1\delta_2s=e^2$ is fixed by the same calculation.
In the mixed case the coefficient-extraction rule for $dx/x$ selects
even powers of $x$ and that for $dy$ selects odd powers of $y$;
it gives the first row. The retained exact program
\path{code/ppav_nonordinary_covector_checks.py} expands and verifies
all twelve displayed identities over the polynomial ring in
$x,y,a,b,c,e$, with the indicated Laurent differential.

Solving $D=Q(D)$ gives the two-dimensional $\mathbf F_2$-space
spanned by $\delta_2,\delta_1\delta_2$ in the mixed case, and the
one-dimensional $\mathbf F_2$-space spanned by
$\delta_1\delta_2$ in either nilpotent case.
Uniqueness in Lemma~\ref{lem:linear-covectors-any-nil-length}
shows that these are precisely the invariant covectors and that
none vanishes on the stable part.

Finally we check arithmetic descent. For two-rank zero the entire
geometric invariant space has dimension one over $\mathbf F_2$;
every Galois action on it is therefore trivial. In two-rank one the
stable line of the square map on invariant derivations has exactly
one nonzero fixed vector $\delta_1$, so it is defined over the
perfect field $k$. Its nilpotent line is likewise defined over $k$.
Choose $\omega$ over $k$ and the unique $\delta_2$ on that line with
$\omega(\delta_1,\delta_2)=1$.
Both exact operators $\delta_2$ and $\delta_1\delta_2$ are then
defined over $k$. The entire geometric invariant space is arithmetic.
The normality of the maximal-rank open ensures that passage to the
dense smooth-member open preserves the image of every finite cover.
Surjectivity of ambient restriction gives the same images by the
geometrically integral fibres of the associated linear projection;
the stated regularity bound guarantees that surjectivity.
\end{proof}

\begin{remark}[Retained exact checks]
\label{rem:nonordinary-ppav-covector-checks}
\status{Proved under stated hypotheses}
The root certificate ran in $0.09968187000049511$ seconds and is retained
as \path{outputs/root_ppav_nonordinary_covector_checks.json}.
The independent monodromy and obstruction/crystal reruns took
$0.09809181400123634$ and $0.1044188160012709$ seconds respectively.
The program verifies all sixteen local operator identities, including
the ordinary comparison row, and twelve finite illustrations of the
monomial uniqueness used in
Lemma~\ref{lem:linear-covectors-any-nil-length}.
The lemma itself is proved for every exponent and every dimension;
the finite illustrations are not substituted for that proof.
\end{remark}


% Promoted after complete root and reduction/descent proof and source audits.
\subsection{Geometric generic nonordinary principally polarized surfaces}
\label{subsec:generic-nonordinary-ppav-groups-candidate}
\status{Proved and independently checked by the root and reduction/descent agents. Primary source statements and proofs were opened. The scope is geometric generic surfaces and finite-degree open consequences; fixed nongeneric nonordinary surfaces remain a separate question.}

\begin{lemma}[Frame component counts under nonproper specialization]
\label{lem:frame-components-specialization-candidate}
\status{Proved under stated hypotheses}
Let $B$ be an excellent normal integral scheme, let
$U\subset\mathbf P_B(\mathcal W)$ be an open whose geometric
fibres are nonempty, and let $\mathcal V$ be a rank-$r$
$\mathbf F_2$ local system on $U$. Write $G_{\bar b}$ and
$G_{\bar\eta}$ for its geometric images on a geometric fibre
and on the geometric generic fibre. Then
\[
                         |G_{\bar\eta}|\ge |G_{\bar b}|.
\]
If every geometric-fibre image is, up to conjugation, contained
in a fixed subgroup $P\subset\GL_r(\mathbf F_2)$ and one
geometric fibre has image $P$, the geometric generic image is
$P$. Under the same fixed-subgroup bound, if that equality
holds generically, it holds on every
geometric fibre over a nonempty open of $B$.
\end{lemma}

\begin{proof}
Use the full frame cover $Z\to U$, of degree
$|\GL_r(\mathbf F_2)|$. Its geometric fibre over $b$ has
$|\GL_r|/|G_{\bar b}|$ connected components, because the
image acts freely on frames. Choose a finite separable
extension of $k(B)$ defining all components of
$Z_{\bar\eta}$, and normalize $B$ in it. The normalization
$B'\to B$ is finite surjective by excellence. The pullback
$U'$ is a normal integral open of a projective bundle over
$B'$. The finite etale cover $Z'$ is normal, so closures of
its generic components are disjoint open-and-closed
components. Each surjects onto $U'$, since its image is both
open and closed. They all meet every geometric fibre of
$U'\to B'$. Hence the number of geometric generic components
is at most the number on any geometric special fibre,
giving the asserted inequality.

For the last assertion, shrink $B$ so that the normalization
is finite etale, and apply
\cite[Lemma~37.27.5]{Stacks0559} to each component of $Z'$.
Its generic fibre is geometrically irreducible; it is
geometrically normal because it is finite etale over an open
of projective space. The lemma gives an open of $B'$ on
which all these fibres are geometrically irreducible.
Remove from $B$ the finite image of the complementary
closed subset. The remaining open is nonempty and all its
geometric frame fibres have exactly the generic number of
components. Their images therefore have the generic order.
Together with the prescribed subgroup bound, this proves
equality up to conjugation. This argument uses full frames;
it requires no choice of invariant-covector basis over $B$.
\end{proof}

\begin{lemma}[Constant-rank relative unit systems in ppav families]
\label{lem:ppav-family-constant-rank-Lang-candidate}
\status{Proved under stated hypotheses}
Let $B$ be reduced and locally Noetherian, and let $A\to B$
be a family of principally polarized abelian
surfaces in characteristic two, supplied with a relatively
ample line bundle $M$ representing the polarization, and
suppose every geometric fibre has the same two-rank
$f\in\{0,1\}$. Fix $d\ge1$ and put $b=9d^2$.
On the open $U_d$ of smooth curves in $|M^{3d}|$ having
maximal relative stable rank, the relative Lang system is a
finite etale local system of rank
\[
 r=\begin{cases}b-1,&d\text{ odd},\\ b-3+f,&d\text{ even}.
 \end{cases}
\]
Every geometric fibre of $U_d\to B$ is a nonempty integral
open of a projective space. Its geometric image is contained
in $\GL_r(\mathbf F_2)$ in odd degree and in the pointwise
stabilizer of $f+1$ independent covectors in even degree.
\end{lemma}

\begin{proof}
Cohomology and base change for the ample line bundle gives
a rank-$b$ vector bundle of sections, and the universal
hypersurface sequence gives a vector bundle of relative
coherent cohomology of rank $b-1$. The rank formula and
nonemptiness of the smooth maximal-rank open on every
geometric fibre are Theorem~\ref{thm:DE-A3-ppav-ranks}.
All smaller stable ranks are detected by minors of a fixed
high Frobenius iterate: an iterate of exponent at least
$b-1$ has image equal to the stable image on each geometric
coherent fibre.

On $U_d$ the linearization of that iterate has constant rank
$r$, so its image is a vector subbundle, using the usual
invertible-minor presentations. Frobenius restricts to an
isomorphism of its Frobenius pullback with itself, by the
same fibrewise rank statement and Nakayama's lemma. Its
Lang fixed scheme is therefore finite etale and gives the
claimed local system. This constructs a single system over
the family even when the nilpotence index of the complementary
part changes. It avoids identifying the first image with
the stable image.

The last subgroup bounds are the entire invariant-covector
theorem, Theorem~\ref{thm:nonordinary-ppav-all-covectors},
applied to each geometric fibre. In even degree its
dimension is one for $f=0$ and two for $f=1$.
\end{proof}

\begin{lemma}[The needed genus-two strata and their specializations]
\label{lem:ppav-nonordinary-strata-candidate}
\status{Proved under stated hypotheses}
Over $\overline{\mathbf F}_2$, the two-rank-one stratum of
principally polarized abelian surfaces is smooth and
irreducible. It contains the mixed product locus.
Every geometric irreducible component of the two-rank-zero
stratum has a normal one-dimensional parameter scheme
dominating it and containing a point representing the
product principally polarized surface $E_0^2$.
\end{lemma}

\begin{proof}
The primary source is
\cite[Corollary~8.4, printed pp.~607--608; Theorem~11.5 and
Example~11.6, p.~620; Example~12.2, pp.~623--624]{EkedahlVanderGeer2009Strata}.
The four genus-two final strata are, in order, the ordinary
stratum, the two-rank-one stratum, the two-rank-zero
$a$-number-one stratum, and the superspecial stratum.
For the two middle strata used here---the two-rank-one and
the two-rank-zero $a$-number-one strata---the final flag
maps are finite etale of degree one. Indeed their canonical
filtrations are already complete: the canonical-rank-set
formula in
\cite[Example~3.4, p.~589]{EkedahlVanderGeer2009Strata}
is $\{0,f,f+1,\ldots,2g-f-1,2g-f,2g\}$, which contains
every rank from zero to four when $g=2$, $a=1$ and
$f\in\{0,1\}$. Proposition~4.5 and the proof of Lemma~4.6
of the same source, pp.~592--594, identify final flags
with refinements of this canonical conjugate flag, and
identify a conjugate flag with a unique Hodge flag over
an algebraically closed field. Thus the final flag is
unique. Directly, its required line in two-rank
one is the toral line, while in two-rank zero and
$a$-number one the rank-one nilpotent Hodge operator has
the same image and kernel line. Thus the finite
etale cover in Corollary~8.4 has one geometric point in its
fibre and degree one. We do not repeat the broader wording
of Example~12.2 that assigns degree one to all four strata:
its ordinary case contradicts the factor $1+p$ in the table
on p.~623 and the formula for $\gamma_g(w_\varnothing)$ on
p.~624; the superspecial flags likewise need not be unique.
Neither of those two incorrect degree-one cases is used.
Nor do we import the parenthetical semisimple-rank
expression $g-f$ in Example~3.4: the semisimple rank is
$f$, as also stated explicitly before Lemma~4.9 on p.~594.
The canonical-rank-set formula and the refinement argument
are the parts used here.
Corollary~8.4 gives smoothness of the open flag strata
and normality of their closures. Theorem~11.5 and
Example~11.6 give irreducibility of the two-rank-at-most-one
locus, hence of its dense two-rank-one open. These statements
hold in characteristic two: the large-characteristic
condition in Theorem~11.4(ii) concerns the converse
component-counting assertion, not Theorem~11.5.

For the two-rank-zero stratum, take a component of the
closed one-dimensional flag stratum $\overline U_{s_2}$.
It is normal by Corollary~8.4 and its dense open maps
degree one to the corresponding $a$-number-one component.
Each such component meets the zero-dimensional superspecial
flag stratum by
\cite[Proposition~6.1 and its proof, printed p.~601]{EkedahlVanderGeer2009Strata}.
No irreducibility of the full supersingular locus is asserted
or needed. The flag bundle over the abelian moduli space is
projective, and the closed stratum is closed in that bundle.
The forgetful map on a chosen normal curve component is
therefore proper. It is generically degree one onto its
image by the middle-stratum calculation above. A nonconstant map from an integral
curve has zero-dimensional fibres, so this proper map is
quasi-finite and hence finite. It is the normalization of
its image. In particular the superspecial flag point just
obtained is an actual point of this normal parameter curve;
we have not replaced the intersecting ppav component locus
itself by a smooth scheme without a map.

We verify that the underlying principally polarized surface
at such a point is $E_0^2$ in characteristic two. By
Lemma~\ref{lem:ppav-surface-dichotomy}, an indecomposable
principally polarized surface has a smooth genus-two theta
curve and is its Jacobian. A smooth genus-two curve cannot
have Cartier identically zero. Indeed in its function field
the kernel of Cartier on rational differentials is
one-dimensional over the square subfield. If both independent
canonical differentials were killed, their ratio would be a
square. The degree-two canonical map to $\mathbf P^1$
would then be purely inseparable. For a one-variable
function field $K$ over a perfect algebraically closed
field, $[K:K^2]=2$; a purely inseparable degree-two
canonical map would force $K^2=k(\mathbf P^1)$ and hence
genus zero, a contradiction.
Thus an $a$-number-two surface is decomposable. Both
elliptic factors have $a$-number one, hence are supersingular.
Over an algebraically closed field of characteristic two
there is a unique supersingular elliptic isomorphism class,
represented by $E_0:y^2+y=x^3$. This last assertion follows
also by normalizing a Weierstrass equation with $a_1=0$
and $a_3\ne0$ and making successive translations of $x,y$.
The polarization in the decomposable case is the product
polarization, proving the claimed specialization.

All these arguments may be made on a fine prime-to-two
level-five moduli scheme. For two-rank one, choose a
connected etale chart through the desired mixed product;
its smooth component maps openly, hence dominantly, to
the irreducible stratum. We do not assume that a level
cover is irreducible. For two-rank zero, pull back the
normal flag component to that finite etale level cover
and take a component. It still dominates the original
component and meets a point above the superspecial
specialization, since a finite etale component over a
normal connected base is surjective. A local etale
neighbourhood supplies a line bundle representing the
principal polarization: the relative Picard torsor is
smooth and the line bundle can be rigidified at zero.
No symmetric-theta torsor is required in characteristic two.
These changes preserve normality and generic dominance.
They do not change which special product theorem applies:
if $M'$ and $M$ represent the same principal polarization
over an algebraic closure, then $M'\otimes M^{-1}\in\Pic^0(A)$.
The isomorphism $\lambda_M:A\to\Pic^0(A)$ supplies a unique
$a\in A$ with $M'\simeq t_a^*M$. Translation by $a$
identifies every complete system $|(M')^{3d}|$ with
$|M^{3d}|$ and intertwines the curve cohomology, the ambient
restriction and Frobenius. Thus their relative local systems
and images agree. In particular a locally chosen Picard
section at the special point may replace the explicit
product theta line bundle without altering any seed group.
\end{proof}

\begin{theorem}[Generic full images on nonordinary ppav surfaces]
\label{thm:generic-nonordinary-ppav-groups-candidate}
\status{Proved under stated hypotheses}
Let $(A,M)$ be a geometric generic principally polarized
abelian surface on a two-rank-$f$ stratum, $f\in\{0,1\}$,
in characteristic two. For $f=0$ this means the geometric
generic point of any irreducible component.
For every $d\ge1$, the relative system on the maximal-rank
smooth open of the complete linear system $|M^{3d}|$ has
\[
 G_{\mathrm{geom},d}=
 \begin{cases}
 \GL_{9d^2-1}(\mathbf F_2),&d\text{ odd},\\[2pt]
 \operatorname{Hom}(\mathbf F_2^{f+1},\mathbf F_2^{9d^2-4})
       \rtimes\GL_{9d^2-4}(\mathbf F_2),&d\text{ even}.
 \end{cases}
\]
For any perfect field of definition with such geometric
generic moduli, the arithmetic image equals the geometric
image. For the embedding by $M^3$, the same assertion for
ambient degree-$d$ hypersurfaces holds whenever ambient
restriction is surjective, in particular for
$d\ge\max(1,\operatorname{reg}(\mathcal I_A)-1)$.
\end{theorem}

\begin{proof}
For $f=1$, base change the stratum to an algebraically
closed field containing a transcendental ordinary elliptic
modulus. Choose its mixed product point and a normal
etale chart through it as in the preceding lemma.
Every degree at that point has the asserted full group
by Theorem~\ref{thm:generic-mixed-product-all-groups-candidate}.
The relative ranks on the entire stratum are the same as
at this point. The family Lang lemma and the full frame
component-count lemma therefore give at least that group
order at the geometric generic point. The entire
invariant-covector theorem gives precisely that full
stabilizer as an upper bound. The two bounds coincide.
Dominance of the chart identifies this geometric generic
group with that of the original moduli stratum; an
algebraically closed extension does not change the image
of a finite etale local system.

For $f=0$, use the normal parameter component and its
superspecial product point supplied by the preceding
lemma. The full rank is constant at every geometric
point, including this specialization. The product groups
in every degree are Theorem~\ref{thm:ss-product-all-full-groups}.
The same frame-component inequality and invariant-covector
upper bound prove the assertion at the generic point of
each component. This use of a normalization is precisely
why singularities or multiple branches of a supersingular
coarse moduli locus cause no gap.

In both cases the entire geometric invariant space is
arithmetic over a perfect field, by
Theorem~\ref{thm:nonordinary-ppav-all-covectors}.
Its pointwise stabilizer is the displayed geometric image,
so arithmetic equality follows. The ambient restriction
claim uses only the stated surjectivity and the integral
fibres of a linear projection; it does not assume
characteristic-two projective normality of the cubic
theta embedding.
\end{proof}

\begin{corollary}[Open and finite-field consequences for finitely many degrees]
\label{cor:generic-nonordinary-ppav-finite-degrees-candidate}
\status{Proved under stated hypotheses}
Fix a finite set of positive degrees and one of the strata
or components in the preceding theorem. On a nonempty
Zariski open of that stratum or component, every geometric
surface has the displayed full groups in all the chosen
degrees. Consequently there are examples defined over
finite extensions of $\mathbf F_2$ with those groups in
all the chosen complete linear systems.
No claim that one finite-field point works for every
degree follows from this corollary.
\end{corollary}

\begin{proof}
For a single degree, apply the last assertion of the
frame-component lemma to the universal family on a
normal chart. After a finite separable extension defining
its geometric generic frame components, shrink to make
that cover finite etale and make every component's fibres
geometrically irreducible using Stacks~\cite{Stacks0559}.
The resulting full-image property descends to an open
of the original chart and then to an open of the moduli
component: an etale map is open, and the property depends
only on the geometric polarized surface. For a normalized
two-rank-zero curve, first restrict to the dense normal
open of its image, where the normalization is an
isomorphism. Intersect the finitely many resulting dense
opens for the prescribed degrees.

The moduli charts, universal families, finite covers and
their finitely many open conditions descend from
$\overline{\mathbf F}_2$ to a finite field. A nonempty
finite-type scheme over a finite field has a closed
point with finite residue field. Its polarized surface
gives the asserted examples, after a further finite
extension if needed to choose a representing line bundle.
Only finitely many opens are intersected; no countable
intersection or reversed specialization is used.
\end{proof}


\subsection{Some inseparable product polarizations: sharp rectangular ranks}
\label{sec:ss-rectangular-ranks}

\begin{lemma}[Cartier bounds from one or two even factors]
\label{lem:ss-rectangular-even-bound}
\status{Proved; independent full proof audit passed.}
Let $E:y^2+y=x^3$ over an algebraically closed field of characteristic
two, let $A=E\times E$, and put
$N=\OO_E(mO)\boxtimes\OO_E(nO)$, where $m,n\ge3$.
On $W=H^0(A,N)$ define the root-removed Cartier matrix $C_s$ by
the Cartier trace of $su\,dx_1\wedge dx_2$, for $s,u\in W$.
The generic stable rank on $H^2(A,N^{-1})$ is at most
$mn-2$ if $m$ is even, and at most $mn-3$ if both $m,n$ are even.
For every smooth divisor $Y\in|N|$, the stable rank on its relative
coherent cohomology equals the stable rank on $H^2(A,N^{-1})$.
\end{lemma}
\begin{proof}
The invariant derivation on $E$ is
$\partial x=1$, $\partial y=x^2$, and $\partial^2=0$.
If $m$ is even, the line bundle $\OO_E(mO)$ is pulled back by
relative Frobenius from $\OO_{E^{(2)}}((m/2)O)$.  Its canonical
connection therefore makes $\partial_1$ a global $k$-linear
endomorphism of $W$.  In horizontal local trivializations the
following identities are Cartier identities of rational differential
forms and hence identities of regular twisted sections:
\[
                  C_s(s)=0,\qquad C_s(\partial_1s)=0.
\]
For the first, $C(s^2\,dx_1\wedge dx_2)=sC(dx_1\wedge dx_2)=0$.
For the second, write $s=f$ locally.  Then
$f\partial_1f\,dx_1\wedge dx_2=f\,df\wedge dx_2$.
Both factors on the right are closed, so multiplicativity of Cartier
and $C(dx_2)=0$ give its vanishing.  The horizontal transition in
the first factor and the transition depending only on the second
coordinate make $\partial_1s$ a section of $N$; this is why parity
is required for this particular global relation.

Generically $s$ and $\partial_1s$ are linearly independent.
Indeed, if $\partial_1s=cs$ with $c\in k$, applying $\partial_1$
gives $c^2s=0$, hence $c=0$.  The operator is nonzero: the rational
section $x_1$ belongs to $W$ and has nonzero derivative, using the
horizontal rational trivialization with pole divisor $mO$.
If both factors are even, the same argument gives
$C_s(\partial_2s)=0$.  The three sections
$s,\partial_1s,\partial_2s$ are generically independent, as one sees
already from $s=x_1x_2$: they are $x_1x_2,x_2,x_1$.
The corresponding matrix-minor vanishing holds on a dense open and
therefore everywhere.  Linear rank bounds stable rank, proving both
inequalities.

The divisor sequence and ample vanishing give a Frobenius-compatible
exact sequence
\[
0\longrightarrow H^1(A,\OO_A)\longrightarrow H^1(Y,\OO_Y)
 \longrightarrow H^2(A,N^{-1})\longrightarrow H^2(A,\OO_A)
 \longrightarrow0.
\]
Frobenius on both ambient terms is zero, because $E$ is supersingular.
Passing to stable parts therefore identifies the stable part of the
full curve cohomology, of its relative quotient, and of
$H^2(A,N^{-1})$.  The latter Frobenius is Serre dual to the indicated
Cartier matrix, with the usual Frobenius twist.  These statements also
hold for reduced divisors used below: the same exact divisor sequence
and the same vanishing apply.
\end{proof}

\begin{lemma}[An exact rectangular base and its odd norm images]
\label{lem:ss-rectangular-base-norm}
\status{Proved; independent full proof audit and independent exact matrix reconstruction passed.}
For $N_0=\OO_E(6O)\boxtimes\OO_E(3O)$, a general smooth divisor
has genus $19$ and full stable rank $16$.
If $a,b$ are positive odd integers, a general smooth divisor of
$\OO_E(6aO)\boxtimes\OO_E(3bO)$ has full stable rank $18ab-2$.
\end{lemma}
\begin{proof}
We first supply a finite exact certificate for the base.  Use the bases
$(1,x,x^2,x^3,y,xy)$ and $(1,x,y)$, with tensor index $3i+j$.
Over $\mathbf F_2$ take the coefficient vector encoded by the integer
$74594$, its $j$-th coefficient being its $j$-th binary digit.
The eighteen columns of the Cartier matrix, encoded in the same way,
are
\[
\begin{gathered}
4106,12,16432,4177,40,16660,32848,96,131456,\\
33416,320,133280,12424,49472,82080,65664,164608,66560.
\end{gathered}
\]
Binary Gaussian elimination gives
$\operatorname{rank}C_s=\operatorname{rank}C_s^2=16$.
The exact code, both sampled inputs (including the rejected one),
complete matrix and output are
\path{code/supersingular_rectangular_rank.py} and
\path{outputs/supersingular_rectangular_rank.json}.
The formula used to build every column is
\[
 C(x^r dx)=\begin{cases}0&r\text{ even},\\x^{(r-1)/2}dx&r\text{ odd},\end{cases}
\quad
 C(x^r y\,dx)=\begin{cases}x^{r/2+1}dx&r\text{ even},\\
                       x^{(r-1)/2}y\,dx&r\text{ odd},\end{cases}
\]
with $y^2=y+x^3$ and the tensor product on the surface.
Thus this is a reproducible certificate, not an inference from
sample frequency.  Equality of the first two ranks makes the image
stable and proves stable rank $16$ at this coefficient vector.
The corresponding nonzero minors of iterated Frobenius define a
nonempty open.  The upper bound in
Lemma~\ref{lem:ss-rectangular-even-bound} is $16$, and $N_0$ is very
ample.  Bertini's nonempty smooth open intersects this rank open.
Adjunction gives genus $1+\chi(N_0)=19$.

For each positive odd integer $c$ there is an \'etale endomorphism of
$E$ of degree $c$.  Indeed Lemma~\ref{lem:DE-A3-ss-norm} constructs
an endomorphism of every positive degree by explicit quaternion norms
and the four-square identity.  An endomorphism of odd degree in
characteristic two is automatically \'etale.
Choose such endomorphisms of degrees $b,a$ on the first and second
factors, respectively, and let $\phi:A\to A$ be their product.
It is \'etale of degree $ab$.  Pushforward of a horizontal or vertical
divisor gives
\[
     \operatorname{Nm}_\phi(N_0)\equiv
                  \OO_E(6aO)\boxtimes\OO_E(3bO)=N_{a,b}.
\]

We prove that a general source curve has only nodes after applying
$\phi$, including kernels whose projections are not injective.
For a nonzero kernel element $g$, impose the two values
$s(P)=s(P+g)=0$.  They are independent because $N_0$ is very ample.
If $g_1\ne0$, degree six on the first elliptic factor separates the
two first jets, and degree three on the second separates a first
jet at one point.  Tensor products therefore give all six surface
first-jet coordinates at $P,P+g$.  Tangency imposes one further
nonzero determinant equation after the two value equations.
If $g_1=0$, then $g_2\ne0$.  Degree three on the second factor has
rank three on the length-four scheme $2P_2+2(P_2+g_2)$: its kernel
is zero because a nonzero degree-three section cannot vanish to
total order four.  Its two value functionals are independent.
The two derivatives in the first factor are independently supplied
by the two second-factor value functionals.  After imposing the
values, the two second-direction derivatives span a nonzero line;
otherwise a degree-three section would again vanish to order four.
The tangent determinant is consequently a nonzero polynomial in
these derivative parameters.  In either case, the incidence of
tangencies has dimension at most $\dim|N_0|-1$ after allowing the
two-dimensional point $P$ to vary.  There are finitely many $g$.

For two distinct nonzero kernel elements $g,h$, the three value
conditions at $P,P+g,P+h$ are independent unless all three first
coordinates coincide and their three second coordinates form a
degree-three divisor linearly equivalent to $3O$.
To verify this, separate the at most three distinct first coordinates
using the degree-six bundle, then separate the at most two second
coordinates within each repeated first-coordinate class using the
degree-three bundle.  Only the class of size three remains; on an
elliptic curve its value rank drops from three to two precisely when
the sum of its three points is zero.
That exceptional condition is $3P_2+g_2+h_2=0$, so its point locus
has dimension one, because $[3]$ is finite \'etale.  The ordinary
rank-three incidence over a dimension-two locus and the exceptional
rank-two incidence over a dimension-one locus both have dimension
at most $\dim|N_0|-1$.  Thus a general source member has no triple
identifications.  A higher identification contains a triple and is
also excluded.  The tangency calculation excludes a translated
component equal to the source, hence $\phi|_{C'}$ is birational onto
its image.  Its image is integral and nodal, because $\phi$ is \'etale
and all remaining pairs are transverse.

Choose $C'$ also in the smooth rank-$16$ open.  The image has
arithmetic genus $18ab+1$, normalization genus $19$, and exactly
$18(ab-1)$ nodes.  The normalization sequence contributes a
Frobenius-bijective vector space of that dimension.  Its full stable
rank is therefore $16+18(ab-1)=18ab-2$.
The two ample line bundles in the numerical equality above differ
by an element of $\Pic^0(A)$.  Surjectivity of the polarization
$\lambda_{N_{a,b}}:A\to A^\vee$ gives a target translation putting
the nodal image in the exact system $|N_{a,b}|$.  No separability of
this polarization, and no two-torsion translation, is needed for
this step.  Translation preserves every rank and singularity used.
Openness of iterated Frobenius rank and the nonempty smooth-member
open give the same lower bound for a general smooth divisor.
Lemma~\ref{lem:ss-rectangular-even-bound} gives the matching upper
bound, proving the result.
\end{proof}

\begin{theorem}[All degrees for some even product types]
\label{thm:ss-rectangular-all-degrees}
\status{Proved; independent full proof audit passed.}
Let $E:y^2+y=x^3$ over a perfect field of characteristic two, and
let $A=E\times E$, and let $a,b$ be positive integers.  Put
$M=\OO_E(aO)\boxtimes\OO_E(bO)$ and $L=M^3$.
The relative coherent dimension in $|L^d|$ is $9abd^2-1$.
If exactly one of $a,b$ is even and that even integer has
$2$-adic valuation one, then for every $d\ge1$ the geometric generic
relative stable rank is
\[
 r_d=\begin{cases}9abd^2-2&d\text{ odd},\\
                  9abd^2-3&d\text{ even}.
      \end{cases}
\]
If both $a,b$ are even, it is $r_d=9abd^2-3$ for every $d\ge1$.
For the embedding by $L$, these formulas also hold for ambient
hypersurfaces whenever
$d\ge\max\{1,\operatorname{reg}(\mathcal I_A)-1\}$.
\end{theorem}
\begin{proof}
Pass to an algebraic closure; geometric generic ranks are unchanged.
In the first case interchange the factors if necessary and write
$a=2a'$ with $a',b$ odd.  For odd $d$, apply
Lemma~\ref{lem:ss-rectangular-base-norm} with parameters $a'd,bd$.
It gives $18a'bd^2-2=9abd^2-2$.

Whenever $ad$ and $bd$ are both even, decompose the target divisor
class as
\[
 (3ad,3bd)=(3,3)+(3(ad-1),3(bd-1)).
\]
Both summands are cubes of ample product bundles with odd Euler
characteristic.  Theorem~\ref{thm:DE-A3-separable-ranks} supplies
smooth members in each system with full curve defect two.
Choose them simultaneously in their sharp opens and transverse
to each other; very ampleness and Bertini make these conditions
nonempty opens.  Their union is connected, since their intersection
number is positive.  If their genera are $g_1,g_2$ and they meet in
$q$ transverse points, normalization gives genus $g_1+g_2+q-1$
and stable rank $(g_1-2)+(g_2-2)+q-1$.  Thus the union has full
defect four and stable rank $9abd^2-3$.
The rank open again meets the smooth-member open in the target
linear system.  The both-even bound of
Lemma~\ref{lem:ss-rectangular-even-bound} is exactly $9abd^2-3$,
so equality holds.  This proves all even $d$ in the first case and
every $d$ in the second.

The relative dimension follows from surface Riemann--Roch and ample
vanishing.  Finally the indicated regularity bound makes ambient
restriction onto $H^0(A,L^d)$ surjective; the induced parameter map
has integral geometric fibres, preserving its geometric generic
rank.  The remaining product types with one even factor of higher
$2$-adic valuation are settled by
Theorem~\ref{thm:DE-A3-all-ss-rectangles} below. Arbitrary
inseparable polarizations on a nonproduct surface remain outside
these product-polarization results.
\end{proof}


\subsection{A uniform Cartier calculation for every even product type}
\label{sec:DE-A3-all-ss-rectangles}

\begin{lemma}[A rank-two Frobenius bundle and its endomorphism quotient]
\label{lem:DE-A3-ss-Frobenius-endomorphisms}
\status{Proved; root, reduction-descent, and monodromy full independent proof and certificate audits passed.}
Let $E/\mathbf F_2$ be $y^2+y=x^3$, let $O$ be its point at infinity,
and identify its Frobenius twist with $E$ by this model.  Write $F$ for
relative Frobenius and $\partial x=1$, $\partial y=x^2$ for the invariant
vector field.  For $q>0$ put
\[
 V=F_*\OO_E(2qO),\qquad H=\OO_E(qO).
\]
Then $V\simeq U\otimes H$, where $U=F_*\OO_E$ is the nonsplit extension
of $\OO_E$ by itself.  There is a perfect alternating pairing
$V\otimes V\to\det V\simeq\OO_E(2qO)$, given by Cartier trace with
respect to $dx$.  Moreover
\[
 h^0(E,\End V)=2,\qquad
 h^0\bigl(E,\End V/\OO_E\operatorname{id}\bigr)=2.
\]
The latter two-dimensional space is spanned by the classes of
multiplication by the source function $x$ and of $\partial$.
\end{lemma}
\begin{proof}
The Cartier formula $C(dx)=0$ makes Frobenius zero on $H^1(E,\OO_E)$
by Serre duality.  In
$0\to\OO_E\to F_*\OO_E\to B^1_E\to0$, the quotient is a line
bundle of degree zero with one nonzero global section.  It is therefore
$\OO_E$.  The extension is nonsplit because $h^0(F_*\OO_E)=1$.
Projection formula gives $V=U\otimes H$ and $\det V=H^2$.
We use Cartier in the convention of
\cite[0, Theorem~2.1.9, p.~515]{Illusie1979DRW}, already opened and
checked in the crystalline comparison above.
The pairing in the rational basis $(1,x)$ is the wedge pairing:
$C(dx)=C(x^2dx)=0$, while $C(xdx)=dx$.  It is thus alternating and
perfect, and the Cartier transformation law makes it global.

The dual of the displayed nonsplit extension is again nonsplit, so
$\dim\Hom(U,\OO_E)=1$.  Applying $\Hom(U,-)$ to the extension,
the identity maps to its nonzero quotient homomorphism.  Consequently
$\dim\End(U)=2$.  The identity and the nonzero nilpotent endomorphism
$\partial$ span this space; $\partial^2=0$.  Both have trace zero
in characteristic two.  The trace pairing makes $\End(U)$ self-dual.
Thus the map $H^1(\OO_E)\to H^1(\End(U))$ induced by the scalar
inclusion is dual to the zero map on global traces.  The cohomology
sequence of $0\to\OO_E\to\End(U)\to\End(U)/\OO_E\to0$ now gives
the second dimension assertion.

The derivative is a global endomorphism of $V$: the line bundle
$\OO_E(2qO)=F^*H$ has horizontal square transition functions.
Multiplication by $x$ gives a global class modulo scalar endomorphisms,
although multiplication by $x$ itself has a pole at $O$.  Indeed, for
the uniformizer $t=x/y$ at $O$ the equation of $E$ gives the exact identity
\[
                    x+t^{-2}=y/x^2.
\]
The right side is regular at $O$, and multiplication by $t^{-2}$ is a
scalar endomorphism over the Frobenius base because it is a square.
Away from $O$, multiplication by $x$ is regular.  In the rational basis
$(1,x)$, with $X$ denoting the coordinate on the Frobenius base, the two
classes have representatives
\[
 \begin{pmatrix}0&X\\1&0\end{pmatrix},\qquad
 \begin{pmatrix}0&1\\0&0\end{pmatrix}.
\]
They are linearly independent modulo scalars.  The preceding dimension
calculation shows that they span all global classes.  This argument
uses no classification of vector bundles on elliptic curves.
\end{proof}

\begin{lemma}[A uniform sparse section]
\label{lem:DE-A3-ss-uniform-sparse}
\status{Proved; root, reduction-descent, and monodromy full independent proof and certificate audits passed.}
For every positive integer $q$ divisible by six, put
$N_q=\OO_E(2qO)\boxtimes\OO_E(3O)$ on $E\times E$.
In coordinates $(x,y)$ and $(z,w)$ take the section
\[
\begin{split}
 s_q&=A+Bz+Cw,\\
 A&=x+P\{x^2+(1+x)y\},\\
 B&=x+Py^2,\\
 C&=1+Py,\\
 P&=x^{q-3}.
\end{split}
\]
All coefficients lie in $\mathbf F_2$.  The root-removed Cartier
endomorphism $C_{s_q}$ of $H^0(E\times E,N_q)$ satisfies
\[
             \rank C_{s_q}=\rank C_{s_q}^2=6q-2.
\]
In particular a general smooth member of $|N_q|$ has full and relative
stable rank $6q-2$.
\end{lemma}
\begin{proof}
The pole orders of $A,B,C$ at $O$ are respectively $2q-1$, $2q$, and
$2q-3$, so the displayed functions are sections as asserted.  Set
$u=x^{q/2-2}$; then $P=xu^2$.  Reduction by $y^2+y=x^3$ gives
\begin{equation}
 (A,B,C)=a^2+xb^2,
 \quad
 \begin{cases}
 a=(xu(x+y),\ 0,\ 1+x^2u),\\
 b=(1+u(x+y+x^2),\ 1+uy,\ uy).
 \end{cases}                                                   \label{eq:DE-A3-ss-pbasis}
\end{equation}
Here the squares are componentwise.  For polynomials over $\mathbf F_2$,
$a,b$ are precisely the two coordinate rows of the three sections of
$V=F_*\OO_E(2qO)$ in the rational basis $(1,x)$, after the indicated
identification of the Frobenius base with $E$.

Let $f_0=A,f_1=B,f_2=C$ denote those three sections of $V$.  On the
Frobenius base consider
\[
\begin{aligned}
 \Psi:V^3&\longrightarrow(\det V)^3,\\
 (v_0,v_1,v_2)&\longmapsto
 \begin{pmatrix}
 f_0\wedge v_1+f_1\wedge v_0\\
 f_0\wedge v_2+f_2\wedge v_0\\
 f_1\wedge v_2+f_2\wedge v_1
 \end{pmatrix}.
\end{aligned}
\]
The second-factor Cartier products in the basis $(1,z,w)$ are
$C(z\,dz)=dz$, $C(w\,dz)=z\,dz$, and $C(zw\,dz)=w\,dz$;
the three square products have Cartier zero.  Consequently $H^0(\Psi)$
is exactly $C_{s_q}$, in the tensor basis with indices $3i+j$.
This identifies the sheaf calculation with the required operator,
including its ordering.

We determine the rank-drop scheme of the three sections of $V$.
Put $\beta=1+uy$.  If $\beta\ne0$, the two rows in
\eqref{eq:DE-A3-ss-pbasis} have rank one precisely when
\[
                 xu(x+y)=0,\qquad 1+x^2u=0.
\]
The second equation forces $x\ne0$.  The first then gives $y=x$;
the elliptic equation gives $x^2+x+1=0$.  Conversely, because
$q/2$ is divisible by three, both points
\[
          Q_\zeta=(\zeta,\zeta),\qquad
          Q_{\zeta^2}=(\zeta^2,\zeta^2),\qquad \zeta^2+\zeta+1=0,
\]
satisfy these equations and have $\beta\ne0$.
If $\beta=0$, then $uy=1$, and the remaining minor is
\[
 a_0b_2+a_2b_0=1+u.
\]
Indeed $1+u=x^3u^2$ follows from the elliptic equation and $uy=1$.
Vanishing of this minor would force $u=1$ and $x=0$, which contradicts
$u=x^{q/2-2}$ with $q\ge6$.  Thus there are no other affine rank drops.
At infinity the leading poles of $A$ and $B$ have different parities,
$2q-1$ and $2q$, and their leading coefficients are one.  They give
an invertible two-by-two matrix in the local Frobenius basis
$(t^{-2q},t^{-2q+1})$.  There is no rank drop there.
At either $Q_\lambda$, $\lambda\in\{\zeta,\zeta^2\}$, the factor
$x+y$ has a simple zero: $\partial(x+y)=1+x^2=\lambda\ne0$.
It follows that the rank-drop ideal is the maximal ideal at each of
these two points, with length one.

Write $K=\ker\Psi$.  Away from these two points, the map
\[
             \End^0(V)\longrightarrow K,\qquad
                    T\longmapsto(Tf_0,Tf_1,Tf_2)
\]
is an isomorphism; here $\End^0$ means trace zero.  At a rank-drop
point choose a uniformizer $t$ and bases so that
$(f_0,f_1,f_2)=(t e_1,e_2,0)$.
Writing $v_0=(a,b)$, $v_1=(c,d)$, $v_2=(e,f)$ gives
\[
                    \Psi(v_0,v_1,v_2)=(a+td,tf,e).
\]
Its kernel consists of $((td,b),(c,d),0)$.
The trace-zero endomorphisms give the submodule in which $b$ is
divisible by $t$.  Thus $K/\End^0(V)$ has length one at each point,
and $\operatorname{coker}\Psi$ also has length one at each point.
In particular $\deg K=2$.

The perfect trace pairing identifies
$(\End^0 V)^\vee=\End(V)/\OO_E\operatorname{id}$.
The local calculation shows that a class in this dual lies in $K^\vee$
exactly when it annihilates the nilpotent endomorphism mapping the
missing direction $e_1$ to the image line $k e_2$.  Equivalently its
representative preserves that image line.  At $Q_\lambda$ the coordinate
row $a$ vanishes, so the image line in the rational basis $(1,x)$ is
$k x$.  By Lemma~\ref{lem:DE-A3-ss-Frobenius-endomorphisms}, a global
class is $\alpha[\operatorname{mult}_x]+\beta[\partial]$, with
$\alpha,\beta\in k$.  The two matrices displayed there show that
preserving this line imposes
\[
                         \alpha\lambda+\beta=0.
\]
The conditions for $\lambda=\zeta$ and $\zeta^2$ are independent.
Hence $H^0(E,K^\vee)=0$.  Serre duality and Riemann--Roch give
$h^1(K)=0$ and $h^0(K)=\deg K=2$.
The two bundles $V$ and $\det V$ have no $H^1$, so the exact sequences
for $\Psi$ identify the cokernel of $H^0(\Psi)$ with the two-dimensional
space of values of its torsion cokernel.  Its rank is therefore $6q-2$.

It remains to check stable rank, rather than merely first rank.
At $Q_\lambda$ the row $b$ is $(0,\lambda,\lambda^2)$, so the local
cokernel functional on an output triple $(U,V,W)$ is
\[
              \ell_\lambda(U,V,W)
                   =\lambda^2 U(Q_\lambda)+\lambda V(Q_\lambda).
\]
The global kernel contains $s_q$ and $\partial_1s_q$, by the Cartier
identities of Lemma~\ref{lem:ss-rectangular-even-bound}.
Their triples of values at $Q_\lambda$ are
\[
              (A,B,C)=(0,1,\lambda^2),\qquad
              (\partial A,\partial B,\partial C)
                        =(0,\lambda^2,\lambda).
\]
Thus the map from the kernel to the cokernel, in these bases, is
\[
 \begin{pmatrix}
    \ell_\zeta(s_q)&\ell_\zeta(\partial_1s_q)\\
    \ell_{\zeta^2}(s_q)&\ell_{\zeta^2}(\partial_1s_q)
 \end{pmatrix}
       =\begin{pmatrix}\zeta&1\\\zeta^2&1\end{pmatrix},
 \qquad \det=1.
\]
In particular the two kernel vectors are independent and their span
meets the image of $C_{s_q}$ trivially.  Since all matrix coefficients
are in $\mathbf F_2$, the Frobenius twists in the second iterate leave
this matrix unchanged.  Therefore its image is stable and
$\rank C_{s_q}^2=\rank C_{s_q}=6q-2$.
The nonzero iterated-rank minors give a nonempty rank open in $|N_q|$.
The first-degree and second-degree line bundles are very ample, so
Bertini gives a nonempty smooth open.  The two opens intersect.
Lemma~\ref{lem:ss-rectangular-even-bound} supplies the matching upper
bound and the identification with relative stable rank.
\end{proof}

\begin{theorem}[All rectangular product polarizations, without a valuation restriction]
\label{thm:DE-A3-all-ss-rectangles}
\status{Proved; root, reduction-descent, and monodromy full independent proof and certificate audits passed.}
Let $E:y^2+y=x^3$ over a perfect field of characteristic two, let
$a,b$ be arbitrary positive integers, and put
\[
 A=E\times E,\qquad M=\OO_E(aO)\boxtimes\OO_E(bO),\qquad L=M^3.
\]
Define $\epsilon(c)=1$ when $c$ is even and $\epsilon(c)=0$ when it
is odd.  For every $d\ge1$ the geometric generic relative stable rank
of a smooth member of $|L^d|$ is
\[
                  r_d=9abd^2-1-\epsilon(ad)-\epsilon(bd).
\]
Its relative coherent dimension is $9abd^2-1$ and its genus is
$9abd^2+1$.  For the embedding by $L$, the same formula applies to
ambient degree-$d$ hypersurfaces whenever
$d\ge\max\{1,\operatorname{reg}(\mathcal I_A)-1\}$.
\end{theorem}
\begin{proof}
Work over an algebraic closure.  If both $ad,bd$ are odd, the separable
polarization theorem, Theorem~\ref{thm:DE-A3-separable-ranks}, gives
$9abd^2-1$.  If both are even, the already proved transverse union
argument in Theorem~\ref{thm:ss-rectangular-all-degrees} gives
$9abd^2-3$.  Thus suppose exactly one is even, say $ad$.
If $v_2(ad)=1$, Lemma~\ref{lem:ss-rectangular-base-norm} gives
$9abd^2-2$.
If $4\mid ad$, put $q=3ad/2$.  It is divisible by six.
Lemma~\ref{lem:DE-A3-ss-uniform-sparse} gives a sharp smooth curve
in $|\OO_E(3adO)\boxtimes\OO_E(3O)|$.
Choose an odd-degree \'etale endomorphism of the first factor of degree
$bd$, and the identity on the second.  Its norm has numerical class
$(3ad,3bd)$.  The complete pair-jet and triple-value incidence proof
of Lemma~\ref{lem:ss-rectangular-base-norm} applies unchanged: the first
factor now has degree at least twelve, and the second still has degree
three.  It uses only separation in degrees at least six and three.
A general sharp source curve therefore has integral nodal image.
Its genus increases from $9ad+1$ to $9abd^2+1$ by exactly
$9ad(bd-1)$ nodes, so its stable rank increases by the same number,
from $9ad-2$ to $9abd^2-2$.
A target translation puts the image in the exact system, because the
polarization of the target ample bundle is surjective on geometric
points.  The smooth rank open then gives the stated lower bound.
The upper bound is Lemma~\ref{lem:ss-rectangular-even-bound}.
This proves every parity case and every valuation.
The dimension, genus, and ambient restriction assertions follow from
the same Riemann--Roch, adjunction, and regularity arguments already
proved in Theorem~\ref{thm:ss-rectangular-all-degrees}.
\end{proof}

\begin{remark}[Scope and retained verification]
\label{rem:DE-A3-all-ss-rectangles-checks}
\status{Proved; root, reduction-descent, and monodromy full independent proof and certificate audits passed.}
The uniform sheaf argument, including the two local modifications and
the invertible kernel--cokernel pairing, is the proof of the infinite
family.  The preceding finite sparse-template searches are retained
as discovery records, including every rejected template and every
nonsharp degree.  In particular template $62556$ fails at many
integers not divisible by six; no extrapolation from its successful
doubling sequence is used.  The exact verification program
\path{agents/obstructions_crystals/code/ss_uniform_rectangular_checks.py}
and its adjacent output check the polynomial identities, the two
geometric points, and eight independently reconstructed finite matrices.
The author run passed in $0.1617353309993632$ seconds; the root reviewer
read the entire program and reran it, preserving the author output,
in $0.154388964001555$ seconds.  A second independent full proof and
program audit passed, with an isolated rerun in
$0.1660513059978257$ seconds.  This result closes the product-polarization valuation
restriction, while making no assertion for arbitrary inseparable
polarizations on a nonproduct supersingular surface or for arbitrary
smooth projective varieties.
\end{remark}

